\documentclass[10pt,a4paper,twoside,bg=print,twocolumn,openany,nodeprecatedcode]{dndbook}
\usepackage[utf8]{inputenc}
\usepackage{tabulary}
\usepackage[normalem]{ulem}
\usepackage{color,soul}
\usepackage{float}
\usepackage{caption}
\usepackage{wrapfig}
\usepackage{tocloft}
\usepackage[toc]{multitoc}
\usepackage[hidelinks]{hyperref}
\raggedbottom
\extrafloats{2000}
\cftsetindents{section}{0em}{0em}
\cftsetindents{subsection}{0em}{0em}
\cftsetindents{subsubsection}{0.2em}{0em}
\setcounter{tocdepth}{1}
\renewcommand*{\multicolumntoc}{3}
\setlist[description]{nolistsep,listparindent=3pt,topsep=6pt}
\date{}
\begin{document}
\frontmatter
\tableofcontents
\mainmatter
\chapter{The Basics}
Dungeons \& Dragons is a game in which you and your friends take on roles and tell a shared story. While the \textit{Player's Handbook} teaches you how to play the game and how to create characters who are the heroes of the story, the \textit{Dungeon Master's Guide} is written for the player who presides over the game and makes sure everyone is having fun. This player is the Dungeon Master, or DM. Being a Dungeon Master is a fun, empowering, and rewarding experience, and this chapter walks you through the basics.

\section{What Does a DM Do?}
The DM gets to play many fun roles:


\begin{description}
\item[Actor.]
The DM plays the monsters, choosing their actions and rolling dice for their attacks. The DM also plays all the people the characters meet.
\item[Director.]
Like the director of a movie, the DM decides (and describes) what the players' characters encounter in the course of an adventure. The DM is also responsible for the pace of a play session and for creating situations that facilitate fun.
\item[Improviser.]
A big part of being the DM is deciding how to apply the rules as you go and imagining the consequences of the characters' actions in a way that will make the game fun for everyone.
\item[Referee.]
When it's not clear what ought to happen next, the DM decides how to apply the rules.
\item[Storyteller.]
The DM crafts adventures, setting situations in front of the characters that entice them to explore and interact with the game world.
\item[Teacher.]
It's often the DM's job to teach new players how to play the game.
\item[Worldbuilder.]
The DM creates the world where the game's adventures take place. Even if you're using a published setting, you get to make it yours.
\end{description}

\subsection{DM Tips}
The most important part of being a good DM is facilitating the fun of everyone at the table. Keep these tips in mind to help things go smoothly.


\begin{description}
\item[Embrace the Shared Story.]
D\&D is about telling a story as a group, so let the other players contribute through the words and deeds of their characters. Encourage players to engage by asking them what their characters are doing.
\item[It's Not a Competition.]
The DM isn't competing against the other players. It's your job to provide fun challenges and keep the story moving.
\item[Be Fair and Flexible.]
Treat your players in a fair, impartial manner. The rules help you do this, but when you need to act as referee, try to make decisions that ensure everyone is having fun.
\item[Communicate with Your Players.]
Open communication is essential to a successful D\&D game. Many problems can be solved or even prevented with honest conversation. Ask questions and solicit feedback after or between sessions.
\item[It's OK to Make Mistakes.]
If you overlook or misrepresent something, correct yourself and move on. No one expects you to memorize every rule or detail. Even if you don't realize your mistake until after a game session is over, it's OK to acknowledge the mistake at the start of the next session and make adjustments moving forward.
\end{description}

\begin{DndSidebar}[float=!b]{What's New in the 2024 Version?}
This is the 2024 version of the fifth edition \textit{Dungeon Master's Guide}. Much of the book has been reorganized, expanded, and rewritten from the 2014 version, and the versions of things in this book replace versions from older books. Here are a few highlights:




\begin{description}
\item[Sound Advice.]
Every chapter (but especially chapters 1, \textit{2}, \textit{4}, and \textit{5}) has new advice for Dungeon Masters of all experience levels.
\item[Tracking Sheets.]
Helpful sheets throughout the book give you tools to plan your game and keep track of your campaign. These sheets are also available for download in \textit{appendix C}.
\item[Encounter-Building Assistance.]
The rules for estimating the difficulty of combat encounters have changed, as you'll see in \textit{chapter 4}.
\item[Ready-Made Elements.]
Sample adventures in \textit{chapter 4}, a \textit{campaign setting} in \textit{chapter 5}, and new maps in \textit{appendix B} make it easier to run a game right away.
\item[Expanded and Revised Magic Items.]
\textit{Chapter 7} is packed with new magic items and old ones that have been revised.
\item[Bastions.]
\textit{Chapter 8} has rules that allow player characters to build, maintain, and enjoy their own strongholds.
\item[Lore Glossary.]
In \textit{appendix A}, a helpful glossary explains many of the iconic people and locations found throughout the D\&D multiverse.
\end{description}



\end{DndSidebar}

\section{Things You Need}
What you need to play hasn't changed much since the game's first publication in 1974.

\subsection{Rulebooks}
As the Dungeon Master, you need this book plus the \textit{Player's Handbook} (which contains most of the rules of the game) and the \textit{Monster Manual}. Your players need access to the \textit{Player's Handbook}, too, but they can share as needed.

Let players know beforehand what books (other than the \textit{Player's Handbook}) they can reference during a playing session. For example, it's not appropriate for players to look up a monster in the \textit{Monster Manual} (or the equivalent digital tool) while fighting that monster. If you're running a published adventure, players should avoid reading that adventure so they don't spoil any surprises.

\subsection{A Dungeon Master}
One player has the special role of Dungeon Master.

Some people love being the DM all the time, while others can end up feeling trapped as the "forever DM" for their gaming group. The "\textit{Group Size}" section in \textit{chapter 2} discusses possibilities for sharing the role of Dungeon Master among multiple players in a group.

\subsection{Players}
Players who aren't the Dungeon Master take on the roles of the heroes, also known as the characters or the adventurers.

D\&D plays best with four to six players in addition to the DM, but it's possible to run a game with fewer or more adventurers. See the "\textit{Group Size}" section in \textit{chapter 2} for advice on doing so.

\subsubsection{Finding Players}
Where do you find players? Here are a handful of suggestions:


\begin{itemize}
\item Game or hobby stores (the Store Locator on the Wizards of the Coast website can help you find stores near you that host D\&D events)
\item Friends, family, community members, and work colleagues who enjoy gaming or fantasy
\item Gaming clubs at schools
\item Social media and online messaging sites
\item Gaming conventions
\end{itemize}

\subsection{A Place to Play}
The bare minimum of space you need to play D\&D is room for everyone in your group to gather and participate.

When choosing the space you'll be playing in, enlist your players' help. Think about any accessibility needs you or they might have. Some players might have difficulty with low light, background music, strong odors, cramped spaces, or specific allergens. Accommodate what you can; communicate what you can't as early as possible.

If possible, play in an area with minimal visual or auditory distractions. Favor surroundings that reinforce your desired atmosphere and have little non-player traffic. If space is shared, reserve the space in advance.

You can also play D\&D anywhere you might come together in an online space, from a group video call to a sophisticated virtual tabletop.

\begin{DndSidebar}[float=!b]{Scheduling Games}
Sometimes the hardest thing about running a game is finding a time when everyone can play. Some groups play for a few hours every week, while others set aside a whole day once a month. Create a schedule that works best for your group.



For new groups, it often helps to schedule a single-session game (often called a "one-shot") as a way for people to try it out. If everyone has a great time at that one session, it can be easier to get them to make a long-term commitment.



Scheduling conflicts are sometimes inescapable. The "\textit{Group Size}" section in \textit{chapter 2} offers some advice on what to do when a player has to miss a session.



\end{DndSidebar}

\subsection{Dice}
You need a full set of polyhedral dice: d4, d6, d8, d10, d12, and d20. It's helpful to have at least two of each kind. Ideally, each player should also have their own set of polyhedral dice.

Lots of digital dice rollers exist. Simple, browser-based dice rollers are easily found on the internet. Specialized dice apps can be found in app stores, and virtual tabletops typically have dice-rolling functionality built in.

\subsection{Note-Taking Materials}
Everyone needs some way to take notes. During every round of combat, someone needs to keep track of Initiative, Hit Points, conditions, and other information. Players often like to take notes about what happens in the adventure, and at least one of them should record any clues and treasure the characters collect.

\subsubsection{Character Sheets}
Players need some way to record important information about their characters. Plain paper works fine, but players might find official or fan-made character sheets more helpful in organizing the information. A variety of digital character sheets are also available if you're playing online or using digital devices at the table.

\subsubsection{Campaign Journal}
Throughout this book you'll find tracking sheets you can use to make your work as a DM easier. They range from sheets you can use to track NPCs or settlements in your game to trackers you can use to make sure you're giving the adventurers a good number of magic items. These tracking sheets can form the basis of a campaign journal (see \textit{chapter 5}), and they'll help you plan your adventures and build your world. You can scan or photocopy these sheets for your personal use, and you'll find downloadable versions in \textit{appendix C}.

\subsection{Useful Additions}
Various resources can enrich your game and make it more fun. Many of these resources might have digital versions, making computers, tablets, and smartphones essential elements in some D\&D games and for some players.

\subsubsection{DM Screen}
A DM screen shields your books, notes, and die rolls from your players. (See the "\textit{Ensuring Fun for All}" section later in this chapter for more about when and why you might want to hide die rolls.) Most DM screens have art on the outward-facing panels and handy rules information on the inside-facing panels. Others might be made of fancy wood or sculpted to help set the mood for your game.

You don't need a physical screen to hide things if you're playing online, but it can be helpful to have ready access to important information like condition definitions, common actions, and other key rules. Some DMs set up a physical DM screen near their computer screen. A virtual tabletop might have reference information like this built in.

\subsubsection{Adventures and Sourcebooks}
Beyond the three core rulebooks, a plethora of additional content is available from Wizards of the Coast and other publishers. Adventures provide hooks, plots, maps, and encounters you can use in your game. Sourcebooks include things like new character options, new monsters, and inspiration for building your own adventures and campaigns. You can play D\&D without any of these additional products, but many DMs (and players) find them to be exciting additions to the game.

\subsubsection{Battle Grid and Miniatures}
Some DMs use a battle grid and miniatures to run combat encounters, which helps players visualize scenes when playing in person. A vinyl wet-erase mat with a printed grid, a gridded whiteboard, a cutting mat, a large sheet of gridded paper, or a printed poster map—any of these can serve as a battle grid. The grid should be marked in 1-inch squares.

You also need plastic or metal miniatures to represent characters and monsters in the game, but you can use coins, extra dice, paper counters, or even pieces of candy if miniatures are unavailable.

Many software tools designed to facilitate online D\&D play provide a battle grid. Even without such tools, though, many online D\&D games use screen sharing in combination with drawing programs, shared whiteboards, or similar tools as simple battle grids. Some DMs are comfortable with software that allows them to control lighting and show the players exactly what they can see; others find that complex software gets in the way of the game. Use whatever works for you and your group.

\subsubsection{Card Accessories}
Some players and DMs find it helpful to have information available in the form of cards. You can buy (or make) cards with individual spells, magic items, monster stat blocks, rules reference, and similar information for easy reference.

\section{Preparing a Session}
The more you prepare before your game, the more smoothly the game will go—to a certain point. To avoid being either under- or overprepared, use the one-hour guideline below and prioritize what to prepare depending on the time you have available.

\subsection{The One-Hour Guideline}
A D\&D game session usually starts with some out-of-game chatter as everyone settles down to play. Once the session gets underway, most groups can accomplish at least three things during one hour of play, where each "thing" might be any of the following:


\begin{itemize}
\item Explore a location such as a chamber in a castle or a cave
\item Converse with an intelligent creature
\item Reach consensus on a divisive issue
\item Solve a tricky riddle or puzzle
\item Survive a deadly trap
\item Fight a low-difficulty combat encounter
\end{itemize}

A more difficult combat encounter might count as two or three things, and a tense negotiation can use most or all of an hour of play on its own.

\subsection{Preparation Time}
The following guidelines can help you prepare for a session of play using a published adventure.

\subsubsection{One-Hour Preparation}
If you spend one hour each week preparing for your game, follow these steps:


\begin{description}
\item[Step 1.]
Focus on the story of the adventure. Read or reread the adventure's introduction and background information. Create a bulleted list of key plot points to make sure a coherent story unfolds.
\item[Step 2.]
Identify the encounters you want to run, then figure out how likely it is each encounter will get played, categorizing each one as "definite," "possible," or "unlikely."
\item[Step 3.]
Gather any maps you'll need for the definite and possible encounters, then focus the remainder of your prep time on the definite encounters, as outlined below.
\end{description}

For combat encounters, review the monsters' tactics and stat blocks. Note any special rules that apply to the setting of the encounter.

For social interaction encounters, make notes about the nonplayer characters (NPCs) in the encounter—their personalities, goals, and tactics.

For exploration encounters, record any clues or other information the characters should learn, and review any special rules that might come into play in the encounter.


\begin{description}
\item[Step 4.]
Consider how each definite encounter relates to the players' motivations (see the "\textit{Know Your Players}" section in \textit{chapter 2}). Think about elements you can add to interest them. For example, a combat encounter could open with a tense negotiation designed to appeal to players who enjoy social interaction.
\item[Step 5.]
Skim the encounters you flagged as possible.
\end{description}

\subsubsection{Two-Hour Preparation}
With another hour to prepare, add these steps:


\begin{description}
\item[Step 6.]
Carefully review each "possible" encounter.
\item[Step 7.]
Devote any time you have left to creating improvisational aids (see the "\textit{Improvising Answers}" section in \textit{chapter 2}).
\end{description}

\subsubsection{Three-Hour Preparation}
If you have three hours to prepare, add these steps:


\begin{description}
\item[Step 8.]
Skim each "unlikely" encounter.
\item[Step 9.]
Create a new encounter designed to appeal specifically to one player, or alter an existing encounter to relate to the goals and motivations of that player's character. Over the course of several sessions, do this for all your players and their characters.
\end{description}

\section{How to Run a Session}
This section explains how to run a game session; later on, \textit{chapters 4} and \textit{5} detail how to combine sessions into adventures and adventures into campaigns.

\subsection{Recap}
Start each game session after the first with a recap of what happened in the previous session. A recap helps players get back into the story. It also provides important information to players who missed the previous session. You can provide this recap, or you can invite one or more players to deliver the recap instead. Each approach has benefits:


\begin{description}
\item[DM Recap.]
Provide the recap yourself if you have specific information you need to impart or if you want the recap to be concise and focused on what's relevant.
\item[Player Recap.]
Let the players provide the recap if you want to gauge what they think is important or learn more about what they're getting out of the game. If the players miss any important details in their recap, you can interject a reminder.
\end{description}

\subsection{Encounters}
The bulk of a typical D\&D session consists of a series of encounters, similar to how a movie is a series of scenes. In each encounter, there are chances for the DM to describe creatures and places and for characters to make choices. Encounters can involve exploration (interacting with the environment, including puzzles), social interaction with creatures, or combat. The \textit{Player's Handbook} outlines the general rhythm of play in an encounter. The following sections offer more detailed information on how an encounter typically unfolds, in three steps.

\subsubsection{Step 1: Describe the Situation}
As the DM, you decide how much to tell the players and when. All the information the players need to make choices comes from you. Within the rules of the game and the limits of the characters' knowledge and senses, tell players everything they need to know.

\begin{DndSidebar}[float=!b]{}
Published adventures often include text in a box like this, which is meant to be read aloud to the players when their characters first arrive at a location or under a specific circumstance, as described in the text. It usually describes locations so the players know what's happening and have a sense of what their characters' options are.



\end{DndSidebar}

Whether you're running a published adventure or one of your own creation, your initial description of a room or situation should focus on what the characters can perceive. You don't have to reveal every detail at once. Most players begin to lose focus after about three sentences of descriptive text. As characters search rooms, open drawers and chests, and examine things more closely, give players more details about what their characters find.

The "\textit{Narration}" section in \textit{chapter 2} offers more extensive advice and examples of narration.

\subsubsection{Step 2: Let the Players Talk}
Once you're done describing the situation, ask the players what their characters want to do. Note what the players say, and identify how to resolve their actions. Ask them for more information if you need it.

Sometimes the players might give you a group answer: "We go through the door." Other times, individual players might want to do specific things—one might search a chest while another examines a bookshelf. Outside combat, the characters don't need to take turns, but you need to give each player a chance to tell you what their character is doing so you can decide how to resolve everyone's actions. In combat, everyone takes turns in Initiative order.

\subsubsection{Step 3: Describe What Happens}
After the players describe their characters' actions, it's the DM's job to resolve those actions, guided by the rules and the adventure you've prepared. So how do you decide? Think through these possibilities:


\begin{description}
\item[No Rules Required.]
Sometimes, resolving a situation is easy. If an adventurer wants to cross an empty room and open a door, you can just say that the door opens and describe what lies beyond (perhaps referencing your map or notes).
\item[Obstacles to Success.]
A lock, a guard, or some other obstacle might hinder a character's ability to complete a task. In those cases, you typically call for a D20 Test, usually an ability check. For example, a successful Dexterity (Sleight of Hand) check might be needed to pick the lock, while a successful Charisma (Persuasion) check and some coins might be needed to bribe the guard. The "\textit{Resolving Outcomes}" section in \textit{chapter 2} gives more guidance on how to use D20 Tests and other tools to determine the results of characters' actions.
\item[Roleplaying.]
When the players interact with other creatures, roleplay those creatures based on whether they are Friendly [Attitude], Indifferent [Attitude], or Hostile [Attitude]. Improvise based on what you know about the creatures, their knowledge, and their motivations. Then bring these creatures to life as you describe what happens. (See the "\textit{Running Social Interaction}" section in \textit{chapter 2} for more advice.)
\item[One Action at a Time.]
The \textit{rules about actions} in the \textit{Player's Handbook} limit how many things a character can do at once. Keeping those rules in mind can help you adjudicate situations.
\item[Combat.]
In combat, many situations involve attack rolls or saving throws. The rules of combat can help you determine the effectiveness of a character's actions. The "\textit{Running Combat}" section in \textit{chapter 2} offers advice on combat.
\item[Spellcasting.]
If a character casts a spell, you can usually let the player tell you what the spell does and how to resolve it. If questions arise, read the text of the spell yourself—how a spell is supposed to work is usually pretty clear. The general \textit{rules of spellcasting} in the \textit{Player's Handbook} are also essential for resolving a spell's effects.
\item[Exceptions Supersede General Rules.]
General rules govern each part of the game, but the game also includes class features, spells, magic items, monster abilities, and other elements that can contradict a general rule. When an exception and a general rule disagree, the exception wins. For example, it's a general rule that melee weapon attacks use the attacking character's Strength modifier. But if a feature says that a character can make melee weapon attacks using Charisma, that exception supersedes the general rule.
\end{description}

When narrating results, try to give a flavorful description while clearly communicating what's happening in the language of the game. See the "\textit{Narration in Combat}" section in \textit{chapter 2} for more advice and examples.

Describing results often leads to another decision point, which returns the flow of the game to step 1.

\subsection{Passing Time}
The game has a rhythm and flow that includes periods of action and excitement interspersed with lulls. Think of how movies show time passing between scenes. When an encounter ends, you can move on to the next one. You can often gloss over hours of travel with a quick narrative summary (see the "\textit{Travel}" section in \textit{chapter 2} for more advice). Similarly, if a rest period passes uneventfully, tell the players that and move on. Don't make the players spend time discussing which character cooks what for dinner unless they enjoy such descriptions. It's OK to gloss over mundane details and return to the action as quickly as possible.

Expect players to discuss the events of the game, spend time planning, and engage in long conversations in character. You don't need to be involved in those discussions unless they have questions for you. Learn to recognize the times when you can take a break as the DM, and then resume the action as soon as everyone's ready.

\begin{DndSidebar}[float=!b]{Taking Breaks}
When you finish a lengthy combat encounter or a tension-filled scene, or if you need time to think, take a quick break. Give your brain a few moments to refocus, relax, or prepare for the next encounter. It's OK to leave the players in suspense during a break while you figure out the consequences of their actions.



\end{DndSidebar}

\subsection{Ending a Session}
Try not to end a game session in the middle of an encounter. It's difficult to keep track of information such as Initiative order and other round-by-round details between sessions. An exception to this guideline is when you purposely end a session with a cliffhanger, where the story pauses just as something monumental happens or some surprising turn of events occurs. A cliffhanger can keep players intrigued and excited until the next session.

If a player missed a session and you had that player's character leave the party for a while, make sure that there's a way to bring the character back when the player returns. Sometimes a cliffhanger can serve this purpose: the character charges in to help their beleaguered companions.

Allow a few minutes at the end of play for everyone to discuss the events of the session. Ask your players what parts of the session they liked and what they would have liked to see more. Take notes on what happened and the situation at the end of the session so you can refer back to those notes as you prepare the next session.

\section{Example of Play}
These pages present a short example of play, similar to the ones in the \textit{Player's Handbook}, to illustrate how everything outlined in the "How to Run a Session" section works in practice. In this example, the Dungeon Master is running an adventure ("\textit{The Fouled Stream}") from \textit{chapter 4}. The four players are Amy (playing Auro, a Halfling Rogue), Maeve (playing Mirabella, an Elf Wizard), Phillip (playing Gareth, a Human Cleric), and Russell (playing Shreeve, a Goliath Fighter).

The DM starts by asking the players to recap the action of the previous session, most of which consisted of creating characters.


\begin{description}
\item[Jared (as DM):]
Last session, we met our four heroes in the little farming village of High Ery. Who remembers what happened?
\item[Amy:]
We were at a village council meeting about the weird stuff in the river making the fish inedible. We volunteered to investigate.
\item[Russell:]
So we set out and followed the river upstream. At the first fork, we met a treant named Borogrove who pointed us to a cave that was the source of the polluted stream.
\item[Amy:]
Before he wandered off, he gave us a magic acorn, and that's where we ended last week.
\item[Jared:]
So now you're in this gloomy forest. Dry leaves rustle under your feet. You're still beside the stream, which looks murky and unwholesome beneath the shadowy trees. What do you want to do now?
\item[Russell (as Shreeve):]
We continue upstream? \textit{(The others nod.)}
\item[Jared:]
OK, you make your way upstream for about another hour. The farther you go, the murkier and stinkier the water becomes. Rounding a bend, you can see a cave in the hillside ahead of you. The stream tumbles from the cave mouth. There are withered shrubs clumped around the cave, apparently poisoned by this nasty water.
\item[Phillip (as Gareth):]
Into the cave!
\end{description}

{\par \vspace { 9pt plus 3pt minus 3pt } \noindent  \DndFontTableTitle{(1)} \nopagebreak }

\begin{description}
\item[Jared:]
Who's leading the way?
\end{description}


\begin{description}
\item[Russell:]
I'll go first.
\end{description}

{\par \vspace { 9pt plus 3pt minus 3pt } \noindent  \DndFontTableTitle{(2)} \nopagebreak }

\begin{description}
\item[Jared:]
The cave entrance is ten feet wide, with the stream running right down the middle. Do you want to go single file or two abreast?
\item[Phillip:]
I don't love the idea of stepping across the stream. Let's go single file, staying on this side of the water. \textit{(Everyone else agrees.)}
\end{description}


\begin{description}
\item[Jared:]
OK, who's second?
\item[Phillip:]
Gareth will go second.
\item[Maeve (as Mirabella):]
I'll be third in line.
\item[Amy (as Auro):]
I'll make sure nothing's following us.
\item[Jared:]
OK, Shreeve, as you reach the cave mouth, you hear the shrubs rustling.
\item[Russell:]
Oh, I should've checked to make sure nothing was hiding in the shrubs!
\end{description}

{\par \vspace { 9pt plus 3pt minus 3pt } \noindent  \DndFontTableTitle{(3)} \nopagebreak }

\begin{description}
\item[Jared:]
In fact, the shrubs themselves are moving. They're not rooted at all—each one has two little legs and sharp claws! Everyone, roll Initiative.
\item[Russell:]
How many shrubs are attacking?
\item[Jared:]
Six. Auro, what's your Initiative?
\item[Amy:]
I got a 14.
\item[Russell:]
Shreeve goes on 5.
\item[Maeve:]
A natural 20 gives me a 21!
\item[Phillip:]
19 for Gareth.
\end{description}


\begin{description}
\item[Jared:]
Mirabella, you're first. What do you do?
\end{description}

\begin{DndSidebar}[float=!b]{1}
The DM knows something the players don't: the withered shrubs are actually monsters. It's important to establish which characters are closest to the hidden monsters.



\end{DndSidebar}

\begin{DndSidebar}[float=!b]{2}
By asking the players to choose their characters' marching order, the DM cleverly pivots away from the withered shrubs. The players don't realize their characters are in danger, and the DM is waiting for the right time to reveal the hidden monsters.



\end{DndSidebar}

\begin{DndSidebar}[float=!b]{3}
The DM rolls Initiative just once for all six monsters and writes down that they'll go on Initiative count 17. The DM then goes around the table to get each player's Initiative roll. See the "\textit{Running Combat}" section in \textit{chapter 2} for advice about rolling and tracking Initiative.



\end{DndSidebar}

\begin{DndSidebar}[float=!b]{4}
The DM doesn't have the exact positions of the monsters mapped out on a grid, but it's fair to assume that they're clumped close together as they move to attack the characters.



\end{DndSidebar}

\begin{DndSidebar}[float=!b]{5}
It's always fair for the DM to expect players to explain what their spells and abilities do. The DM has enough to keep track of!



\end{DndSidebar}

\begin{DndSidebar}[float=!b]{6}
Asking for the spell's damage allows the DM to roll a saving throw for each monster and mark off the right amount of damage for that one. In this case, though, the monsters have Vulnerability to Fire damage (because they're just dry shrubs) and so few Hit Points that they'll die no matter what they roll.



\end{DndSidebar}

\begin{DndSidebar}[float=!b]{7}
It's not Shreeve's turn, but the DM decides to allow the Goliath Fighter to step in the way of the monster's attack because it gives Shreeve a fun heroic moment. The DM changes the monster's target to Shreeve and makes an attack roll.



\end{DndSidebar}


\begin{description}
\item[Maeve:]
How many of these walking bundles of kindling can I get in a 15-foot Cone?
\end{description}

{\par \vspace { 9pt plus 3pt minus 3pt } \noindent  \DndFontTableTitle{(4)} \nopagebreak }

\begin{description}
\item[Jared:]
There are three on your side of the stream and three on the other side. You can get either group in your Cone.
\end{description}


\begin{description}
\item[Maeve:]
Mirabella puts her thumbs together and wiggles her fingertips. \textit{(Maeve mimics this action.)} Fire shoots out from her fingers, catching the ones on our side of the stream. \textit{Burning Hands}!
\end{description}

{\par \vspace { 9pt plus 3pt minus 3pt } \noindent  \DndFontTableTitle{(4)} \nopagebreak }

\begin{description}
\item[Jared:]
OK, what do I need to do?
\end{description}


\begin{description}
\item[Maeve:]
The shrub things need to make Dexterity saving throws. The DC is 14.
\end{description}

{\par \vspace { 9pt plus 3pt minus 3pt } \noindent  \DndFontTableTitle{(6)} \nopagebreak }

\begin{description}
\item[Jared:]
And how much damage do they take?
\item[Maeve:]
\textit{(Maeve rolls 3d6 for the spell's damage.)} 13 Fire damage if they fail the save, 6 if they succeed.
\item[Jared:]
Magical fire tears through them and leaves smears of ash behind! Anything else, Mirabella?
\end{description}


\begin{description}
\item[Maeve:]
My work here is done. \textit{(She mimes blowing smoke away from her fingertips.)}
\item[Jared:]
Gareth, you're up next.
\item[Phillip:]
Gareth holds his Holy Symbol and utters an imprecation while pointing at the closest shrub and casting \textit{Toll the Dead}. The sound of a bell tolls, and the shrub makes a Wisdom save, DC 14.
\item[Jared:]
Well, I rolled a 1.
\item[Phillip:]
It takes 7 Necrotic damage!
\item[Jared:]
Whatever moisture was in this "bundle of kindling" seems to dry up, and the thing keels over dead. Anything else, Gareth?
\item[Phillip:]
He glares menacingly at the other shrubs.
\item[Jared:]
OK, their turn. One skitters toward Mirabella!
\end{description}

{\par \vspace { 9pt plus 3pt minus 3pt } \noindent  \DndFontTableTitle{(7)} \nopagebreak }

\begin{description}
\item[Russell:]
Can I interject myself between it and Mirabella?
\item[Jared:]
Sure, I'll allow it. You step into the monster's path and... \textit{(The DM makes an attack roll for the monster but rolls a 7, which isn't going to hit.)} It tears at your cloak but fails to wound you. The other one has lost any interest in fighting, and it starts running away. Now it's Auro's turn.
\end{description}


\begin{description}
\item[Amy:]
Auro looks at the one that just attacked Shreeve and pulls out his dagger. I get a 23 to hit!
\item[Jared:]
That hits! What's your damage?
\item[Amy:]
Since Shreeve is next to it, I can use my Sneak Attack! The shrub takes 12 Piercing damage.
\item[Jared:]
It's felled! Mirabella, the last one is running away. Will you let it escape?
\item[Maeve:]
I think Borogrove would be disappointed in us if we let it escape into the woods. I'll cast \textit{Fire Bolt}, getting a 14 to hit.
\item[Jared:]
You nailed it.
\item[Maeve:]
It takes 10 Fire damage!
\item[Jared:]
Yeah, the last shrub is incinerated. Well done!
\end{description}

\begin{DndSidebar}[float=!b]{The Rule of Fun}
D\&D is a game, and everyone should have fun playing it. Everyone shares equal responsibility in moving the game along, and everyone contributes to the fun when they treat each other with respect and consideration: talking through disagreements among players or their characters, and remembering that arguments or mean-spirited squabbles can get in the way of the fun.



People have many different ideas about what makes D\&D fun. The "right way" to play D\&D is the way you and your players agree to and enjoy. If everyone comes to the table prepared to contribute to the game, the entire table is likely to have a wonderful and memorable time.



\end{DndSidebar}

\section{Every DM Is Unique}
The preceding example of play shows how one Dungeon Master might run an encounter, but no two DMs run the game in exactly the same way—and that's how it should be! You'll be most successful as a DM if you choose a play style that works best for you and your players.

\subsection{Play Style}
Here are some questions that can help you define your unique style as a DM and the kind of game you want to run:


\begin{description}
\item[Hack and Slash or Immersive Roleplaying?]
Does the game focus on combat and action or on a rich story with detailed NPCs?
\item[All Ages or Mature Themes?]
Is the game for all ages, or does it involve mature themes?
\item[Gritty or Cinematic?]
Do you prefer gritty realism, or are you more focused on making the game feel cinematic and superheroic?
\item[Serious or Silly?]
Do you want to maintain a serious tone, or is humor your goal?
\item[Preplanned or Improvised?]
Do you like to plan thoroughly, or do you prefer to improvise?
\item[General or Thematic?]
Is the game a mixture of themes and genres, or does it center on a particular theme or a genre such as horror?
\item[Morally Ambiguous or Heroic?]
Are you comfortable with moral ambiguity, such as allowing the characters to explore whether the end justifies the means? Or are you happier with straightforward heroic principles, such as justice, sacrifice, and helping the downtrodden?
\end{description}

\subsection{House Rules}
House rules are new or modified rules you add to your game to make it your own and to enhance the style you have in mind for your game. Before you establish a house rule, ask yourself two questions:


\begin{itemize}
\item Will the rule or change improve the game?
\item Will my players like it?
\end{itemize}

If you're confident that the answer to both questions is yes, give the new rule a try. Present house rules as experiments, and ask your players to provide feedback on them. If you introduce a house rule that isn't fun, remove or revise the rule.

\subsubsection{Recording Rules Interpretations}
If a question about the interpretation of a rule comes up in your game, record how you decide to interpret it. Add that to your collection of house rules so you and the players can reference it when the rule comes up again later.

\subsection{Atmosphere}
Some DMs use music to create an appropriate atmosphere for their game sessions. They might use soundtracks from adventure movies or video games, although classical, ambient, or other music styles can also work well.

Some DMs adjust lighting or use sound effects. Miniatures and dioramas can contribute to the game's atmosphere and help players visualize events. Check with your players, though: some might find music, lighting, or sound effects distracting; might prefer not to be startled by loud noises; or might need to avoid certain lighting effects.

\subsection{Delegation}
If there are parts of the game you prefer not to handle yourself, assign them to players who enjoy them. If you don't want to break your narrative stride by looking up a rule, designate another player to be the rulebook reference expert. If you don't like tracking Initiative, ask another player to do so.

\subsection{Learning by Observing}
One of the best ways to learn how to run a D\&D game is to observe other DMs in action. Another DM can give you a solid foundation for understanding the role—as well as inspire you with cool things you can do in your games.

You can use these questions to help you reflect on a game you observe:


\begin{description}
\item[Beginning the Session.]
How did the DM start the session? Was there a recap?
\item[Body Language.]
What gestures did the DM use when describing a scene? How did the DM's body language change when playing different NPCs?
\item[DM Voice.]
Did the DM use different voices or mannerisms for NPCs? Did the DM change the pitch or tempo of narration in different situations?
\item[Player Participation.]
Did the players participate in the world-building or make decisions that seemed to send the adventure in an unexpected direction? How did the DM handle it?
\item[Rules Adjudication.]
To what extent did the DM lean on the rules to adjudicate outcomes? Did the DM adjudicate situations wisely or in ways that made the game fun to watch?
\item[Three Pillars.]
How much of the session was taken up by combat, exploration, or social interaction?
\item[Tone and Mood.]
How would you describe the tone and mood of the game? Did it change over the course of the session?
\item[Turns of Phrase.]
Were there any words or bits of narration you really liked? (If so, jot them down.)
\end{description}

\textbf{World-building}. What elements of the DM's world or the adventure grabbed your attention?

\section{Ensuring Fun for All}
Ahead of the game, if you haven't done so already, discuss with your players the experience you're all hoping for, as well as topics, themes, and behavior that might spoil someone's enjoyment of the game.

\subsection{Mutual Respect}
Whether you're playing with long-time friends or strangers, it's important to create a foundation of mutual trust. The best games happen when everyone at the table feels safe enough to be themselves, speak up, and get into character.

It's up to everyone to uphold the principles of respect. Difficult conversations often fall on the DM to lead, but they don't have to. If one player's behavior is interfering with everyone else's enjoyment, everyone has a stake in helping to resolve the issue.

\subsubsection{Setting Expectations}
Before you assemble a group around a game table, pitch the adventures you're thinking about running to your prospective players. Note the in-world conflicts that might arise, the setting's overall tone, and the themes you'd like to explore. (The "\textit{Every DM Is Unique}" section earlier in this chapter can help you describe your game to others.)

Telling players what to expect prepares them as they imagine what sorts of characters they could create and launches conversations about content to be embraced and avoided. You don't need to reveal the major plot points or twists in your story, but share the themes you're interested in exploring, the kinds of stories you're inspired by, and which \textit{flavors of fantasy} (outlined in \textit{chapter 5}) interest you. Being transparent with your players allows them to decide if this is a game they want to play, which is best to know before play begins.

Being clear about your expectations and making sure you understand your players' expectations in return can help ensure a smooth game. Take your players' opinions and desires seriously, and make sure they take yours just as seriously. Ideally, you'll find a style of play that suits everyone.

\begin{DndSidebar}[float=!b]{Using the Game Expectations Sheet}
The Game Expectations tracking sheet is a tool you can use to set expectations at the start of a game and ensure the game is fun for everyone.



Before distributing the sheet to players, fill in the two topmost boxes:




\begin{description}
\item[Game Theme and Flavor.]
In this box, broadly describe the direction you envision for your game. See the "Setting Expectations" section for the kinds of information to include here.
\item[Potentially Sensitive Elements.]
If you know that some elements of the game might run up against some players' limits, list those elements in this box. See the "Hard and Soft Limits" section for examples.
\end{description}



Once the above information is added, give a copy of the sheet to each player. Players can fill out their sheets anonymously, but ask each of them to add the following information:




\begin{description}
\item[Limits.]
Using an \textbf{X} for a hard limit or a question mark for a soft limit, indicate any potentially sensitive elements that are problematic. Add any other elements to avoid.
\item[Hopes, Expectations, and Concerns.]
In the last two boxes, share any hopes and expectations for the game, and list any concerns about behavior at the table.
\end{description}



Collect all the sheets, and gather your players' limits into a separate, anonymous document the whole group can access.



\end{DndSidebar}

\subsubsection{Hard and Soft Limits}
Beyond the general themes and flavors of fantasy you're interested in exploring in your campaign, it's important to have a conversation with your players about topics that can be sensitive or uncomfortable. It can be helpful to discuss these topics in terms of soft and hard limits:


\begin{itemize}
\item A soft limit applies to a topic that should be handled carefully, as it might create unwelcome anxiety, fear, or discomfort.
\item A hard limit applies to a topic that should not be mentioned or described.
\end{itemize}

DMs and players can have phobias or triggers that others might not be aware of. Any in-game topic or theme that makes a member of the gaming group feel unsafe (a hard limit) must be avoided. If a topic or theme makes one or more players nervous but they consent to include it in-game (a soft limit), incorporate it with care, if at all, and be ready to quickly veer away from it if needed.

Common in-game limits apply to topics such as intraparty romance, sex, exploitation, racism, enslavement, and violence toward children and animals. Limits can also apply to certain creatures, such as spiders, snakes, rats, and demons. It's also important to discuss limits around what harm might befall characters, including mind-control magic, helplessness, and death.

That said, D\&D is a game that has in-world conflicts and mayhem. Certain core elements of the game are difficult to ignore. For example, taking damage isn't a limit you can work around easily. Similarly, character death is something that happens from time to time, though the game has ways to counteract or avoid it (see "\textit{Death}" in \textit{chapter 3} for suggestions).

\subparagraph{Communicating Limits}
Make sure everyone is comfortable with how the discussion of limits takes place. Players might not want to discuss limits aloud, especially if they're new to roleplaying games or haven't spent a lot of time with other members of the group. One way to alleviate such discomfort is to provide a way for players to share limits anonymously. Everyone can jot down their limits on an anonymous survey, such as the Game Expectations tracking sheet in this chapter.

Compile limits into a list that can be shared with the group. Limits aren't negotiable, and everyone in the group needs to respect them.

The start of a campaign is a great time to have this discussion, but further discussion is warranted each time a new player joins the group or when the campaign has a shift in story or tone. Someone might cross a line and need to be reminded of a limit, or someone might not think to include some of their limits in the initial discussion. Players can also discover new limits as the campaign unfolds. Check in with the group every few sessions to make sure everyone's comfortable with how the game is developing, updating the group's limits as needed.

\subparagraph{Shifting Limits}
Encourage players to bring any additional limits to you, privately or in the moment, so you can add them to the list. Trust that players know their needs best, and update the game accordingly.

\subparagraph{Limits in Play}
Since D\&D is improvisational, the game can go in unexpected directions. It's helpful to have an agreed-on signal that players can use to communicate that a limit has been violated, allowing you to adjust quickly. That signal might be a gesture (such as crossing the arms in an X or raising a palm in a "stop" gesture), a code word or phrase, touching or lifting a designated object, or anything else your group agrees on. Players should also feel safe to say "stop" and pause the game until the issue is resolved. The person who invokes the signals can comment on what they want adjusted but doesn't have to explain why the content is objectionable. The signal shouldn't trigger a debate or discussion: thank the player for being honest about their needs, set the scene right, and move on.

Make it clear to players that if a person isn't comfortable using the signal, they can step away from the game or call for a break to talk to you privately. Players may also give a friend permission to use the signal on their behalf. As the DM, lead by example. Take your players' needs seriously, and use every tool at your disposal to adjust how your shared story plays out.

\subsubsection{Intra-party Conflict}
When there's conflict between characters in an adventuring party, it's usually a sign that one of three things is going on:


\begin{description}
\item[Disruptive Player.]
A player is exhibiting antisocial behavior in the game. How to deal with it is covered in the "\textit{Antisocial Behavior}" section.
\item[Player Conflict.]
Conflicts between characters sometimes surface conflicts between players. These conflicts are best handled away from the gaming table. Encourage the players to resolve their conflict outside the game. If that conflict keeps arising at the game table, you might need to ask them to step away from the campaign for a while or leave the game entirely.
\item[Roleplaying.]
Conflicts between characters aren't always bad. It's OK for characters (and players) to disagree about how to deal with a captured enemy or which side to back in a brewing war. If the disagreement gets too heated, take a break and perhaps discuss, out of character, how the players would like to proceed.
\end{description}

If you can't tell which of these dynamics is in play, have a conversation with the players about it.

\subsection{Respect for the Players}
Your players need to know from the start that you'll run a game that is fun, fair, and tailored for them; that you'll allow each of them to contribute to the story; and that you'll pay attention to them when they take their turns. Your players also count on you to make sure an adventure's threats don't target them personally. Never make players feel uncomfortable or threatened.

\begin{DndSidebar}[float=!b]{Do You Really Do That?}
Can players retract what they just said their characters did? Some DMs take a hard-line position: "If you said it, your character did it." Such a strict position tends to make players much more careful about what they say, which can dampen the atmosphere and discourage humor.



Other DMs let players change their minds freely. This creates a more relaxed mood at the table, which might slow the pace of the game.



A common compromise is to rule that players can retract or change anything their characters did up until the point they learn the consequences of their actions. Once you describe what happens as a result, it's too late for the players to change their minds.



\end{DndSidebar}

\subsubsection{Sharing the Spotlight}
As the DM, don't play favorites. Don't let one player do all the talking, and make sure you check in about what every character is doing, especially during periods of exploration and social interaction, rather than focusing just on one player's character.

Sometimes you'll encounter players who tell other players what their characters should do, claim the best magic items for themselves, bully the other players, and refuse to share the spotlight. Away from the game, point out that the player's behavior is spoiling the fun for others, and ask the player to tone it down. If the player refuses to change this behavior, ask the player to leave the group.

Some problems arise when a player assumes that their particular style of play is superior to others, and they lose patience with encounters tailored to other players' preferences. Remind the impatient player (perhaps away from the table) that you have a group to please, not just one player.

\subsubsection{Tragic Limits}
Some players resist getting invested in the world of the game because they don't want to endure the pain of seeing the people and places they care about threatened or destroyed. Other players gleefully detail a backstory full of beloved NPCs, fully expecting the DM to use those people as bait, tragic victims, and unexpected villains. It's important to understand your players' preferences so you neither alienate the players by callously destroying what they love nor bore them by leaving their backstory out of the campaign story.

When you have antagonists threaten the people and places the characters love, be sure the characters have a chance to stave off the worst outcome. During the game, characters should have the opportunity to avoid or mitigate losses in heroic ways, with tragedy being a consequence of the characters' actions and decisions, not a foregone conclusion. Moments of helplessness in the face of devastating tragedy are better suited for character backstories.

\subsubsection{DM Die Rolling}
Should you hide your die rolls behind a DM screen, or should you roll your dice in the open for all the players to see? Choose either approach, and be consistent. Each approach has benefits:


\begin{description}
\item[Hidden Die Rolls.]
Hiding your die rolls keeps them mysterious and allows you to alter results if you want to. For example, you could ignore a Critical Hit to save a character's life. Don't alter die rolls too often, though, and never let the players know when you fudge a die roll.
\item[Visible Die Rolls.]
Rolling dice in the open demonstrates impartiality—you're not fudging rolls to the characters' benefit or detriment.
\end{description}

Even if you usually roll behind a screen, it can be fun to make an especially dramatic roll where everyone can see it.

\subsubsection{Overly Cautious Players}
Overly cautious players can slow down the game by checking every flagstone, door, and wall in a dungeon for traps and hidden dangers. Sometimes this behavior is a learned response to too many unpleasant surprises in past adventures, and sometimes it's just a manifestation of players' personalities.

Here are some in-game techniques you can use to encourage your players to act boldly:


\begin{description}
\item[Avoid Random Perils.]
Avoid traps and ambushes that feel random and have little importance to the rest of the adventure.
\item[Create Time Pressure.]
Set up a situation where the characters are racing toward a goal or destination. (Use this technique with caution, as time pressure can increase players' anxiety.)
\item[Telegraph Encounters.]
Give players advance warning that an encounter is imminent. Maybe they hear the heavy footfalls of a giant or see a dragon flying overhead before they have to confront it. This can encourage your players to move toward or away from the encounters rather than anxiously anticipating an ambush.
\end{description}

If these in-game techniques don't have the desired effect, have a conversation outside the game with your players about which game elements are causing them to play in an overly cautious way. Come to an agreement that those elements won't appear in your game, as keeping the game moving will result in a better experience.

\subsection{Respect for the DM}
As the DM, you have the right to expect your players to respect you and the effort you put into making a fun game for everyone. The players need to let you direct the campaign (with their input), arbitrate the rules, and settle arguments. And when you're narrating the action of the game, the players should be paying attention.

\subsubsection{Player Die Rolling}
Players should roll their dice in full view of everyone. If a player scoops up their dice before anyone else can see what they rolled, encourage that player to be less secretive.

When a die falls on the floor, do you count it or reroll it? When it lands cocked against a book, do you pull the book away and see where it lands or reroll the die? Work with your players to answer these questions, and record the answers as house rules.

\subsubsection{The Social Contract of Adventures}
You must provide reasonably appealing reasons for characters to undertake the adventures you prepare. (See "\textit{Draw In the Players}" in \textit{chapter 4} for advice on this topic.) In exchange, the players should go along with those hooks. It's OK for your players to give you some pushback on why their characters should want to do what you're asking them to do, but it's not OK for them to invalidate the hard work you've done preparing the adventure by willfully going in a different direction.

If you feel like you're keeping up your end of the bargain but your players aren't, have a conversation with them away from the gaming table. Try to understand what hooks would motivate their characters, and make sure the players understand the work you put into preparing adventures for them.

\subsubsection{Rules Discussions}
Work out a policy about rules discussions at the table. Some groups don't mind putting the game on hold while they discuss different interpretations of a rule. Others prefer to let the DM make a call and continue playing. If you gloss over a rules issue in play, make a note of it and return to the issue later.

Some players like to use the rules to argue against your decisions. While such players can be helpful when you're stuck or make a rules mistake that's easily corrected, players who argue the rules too often can disrupt the flow of the game.

If a player wants to pause play to find a specific rule or reference, you can invite the player to search for it while you and the rest of the players continue the game. That player's character essentially steps out of the game for as long as it takes. Monsters don't attack the character, and the character takes the Dodge action in combat until the player rejoins the group. This solution allows the other players to keep playing instead of letting one player stop the game.

\subsubsection{Character Knowledge}
Encourage players to play their characters within the limits of what the characters know and understand. It can be helpful to maintain the distinction between player and character knowledge by simply asking players, "What do your \textit{characters} think?"

Anachronistic thinking is another potential pitfall. You might need to remind players that their characters don't know how to make things that don't exist in the game world, such as modern firearms or antibiotics, and they don't have the players' understanding of modern science (which might not apply in the game universe anyway).

Similarly, sometimes a player is familiar with the published adventure you're running or knows the \textit{Monster Manual} backward and forward. Encourage the player to keep that knowledge separate from their character's knowledge and allow the other players to discover it through play.

\subsubsection{Antisocial Behavior}
People often play D\&D because it lets them, through their characters, do things they can't do in real life—fight monsters, cast spells, and so on. However, for some players, this means wreaking havoc in towns or betraying their allies. What they want in the game has nothing to do with heroic adventure, but with using the game rules to act out antisocial fantasies.

If this behavior comes up in your game, it might be time to reopen the conversation about the kind of game you want to play. If it's just one player causing the trouble, it's perfectly appropriate to issue an ultimatum: an out-of-control player who wants to continue playing with the group must stop being disruptive and play as part of a team. Don't let players get away with being jerks to the other players using the excuse, "that's what my character would do."

\subparagraph{Evil Characters}
Players who want to play evil characters might be looking to carry out antisocial behavior in the game. If a player asks for permission to play an evil character or comes to the table with one already made, talk to that player about what they have in mind and make sure their plans square with the group's expectations for your game. Sometimes a player wants to explore playing an evil character for perfectly good (and nondisruptive) reasons, and sometimes a whole group decides it might be fun to play evil characters together. These are valid options, as long as everyone's on the same page about how the campaign will go.

\subsubsection{Players Exploiting the Rules}
Some players enjoy poring over the D\&D rules and looking for optimal combinations. This kind of optimizing is part of the game (see "\textit{Know Your Players}" in \textit{chapter 2}), but it can cross a line into being exploitative, interfering with everyone else's fun.

Setting clear expectations is essential when dealing with this kind of rules exploitation. Bear these principles in mind:


\begin{description}
\item[Rules Aren't Physics.]
The rules of the game are meant to provide a fun game experience, not to describe the laws of physics in the worlds of D\&D, let alone the real world. Don't let players argue that a bucket brigade of ordinary people can accelerate a spear to light speed by all using the Ready action to pass the spear to the next person in line. The Ready action facilitates heroic action; it doesn't define the physical limitations of what can happen in a 6-second combat round.
\item[The Game Is Not an Economy.]
The rules of the game aren't intended to model a realistic economy, and players who look for loopholes that let them generate infinite wealth using combinations of spells are exploiting the rules.
\item[Combat Is for Enemies.]
Some rules apply only during combat or while a character is acting in Initiative order. Don't let players attack each other or helpless creatures to activate those rules.
\item[Rules Rely on Good-Faith Interpretation.]
The rules assume that everyone reading and interpreting the rules has the interests of the group's fun at heart and is reading the rules in that light.
\end{description}

Outlining these principles can help hold players' exploits at bay. If a player persistently tries to twist the rules of the game, have a conversation with that player outside the game and ask them to stop.

\begin{DndSidebar}[float=!b]{Knowing the Rules}
You don't have to be an expert on the rules to be a good DM. Of course it's helpful to be familiar with the rules, especially the ones in the \textit{Player's Handbook}, but facilitating fun is more important than implementing the rules perfectly. If you're not sure how to apply the rules in a situation, you can always ask the opinion of the players as a group. It might take a few minutes, but it's usually possible to reach an answer that feels fair to everyone, and that's more important than a "correct" answer.



You don't need to know every spell in the \textit{Player's Handbook} or the features of every class. Set the expectation that players are responsible for telling you what their abilities and spells do.



\end{DndSidebar}

\subsection{Rules for the Virtual Table}
Setting expectations is just as important in a digital environment as in person.

Some groups confine out-of-character jokes, comments, and memes to a text channel, keeping the voice channel focused on the game. But some groups find it distracting to have a separate conversation unfolding in text while the game is going on. Choose an option that works best for your group.

Who moves tokens on a virtual tabletop? Are players expected to use the built-in dice roller, or is it OK to roll physical dice and report the result? The particular technology you're using might dictate answers to these questions or raise other questions you'll need to sort out as you play.

\chapter{Running the Game}
Building on the basics laid out in \textit{chapter 1}, this chapter goes into more depth on running a D\&D game as Dungeon Master.

\section{Know Your Players}
While your players' role is to create characters (the protagonists of the campaign), breathe life into them, and steer the campaign through their actions, your role as Dungeon Master is to keep the players immersed in the world you've created and to give the characters the opportunity to do awesome things.

Knowing what your players enjoy most about the D\&D game helps you create and run adventures that they will enjoy and remember. Once you know which of the following activities each player in your group enjoys, you can tailor adventures to your players' preferences.

It's rare to gather a table of players who all enjoy the same aspects of the game. The trick is to find a balance so everyone can get some enjoyment out of each game session, even if certain encounters don't match their preferences. At best, a group of players is a lot like their characters, in that having different interests and capabilities enables them to handle a broad range of challenges.

\subsection{Acting}
Players who enjoy acting like to embody their characters' personalities, perspectives, and attitudes. They might like dressing up or using their characters' voices while playing. They enjoy social interactions with NPCs, monsters, and their fellow party members.

Engage players who like acting by...


\begin{itemize}
\item Giving them opportunities to develop their characters' personalities and backgrounds.
\item Allowing them to interact regularly with NPCs.
\item Highlighting the roleplaying elements of combat encounters.
\item Incorporating elements from their characters' backstories into your adventures.
\end{itemize}

\subsection{Exploring}
Players who desire exploration want to experience the wonders that a fantasy world has to offer. They want to know what's around the next corner or hill and like to find hidden clues and treasure.

Engage players who like exploration by...


\begin{itemize}
\item Dropping clues that hint at things yet to come.
\item Letting them find things when they take the time to explore.
\item Providing evocative descriptions of exciting environments and using interesting maps and props.
\item Giving monsters secrets for the players to uncover or cultural details for them to learn.
\end{itemize}

\subsection{Fighting}
Players who enjoy fantasy combat like the excitement of battling villains and monsters. They thrive in situations that can best be resolved in combat, favoring bold action over negotiation or investigation.

Engage players who like fighting by...


\begin{itemize}
\item Springing unexpected combat encounters.
\item Vividly describing the havoc their characters wreak with their attacks and spells.
\item Including combat encounters with large numbers of less powerful monsters.
\end{itemize}

\subsection{Instigating}
Players who like to instigate action are eager to make things happen, even if that means taking perilous risks. They would rather rush headlong into danger and face the consequences than cautiously plan their actions.

Engage players who like to instigate by...


\begin{itemize}
\item Allowing their actions to affect the environment.
\item Including things in your adventures to tempt them.
\item Letting their actions put the characters in a tight spot.
\item Including encounters with NPCs who are as feisty and unpredictable as the players are.
\end{itemize}

\subsection{Optimizing}
Players who enjoy optimizing their characters' capabilities like to fine-tune their characters for peak performance by gaining levels, new features, and magic items. They welcome any opportunity to demonstrate their characters' excellence.

Engage players who like optimization by...


\begin{itemize}
\item Using desired magic items as adventure hooks and rewards.
\item Including encounters that let them leverage their characters' most potent abilities.
\item Providing quantifiable rewards, like Experience Points, for noncombat encounters.
\end{itemize}

\subsection{Problem-Solving}
Players who want to solve problems like to scrutinize NPC motivations, untangle a villain's machinations, solve puzzles, and come up with plans.

Engage players who like to solve problems by...


\begin{itemize}
\item Including puzzles and tricky situations that require thinking.
\item Rewarding planning and tactics with in-game benefits.
\item Creating NPCs with complex motives.
\end{itemize}

\subsection{Socializing}
Many groups include players who come to the game primarily because they enjoy the social event and want to spend time with their friends, not because they're especially invested in any part of the actual game. These players want to participate, but they tend not to care whether they're deeply immersed in the adventure, and they don't tend to be assertive or very involved in the details of the game, rules, or story. As a rule, don't try to force these players to be more involved than they want to be.

\subsection{Storytelling}
Players who love storytelling want to contribute to a narrative. They like it when their characters are heavily involved in an unfolding story, and they enjoy encounters that are tied to and expand an overarching plot.

Engage players who like storytelling by...


\begin{itemize}
\item Using their characters' backstories to shape the stories of the campaign.
\item Making sure encounters advance the story.
\item Making their characters' actions steer future events.
\item Giving NPCs characteristics and connections that the adventurers can explore to uncover new adventure opportunities.
\item Including plot elements that call back to decisions the adventurers made earlier.
\end{itemize}

\section{Group Size}
D\&D's rules and published adventures generally assume four to six players plus the DM. The following advice helps you adjust adventures to work for smaller or larger groups.

\subsection{Small Groups}
A group that contains fewer than four players might find combat encounters difficult, especially if the party lacks important capabilities (such as armored characters to stand toe-to-toe with enemies or healing magic to keep everyone alive). You can compensate by reducing the number of monsters in a combat encounter or by giving the party resources they need, such as Potion of Healing.

You can also add party members, as described in the sections that follow.

\subsubsection{DM-Controlled Adventurer}
You can make an adventurer character of your own (sometimes called a "DM PC"—a "Dungeon Master player character") to accompany the party. This is a rewarding way for you to roleplay with your friends while they're exploring your world, but keep in mind that you'll have to run this NPC in combat.

Be sure to keep the players' characters in the spotlight, and don't take away the players' agency by having your character make decisions for the group.

\subsubsection{NPC Party Members}
You can add nonplayer characters (NPCs) to the adventuring party. Use the NPC stat blocks in the \textit{Monster Manual} to represent these supporting characters. If you don't want to run these NPCs yourself, invite one or more of your players to take on an NPC as a secondary character. These NPCs might be apprentices or employees of the adventurers, so it's natural for the main characters to take the lead in exploration and social interactions while the NPCs fade into the background.

See "\textit{Nonplayer Characters}" in \textit{chapter 3} for more information.

\subsubsection{Players with Multiple Characters}
One or more of your players can each play two characters. Running two characters at once is a challenge, so make sure those players are comfortable taking on multiple characters.

This approach works best in a game that's focused on combat, since it fills out a party with combat-capable characters. It can be difficult for a player to roleplay two characters at once. You might suggest that the player focus on roleplaying one character while relegating the other character to a supporting role.

\subsection{Large Groups}
The biggest considerations with large groups are maintaining order at the table and keeping combat from becoming too slow.

\subsubsection{Structured Turns}
If you find yourself in a situation where individual players are having trouble getting a chance to do things during exploration or social interaction, have the characters roll Initiative and act in Initiative order, just as you do in combat. Taking turns ensures that everyone has the chance to do something. Use this approach sparingly, as it can feel artificial and sometimes slows down the game.

\subsubsection{Party Leader}
Consider having the players designate a party leader, who is then the only person who tells you what the group is doing. It becomes the leader's role to work with the rest of the players to find consensus on what the group will do.

\subsubsection{Speeding Combat}
Players who have to wait a long time between their characters' turns in combat are susceptible to distraction. Consider these tips to speed combat with a large group.

\subparagraph{Be Generous with Information}
If you tell the players what the Armor Class of their opponents is, you reduce the steps of interaction needed to resolve an attack. Instead of telling you a number and asking if it hits, a player can simply tell you that an attack hits and how much damage it deals, perhaps adding some narration for good measure (see "\textit{Narration in Combat}" later in this chapter). In the same way, if you know each character's AC, you don't need to ask whether a monster's attack hits.

\subparagraph{Help Players Keep Up}
If a player isn't sure what to do on their turn in combat, help the player decide by offering a quick recap of the state of the battle. How many foes are still standing, and how hurt do they look? What's the most immediate threat to that character?

\subparagraph{Make Initiative Obvious}
Display the Initiative order to your players so they each know when their character's turn is coming up and can think ahead about what their character will do on their next turn. Using Initiative scores (see "\textit{Running Combat}" in this chapter)—and perhaps seating the players in Initiative order—can be helpful with a large group.

\subparagraph{Roll Handfuls of Dice}
Encourage players to roll the dice for their attack rolls and their damage at the same time. You can do the same.

\subsection{Absent Players}
When one of your players is absent, what do you do with that player's character? Consider the following options:


\begin{description}
\item[Fading into the Background.]
Have the character simply fade into the background. This requires everyone to step out of the game world a bit and suspend disbelief, but it might be the easiest solution. Act as if the character were absent, but don't try to come up with any in-game explanation. Monsters don't attack the character, who returns the favor. On returning, the player resumes playing as if the absence never happened.
\item[Narrative Contrivance.]
Decide the character is elsewhere while the rest of the party continues the adventure. Come up with in-game reasons for the character to temporarily leave the party and rejoin later, such as following up on a rumor or reporting back to the party's patron.
\item[Substitute Player.]
With the absent player's consent, have another player run the missing player's character, or run the character yourself if you feel you can do so. Whoever runs the character will need a copy of that character's character sheet and should strive to keep the character alive and use that character's resources wisely.
\end{description}

Give absent characters the same XP that the other characters earned each session, keeping the group at the same level.

Some groups like to work out a policy regarding how many missing players is too many to proceed. For example, your group might play as long as no more than one person is absent. If two or more people can't attend a session, consider playing a short adventure with different characters, and perhaps a different Dungeon Master, or bring out a favorite board game.

\subsection{Incorporating New Players}
When introducing a new player to the group, revisit the group's expectations and limits (see "\textit{Ensuring Fun for All}" in \textit{chapter 1}). Then have the new player create a character who is the same level as the other characters in the adventuring party.

If the new player has never played D\&D before and the rest of the group is higher than level 4, consider taking a short break from the campaign and having everyone play a new level 1 character for a session or two while the new player learns the ropes. (This can also be a good opportunity for another player to take a turn as DM.)

If you're incorporating a new character into the group in the middle of an adventure, work with that character's player to come up with a story hook for how their character joins the group, and make sure the player is happy with the choice. Suggested story hooks include the following:


\begin{description}
\item[Long-Lost Friend.]
The new character is a friend or relative of one of the adventurers. Alternatively, the new character is connected to the adventurers' patron or a member of an organization the other characters are linked to. In either case, the new character has been searching for the group, perhaps bearing important news.
\item[Rescued Prisoner.]
The new character is a prisoner of the foes the other characters are fighting. When rescued, this character joins their group.
\item[Sole Survivor.]
The new character is the sole survivor of an ill-fated group of adventurers. The new character might be able to offer a clue to help the party avoid the same grim fate that befell the other group.
\end{description}

\subsubsection{Special Guest Stars}
The story hooks for incorporating new players can also work for occasions when you want to bring a player into the group for a single session. For example, you might have a friend visiting from out of town who wants to join your game briefly. Or perhaps you have a player you're thinking about adding to the group, but you want to make sure they'll be a good addition. Incorporating an occasional guest player is also a great way to maintain a roster of players as backup in case one of your regular players has to drop out of the game.

\section{Multiple DMs}
Many gaming groups switch DMs from time to time. The following sections describe situations that allow for multiple DMs and ways multiple DMs can add to the group's fun.

\subsection{Occasional Breaks}
Take a break from being the DM if you need to recharge your creative juices, plan out the next arc of your campaign, or finish up the adventure you're working on. By taking a break, you create an opportunity for another player to assume the DM role for a session or two.

If not everyone can make it to a scheduled session, that can also be an opportunity for a different DM to run a short adventure.

\subsection{Variety Series}
Some groups don't want a long campaign with sweeping plotlines; they prefer short, unconnected adventures. With that style of game, different players might take turns as DM for one to three sessions at a time, with each adventure standing as a self-contained story.

\subsection{Concurrent Campaigns}
You and the other DMs in your group can take turns running adventures for a few weeks or months at a time, with your campaign on hold during another DM's turn. Some groups play multiple times each week, with different DMs running their campaigns on different days.

\subsection{Shared World}
Some groups take a large, established campaign setting and divide it up geographically so different DMs can run separate campaigns in the same setting. In theory, characters can travel from one DM's region of the world to another's, creating continuity in the campaign even as characters might be involved in several plotlines.

Rather than dividing a campaign geographically, you and the other DMs in your group could divide it thematically. Using the setting in \textit{chapter 5} of this book as an example, each DM could focus their campaign on one of the three overarching conflicts of that setting. This approach allows the same group of adventurers to sink their teeth into all three overarching conflicts while ensuring that each storyline feels distinct.

\subsection{Joint DMs}
Two or more DMs can share the creation of a single campaign, working together to maintain continuity from session to session and making sure that each DM's adventures advance the larger story of the world and the characters. When players who are also DMs are playing their characters, they shouldn't let their knowledge of the campaign's story influence their characters' actions. Those characters step out of the action when their players take their turns as DM.

Joint DMs can also team up to run each session of a campaign, with each DM focusing on the aspects of the game they most enjoy or the DMs trading focus from session to session. One DM might run combat description and keep a battle moving while the other focuses on miniatures and music. The two DMs can play two different NPCs in a social interaction encounter. Between sessions, they can collaborate or divide up world-building, encounter creation, and other tasks.

\section{Narration}
Using a few time-honored narrative techniques, you can immerse your players in your world and bring the game to life.

\subsection{Lead by Example}
When you roleplay and narrate with enthusiasm, you add energy to the game and draw your players into the world. Encourage the players to describe their characters' actions, then incorporate their narration into your accounts of the characters' successes and failures.

\subsection{Brevity}
Keep your descriptions short and evocative. Focus on the more important information to keep players' interest and to highlight important clues and details. Players need to know about significant features their characters can perceive—especially things like monsters in a room—before they decide what to do. Allow your players to ask follow-up questions, and provide additional description as needed.

\subsection{Atmosphere}
Bring a place to life by adding touches of atmosphere, such as a lingering smell of ash, tiny beetles skittering along the dungeon floor, or blue flowers blossoming in the otherwise desolate and gloomy graveyard. Pick a couple of senses (sight, hearing, smell, touch, or taste) to highlight.

Describe changes in the environment to direct your players' attention. For example, a bird alighting on a gravestone might draw the characters' attention to it.

\subsection{Draw Players' Attention}
Good narration invites the players to examine details of the environment that lead to encounters or important information. Anything you describe with extra, subtle details draws the players' attention. Give them just enough to invite further exploration, but don't create the equivalent of a flashing neon sign reading "This way to adventure!"

When using narration to guide your players, keep the following in mind:


\begin{description}
\item[Distinguish Options.]
When presenting options to players, add details to distinguish the options. Should the characters take the left path or the right path? Perhaps the left path smells of rot and decay, while the faint sound of lapping water comes from the right. These details give players more information to make an informed decision without explicitly telling them where to go.
\item[Don't Limit Options.]
In general, let the players use the information they're given to decide what they want to do. Don't put unnecessary limitations on the characters' actions. That said, it can be helpful (especially with new players) to offer suggestions: "You can go through the door, search the chest, or look down the shaft." Just make sure to finish by saying, "or anything else you can think of!"
\item[Don't Assume Character Actions.]
Don't assume actions on the characters' part. For example, don't say "You step into the room and look up" unless the player has already told you that's what their character is doing.
\end{description}

\subsection{Secrets and Discovery}
In the course of an adventure, the players and their characters will uncover information that was previously unknown to them. Make sure the information they need to complete the adventure is obtainable.

Don't hide important secrets or discoveries in places where the characters aren't likely to uncover them. Make sure they can't miss an important secret or discovery simply by failing an ability check, not talking to the right person, or not looking in the right place.

See also "\textit{Perception}" in this chapter for more advice on hidden secrets in adventures.

\subsubsection{Giving Information to One Player}
When one character separates from the rest of the group, it's usually OK to let the rest of the players know what happens, assuming the separated character will update the rest of the party when they're reunited. You might need to remind the other players that their characters aren't present, so they can't offer advice or information to the lone character.

Sometimes, though, you'll want to give information to just one player. It might be information you think the character won't want to share with the rest of the party, perhaps something related to elements of the character's history that are still secret. In this case, you can use one of these methods to deliver that information:


\begin{description}
\item[Aside to Player.]
Pull the player into another room, or have the other players leave the room. This approach is best if there's a whole scene that plays out with just one character involved. Try to keep these scenes brief so other players don't get bored or feel left out.
\item[Secret Message.]
If you just have a simple piece of information to convey, you can whisper to the player, pass a note to them, or send them a text or a direct message.
\end{description}

\section{Resolving Outcomes}
You decide when a player makes a D20 Test based on what the character is trying to do. Players shouldn't just roll ability checks without context; they should tell you what their characters are trying to achieve, and make ability checks only if you ask them to.

When a situation comes up and you're not sure how to adjudicate it using the rules, use these four questions to help you decide:


\begin{description}
\item[Is a D20 Test Warranted?]
If the task is trivial or impossible, don't bother with a D20 Test. A character can move across an empty room or drink from a flask without making a Dexterity check, whereas no lucky die roll will allow a character with an ordinary bow to hit the moon with an arrow. Call for a D20 Test only if there's a chance of both success and failure and if there are meaningful consequences for failure.
\item[What Kind of D20 Test?]
If a character is actively trying to do something, use an ability check (or an attack roll if the character is trying to hit something). If the character is reactively trying to avoid or resist something, use a saving throw.
\item[Which Ability Does the Test Use?]
Think about which ability has the most influence on a character's chance to succeed on the ability check or saving throw. Refer to the Abilities, Ability Checks, and Saving Throws table for guidance. Also consider whether a skill or tool proficiency might apply to an ability check.
\item[What's the DC?]
Based on how hard you think the task should be, set the DC as follows: 10 for an easy task, 15 for a moderately difficult task, or 20 for a hard task.
\end{description}

The sections that follow offer advice on how to use each kind of D20 Test, when to apply Advantage and Disadvantage, and how to decide what the consequences of success or failure might be.

\begin{DndTable}[header=Abilities, Ability Checks, and Saving Throws]{XXXX}

Ability & Score Measures... & Make an Ability Check To... & Make a Saving Throw To...\\
Strength & Physical might & Lift, push, pull, or break something & Physically resist direct force\\
Dexterity & Agility, reflexes, balance & Move nimbly, quickly, or quietly & Dodge out of harm's way\\
Constitution & Health and stamina & Push your body beyond normal limits & Endure a toxic hazard\\
Intelligence & Reasoning and memory & Reason or remember & Recognize an illusion as fake\\
Wisdom & Perceptiveness and mental fortitude & Notice things in the environment or in creatures' behavior & Resist a mental assault\\
Charisma & Confidence, poise, and charm & Influence, entertain, or deceive & Assert your identity\\
\end{DndTable}

\subsection{Ability Checks}
An ability check is a test to see whether a character succeeds at a task the character has decided to attempt. The Abilities, Ability Checks, and Saving Throws table summarizes what each ability is used for. (Constitution checks are rare, as tests of a character's endurance are usually passive or reactive, making a saving throw more appropriate.)

\subsubsection{Proficiency}
When the rules or a published adventure calls for an ability check, a skill or tool proficiency is often called out: for example, "a character who succeeds on a DC 15 Intelligence (Arcana) check can puzzle out the magic involved." Sometimes the rules allow for any one of two or more proficiencies to apply to a check. When deciding what check a character should make, be generous in determining if the character's Proficiency Bonus comes into play. You might specifically ask for an Intelligence (Arcana) check, or you can ask for an Intelligence check and let the player negotiate with you to see if one of the character's skill or tool proficiencies applies.

\subsubsection{Trying Again}
Sometimes a character fails an ability check and the player wants to try again. In many cases, failing an ability check makes it impossible to attempt the same thing again. For some tasks, however, the only consequence of failure is the time it takes to attempt the task again. For example, failing a Dexterity check to pick a lock on a treasure chest doesn't mean the character can't try again, but each attempt might take a minute.

If failure has no consequences and a character can try and try again, you can skip the ability check and just tell the player how long the task takes. Alternatively, you can call for a single ability check and use the result to determine how long it takes for the character to complete the task.

\subsubsection{Group Checks}
Group checks are a tool you can use when the party is trying to accomplish something together and the most skilled characters can cover for characters who are less adept at the task. To make a group ability check, everyone in the group makes the ability check. If at least half the group succeeds, the whole group succeeds. Otherwise, the group fails.

Group checks aren't appropriate when one character's failure would spell disaster for the whole party, such as if the characters are creeping across a castle courtyard while trying not to alert the guards. In that case, one noisy character will draw the guards' attention, and there's not much that stealthier characters can do about it, so relying on individual checks makes more sense. Similarly, don't use a group check when a single successful check is sufficient, as is the case when finding a hidden compartment with a Wisdom (Perception) check.

Consider using group checks in situations such as the following:


\begin{description}
\item[Research Tasks.]
The characters are trying to learn about an ancient prophecy from an archive. The characters can make a group Intelligence (Investigation) check to find sources in the archive; characters who are knowledgeable about such topics and about research methods in general can tell the other characters the most likely places to direct their search. If the group check is successful, the characters find enough different sources to paint a clear picture of the prophecy; otherwise, their information is incomplete.
\item[Roped Together.]
The characters are tied together as they scale a cliff or cross a rickety rope bridge. If one or two characters fail their checks, the successful characters can stop their fall and prevent disaster, but if more than half the party fails, the whole group falls. You could also apply this idea to something like a long swim, where strong swimmers can help weaker ones.
\item[Social Situations.]
One character offends someone at a diplomatic event, and a noble demands the adventurers be escorted from the premises. The characters can make a group Charisma (Persuasion) check to avoid getting thrown out; they just need a few convincing arguments and the ability to smooth over any gaffes. You could apply this same principle to other Charisma checks using Deception, Intimidation, or Performance.
\end{description}

\subsubsection{Passive Checks}
Ability checks normally represent a character's active effort to accomplish something, but occasionally you need a passive measure of how good a character is at doing a thing. Passive Perception is the most common example. (See "\textit{Perception}" later in this chapter.) You can extend the concept of a passive ability check to other abilities and skills.

For example, if your game features a lot of social interaction, you can record each character's Passive Insight score, calculated in much the same way as Passive Perception: 10 plus all modifiers that normally apply to a Wisdom (Insight) check.

\subsection{Attack Rolls}
In combat, an attack roll is used to determine whether an attack hits.

You can also use attack rolls to resolve noncombat activities that are similar to attacks in combat, such as an archery contest or a game of darts. Assign an Armor Class to the target, decide whether the character is proficient with the weapon used, then have the player make an attack roll. (See also "\textit{Degrees of Success}" in this chapter.)

\subsection{Saving Throws}
In contrast to an ability check, a saving throw is an instant response to an effect and is almost never made by choice. A saving throw makes the most sense when something bad threatens a character and the character has a chance to avoid or resist it.

Most of the time, a saving throw comes into play when an effect—such as a spell, monster ability, or trap—calls for it, telling you what kind of saving throw is involved and providing a DC for it.

In other situations that call for a saving throw, it's up to you to decide which ability score is involved. The Abilities, Ability Checks, and Saving Throws table offers suggestions.

\subsection{Difficulty Class}
You establish the Difficulty Class for an ability check or a saving throw when a rule or an adventure doesn't give you one. Choose a DC from the Typical DCs table based on the task's difficulty.

\begin{DndTable}[header=Typical DCs]{Xc}

Task & DC\\
Very easy & 5\\
Easy & 10\\
Moderate & 15\\
Hard & 20\\
Very hard & 25\\
Nearly impossible & 30\\
\end{DndTable}

The task difficulties are explained below:


\begin{description}
\item[Very Easy.]
Most people can accomplish a DC 5 task with little chance of failure. Unless circumstances are unusual, let characters succeed at such a task without making a check.
\item[Easy, Moderate, and Hard.]
These are the most common difficulties, and you can run your game using only them. A character with a 10 in the associated ability and no proficiency will succeed at an easy task around 50 percent of the time. For a moderate task, a character needs either a higher score or proficiency to have a similar chance of success, whereas a hard task typically requires both to have a similar chance. If you can't decide between two levels of difficulty, choose a DC somewhere in the middle, such as 17 or 18 for a task that is a little easier than "hard."
\item[Very Hard and Nearly Impossible.]
A DC 25 task is almost out of reach for low-level characters, but more reasonable after level 10 or so. Low-level characters have no chance to accomplish a DC 30 task, while a level 20 character with proficiency and a relevant ability score of 20 still needs a 19 or 20 on the die roll to succeed at a task of this difficulty.
\end{description}

If you're setting the DC for a saving throw, don't go lower than 10 or higher than 20. If a creature is the source of the effect forcing a saving throw, use the standard formula for calculating a save DC (see "\textit{Calculated DCs}" below).

\subsubsection{Calculated DCs}
For some ability checks and most saving throws, the rules default to the following formula:

This formula often sets the saving throw DC when a creature is casting a spell or using a special ability, but it can also apply to ability checks that are contests between two creatures. For example, if one creature is holding a door shut, use its Strength modifier and Proficiency Bonus to set the DC for opening the door. When another creature tries to force the door open, the creature makes a Strength (Athletics) check against that DC.

Another way to handle similar situations is to have one creature's ability check set the DC for another creature's check. That's how hiding works, for example: a hiding creature's total Dexterity (Stealth) check sets the DC for Wisdom (Perception) checks made to find the hidden creature.

\subsection{Advantage and Disadvantage}
Advantage and Disadvantage are among the most useful tools in your toolbox. They reflect temporary circumstances that might affect the chances of a character succeeding at a task. Advantage is also a great way to reward a player who shows exceptional creativity in play.

Characters often gain Advantage or Disadvantage through the use of special abilities, actions, spells, or other features of their classes or species. In other cases, you decide whether a circumstance would merit Advantage or Disadvantage.

As described in the \textit{Player's Handbook}, if different circumstances would give both Advantage and Disadvantage in the same situation, the Advantage and Disadvantage cancel out, regardless of how many circumstances would grant Advantage or Disadvantage.

\subsubsection{Advantage}
Consider granting Advantage when...


\begin{itemize}
\item Circumstances not related to a creature's own capabilities provide it with an edge.
\item Some aspect of the environment improves the character's chance of success.
\item A player shows exceptional creativity or cunning in attempting or describing a task.
\item Previous actions (whether taken by the character making the attempt or some other creature) improve the chances of success.
\end{itemize}

\subsubsection{Disadvantage}
Consider imposing Disadvantage when...


\begin{itemize}
\item Circumstances hinder success in some way.
\item Some aspect of the environment makes success less likely.
\item An element of the plan or description of an action makes success less likely.
\end{itemize}

\subsection{Consequences}
As a DM, you can use a variety of approaches when adjudicating success and failure to tailor the game to your liking.

\subsubsection{Success at a Cost}
When a character fails a D20 Test by only 1 or 2, you can offer to let the character succeed at the cost of a complication or hindrance. Such complications can run along any of the following lines:


\begin{itemize}
\item A character gets her sword past an enemy's defenses and turns a near miss into a hit, but she then drops the sword.
\item A character narrowly escapes the full brunt of a \textit{Fireball} spell but has the Prone condition.
\item A character fails to intimidate a kobold prisoner, but the kobold reveals its secrets anyway while shrieking at the top of its lungs, alerting other nearby monsters.
\end{itemize}

By putting the choice of success at a cost in the players' hands, and even letting them choose the setbacks, you can give players more agency in crafting the story of their characters' deeds.

\subsubsection{Degrees of Failure}
Sometimes a failed D20 Test has different consequences depending on the degree of failure. For example, a character who fails to disarm a trapped chest might accidentally spring the trap if the check fails by 5 or more, whereas a lesser failure means the trap wasn't triggered during the botched disarm attempt. Consider adding similar distinctions to other checks. Perhaps a failed Charisma (Persuasion) check means a queen won't help, whereas a failure of 5 or more means she throws the character in the dungeon for such a display of impudence.

\subsubsection{Degrees of Success}
A successful D20 Test can have degrees of success. For example, when characters participate in an archery contest, you might decide that the more an attack roll exceeds the target's AC, the higher the character's score. The archery target might have AC 11, but it has five concentric rings indicating degrees of success. So you could decide that an attack roll of 11 or 12 lands in the outermost ring, a 13 or 14 hits the next ring closer to the center, a 15 or 16 hits the third ring, a 17 or 18 hits the fourth, and a 19 or higher strikes the bull's-eye.

\subsubsection{Critical Success or Failure}
Rolling a 20 or a 1 on an ability check or a saving throw doesn't normally have any special effect. However, you can take such an exceptional roll into account when adjudicating the outcome. It's up to you to determine how this manifests in the game. One approach is to increase the impact of the success or failure. For example, rolling a 1 on a failed attempt to pick a lock might jam the lock, and rolling a 20 on a successful Intelligence (Investigation) check might reveal an extra clue.

For attack rolls, the rules cover what happens on a natural 20 (it's a Critical Hit) or a natural 1 (it always misses). Resist the temptation to add additional negative consequences to a natural 1 on an attack roll: the automatic failure is bad enough. And characters typically make so many attack rolls that they're bound to roll dozens of natural 1s over time. What might seem like an interesting consequence, like breaking the weapon used for the attack, will quickly get tiresome.

\subsection{Improvising Damage}
The Improvising Damage table gives guidelines for determining damage on the fly.

\begin{DndTable}[header=Improvising Damage]{cX}

Damage & Examples\\
1d10 & Burned by coals, hit by a falling bookcase, pricked by a poison needle\\
2d10 & Struck by lightning, stumbling into a firepit\\
4d10 & Hit by falling rubble in a collapsing tunnel, tumbling into a vat of acid\\
10d10 & Crushed by compacting walls, hit by whirling steel blades, wading through lava\\
18d10 & Submerged in lava, hit by a crashing flying fortress\\
24d10 & Tumbling into a vortex of fire on the \textit{Elemental Plane of Fire}, crushed in the jaws of a godlike creature or a moon-size monster\\
\end{DndTable}

The Damage Severity and Level table is a guide to how deadly these damage amounts are for characters of different levels. Cross-reference a character's level with the damage being dealt to gauge the severity of the damage.

\begin{DndTable}[header=Damage Severity and Level]{cXX}

Character Levels & Nuisance & Deadly\\
1–4 & 5 (1d10) & 11 (2d10)\\
5–10 & 11 (2d10) & 22 (4d10)\\
11–16 & 22 (4d10) & 55 (10d10)\\
17–20 & 55 (10d10) & 99 (18d10)\\
\end{DndTable}

Nuisance damage rarely poses a risk of death to characters of the levels shown, but a severely weakened character might be laid low by this damage.

Deadly damage poses a significant threat to characters of the levels shown and could potentially kill such a character that's missing many Hit Points.

\subsection{Improvising Answers}
With a little preparation and a lot of flexibility, you can handle any curveball your players throw at you.

One of the cornerstones of improvisational theater is called "Yes, and..." It's based on the idea that an actor takes whatever the other actors give and builds on that. A similar principle applies as you run sessions for your players. As often as possible, weave what the players give you into your story.

An equally important principle is "No, but..." Sometimes characters can't do what their players want, but you can keep the game moving forward by offering an alternative.

For example, imagine the characters are searching for a lich's lair. A player asks you if there's a mages' guild operating in a nearby city, hoping to find records that mention the lich. This wasn't a possibility you anticipated, and you don't have anything prepared for it. One option is to say yes and use the tools at your disposal to create a suitable mages' guild. By doing this, you reward the player for thinking creatively. Also, the guild can become a great source for adventure hooks.

Another option is to say no, but a solitary mage in town might possess the desired information. This approach rewards the creative player while demanding less work from you.

\subsubsection{Aids to Improvisation}
When you need to make up something on the spot—say, a mages' guild in a town where you hadn't previously planned for one—you have abundant resources to draw on:


\begin{itemize}
\item \textit{Lists of NPC names} (see "\textit{Nonplayer Characters}" in \textit{chapter 3})
\item \textit{Random tables} (such as the ones in the "\textit{Settlements}" section of \textit{chapter 3})
\item \textit{Campaign Journal} (described in \textit{chapter 5})
\item Maps (see \textit{appendix B})
\end{itemize}

\section{Running Social Interaction}
During a social interaction, the adventurers usually have a goal. They want to extract information, secure aid, win someone's trust, escape punishment, avoid combat, negotiate a treaty, or achieve some other objective. Successfully completing the encounter means achieving that goal.

Some DMs run social interaction as a free-form roleplaying opportunity, where dice rarely come into play. Other DMs resolve interactions by having characters make Charisma checks. Most games fall somewhere in between, balancing roleplaying with the occasional ability check.

\subsection{Roleplaying}
You don't need to be a practiced thespian or comedian to create drama or humor through roleplaying. The key is to pay attention to the story elements and characterizations that make your players laugh or feel emotionally engaged and to incorporate those things into your roleplaying.

\subsubsection{NPC Portrayals}
When thinking about how to roleplay an NPC or a monster, consider one or two adjectives that best describe the creature. Knowing the creature's alignment can also help with your portrayal. The classic advice for writers holds true: show, don't tell. For example, rather than describe an NPC as jocular and honest, have the NPC make frequent puns and freely share personal anecdotes.

You can further enhance your portrayal of a creature in the following ways.

\subparagraph{Use Facial Expressions}
Your facial expressions help convey a creature's emotions. Smile, scowl, snarl, yawn, or pout, as appropriate.

\subparagraph{Use Motions and Posture}
Movement and posture can help define an NPC's personality. You might reflect an archmage's displeasure by rolling your eyes and massaging your temples with your fingers. Hanging your head and looking up at the players conveys a sense of submissiveness or fear. Holding your head and chin high conveys confidence.

\subparagraph{Use Voices}
Changing the volume of your voice and borrowing speech patterns from real life, movies, or television can make NPCs distinctive.

\subsubsection{Engaging the Players}
Although some players enjoy roleplaying more than others, social interactions help immerse all players in the game. Consider the following approaches to make an interaction-heavy game session appeal to players of any tastes.

\subparagraph{Appeal to Player Preferences}
Players who like acting (see "\textit{Know Your Players}" in this chapter) thrive in social interactions, so let those players take the spotlight and inspire the other players by their example. However, be sure to tailor aspects of social interactions to fit the other players' tastes too.

\subparagraph{Involve Specific Characters}
If you have players who don't readily get involved in social interactions, you can create situations tailored for their characters. Perhaps the NPC in question is a family member or a contact of a particular adventurer and focuses attention on that character. Some NPCs might pay particular attention to characters with whom they feel kinship.

If a couple of players are doing most of the talking in a social interaction, take a moment now and then to involve someone else. You might have an NPC address another character directly: "And what about your hulking friend? What will you pledge in exchange for my favor?" If a player is less comfortable with roleplaying, you can get them involved by asking them to describe their character's actions during the conversation.

\subparagraph{Use Other Ability Scores}
Consider the following additional possibilities to give characters whose Charisma is not their strong suit a chance to shine:


\begin{description}
\item[Strength.]
An NPC won't talk to the characters until one of them agrees to an arm-wrestling match. Or a strong character needs to bodily prevent the NPC from running away.
\item[Dexterity.]
An NPC is Hostile [Attitude] toward intruders, so the characters must talk from hiding. Or the social interaction provides a distraction that allows a character to get close enough to the NPC to steal something from the NPC's pockets.
\item[Intelligence.]
An NPC's speech is so full of obscure references to a particular area of knowledge that the characters can't use the information they receive until they interpret those obscure facts. Or the NPC refuses to give a direct answer, speaking only in vague hints that the characters must piece together to get the information they seek.
\item[Wisdom.]
An NPC is hiding something important, and the characters must read the NPC's nonverbal cues to understand what's true and what's deception. Or key information is concealed in details around the room where the interaction takes place, which a perceptive character might notice.
\end{description}

\subsection{Attitude}
Each creature controlled by the DM has one of the following attitudes toward the adventurers: Friendly [Attitude], Indifferent [Attitude], or Hostile [Attitude]. The "\textit{Monster Behavior}" section in \textit{chapter 4} offers guidance to help you determine a creature's initial attitude.

Characters can shift a creature's attitude by their words or actions. For example, buying drinks for an Indifferent group of miners might shift their attitude to Friendly. When a shift occurs, describe it to your players. For example, the miners might display their newfound friendliness by imparting some useful information, offering to repay the kind gesture at a future date, or challenging the characters to a friendly drinking contest.

\subsubsection{Ability Checks in Social Interaction}
You decide the extent to which ability checks shape the outcome of a social interaction. A simple social interaction might involve a brief conversation and a single Charisma check, while a more complex encounter might involve multiple ability checks helping to steer the course of the conversation.

\subsubsection{Using the Help Action}
When a character uses the Help action to help another character influence an NPC or a monster, encourage the player of the helpful character to contribute to the conversation or, at the very least, describe what their character is doing or saying to contribute to the other character's success.

\section{Running Exploration}
Traversing a wilderness, searching a dungeon, circumventing an obstacle, finding a hidden object, investigating a strange occurrence, deciphering clues, solving puzzles, and bypassing or disabling traps are all part of exploration.

Not everything in your world needs to be explored painstakingly. For instance, you might gloss over an unimportant journey by telling the players that they spend three uneventful days on the road before reaching the next point of interest.

\subsection{Using a Map}
A map can help you or your players visualize a location or region that the characters are exploring. D\&D maps come in three varieties, with examples of all three found in \textit{appendix B} and on the poster map:


\begin{description}
\item[Dungeon Maps.]
D\&D uses the word "dungeon" loosely to describe any adventure location that has interior spaces to explore (such as a castle, tower, mansion, or subterranean complex). A dungeon map shows passages, chambers, doors, and other important features of a location.
\item[Settlement Maps.]
A map of a settlement often shows terrain (cliffs, trees, streams, and so forth) in addition to buildings, bridges, and other important features.
\item[Wilderness Maps.]
A wilderness map shows roads, rivers, terrain, and other features that might guide the characters on their travels or lead them astray. The area shown on a wilderness map might be as big as a continent or as small as a glade.
\end{description}

Often a map is intended for the DM's eyes only. You can copy portions of a DM's map to share with your players as a visual aid while omitting details that should remain hidden from them. Virtual tabletops often use "fog of war" and similar effects to obscure areas and features on the map that you want to keep hidden from the players.

Maps designed for use with miniatures (see "\textit{Miniatures}" in this chapter) tend to be player facing, revealing nothing that would spoil the adventure.

\subsection{Tracking Time}
If tracking the passage of time is important during exploration, use a time scale appropriate for the situation at hand:


\begin{description}
\item[Rounds.]
In combat and other fast-paced situations, the game relies on 6-second rounds.
\item[Minutes.]
In a dungeon or settlement, movement happens on a scale of minutes. In the Free City of Greyhawk, getting from the Silver Dragon Inn to the wharf takes about 10 minutes, whereas it takes about 1 minute to creep down a 200-foot-long hallway, another minute to check for traps on the door at the end of the hall, and 10 minutes to search the chamber beyond for anything interesting or valuable.
\item[Hours.]
A scale of hours is often appropriate for short wilderness treks. Adventurers eager to reach the lonely tower 20 miles away, at the heart of the forest, can hurry there in 5 hours' time.
\item[Days.]
For long journeys, a scale of days works best. Following the road from Veluna City to the Free City of Greyhawk, the adventurers cover 96 miles in 4 uneventful days before a bandit ambush interrupts their journey.
\end{description}

The \textit{exploration rules} in the \textit{Player's Handbook} give guidelines for determining travel time based on the characters' pace. In most cases, it's fine to estimate that time rather than calculating it down to the minute. Exceptions include situations like these:


\begin{description}
\item[Spell Timer.]
The characters might need to go somewhere or accomplish something before the duration of a spell or similar effect runs out. For example, they might use the \textit{Locate Object} spell to point them in the direction of an item they seek, so you need to know how far they get in the 10 minutes the spell lasts.
\item[Triggered Event.]
An event might occur at a specific time. For example, a door might remain open for 1 minute after the password to open it is spoken, or reinforcements might arrive 2d4 minutes after an alarm is sounded.
\end{description}

If the characters spend time working out a puzzle or talking to an NPC, you can estimate the time spent by keeping track of how much real time passes. Most combat encounters take less than 1 minute (10 rounds), but it's fair to round up to a whole minute in most cases, assuming characters take a few seconds to pull themselves together after a fight.

Use similar principles to track the passage of hours, such as when characters disguise themselves with a \textit{Seeming} spell for 8 hours to infiltrate a stronghold. In this case, it takes a lot of small tasks—or something like a Short Rest—to occupy a full hour.

\subsection{Actions in Exploration}
Most of what characters do during exploration, aside from movement, relates to just a few actions: Search, Study, and Utilize. Characters also often use the Help action to assist each other in these actions. Other actions come up only rarely.

It's seldom necessary to rely on the action rules during exploration, except to remember that a character can do only one thing at a time. A character who's busy taking the Search action to look for a secret door can't simultaneously take the Help action to assist another character who's taking the Study action to find important information in a book.

\subsubsection{Taking Turns}
Often, characters spread out across a room to investigate the elements of the room. (The \textit{exploration example} in \textit{chapter 1} of the \textit{Player's Handbook} shows this dynamic in action.) In such situations, have the characters take turns, though it's usually not necessary to roll Initiative as you would in a combat encounter. Resolve one character's actions before moving to the next.

There's no hard-and-fast rule about how long to spend on each character's activity, but make sure no one is waiting for their turn for too long.

You can build tension in an exploration encounter by shifting focus right before a character makes an ability check or opens a chest, leaving everyone eager to hear what happens next.

\subsubsection{Ability Checks in Exploration}
When a character tries to do something during exploration, you decide whether that action requires an ability check to determine success (as described in the earlier "\textit{Resolving Outcomes}" section).

Certain situations might call for a balance between ability checks and roleplaying. For example, puzzles are an opportunity for players to do some problem-solving, but players can also lean on their characters' talents and attributes to provide direction. A character who succeeds on an Intelligence (Investigation) check might notice a clue that gives the players a hint to the puzzle's solution.

\subsection{Perception}
As the DM, you're the interface between your players and the world of the game. You tell them what their characters perceive, so it's important to make sure you're telling them important information about their surroundings. The Perception skill and Wisdom checks made using it are key tools for you. This section offers guidance to help you use the Perception rules in the \textit{Player's Handbook}.

\subsubsection{When to Call for a Check}
An important time to call for a Wisdom (Perception) check is when another creature is using the Stealth skill to hide. Noticing a hidden creature is never trivially easy or automatically impossible, so characters can always try Wisdom (Perception) checks to do so.

\subparagraph{Using Passive Perception}
Sometimes, asking players to make Wisdom (Perception) checks for their characters tips them off that there's something they should be searching for, giving them a clue you'd rather they didn't have. In those circumstances, use characters' Passive Perception scores instead.

\subparagraph{Using the Investigation Skill}
The Investigation skill applies to situations where a character is using reason and deduction to arrive at a conclusion about something under examination. Investigation applies when characters are trying to figure out how a thing works—how to open a trick door, how to get into a secret compartment, and so on.

Don't use the Investigation skill to determine if a character notices something—that's the purview of Perception. For example, a successful Wisdom (Perception) check allows a character to find a secret door or something that betrays its presence, such as thin seams marking the edges of the door. If the secret door is locked, a successful Intelligence (Investigation) check would allow a character to figure out the trick to opening it—by turning a nearby statue so it faces the door, for example.

\subsubsection{Hidden Things in Adventures}
Secret doors, hidden compartments, concealed traps, and stashed treasures are common elements in adventures. When using such elements, if something is hidden, allow for the possibility that the characters might not find it. It's fine to hide extra treasures or delightful surprises, but don't hide elements that are essential to the characters' success in places where characters might not find them.

Even if the hidden objects aren't essential to the adventure's success, plant hints that clue players in to the idea that there might be something hidden for them to find. Such hints can be subtle (a character hears a strange rattle when opening the desk drawer, suggesting the presence of a hidden compartment in the back or bottom of the drawer) or obvious (clear footprints lead across the room to a blank wall that is actually a secret door). These hints let players discover fun secrets without requiring them to spend extensive time searching every square foot of every room and hallway.

\subsubsection{Perception and Encounters}
If the characters encounter another group of creatures and neither side is being stealthy, the two groups automatically notice each other once they are within sight or hearing range of one another. The Audible Distance table can help you determine the hearing range, and the following sections address visibility. If one group tries to hide from the other, use the rules in the \textit{Player's Handbook}.

\begin{DndTable}[header=Audible Distance]{Xc}

Noise & Distance\\
Trying to be quiet & 2d6 × 5 feet\\
Normal noise level & 2d6 × 10 feet\\
Very loud & 2d6 × 50 feet\\
\end{DndTable}

\subparagraph{Visibility Outdoors}
When traveling outdoors, most characters can see about 2 miles in any direction on a clear day, except where obstructions block their view. That range increases to 40 miles if they are atop a mountain or a tall hill or are otherwise able to look down on the area from a height. Lightly Obscured conditions reduce visibility: rain reduces maximum visibility to 1 mile, and fog reduces it to between 100 and 300 feet.

Outdoor terrain determines the distance at which characters encounter other creatures. The \textit{Travel Terrain} table (see "\textit{Travel}" below) gives suggested encounter distances for different types of terrain.

\subparagraph{Visibility at Sea}
From a ship's crow's nest, a lookout can see things up to 10 miles away, assuming clear skies and a relatively calm sea. Overcast skies reduce that distance by half. Lightly Obscured conditions reduce visibility just as they do on land.

\subparagraph{Visibility Underwater}
Visibility underwater depends on water clarity and the available light. Use the Underwater Encounter Distance table to determine the encounter distances underwater.

\begin{DndTable}[header=Underwater Encounter Distance]{Xc}

Visibility & Encounter Distance\\
Clear water, Bright Light & 60 feet\\
Clear water, Dim Light & 30 feet\\
Murky water or Darkness & 10 feet\\
\end{DndTable}

\subsection{Travel}
The rules in the "\textit{Exploration}" section in the \textit{Player's Handbook} cover the basics of travel on a scale ranging from minutes to days. The tools in this section can add excitement to a longer trek.

\subsubsection{Journey Stages}
It can be helpful to break up a journey into stages, with each stage representing anything from a few hours' journey to ten days or so of travel. A journey might have only a single stage if the trip is a matter of following a clear path to a well-known destination. A journey consisting of three stages makes for a satisfying trek. For example, the characters might travel along a river to the forest's edge (stage 1), follow a trail into the heart of the woods (stage 2), and then search the woods for an ancient ruin (stage 3). A long journey might involve even more stages and occupy several game sessions.

You decide how to break up the journey, though your decision can be shaped by the characters' plan for navigating the journey. When the characters know the route they must take, the stages of the journey should correspond to the way you might give someone directions, as in the example above.

\subparagraph{Planning the Stages}
You can use the accompanying Travel Planner sheet to plan the stages of a journey. (Use multiple copies of the Travel Planner for a journey with more than three stages.)

For each stage, note where it starts and ends, the distance covered, and the predominant terrain. Choose or randomly determine the weather on that stage (see "\textit{Weather}" later in this chapter). Plan one or more challenges for each stage, such as an encounter, an obstacle, a search for something hidden, or a chance of getting lost, as described under "Journey Stage Challenges."

\subparagraph{Running the Stages}
For each stage of the journey, follow these steps in order:


\begin{description}
\item \textbf{Step 1: Set the Pace}. Have the players choose their group's travel pace for the stage: Slow, Normal, or Fast (see "\textit{Travel Pace}"). Based on the length of the stage (in miles) and the group's pace, determine how long this stage takes to complete.
\item \textbf{Step 2: Narrate the Travel}. Describe what happens as the characters complete this stage of their journey. Introduce and resolve any challenges (see "\textit{Journey Stage Challenges}").
\item \textbf{Step 3: Track Food and Water Consumption}. Each creature in the party expends the appropriate amount of food and water for the length of the stage. If the party lacks enough food or water, the characters risk dehydration and malnutrition.
\item \textbf{Step 4: Track Progress}. Track the party's progress at the end of the stage. You might mark their position on a map of the region and note the elapsed time on the Travel Planner.
\end{description}

Depending on how you planned the stages, the end of a stage might mean the characters arrive at a landmark, a waystation, or an adventure location, whether or not it's their final destination.

\begin{DndSidebar}[float=!b]{Journeys without Destinations}
Sometimes, characters travel without a clear path to follow or a clear destination in mind. In such a case, use the grid of your map (squares or hexes) to define the stages of the journey, however many miles each square or hex might be. (This style of play is sometimes called "hex crawling.")



In this kind of wilderness exploration, you can take one of two approaches to travel challenges:




\begin{description}
\item[Sandbox Approach:]
Your map of the area determines what characters find when they enter any particular hex on the map. You might have encounters or obstacles in place for every hex, or they could be spread farther apart.
\item[Random Approach:]
Use tables to randomly determine encounters or obstacles in each hex the characters enter.
\end{description}



Whichever approach you use, running a journey otherwise works the same as described in the rest of the "Travel" section.



\end{DndSidebar}

\subsubsection{Weather}
During each stage of the characters' journey, you can determine what the weather is like by rolling on the Weather table, adjusting for the terrain and season as appropriate. Roll 1d20 three times to determine the temperature, the wind, and the precipitation.

Weather has no significant game effect most of the time, but see "\textit{Environmental Effects}" in \textit{chapter 3} for the effects of extreme weather. Adding weather details to your descriptions of the characters' journey can make it more memorable.

\begin{DndTable}[header=Weather]{cX}

1d20 & Temperature\\
1–14 & Normal for the season\\
15–17 & 1d4 × 10 degrees Fahrenheit colder\\
18–20 & 1d4 × 10 degrees Fahrenheit hotter\\
\end{DndTable}


\begin{DndTable}{cXX}

1d20 & Wind & Precipitation\\
1–12 & None & None\\
13–17 & Light & Light rain or light snowfall\\
18–20 & Strong & Heavy rain or heavy snowfall\\
\end{DndTable}

\subsubsection{Travel Pace}
A group of characters can travel overland at a Normal, Fast, or Slow pace, as described in the \textit{Player's Handbook}. During any journey stage, the predominant terrain determines the characters' maximum travel pace, as shown in the Maximum Pace column of the \textit{Travel Terrain} table. Certain factors can affect a group's travel pace.

\subparagraph{Good Roads}
The presence of a good road increases the group's maximum pace by one step (from Slow to Normal or from Normal to Fast).

\subparagraph{Slower Travelers}
The group must move at a Slow pace if any group member's Speed is reduced to half or less of normal.

\subparagraph{Extended Travel}
Characters can push themselves to travel for more than 8 hours per day, at the risk of tiring. At the end of each additional hour of travel beyond 8 hours, each character must succeed on a Constitution saving throw or gain 1 Exhaustion level. The DC is 10 plus 1 for each hour past 8 hours.

\subparagraph{Special Movement}
If a party can travel at a high Speed for an extended time, as with a spell such as \textit{Wind Walk} or a magic item such as a \textit{Carpet of Flying}, translate the party's Speed into travel rates using these rules:

If the characters are flying or their special movement allows them to ignore Difficult Terrain, they can move at a Fast pace regardless of the terrain.

\subparagraph{Vehicles}
Characters traveling in a vehicle use the vehicle's speed in miles per hour (as shown in \textit{chapter 6} of the \textit{Player's Handbook}) to determine their rate of travel, and they don't choose a travel pace.

\begin{DndTable}[header=Travel Terrain]{XXcccc}

Terrain & Maximum Pace & Encounter Distance & Foraging DC & Navigation DC & Search DC\\
Arctic & Fast* & 6d6 × 10 feet & 20 & 10 & 10\\
Coastal & Normal & 2d10 × 10 feet & 10 & 5 & 15\\
Desert & Normal & 6d6 × 10 feet & 20 & 10 & 10\\
Forest & Normal & 2d8 × 10 feet & 10 & 15 & 15\\
Grassland & Fast & 6d6 × 10 feet & 15 & 5 & 15\\
Hill & Normal & 2d10 × 10 feet & 15 & 10 & 15\\
Mountain & Slow & 4d10 × 10 feet & 20 & 15 & 20\\
Swamp & Slow & 2d8 × 10 feet & 10 & 15 & 20\\
Underdark & Normal & 2d6 × 10 feet & 20 & 10 & 20\\
Urban & Normal & 2d6 × 10 feet & 20 & 15 & 15\\
Waterborne & Special† & 6d6 × 10 feet & 15 & 10 & 15\\
\end{DndTable}

\subsubsection{Narration during Travel}
Just as movies use travel montages to convey long and arduous journeys in a matter of seconds, you can use a few sentences of descriptive text to paint a picture of a journey in your players' minds before moving on. Describe the journey as vividly as you like, but keep momentum by focusing on the most notable details that reinforce the desired mood.

Visual aids can help set the scene for the characters' travels. Image searches on the internet can lead you to breathtaking landscapes (in fact, that's a good phrase to search for). You can spice up your descriptions with truly fantastical elements. For example, a forest might be home to tiny dragonets instead of birds, or its trees might be festooned with giant webs or have eerie green, glowing sap. A single fantastic element within an otherwise realistic and memorable landscape is enough.

Use the landscape to set the mood and tone for your adventure. In one forest, close-set trees might shroud all light and seem to watch the adventurers as they pass. In another, sunlight streams through the leaves above, and flower-laden vines twine up every trunk. Signs of corruption—rotting wood, foul-smelling water, and rocks covered with slimy moss—can be a signal that the adventurers are drawing close to the site of evil power that is their destination or can provide clues to the nature of the threats to be found there.

\subsection{Journey Stage Challenges}
Challenges that adventurers might face during a journey stage include the following, which are discussed in the sections that follow:

\subsubsection{Encounters with Other Creatures}
The Encounter Distance column in the \textit{Travel Terrain} table gives the range at which creatures might become aware of each other while journeying through the wilderness. When staging an encounter, consider these possibilities:


\begin{description}
\item[Ambush.]
Monsters set up an ambush along a route they expect travelers to follow.
\item[Attack from Above.]
Flying monsters swoop down to attack the characters.
\item[Distant Sighting.]
The characters and monsters spot each other from a distance in open terrain.
\item[Found by Chance.]
The characters happen upon monsters that are camping, eating, hunting, basking in the sun, walking along the same trail, or engaged in some other activity.
\item[Pursuit.]
The characters are tracking monsters, or the monsters are tracking them. The encounter begins when the two groups get close enough to interact.
\end{description}

\subsubsection{Foraging}
Characters without water and \textit{Rations} can stave off dehydration and malnutrition by gathering food and water as they travel. A foraging character makes a Wisdom (Survival) check once per journey stage (or once per day if a stage is shorter than a day). The DC is determined by the abundance of food and water in the region, as shown in the Foraging DC column of the \textit{Travel Terrain} table. If multiple characters forage, each character makes a separate check.

A foraging character finds nothing on a failed check. On a successful check, roll 1d6 and add the character's Wisdom modifier to determine how much food (in pounds) the character finds per day of the journey stage, then repeat the roll for water (in gallons).

\begin{DndSidebar}[float=!b]{Do Players Need to Track Rations?}
You might decide that tracking \textit{Rations} is unnecessary in your game. Even if the characters are in a desert, you can assume that a character with proficiency in the Survival skill can find enough food and water to sustain the party. Make sure you work that into your narration of the journey so the player feels good about choosing that skill proficiency. You can also assume the characters can load their mounts with enough Rations for their journey, or they can use magic (such as the \textit{Create Food and Water} spell) to sustain them.



On the other hand, having players track Rations seems appropriate for a more realistic campaign. Characters in such a campaign should approach a long wilderness journey as a challenge in logistics: how many pack animals do they need to carry the food for the journey, and how do they feed the animals?



As always, communicate your expectations about these rules to the players ahead of time. If you don't plan on tracking Rations, tell your players that before they spend an hour purchasing supplies for their journey.



\end{DndSidebar}

\subsubsection{Navigation}
If the characters aren't following an established path or traveling with a landmark in sight, they risk getting lost. Here are some circumstances that can cause a group to lose its way:


\begin{itemize}
\item Branching passages underground
\item Horizon-obscuring terrain, such as dense forest
\item Obscuring weather, such as heavy rain or fog
\item Traveling at night
\item Traveling at sea while unable to see the sky or any familiar land
\end{itemize}

Let the players know when the characters are at risk of getting lost, then have the characters choose one of their number to make a Wisdom (Survival) check against a DC appropriate to the terrain, as shown in the Navigation DC column of the \textit{Travel Terrain} table. Another member of the group can take the Help action to assist this check as normal.

If the check fails, the party goes off course. You decide what this looks like: they might follow the wrong branch of a river, orient themselves to the wrong mountain peak on the horizon, or get turned around in the forest. As a baseline, assume that getting lost extends the length of the current journey stage by 1d6 × 10 percent. It might also affect subsequent stages of the journey.

\subsubsection{Obstacles}
An obstacle is terrain or weather that obstructs the characters' path. Examples include a cliff, a blizzard, or a forest fire. To get past the obstacle, characters might need to backtrack and find an alternate route, or they might need to take shelter until the obstacle goes away. Let the players spend time thinking about a solution, then be generous in adjudicating whether their plan works.

In addition to the chance of a delay (adding a few hours, a day, or a couple of days to the current stage of the journey), here are some other consequences you can impose if characters fail to overcome or bypass an obstacle:


\begin{description}
\item[Combat Encounter.]
The characters encounter one or more Hostile creatures. For example, marching through a burning forest instead of circling around it might prompt an encounter with raging fire elementals.
\item[Damage.]
The characters take damage. For example, a character who tumbles over a waterfall might take Bludgeoning damage. See "\textit{Improvising Damage}" in this chapter for guidelines on determining how much damage is appropriate.
\item[Exhaustion.]
The obstacle fatigues the characters, causing them to gain Exhaustion levels. For example, marching through a blizzard instead of taking shelter might cause each character to gain 1d4 Exhaustion levels.
\item[Another Condition.]
The obstacle imposes another condition on the characters. For example, wading through a fetid swamp rather than skirting around it might impose the Poisoned condition, which lasts until removed by magic.
\end{description}

\subsubsection{Searches}
This challenge often comes up in the last stage of a journey: the characters have to find their destination, which might be an island, an old mine, an ancient ruin, a magical pool, a hag's cottage, or some other feature.

The Search DC column of the \textit{Travel Terrain} table suggests DCs for Wisdom (Perception) checks made to find something in different types of terrain. You can adjust these DCs based on the specific terrain features and the nature of what the characters are trying to find, using the guidelines for setting DCs earlier in this chapter.

\subsubsection{Tracking}
A specific instance of searching on a journey is when adventurers choose their path by following the tracks of other creatures. To track, one or more trackers must succeed on a Wisdom (Survival) check. You might require trackers to make a new check in any of the following circumstances:


\begin{description}
\item[Resting.]
The trackers resume tracking after finishing a Short Rest or Long Rest.
\item[Shifting Weather or Terrain.]
The weather or terrain changes in a way that makes tracking harder.
\item[Terrain Obstacle.]
The trail crosses a river or similar obstacle that allows no tracks.
\end{description}

The DC for the check depends on how well the ground shows signs of a creature's passage. No roll is necessary in situations where the tracks are obvious, such as following an army along a muddy road. Spotting tracks on bare rock is more challenging unless the creature being tracked leaves a distinct trail. Additionally, the passage of time often makes tracks harder to follow. In a situation where there is no trail to follow, you can rule that tracking is impossible.

Use the Search DC column of the \textit{Travel Terrain} table as a starting point for setting the DC for tracking. If you prefer, you can choose a DC based on your assessment of the difficulty—higher if days have elapsed since the creature passed, lower if the creature is leaving an obvious trail such as blood. You can also grant Advantage on the check if there's more than one set of tracks to follow or Disadvantage if the trail passes through a busy area.

On a failed check, the character loses the trail but can attempt to find it again by carefully searching the area. It takes 10 minutes to find a trail in a confined area, such as a series of caverns, or 1 hour outdoors.

\section{Running Combat}
This section builds on the \textit{combat rules} in the \textit{Player's Handbook} and offers tips for keeping the game running smoothly when a fight breaks out.

\subsection{Rolling Initiative}
Combat starts when—and only when—you say it does. Some characters have abilities that trigger on an Initiative roll; you, not the players, decide if and when Initiative is rolled. A high-level Barbarian can't just punch their Paladin friend and roll Initiative to regain expended uses of Rage.

In any situation where a character's actions initiate combat, you can give the acting character Advantage on their Initiative roll. For example, if a conversation with an NPC is cut short because the Sorcerer is convinced that NPC is a doppelganger and targets it with a \textit{Chromatic Orb} spell, everyone rolls Initiative, and the Sorcerer does so with Advantage. If the doppelganger rolls well, it might still act before the Sorcerer's spell goes off, reflecting the monster's ability to anticipate the spell.

\subsubsection{Using Initiative Scores}
You can get to the action of combat more quickly by using Initiative scores instead of rolling. You might decide to use Initiative scores just for characters, just for monsters, or for both.

\subparagraph{Initiative Scores for Characters}
A character's Initiative score is typically 10 plus all modifiers to the character's Initiative roll (including their Dexterity modifier and any special modifiers). If you want your players to use Initiative scores, have them record those scores on their character sheets, and keep your own list of those scores.

\subparagraph{Initiative Scores for Monsters}
A monster's stat block in the \textit{Monster Manual} includes its Initiative score after its Initiative bonus.

\subparagraph{Advantage and Disadvantage}
If a creature has Advantage on Initiative rolls, increase its Initiative score by 5. If it has Disadvantage on those rolls, decrease that score by 5.

\subsection{Tracking Initiative}
The following sections describe different methods for keeping track of who goes when in combat.

\subsubsection{Hidden List}
You can track Initiative on a list your players can't see using any of the following tools:


\begin{itemize}
\item Paper or a notebook behind the DM screen
\item A spreadsheet or document on a laptop or tablet
\item An app on your tablet or phone
\item Index cards for each character and each group of identical monsters, placed in Initiative order in a stack you cycle through
\end{itemize}

A hidden list allows you to track combatants who haven't been revealed yet, and you can use the list as a place to record the current Hit Points of monsters, as well as other useful notes.

If you use this approach, you tell the players when it's their characters' turn. When you call out the character whose turn is starting, consider also mentioning who's next, prompting that character's player to think ahead.

\subsubsection{Open List}
You can track Initiative on a list that is visible to the players using any of the following tools:


\begin{itemize}
\item A whiteboard on a wall or propped up nearby
\item A battle mat you use for miniatures
\item Folded index cards for each character and each group of identical monsters, placed like tents in Initiative order across the top of your DM screen
\item A virtual tabletop program you're using or a group text chat
\item Magnets, clothespins, or an accessory designed to represent the Initiative order spatially
\end{itemize}

An open list makes everyone aware of the order of play. Players know when their characters' turns are coming up so they can plan their actions in advance. An open list also lets the players know when the monsters act in the fight, although you can hold off on adding monsters to the list until they take their first turns.

\subsection{Tracking Monsters' Hit Points}
During a combat encounter, you or a player should track how much damage each monster takes. Most DMs track damage in secret so their players don't know how many Hit Points a monster has remaining.

It helps to have a system to track damage for groups of monsters. If you aren't using miniatures or other visual aids, one way to track your monsters is to assign them unique features. For example, imagine that you're running an encounter with three ogres. You might attach descriptions such as "the ogre with a big scar" and "the ogre with the helmet" to help you and your players track which monster is which. Once Initiative is rolled, jot down each ogre's Hit Points and add notes (and even a name, if you like) to differentiate each one:

\textbf{Krag (ogre w/ scar):} 68

\textbf{Thod (ogre w/ helm):} 71

\textbf{Mur (ogre smeared w/ dirt):} 59

If you use miniatures to represent monsters, one way to differentiate them is to give each one a unique miniature. If you use identical miniatures to represent multiple monsters, you can tag the miniatures with small stickers of different colors or stickers with different letters or numbers on them.

For example, in a combat encounter with three ogres, you could use three identical ogre miniatures tagged with stickers marked A, B, and C, respectively. To track the ogres' Hit Points, you can sort them by letter, then subtract damage from their Hit Points as they take it. Your records might look something like this after a few rounds of combat:

\textbf{Ogre} A: \sout{68} \sout{59} \sout{53} \sout{45} \sout{24} \sout{14} \sout{9} dead

\textbf{Ogre} B: \sout{71} \sout{62} \sout{54} 33

\textbf{Ogre} C: 59

Some DMs prefer to track how much damage a monster has taken, adding to that number as characters deal damage (instead of subtracting from the monster's Hit Points). Adding is generally easier than subtracting, and you can track damage on a visible list of Initiative (such as a whiteboard) without revealing to the players how many Hit Points the monsters have. The tracking might look like this:

\textbf{Ogre} A: \sout{9} \sout{15} \sout{23} \sout{44} \sout{54} \sout{59} dead

\textbf{Ogre} B: \sout{9} \sout{17} 38

\textbf{Ogre} C:

\subsection{Using and Tracking Conditions}
Many rules and features in the game apply conditions to creatures. You can also apply conditions on the fly when it makes sense to do so. For example, the Poisoned condition can reflect a variety of impairments, from influenza to intoxication.

You can track monsters' conditions wherever you track their Hit Points. Players should track any conditions affecting their characters. Character conditions can also be marked on index cards or a whiteboard.

You might also mark index cards or sticky notes with conditions and their effects or use tokens or some other tangible reminder. Then hand the cards, notes, or tokens to players when their characters have a condition. Putting a sticky note with a condition's rules on a player's character sheet can help that player remember the effects of the condition. You can also place tokens or colored plastic rings (the rings from soda bottle caps work well) on a creature's miniature, helping everyone remember which creatures are affected by conditions.

\subsection{Miniatures}
Often, players can rely on your descriptions to imagine where their characters are in relation to their surroundings and their enemies. Certain combat encounters, however, can benefit from having visual aids or physical props, the most common of which are miniatures and a battle grid. Miniatures are typically used in conjunction with model terrain, modular dungeon tiles, or maps drawn on large vinyl mats. Most virtual tabletops for online play simulate miniatures and grids in a digital environment.

The following sections expand on the rules in the \textit{Player's Handbook} for depicting combat \textit{using miniature figures on a grid}.

\subsubsection{Tactical Maps}
You can draw tactical maps with colored markers on an erasable vinyl mat with 1-inch squares or a similar flat surface. Preprinted poster-sized maps, maps assembled from cardboard tiles, and terrain made of sculpted plaster or resin are other options. If you're playing on a virtual tabletop, you can find abundant tactical maps in digital form online.

The most common unit for tactical maps is the 5-foot square, and maps with grids are readily available and easy to create. However, you don't have to use a grid at all. You can track distances with a tape measure, string, rulers, or pipe cleaners cut to specific lengths. Another option is a play surface covered by 1-inch hexagons (often called hexes), which makes movement more flexible while keeping the easy counting of a grid. Dungeon corridors with straight walls and right angles don't map easily onto hexes, though.

\subsubsection{Creature Size and Space}
A creature's size determines how much space it occupies on squares or hexes, as shown in the Creature Size and Space table and the accompanying diagrams.

If the miniature you use for a monster takes up an amount of space different from what's in the table, that's fine, but treat the monster as its official size for all rules. For example, you might use a miniature that has a Large base to represent a Huge giant. That giant takes up less space on the battlefield than its size suggests, but it is still Huge for the purposes of rules like grappling.

\begin{DndTable}[header=Creature Size and Space]{XXX}

Size & Space in Squares & Space in Hexes\\
Tiny & 4 per square & 4 per hex\\
Small & 1 square & 1 hex\\
Medium & 1 square & 1 hex\\
Large & 4 squares (2 by 2) & 3 hexes\\
Huge & 9 squares (3 by 3) & 7 hexes\\
Gargantuan & 16 squares (4 by 4) or more & 12 hexes or more\\
\end{DndTable}

\subsubsection{Areas of Effect}
An area of effect must be translated onto squares or hexes to determine which potential targets are in the area. If the area has a point of origin, choose an intersection of squares or hexes to be the point of origin, then follow its rules as normal. If an area of effect covers at least half a square or hex, the entire square or hex is affected.

\subsubsection{Line of Sight}
To determine whether there is line of sight between two spaces, pick a corner of one space and trace an imaginary line from that corner to any part of another space. If you can trace a line that doesn't pass through or touch an object or effect that blocks vision—such as a stone wall, a thick curtain, or a dense cloud of fog—then there is line of sight.

\subsubsection{Cover}
The accompanying diagrams illustrate cover on squares or hexes. To determine whether a target has cover against an attack or other effect, choose a corner of the attacker's space or the point of origin of an area of effect. Then trace imaginary lines from that corner to every corner of any one square the target occupies. If one or two of those lines are blocked by an obstacle (including a creature), the target has Cover. If three or four of those lines are blocked but the attack or effect can still reach the target (such as when the target is behind an arrow slit), the target has Cover.

On hexes, use the same procedure as above, drawing lines between the corners of the hexagons. The target has Half Cover if one, two, or three lines are blocked by an obstacle, and Three-Quarters Cover if four or more lines are blocked but the attack or effect can still reach the target.

\subsubsection{Diagonal Movement}
The \textit{Player's Handbook} presents a simple method for counting movement and measuring range on a grid of squares: count every square as 5 feet, even if the creature is moving or counting diagonally. While fast in play, this rule breaks the laws of geometry.

If you want more accuracy, use the following rule: the first diagonal square counts as 5 feet, but the second diagonal square counts as 10 feet. This pattern of 5 feet and then 10 feet continues whenever you're counting diagonally, even if the creature moves straight between different bits of diagonal movement. For example, a character might move 1 square diagonally (5 feet), then 3 squares straight (15 feet), and then another square diagonally (10 feet) for a total movement of 30 feet.

\subsection{Tracking Position at Long Range}
If combat erupts between two groups that are hundreds of feet away from each other, try the following techniques to keep track of who's where:


\begin{description}
\item[Note Paper.]
List all combatants on a piece of paper, and keep a running tally of each creature's distance from the party's starting point (the party starts at 0 feet). As the characters advance, increase their numbers; as the monsters advance toward the characters, decrease their numbers.
\item[Adjust the Grid Scale.]
If you're using a battle grid, take a section of that grid and use it to track position, changing the scale so that each square is 30 feet. You don't need to be precise about creatures' positions, just their distance from each other.
\item[Dice as Range Counters.]
Do away with the grid and put miniatures in their relative positions, using dice next to each miniature to show how far they've traveled. You can use percentile dice (or three d10s, with each die representing a digit in a three-digit number, if the encounter begins at a range between 100 and 1,000 feet), or use one or more d20s to show how many 5-foot or 10-foot squares the creature has advanced.
\end{description}

\subsection{Narration in Combat}
Although it's important that the players understand what's going on in terms of the rules, the game can get dull if everyone uses only "gamespeak": "That's an 18 to hit," "You hit; now roll damage," "11 points," and "OK, now we're to Initiative count 13." Instead, use the rules and your knowledge of the scene to help your narration. If 18 is barely a hit, but the 11 points of damage is a bad wound for the enemy, say: "You swing wildly, and the knight brings his shield up just a second too late. Your blade catches him along the jaw, drawing a deep gash. He recoils, bleeding badly!"

As the characters fight monsters, you can reveal information to help the players make good choices, as described in the sections that follow.

\begin{DndSidebar}[float=!b]{Awarding Heroic Inspiration}
As discussed in the \textit{Player's Handbook}, Heroic Inspiration is a reward you can give to characters when their players make the game more fun, exciting, and memorable for everyone at the table. Any player who makes the whole table erupt in laughter, cheers, or howls of surprise probably deserves Heroic Inspiration.



You can also use Heroic Inspiration to reward roleplaying, immersion in the game, and heroism. Use it to incentivize the kind of behavior you want to see in your game, such as acting in character, taking risks, thinking strategically, cooperating well, or embracing the tropes of a particular genre. Make sure your use of Heroic Inspiration is aligned with the expectations you set out at the start of your game (see "\textit{Ensuring Fun for All}" in \textit{chapter 1}).



\end{DndSidebar}

\subsubsection{Loss of Hit Points}
You can give players a sense of how well they're doing against a creature by describing, in narrative terms, how hurt the creature is. For example, if the creature is Bloodied, you might say the creature has visible wounds and appears beaten down. Such information gives the players a sense of progress and might spur them to press the attack. On the other hand, if the characters aren't damaging the creature much, let the players know the creature doesn't look hurt. That might encourage the players to change their plan.

\subsubsection{Abilities, Strengths, and Weaknesses}
Share information with the players about the characteristics of creatures they fight as those characteristics become apparent. For example, if a Wizard casts \textit{Fire Bolt} against a \textbf{Fire Elemental} (a creature that has Immunity to Fire damage), let the players know the spell doesn't seem to bother the creature at all. Players might correctly guess that a creature made of fire probably wouldn't harmed by \textit{Fire Bolt}; feel free to confirm their guesses.

\subsubsection{Actions in Combat}
When a monster takes an action in combat, the players need to have some idea what's going on both in the fictional reality of the game and in terms of the rules of the game. This means that when an enemy with a Crossbow takes the Ready action to cover the area in front of a door, the players should have a pretty good idea that if their characters move in front of that door, the enemy will shoot them. A monster's description in the \textit{Monster Manual} often explains what's happening in the world while the monster is using its special actions. The Describing Actions table has descriptions you can use to explain what's going on when a creature takes one of the common actions available to all creatures.

\begin{DndTable}[header=Describing Actions]{XX}

Action & Description\\
Dash & "Dispensing with attacks, your foe hurries across the room."\\
Disengage & "Careful not to drop its guard, your foe edges away from you."\\
Dodge & "Your foe watches you closely and tries to parry your attacks."\\
Help & "While its ally attacks, your foe darts around, causing a distraction."\\
Magic & "Your foe gestures in a deliberate manner and utters an invocation."\\
Ready & "Your foe seems to be waiting for something, ready to act."\\
\end{DndTable}

You can combine those narrative descriptions with game rules: "Dispensing with attacks, your foe hurries across the room, taking the Dash action."

\subsubsection{Monsters Casting Spells}
It's important that players can tell when their characters' opponents are casting spells, giving the characters the opportunity to cast \textit{Counterspell} or otherwise interfere with the spellcasting.

When a monster casts a spell, check the components it's using and describe its activity appropriately. If the spell has Verbal components, the monster might chant, boldly proclaim, or hiss the mystic syllables of the spell. Somatic components involve the monster moving its hands (or similar appendages) in graceful patterns, shaping them into angular positions, or thrusting them sharply forward. Finally, the monster might be holding a Spellcasting Focus or some other Material component.

Some monsters have the special ability to ignore some or all of a spell's normal components, which might prevent characters from recognizing what the monster is doing. Similarly, when monsters use magical abilities that don't involve casting spells, make sure it's clear to the players that the monster is drawing on its unique magical abilities, not casting a spell their characters could counter.

\subsection{Keeping Combat Moving}
Sometimes even the best-planned combat encounter can turn into a slog, where no one's moving and neither side is hitting or dealing much damage to the other. When that happens, here are a few techniques you can use to get things moving again or bring the encounter to a speedy close.

\subsubsection{Don't Repeat Game States}
When characters do something to change the tactical situation, don't respond by putting things back to the way they were before. For example, if a character takes the Disengage action to move away from a group of monsters, don't respond by having those same monsters chase the character. Move the monsters somewhere else.

\subsubsection{Hasten a Monster's Demise}
If a combat has gone on long enough and the characters' victory is almost certain, you can simply have the monster drop dead. The players don't ever need to know that it still had 15 Hit Points left after the characters' last attack.

\subsubsection{End Hostilities}
Most monsters can see when a fight's not going their way (or not going anywhere at all). A sapient monster might parley with the characters in an effort to get out of the situation alive. Suddenly a combat encounter turns into a social interaction as the monster and characters negotiate an end to their hostilities. A nonsapient monster might play dead to try to get the characters to stop attacking it, only to get up and run away as soon as it has the opportunity. See "\textit{Fight or Flight}" later in this chapter for more suggestions.

\subsubsection{Add a Combatant}
To add excitement to a battle, consider adding a combatant. Maybe a monstrous predator wanders onto the scene where the characters are locked in battle with another foe. Or maybe the noise of the ongoing combat attracts the attention of nearby dungeon denizens. The new combatant might attack both the characters and their foes, or it might join one side or the other. Each time one or more new combatants join the encounter, roll Initiative for them and weave them into the Initiative order.

\subsubsection{Change the Terrain}
Consider changing a battle's terrain to introduce a new element and give combatants reasons to move around. Perhaps a powerful attack or an explosive spell topples a column, shatters a wall, or breaks up the floor. Maybe a fissure opens in the floor, releasing noxious vapors, obscuring smoke, or lava. Magic could tear open the boundaries between planes of existence, unleashing raw elements or other planar energy. Or perhaps a monster's desperation causes wild magic to warp the fabric of reality. You can use the \textit{environmental effects}, \textit{hazards}, and \textit{traps} in \textit{chapter 3} to represent these effects.

\subsubsection{Change the Monster}
You can transform one monster into another to keep a fight interesting. Maybe a \textbf{worg} splits open, and a \textbf{gibbering mouther} spills out to take its place. Or a cultist is consumed in a pillar of infernal flame, and a devil erupts from the ashes. You can also alter a monster's stat block in subtle ways in the middle of combat; for example, you might decide that a monster flies into a frenzy when it's Bloodied, giving it Advantage on its attack rolls—and giving the characters Advantage on their attack rolls against it as well, speeding the fight to an end.

\subsection{Adjusting Difficulty}
Many of the same techniques that help keep combat moving can also be useful in situations where a combat encounter is either harder or easier than you anticipated and you want to adjust it. Monsters might initiate negotiations when they're winning, allowing overmatched characters a chance to surrender or retreat. One of the monsters might switch sides and help the characters defeat its kin (for noble or selfish reasons). A change of terrain can provide characters a chance to escape or give overmatched monsters an edge.

\subsection{Fight or Flight}
Few creatures fight to the death. Nearly all creatures have survival instincts that cause them to reevaluate their tactics in the face of their own destruction. Sapient creatures confronted by obviously more numerous or powerful opponents usually try to avoid battle. But brave, desperate, or devoted creatures might never retreat from a battle.

If you can't decide whether a creature is willing to fight, have it make a DC 10 Wisdom saving throw before Initiative is rolled. You can set the DC higher or lower if you like. On a failed save, the creature either flees or tries to parley with the enemy (see "\textit{Avoiding or Ending a Fight}" below). On a successful save, the creature is willing to fight. When dealing with a group of creatures, the leader makes this saving throw on behalf of the group.

When creatures that are already engaged in battle realize they're likely to lose, they usually try to exit that battle. A monster is likely to flee if either of the following is true:


\begin{itemize}
\item The monster starts its turn Bloodied and more than half its allies are dead or have the Incapacitated condition, while no one is dead or Incapacitated on the other side.
\item The monster starts its turn Bloodied and has the Frightened condition.
\end{itemize}

In those circumstances, you can decide the monster flees, or you can have it make a DC 10 Wisdom saving throw and flee or parley on a failed save. In general, if it is obvious to you that a creature is going to lose, assume it's obvious to that creature as well.

\subsubsection{Avoiding or Ending a Fight}
A creature that wishes to end or avoid a fight has two options:


\begin{description}
\item[Flight.]
The creature can retreat or run away on its turn. Select a destination for the fleeing party, such as a known place of safety (perhaps a room with a door that can be closed and barred). If the opponents pursue, you can use the "\textit{Chases}" section in \textit{chapter 3} to help adjudicate what happens.
\item[Parley.]
A parley is an attempt to settle a conflict nonviolently. The side that wishes to end or avoid combat offers to surrender or proposes some sort of exchange. If one side wants to parley in the middle of combat and the other side agrees, you can suspend the Initiative order for some interaction. If the two sides don't come to an agreement, pick up the Initiative order where it left off.
\end{description}

\section{Character Advancement}
Experience Points (XP) fuel level advancement for player characters and are most often the reward for completing combat encounters.

\subsection{Awarding XP}
Each monster has an XP value based on its Challenge Rating. When adventurers overcome one or more monsters—typically by killing, routing, capturing, or cleverly avoiding them—they divide the total XP value of the monsters evenly among themselves. If the party received substantial assistance from one or more NPCs, count those NPCs as party members when dividing up the XP, since the NPCs made the challenge easier. (See also "\textit{Nonplayer Characters}" in \textit{chapter 3}.)

\subsubsection{Noncombat Challenges}
You decide whether to award XP to characters for overcoming challenges outside combat. If the adventurers complete a tense negotiation with a baron, forge a trade agreement with a guild of surly smiths, or safely navigate the Chasm of Doom, you might decide the characters deserve XP.

As a starting point, use the rules for \textit{building combat encounters} in \textit{chapter 4} to gauge the difficulty of the challenge. Then award the characters XP as if it had been a combat encounter of the same difficulty.

\subsubsection{Milestones}
You can also award XP when characters complete significant milestones. When preparing your adventure, designate certain events or challenges as milestones, as with the following examples:


\begin{itemize}
\item Accomplishing one in a series of goals necessary to complete the adventure.
\item Discovering a hidden location or piece of information relevant to the adventure.
\item Reaching an important destination.
\end{itemize}

When awarding XP, treat a major milestone as a high-difficulty encounter and a minor milestone as a low-difficulty encounter.

\subparagraph{Other Milestone Rewards}
If you want to reward your players for their progress through an adventure with something more than XP and treasure, also give them small rewards at milestone points, such as the following:


\begin{itemize}
\item The adventurers gain the benefit of a Short Rest.
\item Characters recover a Hit Point Die or a level 1 spell slot.
\item Characters regain the use of magic items that have had their limited uses expended.
\end{itemize}

\subsection{Leveling Up}
Some DMs let characters gain the benefits of a new level as soon as the characters have the required XP, which gives the players the joy of using the new features and spells they gain immediately. Other DMs prefer to wait until the characters take a Long Rest or until the end of a session before letting characters level up, which keeps the adventure flowing smoothly and lets players pore over their new options during a lull in the action or between sessions. Do what works best for your group.

If a character levels up outside a Long Rest, the character's current Hit Points and Hit Point maximum both increase by the appropriate number for the new level, and the character gains access to additional abilities and spell slots (if appropriate) without regaining any that are already expended.

\subsubsection{Variant: Training to Gain Levels}
As a variant rule, you can require characters to spend time between adventures training or studying before they gain the benefits of a new level. This variant slows the passage of time in the game world, which can help support a more realistic or gritty tone in your campaign.

If you choose this option, after earning enough Experience Points to attain a new level, a character must train for a number of days before gaining any class features associated with the new level. You can decide whether the character can train independently or requires a trainer.

The training time required depends on the level to be gained, as shown on the Training to Gain Levels table. The training cost is for the total training time.

\begin{DndTable}[header=Training to Gain Levels]{ccc}

Level Attained & Training Time & Training Cost\\
2–4 & 10 days & 20 GP\\
5–10 & 20 days & 40 GP\\
11–16 & 30 days & 60 GP\\
17–20 & 40 days & 80 GP\\
\end{DndTable}

\subsection{Level Advancement without XP}
You can do away with XP entirely and advance characters based on how many sessions they play or when the characters accomplish significant story goals. This method of level advancement can be particularly helpful if your campaign doesn't include much combat or includes so much combat that tracking XP becomes tiresome.

\subsubsection{Session-Based Advancement}
A good rate of session-based advancement is to have characters reach level 2 after the first session of play, level 3 after another session, and level 4 after two more sessions. Then spend two or three sessions for each subsequent level. Above level 10, you can speed the rate of advancement so the characters gain a new level every one or two sessions. This assumes your sessions are about four hours long and include encounters of varying difficulty, ending with a significant milestone as described above. You can adjust the rate if you prefer significantly shorter or longer sessions and to account for how much your group accomplishes in a typical session.

\subsubsection{Story-Based Advancement}
Rather than having characters gain a level after a certain number of sessions, you can instead tie their advancement to accomplishing particular goals in the campaign. When the characters achieve those goals, they level up. Try to plan significant campaign goals so the characters gain levels at about the same rate as for session-based advancement.

\chapter{DM's Toolbox}
Whereas \textit{chapters 1} and \textit{2} teach the essentials of being a Dungeon Master, this chapter provides advice on topics that can surface as you prepare or run a D\&D game session, as well as rules for adventure elements ranging from chases and doors to traps. It also includes guidance on creating new \textit{backgrounds}, \textit{creatures}, \textit{magic items}, and \textit{spells} to amuse your players.

\section{Alignment}
As described in the \textit{Player's Handbook}, alignment is a roleplaying tool. It is a quick way to describe a creature's moral and ethical orientation. Like other elements of the game, it's meant to be a tool to serve you and your game, not a constraint or burden. Alignment can help your game in three ways: as a tool for player characters, as a descriptor of a creature's demeanor, and as a summary of an organization's ethos.

\subsection{Character Alignment}
Some common misconceptions about alignment can cause conflicts between players and DMs. The following sections can help you navigate how player characters interact with alignment.

\subsubsection{Actions Indicate Alignment}
A character might think they're good and profess to believe that senseless slaughter is wrong, but if that character repeatedly engages in senseless slaughter, the character's beliefs aren't what they profess.

Alignment doesn't limit the actions characters can take; rather, the actions they take indicate what their alignment is. It's OK to stray from the tenets of one's alignment now and then, and players can (and should) change their characters' alignments if these alignments no longer describe their characters.

\subsubsection{Good and Evil Can Cooperate}
Good and evil characters can share common goals, though they'll likely use different tactics to pursue those goals.

Imagine two characters—one Lawful Good, the other Lawful Evil—who are both dedicated to stopping monsters from preying on the people of their city. The Lawful Evil character is willing to employ methods (such as bribing or threatening potential witnesses) that the Lawful Good character isn't.

When good- and evil-aligned adventurers coexist in the same party, they're likely to have disagreements as the campaign unfolds. Many players enjoy roleplaying such conflicts, but see "\textit{Ensuring Fun for All}" in \textit{chapter 1} if you run into trouble with evil characters played in a disruptive way.

\subsubsection{Planes and Alignment}
The \textit{Outer Planes} (described in \textit{chapter 6}) are realms where alignment manifests in reality. When creatures explore the Outer Planes, they can experience those realms differently depending on their alignment.

\subsection{Monster Alignment}
Alignment can help you determine how a creature behaves in your game in two simple ways.

\subsubsection{Starting Attitude}
A creature's alignment can help you determine the creature's attitude in an encounter. A Chaotic Evil monster is likely to be Hostile [Attitude], while a Lawful Good one is more likely to have a Friendly [Attitude] attitude, ready to help those in need.

\subsubsection{Personality}
\textit{Chapter 2} of the \textit{Player's Handbook} offers a Alignment and Personality; Personality Traits by Alignment linked to alignment that can inspire you in playing an NPC or another monster in your game.

\subsection{Organization Ethos}
It can be useful to assign an alignment to an organization—including a faction, a guild, or a nation—to describe its ethos. This can help you decide how groups interact with each other.

An organization's ethos doesn't dictate the alignment of its members or even the alignment of its leaders. In fact, a stark difference between a society's ethos and the alignment of its leadership can generate interesting material for adventure. For example, imagine a Neutral Good queen ascending to the throne of a Lawful Evil empire and struggling to reform its institutions.

\section{Chases}
The rules for movement in combat don't translate to every situation. In particular, they can make a potentially thrilling chase seem dull and predictable. Faster creatures always catch up to slower ones, while creatures with the same Speed never close the distance between each other. Use the following rules to introduce random elements that make chases more exciting.

Know the capabilities of the characters in your party before you make a chase an important feature of an adventure. A character with a high Speed or the right spell (such as \textit{Dimension Door}, \textit{Fly}, or \textit{Hold Monster}) can often end a chase before it begins.

\subsection{Beginning a Chase}
A chase requires at least one quarry and at least one pursuer. Any participants not already in Initiative order must roll Initiative as the chase begins. As in combat, each participant in the chase can take one action and move on its turn.

When a chase begins, determine the starting distance between the quarry and the pursuers. Track the distance between them, and designate the pursuer closest to the quarry as the lead. The lead pursuer might change from round to round.

\subsection{Running the Chase}
Participants in the chase are strongly motivated to take the Dash action every round. Pursuers who stop to cast spells and make attacks run the risk of losing their quarry, and a quarry that doesn't take the Dash action is likely to be caught.

\subsubsection{Dashing}
A chase participant can take the Dash action a number of times equal to 3 plus its Constitution modifier (minimum of once). Each additional Dash action it takes during the chase requires the creature to succeed on a DC 10 Constitution saving throw at the end of its turn or gain 1 Exhaustion level. A participant drops out of the chase if its Speed is 0.

\subsubsection{Spells and Attacks}
A chase participant can make attacks and cast spells against other creatures within range.

Chase participants can't normally make opportunity attack against each other, since they are all assumed to be moving in the same direction at the same time. However, participants can still be the targets of Opportunity Attacks from creatures not participating in the chase. For example, adventurers who chase a thief past a gang of ruffians might provoke Opportunity Attacks from the ruffians.

\subsection{Ending a Chase}
A chase ends when one side or the other stops, when each quarry escapes, or when the pursuers are close enough to their quarry to catch it.

If neither side gives up the chase, the quarry makes a Dexterity (Stealth) check on Initiative count 0 each round, after every participant in the chase has taken its turn. If the quarry is never out of the lead pursuer's sight, the check fails automatically. Otherwise, compare the check's total to the Passive Perception scores of the pursuers. If the quarry consists of multiple creatures, they all make the check separately, so it's possible for one quarry to escape while others remain in the chase.

The quarry can gain Advantage or Disadvantage on its check based on the circumstances, as shown in the Escape Factors table.

\begin{DndTable}[header=Escape Factors]{XX}

Factor & Check Has...\\
Many things to hide behind & Advantage\\
A very crowded area & Advantage\\
Few things to hide behind & Disadvantage\\
An uncrowded area & Disadvantage\\
\end{DndTable}

Other factors might help or hinder the quarry's ability to escape, at your discretion. For example, a quarry with a \textit{Faerie Fire} spell cast on it might have Disadvantage on checks made to escape because it's much easier to spot.

If the total of the quarry's check is greater than the highest Passive Perception score of the pursuers, the quarry escapes. If not, the chase continues for another round. Escape doesn't necessarily mean the quarry has outpaced its pursuers. For example, in a city, escape might mean the quarry ducked into a crowd or slipped around a corner, leaving no clue as to where it went.

\subsubsection{Designing Your Own Chase Tables}
Unusual environments might suggest unique chase tables. A chase through the sewers of the Free City of Greyhawk or through the spiderweb-filled alleys of \textit{Menzoberranzan} (a subterranean city teeming with spiders and worshipers of \textit{Lolth}) might inspire you to create your own tables.

\subsection{Splitting Up}
Creatures being chased can split up into smaller groups. This tactic forces pursuers to either divide their forces or allow some of the quarry to escape. If a pursuit splits into several smaller chases, resolve each chase separately. You can keep all the creatures in Initiative order, but track the distances separately for each group.

\subsection{Role Reversal}
During a chase, it's possible for the pursuers to become the quarry. For example, characters chasing a thief through a marketplace might draw unwanted attention from other members of the thieves' guild. As they pursue the fleeing thief, they must also evade the thieves pursuing them. Roll Initiative for the new arrivals, and run both chases simultaneously. Alternatively, the fleeing thief might run into his accomplices. The outnumbered characters might then flee with the thieves in pursuit.

\subsection{Mapping the Chase}
When you plan a chase, draw a rough map that shows the route. Insert obstacles and complications at specific points, especially ones that require the characters to make ability checks or saving throws to avoid slowing or stopping, or use the random tables of complications in the "\textit{Chase Complications}" section to choose obstacles that occur at specific points.

Complications can be barriers to progress or opportunities for mayhem. Characters being chased through a forest by bugbears might spot a wasp nest and slow down long enough to attack the nest or throw rocks at it to enrage the wasps within, thus creating an obstacle for their pursuers.

A map of a chase can be linear or have many branches, depending on the nature of the chase. For example, a mine cart chase might have few (if any) branches, while a sewer chase might have several.

\subsection{Chase Complications}
Unexpected complications make a chase more exciting. The accompanying Urban Chase Complications table and Wilderness Chase Complications table provide several examples. Each participant in the chase rolls 1d12 at the end of its turn. Consult the appropriate table to determine whether a complication occurs. If it does, it affects the next chase participant in the Initiative order, not the participant who rolled the die.

Characters can create their own complications to shake off pursuers or slow their quarry (for example, casting the \textit{Web} spell in a narrow alleyway). Adjudicate these at your discretion.

\begin{DndTable}[header=Urban Chase Complications]{cX}

1d12 & Complication\\
1 & A cart or another large obstacle blocks your way. Make a DC 10 Dexterity saving throw to get past the obstacle. On a failed save, the obstacle counts as 10 feet of Difficult Terrain for you.\\
2 & A crowd blocks your way. Make a DC 10 Strength, Dexterity, or Charisma saving throw (your choice) to navigate through the crowd. On a failed save, the crowd counts as 10 feet of Difficult Terrain for you.\\
3 & A maze of barrels, crates, or similar obstacles stands in your way. Make a DC 10 Dexterity or Intelligence saving throw (your choice) to navigate the maze. On a failed save, the maze counts as 10 feet of Difficult Terrain for you.\\
4 & The ground is slippery with rain, spilled oil, or some other liquid. Make a DC 10 Dexterity saving throw. On a failed save, you have the Prone condition.\\
5 & You encounter a brawl in progress. Make a DC 15 Strength, Dexterity, or Charisma saving throw (your choice) to get past the brawlers unimpeded. On a failed save, you take 2d4 Bludgeoning damage, and the brawlers count as 10 feet of Difficult Terrain for you.\\
6 & You must make a sharp turn to avoid colliding with something impassable. Make a DC 10 Dexterity saving throw to navigate the turn. On a failed save, you collide with something hard and take 1d4 Bludgeoning damage.\\
7–12 & There is no complication.\\
\end{DndTable}

\begin{DndTable}[header=Wilderness Chase Complications]{cX}

1d12 & Complication\\
1 & You pass through a \textbf{Swarm of Insects} (see the \textit{Monster Manual}, with the DM choosing whichever kind of insects makes the most sense). The swarm uses one of its actions, targeting you.\\
2 & A stream or ravine blocks your path. Make a DC 10 Strength or Dexterity saving throw (your choice) to cross the impediment. On a failed save, the impediment counts as 10 feet of Difficult Terrain for you.\\
3 & Make a DC 10 Constitution saving throw. On a failed save, blowing sand, dirt, ash, snow, or pollen causes you to have the Blinded condition until the end of your turn. While you are Blinded in this way, your Speed is halved.\\
4 & A sudden drop catches you by surprise. Make a DC 10 Dexterity saving throw to navigate the impediment. On a failed save, you fall 10 feet.\\
5 & Your path takes you near a patch of \textit{razorvine} (see "\textit{Hazards}" in this chapter). Make a DC 15 Dexterity saving throw or use 10 feet of movement (your choice) to avoid the razorvine. On a failed save, you take 1d10 Slashing damage.\\
6 & A creature native to the area notices you. (The DM chooses a creature appropriate for the terrain.) Make a DC 10 Wisdom or Charisma saving throw (your choice). On a failed save, the creature joins the chase, with you as its quarry.\\
7–12 & There is no complication.\\
\end{DndTable}

\section{Creating a Background}
A character's background represents what the character did prior to becoming an adventurer. Creating a unique background or customizing an existing one from the \textit{Player's Handbook} can reflect the particular theme of your campaign or elements of your world. You can also create a background to help a player craft the story they have in mind for their character.

This section describes, step by step, how you can create backgrounds like the ones in the \textit{Player's Handbook}, tailored for your world and the heroes in it.

\subsection{1: Choose Abilities}
Choose three abilities that seem appropriate for the background:


\begin{description}
\item[Strength or Dexterity.]
These abilities are ideal for a background involving physical exertion.
\item[Constitution.]
This ability is ideal for a background that involves endurance or long hours of activity.
\item[Intelligence or Wisdom.]
One or both abilities are ideal for a background that focuses on cerebral or spiritual matters.
\item[Charisma.]
This ability is ideal for a background that involves performance or social interaction.
\end{description}

\subsection{2: Choose a Feat}
Choose one feat from the Origin category. See the \textit{Player's Handbook} for examples of \textit{Origin feats}.

\subsection{3: Choose Skill Proficiencies}
Choose two skills appropriate for the background. There needn't be a relationship between the skill proficiencies a background grants and the ability scores it increases.

\subsection{4: Choose a Tool Proficiency}
Choose one tool used in the practice of the background or often associated with it.

\subsection{5: Choose Equipment}
Assemble a package of equipment worth 50 GP (including unspent gold). Don't include Martial weapons or armor, as characters get them from their class choices.

\section{Creating a Creature}
Use the approaches and examples in the following sections to build custom creatures for your game.

\subsection{Minor Alterations}
You can change the superficial details of a creature's appearance however you like, and you can alter any of the following pieces of a monster's stat block without impacting its functionality.

\subsubsection{Size and Creature Type}
You can alter a creature's size and creature type as you please. For example, you can use an \textbf{Ogre} stat block for a human bully—just make it a Medium Humanoid instead of a Large Giant.

\subsubsection{Ability Scores}
You can usually change a creature's Intelligence, Wisdom, and Charisma scores freely unless the ability is used for spellcasting. For example, a \textbf{Black Pudding} stat block could represent a sapient alien if you raise its Intelligence and Charisma to 10 or so. It's usually best to leave its Strength, Dexterity, and Constitution scores alone, as changes to these scores can alter a monster's attack bonus, damage, Armor Class, or Hit Points, which in turn can alter its Challenge Rating.

\subsubsection{Languages}
You can change any or all of the languages the creature knows. You might want to add languages if you've made a creature sapient that wasn't before. You can also add or remove telepathy or other forms of communication.

\subsubsection{Proficiencies}
You can give a creature any skill proficiencies you want and give it Expertise if you want it to be very skilled. If you want a creature to be good at hiding, give it Expertise in the Stealth skill. If its keen senses make it an excellent tracker or otherwise adept at finding hidden foes, give it Expertise in the Perception skill. (You can also increase its Wisdom, as noted above.)

You can also swap a monster's saving throw proficiencies. If it has none, you can add one or two.

\subsubsection{Senses}
Blindsight, Darkvision, Tremorsense, and Truesight have no bearing on a creature's Challenge Rating. You can add or remove them freely.

\subsubsection{Spells}
If a stat block has spells, you can replace any of its spells with a different spell of the same level. Avoid replacing a spell that deals damage with one that doesn't and vice versa.

\subsubsection{Attacks}
You can freely change the name and flavor of an attack, as well as its damage type. For example, you can turn an ordinary \textbf{Skeleton} into an ice skeleton that deals Cold damage as it accosts characters with a blade of ice or hurls shards of ice at them.

\subsubsection{Resistances and Immunities}
If a creature doesn't have Resistance or Immunity to one or more damage types, you can give it Resistance or Immunity to one or two damage types. You can also change the damage type of its existing Resistances and Immunities.

\subsection{Traits}
You can add traits to a creature's stat block to communicate aspects of the creature's nature. See the Creature Traits list for sample traits.

You can also use traits from other stat blocks in the \textit{Monster Manual}, provided you don't add traits that alter a creature's Hit Points, confer Temporary Hit Points, or change the amount of damage the creature deals to other creatures.

\subsubsection{Creature Traits}

\begin{description}
\item[Aversion to Fire.]
If the creature takes Fire damage, it has Disadvantage on attack rolls and ability checks until the end of its next turn.
\item[Battle Ready.]
The creature has Advantage on Initiative rolls.
\item[Beast Whisperer.]
The creature can communicate with Beasts as if they shared a common language.
\item[Death Jinx.]
When the creature dies, one random creature within 10 feet of the dead creature is targeted by a \textit{Bane} spell (save DC 13), which lasts for its full duration.
\item[Dimensional Disruption.]
Disruptive energy extends from the creature in a 30-foot Emanation [Area of Effect]. Other creatures can't teleport to or from a space in that area. Any attempt to do so is wasted.
\item[Disciple of the Nine Hells.]
When the creature dies, its body disgorges a Hostile [Attitude] \textbf{Imp} in the same space.
\item[Disintegration.]
When the creature dies, its body and nonmagical possessions turn to dust. Any magic items it possessed are left behind in its space.
\item[Emissary of Juiblex.]
When the creature dies, its body disgorges a Hostile [Attitude] \textbf{Ochre Jelly} in the same space.
\item[Fey Ancestry.]
The creature has Advantage on saving throws it makes to avoid or end the Charmed condition, and magic can't put it to sleep.
\item[Forbiddance.]
The creature can't enter a residence without an invitation from one of its occupants.
\item[Gloom Shroud.]
Imperceptible energy channeled from the \textit{Shadowfell} extends from the creature in a 20-foot Emanation [Area of Effect]. Other creatures in that area have Disadvantage on Charisma checks and Charisma saving throws.
\item[Light.]
The creature sheds Bright Light in a 10-foot radius and Dim Light for an additional 10 feet. As a Bonus Action, the creature can suppress this light or cause it to return. The light winks out if the creature dies.
\item[Mimicry.]
The creature can mimic Beast sounds and Humanoid voices. A creature that hears the sounds can tell they are imitations with a successful DC 14 Wisdom (Insight) check.
\item[Poison Tolerant.]
The creature has Advantage on saving throws it makes to avoid or end the Poisoned condition.
\item[Resonant Connection.]
The creature has a supernatural connection to another creature or an object and knows the most direct route to it, provided the two are within 1 mile of each other.
\item[Siege Monster.]
The creature deals double damage to objects and structures.
\item[Slaad Host.]
When the creature dies, a Hostile [Attitude] \textbf{Slaad Tadpole} bursts from its innards in the same space.
\item[Steadfast.]
The creature has Immunity to the Frightened condition while it can see an ally within 30 feet of itself.
\item[Telepathic Bond.]
The creature is linked psychically to another creature. While both are on the same plane of existence, they can communicate telepathically with each other.
\item[Telepathic Shroud.]
The creature is immune to any effect that would sense its emotions or read its thoughts, as well as to spells from the school of Divination. As a Bonus Action, the creature can suppress this trait or reactivate it.
\item[Ventriloquism.]
Whenever the creature speaks, it can choose a point within 30 feet of itself; its voice emanates from that point.
\item[Warrior's Wrath.]
The creature has Advantage on melee attack rolls against any Bloodied creature.
\item[Wild Talent.]
Choose one cantrip; the creature can cast that cantrip without spell components, using Intelligence, Wisdom, or Charisma as the spellcasting ability.
\end{description}

\section{Creating a Magic Item}
The magic items in \textit{chapter 7} are but a few of the magic treasures that characters can discover during their adventures. You can modify these magic items or create new ones using the guidelines in this section.

Rules for characters \textit{crafting magic items} are in \textit{chapter 7}.

\subsection{Modifying a Magic Item}
You can create a new magic item by tweaking one or more existing ones. Suggestions are provided in the sections that follow.

\subsubsection{Altered Capabilities}
One capability can replace a similar one. For example, a \textit{Potion of Climbing} could become a \textit{Potion of Swimming}.

\subsubsection{Altered Form}
You can alter a magic item's form while leaving its properties intact. For example, you can turn a \textit{Ring of the Ram} into a wand or a \textit{Cloak of Protection} into some other worn object (such as a circlet) without altering the item's properties.

\subsubsection{Altered Damage Types}
An item that deals damage of one type could instead deal damage of another type. For example, a \textit{Flame Tongue} sword could deal Lightning damage instead of Fire damage.

\subsubsection{Combining Items}
You can merge the properties of two magic items of the same rarity into a single item, provided no more than one of them requires Attunement. For example, you could combine the properties of a \textit{Helm of Comprehending Languages} with those of a \textit{Helm of Telepathy} into a single helmet. This makes the item more powerful (and probably increases its rarity), but it won't break your game.

\subsubsection{Special Features and Sentience}
\textit{Chapter 7} has rules for giving magic items interesting histories, minor properties, quirks, and sentience.

\subsection{Creating a New Item}
If modifying an item doesn't quite fulfill the need, you can create one from scratch.

A magic item should either let a character do something they couldn't do before or improve the character's ability to do something they can do already. For example, a \textit{Ring of Jumping} lets its wearer jump greater distances, thus augmenting what a character can already do. A \textit{Ring of the Ram}, however, gives a character the ability to deal Force damage.

The simpler your approach, the easier it is for a character to use the item in play. Giving the item charges is fine, especially if it has several different abilities, but simply deciding that an item is always active or can be used a fixed number of times per day might be easier to manage.

\subsubsection{Power Level}
If you make an item that lets a character kill whatever they hit with it, that item will likely unbalance your game. On the other hand, an item whose benefit rarely comes into play isn't much of a reward.

Use the Magic Item Power by Rarity table as a guide to help you determine how powerful a magic item should be based on its rarity.

\begin{DndTable}[header=Magic Item Power by Rarity]{Xcc}

Rarity & Max. Spell Level & Max. Bonus\\
Common & 1 & —\\
Uncommon & 3 & +1\\
Rare & 5 & +2\\
Very Rare & 8 & +3\\
Legendary & 9 & +4\\
\end{DndTable}

\subparagraph{Maximum Spell Level}
This column of the table indicates the highest-level spell effect the item should confer, in the form of a once-per-day or similarly limited property. For example, a Common magic item might confer the benefit of a level 1 spell once per day (or just once if it's consumable). A Rare, Very Rare, or Legendary magic item might allow its possessor to cast a lower-level spell more frequently.

\subparagraph{Maximum Bonus}
If an item delivers a static bonus to AC, attack rolls, saving throws, or ability checks, this column suggests an appropriate bonus based on the item's rarity.

\subsubsection{Attunement}
Decide whether the item requires a character to be attuned to it to use its properties. Consider the following guidelines.

\subparagraph{Limit Sharing}
If having all the characters in a party pass an item around to gain its lasting benefits would be disruptive, the item should require Attunement.

\subparagraph{Limit Stacking}
If the item grants a bonus that other items also grant, it's a good idea to require Attunement so characters don't try to collect too many of those items.

\section{Creating a Spell}
When creating a new spell, use existing ones as examples. Here are some things to consider:


\begin{description}
\item[Name.]
The spell must have a unique name.
\item[Balance.]
If the spell is so good that a caster would want to use it all the time, it's probably too powerful for its level.
\item[Identity.]
Make sure the spell fits with the identity of those who can cast it. Sorcerers and Wizards don't typically cast healing spells, for example.
\item[Spell Duration, Range, and Area.]
A longer duration, greater range, or larger area can make up for a lesser effect, depending on the spell.
\item[Utility.]
Avoid spells that have very limited use, such as one that works only against Oozes. Few characters will bother to learn or prepare such a spell.
\end{description}

\subsection{Spell Damage}
For any spell that deals damage, use the Spell Damage table to determine approximately how much damage is appropriate given the spell's level. The table assumes the spell deals half as much damage on a successful saving throw or a missed attack. If your spell doesn't deal damage on a successful save, you can increase the damage by 25 percent.

You can use different damage dice than the ones in the table if the average damage is about the same. For example, you could change a cantrip's damage from 1d10 (average 5.5) to 2d4 (average 5), reducing the maximum damage and making an average result more likely.

\begin{DndTable}[header=Spell Damage]{ccc}

Spell Level & One Target & Multiple Targets\\
Cantrip & 1d10 & 1d6\\
1 & 2d10 & 2d6\\
2 & 3d10 & 3d6\\
3 & 5d10 & 6d6\\
4 & 6d10 & 7d6\\
5 & 7d10 & 8d6\\
6 & 10d10 & 11d6\\
7 & 11d10 & 12d6\\
8 & 12d10 & 13d6\\
9 & 15d10 & 16d6\\
\end{DndTable}

\subsection{Healing Spells}
You can also use the Spell Damage table to determine how many Hit Points a healing spell restores. A cantrip shouldn't provide healing.

\section{Curses and Magical Contagions}
A curse is a magical burden that lasts for a specified time or until it is ended by some means. A magical contagion is an adverse effect of magical origin that is contagious by definition.

The following sections discuss curses and magical contagions in detail.

\subsection{Curses}
Every curse has rules that govern it. A curse typically takes one of the forms detailed below.

\subsubsection{Bestow Curse}
The simplest curses are created by the \textit{Bestow Curse} spell. The effects of such curses are limited and can be ended by the \textit{Remove Curse} spell.

\textit{Bestow Curse} provides useful benchmarks for gauging the potency of other curses. A curse that lasts for 1 minute equates to a level 3 spell, while one that lasts until dispelled equates to a level 9 spell.

\subsubsection{Cursed Creatures}
Some monsters are associated with curses, whether as part of their origins or due to their ability to spread curses—werewolves being a prime example.

You decide how a spell like \textit{Remove Curse} affects a creature with accursed origins. For example, you might decide that a mummy was created through a curse and it can be destroyed permanently only by casting \textit{Remove Curse} on its corpse.

\subsubsection{Cursed Magic Items}
Cursed magic items are created deliberately or originate as the result of supernatural events. Such items are detailed in \textit{chapter 7}.

\subsubsection{Narrative Curses}
A curse might manifest during an adventure when a creature's violation of a taboo warrants supernatural punishment, such as breaking a vow, defiling a tomb, or murdering an innocent. Such a curse can have any effects you design, or it might be a customized version of another type of curse discussed in this section.

A creature affected by such a curse should know why they're being punished and be able to learn how to end their curse, likely by symbolically righting the wrong they committed. How a spell like \textit{Remove Curse} affects a curse that's part of your adventure is up to you—the spell might merely suppress the effects of the curse for a time. Regardless, narrative curses should feel like rare, potent magic rooted in the lore of your campaign.

\subsubsection{Environmental Curses}
Some locations are so suffused with evil that anyone who lingers there is burdened with a curse. Demonic Possession is one example of an environmental curse.

\subsubsection{Demonic Possession}
Demonic Possession arises from the chaos and evil of the \textit{Abyss} and commonly besets creatures that interact with demonic objects or linger in desecrated locations, where demonic spirits await victims.

A creature that becomes the target of Demonic Possession must succeed on a DC 15 Charisma saving throw or be possessed by a bodiless demonic entity. Whenever the possessed creature rolls a 1 on a D20 Test, the demonic entity takes control of the creature and determines the creature's behavior thereafter. At the end of each of the possessed creature's later turns, the creature makes a DC 15 Charisma saving throw, regaining control of itself on a success.

After finishing a Long Rest, a creature with Demonic Possession makes a DC 15 Charisma saving throw. On a successful save, the effect ends on the creature. A \textit{Dispel Evil and Good} spell or any magic that removes a curse also ends the effect on it.

\subsection{Magical Contagions}
Alchemists, potion brewers, and areas of wild magic are credited with creating the first magical contagions. An outbreak of such a contagion can form the basis of an adventure as characters search for a cure and try to stop the contagion's spread.

\subsubsection{Rest and Recuperation}
If a creature infected with a magical contagion spends 3 days recuperating, engaging in no activities that would interrupt a Long Rest, the creature makes a DC 15 Constitution saving throw at the end of the recuperation period. On a successful save, the creature has Advantage on saving throws to fight off the magical contagion for the next 24 hours.

\subsubsection{Example Contagions}
The following examples show how magical contagions can work. Feel free to alter the saving throw DCs, effects, and other characteristics of these contagions to suit your campaign.


\begin{itemize}
\begin{multicols}{3}
\item Cackle Fever
\item Sewer Plague
\item Sight Rot
\end{multicols}

\end{itemize}

\section{Death}
Adventures involve risk, with consequences that can be as catastrophic as the death of a single character or an entire group. Given the degree to which players get attached to their characters, character death can be an emotionally charged situation. It might even be a hard limit for some players (see "\textit{Ensuring Fun for All}" in \textit{chapter 1}), so it's worth having a conversation about how to handle character death at the start of a new game.

\subsection{Death Must Be Fair}
The best way to avoid hard feelings connected to the death of a beloved character is to make sure the players know you're being fair. Keep these principles in mind:


\begin{description}
\item[Don't Cheat in the Monsters' Favor.]
Rolling dice in front of the players when a situation is especially deadly is one way to communicate that you're not cheating in the monsters' favor or singling out a single character for punishment.
\item[Don't Make It Personal.]
Don't punish a character for a player's behavior or some personal grudge. That's probably the quickest way to undermine your players' trust in you as DM and as a fair arbiter of the rules.
\item[Provide Fair Warning.]
Let characters face the consequences of their foolish actions, but make sure you give enough cues for the players to recognize self-destructive actions. You might want to ask a player, "Are you sure?" before committing a character to a potentially fatal course of action.
\item[Fair Encounters.]
Your players have to know that you're fair in designing encounters. It's fine to throw tough encounters at them and sometimes to let them face monsters they can't beat. But it's not fair if the players have no way to know they can't win the fight or have no way to escape.
\end{description}

\subsection{Scaling Lethality}
You can adjust the lethality of your campaign using the \textit{encounter-building guidelines} in \textit{chapter 4}. If your players enjoy games that test their characters to the utmost and are prepared to create new characters at a moment's notice, consider using high-difficulty encounters over and over, with little opportunity for rests between encounters, to create a more lethal adventure. Conversely, using only low-difficulty encounters is less likely to lead to character death, especially if characters have ample opportunity to rest during the adventure.

\subsubsection{Defeated, Not Dead}
If you and your players agree to avoid character death in your game, you might consider an alternative: a character who would otherwise die is instead "defeated." The following rules apply to a defeated character.

\subparagraph{Comatose}
The character has 1 Hit Point and the Unconscious condition. The character can regain Hit Points as normal, but the character remains Unconscious until they are targeted by a \textit{Greater Restoration} spell or experience a sudden awakening (see below).

\subparagraph{Sudden Awakening}
After finishing a Long Rest, the character makes a DC 20 Constitution saving throw. On a successful save, the Unconscious condition ends on the character. On a failed save, the condition persists.

\subsection{Death Scenes}
When a character is reduced to 0 Hit Points, the player sometimes has to sit out one or more rounds of combat with nothing to do but roll Death Saving Throw. One way to keep a player involved in the game is to prompt some roleplaying along with each Death Save. You might ask the player to describe a memory that surfaces in the character's mind while hovering near death. Consider these possibilities:


\begin{description}
\item[On a Successful Death Save.]
A memory that inspires hope and courage. A beloved person who would urge the character to cling to life. A thought of something to live for. A favorite childhood memory.
\item[On a Failed Death Save.]
A memory that stirs up shame or grief. A beloved person who is already dead, beckoning the character to join them. An experience of weariness or despair.
\end{description}

You can also reward a player who describes a memory or something else occupying the dying character's thoughts with Advantage on the Death Save.

When a character dies, either from failed Death Saves or from an effect that kills the character outright, consider giving the player some ownership over the character's final moments by asking what the character's last words are or how the character greets death.

\subsection{Dealing with Death}
When a character dies, consult with the players to decide what happens next. Some players are perfectly happy to make new characters, especially when they're eager to try out new options. A new party member should start at the same level as the other characters in the party and have gear of similar value.

It's also possible for dead characters to be brought back to life. The most common way is through spells such as \textit{Revivify} and \textit{Raise Dead}. It's up to you to decide how easy it is for characters to access those spells if they can't cast them. The \textit{Player's Handbook} offers suggested prices for spellcasting services.

\subsection{What If Everyone Dies?}
Misadventure can wipe out an entire group. (You'll sometimes hear players refer to this as a "total party kill" or "TPK.") Such a catastrophe doesn't have to end the whole game—rather, it presents an opportunity to take the game in a new direction. Consider these possibilities.

\subsubsection{A Fresh Start}
Everyone makes new characters, and the campaign starts anew. This might be the most drastic option, but it allows for new stories and fresh character dynamics.

\subsubsection{Divine Council}
The characters find themselves before a council of deities who are arguing about the characters' fate. The characters must convince the council to return them to life.

\subsubsection{Escape from the Underworld}
The dead characters wake up in \textit{Hades} (see \textit{chapter 6}) and must find a way to escape the grim underworld and return to the world of the living.

\subsubsection{Imprisoned}
The characters wake up in cells, kept alive and imprisoned by their foes for some purpose.

\subsubsection{Raised by Another}
A powerful individual finds the adventurers' bodies and has them raised from the dead, putting the adventurers in the debt of their rescuer. What if the adventurers wake up decades after their death, returned to life by a \textit{Resurrection} spell cast by someone who believed they had an important role to play in this future era?

\subsubsection{Rescue Mission}
The players create new, temporary characters who are tasked with retrieving the bodies of the fallen heroes, so they can be raised from the dead or given proper burials. If the dead characters have Bastions (see \textit{chapter 8}), the stand-in party could consist of hirelings from those Bastions.

\section{Doors}
Adventurers interact with doors often in a D\&D campaign. This section gives rules for most of the doors the adventurers encounter.

\subsection{Common Doors}
The Doors table provides the AC and Hit Points for common doors, which are Medium objects.

With the Utilize action, a creature can try to force open a door that is barred or locked, doing so with a successful Strength (Athletics) check. The table provides the DC of the check. For bigger doors, double or triple the Hit Points and increase the DC of the check by 5.

\begin{DndTable}[header=Doors]{Xccc}

Door & AC & HP & DC to Open\\
Glass door & 13 & 4 & 10\\
Metal door & 19 & 72 & 25\\
Stone door & 17 & 40 & 20\\
Wooden door & 15 & 18 & 15\\
\end{DndTable}

\subsubsection{Barred Door}
A barred door has no lock. A creature on the barred side of the door can take the Utilize action to lift the bar from its braces, allowing the door to be opened.

\subsubsection{Locked Door}
Characters who don't have the key to a locked door can try to pick the lock using \textit{Thieves' Tools}. The Lock Complexity table tells you how long it takes to try to pick a lock based on its complexity. At the end of that time, the character picks the lock by making a successful Dexterity (Sleight of Hand) check using Thieves' Tools. The DC is determined by the lock's quality, as shown in the Lock Quality table.

\begin{DndTable}[header=Lock Complexity]{Xc}

Complexity & Time\\
Simple & 1 action\\
Complex & 1 minute\\
\end{DndTable}

\begin{DndTable}[header=Lock Quality]{Xc}

Quality & DC to Unlock\\
Inferior & 10\\
Good & 15\\
Superior & 20\\
\end{DndTable}

\subsection{Secret Doors}
A secret door is crafted to blend into the wall that surrounds it. Sometimes faint cracks in the wall or scuff marks on the floor betray the secret door's presence. Other than the fact that it's hidden, a secret door is similar to a common door.

With the Search action, a character can search for a secret door along a 10-foot-square section of wall and make a Wisdom (Perception) check. On a successful check, the character finds any secret door hidden in that section of wall as well as the mechanism to open the door. The DC of the check depends on how well the secret door is hidden, as shown in the Secret Doors table.

You can instead call for an Intelligence (Investigation) check if the challenge involves deducing that a door is present from noticeable clues, rather than spotting those clues in the first place. See "\textit{Perception}" in \textit{chapter 2} for more advice.

\begin{DndTable}[header=Secret Doors]{Xc}

Secret Door & DC to Detect\\
Barely hidden secret door & 10\\
Standard secret door & 15\\
Well-hidden secret door & 20\\
\end{DndTable}

\subsubsection{Secret Door Etiquette}
Adventurers often fail to locate secret doors. For this reason, don't hide important treasures or locations behind secret doors unless you're comfortable with the characters not finding them, and don't risk letting your adventure grind to a halt because the only path forward is hidden behind a secret door.

\subsection{Portcullises}
Typically made of iron or wood, a portcullis blocks a passage or an archway until it is raised into the ceiling by a winch and chain. Creatures within 5 feet of a lowered portcullis can make ranged attacks or cast spells through it, and they have Cover against attacks, spells, and other effects originating from the opposite side. A portcullis can also be attacked and destroyed, using the AC and Hit Points of a metal door (if iron) or a wooden door (if wood).

Winching a portcullis up or down requires the Utilize action. If a creature can't reach the winch (usually because it's on the other side of the portcullis), lifting the portcullis requires the Utilize action and a successful Strength (Athletics) check. The DC of the check depends on the type of portcullis, as shown in the Portcullises table.

\begin{DndTable}[header=Portcullises]{Xcc}

Portcullis Size & Iron DC & Wood DC\\
Medium (8 ft. tall × 5 ft. wide) & 20 & 15\\
Large (10 ft. tall × 10 ft. wide) & 25 & 20\\
Huge (20 ft. tall × 15 ft. wide) & 30 & 25\\
\end{DndTable}

\section{Dungeons}
Some dungeons are old strongholds abandoned by the folk who built them. Others are natural caves or lairs carved out by monsters. Dungeons attract cults, groups of monsters, and reclusive creatures. Because of their varied origins and purposes, dungeons have a range of distinctive qualities. For example, a dungeon that serves as a stronghold for hobgoblin soldiers has a different mood and features than an ancient temple inhabited by the yuan-ti.

You can use the Dungeon Quirks table to add distinctive character to a dungeon you're creating or one in a published adventure. The quirks on the table reflect the characteristics of a dungeon's creator, its intended purpose, its location, or some (often catastrophic) event in its history. You can use a single quirk or combine quirks as you see fit, and roll or choose a result that inspires you.

\begin{DndTable}[header=Dungeon Quirks]{cX}

1d100 & Quirk\\
01–02 & Abandoned after internal strife devastated its population\\
03–04 & Abandoned because the site was cursed by a god or other powerful entity\\
05–06 & Abandoned by its original creators when a plague spread through the dungeon\\
07–09 & Amazingly well preserved ancient city inside a dome encased in volcanic ash, submerged underwater, or entombed in desert sands\\
10–12 & Built as a fortress guarding a mountain pass\\
13–15 & Built as a maze, either to protect treasure from intruders or as a gauntlet where prisoners were hunted by monsters\\
16–18 & Built as a stronghold but abandoned after it fell to invaders\\
19–21 & Built as a treasure vault to protect powerful magic items and great wealth\\
22–23 & Built atop a cloud\\
24–26 & Built beneath a city in catacombs or sewers\\
27–29 & Built beneath or on top of a mesa or several connected mesas\\
30–32 & Built by a religious group to serve as a temple and linked to the energy of other planes of existence\\
33–35 & Built by dwarves and decorated with enormous dwarven faces that have been defaced by its current inhabitants\\
36–38 & Built in a volcano\\
39–40 & Built in or among the branches of a tree\\
41–43 & Built to house a planar portal but abandoned when creatures or energy from the other side of the portal seeped into the dungeon\\
44–46 & Carved into a meteorite (before or after it fell to earth)\\
47–49 & Carved into a sheer cliff face\\
50–52 & Caverns carved by a beholder's disintegration eye ray, with unnaturally smooth walls and vertical shafts connecting different levels\\
53–55 & Contains something that led to the downfall of its creators or inhabitants\\
56–58 & Dug as a burrow by a monster that might still live inside\\
59–61 & Entrance concealed behind a waterfall\\
62–64 & Floating on the sea\\
65–66 & Intended as a death trap to eliminate any creature that enters, perhaps to guard a treasure or to harvest souls for a necromantic rite\\
67–69 & Intended as a tomb\\
70–72 & Long known as the site of a great miracle or another auspicious event\\
73–75 & Made by amphibious creatures (such as kuo-toa or aboleths), using water to protect the innermost reaches from air-breathing intruders\\
76–78 & Made by a powerful spellcaster (perhaps a lich) as a site for magical research and experimentation\\
79–81 & Made by giants at a vast scale\\
82–84 & Natural caverns featuring a range of strikingly beautiful rock and crystal formations\\
85–87 & On an island in an underground sea\\
88–90 & On the back of a Gargantuan creature\\
91–93 & Originally constructed as a mine but abandoned when tunnels connected to dangerous Underdark tunnels\\
94–96 & Secreted away in a demiplane or in a pocket dimension\\
97–98 & Slowly abandoned as its creators died out or migrated away\\
99–00 & Transformed by multiple events or disasters over the course of centuries\\
\end{DndTable}

\subsection{Mapping a Dungeon}
A dungeon can range in size from a few chambers to a huge complex of rooms and passages extending hundreds of feet. The adventurers' goal often lies as far from the dungeon entrance as possible, forcing characters to delve deeper underground or push farther into the heart of the complex.

A dungeon is usually mapped on a grid like graph paper, with each square on the paper representing an area of 5 feet by 5 feet. \textit{Appendix B} shows several examples. If you play with miniatures on a grid, this scale makes it easy to transfer your map to a battle grid.

\subsubsection{Mapping Principles}
As you draw your map, keep the following in mind.

\subparagraph{Asymmetry}
Asymmetrical rooms and map layouts make a dungeon interesting and unpredictable.

\subparagraph{Three-Dimensional Layout}
Stairs, ramps, lifts, platforms, ledges, balconies, pits, and other changes of elevation make a dungeon interesting and make combat encounters in those areas challenging.

\subparagraph{Multiple Pathways}
Add multiple entrances and exits—to the dungeon as a whole and to individual rooms. By offering multiple paths the characters can follow, you present meaningful decision points to the players.

\subparagraph{Wear and Tear}
If you'd like to show wear and tear caused by time or the elements, collapsed passages can be commonplace, cutting off formerly connected sections of the dungeon from each other. Past earthquakes might have opened chasms within a dungeon, splitting rooms and corridors to make interesting obstacles.

\subparagraph{Natural Features}
Many dungeons include natural features. An underground stream might run through the middle of a stronghold, causing variation in the shapes and sizes of rooms and necessitating features such as bridges and drains.

\subparagraph{Secrets}
Add secret doors and secret rooms to reward players who take the time to search for them. For each door and room, consider their original purpose: were secret doors a defense against invaders, or do denizens of the dungeon scheme to keep secrets from each other? Secrets can help you develop the story of a dungeon.

\subsection{Designing Dungeon Rooms}
Keep the following things in mind when designing a dungeon room:


\begin{description}
\item[Ceiling Support.]
Underground chambers are prone to collapse, so many rooms—particularly large ones—have arched ceilings or pillars to support the weight of the rock overhead.
\item[Decoration.]
Most sapient creatures decorate their lairs. Statues, bas-reliefs, murals, and mosaics often adorn dungeon rooms. Equally common are scrawled messages, marks, and maps left behind by others who have passed through the area. Some of these marks are simply graffiti, while others may be useful to adventurers who examine them closely.
\item[Exits.]
Creatures that can't open doors can't make a lair in a sealed room without some sort of external assistance. Strong creatures without the ability to open doors smash them down if necessary. Burrowing creatures might dig their own exits.
\end{description}

Common dungeon rooms fall into the broad categories described below.

\subsubsection{Crypts}
Although it sometimes resembles a vault, a crypt can also be a series of individual rooms, each with its own sarcophagus, or a long hall with recesses on either side to hold coffins or bodies.

Crypt builders who are worried about undead rising from the grave lock and trap crypts from the outside—making the crypts easy to get into but difficult to exit. Other builders worried about tomb robbers make their crypts difficult to get into. Some builders make both entry and exit difficult, just to be safe.

\subsubsection{Guard Posts}
Sapient, social denizens of the dungeon generally guard the entrances to their shared spaces. A guard post may just be a room with a table where bored sentries play a dice game, or it might be a pair of iron golems backed up by spellcasters hiding in balconies overhead.

When you design a guard post, decide how many guards are on duty, note their Passive Perception scores, and decide what they do when they notice intruders (see "\textit{Monster Behavior}" in \textit{chapter 4}). Some will rush headlong into a fight, while others will negotiate, sound an alarm, or flee to get help.

\subsubsection{Living Quarters}
Most creatures have a lair where they can rest, eat, and store their treasure. Living quarters commonly include beds (if the creatures sleep), possessions (both valuable and mundane), and some sort of food preparation area (anything from a well-stocked kitchen to a firepit to a hunk of rotting meat).

\subsubsection{Natural Subterranean Areas}
Built dungeons often intersect with natural caverns, grottoes, and passages that are home to subterranean creatures, strange rock formations, pools of water, molds, fungi, and bioluminescent moss.

\subsubsection{Shrines}
Any sapient creature might have some place dedicated to worship. Depending on the creature's resources and piety, such a shrine can be humble or extensive. Adventurers are likely to encounter priests, cultists, and similar creatures there, and wounded monsters might flee to a shrine to seek healing.

\subsubsection{Vaults}
A vault contains treasure and is usually sealed behind a locked or secret door. Many vaults are further protected by magic, monsters that can survive without food and water, and traps (see "\textit{Traps}" in this chapter).

\subsubsection{Work Areas}
Sapient creatures often have laboratories, workshops, libraries, forges, and studios. Because such areas tend to contain valuable equipment, their doors are often locked and sometimes even warded by \textit{Glyph of Warding} spells and similar effects.

\subsection{Dungeon Decay}
The States of Ruin table can help you determine the general conditions of a dungeon area.

\begin{DndTable}[header=States of Ruin]{cX}

1d6 & Features\\
1 & Perilous. The area is dangerously worn and prone to collapse. Any impacts or damage to the structure, including from spells and other areas of effect, have a 50 percent chance of causing a collapse.\\
2 & Crumbling. Areas within the dungeon section are choked with rubble and have a 50 percent chance of being Difficult Terrain. Cover and hiding places are plentiful.\\
3 & Neglected. One dungeon hazard—such as \textit{brown mold}, \textit{green slime}, or \textit{yellow mold} (see "\textit{Hazards}" in this chapter)—is abundant.\\
4 & Abandoned. Most of the dungeon is deserted. Dexterity (Stealth) checks have Disadvantage because any sounds stand out as unusual.\\
5 & Secure. Ability checks made to break down doors, open locks, or carry out similar activities have Disadvantage.\\
6 & Thriving. The dungeon is heavily populated. Any loud noises draw the attention of nearby creatures.\\
\end{DndTable}

\section{Environmental Effects}
Characters crossing a frigid tundra might suffer the effects of extreme cold, while a visit to a cloud giant's castle might subject characters to the effects of high altitude. The following sections provide rules for handling these and other environmental effects.


\begin{itemize}
\begin{multicols}{3}
\item Dead Magic Zone
\item Deep Water
\item Extreme Cold
\item Extreme Heat
\item Frigid Water
\item Heavy Precipitation
\item High Altitude
\item Planar Effects
\item Slippery Ice
\item Strong Wind
\item Thin Ice
\item Wild Magic Zone
\end{multicols}

\end{itemize}

\section{Fear and Mental Stress}
Due to the nature of their vocation, adventurers tend to be less susceptible to fear and mental stress than common folk. Whereas a farmer might flee in terror from a bear or an apparition, adventurers are made of sterner stuff. That said, certain creatures and game effects can terrify or fray the mind of even the most stalwart adventurer.

If you plan to use any of these rules, discuss them with your players at the start of the campaign. See "\textit{Ensuring Fun for All}" in \textit{chapter 1}.

\subsection{Fear Effects}
Whenever the characters encounter something that is supernaturally frightful, use the Frightened condition as the baseline effect. Fear effects typically require a Wisdom saving throw, with a save DC set based on how terrifying the situation is. The Sample Fear DCs table provides some examples.

\begin{DndTable}[header=Sample Fear DCs]{Xc}

Example & Save DC\\
When the characters open a sarcophagus, a harmless yet terrifying apparition appears. & 10\\
A character triggers a magical trap that creates an illusory manifestation of that character's worst fears, visible only to that character. & 15\\
A portal to the Abyss opens, revealing a nightmarish realm of torment and slaughter. & 20\\
\end{DndTable}

Typically, a Frightened creature repeats the saving throw at the end of each of its turns, ending the effect on itself on a success.

At your discretion, a Frightened creature might be subject to other effects as long as the Frightened condition lasts. Consider these examples:


\begin{itemize}
\item The Frightened creature must take the Dash action on each of its turns and uses its movement to get farther away from the source of its fear.
\item Attack rolls against the Frightened creature have Advantage.
\item The Frightened creature can do only one of the following on each of its turns: move, take an action, or take a Bonus Action.
\end{itemize}

\subsection{Mental Stress Effects}
When a character is subjected to an effect that causes intense mental stress, Psychic damage is the best way to emulate that effect.

The Sample Mental Stress Effects table provides a few examples of such effects, with suggested saving throw DCs and damage. Mental stress can usually be resisted with a successful Wisdom save, but sometimes an Intelligence or Charisma save is more appropriate. On a successful save, a character might take half as much damage instead of no damage, at your discretion.

\begin{DndTable}[header=Sample Mental Stress Effects]{Xcc}

Example & Save DC & Psychic Damage\\
A character ingests a hallucinogenic substance that distorts the character's perception of reality. & 10 & 1d6\\
A character touches a fiendish idol that tears at the character's mind, threatening to shatter it. & 15 & 3d6\\
A magical trap flings a character into the Far Realm until the end of that character's next turn. & 20 & 9d6\\
\end{DndTable}

\subsubsection{Prolonged Effects}
Exposure to mental stress can cause prolonged effects. Consider the following possibilities.

\subparagraph{Short-Term Effects}
The character has the Frightened, Incapacitated, or Stunned condition for 1d10 minutes. This condition might be accompanied by alarming behavior or hallucinations. These effects can be suppressed with the \textit{Calm Emotions} spell or removed by the \textit{Lesser Restoration} spell.

\subparagraph{Long-Term Effects}
The character has Disadvantage on some or all ability checks for 1d10 × 10 hours, stemming from an unwillingness or inability to exert a particular set of abilities. The character might feel enervated and unable to exert much Strength, for example, or become so suspicious of others that Charisma checks are more difficult. These effects can be suppressed with the Calm Emotions spell or removed by the Lesser Restoration spell.

\subparagraph{Indefinite Effects}
An indefinite effect is a long-term effect (see above) that lasts until removed by a \textit{Greater Restoration} spell. It can be suppressed by a Calm Emotions spell.

\section{Firearms and Explosives}
Renaissance-era pistols and muskets appear in the \textit{Player's Handbook}. In a campaign involving a crashed spaceship or elements of modern-day Earth, characters might find the items described here.

\subsection{Firearms}
The Firearms table provides examples of modern and futuristic firearms. If you make them available for purchase (perhaps in the fantastical marketplaces of the City of Brass), treat modern items as Rare magic items and futuristic items as Very Rare ones (see \textit{chapter 7}).

\subsubsection{Properties}
Some weapons in the Firearms table have the following properties, in addition to properties described in the \textit{Player's Handbook}.

\subsubsection{Burst Fire}
As an action, you can expend 10 pieces of a Burst Fire weapon's ammunition to spray shots in a 10-foot Cube [Area of Effect] within the weapon's normal range. Each creature in that area must succeed on a DC 15 Dexterity saving throw or take damage. Roll the weapon's damage once, and apply it to each creature that failed the save.

\subsubsection{Reload}
You can make a limited number of shots with a Reload weapon. You must then reload the weapon as an action or a Bonus Action.

\subsubsection{Ammunition}
Firearm Bullets are destroyed upon use in a modern firearm. Futuristic firearms use Energy Cells that become depleted but could possibly be recharged with the proper equipment, at your discretion. An \textit{Energy Cell} weighs 1/2 lb.

\begin{DndTable}[header=Firearms]{XXXcc}

Modern Item & Damage & Properties & Mastery & Weight\\
\textit{Martial Ranged Weapons}\\
\textit{Automatic Rifle} & 2d8 Piercing & AF (Range 80/240; Bullet), BF, RLD (30 shots), 2H & Slow & 8 lb.\\
\textit{Hunting Rifle} & 2d10 Piercing & AF (Range 80/240; Bullet), RLD (5 shots), 2H & Slow & 8 lb.\\
\textit{Revolver} & 2d8 Piercing & AF (Range 40/120; Bullet), RLD (6 shots) & Sap & 3 lb.\\
\textit{Semiautomatic Pistol} & 2d6 Piercing & AF (Range 50/150; Bullet), RLD (15 shots) & Vex & 3 lb.\\
\textit{Shotgun} & 2d8 Piercing & AF (Range 30/90; Bullet), RLD (2 shots), 2H & Push & 7 lb.\\
\end{DndTable}


\begin{DndTable}{XXXcc}

Futuristic Item & Damage & Properties & Mastery & Weight\\
\textit{Martial Ranged Weapons}\\
\textit{Antimatter Rifle} & 6d8 Necrotic & AF (Range 120/360; Energy Cell), RLD (2 shots), 2H & Sap & 10 lb.\\
\textit{Laser Pistol} & 3d6 Radiant & AF (Range 40/120; Energy Cell), RLD (50 shots) & Vex & 2 lb.\\
\textit{Laser Rifle} & 3d8 Radiant & AF (Range 100/300; Energy Cell), RLD (30 shots), 2H & Slow & 7 lb.\\
\end{DndTable}

\subsection{Explosives}
The Explosives table has examples of explosives. If no cost is given for an explosive, it can't typically be bought. If you make these explosives available for purchase, treat them as Rare magic items. Rules for explosives are given below.

\begin{DndTable}[header=Explosives]{XXc}

Item & Cost & Weight\\
\textit{Bomb} & 100 GP & 1 lb.\\
\textit{Dynamite Stick} & — & 1 lb.\\
Fragmentation Grenade & — & 1 lb.\\
\textit{Grenade Launcher} & — & 7 lb.\\
Smoke Grenade & 50 GP & 2 lb.\\
\textit{Gunpowder (keg)} & 250 GP & 20 lb.\\
\textit{Gunpowder (powder horn)} & 35 GP & 2 lb.\\
\end{DndTable}

\subsubsection{Bomb}
As an action, you can light a Bomb and throw it at a point up to 60 feet away, where it explodes. Each creature in a 5-foot-radius Sphere [Area of Effect] centered on that point makes a DC 12 Dexterity saving throw, taking 3d6 Fire damage on a failed save or half as much damage on a successful one.

\subsubsection{Dynamite Stick}
An an action, you can light a Dynamite Stick and throw it at a point up to 60 feet away, where it explodes. Each creature in a 5-foot-radius Sphere [Area of Effect] centered on that point makes a DC 12 Dexterity saving throw, taking 3d6 Force damage on a failed save or half as much damage on a successful one.

It takes 1 minute to bind two or more Dynamite Sticks together so they explode at the same time. Each stick after the first increases the damage by 1d6 (to a maximum of 10d6) and the effect's radius by 5 feet (to a maximum of 20 feet).

It takes 1 minute to rig dynamite with a longer fuse so it explodes after a longer period of time, such as 1 minute or 10 minutes.

\subsubsection{Grenades and Grenade Launchers}
As an action, you can either throw a grenade at a point up to 60 feet away or use a Grenade Launcher to propel the grenade to a point up to 1,000 feet away. The grenade explodes at that point, creating a particular effect in a 20-foot-radius Sphere [Area of Effect].

\subsubsection{Fragmentation Grenade}
Each creature in the Sphere [Area of Effect] makes a DC 15 Dexterity saving throw, taking 17 (5d6) Piercing damage on a failed save or half as much damage on a successful one.

\subsubsection{Smoke Grenade}
The area of the Sphere [Area of Effect] is Heavily Obscured by smoke for 1 minute. A strong wind (such as the \textit{Gust of Wind} spell) disperses the smoke.

\subsubsection{Gunpowder}
Setting fire to a container full of Gunpowder causes it to explode. When a container explodes, each creature in a 10-foot-radius Sphere [Area of Effect] centered on the container makes a DC 12 Dexterity saving throw, taking 10 (3d6) Fire damage (for a powder horn) or 24 (7d6) Fire damage (for a keg) on a failed save or half as much damage on a successful one.

\subsection{Alien Technology}
When adventurers find a piece of technology that isn't from their world or time period, they can deduce what it is with a successful Intelligence (Investigation) check, with the DC depending on the complexity of the item: DC 10 for a relatively simple item like a calculator or a lighter, or DC 20 for a complex item such as a computer, a chainsaw, or a hovercraft. You may require a separate Intelligence (Investigation) check to determine whether a character can activate or operate the technology; a character who has observed the item in use or has operated a similar item either has Advantage on the check or succeeds on the check automatically (your choice).

\section{Gods and Other Powers}
Different deities rule the various aspects of the cosmos and mortal life, sometimes cooperating with each other, sometimes competing. People gather in public shrines to worship gods of life and wisdom or meet in hidden places to venerate gods of deception or destruction.

\subsection{Divine Rank}
The divine beings of the multiverse are often categorized according to their relative cosmic power. Some gods who are worshiped on multiple worlds have a different rank on each world, depending on their influence there.

Greater deities are generally the oldest gods of a pantheon, responsible (at least in myth) for creating or parenting the other gods. Their provinces are major areas of nature and mortal life, such as agriculture, the sun, and death. Greater deities are ultimately beyond mortal understanding, and they're often known by different names across regions, cultures, and worlds. Having no fixed appearance or gender, they can assume whatever forms they like. Occasionally these deities manifest and perform mythic deeds among mortals.

Lesser deities are typically described in myth as the creations, children, or servitors of the greater deities. They govern narrower provinces, such as the activities of mortal life or limited aspects of the natural world. They share the fundamentally ineffable nature of greater gods, but they are more likely to manifest in mortal realms.

Quasi-deities have a divine origin, but they don't receive or answer prayers. They are still immensely powerful beings, and in theory, they could ascend to godhood if they amass enough worshipers. Quasi-deities fall into the following subcategories:

\textbf{Demigods} are divine beings with mortal origin. Some were born mortal and attained godhood, while others were born from the union of a deity and a mortal. Their mortal parentage makes demigods.

\textbf{Titans} are the creations of deities. They might be manufactured on a divine forge, born from the blood spilled by a god, or otherwise brought about through divine will or substance. Some titans, including krakens and the \textbf{tarrasque}, appear in the \textit{Monster Manual}.

\textbf{Vestiges} are deities who have lost nearly all their worshipers and are considered dead from a mortal perspective. Esoteric rituals can sometimes contact vestiges and draw on their latent power.

\subsection{Home Plane and Alignment}
Gods aren't defined by mortal conceptions of alignment, and different mortal worshipers might interpret a god's behavior and teachings through the lens of different alignments. That said, gods tend to live on the \textit{Outer Planes} that most closely match their general alignment tendencies, so it's safe to assume that the teachings of a god who resides in \textit{Pandemonium} (a plane of rampant chaos and evil), encourage behavior that is Chaotic Evil, while a god who resides in \textit{Elysium} (the plane of pure good) encourages Neutral Good behavior.

People can worship a god without obeying that god's tenets or conforming to the god's presumed alignment. People from all walks of life might participate in the annual festival of innocent mischief associated with a trickster god—even people whose alignment is generally lawful and opposed to the trickster's teachings. To stave off disease, good-hearted people might make offerings to appease the wrath of a god associated with plague. Even Cleric characters don't need to have any particular alignment to serve their gods.

\subsection{Gods and Divine Magic}
Divine magic—which includes the spells cast by Clerics, Druids, Paladins, and Rangers—is mediated through beings and forces that are categorized as divine. These can include gods but also include the primal forces of nature, the beneficent power of ancestral spirits, the sacred weight of a Paladin's oath, and impersonal principles or entities such as Fate or the order of the universe. These beings and forces grant characters the power to wield the magic of their planar domains.

For game purposes, wielding divine power isn't dependent on the gods' ongoing approval or the strength of a character's devotion. The power is a gift offered to a select few; once given, it can't be rescinded.

That said, characters' relationships with the divine forces they access to wield their magic, much like Warlocks' relationships with their patrons, are ripe for exploration. A Cleric might accompany every casting of a spell with a litany of complaints directed at the gods. The Paladin in the \textit{Player's Handbook} offers some suggestions for how a player might roleplay a situation where their Paladin has broken their oath. You can also decide how NPCs react to a character whose behavior doesn't square with the ideals implied by the Holy Symbol the character wears.

\subsection{Divine Knowledge}
The \textit{Commune} spell allows its caster to ask a deity (or an agent of the god) yes-or-no questions and receive correct information, and other spells of the Divination school have similar effects. As the \textit{Commune} spell description states, gods aren't necessarily omniscient. But they are tremendously knowledgeable, particularly with regard to their particular areas of influence. A sea god can be reasonably expected to know anything that has happened in or on a sea, for example, and a martial god knows details about wars. Gods can reliably predict the future, at least in the short term (hence their ability to answer spells such as \textit{Augury} and \textit{Divination}). And some gods might be unwilling to reveal their ignorance, choosing to give an unclear answer rather than admit that they don't know the truth.

\subsection{Divine Intervention}
In some campaigns, gods are fond of meddling in mortal affairs, and heroes sometimes call on the gods for aid beyond what divine magic ordinarily provides. The gods sometimes also send dreams, omens, or emissaries to direct mortals along a certain path. Keep these two principles in mind to guide your use of divine intervention in your campaign:


\begin{description}
\item[Don't Eliminate Character Choice.]
The gods can tell characters to do things and even threaten to punish them if they don't do things, but the gods can't control mortal actions.
\item[Don't Eliminate Risk and Danger.]
The intervention of a god should never guarantee success or victory, nor should a god's interference portend immediate defeat. Gods can act to change the balance of an encounter or offer an avenue of escape, but they count on mortal heroes to act like heroes.
\end{description}

With those principles in mind, you might have gods intervene in dire situations in one of these ways:


\begin{description}
\item[Blessings.]
A god might bestow a Blessing (see "\textit{Supernatural Gifts}" in this chapter) to help a character in need.
\item[Emissaries.]
A god might send a Celestial, a Fiend, or some other kind of emissary to aid a character with information, guidance, or combat.
\item[Miracles.]
As the simplest form of miracle, a god can produce the effect of any spell that devotees of that god might cast (typically Cleric or Druid spells). But a god's direct intervention can take any form you choose, often reflecting the god's nature.
\end{description}

\subsection{Creating Religions}
A list of gods is a good starting point, and it can be sufficient to get a campaign started. But you can add more depth to your campaign world by fleshing out more details of religious belief and practice.

\subsubsection{Myths}
Stories about the gods explore their relationships with each other, with the natural world, and with the realm of mortals. Myths might describe familial relationships among the gods, deeds of creation, past interactions with mortals, or battles between gods and other cosmic forces. Given the incomprehensible nature of the gods, these myths might not actually reveal anything about the gods, but they certainly describe people's understanding of their own place in relation to the gods.

\subsubsection{Religious Practice}
People honor multiple gods of a pantheon in different circumstances. A person might burn incense to a hearth or family deity at a kitchen altar in the morning, pray to a deity of the hunt while hunting in the afternoon, and join a communal harvest feast at the temple of an agricultural deity in the evening.

Cities and large towns can host numerous temples dedicated to individual gods important to the community, while smaller settlements might have a single shrine devoted to any gods the locals revere. Temples and shrines outside settlements often mark places where a god (or the manifestation of a god) appeared or caused a miracle. These sites can become the focus of pilgrims who travel long distances to partake in the holy power assumed to linger there.

\begin{DndSidebar}[float=!b]{Build Your Own Pantheon}
Most of the published D\&D settings described in \textit{chapter 5} have their own pantheons of gods. If you're creating your own setting, you can use the list of Greyhawk gods in \textit{chapter 5} or build your own pantheon.



A simple way to build a basic pantheon is to create one god for each of the \textit{Outer Planes} described in \textit{chapter 6}, except for the \textit{Nine Hells} (ruled by the archdevil \textit{Asmodeus}) and the \textit{Abyss} (the domain of demons). So \textit{Arborea} might be the domain of a god who is patron of the arts, celebrated at great feasts, while \textit{Gehenna's} deity might be a greedy, vengeful god worshiped by people of the same bent. If you prefer, you can also put multiple deities on the same plane, so \textit{Arcadia} might be home to twin gods who are patrons of merchants and smiths.



Alternatively, you might decide that your world has only one god (who might be viewed differently by various sects or religions), or one good god and one evil god. Or your world might be alive with spirits great and small, from lesser river spirits to the godlike spirits who inhabit great mountains. Impersonal forces and philosophies can also fill the role of gods in a campaign.



\end{DndSidebar}

\section{Hazards}
The \textit{Player's Handbook} describes common hazards that adventurers encounter, such as falling and dehydration. This section details some more unusual hazards you can add to a location to make it more challenging.

\subsection{Severity and Level}
Each hazard in this section is designated as a nuisance or as deadly for characters of certain levels. A nuisance hazard is unlikely to seriously harm characters of the indicated levels, whereas a deadly hazard can grievously damage characters of the indicated levels.

Use caution when introducing a hazard to characters of a level lower than the hazard's level range. A hazard that is a nuisance at one level range could be deadly to characters in the next-lower range.

\subsection{Example Hazards}
Hazards are presented in alphabetical order.


\begin{itemize}
\begin{multicols}{3}
\item Brown Mold
\item Fireball Fungus
\item Green Slime
\item Inferno
\item Poisonous Gas
\item Quicksand Pit
\item Razorvine
\item River Styx
\item Rockslide
\item Vicious Vine
\item Webs
\item Yellow Mold
\end{multicols}

\end{itemize}

\section{Marks of Prestige}
Sometimes the most memorable reward for adventurers is the prestige they acquire throughout a realm. Their adventures often earn them fame and power, allies and enemies, and titles the adventurers can pass on to their descendants. This section details the most common marks of prestige that adventurers might acquire during a campaign.

The best rewards in an adventure are directly related to the circumstances of the adventure. For example, if a merchant hires the characters to retrieve a family heirloom from a long-abandoned tower, the merchant might give the deed to the tower as a reward.

\subsection{Fortifications}
A fortification is a reward usually given to seasoned adventurers who demonstrate unwavering fealty to a powerful political figure or ruling body, such as a monarch, a knighthood, or a council of wizards. A fortification can be anything from a fortress in the heart of a city to a provincial keep on the borderlands. While the fortification is for the characters to govern as they see fit, the land on which it sits remains the property of the crown or local ruler. Should the characters prove disloyal or unworthy of the gift, they can be asked or forced to relinquish custody of the fortification. In that respect, the fortification is different from the characters' Bastions (described in \textit{chapter 8}). However, you can also use the gift of a fortification as a pretext for the characters acquiring their Bastions.

The individual bequeathing the fortification might offer to pay its maintenance costs for one or more months, after which the characters inherit that responsibility. The type of fortification determines its maintenance costs, as shown in the Maintenance Costs table.

\begin{DndTable}[header=Maintenance Costs]{XX}

Fortification & Cost per Day\\
Fortified outpost or watchtower & 50 GP\\
Keep or small castle & 100 GP\\
Large castle or fortress & 400 GP\\
\end{DndTable}

\subsection{Letters of Recommendation}
A benefactor might provide adventurers with a letter of recommendation rather than payment. Such a letter is usually enclosed in a handsome folio, case, or scroll tube for safe transport, and it usually bears the signature and seal of whoever wrote it.

A letter of recommendation from a person of impeccable reputation can grant adventurers access to NPCs whom they would otherwise have trouble meeting, such as a duke, duchess, viceroy, or monarch. Moreover, carrying such a recommendation on one's person establish a baseline of trust with local authorities.

A letter of recommendation is worth only as much as the reputation of the person who wrote it and offers no benefit where its writer holds no sway.

\subsection{Medals}
Although they are often fashioned from gold and other precious materials, medals have an even greater symbolic value to those who award and receive them. Medals are typically awarded by political figures for acts of heroism, and wearing a medal is usually enough to earn the respect of those who understand its significance.

Different acts of heroism can warrant different kinds of medals. The king of Breland (in the Eberron setting) might award a Royal Badge of Valor (shaped like a shield and made of ruby and electrum) to adventurers for defending Brelish citizens. The Golden Bear of Breland (a medal made of gold and shaped in a likeness of a bear's head, with gems for eyes) might be reserved for adventurers who prove their allegiance to the Brelish Crown.

A medal doesn't offer a specific in-game benefit to one who wears it, but it can affect dealings with NPCs. For example, a character who displays the Golden Bear of Breland is regarded as a hero of the people within the kingdom of Breland. Outside Breland, the medal carries far less weight, except among allies of Breland's king.

\subsection{Parcels of Land}
A parcel of land usually comes with a letter from a local ruler, affirming that the land has been granted as a reward for some service. Such land usually remains the property of the local ruler or ruling body but is lent to a character with the understanding that it can be taken away, especially if the character's loyalty is ever called into question.

Characters who receive a parcel of land are free to build on it and are expected to safeguard it. They may yield the land as part of an inheritance, but they can't sell or trade it without permission from the local ruler or ruling body. If a character already has a Bastion (see \textit{chapter 8}), the parcel of land might surround the Bastion or be close to it.

Parcels of land make fine rewards for adventurers who are looking for a place to settle or who have family or a personal investment in the region where the land is located.

\subsection{Special Favors}
A reward might be a favor the characters can call on at some future date. Special favors work best when the individual granting them is trustworthy. A Lawful Good or Lawful Neutral NPC will do whatever can be done to fulfill an obligation when the time comes, short of breaking laws. A Lawful Evil NPC does the same, but only because a deal is a deal. A Neutral Good or Neutral NPC might pay off favors to protect their reputation. A Chaotic Good NPC is more concerned about doing right by the adventurers, honoring any obligations without worrying too much about personal risk or adherence to the law.

\subsection{Special Rights}
A politically powerful person can reward characters by giving them special rights, which might be articulated in some sort of official document or proclamation. For example, characters might be granted special rights to attack pirate ships or other enemies of the crown, to lead rites or ceremonies in a community, or to negotiate on a ruler's behalf. They might receive a lifetime of free room and board from the grateful citizens of a community or gain the sworn service of local soldiers to assist them as needed.

Special rights last only as long as the legal document dictates, and such rights can be revoked if the adventurers abuse them.

\subsection{Titles}
A politically powerful figure has the ability to dispense titles. A title often comes with a parcel of land (see above). For example, a character might be awarded the title Earl of Stormriver or Countess of Dun Fjord, along with a parcel of land that includes a settlement or region of the same name. Archfey are fond of granting whimsical (and alliterative) titles, such as Chancellor of Chocolates or Grand Duke of Giggles, which might come with minor supernatural gifts (see "\textit{Supernatural Gifts}" in this chapter) rather than land.

A character can hold more than one title, and in a feudal society, those titles can be passed down to (or distributed among) one's children. A character who holds a title is expected to act in a manner befitting that title. By decree, titles can be stripped away if the character fails to meet the obligations and responsibilities that come with the title.

\subsection{Training}
A character might be offered special training. This kind of training isn't widely available and thus is highly desirable.

The character must spend 30 days with the trainer to receive a special benefit. Possible training benefits include the following:


\begin{itemize}
\item The character gains proficiency in a skill.
\item The character gains proficiency with a tool.
\item The character learns a language.
\end{itemize}

\section{Mobs}
This section can help you speed up play when resolving outcomes with large groups of monsters, also known as mobs.

\subsection{Tips}
Follow these tips to smooth a combat encounter with a large number of monsters:


\begin{description}
\item[Damage.]
Use the average damage specified in a monster's stat block.
\item[Hit Points.]
If a spell or attack reduces a monster to a handful of Hit Points, assume the monster is killed or otherwise taken out of the fight.
\item[Monster Mobs.]
Divide a large number of identical monsters into smaller mobs and spread their turns out between the characters' turns. Mobs of five to eight identical creatures work well, but don't have more mobs than there are characters.
\end{description}

\subsection{Average Results}
Whenever you would otherwise make a number of D20 Test Tests} for identical monsters, the Mob Results table can help you determine the number of successful D20 Tests the monsters get without having to roll dice. Follow these steps:


\begin{description}
\item[Step 1.]
Determine the minimum d20 roll the monsters need to succeed on the D20 Test using the following formula:
\end{description}

\textbf{Roll needed} = target number−monster's bonus


\begin{description}
\item[Step 2.]
Find the roll needed on the Mob Results table. If all the monsters have Advantage on the roll (for example, if they're attacking and have the Pack Tactics trait, or if they're making a saving throw against a spell and have the Magic Resistance trait), find the roll needed in the With Advantage column. If all the monsters have Disadvantage (for example, if they're attacking a creature protected by the \textit{Blur} spell), use the With Disadvantage column. Otherwise, use the Normal column.
\item[Step 3.]
Read across the table to find a fractional number of successes you can easily apply to the group of monsters. That's the fraction of monsters that succeed on the D20 Test.
\end{description}

\begin{DndTable}[header=Mob Results]{cccccccc}

1 & 1–4 & 1 & 4/4 & 5/5 & 6/6 & 8/8 & 10/10\\
2 & 5–6 & — & 4/4 & 5/5 & 6/6 & 8/8 & 10/10\\
3 & 7–8 & 2 & 4/4 & 5/5 & 5/6 & 7/8 & 9/10\\
4 & 9 & — & 3/4 & 4/5 & 5/6 & 7/8 & 9/10\\
5 & 10 & 3 & 3/4 & 4/5 & 5/6 & 6/8 & 8/10\\
6 & 11 & — & 3/4 & 4/5 & 5/6 & 6/8 & 8/10\\
7 & 12 & 4 & 3/4 & 4/5 & 4/6 & 6/8 & 7/10\\
8 & 13 & 5 & 3/4 & 3/5 & 4/6 & 5/8 & 7/10\\
9 & 14 & — & 2/4 & 3/5 & 4/6 & 5/8 & 6/10\\
10 & — & 6 & 2/4 & 3/5 & 3/6 & 4/8 & 6/10\\
11 & 15 & 7 & 2/4 & 3/5 & 3/6 & 4/8 & 5/10\\
12 & 16 & — & 2/4 & 2/5 & 3/6 & 4/8 & 5/10\\
13 & — & 8 & 2/4 & 2/5 & 2/6 & 3/8 & 4/10\\
14 & 17 & 9 & 1/4 & 2/5 & 2/6 & 3/8 & 4/10\\
15 & 18 & 10 & 1/4 & 2/5 & 2/6 & 2/8 & 3/10\\
16 & — & 11 & 1/4 & 1/5 & 2/6 & 2/8 & 3/10\\
17 & 19 & 12 & 1/4 & 1/5 & 1/6 & 2/8 & 2/10\\
18 & — & 13 & 1/4 & 1/5 & 1/6 & 1/8 & 2/10\\
19 & 20 & 14–15 & 0 & 1/5 & 1/6 & 1/8 & 1/10\\
20 & — & 16–17 & 0 & 0 & 0 & 0 & 1/10\\
\end{DndTable}

\subsection{Adjudicating Areas of Effect}
When the characters are fighting a large number of monsters, it's not always practical to use miniatures on a battle grid or some other visual aid. So how do you determine how many monsters are affected by the Wizard's \textit{Fireball} spell or some other area of effect?

The Targets in Area of Effect table offers a guideline. To use the table, find the column for the shape of the area, then read down until you find its size. Then check the rightmost column to see about how many creatures are caught in the area. If you imagine that the targets are spread out, decrease the number by 1d3. If they're bunched up, you can increase the number by 1d3. Of course, an area can't encompass more creatures than are present in an encounter.

Your judgment always outweighs these guidelines, and it's fine to err on the side of affecting more creatures. For example, if eight zombies are crowded around a Fighter when the Bard centers a \textit{Shatter} spell on the Fighter's space, the spell's area should definitely engulf all eight zombies, even though according to the table, a 10-foot-radius Sphere [Area of Effect] includes only three creatures.

\begin{DndTable}[header=Targets in Area of Effect]{XXXXc}

10-foot & 5- to 10-foot & 5-foot-radius & — & 1\\
15- to 20-foot & 15-foot & — & 30-foot-long, 5-foot-wide & 2\\
25-foot & — & 10-foot-radius & 30-foot-long, 10-foot-wide or 60-foot-long, 5-foot-wide & 3\\
— & 20-foot & — & 90- or 100-foot-long, 5-foot-wide & 4\\
30-foot & — & — & 60-foot-long, 10-foot-wide or 120-foot-long, 5-foot-wide & 5\\
35-foot & 25-foot & 15-foot-radius & — & 6\\
40-foot & 30-foot & — & 90- or 100-foot-long, 10-foot-wide & 8\\
45-foot & — & — & — & 9\\
50-foot & 35-foot & 20-foot-radius & 120-foot-long, 10-foot-wide & 10\\
55-foot & 40-foot & — & — & 12\\
60-foot & 45-foot & 25-foot-radius & — & 16\\
— & 50-foot & 30-foot-radius & — & 20\\
\end{DndTable}

\subsection{Examples}
The following scenario shows examples of how you as the DM can apply the guidelines described in the rest of this section.

Eight Zombies surround and attack a Fighter. The zombies' attack bonus is +3, and the Fighter's AC is 18, so the roll needed is 15 (18−3). Finding 15 in the "Normal" column and reading across to the "Out of 8" column, the DM gets a result of 2/8—two of the zombies hit. Using the zombies' average damage (4 Bludgeoning damage), the Fighter takes 8 Bludgeoning damage.

After seeing the Fighter mauled by zombies, the Bard casts \textit{Shatter}, centering the spell on the Fighter. (The Bard trusts that the Fighter will succeed on the Constitution saving throw and survive the resulting damage.) The spell affects a 10-foot-radius Sphere [Area of Effect], and the Targets in Area of Effect table suggests that such an area should encompass three zombies. However, the DM decides that all eight zombies (and the Fighter) are affected. The zombies' Constitution saving throw bonus is +3, and the Bard's spell saving throw DC is 16, so the roll needed is 13 (16−3). Finding 13 in the "Normal" column and reading across to the "Out of 8" column, the DM gets a result of 3/8, so three of the zombies succeed on their saving throws.

Seeing a larger crowd of zombies in the distance, the Wizard casts \textit{Fireball}. The spell covers a 20-foot-radius Sphere. The Targets in Area of Effect table suggests that area covers ten zombies, but the DM rules that they're densely packed together and adds 1d3, rolling a 2. So the spell engulfs twelve zombies in its area. The zombies' Dexterity saving throw modifier is −2, and the Wizard's spell save DC is 16, so the roll needed is 18 (16−[−2]). Finding 18 in the Normal column and reading across to the Out of 6 column, the DM gets a result of 1/6. Twelve times 1/6 is 2, so two of the twelve zombies succeed on the save.

\section{Nonplayer Characters}
Nonplayer characters (NPCs) are supporting characters controlled by you, the DM. Examples include the local innkeeper, the sage who lives in the tower on the outskirts of town, and the death knight out to destroy the kingdom.

The \textit{Monster Manual} contains stat blocks you can use to represent NPCs in your game. You can add details to make them distinctive and memorable. For example, your players will have no trouble remembering the no-nonsense blacksmith with the tattoo of the black rose on her right shoulder or the badly dressed musician with the broken nose. NPCs in your game rarely need much more complexity than that.

\subsection{Detailed NPCs}
Flesh out NPCs who play prominent roles in your adventures. You can use the accompanying NPC Tracker to record information as you determine these six elements of your NPC:

\subsubsection{Name}
You'll need a name for any NPC who plays a prominent role in your campaign. You can pick a given name and a surname from any of the accompanying tables; a name can include options from different tables. If you like, you can roll 1d6 to determine which table to choose a name from, then roll 1d12 to get a name. You can also alter or combine names, pull from a book of names, or use a name inspired by a movie or book.

\begin{DndTable}[header=1: Common Names]{cXX}

1d12 & Common Given Name & Common Surname\\
1 & Adrik & Brightsun\\
2 & Alvyn & Dundragon\\
3 & Aurora & Frostbeard\\
4 & Eldeth & Garrick\\
5 & Eldon & Goodbarrel\\
6 & Farris & Greycastle\\
7 & Kathra & Ironfist\\
8 & Kellen & Jaerin\\
9 & Lily & Merryweather\\
10 & Nissa & Redthorn\\
11 & Xinli & Stormriver\\
12 & Zorra & Wren\\
\end{DndTable}

\begin{DndTable}[header=2: Guttural Names]{cXX}

1d12 & Guttural Given Name & Guttural Surname\\
1 & Abzug & Burska\\
2 & Bajok & Gruuthok\\
3 & Bharash & Hrondl\\
4 & Grovis & Jarzzok\\
5 & Gruuna & Kraltus\\
6 & Hokrun & Shamog\\
7 & Mardred & Skrangval\\
8 & Rhogar & Ungart\\
9 & Skuldark & Uuthrakt\\
10 & Thokk & Vrakir\\
11 & Urzul & Yuldra\\
12 & Varka & Zulrax\\
\end{DndTable}

\begin{DndTable}[header=3: Lyrical Names]{cXX}

1d12 & Lyrical Given Name & Lyrical Surname\\
1 & Arannis & Arvannis\\
2 & Damaia & Brawnanvil\\
3 & Darsis & Daardendrian\\
4 & Dweomer & Drachedandion\\
5 & Evabeth & Endryss\\
6 & Jhessail & Meliamne\\
7 & Keyleth & Mishann\\
8 & Netheria & Silverfrond\\
9 & Orianna & Snowmantle\\
10 & Sorcyl & Summerbreeze\\
11 & Umarion & Thunderfoot\\
12 & Velissa & Zashir\\
\end{DndTable}

\begin{DndTable}[header=4: Monosyllabic Names]{cXX}

1d12 & Monosyllabic Given Name & Monosyllabic Surname\\
1 & Chen & Dench\\
2 & Creel & Drog\\
3 & Dain & Dusk\\
4 & Dorn & Holg\\
5 & Flint & Horn\\
6 & Glim & Imsh\\
7 & Henk & Jask\\
8 & Krusk & Keth\\
9 & Nox & Ku\\
10 & Nyx & Kung\\
11 & Rukh & Mott\\
12 & Shan & Quaal\\
\end{DndTable}

\begin{DndTable}[header=5: Sinister Names]{cXX}

1d12 & Sinister Given Name & Sinister Surname\\
1 & Arachne & Doomwhisper\\
2 & Axyss & Dreadfield\\
3 & Carrion & Gallows\\
4 & Grinnus & Hellstryke\\
5 & Melkhis & Killraven\\
6 & Morthos & Nightblade\\
7 & Nadir & Norixius\\
8 & Scandal & Shadowfang\\
9 & Skellendyre & Valtar\\
10 & Thaltus & Winterspell\\
11 & Valkora & Xandros\\
12 & Vexander & Zarkynzorn\\
\end{DndTable}

\begin{DndTable}[header=6: Whimsical Names]{cXX}

1d12 & Whimsical Given Name & Whimsical Surname\\
1 & Cricket & Borogove\\
2 & Daisy & Goldjoy\\
3 & Dimble & Hoddypeak\\
4 & Ellywick & Huddle\\
5 & Erky & Jollywind\\
6 & Fiddlestyx & Oneshoe\\
7 & Fonkin & Scramblewise\\
8 & Golly & Sunnyhill\\
9 & Mimsy & Tallgrass\\
10 & Pumpkin & Timbers\\
11 & Quarrel & Underbough\\
12 & Sybilwick & Wimbly\\
\end{DndTable}

\subsubsection{Stat Block}
Choose a stat block from the \textit{Monster Manual} to represent the NPC's game statistics. You don't need to do this if you don't expect the NPC to engage in combat or use any special abilities (such as casting spells). You can customize the stat block using the guidelines under "\textit{Creating a Creature}" in this chapter to better reflect the NPC you have in mind.

\subsubsection{Alignment}
Choose the NPC's alignment, which can help you sketch the outlines of an NPC's behavior and personality. See the \textit{Player's Handbook} and "\textit{Alignment}" in this chapter for more information.

\subsubsection{Personality}
With the NPC's alignment and ability scores as a starting point, use the guidelines in the \textit{Player's Handbook} to pick a few words that describe the NPC's personality. You can choose or randomly determine one personality trait associated with each element of the NPC's alignment, or with the NPC's highest and lowest ability scores, and combine them to inspire a persona.

For example, if you find the adventurers unexpectedly arguing with a Lawful Neutral guard, you might create a cooperative but laconic guard who is happy to help the adventurers but speaks curtly, hoping to end the conversation as quickly as possible. Or, looking at the \textbf{Imp} stat block in the \textit{Player's Handbook} and noting its highest ability (Dexterity) and its lowest (Strength), you might decide that the little devil is fidgety and indirect, constantly on the move and talking in circles to get to its point

\subsubsection{Appearance}
Briefly describe the NPC's most distinctive physical features. You can start with the basics—skin, hair, and eye colors, as well as the NPC's species. The NPC Appearance table can also help you identify one or two things that stand out about the character's appearance.

\begin{DndTable}[header=NPC Appearance]{cX}

1d12 & Feature\\
1 & Distinctive jewelry\\
2 & Flamboyant, outlandish, formal, or ragged clothes\\
3 & Uses an elegant mobility device (wheelchair, brace, or cane)\\
4 & Pronounced scar\\
5 & Unusual eye color (or two different colors)\\
6 & Tattoos or piercings\\
7 & Birthmark\\
8 & Unusual hair color\\
9 & Bald, or braided beard or hair\\
10 & Distinctive nose (large, bulbous, angular, small)\\
11 & Distinctive posture (stooped or rigid)\\
12 & Exceptionally beautiful or ugly\\
\end{DndTable}

\subsubsection{Secret}
Describe a secret the NPC is trying to hide or protect. The NPC Secrets table provides several ideas.

\begin{DndTable}[header=NPC Secrets]{cX}

1d10 & Secret\\
1 & The NPC is in disguise, concealing their identity or some aspect of their appearance.\\
2 & The NPC is currently planning, executing, or covering up a crime.\\
3 & The NPC (or their family) has been threatened with harm unless the NPC does something.\\
4 & The NPC is under a magical compulsion (perhaps a \textit{Geas} spell or some kind of curse) to behave in a certain way.\\
5 & The NPC is seriously ill or in terrible pain.\\
6 & The NPC feels responsible for someone's death or ill fortune.\\
7 & The NPC is on the brink of financial ruin.\\
8 & The NPC is desperately lonely or harboring an unrequited passion.\\
9 & The NPC nurses a powerful ambition.\\
10 & The NPC is deeply dissatisfied or unhappy.\\
\end{DndTable}

\subsection{Recurring NPCs}
NPCs who keep showing up over the course of a campaign build the sense that the world of the game is a living, breathing place. Whether these NPCs are allies, patrons, friends, or villains, they can deepen players' investment in the world.

You can use different stat blocks in the \textit{Monster Manual}, perhaps with some tweaks, to reflect the same NPC at different times as they grow over the course of a campaign. For example, characters on their very first adventure might face a villain who uses the stat block of a \textbf{Mage Apprentice}, only to have that villain escape and return many adventures later to haunt them as a \textbf{Mage}. Still later, the same villain might reappear as an \textbf{Archmage}. Of course, the trick here is making sure that the villain survives from one adventure to the next, or at least coming up with a plausible way for the villain to return from death. After all, death is rarely the final word for adventurers, so it needn't be for their opponents.

\begin{DndSidebar}[float=!b]{Don't Go Overboard}
Developing nonplayer characters can be fun. That said, use your time wisely. Don't write three pages of backstory for an NPC whose interactions with the adventurers will be over in three minutes. The advice in this section is meant to help you create an interesting character quickly while providing just enough detail.



\end{DndSidebar}

\subsection{NPCs as Party Members}
NPCs might join the adventuring party because they want a share of the loot and are willing to accept an equal share of the risk, or they might follow the adventurers because of a bond of loyalty, gratitude, or love. You can delegate decisions about an NPC's actions to one of the players, especially in combat, but you can override the player's decisions to reflect the NPC's motivations.

When you choose a stat block from the \textit{Monster Manual} for an NPC party member, make sure the NPC doesn't overshadow the player characters. Use a stat block whose Challenge Rating is no higher than half the characters' level. These NPCs don't amass Experience Points and don't become more powerful.

Here are some NPC archetypes that work well as supporting characters in an adventuring party:


\begin{description}
\item[Comic Relief.]
A comic relief NPC helps lighten the mood of an adventure or game session, perhaps with an occasional display of ineptness or a gift for puns.
\item[Curmudgeon.]
A curmudgeon NPC is quick to complain humorously about the characters' terrible choices and bad planning. You can occasionally use this NPC to suggest legitimate courses of action or share insights.
\item[Dutiful Assistant.]
A dutiful assistant NPC is good at carrying equipment and looking after the party's horses and other belongings. Such an NPC might be entirely devoted to their duty, or they might be using this easily overlooked position to pursue goals of their own.
\item[Milquetoast Healer.]
Absent a healer of their own, the characters might love an NPC healer whose personality matters less than the healer's devotion to the party and ability to cast \textit{Cure Wounds} or \textit{Revivify} when needed.
\item[Walking Textbook.]
A walking textbook NPC is knowledgeable about a particular field of expertise and can be a useful source of information, but they can't be relied on to make wise decisions or hold up their end in battle.
\item[Wallflower Warrior.]
A wallflower warrior NPC is good at fading into the background, doesn't usually chat or engage unless approached, and eagerly avoids the spotlight. Their primary purpose is to give monsters another target to attack.
\end{description}

Even useful NPCs can slow down the game or overstay their welcome. Consider having NPC party members stick around for no more than a few game sessions or a single adventure before making their exit. NPCs can benefit from time away from the characters now and then.

\subsubsection{Loyalty}
Loyalty is an optional rule you can use to determine how far an NPC party member will go to protect or assist the characters (even those the NPC doesn't particularly like). An NPC party member who is abused or ignored is likely to abandon or betray the party, whereas an NPC who owes a life debt to the characters or shares their goals might fight to the death for them. You can simply decide on an NPC's loyalty, or you can track a Loyalty Score using the following rules.

\subparagraph{Loyalty Score}
An NPC's loyalty is measured on a numerical scale from 0 to 20. The NPC's maximum Loyalty Score is equal to the highest Charisma score among all adventurers in the party, and its starting Loyalty Score is half that number. If the highest Charisma score changes—perhaps a character dies or leaves the group—adjust the NPC's Loyalty Score accordingly.

\subparagraph{Tracking Loyalty}
Keep track of an NPC's Loyalty Score in secret so that the players won't know for sure whether an NPC party member is loyal or disloyal.

An NPC's Loyalty Score increases by 1d4 if other party members help the NPC achieve a personal goal. Likewise, an NPC's Loyalty Score increases by 1d4 if the NPC is treated particularly well (for example, given a magic weapon as a gift) or rescued by another party member. An NPC's Loyalty Score can never be raised above its maximum.

When other party members act in a manner that runs counter to the NPC's alignment or personality, reduce the NPC's Loyalty Score by 1d4. Reduce the NPC's Loyalty Score by 2d4 if the NPC is abused, misled, or endangered by other party members for purely selfish reasons. A Loyalty Score can never drop below 0.

\subparagraph{Meaning of Loyalty}
An NPC with a Loyalty Score of 10 or higher risks anything to help fellow party members. An NPC whose Loyalty Score is between 1 and 10 is tenuously faithful to the party. An NPC whose Loyalty Score drops to 0 no longer acts in the party's best interests. The disloyal NPC either leaves the party (attacking characters who attempt to intervene) or works in secret to bring about the party's downfall.

\subparagraph{Crew Loyalty and Mutiny}
If the characters own or operate a sailing ship or similar vessel, you can use these rules to track the loyalty of individual crew members or the ship's crew as a whole. If at least half the crew's Loyalty Scores drop to 0 during a voyage, the crew turns Hostile [Attitude] and stages a mutiny. If the ship is berthed, disloyal crew members leave the ship and never return.

\section{Poison}
Given their insidious and deadly nature, poisons are a favorite tool among assassins and evil creatures.

Poisons come in the following four types:


\begin{description}
\item[Contact.]
Contact poison can be smeared on an object and remains potent until it is touched or washed off. A creature that touches contact poison with exposed skin suffers its effects.
\item[Ingested.]
A creature must swallow an entire dose of ingested poison to suffer its effects. The dose can be delivered in food or a liquid. You may decide that a partial dose has a reduced effect, such as allowing Advantage on the saving throw or dealing only half as much damage on a failed save.
\item[Inhaled.]
Poisonous powders and gases take effect when inhaled. Blowing the powder or releasing the gas subjects creatures in a 5-foot Cube [Area of Effect] to its effect. The resulting cloud dissipates immediately afterward. Holding one's breath is ineffective against inhaled poisons, as they affect nasal membranes, tear ducts, and other parts of the body.
\item[Injury.]
Injury poison can be applied as a Bonus Action to a weapon, a piece of ammunition, or similar object. The poison remains potent until delivered through a wound or washed off. A creature that takes Piercing or Slashing damage from an object coated with the poison is exposed to its effects.
\end{description}

\subsection{Purchasing Poison}
In some settings, laws prohibit the possession and use of poison, but an illicit dealer or unscrupulous apothecary might keep a hidden stash. Characters with criminal contacts might be able to acquire poison easily. Other characters might have to make extensive inquiries and pay bribes before they acquire the poison they seek.

\subsection{Harvesting Poison}
A character can attempt to harvest poison from a venomous creature that is dead or has the Incapacitated condition. The effort takes 1d6 minutes, after which the character makes a DC 20 Intelligence (Nature) check using a \textit{Poisoner's Kit}. On a successful check, the character harvests enough poison for a single dose, and no additional poison can be harvested from that creature. On a failed check, the character is unable to extract any poison. If the character fails the check by 5 or more, the character is subjected to the creature's poison.

\subsection{Sample Poisons}
Example poisons are detailed here in alphabetical order. Each poison's description includes the suggested price for a single dose of the poison, its type (contact, ingested, inhaled, or injury), and a description of the poison's debilitating effects.


\begin{itemize}
\begin{multicols}{3}
\item \textit{Assassin's Blood}
\item \textit{Burnt Othur Fumes}
\item \textit{Carrion Crawler Mucus}
\item \textit{Essence of Ether}
\item \textit{Lolth's Sting}
\item \textit{Malice}
\item \textit{Midnight Tears}
\item \textit{Oil of Taggit}
\item \textit{Pale Tincture}
\item \textit{Purple Worm Poison}
\item \textit{Serpent Venom}
\item \textit{Torpor}
\item \textit{Truth Serum}
\item \textit{Wyvern Poison}
\end{multicols}

\end{itemize}

\section{Renown}
Renown is an optional rule you can use to track characters' standing, individually or as a party, within a particular group, such as a faction, an organization, or a community.

A character's or party's Renown Score starts at 0, then increases as characters earn favor and reputation with respect to the group. You can tie benefits to a character's renown, including ranks, titles, and access to resources.

Players track renown separately for each group their characters are associated with. For example, an adventurer might have a Renown Score of 5 with one faction and a Renown Score of 20 with another, based on the character's interaction with each group.

You can use renown over the course of an entire campaign or within a single adventure. At a campaign scale, you might set up factions or guilds that characters can join, individually or as a group, and the characters pursue ranks and rewards by gaining Renown within their organizations. At an adventure level, you might decide that the characters as a group need to earn a Renown Score of 5+ with the council before the council trusts the characters enough to share resources with them.

\subsection{Gaining Renown}
At your discretion, a character or party can increase their renown in the following ways:


\begin{description}
\item[Completing Missions.]
Advancing a group's interests increases a character's Renown Score within that group by 1. Completing a mission specifically assigned by that group or that directly benefits the group increases the character's Renown Score by 2. Hugely significant quests might grant Renown Score increases of 3 or 4 at once.
\item[Group Involvement.]
Once a character has established a Renown Score of 1+ with a group, the character can gain renown by spending time between adventures undertaking minor tasks for the group and socializing with its members. After doing so for a number of days equal to 10 times the character's current Renown Score, the character's Renown Score increases by 1.
\end{description}

\subsection{Benefits of Renown}
Use these guidelines when determining the benefits of increasing renown.

\subsubsection{Recognition}
A character who has a Renown Score of 3+ with a group is a respected member of that group. Other members of the group are Friendly toward the character by default and provide the character with lodging and food in dire circumstances.

\subsubsection{Rank}
Some groups have hierarchies that characters can ascend as they improve their Renown Scores. Other groups have positions of honor that characters can apply for if their Renown Score is high enough. Characters can earn promotions as their Renown Scores increase. You can establish certain Renown Score thresholds as prerequisites (though not necessarily the only prerequisites) for advancing in rank. You can set these thresholds however you like, creating ranks and titles for the groups in your campaign.

\subsubsection{Perks}
Earning renown within a group might come with certain benefits. A character with a Renown Score of 3+ might gain access to a reliable contact, a safe house, or a discount on adventuring gear. With a Renown Score of 10+, a character might gain access to Potions and Scrolls, the ability to call in a favor, or backup on dangerous missions. A character whose Renown Score rises to 50 might be able to call on a small army, acquire a Rare magic item, gain access to a helpful spellcaster, or assign special missions to members of lower status.

\subsection{Losing Renown}
Disagreements with members of a group aren't enough to cause a loss of renown within that group. However, serious offenses committed against the group or its members can result in a loss of renown and rank within the organization. The extent of the loss depends on the infraction and is left to your discretion. A character's Renown Score with a group can never drop below 0.

\subsection{Level-Based Renown}
If you want to use the benefits of renown without tracking Renown Scores, you can use a character's level as a shorthand for the character's Renown Score with a group, assuming the character has worked with or for that group for most of the character's career. The Level-Based Renown table shows equivalencies between Renown Score and character level.

\begin{DndTable}[header=Level-Based Renown]{cX}

Renown Score & Character Level\\
1 & 1\\
3 & 3\\
10 & 5\\
25 & 11\\
50 & 17\\
\end{DndTable}

\section{Settlements}
Your campaign world is likely to include settlements that characters can visit. The characters might even adopt one of these settlements as a home base, in or near which they can build their Bastions when they are of high enough level to do so (see \textit{chapter 8}).

The Settlements by Size table provides population ranges for villages, towns, and cities as well as the value of the most expensive item the settlement is likely to have for sale. Adjust these numbers as you wish to account for special circumstances. For example, a \textit{Potion of Healing} (which costs 50 GP) is too expensive an item to purchase in most villages, but a village that happens to have an alchemist, an herbalist, or a potion brewer might have one or more such potions for sale.

\begin{DndTable}[header=Settlements by Size]{XXX}

Settlement & Population Range & Max. GP Value\\
Village & Up to 500 & 20 GP\\
Town & 501–5,000 & 2,000 GP\\
City & 5,001 and higher & 200,000 GP\\
\end{DndTable}

\begin{DndSidebar}[float=!b]{Do I Need a Settlement Map?}
A settlement doesn't always require a map. Simply describing the settlement to your players is usually sufficient. But if it's important for the players to know where certain buildings or other locations are in the settlement, having a map is helpful.



For an example of a settlement map, see the Crossroads Village map in \textit{appendix B}.



\end{DndSidebar}

\subsection{Settlement Tables and Tracker}
The following tables allow you to flesh out details about a settlement. You can use the accompanying Settlement Tracker to record important information about a village, town, or city in your campaign.

\begin{DndTable}[header=Defining Traits]{cX}

1d20 & Trait\\
1–2 & Fortified outer wall\\
3–4 & Lots of gardens, parks, and greenery\\
5–6 & Lots of mud, filth, and litter\\
7–8 & Sprawling cemetery\\
9–10 & Lingering fog\\
11–12 & Noise and smoke from smithies and forges\\
13 & Canals and bridges\\
14 & Cliffs on one or more sides\\
15–16 & Clean streets and well-maintained buildings\\
17–18 & Ancient ruins within the settlement\\
19–20 & Impressive structure (such as a keep, temple, circle of standing stones, or ziggurat)\\
\end{DndTable}

\begin{DndTable}[header=Claims to Fame]{cX}

1d20 & Claim to Fame\\
1 & Delicious food\\
2 & Rude people\\
3 & Friendly folk\\
4 & Artists or writers\\
5 & Great hero/savior\\
6 & Flowers\\
7 & Seasonal festival\\
8 & Hauntings\\
9 & Spellcasters\\
10 & Decadence\\
11 & Piety\\
12 & Gambling\\
13 & Godlessness\\
14 & Education\\
15 & Wines\\
16 & High fashion\\
17 & Political intrigue\\
18 & Powerful guilds\\
19 & Patriotism\\
20 & Ancient ruins\\
\end{DndTable}

\begin{DndTable}[header=Current Calamities]{cX}

1d12 & Calamity\\
1 & Monsters infest the settlement.\\
2 & A key figure died; murder is suspected.\\
3 & War brews between rival guilds or gangs.\\
4 & A plague or famine sparks riots.\\
5 & Monsters attack anyone who approaches or leaves the settlement.\\
6 & Trade disputes cause economic hardship.\\
7 & A natural disaster threatens the settlement.\\
8 & A prophecy of doom has residents on edge.\\
9 & Locals are being drafted to fight in a war.\\
10 & Political or religious strife threatens violence.\\
11 & The settlement is under siege.\\
12 & Scandal threatens powerful local families.\\
\end{DndTable}

\begin{DndTable}[header=Local Leaders]{cX}

1d12 & Leader\\
1 & Respected, fair, and just leader or council\\
2 & Feared tyrant\\
3 & Coward manipulated by others\\
4 & Illegitimate leader causing civil unrest\\
5 & Powerful monster\\
6 & Mysterious, anonymous conspirators\\
7 & Contested leadership (with open fighting)\\
8 & Acrimonious council unable to make decisions\\
9 & Doltish lout\\
10 & Dying leader (with disputed succession)\\
11 & Iron-willed and respected leader or council\\
12 & Religious leader or council\\
\end{DndTable}

\begin{DndTable}[header=Tavern Names]{cXX}

1d20 & First Part & Second Part\\
1 & The Golden & Lyre\\
2 & The Silver & Dolphin\\
3 & The Beardless & Dwarf\\
4 & The Laughing & Pegasus\\
5 & The Dancing & Hut\\
6 & The Gilded & Rose\\
7 & The Stumbling & Stag\\
8 & The Wolf and & Duck\\
9 & The Fallen & Lamb\\
10 & The Leering & Demon\\
11 & The Drunken & Goat\\
12 & The Wine and & Spirit\\
13 & The Roaring & Horde\\
14 & The Frowning & Jester\\
15 & The Barrel and & Bucket\\
16 & The Thirsty & Crow\\
17 & The Wandering & Satyr\\
18 & The Barking & Dog\\
19 & The Happy & Spider\\
20 & The Witch and & Dragon\\
\end{DndTable}

\begin{DndTable}[header=Random Shops]{cX}

1d20 & Type\\
1 & Pawnshop\\
2 & Apothecary\\
3 & Grocer\\
4 & Delicatessen\\
5 & Potter\\
6 & Undertaker\\
7 & Bookstore\\
8 & Moneylender\\
9 & Armorer\\
10 & Chandler\\
11 & Smithy\\
12 & Carpenter\\
13 & Weaver\\
14 & Jeweler\\
15 & Baker\\
16 & Mapmaker\\
17 & Tailor\\
18 & Ropemaker\\
19 & Mason\\
20 & Scribe\\
\end{DndTable}

\section{Siege Equipment}
Siege equipment includes objects designed to assail castles and other walled fortifications. Most siege weapons require creatures to move them, as well as to load, aim, and fire them.


\begin{itemize}
\begin{multicols}{3}
\item Ballista
\item Cannon
\item Flamethrower Coach
\item Lightning Cannon
\item Mangonel
\item Ram
\item Siege Tower
\item Suspended Cauldron
\item Trebuchet
\end{multicols}

\end{itemize}

\section{Supernatural Gifts}
A supernatural gift is a special reward granted by a being or force of great magical power. Supernatural gifts come in two forms:


\begin{description}
\item[Blessing.]
A Blessing is usually bestowed by a god or a godlike being.
\item[Charm.]
A Charm is usually the work of a powerful spirit, a magical location, or a mythic creature.
\end{description}

Unlike a magic item, a supernatural gift isn't an object and doesn't require Attunement.

\subsection{Blessings}
A character might receive a Blessing from a deity for doing something truly momentous—an accomplishment that catches the attention of both gods and mortals. A Blessing is an appropriate reward for one of the following accomplishments:


\begin{itemize}
\item Restoring a god's most sacred shrine
\item Foiling an apocalyptic plot by a god's enemies
\item Helping a god's favored servant complete a quest
\end{itemize}

An adventurer might also receive a Blessing in advance of a perilous quest. For example, a Paladin could receive one before setting out on a quest to slay a terrifying lich that is responsible for a magical plague sweeping the land.

A character should receive a Blessing only if it is useful to that character, and some Blessings come with expectations on the part of the benefactor. A god might give a Blessing for a particular purpose, such as recovering a holy person's remains or toppling a tyrannical empire, and could revoke the Blessing if a character fails to pursue that purpose or acts counter to it.

A character retains the benefit of a Blessing forever or until it is taken away by the being who granted it.

There is no limit on the number of Blessings a character can receive, but it should be rare for a character to have more than one at a time. Moreover, a character can't benefit from multiple instances of a Blessing at the same time. For example, a character can't benefit from two instances of the \textit{Blessing of Health} at once.

You can easily create more Blessings by mimicking the properties of a Wondrous Item.


\begin{itemize}
\begin{multicols}{3}
\item Blessing of Health
\item Blessing of Magic Resistance
\item Blessing of Protection
\item Blessing of Understanding
\item Blessing of Valhalla
\item Blessing of Weapon Enhancement
\item Blessing of Wound Closure
\end{multicols}

\end{itemize}

\subsection{Charms}
Charms can be received in many different ways. For example, a Wizard who finds an eldritch secret in a dead archmage's spellbook might be infused with the magic of a Charm, as might a character who solves a sphinx's riddle or drinks from a magical fountain. Mythic creatures sometimes grace their allies with Charms, and some explorers find themselves bearing a Charm after discovering a long-lost location drenched in primeval magic.

Some Charms can be used only once; others can be used a specific number of times before vanishing. If a Charm lets a character cast a spell, the character can do so without expending a spell slot or providing any spell components. Unless otherwise stated, the spell uses its normal casting time, range, and duration; if the spell requires Concentration, the character must concentrate.

A Charm can't be removed from a creature by anything short of divine intervention or a \textit{Wish} spell. A character can't benefit from multiple instances of a Charm at the same time.

A typical Charm mimics the effects of a Potion or spell, so it is easy to create more Charms of your own.


\begin{itemize}
\begin{multicols}{3}
\item Charm of Animal Conjuring
\item Charm of Darkvision
\item Charm of Feather Falling
\item Charm of Heroism
\item Charm of Restoration
\item Charm of the Slayer
\item Charm of Vitality
\end{multicols}

\end{itemize}

\section{Traps}
Traps should be used sparingly, lest they lose their charm. A hidden pit can be a fun surprise, but too many traps in an adventure can lead players to become overly cautious, which slows down the game.

The best traps are fleeting distractions that skilled characters can overcome in a short amount of time or deadly puzzles that require quick thinking and teamwork to overcome. Traps that are undetectable and inescapable are rarely fun.

\subsection{Parts of a Trap}
The description of a trap includes the following parts after the trap's name:


\begin{description}
\item[Severity and Levels.]
A trap is designated as a nuisance or as deadly for characters of certain levels. A nuisance trap is unlikely to kill or seriously harm characters of the indicated levels, whereas a deadly trap can grievously damage characters of the indicated levels.
\item[Trigger.]
Traps are often set to go off when a creature enters an area or touches an object. Examples of triggers include stepping on a pressure plate, crossing a trip wire, turning a doorknob, or using the wrong key in a lock.
\item[Duration.]
Some traps have durations expressed in rounds, minutes, or hours. Others specify that their effects last until the trap is destroyed or dispelled. If a trap's duration is instantaneous, its effect is resolved instantly.
\end{description}

Use caution when introducing a trap to characters of a level lower than the trap's level range. A trap that is a nuisance at one level range could be deadly to characters of a lower level range.

\subsection{Example Traps}
Traps are presented in alphabetical order.


\begin{itemize}
\begin{multicols}{3}
\item Collapsing Roof
\item Falling Net
\item Fire-Casting Statue
\item Hidden Pit
\item Poisoned Darts
\item Poisoned Needle
\item Rolling Stone
\item Spiked Pit
\end{multicols}

\end{itemize}

\subsection{Building Your Own Traps}
When designing your own traps, use the Building a Trap table to determine an appropriate total amount of damage for the trap to deal based on its level and severity. If the trap also applies a condition, consider reducing the damage.

If the trap requires an attack roll or allows a saving throw, use the appropriate columns on the table to determine the attack bonus or an appropriate save DC.

\begin{DndTable}[header=Building a Trap]{ccccccc}

1–4 & 4 & 10–12 & 5 (1d10) & 8 & 13–15 & 11 (2d10)\\
5–10 & 4 & 12–14 & 11 (2d10) & 8 & 15–17 & 22 (4d10)\\
11–16 & 4 & 14–16 & 22 (4d10) & 8 & 17–19 & 55 (10d10)\\
17–20 & 4 & 16–18 & 55 (10d10) & 8 & 19–21 & 99 (18d10)\\
\end{DndTable}

\chapter{Creating Adventures}
Whether you're creating your own adventures or using published ones, this chapter helps you create fun and memorable experiences for your players.

Creating an adventure involves blending scenes of exploration, social interaction, and combat into a unified whole that meets the needs of your players and your campaign. The basic elements of good storytelling should guide you throughout this process, joining the encounters into a coherent story.

\section{Step-by-Step Adventures}
Follow these steps to create an adventure:


\begin{description}
\item[Step 1:]
Lay Out the Premise. Determine the situation or conflict that underscores the adventure. Also think about the adventure's setting and what is unique and fun about it.
\item[Step 2:]
Draw In the Players. Think about how the characters will get drawn into the situation you've established. Consider how the adventure might tie in with the characters' goals.
\item[Step 3:]
Plan Encounters. Determine the encounters or events that take the characters from the beginning of the adventure to the end.
\item[Step 4:]
Bring It to an End. How do you expect the adventure will end? Think about possible endings as well as rewards for the characters.
\end{description}

The rest of this chapter offers inspiration and advice for each of these four steps.

\begin{DndSidebar}[float=!b]{Using Published Adventures}
A published adventure includes a pregenerated scenario with the maps, NPCs, monsters, and treasures you need to run it. This allows you to focus your preparation time on plot developments that arise from the characters' actions.



You can adjust a published adventure so it better suits your campaign and appeals to your players. For example, you can replace the villain of an adventure with one the players have already encountered in your campaign, or add details from your campaign setting so the adventure involves your players' characters in ways that the adventure's designer never could have imagined.



Published adventures also provide inspiration for your own adventures. You can even take a part of an adventure and incorporate it into a different one. For example, you might use a map of a temple but repopulate it with monsters of your choice, or you might use a chase sequence as a model for a pursuit scene in your campaign.



\end{DndSidebar}

\section{Lay Out the Premise}
An adventure shares many of the features of a novel, a movie, an issue of a comic, or an episode of a TV show. Comic series and serialized TV dramas are particularly good comparisons because of the way individual adventures are limited in scope but blend together (to some degree) to create a larger narrative. If an adventure is a single episode or season of a series, a campaign is the series as a whole.

But while it's worthwhile to compare an adventure to these other forms of storytelling, remember that an adventure isn't a complete story until you play it. Your players are coauthors of the story with you, and the events of the story shouldn't be predetermined; the actions of the players' characters have to matter. For example, if a major villain shows up before the end of the adventure, the adventure should allow for the possibility that the heroes defeat that villain. Otherwise, players can feel as if they've been railroaded—set onto a course that has only one destination or outcome, no matter how hard they try to change it.

You might find it helpful to think about an adventure not as a narrative that arcs from beginning to end with little chance for deviation, but more in terms of situations that you are presenting to the characters. The adventure unfolds organically from the players' responses to the situations you present.

\begin{DndSidebar}[float=!b]{Guide Rails and Railroads}
Players need to feel like they're in control of their characters, the choices they make matter, and what they do has some effect on the outcome of the adventure and on the game world. Keep that in mind as you're planning adventures. If your adventure relies on certain events, plan for multiple ways they might come about, or be prepared for clever players to prevent those events from happening as you expect. Otherwise, your players might end up feeling railroaded.



On the flip side, players sometimes willfully disregard the adventure hooks you put in front of them and go entirely off the rails. See "\textit{Respect for the DM}" in \textit{chapter 1} and "\textit{Draw In the Players}" later in this chapter for advice about dealing with this situation.



One way to give players impactful choices is to keep multiple adventure possibilities available to them at the same time. If the characters have two or three things they can investigate or pursue, they have a meaningful choice. And if whatever threads they don't investigate turn into bigger problems, you've clearly demonstrated that their decisions matter.



\end{DndSidebar}

\subsection{Adventure Premise}
One way to start brainstorming an adventure is to imagine a situation that might pique the characters' interest. For some D\&D players, a rumor of a dungeon filled with treasure is enough of a premise to launch an adventure. A brewing war between two small nations, the death of a leader and the accession of a new one, a migration of dangerous monsters, the appearance of a comet, and the opening of a portal to another plane of existence are other situations that could lead to adventure.

A simple premise might also boil down to "a magic item that the characters want is hidden away in a dungeon." Browsing the magic items in \textit{chapter 7} can inspire you to create a simple adventure seed.

\subsection{Adventure Conflict}
A premise can be a good starting point, but before you can turn it into an adventure, it needs a conflict worthy of the heroes' attention. The conflict might be driven by a single villain or monster, a villain with lackeys, an assortment of monsters, or an evil organization. But it need not involve the forces of evil; it could be a rivalry or disagreement between two families, organizations, or nations; a looming natural (or magical) disaster; or even conflict within the adventuring party about how to pursue the characters' goals.

Given a premise of a dungeon filled with treasure, what conflict awaits the characters when they enter the dungeon? That might be as simple as "hostile monsters want to eat the characters" or "two rival factions of monsters inhabit the dungeon." It might also involve a rival or a villain hoping to plunder the dungeon first, a rumbling volcano threatening to erupt and bury the dungeon, or two rival families claiming ownership of the treasure left behind in the dungeon by their ancestors.

If you're stuck, browse through the \textit{Monster Manual} until you find a monster that inspires you.

\subsection{Adventure Situations by Level}
Use the tables in this section to inspire adventure ideas for characters of different levels, with the range of possible levels grouped into four tiers. You can roll on the tables and see if the result sparks your imagination or read the entries on the tables until you find something that grabs you.

\subsubsection{Levels 1–4: Local Heroes}
The fate of a village might depend on the abilities of fledgling adventurers. These characters navigate dangerous terrain and explore haunted crypts, where they might fight ferocious wolves, giant spiders, evil cultists, flesh-eating ghouls, and ruthless brigands.

\begin{DndTable}[header=Levels 1–4 Adventure Situations]{cX}

1d20 & Situation\\
1 & A dragon wyrmling has gathered a band of kobolds to help it amass a hoard.\\
2 & Wererats living in a city's sewers plot to take control of the governing council.\\
3 & Bandit activity signals efforts to revive an evil cult long ago driven from the region.\\
4 & A pack of gnolls is rampaging dangerously close to local farmlands.\\
5 & A rivalry between two merchant families escalates from mischief to mayhem.\\
6 & A new sinkhole has revealed a long-buried dungeon thought to hold treasure.\\
7 & Miners discovered an underground ruin and were captured by monsters living there.\\
8 & An innocent person is being framed for the crimes of a shape-shifting monster.\\
9 & Ghouls are venturing out of the catacombs at night.\\
10 & A notorious criminal hides from the law in an old ruin or abandoned mine.\\
11 & A contagion in a forest is causing spiders to grow massive and become aggressive.\\
12 & To take revenge against a village for an imagined slight, a necromancer has been animating the corpses in the village cemetery.\\
13 & An evil cult is spreading in a village. Those who oppose the cult are marked for sacrifice.\\
14 & An abandoned house on the edge of town is haunted by Undead because of a cursed item in the house.\\
15 & Creatures from the Feywild enter the world and cause mischief and misfortune among villagers and their livestock.\\
16 & A hag's curse is making animals unusually aggressive.\\
17 & Bullies have appointed themselves the village militia and are extorting money and food from villagers.\\
18 & After a local fisher pulls a grotesque statue from the sea, aquatic monsters start attacking the waterfront at night.\\
19 & The ruins on the hill near the village lie under a curse, so people don't go there—except a scholar who wants to study the ruins.\\
20 & A new captain has taken charge of a band of pirates or bandits and started raiding more frequently.\\
\end{DndTable}

\subsubsection{Levels 5–10: Heroes of the Realm}
At this tier, characters undertake adventures that might determine the fate of a region. These adventurers venture into fearsome wilds and ancient ruins, where they confront giants, hydras, golems, devils, demons, and mind flayers. They might also face a young dragon that has just established a lair.

\begin{DndTable}[header=Levels 5–10 Adventure Situations]{cX}

1d20 & Situation\\
1 & A group of cultists has summoned a demon to wreak havoc in the city.\\
2 & A rebel lures monsters to the cause with the promise of looting the king's treasury.\\
3 & An evil Artifact has transformed a forest into a dismal swamp full of horrific monsters.\\
4 & An Aberration living in the Underdark sends minions to capture people from the surface to turn those people into new minions.\\
5 & A monster (perhaps a devil, slaad, or hag) is impersonating a prominent noble to throw the realm into civil war.\\
6 & A master thief plans to steal royal regalia.\\
7 & A golem intended to serve as a protector has gone berserk and captured its creator.\\
8 & A conspiracy of spies, assassins, and necromancers schemes to overthrow a ruler.\\
9 & After establishing a lair, a young dragon is trying to earn the fear and respect of other creatures living nearby.\\
10 & The approach of a lone giant alarms the people of a town, but the giant is simply looking for a place to live in peace.\\
11 & An enormous monster on display in a menagerie breaks free and goes on a rampage.\\
12 & A coven of hags steals cherished memories from travelers.\\
13 & A villain seeks powerful magic in an ancient ruin, hoping to use it to conquer the region.\\
14 & A scheming aristocrat hosts a masquerade ball, which many guests see as an opportunity to advance their own agendas. At least one shape-shifting monster also attends.\\
15 & A ship carrying a valuable treasure or an evil Artifact sinks in a storm or monster attack.\\
16 & A natural disaster was actually caused by magic gone awry or a cult's villainous plans.\\
17 & A secretive cult uses spies to heighten tensions between two rival nations, hoping to provoke a war that will weaken both.\\
18 & Rebels or forces of an enemy nation have kidnapped an important noble.\\
19 & The descendants of a displaced people want to reclaim their ancestral city, which is now inhabited by monsters.\\
20 & A renowned group of adventurers never returned from an expedition to a famous ruin.\\
\end{DndTable}

\subsubsection{Levels 11–16: Masters of the Realm}
The fate of a nation or even the world depends on the characters at this tier. These adventurers explore uncharted regions and delve into forgotten dungeons, where the characters confront terrible schemers of the Lower Planes, cunning rakshasas and beholders, and hungry purple worms. They might encounter and even defeat a powerful adult dragon. At this tier, they broker peace between nations or lead them into war, and their formidable reputations attract the attention of powerful foes.

\begin{DndTable}[header=Levels 11–16 Adventure Situations]{cX}

1d12 & Situation\\
1 & A portal to the Abyss opens in a cursed location and spews demons into the world.\\
2 & A band of hunting giants has driven its prey—enormous beasts—into pastureland.\\
3 & An adult dragon's lair is transforming an expanse into an environment inhospitable to the other creatures living there.\\
4 & A long-lost journal describes an incredible journey to a hidden subterranean realm full of magical wonders.\\
5 & Cultists hope to persuade a dragon to undergo the rite that will transform it into a dracolich.\\
6 & The ruler of the realm is sending an emissary to a hostile neighbor to negotiate a truce, and the emissary needs protection.\\
7 & A castle or city has been drawn into another plane of existence.\\
8 & A storm tears across the land, with a mysterious flying citadel in the eye of the storm.\\
9 & Two parts of a magic item are in the hands of bitter enemies; the third piece is lost.\\
10 & Evil cultists gather from around the world to summon a monstrous god or alien entity.\\
11 & A tyrannical ruler outlaws the use of magic without official sanction. A secret society of spellcasters seeks to oust the tyrant.\\
12 & During a drought, low water levels in a lake reveal previously unknown ancient ruins that contain a powerful evil.\\
\end{DndTable}

\subsubsection{Levels 17–20: Masters of the World}
At this tier, adventures have far-reaching consequences, possibly determining the fate of millions on the Material Plane and even places beyond. Characters traverse otherworldly realms and explore demiplanes and other extraplanar locales, where they fight terrible balor demons, titans, archdevils, liches, ancient dragons, and even manifestations of the gods.

\begin{DndTable}[header=Levels 17–20 Adventure Situations]{cX}

1d10 & Situation\\
1 & An ancient dragon is scheming to destroy a god and take the god's place in the pantheon. The dragon's minions are searching for Artifacts that can summon and weaken this god.\\
2 & A band of giants drove away a metallic dragon and took over the dragon's lair, and the dragon wants to reclaim the lair.\\
3 & An ancient hero returns from the dead to prepare the world for the return of an equally ancient monster.\\
4 & An ancient Artifact has the power to defeat or imprison a rampaging titan.\\
5 & A god of agriculture is angry, causing rivers to dry up and crops to wither.\\
6 & An Artifact belonging to a god falls into mortal hands.\\
7 & A titan imprisoned in the Underdark begins to break free, causing terrible earthquakes that are only a hint of the destruction that the titan will cause if it is released.\\
8 & A lich tries to exterminate any spellcasters that approach the lich's level of power.\\
9 & A holy temple was built around a portal leading to one of the Lower Planes to prevent any evil from passing through in either direction. Now the temple has come under siege from both directions.\\
10 & Five ancient metallic dragons lair in the Pillars of Creation. If all these dragons are killed, the world will collapse into chaos. One has just been slain.\\
\end{DndTable}

\subsection{Adventure Setting}
Many D\&D adventures revolve dungeons—interior spaces such as great halls and tombs, subterranean monster lairs, mazes riddled with traps, natural caverns extending for miles beneath the surface, and ruined castles. The "\textit{Dungeons}" section in \textit{chapter 3} can help you craft a dungeon environment for an adventure.

Of course, not every adventure takes place in a dungeon. A wilderness trek across the desert or a harrowing journey into the jungle can be an exciting adventure in its own right. Outdoors, dragons wheel across the sky in search of prey, fierce warriors pour forth from grim fortresses to wage war against their neighbors, ogres and trolls plunder farmsteads for food, and monstrous spiders drop from web-shrouded trees.

Adventures can also take place in cities, towns, and villages, which are often no less dangerous than dungeons or the wilds. The "\textit{Settlements}" section in \textit{chapter 3} can help you create a settlement where an adventure can take place.

\begin{DndSidebar}[float=!b]{Writing for Yourself}
When you're preparing an adventure to run for your friends, you don't need to write hundreds of pages describing each location in exhaustive detail. You can run a game with no more written notes than you'll find in one of the short adventures at the end of this chapter.



\end{DndSidebar}

\subsubsection{Adventure Maps}
An adventure location almost always benefits from a map, and the more thoughtfully constructed the map is, the more fun players are likely to have as their characters explore the location.

\subparagraph{Maps and Adventure Structure}
An adventure map can take many forms—from a detailed dungeon map that shows the dimensions and contents of every room to a rough outline of how one encounter might lead to another, depending on the route the characters choose. Whatever form your map takes, it functions as a flowchart since each decision point (a branch in a corridor, a room with multiple exits) leads to new decision points. If the characters leave a room by the north door, you check your map and determine it leads them into the great hall, lined with pillars, where the fire giant king holds court. If they leave by the secret door to the southeast, you check the map and follow the secret tunnel as it winds to the hidden vaults below the great hall.

\subparagraph{Sample Maps}
\textit{Appendix B} contains maps you can use for your adventures or as inspiration for your own maps. You can also modify those maps to fit the details of the location you have in mind.

\subparagraph{Map Inspiration}
The internet is a great place to find adventure maps that have been made available, as well as real-world building floor plans and city maps and other images that can inspire your mapmaking.

\subsubsection{Bringing a Location to Life}
An inhabited adventure location has its own ecosystem. The creatures that live there need to eat, drink, breathe, and sleep. Predators need prey, and intelligent creatures search for lairs offering the best combination of air, food, water, and security. Keep these factors in mind when designing an adventure location. If the site has an internal logic, adventurers can use their understanding of that logic to make informed decisions.

For example, characters who find a pool of fresh water in a dungeon might infer that many of the creatures inhabiting the dungeon come to that spot to drink. The adventurers might set an ambush at the pool. Likewise, closed or locked doors can restrict the movement of some creatures. A dungeon infested with carrion crawlers or stirges would need open passages so that these creatures can move about to find food.

\subsubsection{Adventure Inhabitants}
The monsters in any adventure location are more than a collection of random creatures that happen to live near one another. Fungi, natural animals, scavengers, and predators can coexist in a complex ecology, alongside sapient creatures who share living space through some combination of negotiation and domination.

Each creature's entry in the \textit{Monster Manual} indicates the terrain types where that creature is most often found, and that book also includes tables listing the creatures commonly found within each type of terrain. Using that information, you can decide which creatures inhabit an adventure location within a particular environment. You can choose a range of creatures, from common vermin to sapient inhabitants and terrifying predators, and decide how they live together.

\subparagraph{Factions}
Particularly in larger areas, groups of creatures might compete for resources. When these groups consist of sapient creatures, opportunities abound for the adventurers who enter those areas. Characters might ally with one group or play groups against each other to reduce the threat of the more powerful monsters. For example, in a dungeon inhabited by mind flayers and the grimlocks they rule, the adventurers might try to incite the grimlocks to revolt against the mind flayers.

Bring the NPC leaders of groups to life as described in the "\textit{Nonplayer Characters}" section in \textit{chapter 3}, fleshing out their personalities and goals. Then use those elements to decide how those leaders respond to adventurers.

\section{Draw In the Players}
If an adventure situation directly affects the characters or the people and places they care about, that is often enough motivation for the characters to get involved. (However, see "\textit{Respect for the Players}" in \textit{chapter 1} for advice about harming the people and places characters love.)

If the adventure situation doesn't have an obvious impact on the characters or the people or things they care about, you can use other techniques to draw in the players. These are best tailored to the motivations of your players and their characters. For example, some adventuring groups are noble heroes who respond without hesitation to the pleas of innocent villagers crying for help; other groups are hardened mercenaries who respond only to offers of payment. Some groups are devoted to gods, rulers, or other patrons who might send them on quests, either directly or through intermediaries.

\begin{DndSidebar}[float=!b]{Subvert Clichés}
As you populate your world with interesting supporting characters, consider the following:




\begin{description}
\item[Avoiding Stereotypes.]
Show how multiple people from the same culture are different. Don't use a real-world accent in a disparaging way.
\item[Beautiful Diversity.]
Feature members of different genders, ethnicities, and sexualities, as well as people with varied beliefs, capabilities, roles, professions, interests, and outlooks.
\item[Fresh Spin.]
Whenever possible, put a fresh spin on a familiar trope. The mysterious figure who presents adventurers with a quest on behalf of the king might be the king in disguise. The wizard in the tower might be a projected illusion created by a band of thieves to guard their loot.
\item[Relatability.]
Treat NPCs as real people with real motivations. Put yourself in their shoes. What would you do?
\end{description}



\end{DndSidebar}

\subsection{Adventure Patrons}
Many adventures begin with a patron asking the characters to undertake a quest or mission, offering a reward in exchange for this service.

Take the time to flesh out an NPC who serves as a patron. Once in a while, it can be interesting for the characters' patron to betray them. Pulling that trick more than once in a campaign, though, is likely to make the players unwilling to trust any future patrons and possibly suspicious about any adventure hooks you put in front of them.

The Patron Hooks table offers some suggestions for ways a patron can lead characters to an adventure situation. The "\textit{Campaign Start}" section in \textit{chapter 5} offers some more suggestions for patrons.

\begin{DndTable}[header=Patron Hooks]{cX}

1d6 & Hook\\
1 & A town crier announces that someone is hoping to hire adventurers.\\
2 & Someone the characters want to impress or need a favor from asks them to deal with the adventure situation.\\
3 & When the characters arrive in a new city, they find a job board where someone has posted in search of adventurers.\\
4 & A wealthy patron who is aware of the adventurers' accomplishments writes to them, offering to pay them for their talents.\\
5 & A citizen in need, who has learned of the adventurers' accomplishments and kindness, travels miles to find them and implore them for help.\\
6 & The adventurers are arrested (on valid or invented charges) and offered a chance to escape punishment by completing a quest.\\
\end{DndTable}

\subsection{Supernatural Hooks}
Celestial omens, vivid dreams, or other magical phenomena can point characters to the adventure situation and suggest a course of action. The Supernatural Hooks table offers some suggestions.

\begin{DndTable}[header=Supernatural Hooks]{cX}

1d6 & Hook\\
1 & The characters all have a vivid dream that foreshadows elements of the adventure.\\
2 & While preparing spells, one character receives a quest from a god or patron.\\
3 & A fortune teller's reading for one of the characters points to a quest and offers hints about challenges that lie ahead.\\
4 & Flames, clouds, smoke, or huge flocks of birds take distinct shapes that portend the adventure situation.\\
5 & Animals or animated objects speak clearly to direct the adventurers toward the situation.\\
6 & Someone who died returns as a ghost and haunts the characters. The ghost prompts the characters to investigate the cause of the ghost's death and put it to rest.\\
\end{DndTable}

\subsection{Happenstance Hooks}
Sometimes, characters just happen on an adventure through sheer coincidence—or at least what appears to be coincidence (which might actually involve divine or other supernatural intervention). The Happenstance Hooks table provides some ideas.

\begin{DndTable}[header=Happenstance Hooks]{cX}

1d6 & Hook\\
1 & The characters find a letter describing the adventure situation.\\
2 & The characters are on an unrelated quest, such as searching for a particular magic item, that leads them into the adventure situation.\\
3 & The adventure situation disrupts a festival or ceremony that the characters are attending.\\
4 & A magical mishap places the characters in the adventure situation.\\
5 & While traveling in a caravan or aboard a ship, the characters befriend an NPC who has news about the adventure situation.\\
6 & The characters are attacked after being mistaken for another group of adventurers. They learn about the adventure situation from a clue left behind by their attackers.\\
\end{DndTable}

\section{Plan Encounters}
Encounters are the individual scenes in the larger story of your adventure. Reduced to fundamentals, an encounter is an objective with an obstacle. It accomplishes one or more of the following:


\begin{itemize}
\item Moving characters closer to achieving a goal
\item Frustrating the characters' progress toward a goal
\item Revealing new information
\end{itemize}

\subsection{Character Objectives}
The following objectives can be foundations for encounters. Although these objectives focus on a single encounter during an adventure, using the same objective in multiple encounters allows you to combine these encounters into a larger obstacle or problem the adventurers must overcome.

\subsubsection{Make Peace}
The characters must convince two opposing groups (or their leaders) to end the conflict that embroils them. As a complication, the characters might have enemies on one or both of the opposing sides, or some other group or individual might be instigating the conflict to further its own ends. An encounter aimed at making peace might involve only social interaction, perhaps with the threat of combat if negotiations go poorly. It could also begin as a combat encounter, with the characters trying to stop the fighting and get the parties talking to each other.

\subsubsection{Protect an NPC or Object}
The characters must act as bodyguards or protect some object in their custody. As a complication, an NPC under the party's protection might be cursed, panicked, unable to fight, or apt to risk the lives of the adventurers through dubious decisions. The object the adventurers have sworn to protect might be sentient, cursed, or difficult to transport. Such an encounter might be a combat encounter or an exploration encounter, with either Hostile [Attitude] monsters or a dangerous environment threatening the NPC or object the characters are trying to protect. If the characters are protecting an NPC, this objective can add an element of social interaction to a combat or exploration encounter.

\subsubsection{Retrieve an Object}
The adventurers must gain possession of a specific object in the area of the encounter, often with a time limit. This might be a combat encounter, with monsters protecting the object, or an exploration encounter, with traps or hazards preventing access to the object. As a complication, enemies might desire the object as much as the adventurers do, adding a combat element to an exploration encounter or complicating a combat.

\subsubsection{Run a Gauntlet}
The adventurers must pass through a dangerous area. As with retrieving an object, reaching the exit is a higher priority than killing opponents in the area. A time limit adds a complication, as does a decision point that might lead characters astray. This might be an exploration encounter, with traps and hazards as complications, or a combat encounter against Hostile [Attitude] monsters.

\subsubsection{Sneak In}
The adventurers need to move through the encounter area without their enemies noticing. This is typically an exploration encounter, but if the characters are detected, a combat encounter or social interaction might result.

\subsubsection{Stop a Ritual}
The plots of evil cult leaders, malevolent spellcasters, and Fiends often involve rituals that must be foiled. In a combat encounter, characters might have to fight their way through evil minions before attempting to disrupt the ritual's magic. This could also drive an exploration encounter, where the challenge is getting to the place where the ritual is occurring, or a social interaction encounter, where the objective is convincing the ritual leaders to stop their rite. As a complication, the ritual might be close to completion when the characters arrive, imposing a time limit. The ritual's completion might have immediate consequences, too.

\subsubsection{Take Out a Single Target}
A villain the characters seek to defeat is surrounded by minions powerful enough to kill the adventurers. The characters can flee and hope to confront the villain another day, or they can try to fight their way through the minions. As a complication, the minions might be innocent creatures under the villain's control. Killing the villain means breaking that control, but the adventurers must endure the minions' attacks until the villain falls.

\subsection{Keeping the Adventure Moving}
Make sure your players have clear objectives they can pursue at every stage of the adventure. Three simple techniques can ensure that the players understand the task at hand and how to pursue it:


\begin{description}
\item[Adviser NPCs.]
A helpful NPC in a social interaction can offer advice and suggestions to the characters. Such an NPC might be the patron who initially sent the characters on the adventure, someone they met along the way, or a character's contact. When you're planning an adventure, include NPCs who can fill this role.
\item[Evil Intrusion.]
If things start grinding to a halt, have the characters encounter a minion or monster connected to the adventure's main threat. At the end of the encounter, perhaps the characters find information that gets them back on track. Plan one or two encounters like this ahead of time.
\item[The DM's Role.]
If the characters can't figure out how to solve an encounter or aren't sure what to do next, you can remind the players of things their characters have already learned or call for Intelligence (Investigation) or similar checks to see if their characters can remember and connect things that the players might be missing.
\end{description}

\subsection{Something for Everyone}
As described in the "\textit{Know Your Players}" section in \textit{chapter 2}, players have different tastes in the activities they enjoy in the game. An adventure needs to account for the different players and characters in your group to draw them into the story.

An adventure that includes a balance of exploration, social interaction, and combat is likely to appeal to a wide breadth of players. But an adventure you create for your home campaign doesn't have to appeal to every possible player interest—only to the players at your table.

You can design encounters that appeal to multiple player motivations. Imagine a fight pitting the characters against a gang of gnolls, delighting the players who enjoy fighting. Then a young dragon wanders into the middle of the fight. Suddenly the fight can swing one of two ways: the dragon could help the gnolls against the party or help the party against the gnolls. It's up to the players who thrive on acting to persuade the dragon to help the party.

\subsection{Multiple Ways to Progress}
Make sure there are multiple ways the characters can progress through the adventure at any point, so if they miss one way, they have an alternative. Plan opportunities for the adventure to move forward even when the characters fail. Use challenges with a single path to success only as chances for the characters to obtain extra rewards.

\subsection{Social Interaction Encounters}
The "\textit{Running Social Interaction}" section in \textit{chapter 2} offers advice for handling social interaction encounters and can help you craft these encounters. The "\textit{Nonplayer Characters}" section in \textit{chapter 3} is also essential for creating these encounters.

\subsection{Exploration Encounters}
An encounter centered on exploration might involve the characters trying to disarm a trap, find a secret door, or discover something about the adventure location. An exploration encounter could also involve the characters spending a day crossing a rolling plain or traversing vast caverns.

The "\textit{Running Exploration}" section in \textit{chapter 2} can help you craft these encounters as well as run them. Various sections in \textit{chapter 3} can also help you detail obstacles and dangers in an exploration encounter: see "\textit{Chases}," "\textit{Curses and Magical Contagions}," "\textit{Doors}," "\textit{Environmental Effects}," "\textit{Hazards}," "\textit{Poison}," and "\textit{Traps}" in particular.

\subsection{Combat Encounters}
The following features can make a combat encounter more interesting or challenging:


\begin{description}
\item[Changes in Elevation.]
Terrain features that provide a change of elevation (such as stacks of empty crates, ledges, and balconies) reward clever positioning and encourage characters to jump, climb, fly, or teleport.
\item[Defensive Positions.]
Enemies in hard-to-reach locations or defensive positions force characters who normally attack at range to move around.
\item[Hazards.]
The "\textit{Hazards}" section in \textit{chapter 3} describes dangerous features, such as patches of green slime, that characters or their enemies can use to their advantage.
\item[Mixed Monster Groups.]
When different types of monsters work together, they can combine their abilities—just like characters with different classes and origins. A diverse force is more powerful.
\item[Reasons to Move.]
Use features that encourage characters and their enemies to move around, such as chandeliers, kegs of gunpowder or oil, and rolling stone traps.
\end{description}

\subsubsection{Combat Encounter Difficulty}
Use the following guidelines to create a combat encounter of a desired level of difficulty.

\subparagraph{Step 1: Choose a Difficulty}
Three categories describe the range of encounter difficulty:


\begin{description}
\item[Low Difficulty.]
An encounter of low difficulty is likely to have one or two scary moments for the players, but their characters should emerge victorious with no casualties. One or more of them might need to use healing resources, however. As a rough guideline, a single monster generally presents a low-difficulty challenge for a party of four characters whose level equals the monster's CR.
\item[Moderate Difficulty.]
Absent healing and other resources, an encounter of moderate difficulty could go badly for the adventurers. Weaker characters might get taken out of the fight, and there's a slim chance that one or more characters might die.
\item[High Difficulty.]
A high-difficulty encounter could be lethal for one or more characters. To survive it, the characters will need smart tactics, quick thinking, and maybe even a little luck.
\end{description}

\subparagraph{Step 2: Determine Your XP Budget}
Using the XP Budget per Character table, cross-reference the party's level with the desired encounter difficulty. Multiply the number in the table by the number of characters in the party to get your XP budget for the encounter.

\begin{DndTable}[header=XP Budget per Character]{cccc}

1 & 50 & 75 & 100\\
2 & 100 & 150 & 200\\
3 & 150 & 225 & 400\\
4 & 250 & 375 & 500\\
5 & 500 & 750 & 1,100\\
6 & 600 & 1,000 & 1,400\\
7 & 750 & 1,300 & 1,700\\
8 & 1,000 & 1,700 & 2,100\\
9 & 1,300 & 2,000 & 2,600\\
10 & 1,600 & 2,300 & 3,100\\
11 & 1,900 & 2,900 & 4,100\\
12 & 2,200 & 3,700 & 4,700\\
13 & 2,600 & 4,200 & 5,400\\
14 & 2,900 & 4,900 & 6,200\\
15 & 3,300 & 5,400 & 7,800\\
16 & 3,800 & 6,100 & 9,800\\
17 & 4,500 & 7,200 & 11,700\\
18 & 5,000 & 8,700 & 14,200\\
19 & 5,500 & 10,700 & 17,200\\
20 & 6,400 & 13,200 & 22,000\\
\end{DndTable}

\subparagraph{Step 3: Spend Your Budget}
Every creature has an XP value in its stat block. When you add a creature to your combat encounter, deduct its XP from your XP budget to determine how many XP you have left to spend. Spend as much of your XP budget as you can without going over. It's OK if you have a few unspent XP left over. Examples are given below:


\begin{description}
\item[Example 1.]
A low-difficulty encounter for four level 1 characters has an XP budget of 50 × 4, for a total of 200 XP. With that, you could build any of the following encounters:
\end{description}


\begin{itemize}
\item 1 Bugbear Warrior (200 XP)
\item 2 Giant Wasps (100 XP each), for 200 XP total
\item 6 Twig Blights (25 XP each), for 150 XP total
\end{itemize}


\begin{description}
\item[Example 2.]
A moderate-difficulty encounter for five level 3 characters has an XP budget of 225 × 5, for a total of 1,125 XP. With that, you could build either of these encounters:
\end{description}


\begin{itemize}
\item 2 Nothics (450 XP each) and 9 Stirges (25 XP each), for 1,125 XP total
\item 1 \textbf{Wight} (700 XP), 1 \textbf{Warhorse Skeleton} (100 XP), and 6 Skeletons (50 XP each), for 1,100 XP total
\end{itemize}


\begin{description}
\item[Example 3.]
A high-difficulty encounter for six level 15 characters has an XP budget of 7,800 × 6, for a total of 46,800 XP. With that, you could build this encounter:
\end{description}


\begin{itemize}
\item 2 Adult Red Dragons (18,000 XP each) and 2 Fire Giants (5,000 XP each), for 46,000 XP total
\end{itemize}

\subsubsection{Troubleshooting}
When creating and running combat encounters, keep the following in mind.

\subparagraph{Many Creatures}
The more creatures in an encounter, the higher the risk that a lucky streak on their part could deal more damage to the characters than you expect. If your encounter includes more than two creatures per character, include fragile creatures that can be defeated quickly. This guideline is especially important for characters of level 1 or 2.

\subparagraph{Adjustments}
A player's absence might warrant removing creatures from an encounter to keep it at the intended difficulty. Also, die rolls and other factors can result in an encounter being easier or harder than intended. You can adjust an encounter on the fly, such as by having creatures flee (making the encounter easier) or adding reinforcements (making the encounter harder).

\subparagraph{CR 0 Creatures}
Creatures that have a CR of 0, particularly ones that are worth 0 XP, should be used sparingly. If you want to include many CR 0 critters in an encounter, use swarms from the \textit{Monster Manual} instead.

\subparagraph{Number of Stat Blocks}
The best combat encounters often pair one kind of creature with another, such as fire giants paired with hell hounds. Be mindful of the number of stat blocks you need to run the encounter. Referencing more than two or three stat blocks for a single encounter can be daunting, particularly if the creatures are complex.

\subparagraph{Powerful Creatures}
If your combat encounter includes a creature whose CR is higher than the party's level, be aware that such a creature might deal enough damage with a single action to take out one or more characters. For example, an \textbf{Ogre} (CR 2) can kill a level 1 Wizard with a single blow.

\subparagraph{Unusual Features}
If a monster has a feature that lower-level characters can't easily overcome, consider not adding that monster to an encounter for characters whose level is lower than the monster's Challenge Rating.

\subsection{Monster Behavior}
The attitudes, motivations, and behavior of the monsters in an encounter help determine how a social interaction plays out (and whether it might erupt into combat) and influence the course of combat.

\subsubsection{Initial Attitudes}
A published adventure typically notes or implies whether a creature's initial attitude toward the adventurers is Friendly [Attitude], Indifferent [Attitude], or Hostile [Attitude]. In an encounter you've created, you can decide that starting attitude, or you can randomly determine it using the Initial Attitude table.

\begin{DndTable}[header=Initial Attitude]{cX}

1d12* & Initial Attitude\\
1—4 & Hostile\\
5—8 & Indifferent\\
9—12 & Friendly\\
\end{DndTable}

\subsubsection{Monster Personality}
If an encounter involves a significant individual, use the guidance in the "\textit{Nonplayer Characters}" section of \textit{chapter 3} to flesh out the details of that individual's personality and aims. For a group of nameless monsters, you can decide on a personality based on the monsters' entries in the \textit{Monster Manual}, or you can use the Monster Personality table to inform how you portray the monsters and their actions. It's simplest to assign the same personality traits to an entire group of monsters in an encounter. For example, one bandit gang might be an unruly mob of braggarts, while the members of another gang are always on edge and ready to flee at the first sign of danger.

\begin{DndTable}[header=Monster Personality]{cX}

1d8 & Personality\\
1 & Cowardly; surrenders easily\\
2 & Greedy; wants treasure\\
3 & Boastful; makes a show of bravery but runs from danger\\
4 & Disorderly; poorly trained and easily rattled\\
5 & Fanatical; ready to die fighting\\
6 & Brave; stands firm against danger\\
7 & Jocular; taunts enemies\\
8 & Orderly; difficult to rattle\\
\end{DndTable}

\subsubsection{Monster Relationships}
Encounters with groups of monsters can be more interesting if rivalries, hatreds, or attachments exist among the monsters in the group. The death of a much-revered leader might throw its followers into a frenzy. On the other hand, a monster might flee if its hunting companion is killed, or a mistreated toady might be eager to surrender and betray its boss in return for its life. You can use the Monster Relationships table to inspire such relationships within a monster group.

\begin{DndTable}[header=Monster Relationships]{cX}

1d6 & Relationship\\
1 & Two monsters have a bitter rivalry; each wants the other to suffer.\\
2 & One monster, bullied by the others, hangs back and flees at the first opportunity.\\
3 & One monster is revered or even worshiped by the others, who will die for it.\\
4 & One monster is admired by the group; its allies try to impress or help it.\\
5 & One monster cares only for itself and not the rest of the group.\\
6 & One monster bullies the others; it forces them into danger, but they want it defeated.\\
\end{DndTable}

\subsubsection{Reactive Tactics}
A great way to make an adventure location feel alive—particularly an organized base or stronghold—is to allow its denizens to react to the presence of the adventurers. Once they are aware of trespassers, sapient creatures might either fortify their own locations or leave those locations to assist colleagues and expel the invaders.

Take a copy of the adventure map, and pencil in the locations of all its inhabitants to give yourself a sense of where they're located relative to each other. When the adventurers engage in combat or any other noisy activity, assume that nearby creatures hear the noise and are alerted to the adventurers' presence. (Creatures that can't hear might be alerted by vibrations or other sensory cues.) Once alerted, a creature has several options:


\begin{description}
\item[Ambush.]
The creature leaves its current location and takes up a position near the adventurers' location, hoping to catch the adventurers unaware. Ambushers try to hide and, once hidden, take advantage of any opportune moment to attack.
\item[Fortify.]
The creature attempts to fortify its location by using furniture or heavy objects to block doors or entryways. Increase the DC to force open a blocked door by 3, and hastily blocked passageways are Difficult Terrain. A creature with a ranged attack that selects this option seeks cover in its location.
\item[Hide/Flee.]
If possible, the creature hides in its current location, hoping to avoid any adventurers that enter. If there is nowhere to hide, the creature flees to a location farther from the adventurers or flees to any nearby location occupied by its allies. If a fleeing creature reaches allies, those allies are immediately alerted to the presence of the adventurers; determine how those allies react.
\item[Investigate.]
The creature rushes to the sound of the disturbance to investigate, possibly joining any battle that is underway. A creature that passes near an ally while moving to investigate a disturbance might ask that ally to accompany it.
\end{description}

If you can't decide what an alerted creature should do, have it make a DC 10 Wisdom saving throw. On a successful save, the creature either investigates or lays an ambush; on a failed save, it hides or fortifies its location. When dealing with a group of creatures, the leader makes this saving throw on behalf of the entire group.

As creatures employ reactive tactics, make notes about their new locations on your adventure map.

\subsubsection{Prepared Defenders}
Sapient creatures that have reason to believe their lair is likely to be invaded might set up a defense. Reasons to set up a defense include the following:


\begin{itemize}
\item Adventurers invaded the lair recently and retreated.
\item Scouts affiliated with the lair's denizens noticed the adventurers heading toward the lair.
\item A spy in a nearby settlement overheard or discovered the adventurers' plans and alerted the lair's occupants.
\end{itemize}

Prepared defenders commonly use one or more of the following tactics:


\begin{description}
\item[Ambushes and Barricades.]
Some defenders might move from their original locations to locations where they can hide near critical passageways. Defenders might also use furniture and debris to block off passages, hoping to channel invaders toward prepared strongpoints.
\item[Sentries and Alarms.]
Some creatures might move from their keyed locations to locations that allow them to monitor entrances to the lair. If possible, these creatures might equip themselves with an alarm such as a horn or improvised gong (an empty kettle, perhaps). At the first sight of intruders, they raise the alarm, alerting nearby allies so those comrades can employ the reactive tactics described above.
\item[Traps.]
The defenders might place additional traps, such as falling nets, throughout the lair (see "\textit{Traps}" in \textit{chapter 3}).
\end{description}

If the denizens of an adventure location employ any of these tactics, update your map as appropriate.

\subsection{Encounter Pace and Tension}
A good story hooks you in with an interesting introduction, builds tension steadily throughout the story, then reaches a climactic conclusion. It's not always easy to mimic that structure in an adventure where the players control their characters' actions, but you can use the encounters you plan to build tension toward a climax.

Each encounter in an adventure is an opportunity to make the characters' situation more complex and urgent, with more significant consequences. Successive encounters raise the tension in an adventure naturally, as characters spend their limited resources. Variety also contributes to a sense of escalating tension. Build variety into your encounters in three ways:


\begin{description}
\item[Vary Encounter Type.]
Use a mix of social interaction, exploration, and combat encounters. Different types of encounters provide different amounts of tension (generally, combat encounters offer the most), but they also feel very different and can have drastically different stakes.
\item[Vary Encounter Difficulty.]
Include encounters that offer low, moderate, and high difficulty. A mix of low- and moderate-difficulty encounters early in the adventure can lead to a climactic high-difficulty fight, perhaps against the adventure's primary villain or another threat.
\item[Vary Threats.]
Build encounters using different threats. If the characters are delving into a kuo-toa temple and therefore expect many encounters to include various kuo-toa, look for opportunities to include different monsters that might serve as guards, pets, or allies to the kuo-toa. Include a variety of hazards, environmental dangers, and traps in exploration encounters, and use NPCs with different personalities and different goals in social interaction encounters.
\end{description}

\subsubsection{Urgency and Rests}
While successive encounters increase tension, taking a Short Rest relaxes the tension somewhat, as characters have a chance to replenish some of their resources. In many adventures, though, the characters and their players have a sense that they need to act quickly to deal with the situation presented by the adventure. This creates tension between the need to rest and the sense that things are getting worse while the characters are resting.

You can influence the pace and tension of your adventure by determining where and when the characters can rest. If the characters are exploring a vast dungeon, consider scattering a few small rooms with only one door, where the characters can bar the door and reasonably expect to spend an hour or even a night resting in safety. On the flip side, cautious characters might try to take a Short Rest between every encounter, never really straining their resources. It's OK to interrupt those rests once in a while to maintain a sense of tension or to heighten the urgency, making it clear that even an hour spent resting could jeopardize their chances of success.

\subparagraph{Easing Up}
It is possible to dial up the urgency to the point that the players feel they don't have time to investigate the interesting details they encounter in the world. When this happens, consider using a helpful NPC to take some of the pressure off. A wise elder might advise them that the situation is not as urgent as they fear, a whimsical Fey being might use magical mischief to force them to slow down, or a kindly Celestial could tell them they're taking the concerns of the mortal world just a bit too seriously.

\subsubsection{Random Encounters}
Random encounters are randomly determined encounters that don't occur in a fixed location. The options are often presented in a table. When a random encounter occurs, you roll a die and consult the table to determine what the party encounters. Sample random encounter tables appear in the adventure examples later in this chapter. Similar tables appear in many published adventures and rulebooks, and you can easily create your own by following these examples.

Handled well, random encounters can serve a variety of useful purposes.

\subparagraph{Create Urgency}
Wandering monsters encourage characters to keep moving and to find a safe place to rest. (Sometimes you can create a sense of danger and urgency by rolling dice behind your DM screen, even without an actual encounter!)

\subparagraph{Drain Character Resources}
By draining the party's Hit Points and spell slots, leaving the adventurers feeling underpowered and vulnerable, random encounters can build tension in an adventure.

\subparagraph{Establish Atmosphere}
Thematic links among monsters appearing in random encounters create a tone and an atmosphere that define the environment the characters are exploring. For example, an encounter table filled with bats, wraiths, giant spiders, and zombies creates a sense of horror and suggests the possibility of even more terrifying foes.

\subparagraph{Provide Assistance}
Some random encounters can benefit the characters instead of hindering or harming them. Helpful creatures might give the characters useful information or assistance.

\subparagraph{Reinforce Campaign Themes}
Random encounters can remind the players of the major themes and conflicts in your campaign. For example, if a war between two nations is a major conflict in your campaign, you might design random encounter tables to reinforce the ever-present threat of that conflict, including encounters with bedraggled troops returning from battle, refugees fleeing invaders, lone messengers riding for the front lines, enemy war parties, and so on.

\subparagraph{A Life of Its Own}
Sometimes a random encounter that starts off incidental can become important to the story. A random encounter with a wandering ogre might end with the ogre offering to help the characters get where they need to go, or the ogre might have something in a nearby den that is significant to the adventure story—a prisoner, a stolen item, or important information. Or the players might find the ogre's personality delightful, prompting you to make the ogre a more important part of the story.

\section{Bring It to an End}
The climactic ending of an adventure fulfills the promise of all that came before. The best climax is one the players see coming, so if a dragon is the mastermind behind all the nefarious activity happening in an adventure, having the dragon's minions mention the nature of their boss sets up the coming climactic encounter.

Although the climax must hinge on the successes and failures of the characters up to that moment, the Adventure Climax table can provide suggestions to help you shape the end of your adventure.

\begin{DndTable}[header=Adventure Climax]{cX}

1d10 & Climax\\
1 & The adventurers confront a villain and a group of minions in a battle to the finish.\\
2 & The adventurers chase a villain while dodging obstacles designed to thwart them, leading to a final confrontation in the villain's refuge.\\
3 & The actions of the adventurers or a villain result in a cataclysmic event that the adventurers must escape.\\
4 & The adventurers race to the site where a villain is bringing a master plan to its conclusion, arriving just as that plan is about to be completed.\\
5 & A villain and two or three lieutenants perform separate rites in a large room. The adventurers must disrupt all the rites.\\
6 & An ally betrays the adventurers as they're about to achieve their goal. (Use this climax carefully, and don't overuse it.)\\
7 & A portal opens to another plane of existence. Creatures on the other side spill out, forcing the adventurers to close the portal while dealing with a villain at the same time.\\
8 & The dungeon begins to collapse while a villain attempts to escape in the chaos.\\
9 & The adventurers must choose whether to pursue a fleeing villain or save an NPC they care about or a group of innocents.\\
10 & Just when the characters think the main threat is defeated, it transforms into a different monster or a more powerful form.\\
\end{DndTable}

\subsection{Denouement}
In most stories, there's a period after the climax, in which loose plot threads are tied up and everything is explained. An adventure might also include this kind of denouement: time to discover what treasure is in the dragon's hoard, an award ceremony where the queen gives medals to the victorious heroes, or even a time to mourn adventuring companions who didn't survive the battle.

The denouement can also be an opportunity for the players to identify loose threads that haven't been tied up—threads that can lead them into the next adventure. The "\textit{Episodes and Serials}" section in \textit{chapter 5} offers suggestions for weaving these connecting threads.

\section{Adventure Rewards}
For some characters, the prospect of material reward is their primary reason for going on adventures. For others it's a welcome added benefit to pursuing their other goals.

\textit{Chapter 7} describes different kinds of treasure, but see also "\textit{Marks of Prestige}" in \textit{chapter 3} for other rewards you might use.

The following sections describe how treasure is typically dispersed in an adventure.

\subsection{Individual Treasure}
Characters might find small amounts of treasure in the pockets, pouches, or personal stashes of individual monsters. Even if a monster doesn't intentionally collect treasure, characters might find scattered coins and other monetary treasure left behind by the monster's previous victims.

You can use the Random Individual Treasure table to determine how much treasure a single monster has based on its Challenge Rating (CR). The table includes the average total in parentheses, which you can use instead of rolling. To determine the total amount of treasure for a group of similar creatures, you can roll once and multiply the total by the number of creatures in the group.

\begin{DndTable}[header=Random Individual Treasure]{cX}

CR & Treasure\\
0–4 & 3d6 (10) GP\\
5–10 & 2d8 × 10 (90) GP\\
11–16 & 2d10 × 10 (110) PP\\
17+ & 2d8 × 100 (900) PP\\
\end{DndTable}

\subsection{Treasure Hoards}
Adventurers sometimes discover large caches of treasure, the accumulated wealth of a large group of creatures or the belongings of a single powerful creature that hoards valuables. The Random Treasure Hoard table can help you create such a cache. When determining the contents of a hoard belonging to one monster, use the row for that monster's Challenge Rating (CR). When the hoard belongs to a large group of monsters, use the CR of the monster that leads the group. Each row includes average results for monetary treasure, which you can use instead of rolling. To create a hoard for a monster that is particularly fond of amassing treasure (such as a dragon), you can roll twice on the table or roll once and double the total.

As a rough benchmark, aim to roll on the Random Treasure Hoard table about once per game session. Use the guidelines in \textit{chapter 7} to determine which magic items are in the hoard (see "\textit{Awarding Magic Items}" and "\textit{Random Magic Items}").

\begin{DndTable}[header=Random Treasure Hoard]{cXc}

CR & Monetary Treasure & Magic Items\\
0–4 & 2d4 × 100 (500) GP & 1d4-1\\
5–10 & 8d10 × 100 (4,400) GP & 1d3\\
11–16 & 8d8 × 10,000 (36,000) GP & 1d4\\
17+ & 6d10 × 10,000 (330,000) GP & 1d6\\
\end{DndTable}

\subsection{Quest Rewards}
Sometimes, characters are paid for completing a quest. To determine a suitable quest reward, roll once on the Random Treasure Hoard table, using the characters' level for the Challenge Rating (CR).

\subsection{Monster Treasure Preferences}
The \textit{Monster Manual} gives treasure preferences for monsters in that book. These preferences are categorized as follows:


\begin{description}
\item[Any.]
The monster has a treasure hoard, the contents of which you can determine by rolling on the Random Treasure Hoard table. Monetary treasure can take the form of \textit{coins}, \textit{trade bars}, \textit{trade goods}, \textit{gems}, or \textit{art objects} (all described in \textit{chapter 7}). Magic items can belong to any treasure theme or category (see "\textit{Treasure Themes}" and "\textit{Magic Item Categories}" in \textit{chapter 7}).
\item[Individual.]
The monster doesn't have a treasure hoard; however, it might have monetary treasure, which you can determine by rolling on the \textit{Random Individual Treasure} table. This treasure can take the form of coins, trade bars, trade goods, gems, or art objects (all described in \textit{chapter 7}).
\item[Treasure Theme (Arcana, Armaments, Implements, or Relics).]
The monster has a treasure hoard skewed toward a particular theme (see "\textit{Treasure Themes}" in \textit{chapter 7}). You can determine the size of the hoard by rolling on the \textit{Random Treasure Hoard} table. If the hoard contains magic items, use the \textit{guidelines and tables} in \textit{chapter 7} to determine each one.
\end{description}

\section{Adventure Examples}
This section contains example adventures that demonstrate the principles described throughout the chapter. Each provides enough information for you to run a one-session adventure, with the help of the maps in \textit{Appendix B} and the monster stat blocks in the \textit{Monster Manual}.

Each adventure in this section includes the following information:


\begin{description}
\item[Title.]
An adventure title can help you organize your campaign notes, and if you share the title with your players, it can set the tone for what's ahead.
\item[Character Level.]
Each adventure specifies the level of characters it's aimed at. The difficulty of encounters in each adventure is tailored for four characters of that level. You can use adventures for characters of higher or lower level or for larger or smaller groups. However, the encounters might be easier or harder than you expect unless you adjust them.
\item[Situation.]
Each adventure lays out what's going on—the situation that the adventurers are called on to deal with. See "\textit{Lay Out the Premise}" earlier in this chapter.
\item[Hook.]
Each adventure offers one way to draw characters in to the adventure. See "\textit{Draw In the Players}" earlier in this chapter.
\item[Encounters.]
The rest of each adventure description is a series of encounters. The text describes the location where the encounter occurs, often pointing to the maps in \textit{Appendix B}, and any triggering event that might provoke the encounter. Monster names in \textbf{bold} point you to the stat blocks in the \textit{Monster Manual}. Some encounters also specify treasure the characters might find.
\end{description}

Use your imagination to bring the locations and encounters to life, and build on the ideas the players bring to each encounter. Alter these outlines freely to suit your tastes—and those of your players—and your ideas for your campaign.

\begin{DndSidebar}[float=!b]{A Starting Campaign}
You can use the adventures in this section to get a new campaign off the ground. These adventures are linked to locations near the \textit{Free City of Greyhawk}, as described in \textit{chapter 5}. You can run the first three adventures in sequence, having the characters gain a level after each adventure. They might return to their home base in the city between adventures, or they could travel to Greyhawk after they complete "\textit{The Winged God}." Use the encounters and interactions the characters have in these early levels, and the situations that interest your players, to plan later adventures.



\end{DndSidebar}

\subsection{The Fouled Stream}
\textit{Adventure for Level 1 Characters}


\begin{description}
\item[Situation.]
An alien fungus in a cave is polluting the stream that flows past the village of High Ery. The fungus has spawned vile creatures in and around the cave.
\item[Hook.]
The folk of High Ery are noticing fungal growths on the riverbanks and a layer of scum on the water. The characters might live in the village, or a contact in the Free City of Greyhawk might ask them to investigate.
\end{description}

\subsubsection{Encounters}
The adventure consists of these encounters.

\subparagraph{The First Fork}
A mile upstream from the village, a stream flows into the river from a little wood on the river's south side. Characters can tell that this stream is the source of the pollution.

\subparagraph{Journey Upstream}
Borogrove, a kindly \textbf{Treant}, keeps watch over the wood and meets the characters as they follow the polluted stream. He knows the source of the corruption is inside a cave that the stream spills out of. He gives the characters a magic acorn. If swallowed, the acorn conveys the benefits of a \textit{Potion of Healing} and the \textit{Lesser Restoration} spell.

\subparagraph{Twig Blights}
Just outside the cave, the characters encounter six Twig Blights.

\subparagraph{Corrupted Cave}
Use the Underdark Warren map in \textit{Appendix B} for the corrupted cave. Ignore the secret door and the inner chambers behind it. Close off the tunnels leading off the map to the south, east, and north. The characters enter the cave in the southeast, following the stream. The cave's main features and inhabitants are as follows:


\begin{description}
\item[Entrance.]
A \textbf{Shrieker Fungus} just inside the cave entrance alerts the inhabitants to the characters' arrival. On watch near the entrance and quick to respond to the shriekers' cry are four Bullywug Warriors who have fungal growths on them.
\item[Berserk Bear.]
In a side cave to the southeast is a \textbf{Brown Bear} that drank from the stream. It's upset because the water made it ill. If the characters can make it eat Borogrove's acorn or otherwise rid it the Poisoned condition, the bear recovers immediately and leaves them alone.
\item[Ooze's Lair.]
At the north end of the stream are a Psychic Gray Ooze and six Stirges. After defeating these creatures, the characters can destroy the brain-like fungus in the water, which is the source of the corruption. If they do, each character earns a bonus 100 XP.
\end{description}

\subparagraph{Journey Home}
As they leave the wood, the characters encounter Borogrove again. If they used his acorn, he gives them another one. If they purified the source of the stream, he gives them a \textit{Staff of Flowers} in gratitude.

\subsection{Miner Difficulties}
\textit{Adventure for Level 2 Characters}


\begin{description}
\item[Situation.]
After miners dug into an Underdark tunnel, a \textbf{Hook Horror} found its way into the mine and became trapped. It has eaten a few miners, and the others are too terrified of the echoing clicks in the mine to hunt down the predator.
\item[Hook.]
The mayor of the village of Blackstone, Kristryd Splitanvil (a Lawful Good, dwarf \textbf{Tough}), hires the adventurers—perhaps because of how adeptly they handled the situation in "\textit{The Fouled Stream}"—to deal with the monster in the mines. She offers a precious topaz worth 500 GP to adventurers who kill the creature or drive it away.
\end{description}

\subsubsection{Encounters}
The adventure consists of these encounters.

\subparagraph{Exploring the Mine}
Use the Mine map in \textit{Appendix B}. In the weeks since the miners abandoned the place, pests have flourished in the tunnels. Each time the characters enter a distinct area of the mine, roll on the following table.


\begin{DndTable}{cX}

1d6 & Encounter\\
1 & Four Violet Fungi and one \textbf{Rust Monster}\\
2 & One \textbf{Giant Spider} and two Swarms of Insects (spiders)\\
3 & One \textbf{Darkmantle} and three Piercers\\
4 & A patch of \textit{yellow mold} (see "\textit{Hazards}" in \textit{chapter 3}) on a miner's remains\\
5 & One \textbf{Gelatinous Cube}\\
6 & Sounds of Terror (see below)\\
\end{DndTable}

\subparagraph{Sounds of Terror}
The first time this encounter occurs, the characters hear eerie clacking and scraping noises echoing in the mine shafts and notice gouges in the walls. The second time, they find the source of those sounds and markings—the \textbf{Hook Horror}. The creature is hungry, but what it really wants is to find its way back to the Underdark. It retreats from a fight that is going badly for it.

\subparagraph{Underdark Connection}
The tunnel in the southeast corner of the bottom level of the mine is where the miners accidentally connected to an Underdark tunnel. The hook horror entered the mine through a hole in the wall, but the hole closed behind it in a cascade of rubble. If the characters clear away the rubble from the hole, they hear clacking and scraping from the other side. The hook horror hears it too and hurries here to rejoin its kin. As long as the characters stay out of its way, it ignores them as it scrambles back through the hole, never to plague the mine again. Each character then earns a bonus 200 XP.

\subsection{The Winged God}
\textit{Adventure for Level 3 Characters}


\begin{description}
\item[Situation.]
A few weeks ago, a \textbf{Red Dragon Wyrmling} drove a band of kobolds out of their warren to claim the place as its lair. Now some of the kobolds are causing trouble in the Cairn Hills. They're raiding merchants, hoping the dragon will allow them to return home as its loyal servants.
\item[Hook.]
A merchant named Nondy Barducks (a Lawful Neutral, gnome \textbf{Commoner}) hires the characters to escort his wagon to the remote mining village of Diamond Lake, which happens to be near the dragon's new lair. Nondy was robbed by kobolds on his last trip, and he wants protection this time. He offers to pay each character 150 GP.
\end{description}

\subsubsection{Encounters}
The adventure consists of these encounters.

\subparagraph{Kobold Bandits}
Along the road, the wagon is surrounded by eight Kobold Warriors (Neutral) who demand that the merchant surrender his goods. In combat, the kobolds shout things like "For the Winged God!" and "Fight to reclaim our home!" If four of them fall in battle, the remaining kobolds try to flee. Any captured kobold explains the situation.

\subparagraph{Kobold Supplicants}
If the characters continue on their way without pursuing the kobolds, twelve Kobold Warriors (including any survivors of the first encounter) and six Winged Kobolds (all Neutral) approach the wagon. These kobolds humbly ask the adventurers to help them. They promise to return the goods they stole from Nondy if the adventurers drive off the dragon.

\subparagraph{Kobold Camp}
If the adventurers follow fleeing kobolds, they can find the kobolds' camp on a nearby outcropping. The kobolds there don't fight, though, instead begging the adventurers to help them (as in "Kobold Supplicants" above).

\subparagraph{Dragon's Lair}
Use the Volcanic Caves map in \textit{Appendix B} for the dragon's lair, but close off passages to keep the lair small. Near the entrance, the characters encounter a gang of four Magma Mephits and three Smoke Mephits, drawn to the lair by the magic of the dragon.

In the inner cave, the troublesome \textbf{Red Dragon Wyrmling} rests on its little hoard:


\begin{itemize}
\item Crate holding Nondy's stolen goods (worth 400 GP)
\item 4,200 CP, 2,000 SP, and 180 GP
\item Seven gemstones worth 50 GP each
\item \textit{Potion of Healing}
\item \textit{Rope of Climbing}
\item Two Spell Scroll (\textit{Alarm} and \textit{Comprehend Languages})
\end{itemize}

If the characters defeat or drive off the dragon, each of them earns a bonus 400 XP.

\subsection{Horns of the Beast}
\textit{Adventure for Level 5 Characters}


\begin{description}
\item[Situation.]
A long-forgotten ruin is rumored to hold a fiendish Artifact called the Horns of the Beast. A villain hopes to claim the Artifact and put it to terrible use.
\item[Hook.]
An unassuming human merchant named Melchis (secretly a Chaotic Evil \textbf{Fiend Cultist} devoted to \textit{Iuz}) hires the characters to escort him on an expedition to find an ancient temple lost in the jungle. He offers to pay them a total of 2,000 GP in trade bars—half when they reach the temple and half when they safely return to civilization—and promises to support them with what he claims to be "limited magical ability."
\end{description}

\subsubsection{Encounters}
You can use the poster map of the world of Greyhawk for this adventure; the temple is located near the southern edge of the map, in the Amedio Jungle. The journey unfolds in three stages (see "\textit{Travel}" in \textit{chapter 2}), culminating in the discovery of the ruins and the artifact.

\subparagraph{Stage 1}
Melchis hires a ship to carry the party to the Amedio Jungle. Use the map to determine how long the sea voyage takes, depending on where the adventurers begin the trip and figuring that the ship covers about 1½ hexes per day. Near the end of the trip, as the ship crosses Jeklea Bay, it's attacked by a group of Hostile [Attitude] sahuagin, including two Sahuagin Priests, six Sahuagin Warriors, and a \textbf{Water Elemental}. (Assume the ship's crew stays out of the way of these terrifying monsters and lets the characters and Melchis deal with this threat.) Stage 1 ends when Melchis and the characters row a launch ashore.

\subparagraph{Stage 2}
Melchis leads the characters into the jungle, aiming for the shore of the lake. This stage covers about 180 miles through the dense forest. Each day of the journey, roll once on the following table to determine what the characters encounter on their journey, if anything.


\begin{DndTable}{cX}

1d20 & Encounter\\
1–14 & No encounter\\
15 & An Indifferent [Attitude] \textbf{Giant Ape} protects its territory; its primary concern is getting the party to leave.\\
16 & A Hostile [Attitude] \textbf{Tyrannosaurus Rex} is on the hunt and tries to eat the characters.\\
17 & Three Allosauruses are hunting in the jungle; they are Hostile [Attitude] and treat the party as prey.\\
18 & Two Ankylosauruses tromp through the forest nearby. They are territorial and aggressive but Indifferent [Attitude], and they won't pursue a fleeing party.\\
19 & A band of humans, including a \textbf{Warrior Veteran} and eight \textbf{Warrior Infantry}, watch the party. They are Indifferent [Attitude]; they live in the jungle and aren't used to seeing other people.\\
20 & Four Minotaurs of Baphomet prowl the jungle looking for Humanoids they can capture and bring back to the temple.\\
\end{DndTable}

\subparagraph{Stage 3}
Once the characters reach the edge of the lake, Melchis leads them southwest along the shore until they find the ruin. This stage covers 90 miles through coastal terrain. The second day of the journey, \textit{heavy rain} obscures vision and creates \textit{quicksand pits} (see "\textit{Environmental Effects}" and "\textit{Hazards}" in \textit{chapter 3}). The front rank of the party might fall into a quicksand pit while two Giant Crocodiles attack at the same time. The rest of the journey passes without incident.

\subparagraph{The Ruins}
For the ruins, use the western half of the ground floor of the Dungeon Hideout map in \textit{Appendix B}, and ignore the stairs leading to the lower levels. Dwelling in the ruins are six Minotaurs of Baphomet who are Hostile to all intruders. They are spread out around the ruins but come quickly when they hear combat.

\subparagraph{The Artifact}
The Horns of the Beast—a jagged crown made from the horns of demons and wild animals—rests on a pedestal in the northwest corner of the ruins. Melchis immediately attempts to seize the Artifact and place it on his head. If he does, he is transformed into a \textbf{Hezrou} and tries to kill the characters. A character who dons the Artifact is cursed with \textit{Demonic Possession} (see "\textit{Curses and Magical Contagions}" in \textit{chapter 3}). The Artifact can't be removed from the character's head until a \textit{Remove Curse} spell is cast on the character. If the characters defeat Melchis and search him, they find a \textit{Bag of Holding} containing twelve 5-pound gold trade bars (worth 250 GP each) and a \textit{Spell Scroll} of \textit{Teleportation Circle}. The scroll also contains the sigil sequence for a permanent teleportation circle. At the end of the adventure, each character earns a bonus 1,000 XP.

\subparagraph{Destroying the Artifact}
The Horns of the Beast can be destroyed only by dropping it in the \textit{River Oceanus}, which flows through the Upper Planes (see \textit{chapter 6}). A character can learn this by casting \textit{Identify} on the Artifact.

\subsection{Boreal Ball}
\textit{Adventure for Level 7 Characters}


\begin{description}
\item[Situation.]
The Baron of the Boreal Ball, a minor noble of the Feywild, holds an unending ball in his ice palace. Behind the revelry, the ball is the scene for schemes and intrigues.
\item[Hook.]
The adventurers receive a magical invitation to the Boreal Ball that teleports the group to the ball at the appointed time.
\end{description}

\subsubsection{Encounters}
The adventure plays out over three dances. During each one, the characters can decide whether they dance, mingle, watch the dancers, or engage in some other activity. (The Manor map in \textit{Appendix B} can serve as a floor plan for the baron's palace.) Each character has the opportunity to earn renown for the party (see "\textit{Renown}" in \textit{chapter 3}) by making a positive impression on whomever they're interacting with—a potential Renown Score increase of 1 per character per dance. These guests pay particular attention to the characters:


\begin{itemize}
\item Cannifer is a \textbf{Satyr Revelmaster} who is used to being the center of attention at every ball he attends and is thus Hostile [Attitude] to the adventurers.
\item Darisis a Friendly [Attitude], fun-loving \textbf{Dryad} who wants to be seen dancing with all the adventurers.
\item Fidget is a playful \textbf{Pixie} who is Indifferent [Attitude] and pesters the adventurers with pranks throughout the evening, without causing any actual harm.
\item Granny Snailtongue is a Hostile [Attitude] \textbf{Green Hag} who sees the adventurers as potentially useful tools. She offers to help them at every turn in hopes of putting them in her debt.
\item Raxas Albrethin is an arrogant, Chaotic Neutral drow \textbf{Mage} who is initially Hostile [Attitude] and wants to see the adventurers humiliated. However, once the characters' Renown Score reaches 6, Raxas admits he misjudged them and becomes Friendly [Attitude].
\end{itemize}

\subparagraph{Rude Interruption}
After the second dance, the ball is interrupted by a Neutral Evil \textbf{Hobgoblin Warlord} named Varka, who is accompanied by a \textbf{Hobgoblin Captain} and five Hobgoblin Warriors. The hobgoblins are offended by the adventurers' presence and attack them. If the characters defeat the hobgoblins, their Renown Score increases by 2.

\subparagraph{Conclusion}
After three dances, the Baron of the Boreal Ball appears. If the party's Renown Score is at least 6, he bestows on each character a Charm of the Boreal Ball. This charm (see "\textit{Supernatural Gifts}" in \textit{chapter 3}) allows a character who has it to cast the level 3 version of the \textit{Ice Knife} spell. Once used, the Charm vanishes. Each character also earns a bonus 1,700 XP.

\chapter{Creating Campaigns}
If encounters are the building blocks of a D\&D adventure, then adventures are the building blocks of a D\&D campaign, for a campaign is what you get when you string two or more adventures together. A campaign setting is the world in which those adventures take place—both a backdrop for your adventures and a hotbed of conflicts and personalities that can inspire and drive adventures.

\section{Step-by-Step Campaigns}
Follow these steps to create a campaign:


\begin{description}
\item[Step 1:]
Lay Out the Premise. Consider the core conflicts driving the campaign, and choose a setting that reinforces the themes and tone you hope to evoke.
\item[Step 2:]
Draw In the Players. Start your campaign in a memorable way. Determine how the characters get drawn into events and how the characters' goals and ambitions might come into play.
\item[Step 3:]
Plan Adventures. Consider the smaller conflicts that make up the larger conflicts of the campaign, and devise fun quests that help drive the story. Flesh out the antagonists, the important locations, and the elements that link the adventures together.
\item[Step 4:]
Bring It to an End. Think about how the campaign might end and what level you expect the characters to be when the campaign wraps up.
\end{description}

You might have noticed that these steps are similar to the "\textit{Step-by-Step Adventures}" list at the start of \textit{chapter 4}. In many ways, a campaign is just an adventure writ large. In an ongoing campaign, one adventure flows naturally into the next.

Later sections of this chapter offer inspiration and advice for each of these four steps. The chapter concludes with a campaign example.

\section{Your Campaign Journal}
At the start of any campaign, there's a buzz of excitement as you and your players look forward to creating a new world together—one full of adventure and promise. Every game session is a chance for you to show off more of the campaign setting and deepen your players' investment in it.

If your campaign lasts for months or years, sustaining that high level of excitement—yours as well as your players'—takes effort. An important tool to help you keep interest in the campaign high is a campaign journal, a collection of notes from past sessions. Use your journal to refresh your memory on events that transpired early in the campaign and bring closure to unresolved conflicts and mysteries.

\subsection{Keeping a Journal}
A campaign journal documents the progression of your campaign, from the first game session to the last. Your journal can take whatever form works best for you. It might be a physical notebook; a binder of loose notes, maps, and tracking sheets; a wiki; or a collection of files on your computer. Journal entries are best organized by date or game session. (Some DMs prefer the term "episode" to "game session," but the terms are interchangeable.)

A sample Campaign Journal page is provided. Make copies of it, or use it as inspiration for your own journal pages.

\subsection{Using Your Journal}
Use your journal to plan out your next game session (see "\textit{Preparing a Session}" in \textit{chapter 1}). Then, when the game session is over, use the journal to capture anything else of importance that might have bearing on future sessions, such as the name of an NPC you created on the fly or a critical piece of information the characters learned.

During a game session, you can use your campaign journal to quickly recall a piece of information you've forgotten (such as the name of a character's mule) or to jot down things you want to remember later (such as the name of a tavern). In this way, the journal becomes a living chronicle of the campaign in flight.

\subsubsection{Foreshadowing}
Foreshadowing is a storytelling technique that never goes out of style. Players love it when something happens in a game session that hearkens to some event from an earlier session.

Foreshadowing is about planting seeds early so you can reap the rewards later. Having an up-to-date campaign journal makes foreshadowing easier because you can reread your notes from earlier game sessions and identify things that could resurface in upcoming sessions, giving past events greater weight or a bigger payoff. Consider the following example.

The characters find the dead body of an unidentified halfling adventurer. A search of the body yields a cameo necklace containing the portrait of another halfling. A character decides to keep the cameo, which was intended as a bit of embellishment. You make a note of it in your journal. Months later, while planning a future session, you flip through the journal and are reminded of the cameo. It inspires you to plan a chance encounter with another halfling, whom the characters might recognize as the one depicted in the cameo. What happens if the characters return the cameo to this halfling? This halfling could be tied to a bigger plot or have information that could help the characters resolve some conflict. Suddenly, a minor trinket foreshadows bigger events to come.

\subsubsection{Adventure Stockpile}
Besides tracking each session of your campaign, keep a list of adventure ideas. Even if you don't end up using every adventure idea, having a stockpile will keep you ready for whatever your players throw at you, and you can even borrow pieces of various ideas to incorporate into future adventures. Not every adventure needs to build on earlier plots; a good stand-alone adventure tucked in the middle of a serialized campaign can be a welcome change of pace for you and your players.

\section{Campaign Premise}
Everything outlined about the story of an adventure in \textit{chapter 4} is true of a campaign's story as well: a campaign is like a series of comics or TV shows, where each adventure (like an issue of a comic or a TV episode) tells a self-contained story that contributes to the larger story. Just like with an adventure, a campaign's story isn't predetermined, because the actions of the players' characters will influence how the story plays out.

\subsection{Campaign Characters}
The characters are the focus of every D\&D adventure, and their players are your partners in developing their characters' epic journeys.

By working with your players to understand what excites them most, you can craft stories they want to see their characters star in. You can also more effectively draw players into adventure plots (see "\textit{Draw In the Players}" in \textit{chapter 4}) if you understand what motivates both them and their characters.

\subsubsection{Player Input}
It's not up to you to create every aspect of a D\&D campaign. Players contribute through their characters' actions and by directly sharing what they want to see in a campaign. You can learn about your players' preferences in two ways:


\begin{description}
\item[Direct Input.]
Ask your players what they want to do in a campaign. Regularly inquire about how they think the campaign is going, what they'd like to experience more of, and what elements they'd like to explore further. After a session concludes and between sessions are great times to ask players for thoughts about the campaign.
\item[Indirect Input.]
The choices a player makes, starting at character creation, can indicate what they want to see in the game. For example, a Rogue player likely wants opportunities for subtlety or skulduggery, while a Barbarian player likely craves combat. Take note of what encounters players are enthusiastic about, and seek ways to help the players' characters shine.
\end{description}

\subsubsection{Character Arcs}
Like most protagonists in film and literature, D\&D adventurers face challenges and change through the experience of overcoming them. By incorporating each character's motivations into your adventures and setting higher stakes through play, you'll help characters grow in exciting ways. You can use the \textit{DM's Character Tracker} sheet to keep track of key information about each character. See "\textit{Getting Players Invested}" in this chapter for more ideas.

\subparagraph{Character Motivations}
For each character, think about what motivates them to adventure. Motivations generally fall into the following categories:


\begin{description}
\item[Goal.]
A character's goal is a short-term reason for the character to adventure. At the start of a campaign, this might be a desire for treasure, a thirst for excitement, or some need from a character's backstory. As characters continue to adventure, they'll find different goals to pursue, such as finding a lost relic, honoring an ancestor, avenging a fallen mentor, or defeating a villain.
\item[Ambition.]
A character's ambition is a broad, personal aspiration the character hopes to achieve through a lifetime of adventuring. A character might dream of becoming a legendary knight or bringing peace to their homeland. Ambitions might be unrelated to the character's current goal.
\item[Quirks and Whims.]
Quirks and whims are a character's preferences, impulses, or other traits. They often emerge during play, such as a character's tendency to one-up a rude innkeeper or their oft-expressed fondness for displacer beast fur.
\end{description}

Players often reveal their characters' motivations through play. If you're uncertain or a character's motivations seem to have changed, it's OK to ask players for clarification.

\subparagraph{Family, Friends, and Foes}
A character's origin (species and background) implies some amount of backstory, suggesting the character's family and what the character did before becoming an adventurer. Take note of specific background characters—friends, foes, family members, and others—who might appear in the campaign.

Should these background characters become important to the campaign, work with the player to develop them in detail. Revealing a character's lost sibling or childhood rival midcampaign should be handled carefully to avoid straining credulity. Make sure a player is comfortable with new developments about their character before introducing them.

\subparagraph{Character-Focused Adventures}
Adventures should occasionally highlight character motivations or elements of their backstory. Here are a few examples of character-focused adventures:


\begin{itemize}
\item A rival from a character's past shows up to settle a grudge.
\item A sneaky character puts their skills to the test by leading the rest of the party to conduct a heist.
\item A character learns the location of a magic item needed to save their hometown.
\item A spellcasting character must undertake a trial to join an exclusive group of spellcasters.
\end{itemize}

Any adventure that focuses on a single character should incentivize the whole party to participate—even if just to help their companion. Avoid focusing adventures on one character too often, and look for opportunities to have character-focused adventures for each character from time to time.

\subparagraph{Setting New Goals}
Characters can change their goals whenever they please, but you can encourage them to do so by giving them significant victories roughly every 5 levels. When characters accomplish their goals, consider the following questions:


\begin{itemize}
\item How does completing this goal create a new challenge?
\item How is this victory only part of what the character wants to achieve?
\item Who might be upset by the character completing this goal?
\item What is a reward the character will be excited to receive that also moves them closer to their ambition?
\end{itemize}

Use the answers to these questions to develop new character goals and to inspire further adventures.

\subparagraph{Building on the Characters' Actions}
Sometimes it can be fun to let the players steer the campaign by having their characters' actions dictate future adventures. For example, if the characters buy a tavern using the treasure they've amassed, you can adjust the campaign so that the tavern has a role in future adventures. One adventure might involve a competitor trying to put the characters' tavern out of business. Another might use the tavern as the setting for a murder mystery.

\subsection{Campaign Conflicts}
One way to ensure your campaign's longevity is to come up with three compelling conflicts you can create adventures around. Introduce these conflicts early in the campaign. As the campaign unfolds, focus adventures on different conflicts to keep the players' excitement high.

Use the Campaign Conflicts tracking sheet to record your campaign's conflicts (with room to add details or notes). A conflict can be as big or as small as you like, and it's nice to have at least one conflict that can be resolved quickly. Each conflict should involve the adventurers against some antagonistic force, though you can also create conflicts between two powerful forces without necessarily knowing which force (if either) the adventurers will align themselves with. The "Flavors of Fantasy" section below provides examples of conflicts that reinforce particular themes.

If a conflict reaches a satisfying end before the end of the campaign, create a new conflict to replace it. You can also replace conflicts that don't resonate with your players as well as conflicts you're having trouble building adventures around.

\subsubsection{Conflict Arcs}
In the same way you think about character arcs over the course of a campaign, think about how each conflict might manifest over the course of the campaign. How do the characters first encounter the conflict? How does the conflict develop over time? What might a climactic ending to that conflict look like?

One helpful way to structure a conflict arc is to use the \textit{tiers of play} described in \textit{chapter 4}. Levels 5, 11, and 17 represent milestones in character power and capabilities, and they can also be story milestones in the arc of your campaign. The shift from one tier to another is an ideal time to wrap up a campaign conflict and introduce a new one that has a broader reach and represents a greater threat. The threshold of a new tier can also be an opportunity for characters to realize the scale of a conflict they've been dealing with—to realize, for example, that the bandits they fought throughout their first four levels are merely puppets of an enemy nation they must confront in the second tier.

The "\textit{Greyhawk}" section in this chapter has examples of conflict arcs.

\subsection{Flavors of Fantasy}
Your D\&D campaign might be inspired by a particular flavor of fantasy, several of which are discussed in the sections that follow. Any of these fantastical subgenres can be informed and inspired by the cultures, myths, legends, and fantasies of any culture: an epic fantasy campaign could draw on French romances or Chinese wuxia stories, a mythic fantasy campaign could be based on Greek myth or the \textit{Epic of Gilgamesh}, and so on.

\subsubsection{Heroic Fantasy}
Heroic fantasy features adventurers bringing magic to bear against monstrous threats—the default subgenre presented in the core D\&D rulebooks.

\subparagraph{Heroic Fantasy Conflicts}
Heroic fantasy campaigns often revolve around delving into ancient dungeons in search of treasure or to destroy monsters or villains. Consider conflicts like these to drive the action of a campaign:


\begin{description}
\item[Evil Cult.]
Wicked cultists infiltrate a peaceful realm to free an ancient evil entity trapped in a dungeon. Releasing the entity would surely spell the realm's doom.
\item[Fungal Plague.]
To protect a primeval forest from the encroachment of hunters and settlers, druids unleash a fungal plague that quickly gets out of hand.
\item[Old Enemy.]
An elusive villain who plagued the characters years ago resurfaces, giving the characters a chance to finally bring the villain to justice.
\end{description}

\subsubsection{Sword and Sorcery}
A sword-and-sorcery campaign features a grim world of evil spellcasters and decadent cities, where the protagonists are often motivated more by greed and self-interest than by altruistic virtue.

\subparagraph{Sword-and-Sorcery Conflicts}
In this flavor of campaign, magic-users often symbolize the decadence and corruption of civilization, and mages are the classic villains of these settings. Magic items are therefore rare and often dangerous. Consider conflicts like these to drive the campaign:


\begin{description}
\item[Evil Adventurers.]
An evil band of experienced adventurers wields power and influence to oppress hapless folk.
\item[Evil Weapon.]
A knight under the influence of a sentient, evil weapon terrorizes a peaceful realm. Cultists worship and protect this weapon, which must be seized and destroyed to end the threat.
\item[Forgotten Dynasty.]
The long-lost seat of a forgotten dynasty rises from the sea or the desert sands, and its people launch a campaign of conquest.
\end{description}

\subsubsection{Epic Fantasy}
An epic fantasy campaign emphasizes the conflict between good and evil, with the adventurers on the side of good. These heroic characters are driven by a higher purpose than selfish gain or ambition. Characters might struggle with moral quandaries, fighting the evil tendencies within themselves as well as the evil that threatens the world. And the stories of these campaigns often include an element of romance: tragic affairs between star-crossed lovers, passion that transcends even death, and chaste adoration between knights and nobles.

\subparagraph{Epic Fantasy Conflicts}
Conflicts like these highlight the themes of an epic fantasy campaign:


\begin{description}
\item[Apocalypse.]
A prophecy predicts the end of the world unless the adventurers intervene. Apocalypse cultists oppose the characters at every turn.
\item[Dragon Tyrant.]
An evil and powerful dragon moves into the region, upsetting the ecology and demanding tribute from nearby settlements.
\item[The Foe Time Forgot.]
An evil foe believed long dead emerges from the Feywild, alive and well after being lost in time. This foe seeks revenge against the descendants of long-dead enemies.
\end{description}

\subsubsection{Mythic Fantasy}
A mythic fantasy campaign draws on the themes and stories of ancient myth and legend, from Gilgamesh to Cú Chulainn. Adventurers attempt mighty feats of legend, aided or hindered by the gods or their agents—and the characters might have divine ancestry themselves. The monsters and villains they face might have a similar origin. The chimera in the dungeon isn't just a random beast but the product of a divine curse.

\subparagraph{Mythic Fantasy Conflicts}
Conflicts like these highlight the themes of a mythic fantasy campaign:


\begin{description}
\item[Divine Trials.]
Seeking a gift or favor from the gods, the adventurers undertake a series of trials that lead them to the realms of the gods, where the adventurers can plead their case.
\item[Divine Wrath.]
After a temple is sacked, a vengeful god sends an escalating series of woes upon a kingdom until the temple's relics are returned.
\item[Giants.]
An enormous castle on a cloud settles over the land. The characters can battle the giants living there or try to broker a lasting peace.
\end{description}

\subsubsection{Supernatural Horror}
If you want to put a horror spin on your campaign, the \textit{Monster Manual} is full of creatures that suit a storyline of supernatural horror. An essential element of such a campaign is an atmosphere of dread, created through careful pacing and evocative description. Your players contribute too; they must be willing to embrace the mood.

Whether you want to run a full-fledged horror campaign or a single creepy adventure, discuss your plans with the players ahead of time. Horror can be intense and personal, and not everyone is comfortable with such a game. (The advice on discussing limits under "\textit{Ensuring Fun for All}" in \textit{chapter 1} is particularly important for a horror game.)

\subparagraph{Supernatural Horror Conflicts}
A supernatural horror campaign often features Undead or demonic foes whose evil transcends the merely mortal. Consider conflicts like these to drive the campaign:


\begin{description}
\item[The Faceless Lord.]
\textit{Juiblex}, the Faceless Lord, oozes out of the \textit{Abyss} and into the \textit{Underdark}. The characters hear from subterranean folk who need help defeating the demon lord and its minions.
\item[School of Necromancy.]
Vampires open a college of necromancy, attracting evil necromancers who need fresh corpses for their studies. An order of vampire hunters seeks the characters' help.
\item[Undying Monarch.]
A venerable monarch clings to power by worshiping \textit{Orcus} and becoming a \textbf{lich}.
\end{description}

\subsubsection{Intrigue}
Political intrigue, espionage, sabotage, and similar cloak-and-dagger activities can provide the basis for an exciting campaign. In this kind of game, the characters might care more about skill proficiencies and making friends in high places than about attack spells and magic weapons. Social interaction takes on greater importance than combat. Make sure your players know ahead of time that you want to run this kind of campaign. Otherwise, a player might create a combat-focused character, only to feel out of place among diplomats and spies.

\subparagraph{Intrigue Conflicts}
Conflicts like these are ripe for an intrigue campaign:


\begin{description}
\item[Feuding Fiefs.]
Two fiefs or settlements have been feuding for years. The characters are drawn into the ongoing feud after helping one side.
\item[Royal Rivals.]
The sudden death of a sovereign plunges a kingdom into chaos when the rightful heir is challenged and threatened by rivals.
\item[Scheming Adviser.]
After a monarch takes an interest in the characters, they become targets of the monarch's most trusted adviser, who is scheming to become the true power in the realm.
\end{description}

\subsubsection{Mystery}
A mystery-themed campaign puts the characters in the role of investigators, perhaps traveling from town to town to crack tough cases that local authorities can't handle. Such a campaign emphasizes puzzles and problem-solving in addition to combat prowess. An adventure composed of nothing but puzzles can become frustrating, so be sure to mix up the kinds of encounters you present.

\subparagraph{Mystery Conflicts}
A mystery might set the stage for the whole campaign. The characters might uncover clues to this mystery from time to time, while individual adventures might be only tangentially related to it. Consider conflicts like these for a mystery campaign:


\begin{description}
\item[Criminal Syndicate.]
A many-headed criminal syndicate seeks economic and political power. The syndicate has spies everywhere, including among the adventurers' families or friends.
\item[Shape-Shifting Assassins.]
A secret association of doppelgangers or other shape-shifters slowly assassinates prominent figures one by one.
\item[To Catch a Thief.]
An extraordinary thief steals only the most valuable jewelry and works of art. The characters might become a target of the thief when they acquire a priceless treasure.
\end{description}

\subsubsection{Swashbuckling}
The swashbuckling adventures of pirates and musketeers make for a dynamic campaign in which dashing, charming heroes weave their way through palace intrigues and leap from balconies onto waiting horses to escape dogged pursuers. In a swashbuckling campaign, the characters typically spend a lot of time in cities, in royal courts, and aboard seafaring vessels. Nevertheless, the heroes might end up in classic dungeon situations, such as escaping from a prison cell block or searching storm sewers to find a villain's hidden chambers.

\subparagraph{Swashbuckling Conflicts}
Conflicts like these highlight the themes of a swashbuckling campaign:


\begin{description}
\item[Inherited Antagonists.]
A character inherits a magic item from a deceased relative, unaware that this relative's enemies are after the item.
\item[Pirates and Privateers.]
A new monarch cracks down on piracy by commissioning privateers and naval officers to hunt pirate ships.
\item[The Waking Deep.]
A monstrous horror slumbering in the depths of the ocean stirs, driving minions such as \textbf{sahuagin}, merrows, or dragon turtles to attack seafaring vessels.
\end{description}

\subsubsection{War}
A campaign focused on warfare centers on heroes whose actions turn the tide of battle. The characters carry out specific missions: capture a magical standard that empowers undead armies, gather reinforcements to break a siege, or cut through the enemy's flank to reach a demonic commander. The party might also support the larger army by holding a strategic location until reinforcements arrive, killing enemy scouts, or cutting off supply lines. Information-gathering and diplomatic missions can supplement combat-oriented adventures.

\subparagraph{War Conflicts}
Conflicts like these highlight the themes and flavor of a war campaign:


\begin{description}
\item[Freedom Fighters.]
Poorly armed and disorganized subjects of a tyrant revolt.
\item[Invaders.]
A militaristic nation invades its benevolent neighbors.
\item[Pawns in a Game.]
A war rages on for decades, its original cause all but forgotten. The people caught up in it strive to find meaning and purpose in a bleak and violent world.
\end{description}

\subsubsection{Crossing the Streams}
Deep in D\&D's roots are elements of science fiction and science fantasy as well as a wide-ranging collection of fantasy inspiration, and your campaign might draw on those sources as well. You can send your characters hurtling through a magic mirror to Lewis Carroll's Wonderland, put them aboard a ship traveling between the stars, or set your campaign in a far-future world where laser weapons (see "\textit{Firearms and Explosives}" in \textit{chapter 3}) and Wand of Magic Missiles exist side by side.

\subparagraph{Crossing the Streams Conflicts}
Conflicts like these create opportunities for crossing the streams:


\begin{description}
\item[Beyond the Magic Mirror.]
A mysterious mirror in a strange dungeon is a portal into a different world where whimsical tales unfold—or perhaps some version of the modern world.
\item[Gamma World.]
The characters inhabit a post-apocalyptic wasteland that is largely medieval in feel, but isolated outposts still hold futuristic technology from before the cataclysm.
\item[Invaders from Wildspace.]
Spaceships land on the characters' world and disgorge hostile creatures armed with advanced technology.
\end{description}

\subsection{Campaign Setting}
Just like an \textit{adventure's setting} (as described in \textit{chapter 4}), a campaign setting is an essential part of a campaign's premise, shaping the kinds of stories that unfold there.

As the DM, you have two options when choosing a campaign setting:


\begin{itemize}
\item Use a published campaign setting.
\item Create your own campaign setting.
\end{itemize}

Whether you create your own campaign setting or use a published one, the world of your game is always your own. You can customize it to suit your tastes and those of your players.

\subsubsection{Using a Published Setting}
One advantage of using a published campaign setting is that much of the world-building is done for you. However, this means your players might know as much about the setting as you do. You can get around this by changing key aspects of the setting to better serve your needs, which has the added benefit of challenging your players' expectations.

The D\&D Settings table describes several established campaign settings.

\begin{DndTable}[header=D&D Settings]{XX}

Setting & Description\\
Dark Sun & Heroes make their mark on a postapocalyptic world defiled by magic and forsaken by the gods.\\
Dragonlance & The forces of good battle the evil queen of dragons and her armies in the world-shaking War of the Lance.\\
\textit{Eberron} & In the aftermath of a deadly war, magically advanced nations rebuild as a cold war threatens lasting peace.\\
\textit{Exandria} & Heroes make names for themselves in the world made popular by the streaming show Critical Role.\\
Forgotten Realms & Larger-than-life heroes and villains struggle to determine the fate of the world as they explore the ruins and dungeons of fallen kingdoms and long-forgotten empires.\\
\textit{Greyhawk} & As tensions rise among warring nations, heroes plunder dungeons to gain the magic and might they need to defeat the growing forces of evil.\\
\textit{Planescape} & Sigil, the City of Doors, is where heroes begin to explore the wonders of the D\\&D multiverse and its many planes of existence.\\
\textit{Ravenloft} & Heroes are drawn into the gloomy Domains of Dread—cursed realms ruled by evil lords—and must find a means of escape.\\
\textit{Ravnica}* & In a world-spanning city, ten disparate factions draw heroes into a web of adventure and danger.\\
\textit{Spelljammer} & Travel among the stars on a spelljamming ship, and visit worlds floating in the majestic oceans of Wildspace.\\
\textit{Strixhaven}* & Strixhaven, a school of magic, serves as a hub of learning and adventure.\\
\textit{Theros}* & Heroic destinies wait to be fulfilled in this setting inspired by the myths of ancient Greece.\\
\end{DndTable}

\subsubsection{Creating Your Own Setting}
One advantage of creating your own world is it can be whatever you want it to be. Your players will never know more about the world than you do, which can be both a comfort to you and a source of wonder to your players. Moreover, you don't need to memorize any source material about the campaign setting, other than what you create for yourself.

Whether you create a setting from scratch or borrow elements from established settings, the result needs to resonate with your players. As you create your world, ask your players what settings and genres they enjoy, then use those sources for inspiration to create compelling locations, memorable inhabitants, exciting conflicts, and an internal logic that will resonate with your players.

\subparagraph{Five Questions to Consider}
As you contemplate a new campaign setting, think about your answers to the following questions:


\begin{description}
\item[What's Your Campaign Setting Called?]
Choose an evocative name for your setting. It can be a word or phrase that reflects the theme and tone of the game, or just a made-up name that sounds cool to you. Keep a running list of ideas as you decide on other aspects of your setting.
\item[What Factions and Organizations Are Prominent?]
Nations, temples, guilds, orders, secret societies, and colleges shape the social fabric of the setting. What organizations or societal groups play an important part in your setting? Which ones might be involved in the lives of player characters as patrons, allies, or enemies? What organizations can characters join, becoming part of something larger than themselves?
\item[How Common Is Magic?]
Spellcasters and magic item shops might be common, rare, or practically nonexistent in your world. How readily available are spells such as \textit{Lesser Restoration}, \textit{Raise Dead}, and \textit{Teleportation Circle}? Is magic so widespread that it's part of daily life, or so rare that it conjures all sorts of superstitions?
\item[What Mysteries Does the World Hold?]
Every campaign setting has mysteries: a fabled land across the sea, a grim forest hiding a terrible secret, restless spirits haunting a ruined keep for reasons unknown, an ancient dungeon built for a forgotten purpose, and so on. Dream up as many mysteries as you wish—you never know which ones will seize your players' imaginations and become central to the campaign—and record them in your campaign journal.
\item[What Roles, If Any, Do the Gods Play?]
What greater gods, lesser gods, and quasi-deities are present or worshiped in your world? If there are gods, how involved are they in the world? Are they distant and detached beings, or do they appear before their worshipers and meddle in mortal affairs?
\end{description}

\section{Campaign Start}
With your campaign journal in hand and the basic premise of your campaign (characters, conflicts, and setting) in mind, it's time to consider how to begin the campaign.

\subsection{Session Zero}
At the start of a campaign, you and your players can run a special session—called session zero because it comes before the first session of play—to establish expectations, share ideas, and discuss house rules, with the goal of ensuring the game is a fun experience for everyone involved. The "\textit{Ensuring Fun for All}" section in \textit{chapter 1} covers some of the most important groundwork you need to establish at the start of a new campaign.

Often session zero includes building characters together. As the DM, you can help players during character creation by advising them on which options best suit the campaign.

\subsubsection{Character Creation}
When players are choosing their characters' classes and origins, you can restrict options that are unsuitable for the campaign.

Encourage the players to choose different classes so that the adventuring party has a range of abilities. It's less important that the party include multiple backgrounds or species; sometimes it's fun to play an all-Dwarf party or a troupe of adventuring Entertainers.

The origins the players choose define who their characters were before becoming adventurers. Think about how the characters' backgrounds might inform adventures in your campaign. For example, if a player chooses the Criminal background, help the player flesh out their character's criminal past, and use that information when building relevant storylines into the larger campaign.

\subparagraph{Starting Level}
What level are the characters when they start? Many D\&D campaigns start the characters at level 1. If you want the characters to be a bit more resilient and your players are experienced, start the campaign at level 3 instead. (See the \textit{Player's Handbook} for rules on \textit{starting at higher levels}.)

\subsubsection{Bringing the Party Together}
During session zero, help the players come up with explanations for how their characters know each other and have some sort of history together, however brief that history might be. To get a sense of the party's relationships, here are some questions you can ask the players as they create characters:


\begin{itemize}
\item Are any of the characters related to each other?
\item What keeps the characters together as a party?
\item What does each character like most about each member of the party?
\item Does the group have a patron—an individual or organization that points them toward their adventures?
\end{itemize}

If the players are having trouble coming up with a story for how their characters met, you can suggest the following options.

\subparagraph{Bonding Event}
Some bonding event (such as a wedding, a festival, or a funeral) brings the characters together, whereupon they quickly discover a shared sense of purpose.

\subparagraph{Happenstance}
Someone puts out a call for adventurers to complete a quest, and the characters answer the call. Alternatively, all the characters could meet by accident, only to discover they're headed to the same place, or they could find themselves trapped together.

\subparagraph{Mutual Acquaintance}
The characters are introduced to one another by a mutual NPC acquaintance whom they all trust. This shared acquaintance could serve as a patron for the party—perhaps a representative of an organization (an academy, a criminal syndicate, a guild, a military force, or a religious order), a politically powerful person (an aristocrat or even a sovereign), or a magical creature like a sphinx or a dragon.

\subparagraph{Shared History}
The characters grew up in the same place and have known one another for years. Despite their different backgrounds and training, they're already good friends.

\subparagraph{Tavern Gathering}
The characters meet in a tavern over mugs of ale and decide to embark on a life of adventure together—a tried and true trope!

\subsubsection{Setting the Stage}
Session zero is a great time to share basic information about the campaign with your players. Such information typically includes the following:


\begin{description}
\item[Starting Location Details.]
Your players need basic information about the place where the characters are starting, such as the name of the settlement, important locations in and around it, and prominent NPCs they'd know about (see "Starting Location").
\item[Key Events.]
Describe any current or past events that help frame the campaign. For example, the campaign might start on the heels of a great war or on the day of a festival. Describing key events helps set the mood and prepare players for upcoming adventures.
\item[House Rules.]
If you're using any \textit{house rules} (as discussed in \textit{chapter 1}), or adopting any of the variant rules presented in this or any other book, let your players know about them.
\end{description}

Remember, you'll always know more about your campaign world than the players do. Having spent all their lives in this world, though, the characters also know more than their players do. Fill in the basics of what the characters should know anytime that information matters to their adventures.

\subsection{Starting Location}
Begin your campaign in a location you can detail, such as a village, a neighborhood in a larger city, an outpost, or a roadside tavern. Be prepared to give players enough information about that location to help them figure out what ties, if any, their characters have to it. Once you have this campaign hub fleshed out, create one or two local attractions that might serve as adventure locations, such as a haunted house on the outskirts of town or a dungeon complex tucked in the nearby hills.

If you're using a published campaign setting, pick any location in that setting and develop it as you like. A published setting or adventure might give you all the details you need. \textit{The Free City of Greyhawk}, described later in this chapter, is an ideal starting location and illustrates the kinds of things to consider as you detail a starting location.

If you're building your own setting, start small by detailing only this starting area. The rest of your setting can remain undeveloped for now. Don't spend too much time fleshing out the geopolitical landscape of your world or locations the adventurers aren't likely to visit right away; save those fun tasks for when you and your players have a better sense of where the campaign is headed.

\subsection{First Adventure}
If you're using a published adventure to launch your campaign, use the character hooks in that adventure to bring the characters from their starting location to the adventure's action. Many campaigns begin with a published adventure and then develop organically as the characters explore beyond the scope of the adventure.

If you're creating your own adventure for the start of your campaign, refer to the advice in \textit{chapter 4}. Keep the first adventure relatively short and simple, allowing plenty of time for the characters to get to know each other as the players roleplay. What's most important is that they begin to feel like an adventuring party and get comfortable with their abilities. The full scope of the campaign can unfold to them later.

\section{Plan Adventures}
A D\&D campaign is like a garden. Each new adventure plants new seeds in the garden, which requires regular tending lest it run wild. Over time, your campaign will grow and flourish in ways you expected and in ways that will surprise you. You might need to weed out elements that aren't resonating with your players while planting new elements to tantalize them.

Most D\&D campaigns grow organically, rather than having all their elements set in stone from the get-go. From time to time, the characters' decisions will require you to improvise and create new campaign elements on the fly. For example, a new location might need to be developed to address the needs of the unfolding story, or certain NPCs might need fleshing out at a moment's notice. Other parts of this book, such as the "\textit{Nonplayer Characters}" and "\textit{Settlements}" sections in \textit{chapter 3}, can help you expand your campaign quickly.

\subsection{Episodes and Serials}
There are two basic ways to think about how adventures fit together in your campaign: as distinct episodes or as a serialized story. If you're not sure which type of campaign to run, ask your players what they prefer. If your players have different preferences, you can intersperse episodic, stand-alone adventures among serialized adventures to break up the bigger story.

\subsubsection{Episodes}
An episodic campaign is a campaign in which the component adventures don't combine to form an overarching story. Episodic adventures are stand-alone quests, and the villains who appear in one adventure rarely resurface to trouble the characters again. If your game group plays infrequently, an episodic campaign might be ideal because the players can enjoy the current adventure even if they've forgotten the details of earlier adventures.

\subparagraph{Starting a New Episode}
In an episodic campaign, the start of a new adventure doesn't necessarily have any connection to the end of the last one. The action might pick up immediately after the end of the previous adventure, but it might instead begin weeks, months, or years after the last adventure, allowing interim events to unfold while the characters take a break from adventuring.

\subsubsection{Serials}
A serialized campaign is one continuous story broken up into smaller parts that flow naturally from one to the next. It often has one or more overarching threats, and the outcome of one adventure can affect how the rest of the campaign unfolds. If your game group meets regularly and often, a serialized campaign allows you to keep your players guessing what will come next as the campaign builds toward a satisfying conclusion.

\subparagraph{Linking Adventures}
In a serialized campaign, make connections between the end of one adventure and the start of the next to help it feel like a connected story. Sometimes you can simply continue the current storyline with new locations to explore and new threats to overcome. Alternatively, you can use the Adventure Connections table to inspire a link from one adventure to the next. The table suggests things you can do near the end of one adventure to lead characters into the next one.

\begin{DndTable}[header=Adventure Connections]{cX}

1d6 & Adventure Connection\\
1 & Introduce a person, an object, or information that the characters need to transport safely to a location involved in the new adventure.\\
2 & Have a major villain flee to a location that features in the new adventure. The characters might be able to pursue the villain, or they might have to search for clues about where the villain has gone.\\
3 & Introduce clues suggesting that a villain or another NPC in this adventure is part of a larger group—a group that features prominently in the new adventure.\\
4 & Introduce a villainous group that's featured in the new adventure by having its agents spy on or interfere with the characters' activities.\\
5 & Have travelers bring news of events transpiring elsewhere, leading characters toward the new adventure.\\
6 & Give the characters a treasure that's wrapped in mystery they'll need to unravel in the new adventure.\\
\end{DndTable}

\subsection{Getting Players Invested}
To get your players excited about and invested in your campaign, create a setting that features people and places they recognize and where their characters' choices matter.

The following sections suggest ways to help you create a world your players will be excited to explore.

\subsubsection{Recurring Elements}
When characters form relationships—friendships, business arrangements, or even lasting antagonism—with the people and places of your setting, those people and places stick in the players' minds. Introduce opportunities to forge these lasting relationships early and often.

Consider featuring recurring elements such as these in your game:


\begin{description}
\item[Community.]
Introduce a small group or community the characters can think of as their people, like a village, neighborhood, guild, or crew.
\item[Home Base.]
Give the characters a place to call home, such as a tavern, a hideout, or a ship. Bastions, as presented in \textit{chapter 8}, are ideal home bases for characters.
\item[Prominent Friend.]
Create a supportive NPC whom the characters can trust and turn to when they need help, such as a local leader, an innkeeper, a patron, a retired adventurer, or a family member.
\item[Friendly Resources.]
Provide experts or institutions that can assist the characters, like a temple that can provide healing or a learned sage who can help solve mysteries.
\item[Likable Villain.]
Craft a villain who has at least one likable or redeeming quality the characters can appreciate—ideally a villain who isn't preoccupied with killing or harming the characters.
\end{description}

As your campaign continues, introduce new people and locations, and bring back favorites from earlier in the campaign for the occasional cameo.

\subsubsection{Player Favorites}
It's often easier to describe people and places that are hostile or frightening than it is to detail a feature you want characters to love. How can you know what rustic scene will make a character associate a place with home or what personality quirk will remind a character of their favorite mentor? You can ask a character's player directly, but instead consider handing over your narrative reins and letting a player describe the perfect detail.

For example, say you have a peaceful village you plan to feature across several adventures. You hope the characters will connect with the place and treat it as home. As the characters enter the community, they smell something amazing. At this point, you could describe something you think smells good or something you think a character would like. Or you can ask a player, "The smell of something amazing drifts from around the corner. What is it?" Whatever the player's answer—cinnamon rolls from a nearby baker, firework charges being prepared for a celebration, or anything else—becomes part of the village, and the player has added an important detail to the location.

You can use player input whenever you want to pinpoint something meaningful to the characters and their players. Consider asking players questions like these whenever you want to describe something in an impactful way:


\begin{itemize}
\item The tavern owner brings out your favorite dish—cooked to perfection. What's the dish, and what makes this one remarkable?
\item The curio shop is selling a trinket that reminds you of one of your family members. What's the trinket, and who does it remind you of?
\item The local children are playing a game you played in your hometown. What is it?
\item The young pickpocket reminds you of someone you once knew. Who?
\item From the animate mass of murderous dolls scrambles a figure that reminds you of your favorite childhood toy. What is it?
\end{itemize}

Questions such as these don't need to draw on warm memories. Having players describe what unsettles or disgusts their characters can make menacing encounters more impactful as well. In any case, take note of interesting character details that your players share, and record them in your campaign journal, as these details might be useful inspiration for later adventures.

\subsubsection{Acknowledge the Incredible}
Adventurers are, by their nature, remarkable. Even at level 1, they perform miraculous deeds and possess qualities that set them apart from common folk. Reinforce this in your game. NPCs don't need to gush over the characters, but the characters' reputations as heroes, problem-solvers, or wonderworkers should be cemented early and develop throughout a campaign.

During every session, look for opportunities to make the characters feel like the stars of the story, and try to answer one or more of the following questions:


\begin{itemize}
\item How are the characters the perfect people to solve a problem?
\item How are the characters' talents highlighted during the adventure?
\item What stories do NPCs know of the characters' past exploits?
\item How might an NPC comment on a character's abilities or recognize that they're special?
\end{itemize}

\subsubsection{Break Episodes}
It's easy to get caught up in a story with dramatic stakes, pitting characters against mounting threats. But every so often, at least once every three to five levels, give the characters a break—a low-stakes session or adventure that has nothing to do with the overarching plot or broader perils.

A break episode can be an opportunity for the characters to reflect on the events of the ongoing campaign, explore the nuances of the world, and further develop the relationships between them in a more relaxed setting. Give the group space to breathe, note developments you want to highlight later, then continue with your adventures.

Consider these ideas for a break episode.

\subparagraph{Bastions Episode}
The characters take a break from adventuring to tend to their Bastions (see \textit{chapter 8}), with players taking one or more Bastion turns and describing what happens.

\subparagraph{Carnival Episode}
A carnival tempts the characters with magical attractions, games, and prizes. The Witchlight Carnival, described in The Wild Beyond the Witchlight, is one such carnival.

\subparagraph{Creature Comedy}
The characters encounter monsters with a comedic flavor—such as flumphs, pixies, faerie dragons, or chatty mimics—in a situation that leads to mischief and humor rather than combat.

\subparagraph{Missing Pet Episode}
Someone's pet is missing. The characters must search a settlement and connect with locals to help find it.

\subparagraph{Shopping Episode}
A friendly NPC asks the characters to help shop for someone's birthday.

\subparagraph{Special Event Episode}
The characters are invited to a sporting event, holiday celebration, fancy dinner, or ball.

\subparagraph{Vacation Getaway}
The characters relax on a quiet beach, enjoy the comforts of a grateful noble's villa, withdraw to a serene monastery, or while away the hours in a fairy hot spring.

\subsection{Time in the Campaign}
Most conflicts in a D\&D campaign take weeks or months of in-world time to resolve. A typical campaign concludes within a year of in-world time unless you allow the characters to enjoy lengthy periods of quiet time between adventures.

If you don't want to track the passage of days, weeks, and months, you might instead track the passage of time using seasons and seasonal festivals. The answer to the question "When does this adventure take place?" can be as simple as "in the winter" or "during the fall harvest festival."

\subsubsection{Timed Events}
Extraordinary events coinciding with certain times of year make for great adventure opportunities. Perhaps a ghostly castle appears on a certain hill on the winter solstice every year, or every thirteenth full moon is blood red and fills werewolves with a particularly strong bloodlust. The appearance of a comet in the sky might portend all manner of significant events. The festivals of the gods can serve as opportunities to launch adventures, especially if the gods themselves are involved.

\section{Ending a Campaign}
A campaign's ending should conclude the last of the major conflicts and tie up most of the threads of its beginning and middle. (It's OK to leave some loose ends for characters to explore in the next campaign.) You don't have to take a campaign all the way to level 20 for it to be satisfying; wrap up the campaign whenever the story reaches its natural conclusion.

Allow time near the end of your campaign for the characters to finish up any personal goals. Their stories need to end in a satisfying way, just as the campaign story does. Ideally, some of the characters' individual goals will be fulfilled by the final adventure. Give characters with unfinished goals a chance to finish them before the very end.

Once your campaign has ended, a new one can begin. If you intend to run a new campaign for the same group of players in the same setting, using their previous characters' actions as the basis for legends is one way to invest your players in the new campaign. Let the new characters experience how the world has changed because of the actions or accomplishments of the previous campaign's characters. In the end, though, the new campaign is a new story with new protagonists. They shouldn't have to share the spotlight with the heroes of days gone by.

\subsection{Ending Sooner Than Expected}
Sometimes you run out of ideas for your campaign, or it gets so sidetracked that you have no idea how to bring it to a satisfying conclusion. You might just not feel excited about it anymore, or you might be so excited with ideas for a new campaign that you can't focus on the current one. Any of these might signal the end of your campaign.

The best way forward when you want to end a campaign is to talk to your players about it. If you're not excited about the game anymore, it's quite possible they're not either, and you can change or end the campaign to everyone's satisfaction. Consider the following possibilities:


\begin{description}
\item[Player Input.]
If you're running out of ideas for your campaign, your players might be more than happy to supply you with some. Find out what they'd like to have happen if the campaign continues. They might give you all the inspiration you need!
\item[Switch DMs.]
One of your players might have so many ideas about the future of the campaign that they're willing to take over as the DM. You can either take over that player's character or make a new one of your own. Let go of your plans for where the story was going, and allow the new DM to have creative control.
\item[Transport the Characters.]
If you or another DM wants to start up a campaign in a new setting but the players don't want to make new characters, consider having the characters travel through a portal to a new world.
\item[Arrange a Grand Finale.]
Sometimes an end to the campaign is the right answer. Look for ways to end the campaign with a bang, even if it's earlier than you originally planned. Flip through your campaign journal to see if there are forgotten elements you can resurface for one last hurrah.
\end{description}

\section{Greyhawk}
Greyhawk is a D\&D setting you can use as the backdrop for your campaign or as a model you can reference while creating your own setting. Important aspects of Greyhawk are described herein so that you can make it your own, expanding or altering it however you wish.

Greyhawk is the invention of Gary Gygax, one of the D\&D game's original creators. Gary based many of D\&D's earliest adventures in this home-brewed setting. The version of Greyhawk presented here is largely based on \textit{The World of Greyhawk} gazetteer, published in 1980.

\begin{DndSidebar}[float=!b]{Poster Map}
Included is a poster map showing the lands of Eastern Oerik on one side (each hex on the map is equivalent to 30 miles) and the Free City of Greyhawk on the other. These locations are described in the "\textit{Free City of Greyhawk}" and "\textit{Greyhawk Gazetteer}" sections later in this chapter.



\end{DndSidebar}

\subsection{Important Names}
A handful of important names are defined below:

\textbf{Oerth} (pronounced \textit{orth} or \textit{oyth}) is the world of Greyhawk. It has four continents, four oceans, and a plethora of islands and seas.

\textbf{Oerik} (pronounced \textit{or}-ick or \textit{oy}-rick) is one of Oerth's continents.

\textbf{Eastern Oerik}, the vast region explored in this chapter, is home to many powerful nations and some of the D\&D game's most famous dungeons and adventurers.

\textbf{The Flanaess} (pronounced flah-\textit{nay}-ess or \textit{flay}-nayz) is another name for Eastern Oerik and means "land of the Flan." The region's first human settlers and their descendants are known as the Flan.

\textbf{Greyhawk} is an independent city in Eastern Oerik that attracts large numbers of adventurers. Greyhawk doubles as the name of the campaign setting.

\subsection{Greyhawk's Premise}
The year is 576 CY (Common Year). Evil is ascendant across the lands of Eastern Oerik. If something isn't done to curtail the growing threat, Eastern Oerik will fall to tyrants, evil dragons, and monstrous hordes. Heroes are needed to bring hope to the people of the Flanaess. Even if the heroes die trying, the legends of their exploits will live on!

Adventuring parties from the Free City of Greyhawk and other settlements trek across the vast wilderness of Eastern Oerik, slaying monsters and exploring dungeons to find magic items the adventurers can use to defend their homeland and take the fight to their enemies.

\subsubsection{Greyhawk Conflicts}
Although Greyhawk lends itself well to any D\&D adventure you might want to run, the default setting features conflicts with three major villainous groups: chromatic dragons, Elemental Evil cults, and Iuz and his followers. You can replace one or more of these conflicts with ones of your devising or with ones from the "\textit{Flavors of Fantasy}" section earlier in this chapter.

If you use these conflicts, look for opportunities in your adventures to introduce creatures in service to the three villainous groups. Give goals to these villains that bring their operatives into conflict with the player characters.

The three major conflicts and the goals of the villainous groups are described below.

\subparagraph{Chromatic Dragons}
Evil chromatic dragons dwell in the wilds of Eastern Oerik. For years, adventurers have kept these evil dragons at bay, sometimes with the help of benevolent metallic dragons. Lately, the chromatic dragons have grown restless, their dreams invaded by the whispers of \textit{Tiamat}, who is trapped in the \textit{Nine Hells}. The five-headed queen of dragons believes her escape is nigh, and from the depths of her prison, she commands her kin to go forth and claim the world of Oerth for themselves. Only the greatest among them will live to become her consorts.

\subparagraph{Goals of Chromatic Dragons}
Fortify their lairs to safeguard their treasure hoards; strike out across Eastern Oerik, raiding poorly defended settlements and stealing cattle; demand tribute in the form of food or treasure; and destroy territorial rivals (draconic or otherwise).

\subparagraph{A Chromatic Dragon Arc}
The conflict between adventurers and chromatic dragons might follow this broad outline:


\begin{description}
\item[Levels 1–4.]
Consider introducing this conflict as the adventurers reach level 3 or 4, with the adventurers confronting an aggressive chromatic dragon wyrmling. (You can use the adventure "\textit{The Winged God}" from \textit{chapter 4}.)
\item[Levels 5–10.]
The adventurers might face a handful of ambitious young chromatic dragons, without hinting at a more significant conflict.
\item[Levels 11–16.]
It eventually becomes clear that the behavior of the adult dragons the characters face isn't normal. The characters might get involved in one dragon's schemes to undermine or overthrow another, or the characters might hear whispers of the dragons' dream of liberating Tiamat.
\item[Levels 17–20.]
The conflict reaches its worldshattering conclusion, with ancient dragons threatening nations and clashing with each other in devastating battles. The campaign might end with Tiamat herself appearing in the Flanaess—perhaps emerging from the Riftcanyon (see "\textit{Mysteries of Greyhawk}" in this chapter) or from the depths of the Nyr Dyv.
\end{description}

\subparagraph{Elemental Evil}
"Elemental Evil" is the name given to a host of destructive, extraplanar entities—demon lords, evil elemental princes, and elder gods—who ravaged the world of Oerth long ago. Many of these entities are now trapped in dungeons, with cults and monsters seeking to free and serve them. Adventurers are the only ones equipped to keep these malign entities from escaping their subterranean prisons.

\subparagraph{Goals of Elemental Evil}
Search for a demon lord, an elemental prince, or an elder god trapped in a dungeon; build a stronghold above or near the dungeon; drive other inhabitants out of the region; and use a special magic item or ritual to free whatever is trapped in the dungeon.

\subparagraph{An Elemental Evil Arc}
Two published adventures have explored campaign arcs centered around Elemental Evil. The Temple of Elemental Evil (published in 1985) is set in the world of Greyhawk, in the Kron Hills between Verbobonc and Celene. It begins in the unremarkable village of Hommlet, with characters slowly discovering that agents of Elemental Evil have taken up residence in the nearby ruins and are working to rebuild their great temple. The rest of the adventure focuses on exploring the ruined Temple of Elemental Evil and dealing with the varied evil factions and forces within (including agents of Iuz, described below).

Princes of the Apocalypse, inspired by \textit{The Temple of Elemental Evil}, presents an alternative arc for an Elemental Evil–themed campaign:


\begin{description}
\item[Levels 1–4.]
The characters discover the villainous activity of four elemental cults.
\item[Levels 5–10.]
The characters strike at the four headquarters of these evil cults while investigating the cults' activities in the surrounding region. While the characters are battling one cult in its headquarters, the other three cults might still be wreaking havoc nearby, forcing the characters to divide their attention.
\item[Levels 11–16.]
Finally, the characters discover an ancient Temple of the Elder Elemental Eye deep beneath the cults' separate temples, and they strive to contain the damage wrought by the cults' activities and thwart the cults' evil leaders before these leaders unleash an apocalypse.
\end{description}

Though this adventure is set in the world of the Forgotten Realms, it includes notes on how you might transplant it into Greyhawk or any other setting.

\subparagraph{Iuz the Evil}
North of the Free City of Greyhawk, a demigod named Iuz has reclaimed the vast tract of land he lost after being imprisoned under Castle Greyhawk by the archmage Zagig Yragerne. The newly freed Iuz aims to lay waste to the kingdoms, steadings, temples, and outposts of his rivals. To that end, Iuz's spies are searching for powerful Artifacts they can use to ensure victory, while evil creatures spawn in Iuz's homeland and threaten neighboring realms. Adventurers can thwart Iuz by keeping evil Artifacts out of his hands and defeating the vile creatures that serve him.

\subparagraph{Goals of Iuz}
Install loyal operatives in settlements across Eastern Oerik, search libraries and vaults for lore pertaining to ancient and powerful magic, scour dungeons for lost Artifacts and other magic items, secure such items, and use magic and monsters to conquer rival nations.

\subparagraph{An Iuz Arc}
The conflict between adventurers and Iuz might follow this broad outline:


\begin{description}
\item[Levels 1–4.]
Early in their adventuring careers, the characters might face what appear to be ordinary toughs who are disrupting mining operations near the Free City of Greyhawk (see "\textit{Beyond the City Walls}" in this chapter), only to discover these toughs are agents of some greater villain. The identity of this villain remains a mystery—for now. If you use the adventure "\textit{Miner Difficulties}" from \textit{chapter 4}, NPCs speaking to the characters might assume the trouble in the mine is related to these toughs and their bullying.
\item[Levels 5–10.]
You might use the adventure "\textit{Horns of the Beast}" from \textit{chapter 4} to introduce an agent of Iuz to the characters. After their return from that expedition, they start having unpleasant encounters with the City Watch in Greyhawk. Eventually, they discover that Captain-General Sental Nurev is being manipulated by the leaders of Stoink, a petty fief in the Bandit Kingdoms. When the characters undertake an expedition into that dangerous realm to confront Stoink's leaders and free the captain-general's captive brother, they discover that the villains were agents of Iuz.
\item[Levels 11–16.]
Iuz and the Horned Society launch an all-out invasion into the Shield Lands, overwhelming its defenses and moving toward Furyondy. The characters might have adventures to muster forces in surrounding lands and bring them to Furyondy's defense or hinder Iuz's advance.
\item[Levels 17–20.]
Finally, the characters discover that Iuz's assault is merely a cover to distract the southern realms from his true aim: retrieving the \textit{Eye and Hand of Vecna} from an ancient keep on Lake Quag. The characters confront Iuz at the shores of the lake, perhaps facing a terrible choice: Will they wield the power of Vecna to stop Iuz, or will they risk Iuz wielding that awful might against them?
\end{description}

\subsubsection{The Greyhawk Setting}
The planet Oerth is at the very center of a Wildspace system called Greyspace. (See \textit{chapter 6} for more information about \textit{Wildspace}). Oerth has two moons: Luna (a great white moon, also called the Mistress) and Celene (a smaller blue moon, also called the Handmaiden). Greyspace's sun orbits Oerth, rather than the other way around.

The sun takes 360 days to travel once around Oerth. Luna waxes and wanes in fixed cycles of 28 days each, upon which the months are based, while Celene follows a path that has full moons only four times each year, coinciding with four lunar festivals.

\subparagraph{Months and Festivals}
The standard year is 360 days long and consists of twelve twenty-eight-day months (each month divided into four seven-day weeks) and four six-day lunar festivals (Needfest, Growfest, Richfest, and Brewfest). The midwinter festival of Needfest is considered the start of the year. The diagram here shows the months and festivals that make up a year.

\begin{DndTable}[header=Days of the Week]{ccccccc}

Starday & Sunday & Moonday & Godsday & Waterday & Earthday & Freeday\\
(Saturday) & (Sunday) & (Monday) & (Tuesday) & (Wednesday) & (Thursday) & (Friday)\\
Day of Worship & Day of Rest\\
\end{DndTable}

\begin{DndSidebar}[float=!b]{Your World's Calendar}
What does your campaign's calendar look like? The more your campaign calendar resembles the one that's familiar to you and your players, the easier it will be to remember and use. Familiar names for months and days of the week lend your campaign a wonderful simplicity that many players will appreciate. A calendar that uses ten-day weeks or names such as "Moonday" and "Coldeven" is harder for players to internalize but reminds them they're in a fantasy world. The Greyhawk calendar has twelve months and seven-day weeks like the Gregorian calendar, but gives unique names to months and days and introduces festivals that fall outside the calendar's months, giving it a fantastical feel.



\end{DndSidebar}

\subparagraph{Factions and Organizations}
In addition to the many political entities that dot the lands of the Flanaess and the temples of its many gods, several organizations operate across national borders in pursuit of their goals. Some of these organizations could serve as patrons or allies of adventurers in a Greyhawk campaign, while others might appear as villains. Some might even accept adventurers as members.

\subparagraph{Circle of Eight}
Some of the greatest spellcasters of the world of Greyhawk form the Circle of Eight, a group dedicated to preserving balance in the world. The group's general aim is to prevent any single country, faction, or other organized group from becoming too powerful and overwhelming others. The membership of the Circle of Eight is secret but includes \textit{Mordenkainen} (the strategist behind the group), \textit{Bigby}, \textit{Jallarzi Sallavarian}, \textit{Otiluke}, and \textit{Otto}.

\subparagraph{Knights of the Watch}
The order of the Knights of the Watch originated as a military force protecting the lands of Bissel, Gran March, Geoff, and Keoland from hostile neighbors to the north and west (particularly Ket and the Ulakandar). Though Watchers maintain strongholds along the border with Ket, most of their energy is spent defending against giants and dragons in the western mountains. The Watchers are sworn to an ascetic and disciplined code, and they train rigorously to the exclusion of personal property or other attachments.

\subparagraph{Order of the Hart}
The knights of the Order of the Hart were organized to preserve the freedom of the states of Furyondy, Veluna, and Highfolk against the threats of bandits and hostile neighbors. These nations have little centralized authority or military power, so the knights have historically served as a first line of defense against these varied threats. In recent years, they have mobilized against the rising threats of \textit{Iuz} and Elemental Evil, forcing them to broaden the scope of their operations into neighboring realms where these evils are active.

\subparagraph{Scarlet Order}
The Scarlet Order is a monastic order of Suloise militarists whose spies and assassins have infiltrated many courts and castles throughout the Flanaess, ready to strike. The leader of the order is a seemingly immortal being known as the Father of Obedience, Korenth Zan. He is rumored to be a Suloise monk who walked the lands of Oerik long before the Rain of Colorless Fire destroyed the Suloise Empire. Others claim Korenth is a red dragon—a former consort of \textit{Tiamat} who became trapped in human form. Whatever the true story, the Father of Obedience is revered by all who pledge their lives to the Scarlet Order. His goals—and, by extension, the order's goals—are shrouded in mystery and could one day tilt the balance of power across the whole of Eastern Oerik.

\begin{DndSidebar}[float=!b]{Adventurers and Organizations}
Factions and organizations aimed at player characters can connect adventurers to your world, providing ties to key NPCs and a clear agenda beyond individual gain. In the same way, villainous organizations create an ongoing sense of menace beyond the threat of solitary foes.



Having different characters tied to different factions can create interesting situations at the gaming table, as long as those factions have similar goals and don't work in opposition to one another all the time. Adventurers representing different factions might have competing interests or priorities while they pursue the same goals.



Adventurer organizations are also a great source of special rewards beyond Experience Points and treasure. Increased standing in an organization might come with concrete benefits such as access to an organization's information, equipment, magic, and other resources. See "\textit{Renown}" in \textit{chapter 3} for rules you can use to track characters' standing in an organization.



\end{DndSidebar}

\subparagraph{Magic in Greyhawk}
In the world of Greyhawk, as in most D\&D worlds, magic is widespread but still wondrous and sometimes frightening. People everywhere know about magic, and most people see evidence of it at some point in their lives. Magic permeates the cosmos and moves through the ancient possessions of legendary heroes, the mysterious ruins of fallen empires, those touched by the gods, creatures born with supernatural power, and individuals who study the secrets of the multiverse. Histories and fireside tales are filled with the exploits of those who wield magic.

What normal folk know of magic depends on where they live and whether they know people who practice magic. Citizens of an isolated hamlet might not have seen true magic used for generations except the strange powers of the old hermit living in the nearby woods, which they regard with suspicion and mention only in whispers.

By contrast, magic is common enough in the Free City of Greyhawk that the Guild of Wizardry teaches magic and sells spellcasting services. Extensive codes of law govern the use and abuse of magic. The law treats magical coercion as a major crime, and punishes the public use of magic in situations that could harm people or property.

\subparagraph{Mysteries of Greyhawk}
Eastern Oerik is a realm of many mysteries, several of which are described below.

\textbf{Bat-Folk of Hepmonaland}. Separated from Eastern Oerik by the Tilva Strait, Hepmonaland is a relatively small continent that few people of the Flanaess known much about. Those who have explored the north spur of Hepmonaland report dense rainforests, severe tropical storms, steamy wetlands, and a fetid swamp (called the Pelisso Swamp). Adventurers are sometimes lured into Hepmonaland's rainforests by ancient ruins, including tombs and shrines left behind by an ancient civilization of bat-like humanoids whose history is largely forgotten.

\subparagraph{Devastating Magic}
Almost a thousand years ago, the war between the Baklunish and Suloise empires came to a horrific end. The Baklunish people who lived in what is now the Dry Steppes called down a rain of colorless fire that burned all living things, ignited the landscape, and reduced the Suloise lands to ashes, creating the Sea of Dust. In retaliation, Suloise survivors invoked their own magic to devastate the Baklunish lands. What magic was responsible for the Rain of Colorless Fire and the Invoked Devastation? What would happen if such magic fell into the wrong hands today?

A central portion of the Dry Steppes, where the seat of the Baklunish empire stood, is said to remain pleasant and rich, roamed by Baklunish nomads. The former Suloise capital, by contrast, in the heart of the Sea of Dust, is beset by howling winds, terrible dust storms, and rains of volcanic ash and cinders from the nearby Hellfurnaces.

\subparagraph{Land of Black Ice}
Those who have ventured far north of the Burneal Forest tell of a strange phenomenon. Instead of normal stark-white snow and translucent blue-white ice, there is an endless landscape of deep-blue ice partially covered in snow. Strange arctic monsters prowl these fields of dark ice. Stranger still, a verdant land is rumored to exist beyond the ice, where the sun never sets.

\subparagraph{Riches of the Bright Desert}
The Bright Desert, walled off from the rest of the Flanaess by the monster-infested hills of the Abbor-Alz and the aptly named Gnatmarsh, is supposedly filled with copper, silver, gold, and precious stones. The harsh climate, wildly varying temperatures, and hostile inhabitants discourage exploration. Expeditions have attempted to penetrate the Bright Desert and extract its riches, but none have ever returned.

\subparagraph{Riftcanyon}
Between the Bandit Kingdoms and the Shield Lands stretches a deep canyon, ten miles wide at the ends, thirty miles wide at its midsection, and 180 miles long. The Riftcanyon, which is more than a mile deep, is home to at least one blue dragon and has tunnels near its base that lead to the Underdark.

\subparagraph{White Plume Mountain}
Situated just south of the Riftcanyon, the ever-smoking White Plume Mountain has always been a subject of superstitious awe to the neighboring villagers. People still travel many miles to gaze upon this natural wonder, though few dare to approach it closely, as it is reputed to be the haunt of demons and devils. The occasional disappearance of those who stray too close to the Plume reinforces this belief.

White Plume Mountain is detailed in \textit{Tales from the Yawning Portal}.

\subparagraph{Gods of Greyhawk}
The Gods of Greyhawk table shows many of the most popular deities worshiped in the Flanaess. Greater gods and demigods are marked as such; the others are lesser gods. Many other deities and demigods are also worshiped in the Flanaess, beyond those shown on the table. Some deities of Greyhawk have also transcended their origin on this world to impact the broader multiverse. Two of these, \textit{Tharizdun} and \textit{Vecna}, are described in \textit{appendix A}.

The greater gods of Greyhawk rarely get directly involved with happenings on Oerth. Lesser gods are more likely to manifest in some form on the Material Plane and interact with their worshipers. Cuthbert, for example, is well known for appearing in mortal guise, appearing as a dirt-covered farmer, a wanderer robed in brown and green, or an elderly tinker. And of the many quasi-deities that appear on Oerth, most prominent among them is Iuz, a demigod who rules his own nation in the Flanaess.

\begin{DndTable}[header=Gods of Greyhawk]{XXXX}

Name and Epithet & Home Plane & Typical Worshipers & Symbol\\
Beory, Heart of Oerth* & \textit{Material Plane} & Farmers, herders & Green disk\\
Berei of the Hearth & \textit{Bytopia} & Families, farmers & Sheaf of wheat stalks\\
Boccob the Uncaring, Archmage of the Gods* & \textit{Outlands} & Sages, spellcasters, seers & Eye within a pentagram\\
Celestian, the Far Wanderer & \textit{Astral Plane} & Wanderers, astronomers & Arc of seven stars inside a circle\\
Cuthbert of the Cudgel & \textit{Arcadia} & Practical, honest folk & Circle at the center of a starburst of lines\\
Ehlonna of the Forests & \textit{Beastlands} & Hunters, foragers & Unicorn horn\\
Erythnul, the Many & \textit{Pandemonium} & Raiders, bandits, berserkers & Blood drop\\
Fharlanghn, the Dweller on the Horizon & \textit{Outlands} & Travelers & Circle crossed by a curved horizon line\\
Heironeous the Invincible & \textit{Mount Celestia} & Knights, soldiers & Lightning bolt\\
Hextor, Scourge of Battle & \textit{Acheron} & Soldiers, tyrants & Six down-pointing arrows in a fan\\
Incabulos, the Black Rider* & \textit{Hades} & Necromancers, those who seek to ward off illness & Reptilian eye within a horizontal diamond\\
Istus, Weaver of Our Fate* & \textit{Mechanus} & Seers, advisers & Spindle with three strands\\
Iuz the Evil† & \textit{Material Plane} & His subjects and allies & Grinning human skull\\
Kord, the Brawler & \textit{Ysgard} & Athletes, berserkers & Spears and maces radiating from a point\\
Nerull, the Reaper* & \textit{Carceri} & Murderers, necromancers & Skull with a scythe\\
Obad-Hai, the Shalm & \textit{Outlands} & Hunters, gatherers, hermits & Oak leaf and acorn\\
Olidammara, the Laughing Rogue & \textit{Ysgard} & Revelers, gamblers, pranksters & Laughing mask\\
Pelor, the Radiant Sun* & \textit{Elysium} & Healers, the compassionate & Sun\\
Pholtus of the Blinding Light & \textit{Arcadia} & Judges, lawyers, arbiters & Silver sun partially eclipsed by a crescent moon\\
Ralishaz, the Unlooked For & \textit{Limbo} & Gamblers & Three bone fate-casting sticks\\
Rao, the Mediator* & \textit{Mount Celestia} & Mediators, sages, scientists & White heart\\
Syrul Oathbreaker & \textit{Gehenna} & Liars, charlatans, traitors & Forked tongue\\
Tharizdun, the Eater of Worlds & Imprisoned in a demiplane & Nihilistic cultists & Spiral rune\\
Trithereon, the Summoner & \textit{Arborea} & Rebels, individualists & Triskelion\\
Ulaa, the Bejeweled & \textit{Arcadia} & Miners, jewelers, quarriers & Ruby-hearted mountain\\
Vecna, the Whispered One & Unknown & Necromancers, undead, those who keep or unearth secrets & An eye in the palm of a left hand\\
Wee Jas, the Witch & \textit{Acheron} & Spellcasters, advisers & Red skull in front of fireball\\
\end{DndTable}

\subsection{Free City of Greyhawk}
Would-be heroes are drawn to the Free City of Greyhawk by promises of adventure. The city is rife with opportunities for peril and plunder.

The city stands on the eastern banks of the Selintan River. The river flows south from the Nyr Dyv (the Lake of Unknown Depths) down to Woolly Bay and remains easily navigable for its entire length.

Once a frontier hub of the Great Kingdom of Aerdy, Greyhawk proclaimed itself free and independent seventy-eight years ago, claiming the Selintan basin as its territory. Adventurers drawn to the nearby ruins of Castle Greyhawk have provided a steady influx of cash to the city in the years since.

Throughout this section, if a creature's name appears in \textbf{bold} type, you'll find that creature's stat block in the \textit{Monster Manual}. If a creature's alignment isn't specified, you can decide what it is.

\subsubsection{Start Here}
The City of Greyhawk is a great starting point for a D\&D campaign for many reasons, as discussed in the sections that follow.

\subparagraph{Adventure Hooks}
The city contains plenty of rumors, local legends, and quest givers, any of which could point characters to their next adventure. The \textit{sample adventures} in \textit{chapter 4} can all begin in the Free City of Greyhawk.

\subparagraph{Bastion Friendly}
There are ample places within the city and on the city's outskirts where adventurers can build Bastions (see \textit{chapter 8}).

\subparagraph{Key Conflicts}
Two of the three central conflicts of the Greyhawk setting—the threats of Elemental Evil and Iuz—are the source of major tension and intrigue in the Free City of Greyhawk. The third conflict, involving evil dragons, is a looming threat nearby, particularly in the Mistmarsh, the hills of the Abbor-Alz, and the Bright Desert.

\subparagraph{Local Hurly-Burly}
Greyhawk has a frontier spirit atypical for a settlement of its size. The locals are a tough and rowdy lot. Adventurers seeking action don't need to look far, as the city contains more than its fair share of troublemakers.

\subparagraph{Nearby Attractions}
North of the city are the Cairn Hills, which are known to have tombs and dungeons hidden among them. Nearby are a forest, a swamp, and a desert. The monsters that haunt these areas tend to be weak—perfect for testing the mettle of low-level adventurers.

\subparagraph{Port in the Storm}
The city provides a place to rest, heal, acquire information, and procure magic items. Adventurers looking to visit distant lands can book passage on ships docked at the wharf.

\subparagraph{Trade Hub}
Adventurers can buy gear and sell their hard-won loot in the city's shops and markets.

\subsubsection{How to Use the City}
A bustling city like Greyhawk can serve the following important functions in a campaign.

\subparagraph{Background Connections}
Use the backgrounds of the characters to connect them to people and places in the city. These connections help the players feel like Greyhawk is their characters' home—or will quickly become their new home. A character who was born and raised in Greyhawk might have known one of the city's prominent figures for many years, while someone who has just arrived in the city might have a mutual friend with that connection or might carry a letter of introduction recommending them to that person. (Each location detailed in the "\textit{City Locations}" section includes potential connections for two character backgrounds.)

\subparagraph{City Activities}
Greyhawk is an ideal place for activities that support adventuring. There's endless opportunity for social interaction in such a bustling place, as well as places where characters can rest and recuperate between adventures, acquire new adventuring gear, and spend their gold.

\subparagraph{Home Base}
As a home base for characters, Greyhawk can serve as a place to live, train, and recuperate between adventures. As described under "\textit{Getting Players Invested}" in this chapter, Greyhawk offers a host of potential friends, rivals, villains, and resources. Use the people and locations mentioned in this chapter as a starting point for fleshing out characters' connections to the city, and work with your players to develop those connections. Choose one neighborhood of the city (see "\textit{City Neighborhoods}" below) as a focus for the characters and their activities.

\subparagraph{Urban Adventures}
A city isn't just a place to spend time between adventures—plenty of adventures happen within the city walls. From wererats in the sewers to scheming bureaucrats in the halls of power, dangers lurk around every corner.

\subsubsection{City Overview}
The city is yours to make your own. A few important features and locations are described in the sections that follow, but otherwise flesh out the city as you and your players see fit.

\subparagraph{City Government}
The Free City of Greyhawk is ruled by a council called the Directing Oligarchy, made up of sixteen coequal rulers. This council elects its chief officer, the lord mayor—a position currently held by a human \textbf{Spy Master} (Lawful Neutral) named Nerof Gasgol. The other directors include the captain-general and constable of the City Watch, several guild masters, priests of Boccob and Rao, the inspector of taxes, and a few influential magic-users with ties to politically active secret societies. Several of these directors represent criminal or unsavory interests, including Nerof Gasgol himself, who achieved his position and wealth as the owner of a notorious gambling den.

\subparagraph{City Watch}
The City Watch is a standing garrison of some eight hundred Guards and Veteran Warriors. Bolstering these defenders are Mages from the city's Guild of Wizardry, as well as Priests from local temples.

The captain-general and constable of the City Watch are stationed at the Grand Citadel (see "\textit{City Locations}").

\subparagraph{City Walls}
A 30-foot-high stone wall winds like a snake around the city. Two other walls, identical in height to the outer wall, separate the city into its three great sections. Access to the wall tops can be gained via lifts in each gatehouse. In addition, along the inside base of the outer wall are secret compartments at 300-foot intervals, each one containing a 30-foot-tall wooden ladder. All members of the City Watch know the locations of these secret ladders, which, in an emergency, can be pulled out and used by city defenders to quickly reach the parapets.

The walls are patrolled regularly. During daytime, the typical patrol is one sentry (a \textbf{Guard}) placed every 300 feet along the top of the wall. At night, the guard patrol is quadrupled, with two sentries posted together every 150 feet along the wall. Also at night, torches light the wall top at 150-foot intervals between the guards so each sentry station is 75 feet from a torch in each direction.

\subparagraph{City Gates}
Each city gate consists of a pair of iron-reinforced wooden doors that can be barred from the inside. These heavy doors are backed by a massive portcullis of iron bars. A very small child might be able to squeeze between the bars, but not a youth or even an adult halfling. The city's portcullises are usually left open even when the gates are closed.

Each gate is contained within a small gatehouse flanked by a pair of towers. The tower tops and connecting blockhouse are equipped with arrow slits and holes for pouring boiling oil straight down onto invaders. Each gatehouse tower connects to the city through a door in its base and to the wall top by a door in its side. The towers contain three platforms, beginning at the top of the wall and extending upward. Each of these can shelter and provide a firing platform for up to forty archers.

Three of the city's gates typically remain open throughout the day and night: the Highway Gate (the grand entrance to the city), the Cargo Gate (used primarily by traders and merchants), and the Garden Gate (one of the city's two inner gates). The remaining gates are closed from dusk until dawn, and a visitor must produce a written message from the lord mayor of Greyhawk, the captain-general of the City Watch, or a head of state to be allowed through. In the latter case, the guards ensure the traveler is harmless before opening the gates.

Those passing through open gates aren't asked to explain their business, nor are they detained or turned back unless they are recognized as known fugitives. Wagons and carts might get searched if they trigger the guards' suspicions, but most vehicles are waved through without inspection. The guards keep a daily roster of who and what pass through their gates.

\subparagraph{Crime}
The Free City of Greyhawk is home to many thieves, vandals, charlatans, and hooligans. Crimes are divided into three categories.

\subparagraph{Petty Crime}
Public unarmed brawling, pickpocketing, vandalism, and other crimes that cause up to 50 GP in property damage are petty crimes. The perpetrator pays a fine of 2d10 GP or works to provide restitution.

\subparagraph{Minor Crime}
The category of minor crimes includes armed assault (defined as any nonfatal attack made with a weapon or damaging spell) and property crimes that cause between 50 and 250 GP in damages. The perpetrator must pay a fine of at least 100 GP and is sentenced to 1d6 years in prison.

\subparagraph{Major Crime}
Crimes more severe than those described above—including murder, bribery or impersonation of a city official, and magical coercion—are major crimes. The criminal faces 2d10 years of imprisonment, the death penalty, or permanent exile. A city magistrate decides which punishment is appropriate.

\subparagraph{Religion}
The city has temples and shrines dedicated to various gods. Religious practices that are certifiably evil aren't tolerated, however. When an evil sect is discovered in the city, its wealth is confiscated, its leaders are put to death, and all other members are banished from the city for life.

See the "\textit{Gods of Greyhawk}" table for many deities worshiped by the inhabitants of the Free City.

\subsubsection{City Neighborhoods}
The Free City of Greyhawk is split into three main sections by two internal walls running west to east. The northern section is home to the High Quarter and the Garden Quarter, where the wealthiest folk of the city reside. The central section is home to the River Quarter, Clerkburg, the Artisans' Quarter, and the Foreign Quarter. The southern portion, known as the Old City, includes the poorer and rowdier neighborhoods of the Slum Quarter and the Thieves' Quarter.

Brief descriptions of the city's neighborhoods are presented below:


\begin{description}
\item[Artisans' Quarter.]
The Artisans' Quarter is built around a large marketplace. The finest artisans live and work here, and the city's trade guilds are headquartered here.
\item[Clerkburg.]
Clerkburg is the university district of Greyhawk, with dozens of schools and colleges and the businesses that support them. Temples line the appropriately named Street of Temples in the southeast corner of the district.
\item[Foreign Quarter.]
The Foreign Quarter is among the most multicultural districts of the city, and it boasts fine apartments and restaurants.
\item[Garden Quarter.]
The Garden Quarter is an extravagant neighborhood similar to the High Quarter, but the mansions aren't quite as ornate, the estates aren't as large, and it's not as gaudy.
\item[High Quarter.]
Palaces, temples, mansions, and gardens fill the posh High Quarter. Extravagant architecture and wide-open spaces define this quarter.
\item[River Quarter.]
The River Quarter encompasses taverns and entertainment venues, as well as the wharves along the Selintan River outside the city wall. Because it's a hub of trade, it's the most diverse, multicultural part of the city.
\item[Slum Quarter.]
The Slum Quarter is the poorest, most desperate region within Greyhawk's walls, full of crime-ridden apartments.
\item[Thieves' Quarter.]
The buildings of the Thieves' Quarter are slightly less run down than their Slum Quarter equivalents, and its people are marginally better off.
\end{description}

\subsubsection{City Locations}
The locations detailed here can serve as a good starting point for your campaign. Use them as examples when fleshing out new locations for your game.

\subparagraph{Black Dragon Inn}
\begin{DndSidebar}[float=!b]{}
This three-story, slightly run-down inn is situated in the heart of the city. A sign carved to resemble the grinning visage of a black dragon hangs over the front door. A stable is located behind the inn.



\end{DndSidebar}

The Black Dragon Inn in Clerkburg has good food and affordable rooms. The inn's stable can hold up to a dozen steeds.

The inn's proprietor is Miklos Dare, a human \textbf{Warrior Veteran} (Chaotic Good) who loves to recount his heroic exploits in the Battle of Emridy Meadows seven years ago, when warriors from across the Central Flanaess united to drive the forces of wickedness from the Temple of Elemental Evil (see "\textit{Central Flanaess}" in this chapter). A red-bearded bear of a man with a prosthetic leg, Miklos is affable and proud. His friendly rivalry with Olaf and Sivan, the proprietors of the Silver Dragon Inn just up the street, is the talk of the city. Olaf and Sivan recently hired a mage to make Miklos's black dragon sign drool acid, much to the chagrin of visitors entering and leaving the Black Dragon. Miklos is itching to pull a similar prank of his own.

\subparagraph{Character Backgrounds}
An adventurer with the Soldier background might have a tie to Miklos, perhaps having fought alongside him at the Battle of Emridy Meadows. A character with the Wayfarer background might know Miklos as a generous man who gives away food and sometimes even lodging to people in need.

\subparagraph{Reasons to Visit}
Adventurers might visit the Black Dragon Inn for one of the following reasons:


\begin{description}
\item[Eavesdropper's Paradise.]
Many clandestine meetings occur at the Black Dragon. Adventurers eavesdropping on private conversations might overhear tantalizing rumors or uncover valuable information.
\item[Information Source.]
If the adventurers let Miklos tell stories of his past exploits or agree to help him play a prank on his rivals, he can steer them toward new adventure opportunities. He's also quite familiar with the nature of Elemental Evil.
\item[Place to Stay.]
The Black Dragon is close to the city's central marketplace. A traveler can sleep in a common room for 2 SP per night or secure a private room for 5 SP per night. A luxury suite costs 2 GP per night.
\end{description}

\subparagraph{Grand Citadel}
\begin{DndSidebar}[float=!b]{}
A many-towered fortress looms above all quarters of the city from its position atop a low rise. Its outer walls, darkened by soot and smoke, could use a good scrubbing.



\end{DndSidebar}

The grand edifice at the northern end of the High Quarter, simply called "the Citadel" by the city's inhabitants, contains barracks for the City Watch, the offices of the captain-general, the city's treasury, and a large store of armaments for the emergency citizen militia. The Citadel also contains a prison where the city's most hardened criminals are incarcerated.

The captain-general of the City Watch is Sental Nurev, a tall, human \textbf{Warrior Veteran} (Neutral Good) with thinning blond hair and a mustache. Sental is usually incorruptible, but he is under great stress. The rulers of Stoink, a fortified town in the Bandit Kingdoms, have captured Sental's brother Sarek and are forcing the captain-general to provide information about Greyhawk's defenses and local politicians. Sental gives this information to a human \textbf{Spy} (Chaotic Evil) who stays at the Black Dragon Inn under the false name Skanda Drond. Sental is unaware that the bandit lords of Stoink are pawns of Iuz, whose dreams of conquest extend to the Free City of Greyhawk and far beyond.

The city's constable—who serves as second-in-command to the captain-general, manager to the members of the watch, and a member of the Directing Oligarchy—is a compassionate \textbf{Priest} of Pelor named Derider Fanshen (Neutral Good). Her kindness and talent for healing make her well loved among the watch, and as a former adventurer, she is sympathetic to adventurers' needs. She's unaware of Sental's compromised position.

\subparagraph{Character Backgrounds}
Adventurers with the Criminal or Guard background might have a connection to the Grand Citadel involving a past run-in with the law or past service on the watch.

\subparagraph{Reasons to Visit}
Adventurers might be drawn to the Grand Citadel for one of the following reasons:


\begin{description}
\item[Appointment.]
The adventurers have an appointment to speak with Sental Nurev, perhaps because they need help freeing a companion who was arrested for a crime or because they wish to report a threat to the city.
\item[Break-In or Breakout.]
The adventurers are hired to break into the Citadel's treasury vault or break someone out of the Citadel's prison.
\item[Imprisonment.]
The adventurers are imprisoned in the Citadel for some heinous crime.
\end{description}

\subparagraph{Great Library}
\begin{DndSidebar}[float=!b]{}
The front of this building is a grand sweep of granite walls and tall columns. A wide ramp leads to a pair of massive doors flanked by stone-carved dragons. Inside, it's cool and musty.



\end{DndSidebar}

Weapons and armor aren't permitted in the Great Library. If anyone wearing armor or carrying a visible weapon tries to enter the library, or if a thief is spotted trying to leave the library with one or more stolen books, the stone-carved bronze dragons flanking the entrance animate and attack. These statues are Stone Golems.

Abra Saghast, a crusty and irascible dragonborn sage, serves as the head librarian. Abra, an \textbf{Archmage} (Chaotic Good), has bright-green eyes, and her bronze scales are tinged with aquamarine blue. She typically wears a patchwork robe.

Abra sits behind a high desk in the main hall. Six open archways lead from the main hall to wings where the bulk of the library's books are shelved, free for visitors to peruse (but not remove from the library). The library has several sages and scribes under contract to write books, mostly detailing current affairs in the city.

An iron door leads to a hallway behind the head librarian's desk. Three scribes (Mages) labor here and act as sentries, for next to their desks are three locked, iron doors to the library's vaults. \textit{Arcane Lock} spells seal these doors, beyond which are repositories for the library's most valuable or scandalous works. Next to each scribe's desk is a pull cord hanging through a hole in the ceiling. A tug on any one of these cords releases a homing pigeon from a loft above the library. It takes the bird 1 minute to find and alert an \textbf{Archmage}, who teleports to the main hall of the library to investigate.

\subparagraph{Character Backgrounds}
Adventurers with the Sage or Scribe background might have a connection to the Great Library and its proprietor.

\subparagraph{Reasons to Visit}
Adventurers might visit the Great Library for one of the following reasons:


\begin{description}
\item[Research.]
Adventurers searching for a specific book or more information about a specific topic might find what they're looking for in the library.
\item[Spellbooks.]
The adventurers might need to purloin one of the many spellbooks kept in the library, necessitating a carefully planned heist.
\item[Spell Scrolls.]
Adventurers can commission the scribes to create a \textit{Spell Scroll} that bears a Wizard spell of level 5 or lower. See the \textit{Player's Handbook} for the time required to \textit{craft a scroll}; the scribes charge double the cost shown there.
\end{description}

\subparagraph{High Tower Inn}
\begin{DndSidebar}[float=!b]{}
Conveniently located near the Selintan River, this inviting inn is distinguished by its tall tower, which is pointed at the top like the hat of an eccentric wizard. The clientele is notably wealthy, but the inn itself isn't at all ostentatious.



\end{DndSidebar}

The High Tower Inn's human proprietor, Erlynn Goodfellow, is a soft-spoken, middle-aged, pot-bellied \textbf{Mage} (Lawful Good) with gray hair, bright-blue eyes, and platinum-rimmed spectacles. She dabbled in adventuring before realizing she had little taste for danger and her life's calling might involve more sedentary pursuits. Few guests know of Erlynn's magical abilities, as she rarely casts spells in front of strangers.

The High Tower, located in the Garden Quarter, is a favorite haunt for some of the city's most famous wizards, including \textit{Otto} and \textit{Jallarzi}.

\subparagraph{Character Backgrounds}
Adventurers with the Merchant or Noble background might have a connection to the High Tower Inn, which caters to people of means.

\subparagraph{Reasons to Visit}
Adventurers might visit the High Tower for one of the following reasons:


\begin{description}
\item[Information Source.]
Erlynn knows all the local rumors. If she trusts the adventurers, Erlynn can direct them to an NPC who needs their services.
\item[Place to Stay.]
The adventurers need a place to stay, and the High Tower boasts comfortable quarters, ample supplies of wine and ale, and spicy food. The establishment has six private sleeping chambers, each of which rents for 3 GP per night. Three of these guest rooms are in the tower, each one on its own floor.
\item[Spellcaster.]
The adventurers might have business with a powerful spellcaster staying at the inn.
\end{description}

\subparagraph{Silver Dragon Inn}
\begin{DndSidebar}[float=!b]{}
This grandiose, multistory inn sports a wooden sign bearing the words "Silver Dragon Inn" in fancy silver script, the S shaped like a silver dragon. A more modest sign next to the front door reads, "No metal armor. Shields and weapons must be checked at the door."



\end{DndSidebar}

The grand Silver Dragon Inn, located in the Foreign Quarter, is often the first place sought by new arrivals to the city. The prices are average, but the food servings are huge. The inn's menu includes spicy bean dishes, seafood delicacies of the Wild Coast, and rice and vegetable entrées.

Weapons larger than daggers must be checked at the door, together with shields. Customers wearing metal armor aren't admitted. Two bouncers (Neutral Tough Bosses) stand at the door, politely enforcing the rule.

The inn's married human proprietors, Olaf Al-Azul (Chaotic Good \textbf{Warrior Veteran}) and Sivan Al-Azul (Chaotic Neutral \textbf{Assassin}), speak multiple languages and use humor to raise spirits and diffuse tensions. Olaf can almost always break up a fight before it starts, generally with a round of drinks for the instigators. Sivan is quiet and introspective, but he always keeps a hilarious joke or cutting remark at the ready.

\subparagraph{Character Backgrounds}
Adventurers with the Artisan or Entertainer background might do business with the Silver Dragon Inn.

\subparagraph{Reasons to Visit}
Adventurers might visit the Silver Dragon for one of the following reasons:


\begin{description}
\item[Meeting with Foreign Dignitaries.]
Foreign dignitaries come here to enjoy the Silver Dragon's food and accommodations. Characters who plan to visit distant lands might connect with these esteemed visitors.
\item[Place to Stay.]
The Silver Dragon is centrally located in the city and has twenty-four guest rooms on the upper floors. A traveler can sleep in a common room for 1 SP per night, a private room for 3 SP per night, or a luxury suite for 1 GP per night.
\item[Security.]
The Silver Dragon Inn prides itself on being a safe stop for visitors. Its proprietors and bouncers are trained to deal with trouble without the support of the City Watch.
\end{description}

\subparagraph{Temple of the Far Horizon}
\begin{DndSidebar}[float=!b]{}
Hidden among the city's grander temples is a quiet, modest house of worship with clay-tiled rooftops, a corner bell tower, and well-tended vegetable gardens. Sick and hungry folk gather in short lines outside as they wait for priests to attend to their needs.



\end{DndSidebar}

Situated in the Garden Quarter, this temple is dedicated to Fharlanghn, a god favored by travelers and mercenaries. The Priests who staff the temple offer nourishment, rest, and healing to those in need, day and night. Several small rooms are maintained for guests, and simple, hot meals are free to all visitors.

\subparagraph{Character Backgrounds}
Adventurers with the Guide or Sailor background might have a connection to the temple, which offers help to travelers.

\subparagraph{Reasons to Visit}
Adventurers might visit the temple for one of the following reasons:


\begin{description}
\item[Adventurers Wanted.]
The priests keep tabs on threats in the region around the city. They're paying close attention to rumors of dragon activity in the nearby Cairn Hills, and they're looking to hire adventurers to investigate these rumors.
\item[Healing.]
Adventurers can purchase Potion of Healing for 50 GP each, and the temple's priests have 1d4 such potions in stock on any given day. The priests also have \textit{Cure Wounds} and \textit{Lesser Restoration} spells prepared and customarily cast them for free. For more powerful magic, such as \textit{Greater Restoration} and \textit{Raise Dead} spells, the priests direct the adventurers to the Temple of the Radiant Sun.
\item[Safe Travels.]
By making a small donation to the temple, adventurers increase the likelihood of safe travel to their next destination.
\item[Teleportation Circle.]
Though it isn't the only permanent teleportation circle in the city, the circle within the Temple of the Far Horizon is the easiest to access. The priests allow free access to the teleportation circle in either direction. For 2,000 GP, the chief priest will cast the \textit{Teleportation Circle} spell to open a connection to another permanent circle on the Material Plane.
\end{description}

\subparagraph{Temple of the Radiant Sun}
\begin{DndSidebar}[float=!b]{}
This copper-roofed temple has a gold-inlaid symbol of the sun above its double-door entrance. During the day, sunlight shines through high windows to illuminate the temple's interior, which is adorned with golden draperies.



\end{DndSidebar}

This temple, dedicated to serving the god Pelor in the heart of the Garden Quarter, opens at dawn and closes at dusk. In a sanctuary in the heart of the temple, Priests conduct daily morning rites, as well as all-day observances every Godsday.

Sarana, the temple's \textbf{Archpriest} (Neutral Good), is a middle-aged, human woman wearing a sun-shaped headdress and yellow-and-gold robes. She is never seen in public without her \textit{Staff of Healing}. Sarana has straw-colored hair, green eyes, and a forgiving nature.

\subparagraph{Character Backgrounds}
Adventurers with the Acolyte background might have served in the Temple of the Radiant Sun, while those with the Farmer background might seek it out as a place for blessing.

\subparagraph{Reasons to Visit}
Adventurers might visit the temple for one of the following reasons:


\begin{description}
\item[Healing.]
The temple sells Spell Scroll of \textit{Greater Restoration} for 3,200 GP apiece and Spell Scroll of \textit{Remove Curse} for 300 GP apiece, and the priests have 1d3 copies of each scroll in stock on any given day. The priests also have \textit{Cure Wounds} and \textit{Lesser Restoration} spells prepared, which they customarily cast for free.
\item[Raise Dead.]
Archpriest Sarana is one of a handful of people in the Free City of Greyhawk who can cast the \textit{Raise Dead} spell, but she needs the requisite 500 GP diamond to do so. Sarana can recommend a jeweler who sells diamonds of sufficient value. Before agreeing to cast the spell, Sarana casts \textit{Zone of Truth} and asks questions about the deceased individual to make sure she's not returning to life someone who should stay dead.
\item[Service to the Greater Good.]
The temple might call upon the adventurers to perform good acts in the city or abroad. Sarana is particularly vigilant about the threat of Elemental Evil, since she was involved in the battle at the Temple of Elemental Evil seven years ago. In exchange for their service, the characters and their companions are entitled to a 50 percent discount on goods purchased at the temple.
\end{description}

\subparagraph{Unearthed Arcana}
\begin{DndSidebar}[float=!b]{}
This quaint, two-story shop has a sign depicting a white-bearded human wizard holding a staff that has a copper ball affixed to its tip. Displayed in the store's window box are various potions, scrolls, wands, and wondrous oddities.



\end{DndSidebar}

Magic items are bought and sold in Unearthed Arcana, a quaint shop in Clerkburg. Magical wards render the store's windows and doors shatterproof, and no one can use magic to enter or leave the shop without the consent of its proprietor, Morley, whose quarters take up the second floor.

Morley is an \textbf{Adult Copper Dragon} (Chaotic Good) who spends his days shape-shifted into a talkative, alert, white-bearded human mage wearing a pointed hat, frayed robes, and pointed slippers. Only a few people in the city—including the esteemed local members of the Circle of Eight, Jallarzi Sallavarian and Otto—know Morley's true form.

Morley is one of the city's secret weapons, ready to repel invaders or break a siege should the need arise. The dragon has a soft spot for adventurers who risk their lives for good causes. He occasionally loans magic items free of charge to valorous heroes who can't afford them, on the condition that the items be returned to him as soon as they're no longer needed.

\subparagraph{Character Backgrounds}
Adventurers with the Charlatan or Hermit background might have a connection to Unearthed Arcana, as Morley has a variety of unusual interests.

\subparagraph{Reasons to Visit}
Adventurers might visit Unearthed Arcana for one of the following reasons:


\begin{description}
\item[Buying and Selling Magic Items.]
Morley buys and sells magic items at \textit{standard prices} (see \textit{chapter 7}). Although he keeps a few magic items in the shop to catch the eye, most of his inventory is stored in extradimensional vaults only he can access. The shop sells many Common, Uncommon, and Rare magic items—mainly potions, rings, rods, staffs, wands, and wondrous items. Morley has access to a few Very Rare and Legendary magic items as well.
\item[Free Loan.]
A benefactor arranges for Morley to loan the characters a magic item to help them complete a quest. Before giving them the item, Morley asks they return it in pristine condition.
\item[Magic Item Identification.]
Morley can cast the \textit{Identify} spell at will. He charges 50 GP for each casting of the spell.
\end{description}

\subsubsection{Beyond the City Walls}
The City of Greyhawk and Environs map shows the lands around the Free City of Greyhawk. Locations on the map are presented below as places where adventures can happen:


\begin{description}
\item[Blackfair Manor.]
Typical of several manor houses and keeps scattered across the Plain of Greyhawk, Blackfair Manor was founded by a distinguished cavalry commander. Its stable is the most famous source of fast, durable warhorses across the breadth of the Flanaess, drawing shrewd shoppers from Greyhawk and beyond. The manor house is surrounded by farms, extensive pastureland, and a small village with a mill, taverns, a smithy, and a saddlery.
\item[Blackstone.]
See "Mining Towns" below.
\item[Blackwall Keep.]
One of two new keeps built to keep an eye on the Mistmarsh, Blackwall Keep is a strong, stone tower with a horse corral surrounded by a wooden stockade. Soldiers from Greyhawk garrison the keep and venture out from it to patrol the northern edge of the swamp.
\item[Cairn Hills.]
Hidden among the hills north and east of the city are ancient tombs and half-buried ruins that attract adventurers, bandits, cultists of Elemental Evil, and monsters.
\item[Castle Greyhawk.]
Travelers who follow the Selintan River westward from the city come to a stone bridge. From there, they must travel several miles northeast to reach the ruins of Castle Greyhawk. Built by the archmage Zagig Yragerne and abandoned with his demise, the ruins (and the many-leveled dungeon below) are a powerful draw to adventurers who seek wealth, glory, and magical might. All manner of marvels are said to fill the ruins, including numerous portals to other planes.
\item[Diamond Lake.]
See "Mining Towns" below.
\item[Elmshire.]
This sleepy town with a sizable halfling population lies on the shore of the Nyr Dyv. Fishing boats crowd the wharves. The townsfolk welcome peaceful visitors—particularly adventurers who can help fend off monsters, bandits, and cults of Elemental Evil lurking in Cairn Hills.
\item[Ery Villages.]
The villages of High Ery and Erybend are populated largely by farmers who send most of their produce for sale in Greyhawk. The two villages are both dominated by two prominent families, the Fairheights and the Witherwinns. The families have a complicated history including abundant instances of feuding and intermarriage, as well as a catalog of lesser slights and favors.
\item[Ford Keep.]
A ferry crossing allows traffic from Greyhawk and regions to the south to cross the Selintan River at its first major bend. The lord mayor of Greyhawk built Ford Keep here to protect the crossing from bandits.
\item[Gorge of the Selintan.]
Soaring cliffs flank the Selintan River for nearly ten miles. Spanning this gorge, 800 feet above the river, is a stone arch bridge sculpted to look like an extension of the natural bedrock. The bridge allows easy travel between Greyhawk and Grossettgrottell. As a defensive measure, the bridge's gnome architects hid an iron pin somewhere in the bridge; if this pin is removed, the entire structure collapses.
\item[Grossettgrottell.]
An industrious community of gnome miners and foragers lives in this network of hewn tunnels and natural caverns. The gnomes trade gemstones and rare fungi in exchange for help repelling monsters from the Underdark.
\item[Marsh Keep.]
Like Blackwall Keep, Marsh Keep is newly built and watches over the Mistmarsh. The Dwarfwalk road leads east from the tower to Greysmere, a quarrying town with a large population of dwarves.
\item[Mining Towns.]
Blackstone, Diamond Lake, and Steaming Springs are small mining towns governed and protected by the Free City of Greyhawk. The city frequently dispatches adventurers to quell threats to the towns' miners and mining operations, which of late includes agents of Iuz intent on destabilizing the city's economy.
\item[Mistmarsh.]
This vast swamp holds the half-sunken ruins of an ancient city that is now shrouded by fog and guarded by a family of black dragons. Will-o'-wisps lure prey to the ruins, where doom awaits.
\item[Peculiar Manor.]
Like other manors in the Plain of Greyhawk, Peculiar Manor was established by a now-forgotten hero of an ancient war. However, a few years ago it was purchased by retired adventurers from Ekbir, Sanjarah and Chetna Mohsin. The Mohsins brew an extraordinary ale they call Old Peculiar, which lends its name to the manor.
\item[Steaming Springs.]
See "Mining Towns" above.
\item[Stone Bridge.]
A small garrison in the fort of Stone Bridge keeps careful watch over river traffic approaching Greyhawk.
\item[Stonefort.]
A garrison of soldiers from Greyhawk watches over the southern Nyr Dyv from the high battlements of Stonefort. The fort also guards a gravel quarry.
\item[Tokhel Castle.]
This blasted ruin stands on a rocky promontory above the Nyr Dyv. A \textit{dead magic zone} (see "\textit{Environmental Effects}" in \textit{chapter 3}) encompasses most of the ruin, and monsters guard whatever secrets the ancient castle and its dungeons might hold.
\item[Two Ford.]
The small village of Two Ford relies on river trade to supply its inns, smithies, and merchants. Ore from the mining towns is also traded here, as it is easier to transport it by river than overland.
\end{description}

\subsection{Greyhawk Gazetteer}
The poster map in this book shows the entire region of the Flanaess, with the Free City of Greyhawk near the center. As characters venture beyond the confines of the city and its surrounding lands, you can use the map and the information on these pages to inspire your own adventures and world details.

\subsubsection{The Big Picture}
To understand the Greyhawk of today (the year 576 in the Common Year), it is helpful to picture the Flanaess about 200 years ago. At that time, the Great Kingdom of Aerdy stretched from the Vast Swamp to the Rakers, and from the Solnor Ocean to the Yatil and Lortmil Mountains. Between the Lortmil Mountains and the Crystalmists was the Kingdom of Keoland, mimicking the Great Kingdom in its imperialist approach. To the west of the mountains, the Baklunish nations—survivors of the Invoked Devastation—stood much as they do now. And to the north, nomadic peoples (the Chakyik, the Wegwiur, the rovers of the Hunting Lands) and the North Kingdoms were, as now, independent of the politics of the south.

The contraction of the two southern realms over the last two centuries is the primary force that has shaped the modern Flanaess. First, the provinces of Perrenland, Veluna, and Furyondy—far from the overking's throne in Rauxes—declared their independence from the Great Kingdom. When Nyrond joined them (in 356 CY), the stage was set for the slow disintegration of the Great Kingdom and the ongoing rebellion that dominates the eastern part of the Flanaess to this day.

The beginning of Aerdy's decline marked the high-water mark of Keoland's expansionist policies, as it held sway from the Pomarj to the borders of Ket. Aided by Celene, rebels in the Ulek region strove to curb the kingdom's warlike ways, and the accession of King Tavish IV in 453 CY marked a dramatic shift in royal policy. The Yeomanry and the Ulek states were granted autonomy, the Gran March and Sterich became semi-independent while remaining loyal to the crown, and the diverse peoples of this region coexist in relative peace once again.

With this big picture in mind, you can think of the Flanaess beyond the Free City of Greyhawk as five major regions, each with its own store of adventure possibilities waiting to be explored:


\begin{description}
\item[Central Flanaess.]
Diverse peoples clash against Iuz and the forces of Elemental Evil.
\item[Eastern Flanaess.]
The remnants of the Great Kingdom struggle to determine the fate of the lands in the overking's wicked clutches.
\item[Northern Flanaess.]
In vast wilderness expanses populated by indigenous folk, one's mettle is tested by the environment and roving dragons.
\item[Old Keoland.]
The former provinces of Keoland contend against evil monsters from the western mountains, including dragons and giants.
\item[Western Flanaess.]
The Baklunish peoples and nations navigate complex political relationships.
\end{description}

\begin{DndSidebar}[float=!b]{Living History}
Greyhawk has a long history of ancient empires and more recent wars, but the only relevant details of this history are those that feature in your adventures. Highlight significant in-world details by revealing them in the course of your adventures. Use the following techniques to share lore with your players:




\begin{description}
\item[Echoes of the Past.]
Features like a crater in the side of a mountain, a defaced statue in the town square, or a holiday celebrating a local hero provide concrete touchstones to past events. The characters might learn that a ruin they're exploring was destroyed in a catastrophic battle or natural disaster.
\item[Historical Records.]
Written historical details might appear anywhere in an adventure: glyphs on ancient dungeon walls, books in a library, files in a royal vault, or tapestries depicting key events. Use such set dressing to share important details. Summarize what lengthy works say, and focus on the most plot-worthy parts.
\item[Scholarly Expertise.]
Characters who have proficiency in the Arcana, History, or Religion skill can be fonts of useful information. When it would be helpful for a group to know something about the setting, ask such characters to make an Intelligence check using the relevant skill, then share plot-relevant details if their roll warrants it.
\item[Magical Footnotes.]
Spells such as \textit{Contact Other Plane}, \textit{Legend Lore}, and \textit{Speak with Dead} allow characters to learn information while leaving you control of the particulars.
\item[Primary Sources.]
Personify the past through a tragic ghost, an otherworldly guardian, an artificial intelligence, an ancient sage, or another long-lived individual. Such NPCs give you a way to share relevant information and respond to questions from the party. If the characters miss an important detail, this NPC can reinforce details in a way books and recollected facts can't.
\end{description}



Any one of these methods is useful for revealing a few details. You can combine them to share nuanced histories and help players feel like they're digging into a rich and realistic history.



\end{DndSidebar}

\subsubsection{Central Flanaess}
The rich soil and pleasant climate of the region between the Nyr Dyv and the Yatil Mountains—combined with healthy trade relations between these realms and their neighbors to the east, south, and west—make this a strong and prosperous region.

\begin{DndTable}[header=Central Flanaess Locations]{XXX}

Location & Ruler & Description\\
Celene & Queen Yolande (elf) & Elven monarchy with large gnome and halfling populations\\
Dyvers (Free City) & Magister Thymantia Gortoz (aasimar) & Important port and trading center with a powerful navy\\
Furyondy, the Kingdom of & King Belvor IV (human) & Former province of the Great Kingdom, among the first to claim independence\\
Highfolk (Free City) & Mayor Talisyr (appears as an elf) & Fortified city with large population of elves; the mayor is a disguised adult silver dragon\\
Horned Society, the & Nine hierarchs (mostly humans and hobgoblins) & Theocracy ruled by devil worshipers allied with Iuz, supported by mercenaries enforcing their tyrannical rule\\
Iuz & Iuz (cambion demigod) & The monster-infested domain of the demonic dictator, steeped in wickedness\\
Nyr Dyv & — & The Lake of Unknown Depths; home to barge dwellers\\
Perrenland & Voorzitter Yrenda Schwartzen (human) & Fiercely independent confederation of canons\\
Pomarj, the & — & Lawless peninsula; home to bandits and marauders\\
Shield Lands, the & Various allied nobles & Independent alliance of nobles protected by the Knights of Holy Shielding, led by Knight Commander Aleshh Kaarth (dragonborn) and fortified by Furyondy and Urnst\\
Veluna, the Archclericy of & Canon Hazen (human) & Theocracy ruled by priests of Rao, a divine beacon of justice and hope\\
Verbobonc (Free City and Viscounty) & Viscountess Wilfrick Rejjin (human) & Vassal state of Veluna; site of the Temple of Elemental Evil\\
Wild Coast, the & Various burgomasters, lord mayors, and others & Free territory with self-governing settlements; haven for outcasts and dissidents\\
\end{DndTable}

\subparagraph{Battle of Emridy Meadows}
Seven years ago, knights and soldiers from Furyondy, the Archclericy of Veluna, the Viscounty of Verbobonc, and the elven kingdom of Celene formed an alliance to repel an evil horde that had gathered in the grassy fields south of the Velverdyva River. This clash of armies—arguably the greatest seen in Eastern Oerik—was called the Battle of Emridy Meadows. The forces of evil were smashed, and their remnants were driven back into the dungeons under their stronghold, the Temple of Elemental Evil. The forces of good, under the command of Prince Thrommel IV of Furyondy, besieged the temple, which fell in a fortnight. Only a few of the temple's wicked leaders escaped, and it is suspected that these individuals were responsible for the subsequent kidnapping of the prince.

Prince Thrommel was engaged to marry Lady Jolene, a priest from a prominent noble family of Veluna. Their marriage would have united Furyondy and Veluna as a single entity, with the canon of Veluna ruling in matters spiritual and the king of Furyondy ruling in matters temporal. This combined state, with its powerful elf allies in Celene, could wage a steady war against the evil plaguing Eastern Oerik. The prince's disappearance has stalled these plans.

\subparagraph{The Rise of Iuz}
Iuz is the offspring of the demon lord Graz'zt and a human archmage named Iggwilv. For ages, he ruled the lands from the Howling Hills to Lake Whyestil, naming his domain after himself. These lands are so despoiled and dangerous that the otherwise fierce nomads of the Hunting Lands and Wegwiur pass through the Cold Marshes rather than enter the merest edge of Iuz's realm.

Iuz's evil reign was interrupted by a sixty-five-year imprisonment in the dungeon under Castle Greyhawk. During his absence, the Kingdom of Furyondy and its allies prospered, while the land of Iuz was overrun with evil bandits and monsters. Iuz's absence turned him into a legend and attracted a host of new followers, whose misplaced faith invested him with the power of a demigod.

Upon winning his freedom, Iuz had no trouble reclaiming his homeland. He forged tenuous alliances with the leaders of the Bandit Kingdoms and the Horned Society, whom he controls through terror. With their aid, he aims to destroy his neighbors and lay waste to the Free City of Greyhawk.

Since the resurgence of Iuz, the northern quarter of the Vesve Forest and the eastern part of the Howling Hills have become filled with marauders and monsters. While the Wegwiur battle Iuz's forces in the Howling Hills, scouts and troops from Furyondy join forces with Highfolk's defenders to drive out the Vesve Forest's evil inhabitants.

\subparagraph{Central Flanaess Culture}
The culture of the Central Flanaess is a result of the long imposition of the Great Kingdom's rule over a variety of peoples living in close proximity. These peoples, by and large, share the Great Kingdom's practical, hardworking values, and they rely on the family and local community, rather than the might of nations and armies. They have a strong egalitarian streak unlike the Great Kingdom's strict social hierarchy, and (beyond the domains of Iuz and the Horned Society) they have little tolerance for would-be tyrants or aloof nobility. Amid a large number of free cities and confederations, the monarchies of Furyondy and Celene are far more democratic in practice than those in other regions.

Typical dress in the Central Flanaess includes a tunic of varying length, sometimes worn with close-fitting trousers. A cape or cloak, usually featuring patterns of ovals or diamonds, completes the ensemble. The cuisine of the Central Flanaess uses rice and potatoes alongside cheese and meat that is typically boiled or roasted.

\subparagraph{Central Flanaess Adventures}
The dual threats of Iuz and Elemental Evil present abundant opportunities for adventure in the Central Flanaess (see "\textit{Greyhawk Conflicts}" in this chapter). This region is particularly appropriate for campaigns flavored with epic fantasy, supernatural horror, or war (see "\textit{Flavors of Fantasy}" in this chapter). This region is also home to many of the most famous dungeons and ruins of Greyhawk, including those described in the sections that follow.

\subparagraph{Ghost Tower of Inverness}
Ages ago, an archmage raised the mighty fortress of Inverness from the very rock of the Abbor-Alz. In the great inner tower of the keep, he hid his most prized possession: the Soul-Gem. Legend says this great white diamond fell from the sky and glowed with the brilliance of the sun, and its magic could drag mortal souls screaming from their bodies and trap them. The fortress was ruined long ago, but on foggy nights the great central tower still appears.

\subparagraph{Lost Caverns of Tsojcanth}
The archmage \textit{Iggwilv} is said to have acquired much of her prowess from the hidden magic she discovered within the Lost Caverns of Tsojcanth (pronounced SAWJ-kahn). In these caverns she conducted experiments and rituals to increase her powers. One of these rituals led to her downfall, though, when she accidentally freed the demon lord \textit{Graz'zt} from the prison where she had bound him. Though Graz'zt fled to the \textit{Abyss}, Iggwilv was weakened and forced to abandon the caverns, but a secret cache of her treasure is said to remain. The Lost Caverns of Tsojcanth are detailed in Quests from the Infinite Staircase.

\subparagraph{Maure Castle}
Maure Castle is a forlorn and foreboding place surrounded by boggy ground and scraggly trees. Rumors suggest that a tome holding eldritch lore of unimaginable evil is held within, guarded by a powerful demon.

\subparagraph{Temple of Elemental Evil}
The Temple of Elemental Evil, built long ago, spawned hordes of bloodthirsty monsters that ravaged the lands between Celene and Veluna. As far as anyone in the area knows, the temple is currently abandoned and has not posed a threat since the Battle of Emridy Meadows in 569 CY, but the forces of good in the region keep vigilant against any sign that the temple and its cults might arise once more.

\subsubsection{Eastern Flanaess}
Once a powerful force for order and good, the Great Kingdom of Aerdy has declined over the last century to a state of utter decadence. The reigning overking—Ivid V, patriarch of the House of Naelax—is rumored to have fiendish advisers as well as a noble court infested with evil. Ruling from the Malachite Throne in Rauxes, Ivid commands an unmatched army currently embroiled in two wars at once: one against the Kingdom of Nyrond and the Prelacy of Almor, and the other against the Iron League (consisting of Idee, Irongate, the Lordship of the Isles, Onnwal, and Sunndi). To pay for these costly wars, the overking has imposed heavy taxes on his subjects, further diminishing his popularity.

Aerdiaak, Ahlissa, Medegia, and Rel Astra are provinces and fiefs of the Great Kingdom. The Sea Baronies are vassal states that provide most of the kingdom's navy.

\begin{DndTable}[header=Eastern Flanaess Locations]{XXX}

Location & Ruler & Description\\
Aerdiaak & Herzog Varz Grenell (human) & North province of the Great Kingdom, ruled by a cousin of the overking; its court is rife with debauchery and intrigue\\
Ahlissa & Herzogin Seprenna Calyn (human) & South province of the Great Kingdom, ruled by a cousin of the overking; embroiled in war with the Iron League\\
Almor, the Prelacy of & Prelate Xanther Klimstyn (human) & Theocracy ruled by a priest of Pelor who declared independence when the Great Kingdom descended into evil\\
Bone March, the & — & Fallen territory of the Great Kingdom, now held by armies from Almor and Nyrond\\
Celadon Forest & — & Ancient forest protected by druidic circles and fey\\
Flinty Hills and Gamboge Forest & — & Home to several independent communities with no great love for Nyrond or the Pale\\
Great Kingdom, the & Overking Ivid V (human) & Unspeakably evil monarchy\\
Idee & Count Vasiliek Donsten (human) & Independent fiefdom; member of the Iron League\\
Irongate (Free City) & Mayor Unthera Selvich (appears as a dwarf) & Thriving metropolis; member of the Iron League; the mayor is a disguised adult bronze dragon\\
Lordship of the Isles, the & Princess Ronthal III (human) & Independent principality; member of the Iron League\\
Medegia, the See of & Holy Censor Starvik Jerel (human) & Theocratic fiefdom ruled by a priest whose power is rumored to come from pacts with archdevils\\
Nyrond, the Kingdom of & King Dunstan I (human) & Center of resistance to the Great Kingdom\\
Onnwal & The Raven of Onnwal, Zyl Grayshadow (dwarf) & Independent state; member of the Iron League\\
Pale, the Theocracy of the & Supreme Prelate Ogon Tillit (human) & Theocracy ruled by a priest of Pholtus\\
Rel Astra, City of & Constable Mayor Drax (orc) & Independent fief plotting in secret against the Great Kingdom, hoping to ally with Medegia or the Sea Baronies\\
Sea Baronies, the & Four sea barons, including High Admiral Kalashe Asperdi (human) & Independent island fiefdoms that serve as the Great Kingdom's navy\\
Shar, the Hidden Empire of & Father of Obedience Korenth Zan (human?) & Isolated order of Suloise militants whose spies operate across the Flanaess\\
Spindrift Isles, the & The Council of Five (on the northern islands) and the Council of Seven (on the southern island) & Independent islands that keep watchful eyes on aggressive island neighbors\\
Sunndi & Steward Valenta (elf) & Independent fiefdom; member of the Iron League\\
Tenh, the Duchy of & Duchess Ehliyah Raynar III (human) & Independent fiefdom allied with Nyrond for defense against Iuz\\
Urnst, the County of & Countess Belissica Gellor (human) & Independent fiefdom\\
Urnst, the Duchy of & Duke Jalken Lorinar (human) & Independent fiefdom\\
Vast Swamp, the & — & Morass separating the Tilvanot Peninsula and Shar from the rest of the East\\
\end{DndTable}

\subparagraph{Shar}
Fanatical Suloise militarists called the Scarlet Order founded the Hidden Empire of Shar, which is closed to outsiders. The order controls the peninsula west of the Tilva Strait, as far north as the Vast Swamp. See "\textit{Factions and Organizations}" in this chapter for more about the \textit{Scarlet Order}.

\subparagraph{Eastern Flanaess Culture}
The culture of the Eastern Flanaess is largely that of the ancient Aerdi tribe of humans that conquered the region and established the Great Kingdom almost 800 years ago. The Aerdi valued common sense, hard work, and knowing one's place in a strict social order. Having claimed their position of power through conquest, they put great emphasis on military power and martial skill. These values persist throughout the region, reinforced in the Great Kingdom by strict laws and even stricter social mores.

These same values persist in many of the lands resisting the overking's reign. Nyrond and Almor, in particular, share the stratified social structure of the Great Kingdom, with their king and prelate remaining distant from the common people they rule. The independent states of the Iron League are more egalitarian, sharing that trait with the peoples of the Central Flanaess.

Typical clothing in the Eastern Flanaess is a tunic of varying length, often worn with trousers, with a cape or cloak. The fabrics of the east are often patterned with checks or plaids, with different patterns often relating to the wearer's lineage. Eastern cuisine pairs rice and potatoes with a variety of meats, especially seafood.

\subparagraph{Eastern Flanaess and Its Neighbors}
The Duchy and County of Urnst bridge the regions of the Central and Eastern Flanaess. Once part of the Great Kingdom's province of Nyrond, they declared their independence from the Great Kingdom and the new Kingdom of Nyrond at the same time, achieving their separation from Nyrond with minimal bloodshed. While the people of Urnst distrust the king of Nyrond, they don't hate him like they do the overking.

The proximity of the Nyr Dyv, the Cairn Hills, and the Shield Lands means the Urnst lands can't ignore the rising threat of Iuz or the politics of the Free City of Greyhawk. At the same time, Nyrond stands as a buffer between Urnst and the Great Kingdom, but the overking's threat still looms. The duke and countess of Urnst believe that a united Urnst will stand more strongly against pressures from east and west, which they hope to achieve through the marriage of Countess Belissica Gellor to Byron Lorinar, eldest son of Duke Jalken Lorinar.

\subparagraph{Eastern Flanaess Adventures}
The story of the Eastern Flanaess is a tale of scrappy rebels—Nyrond, Almor, and the Iron League—defying the overwhelming power of a corrupt and decadent empire. This story lends itself to campaigns exploring themes of supernatural horror (in the fiend-haunted courts of the overking), swashbuckling (in the cities across the region as well as the eastern seas), and war (see "\textit{Flavors of Fantasy}" in this chapter).

\subparagraph{Havens of Unrest}
Those who despise and challenge the Great Kingdom's oppression—outlaws both good and bad—find refuge in borderlands just beyond the reach of the overking's soldiers. These include the woods and swamp near Rel Astra (the eastern reaches of the Grandwood Forest and the Lone Heath) and the Hestmark Highlands east of Sunndi. The outlaws in the Grandwood include significant numbers of elves and halflings as well as humans, while those in the Hestmark Highlands include dwarves and gnomes.

\subparagraph{Tomb of Horrors}
Deep within the Vast Swamp is the sinister Tomb of Horrors—a labyrinthine crypt filled with terrible traps, strange and ferocious monsters, and rich and magical treasures. Somewhere within rests the demilich \textit{Acererak}, who ruled much of the region long ago. The demilich is said to take perverse pleasure in devouring adventurers' souls.

The Tomb of Horrors is detailed in \textit{Tales from the Yawning Portal}.

\subparagraph{Troll Fens}
The chilly mists of the Troll Fens, located against the shoulders of the Griff Mountains and the Rakers at the head of the Yol River, cloak a place of unnameable horrors. As the name implies, the fens are infested with particularly large and vicious trolls. The Pale carefully hedges the place with watchtowers and keeps, and strong patrols ride the verges of the southern end of the Troll Fens to watch for unwelcome visits from the monsters dwelling within.

\subsubsection{Northern Flanaess}
The northern region of the Flanaess includes three distinct areas populated by different peoples: the Baklunish horse riders of the Chakyik and the Wegwiur, the Suloise people of the North Kingdoms, and the Flan nomads of the Hunting Lands.

\begin{DndTable}[header=Northern Flanaess Locations]{XXX}

Location & Ruler & Description\\
Arn, the Archbarony of & Archbaron of Arn (identity unknown) & Remote and little-known region located near a ruined castle with monster-filled dungeons\\
Bandit Kingdoms, the & Four to six bandit lords & Feuding kingdoms ruled by greedy bandit lords with private armies\\
Barren Wastes, the & — & Harsh, despoiled land where dragons roam and sometimes go to die\\
Chakyik & Lord Agul Krusef (human) & Land of the Tiger Nomads—horse riders with scattered trading outposts\\
Hunting Lands, the & Overlord-Protector Yhareen Sakarr (tiefling) & Home to Flan nomads, known to their neighbors as the Rovers of the Barrens\\
North Kingdom of the Cruski, the & Queen Tharla of the Cruski (human) & Fierce, seafaring berserkers of the North Kingdoms\\
North Kingdom of the Fruzti, the & King Hundgred of the Fruzti (human) & Weakest of the three North Kingdoms, having suffered great losses battling in the Bone March\\
North Kingdom of the Schnai, the & Queen Ingrid of the Schnai (human) & Strongest and most populated of the North Kingdoms\\
Ratik, the Barony of & Baron Lexnol Haarkof (human) & Former province of the Great Kingdom trying to ally with the North Kingdoms\\
Stonefist, The Hold of & King Sevvord Redbeard of the Hold (human) & Monarchy founded by a bandit leader who attracted malcontents from many nations\\
Wegwiur & Wolf-Mother Bargra Yefkos (human) & Land of the Wolf Nomads—horse riders engaged in war against Iuz\\
\end{DndTable}

\subparagraph{Baklunish Nomads}
The Chakyik and Wegwiur—called Tiger Nomads and Wolf Nomads, respectively, by their neighbors—are horse riders of Baklunish descent who dwell on the steppes north of the Yatil Mountains and Lake Quag. The climate in the steppes and pine forests ranges from cool to frigid. Both peoples maintain scattered trading outposts that welcome visitors from neighboring and distant lands.

The steppe nomads have a rich storytelling tradition that reinforces a strong sense of clan identity and family line within the clan. Their tales include stories not only of heroes within their clans, but also of heroic horses, and the nomads trace equine lineages as carefully as their own. These nomads maintain the traditions of their people, many of which can be traced back to the ancient Baklunish empire.

The Baklunish nomads favor bright pastel colors in gowns and robes. When traveling or at war, though, they prefer more rugged gear of leather and hide.

\subparagraph{The North Kingdoms}
Three kingdoms of related peoples occupy the Thillonrian Peninsula in the northeast of the Flanaess—a beautiful subarctic landscape of high mountains, coniferous forests, and deep fjords. The kingdoms are named for the three distinct tribal lines that inhabit them: the Cruski (whose name means "ice"), the Fruzti ("frost"), and the Schnai ("snow").

The Schnai are strong and numerous. When Queen Ingrid of the Schnai has a mind to raid the isles of the Sea Baronies or the coasts of Aerdiaak and the Great Kingdom, she calls upon the king of the Fruzti and the queen of the Cruski to join her army. King Hundgred of the Fruzti has no choice but to honor his oath to the Schnai, while Queen Tharla of the Cruski rarely turns down a chance to attack her enemies. At other times, the Cruski raid the Fruzti, the Schnai, or the Hold of Stonefist.

As their distinct kingdoms suggest, the people of the North Kingdoms value their connection to their clan lineage. They preserve a love of learning from their distant ancestors of the Suel Imperium, and they value knowledge of the natural world as highly as they do the skills of hunting, sailing, and warfare. Their clothing includes shirts and pants made of wool, augmented with furs, capes, mittens, and warm boots. They often wear large pins, brooches, or emblems in their cloaks as a sign of wealth or accomplishment.

\subparagraph{The Hunting Lands}
The People of the Hunting Lands (called Rovers of the Barrens by their neighbors) have a history of raiding the outskirts of Furyondy, the Bandit Kingdoms, Tenh, and Wegwiur, which hasn't won them many allies. The nomads' legendary dominance of the north has faded, as the forces of Iuz and the Horned Society wage steady war against them while raiders from the Hold of Stonefist prey on the Hunting Lands farther east. Many of the Rovers' mightiest warriors—the Wardogs—have perished in battles against all these relentless foes.

The people of the Hunting Lands value a close connection to the natural world. They view nature as an entity to be respected, not controlled, and their myths and legends teach the value of accepting nature's bounty as a gift that evokes gratitude. They wear clothes made entirely of animal skins, including belts, capes, robes, and slippers, and decorate their skin with paints and tattoos.

\subparagraph{Northern Flanaess and Its Neighbors}
The regions of the north exist on the fringes of other regions of the Flanaess. Three realms are the primary points of intersection between the Northern Flanaess and neighboring areas.

\subparagraph{Bandit Kingdoms}
The Bandit Kingdoms is a lawless frontier between the Hunting Lands in the north, the Horned Society and the Shield Lands in the Central Flanaess, and the Duchy of Tenh in the east. No single bandit lord is powerful enough to conquer the whole territory, and the combined strength of all is often required to defend against retaliation by neighboring states for the bandit lords' aggression. At least one of the bandit lords, Renfus the Mottled (ruler of Stoink), is wholly in the service of Iuz.

\subparagraph{Ratik}
As a former province of the Great Kingdom, Ratik rides the boundary between the northern and eastern regions of the Flanaess. Without the protection of the Great Kingdom, Ratik has been forced to defend itself against frequent raids from the North Kingdoms and the Sea Baronies, as well as attacks from mountain-dwelling monsters. Baron Lexnol Haarkof's emissaries hope to forge an alliance with the North Kingdoms and redirect the berserkers' aggressions toward the Hold of Stonefist.

\subparagraph{Wegwiur}
Wolf-Mother Bargra Yefkos of the Wegwiur Hordes is preoccupied with the threat of Iuz, and she meets frequently with clan leaders and Perrenlander mercenaries to strategize. The Wegwiur consider their territory to extend to the Dulsi River, so they fiercely defend the western Howling Hills from the incursion of the hideous monsters that serve Iuz. Several large battles between Wegwiur and the forces of Iuz have taken place in that area.

\subparagraph{Northern Flanaess Adventures}
The cold north is an ideal location for a campaign featuring themes of \textit{sword-and-sorcery fantasy} (see "\textit{Flavors of Fantasy}" in this chapter). The peoples of the Northern Flanaess battle giants, dragons, and other horrific monsters in equally dangerous environments, while remaining suspicious of the decadence of the cities and nations of the south.

\subsubsection{Old Keoland}
United by their shared history as part of the ancient Kingdom of Keoland, the marches and fiefdoms between the Lortmil Mountains and the higher mountains to the west gather diverse populations of many different species in relative peace with each other. Even the long-standing feud between Keoland and the Hold of the Sea Princes might be drawing to an end under the leadership of Keoland's current ruler, King Kimbertos Skotti. The region enjoys a warm, mild climate but faces many threats from monstrous foes.

\begin{DndTable}[header=Old Keoland Locations]{XXX}

Location & Ruler & Description\\
Bissel, the March of & Margrave Imran Rendulkar (human) & Bone of contention between Keoland, Veluna, and Ket\\
Geoff, the Grand Duchy of & Grand Duchess Owena Blackthorn (human) & Isolated fiefdom with a long history of battling giants in the nearby mountains\\
Gran March, the & Commandant Magnus Onyxbeard (dwarf) & Nominal vassal of Keoland and ally with Bissel; the commandant is elected from among the March's noble houses\\
Hold of the Sea Princes, the & Prince Zygmund III of Monmurg (human) & Independent oligarchy of sea traders founded by buccaneers; now a powerful naval force\\
Keoland, the Kingdom of & King Kimbertos Skotti (human) & Heart of the Old Keoland region, surrounded by friendly neighbors that swear fealty to Keoland's monarch\\
Lortmil Mountains, the & — & Natural border between the Old Keoland region and the Central Flanaess\\
Sterich, the March of & Marquise Quercha Emondav (human) & Nominal vassal state of Keoland, but its ruler is more like a sister than a vassal to the king of Keoland\\
Ulek, the County of & Countess Lewenn Richfield (human) & Former vassal of Keoland\\
Ulek, the Duchy of & Duke Grenowin (elf) & Former vassal of Keoland with a large population of elves\\
Ulek, the Principality of & Princess Olynn Corond (dwarf) & Fiefdom with a significant navy; its princess commands the respect of many dwarves beyond Ulek\\
Valley of the Mage, the & The Mage of the Valley (identity unknown) & Secluded refuge of an ancient archmage; current inhabitants unknown\\
Yeomanry, the & Freeholder Vyndi Skyspear (goliath) & Independent republic governed by an elected freeholder\\
\end{DndTable}

\subparagraph{Old Keoland Culture}
Old Keoland is a diverse region of the Flanaess where different cultures have mingled for many centuries. Keolish folk often garden, maintain close family ties, have a down-to-earth nature, and love storytelling.

The clothes worn in Old Keoland tend toward loose-fitting shirts and wide-legged pants, voluminous cloaks for cold or wet weather, and sturdy boots. This region's cuisine represents a fusion of Central Flan dishes of rice, potato, and meat with some spices and seasonings brought from the west, creating unique flavors.

\subparagraph{Old Keoland and Its Neighbors}
The long chain of the Lortmil Mountains forms a natural barrier between Old Keoland and the region of the Central Flanaess. The mountains contain some of the richest gem and precious metal deposits in Eastern Oerik. The humans, dragonborn, dwarves, gnomes, halflings, goblinoids, goliaths, and orcs that live in these mountains and their foothills are subjects of the realms that surround the mountain range, but they often band together to deal with greater threats on both sides of the mountains.

\subparagraph{Old Keoland Adventures}
One reason for the amicable relations among the nations of Old Keoland is the danger posed by dragons, giants, and other monsters found throughout the region. That makes this region particularly well suited to a heroic fantasy campaign (see "\textit{Flavors of Fantasy}" in this chapter). The most dangerous places include those described below.

\subparagraph{Barrier Peaks}
These forbidding highlands are home to strange monsters. "Expedition to the Barrier Peaks," an adventure in Quests from the Infinite Staircase, explores the origin of these monsters.

\subparagraph{Crystalmist Mountains}
Giants, white dragons, and other monsters frequently descend from the Crystalmist Mountains into Geoff and Sterich, searching for food and plunder. An ancient tunnel stretches from the western end of the Yeomanry into the Sea of Dust, attracting many adventurers to explore its lengths.

\subparagraph{Dim Forest}
Though elves inhabit the western part of the Dim Forest, the eastern part is wild and prowled by monsters, including green dragons.

\subparagraph{Hellfurnaces}
The Hellfurnaces are a volcanically active extension of the Crystalmist Mountains populated with threats including fire giants and red dragons. Beneath the mountains are labyrinths that connect to the Underdark, wherein lie hidden cities, strongholds, and temples harboring terrible evil.

\subparagraph{Jotens}
Hill giants, manticores, and wyverns from the Jotens regularly threaten the tranquility of both Sterich and the Yeomanry.

\subparagraph{Rushmoors}
Hungry black dragons, otyughs, and other monsters haunt the Rushmoors.

\subsubsection{Western Flanaess}
Survivors of the Invoked Devastation that destroyed the ancient Baklunish empire settled the temperate prairies, forests, and coastal lands of the Western Flanaess about a thousand years ago. Largely separated from the rest of the Flanaess by the Yatils, the Barrier Peaks, and the Crystalmist Mountains, these realms are a stronghold of Baklunish cultures.

The nations of Ekbir, Tusmit, and Zeif represent the heart of the region, and two rivers—the Blashikmund and the Tuflik—form natural borders between them. Although the nations currently enjoy peaceful relations, Tusmit profits by playing its political neighbors against each other—Ekbir against Zeif, Zeif against the Ulakandar nomads, the Ulakandar nomads against Ket, and so forth. Pasha Qharlan Sylba of Tusmit is careful to keep his name well clear of these schemes so he can avoid embroiling Tusmit in open warfare. But Tusmit's spies are currently causing discord by spreading rumors that Zeif is planning to invade Ekbir. Ekbir's sultan believes the rumors are true and is readying his army.

Sultan Naxas Murad of Zeif is a reclusive man, a great philosopher, and a stern father figure to the rulers of Ekbir and Tusmit. Over the years, advisers and family members have urged Naxas to expand Zeif's borders through military conquest, but he refuses to do so, citing failed land grabs by kingdoms through history as proof that imperial expansion across the Flanaess rarely ends well.

\begin{DndTable}[header=Western Flanaess Locations]{XXX}

Location & Ruler & Description\\
Dry Steppes, the & — & Desert where the Baklunish empire once stood\\
Ekbir, the Sultanate of & Sultan Xargun II (aasimar) & Monarchy in an uneasy peace with its neighbors, bracing for a rumored invasion from Zeif\\
Ket & Beygraf Zoltana Lhaz (human) & Crossroads region\\
Plains of the Ulakandar, the & Various Ulakandar clan leaders & Land of the Ulakandar nomads, who roam between the Dry Steppes and the border of Zeif\\
Sea of Dust, the & — & Wasteland where the Suel Imperium once stood\\
Tusmit, the Pashalik of & Pasha Qharlan Sylba (human) & Monarchy profiting by playing its neighbors against each other\\
Ull & Orakhan Drasika Borinok (human) & Independent fiefdom founded by Ulakandar nomads who settled the land\\
Zeif, the Sultanate of & Sultan Naxas Murad (human) & Monarchy ruled by a reclusive philosopher who resists his advisers' call to imperial expansion\\
\end{DndTable}

\subparagraph{Western Flanaess Culture}
The culture of the Western Flanaess preserves many of the ways and traditions of the ancient Baklunish empire. Enormous value is placed on the virtues of hospitality and generosity, particularly almsgiving and pious donations to temples and clergy. Since the fall of their ancient empire, the people of the Western Flanaess have demonstrated more interest in trade—as a way of amassing power and wealth, but also as a means of connecting and coexisting with neighbors—than in imperial expansion or military domination.

The clothing favored by the people of the Western Flanaess features bright patterns in vibrant colors, worn in flowing gowns, robes, and long coats worn with breeches. Their cuisine is renowned for its rich flavors and spiciness.

\subparagraph{Western Flanaess and Its Neighbors}
Ket is the crossroads between the Western Flanaess and the rest of the continent, nestled between the Yatil Mountains and Barrier Peaks. Though it was founded by Baklunish settlers, Ket's proximity to both Bissel (and the lands of Old Keoland) and Veluna (and the Central Flanaess) makes it a vibrant, multicultural land rich from extensive trade. Its culture reflects the breadth of its trade, including sizable populations of dragonborn and tieflings and displaying a fusion of Baklunish and Oeridian influence. This mixture is visible in the title of its ruler, which is a combination of a Baklunish title (bey) and an Oeridian one (graf). The beygraf is a noble chosen by the lords of Lopolla's wealthiest, most influential families. Many of these lords also serve as generals in the Kettish military. The current beygraf, Zoltana Lhaz, is a skilled diplomat who so far has balanced the interests of different forces both inside her nation and among her neighbors.

\subparagraph{Western Flanaess Adventures}
The political and mercantile intrigue among Ekbir, Tusmit, and Zeif provides abundant adventure opportunities for characters in the Western Flanaess. A campaign focused around intrigue or mystery (see "\textit{Flavors of Fantasy}" in this chapter) works particularly well in this region.

Of course, the Western Flanaess has its fair share of monsters, dungeons, and ruins as well. The ruins of the Baklunish Empire in the Dry Steppes, and those of the Suel Imperium in the Sea of Dust, attract plenty of adventurers as well as villains hoping to claim the magical knowledge that caused the terrible catastrophes leading to the fall of those empires. Ket and Ull, too, suffer from the depredations of the monsters in the Barrier Peaks and the Yatil Mountains, just as their eastern neighbors do.

\chapter{Cosmology}
The worlds of D\&D are part of an immense cosmos. Most campaigns and adventures play out on worlds on the Material Plane. The rest of the multiverse consists of different planes of existence defined in relation to the Material Plane.

The planes of existence are strange and often dangerous environments undreamed of in the natural world. Adventurers can stroll along streets of fire, test their mettle on battlefields where the fallen are resurrected with each dawn, and behold the terrifying majesty of the Lady of Pain as she floats above the streets of the ring-shaped city at the center of the multiverse.

\section{The Planes}
The planes of existence are realms of myth and mystery. They're not simply other worlds, but dimensions formed and governed by spiritual and elemental principles. They fall into the following categories:


\begin{description}
\item[Material Realms.]
Most D\&D worlds are located on the Material Plane, which has two planar echoes: the Feywild and the Shadowfell.
\item[Transitive Planes.]
The Ethereal Plane and the Astral Plane are boundless realms that provide passage between other planes of existence.
\item[Inner Planes.]
The four Elemental Planes (Air, Earth, Fire, and Water), plus the Para-elemental Planes between them, are the Inner Planes.
\item[Outer Planes.]
Seventeen Outer Planes correspond to the nine alignments and shades of philosophical difference between them.
\item[Positive and Negative Planes.]
These two planes enfold the rest of the cosmology, providing the raw forces of life and death that underlie all existence in the multiverse.
\end{description}

\subsection{The Great Wheel}
The default D\&D cosmology includes more than two dozen planes, detailed in this chapter. The most common understanding of these planes visualizes them as a group of concentric wheels, with the Material realms at the center. The Inner Planes form a wheel around the Material Plane, enveloped in the Ethereal Plane. Then the Outer Planes form another wheel around and behind (or above or below) that one, arranged according to alignment, with the Outlands linking them all.

Since the primary way of traveling from plane to plane is through magical portals, the spatial relationship between different planes is largely theoretical. No being in the multiverse can look down and see the planes arranged like a diagram in a book. No mortal can verify whether Mount Celestia is sandwiched between Bytopia and Arcadia; rather, this theoretical positioning is based on the philosophical shading among the three planes and the relative importance they give to law and good.

\subsubsection{Other Configurations}
For your campaign, you can use a different model of the planes. Here are several examples:


\begin{itemize}
\item Planes situated among the roots and branches of a great cosmic tree (literally or figuratively)
\item Material Realms suspended between two other realities: the Astral Realms (the Astral Plane and the Outer Planes) above and the Elemental Realms (the Inner Planes) below
\item A cosmology with fewer planes: a Material Plane; the Transitive Planes; a single undifferentiated Elemental Plane, where all four elements churn in chaos; an Overheaven, where good deities and Celestials dwell; and an Underworld, where evil deities and Fiends reside
\item Planes arranged in a complex system of orbits, with planes exerting greater influence on the Material Plane the closer they draw to it
\end{itemize}

\subsection{Material Realms}
The Material Plane is where the philosophical and elemental forces of the other planes of existence collide in the jumbled existence of mortal life and matter. It is a thoroughly magical place, reflected in the two planes that share its central place in the multiverse.

The Feywild and the Shadowfell are parallel dimensions occupying the same cosmological space as the Material Plane. The landscapes of these three planes are similar, but those of the Feywild are more marvelous and whimsical, while those of the Shadowfell are more bleak and ominous. Passage between the Material Plane and these other realms is sometimes effortless, even accidental. Adventurers might enter a grove of trees on the Material Plane and suddenly find themselves in a lush, colorful forest on the Feywild or a grim wood of dead trees on the Shadowfell.

\subsection{Inner Planes}
The Inner Planes surround the Material Plane and its echoes, providing the raw elemental substance from which all worlds were made. The four Elemental Planes—Air, Earth, Fire, and Water—form a ring around the Material Plane. The border regions between these planes are sometimes described as distinct planes in their own right: the Para-elemental Planes.

These realms exemplify the physical essence and elemental nature of air, earth, fire, and water. The entire substance of the Elemental Plane of Fire, for example, is suffused with the fundamental nature of fire: energy, passion, transformation, and destruction. Even objects of solid brass or basalt seem to dance with flame in a manifestation of the vibrancy of fire's dominion.

At their innermost edges, where they are conceptually closest to the Material Plane, the four Elemental Planes and the four Para-elemental Planes resemble places on the Material Plane. The four elements mingle together as they do on the Material Plane, forming land, sea, and sky. But the dominant element strongly influences the environment, altering those locations' fundamental qualities.

The inhabitants of this inner ring include aarakocra, azers, dragon turtles, gargoyles, genies, lizardfolk, mephits, salamanders, and xorn. Some originated on the Material Plane, and all can travel to the Material Plane (if they have access to the magic required) and survive there.

As the Elemental Planes extend farther from the Material Plane, they become increasingly unstable and hostile. In the outer regions, the elements exist in their purest form: great expanses of solid earth, blazing fire, crystal-clear water, and unsullied air. Any foreign substance is extremely rare; little air can be found in the outer reaches of the Plane of Earth, and earth is all but impossible to find in the outer reaches of the Plane of Fire. These areas are much less hospitable to travelers from the Material Plane than the border regions are. Such regions are little known, so one who mentions the Plane of Fire, for example, usually means the border region.

The outer regions are the domains of creatures formed of the pure elements, including air, earth, fire, and water elementals. These are also the domains of the Elemental Princes of Evil—primordial beings of pure elemental fury.

At the outermost extents of the Elemental Planes, the pure elements dissolve and bleed together into an unending tumult of clashing energies and colliding substance called the Elemental Chaos. Elementals can be found here as well, but they usually don't stay long, preferring the comfort of their native planes.

\subsection{Outer Planes}
If the Inner Planes are the raw matter and energy that make up the multiverse, the Outer Planes provide the direction, thought, and purpose for its construction. These are realms of spirituality and thought, the spheres where Celestials, Fiends, and deities dwell. The plane of Elysium, for example, isn't merely a home for good creatures or where spirits of good creatures go when they die. It is the embodiment of goodness, a spiritual realm where evil can't abide. It is as much a state of being and of mind as it is a physical location.

When discussing anything to do with deities and their realms, the language used must be highly metaphorical. Their actual homes aren't literally places at all but exemplify the idea that the Outer Planes are realms of thought and spirit.

The planes with an element of good in their nature are called the Upper Planes, while those with an element of evil are the Lower Planes. A plane's alignment (as shown in the Outer Planes table) is its essence, and a creature whose alignment doesn't match the plane's alignment experiences a sense of dissonance there. When a good creature visits Elysium, for example, it feels in tune with the plane, but an evil creature feels uncomfortable.

\begin{DndTable}[header=Outer Planes]{XX}

Outer Plane & Alignment\\
\textit{Abyss} & Chaotic Evil\\
\textit{Acheron} & Lawful Evil, Lawful Neutral\\
\textit{Arborea} & Chaotic Good\\
\textit{Arcadia} & Lawful Good, Lawful Neutral\\
\textit{Beastlands} & Chaotic Good, Neutral Good\\
\textit{Bytopia} & Lawful Good, Neutral Good\\
\textit{Carceri} & Chaotic Evil, Neutral Evil\\
\textit{Elysium} & Neutral Good\\
\textit{Gehenna} & Lawful Evil, Neutral Evil\\
\textit{Hades} & Neutral Evil\\
\textit{Limbo} & Chaotic Neutral\\
\textit{Mechanus} & Lawful Neutral\\
\textit{Mount Celestia} & Lawful Good\\
\textit{Nine Hells} & Lawful Evil\\
\textit{Outlands} & Neutral\\
\textit{Pandemonium} & Chaotic Evil, Chaotic Neutral\\
\textit{Ysgard} & Chaotic Good, Chaotic Neutral\\
\end{DndTable}

The Upper Planes are the home of Celestials. The Lower Planes are the home of Fiends. The planes in between host their own unique denizens: for example, modrons are Constructs that inhabit Mechanus, and slaadi are Aberrations that thrive in Limbo.

As with the Elemental Planes, one can imagine the perceptible part of the Outer Planes as a border region, while extensive spiritual regions lie beyond ordinary sensory experience. Even in perceptible regions, appearances can be deceptive. Initially, an Outer Plane might appear hospitable and familiar to natives of the Material Plane. But the landscape can change at the whim of a deity or other powerful forces that dwell on the plane, which can remake the realm completely, erasing and rebuilding existence to better fulfill those forces' needs.

Distance is a virtually meaningless concept on the Outer Planes. A perceptible region of a plane might seem quite small on one visit, and on another trip it can stretch on to what seems like infinity. Adventurers could take a guided tour of the Nine Hells, from the first layer to the ninth, in a single day—if the powers of the Nine Hells desire it. Or it could take weeks for travelers to make a grueling trek across a single layer.

\subsubsection{Layers of the Outer Planes}
Most Outer Planes include a number of distinct realms. These environments are often imagined as a stack of related parts of the same plane, so travelers refer to them as layers. For example, Mount Celestia resembles a sacred mountain with seven great plateaus along its ascent, the Nine Hells is like a pit where the River Styx plunges down through nine tiers, and the Abyss has a seemingly endless number of layers.

Like the planes themselves, the description of layers is highly metaphorical and subject to varying interpretations. The plane of Carceri, for example, has been described as a long series of spherical worlds arranged like beads on a string, with each sphere consisting of six nested spheres—the layers of the plane. This fanciful description is but one attempt to make sense of the distorted geography of a place that isn't even a place in the ordinary sense of the word, but an alternate state of reality.

Most portals from elsewhere reach the first layer of a multilayered plane. This layer is depicted as the top or bottom layer, depending on the plane. As the arrival point for most visitors, the first layer functions like an antechamber for that plane.

\subsubsection{Alignment and the Outer Planes}
The Outer Planes are realms of thought and morality more than merely physical reality, and they can affect visitors on a deeply personal level as well as a physical one.

At your discretion, a creature that spends a long time on an Outer Plane that is not its home plan can begin to take on aspects of that plane's ethos. Visitors to the Upper Planes might feel strange urges to perform deeds of kindness or compassion, while visitors to the Lower Planes might find themselves drawn to acts of cruelty or betrayal. Those who spend time on Mechanus and other lawful planes might feel their ties of loyalty to each other growing stronger, while those who visit Limbo and other chaotic planes might become temporarily more independent or self-absorbed. These tendencies are best handled as DM suggestions and then roleplayed by the players, but you might award Heroic Inspiration to characters who bring these characteristics to life in their characters.

\subparagraph{Planar Dissonance}
Celestials who visit the Lower Planes and Fiends who visit the Upper Planes experience significant discomfort if their visits last more than a few hours. After finishing a Long Rest on a plane that is alien to its nature, a Celestial or Fiend makes a DC 10 Constitution saving throw. On a failed save, whenever the creature makes a D20 Test, the creature must subtract 1d4 from the roll. The effect is cumulative with each failed save and ends when the creature finishes a Long Rest on a plane that isn't opposed to its nature.

\section{Planar Travel}
When adventurers travel to other planes of existence, they undertake a legendary journey in which they might face supernatural guardians and undergo many ordeals. The nature of that journey and the trials along the way depend in part on the means of travel, such as magical portals or spells.

\subsection{Planar Portals}
A portal is a stationary, interplanar connection that links a specific location on one plane of existence to a specific location on another. Some portals function like doorways, appearing as a clear window or a fog-shrouded passage, and interplanar travel is as simple as moving through the portal. Other portals are locations—circles of standing stones, soaring towers, sailing ships, or even whole towns—that exist on multiple planes at once or flicker from one plane to another. Some are vortices, joining an Elemental Plane with a very similar location on the Material Plane, such as a swirling pool of magma in the heart of a volcano (leading to the Plane of Fire) or a maelstrom in the depths of the ocean (leading to the Plane of Water).

Passing through a planar portal can be the simplest way to travel from the Material Plane to a desired location on another plane. Often, though, a portal presents an adventure in itself.

First, the adventurers must find a portal that leads where they want to go. Most portals exist in distant locations, and a portal's location often has thematic similarities to the plane it leads to. For example, a portal to Mount Celestia might be located on a mountain peak.

Second, portals often have guardians charged with ensuring that certain creatures don't pass through. A portal's guardian is typically a powerful magical creature, such as a djinni, a sphinx, a titan, or an inhabitant of the portal's destination plane.

Finally, most portals aren't open all the time, but open only in particular situations or when a certain requirement is met. A portal can have any requirement, but the following are the most common:


\begin{description}
\item[Command.]
The portal functions only if a particular command is given. A command is usually a word that can be invoked in any language (including a signed language). Sometimes the command must be given as a character passes through the portal (which is otherwise a mundane doorway, window, or similar opening). Other portals open when the command is given within 15 feet of themselves, and they remain open for 1d12 minutes.
\item[Key Item.]
The portal functions if the traveler carries a particular object; the item acts much like a key to a door. This item can be a common object or a particular one created for that portal. The city of Sigil above the Outlands is known as the City of Doors because it features an overwhelming number of such item-keyed portals.
\item[Random.]
The portal functions for a random period, then shuts down for a similarly random duration. Typically, such a portal allows 1d6 + 6 travelers to pass through, then closes for 1d6 days.
\item[Situation.]
The portal functions only if a particular condition is met. A situation-keyed portal might open on a clear night, when it rains, or when a certain spell is cast in its vicinity.
\item[Time.]
The portal functions only at particular times on the Material Plane: during a full moon, during the spring equinox or winter solstice, or when the stars are in certain positions. Once it opens, such a portal remains open for a limited time, such as for 3 days following the full moon, for 1 hour, or for 10 minutes.
\end{description}

Learning and meeting a portal's requirements can draw characters into further adventures as they chase down a key item, scour old libraries for commands, or consult sages to find the right time to visit the portal.

\subsection{Spells}
A number of spells allow direct or indirect access to different planes of existence. \textit{Gate} and \textit{Plane Shift} can directly transport adventurers to any other plane, with different degrees of precision. \textit{Etherealness} allows adventurers to enter the Ethereal Plane. And \textit{Astral Projection} lets adventurers project themselves into the Astral Plane and from there travel to the Outer Planes.

\subsection{Traveling the Outer Planes}
Described in the sections that follow are four planar features that connect multiple Outer Planes:


\begin{itemize}
\item \textit{The Infinite Staircase}
\item \textit{The River Oceanus}
\item \textit{The River Styx}
\item \textit{Yggdrasil, the World Tree}
\end{itemize}

Other planar crossings might exist in your campaign, or it might be possible to walk (or journey aboard a wondrous train or similar vehicle) from one plane to another in your cosmology.

\subsubsection{Infinite Staircase}
The Infinite Staircase is an extradimensional staircase that connects the planes. An entrance to the Infinite Staircase usually appears as a nondescript door. Beyond the portal lies a small landing with a stairway leading up and down. The Infinite Staircase changes appearance as it climbs and descends, going from simple stairs of wood or stone to a chaotic jumble of stairs hanging in radiant space, where no two steps share the same gravitational orientation. It includes ramps, hovering platforms, and clockwork conveyor belts along its endless construction. The adventure anthology Quests from the Infinite Staircase provides more details about this planar pathway.

The staircase is home to \textbf{Nafas}, a noble genie created by the planar winds that blow into the expanse through its myriad doors. A distant and benevolent observer, Nafas hears wishes spoken throughout the multiverse—wishes he fulfills with the help of adventurers who happen upon his aeolian palace.

Doors to the Infinite Staircase are often tucked away in dusty, half-forgotten places that no one frequents or pays any attention to. On any given plane, multiple doors might lead to the Infinite Staircase, though entrances aren't common knowledge and are occasionally guarded by devas, sphinxes, yugoloths, and other powerful monsters.

\subsubsection{River Oceanus}
The water of the Oceanus is sweet and fragrant, as befits its headwaters in the Blessed Fields of Elysium. This plane-spanning waterway provides a path through some of the Upper Planes. It flows through each of Elysium's layers, passes through the top layer of the Beastlands, streams across the top layer of Arborea, and finally drains away somewhere in Arborea's second layer.

Though it isn't as far-reaching as the Styx, the Oceanus is still a commonly used path between planes and layers. Trading vessels sail up and down its length, and small towns line its banks. Travelers can usually find a boat to hire somewhere along its shores.

\subsubsection{River Styx}
The River Styx bubbles with grease, foul flotsam, and the putrid remains of battles along its banks. The \textit{ill effects of the Styx} are described under "\textit{Hazards}" in \textit{chapter 3}.

The Styx churns through the top layers of Acheron, the Nine Hells, Gehenna, Hades, Carceri, the Abyss, and Pandemonium. Tributaries of the Styx snake through lower layers of these planes. For example, a tendril of the Styx winds through every layer of the Nine Hells, allowing passage from one layer of that plane to the next.

Sinister ferries float on the waters of the Styx, crewed by pilots skilled in negotiating the unpredictable currents and eddies of the river. For a price, these pilots carry passengers from plane to plane. Some pilots are Fiends, while others are the souls of dead creatures from the Material Plane.

\subsubsection{Yggdrasil, the World Tree}
The World Tree, Yggdrasil, is a cosmic ash tree that spans the Outer Planes and links them to many worlds of the Material Plane. Its roots stretch into the Lower Planes, touching Hades, Pandemonium, and possibly other Lower Planes. Most of its massive trunk rises through the plane of Ysgard, and its branches stretch through the Upper Planes and across the Astral to the Material Plane.

Some legends describe a great tree, a seedling of Yggdrasil, that the god Corellon planted and tended on the First World at the dawn of time. When the First World was destroyed, seeds from this tree scattered into the void and took root to form the worlds of the Material Plane. Thus, many philosophers and naturalists view all trees or even all plants as descendants of Yggdrasil, part of a vast network of plant life across the multiverse.

Planar travelers can climb among the roots and branches of Yggdrasil to travel from plane to plane or world to world. Some creatures position themselves as expert guides to this vast cosmic network of branching pathways, constantly studying the ever-changing paths as the tree continues its eternal growth.

\section{Planar Adventuring}
In real-world myths, legends, and literature, venturing onto other planes of existence isn't simply a matter of visiting an unusual environment inhabited by strange creatures; it's a journey with meaning expressed through layers of symbol and metaphor. Adventuring on the planes of the D\&D cosmology is an opportunity to craft similarly meaningful and profound journeys.

Planar adventures ought to be extraordinary. As adventurers reach the medium to high levels of your campaign, they find more cause to venture beyond the confines of the world they call home and explore the planes. If you are ready for the party to undertake a truly mythic quest, such as to find a lost Artifact or an elusive extraplanar entity, the vast multiverse is yours to explore. As their physical and magical abilities are put to the test, the characters might find their emotions, values, and spirits challenged as well.

\subsection{The Blood War}
Throughout history, the teeming demon hordes of the Abyss and the regimented legions of devils from the Nine Hells have battled for supremacy in the cosmos. On worlds of the Material Plane, those who know of the conflict refer to it as the Blood War—a conflict that has raged for millennia, ravaging the Lower Planes.

The battlefields of the Blood War are concentrated in the Nine Hells and the Abyss, though fighting also takes place on the Material Plane, on the planes between those realms on the Great Wheel, and anywhere else demons and devils congregate. Although the intensity of the conflict waxes and wanes, and the front lines of the war can shift drastically, a moment never goes by when demons and devils aren't battling each other somewhere in the multiverse.

Many demons and devils are obsessed with finding some advantage for their side in the Blood War, and powerful mortal adventurers are sometimes drawn into tangled schemes to that end. Additionally, open conflict between Fiends is a constant threat on most of the Lower Planes, even when it is incidental to the missions that bring characters to those planes. A simple quest to find a treasure lost in Hades becomes much less simple when the treasure is located in the midst of a raging battle!

Natives of the Upper Planes also have roles to play in the Blood War. While most of them are content to watch evil feed on itself, they still take steps to minimize the threat to the rest of the multiverse. Whenever the Blood War spills into a location outside the Abyss and the Nine Hells, angels and other emissaries of the gods—including mortals of exceptional valor—stand ready to intervene.

\subsection{Planar Adventure Situations}
You can use the Planar Adventure Situations table instead of the tables in the "\textit{Adventure Situations by Level}" section in \textit{chapter 4} to inspire adventures that draw characters into the planes of existence. These adventure ideas are most appropriate for characters of level 11+.

\begin{DndTable}[header=Planar Adventure Situations]{cX}

1d10 & Situation\\
1 & When magic fails to revive a dead person, the only solution is to venture to the Outer Planes to find the person's spirit and either release it from some prison or convince the person to return to life.\\
2 & People who venture into the woods keep accidentally wandering into the Feywild or the Shadowfell. They might never return, return with no sense of how much time has passed, or return dramatically changed.\\
3 & A long-dead oracle is the only one who knows how a terrible prophecy might be averted, but the cataclysmic fulfillment of the prophecy has already begun.\\
4 & A god has stopped answering prayers and won't respond to any \textit{Commune} spell.\\
5 & A devil has tricked an angel into meddling in the Blood War, and the angel seeks mortal aid.\\
6 & A ancestor of one of the characters must be convinced to bless the character before the full power of the character's bloodline can be unleashed.\\
7 & A foolhardy knight carried a holy weapon on a doomed mission into the Nine Hells, and the powers of Mount Celestia want the weapon and the knight's remains retrieved.\\
8 & A titan is imprisoned on an Outer Plane. The characters might be trying to stop those who seek to release it, or they might want to release it to help defend the world from a greater threat.\\
9 & To prove themselves worthy of an even greater quest, the characters are sent to slay a horrible monster, win the favor of a powerful planar being, negotiate peace between two warring planar factions, or retrieve a long-lost item on another plane.\\
10 & An item of legend is being sold at auction in Sigil, the City of Brass, or some other planar metropolis.\\
\end{DndTable}

\section{Tour of the Multiverse}
Each plane in the multiverse is described below. The planes are presented in alphabetical order.

\subsection{Abyss}
The Abyss embodies all that is perverse, gruesome, and chaotic. Its virtually endless layers spiral downward into ever more appalling forms.

Each layer of the Abyss boasts a horrific environment that is harsh and inhospitable to mortals. Each layer also reflects the entropic nature of the Abyss. Much of the plane seems to be in a decaying, crumbling, or corroded state, and its corruption affects visitors (see "\textit{Curses and Magical Contagions}" in \textit{chapter 3}).

\subsubsection{Layers of the Abyss}
The layers of the Abyss are numbered based on the sequence of their discovery and cataloging by explorers from Sigil. Thus, the Plain of Infinite Portals is identified as the first layer, Azzagrat encompasses the 45th, 46th, and 47th layers; the Demonweb is the 66th layer; and so on. The Layers of the Abyss table presents several infamous layers; detailed descriptions of these layers follow the table.

\begin{DndTable}[header=Layers of the Abyss]{XX}

Layer & Description\\
The Plain of Infinite Portals & On layer 1, corroded iron fortresses defend routes to lower layers.\\
Azzagrat & Graz'zt's corrupt, decadent city is split across layers 45, 46, and 47.\\
The Demonweb & On layer 66, Lolth's web snares all and hides portals to other planes.\\
Gaping Maw & Layer 88 is a malevolent wilderness surrounding Demogorgon's ocean fortress.\\
Thanatos & On layer 113, an endless graveyard hosts Orcus and the sleepless dead.\\
The Slime Pits (Shedaklah) & Layer 222, a fetid realm of ooze and fungi, obeys the whims of Juiblex and Zuggtmoy.\\
The Death Dells & Yeenoghu and his gnoll servants prowl layer 422—a cruel, desolate realm.\\
The Endless Maze & Layer 600's endless labyrinth turns visitors into Baphomet's prey.\\
\end{DndTable}

\subsubsection{Layer 1: The Plain of Infinite Portals}
This layer is the miserable gateway to the infinite layers of the Abyss. Under a glaring red sun, the rocky ground contains gaping craters that are portals to the other layers of the Abyss. Other portals lead to Pandemonium, Sigil, the gate-town of Plague-Mort in the Outlands, and the Astral Plane, making this layer the best way to escape the horrors of the Abyss. Iron fortresses dot the landscape, homes to petty lords and upstart demons that are as changeable as the Abyss itself.

The portal leading to Plague-Mort is tucked within a fortress called the Broken Reach, ruled by a succubus named Red Shroud. In the Broken Reach, those who can prove their strength and mettle can stay unharmed for a few days at least.

\subsubsection{Layers 45–47: Azzagrat}
The demon lord Graz'zt embodies manipulation and cruelty, tempting mortals with the promise of appalling delights and decadent luxuries. He rules over the realm of Azzagrat, which encompasses three interconnected layers of the Abyss. His seat of power is the fantastic Argent Palace in the city of Zelatar, whose bustling markets and pleasure palaces draw visitors from across the multiverse in search of obscure magical lore and perverse delights. By Graz'zt's command, the demons of Azzagrat present a veneer of civility and courtly comity. However, the so-called Triple Realm holds as much danger as any other part of the Abyss, and planar visitors can vanish without a trace in its mazelike cities and in forests whose trees have serpents for branches.

\subsubsection{Layer 66: The Demonweb}
Lolth is the Demon Queen of Spiders, whose schemes entangle entire civilizations on worlds across the multiverse. Of all demon lords, she might have the most active interest in the worlds of the Material Plane and in the cultists who do her bidding on those worlds, but her interest lies only in domination.

Lolth's layer is an immense network of thick, magical webbing that forms passageways and cocoonlike chambers. Structures, ships, and other objects are caught in the webbing. The webs conceal random portals that snare objects from demiplanes and Material Plane worlds that figure into the schemes of the Spider Queen. Lolth's servants also build dungeons amid the webbing, trapping and hunting Lolth's hated enemies within crisscrossing corridors of web-mortared stone. Far beneath these dungeons lie the bottomless Demonweb Pits where the Spider Queen dwells with her most loyal servants—yochlol demons created to serve her that outrank mightier demons while in the Spider Queen's realm.

\subsubsection{Layer 88: The Gaping Maw}
The Sibilant Beast and the self-styled Prince of Demons, Demogorgon yearns for nothing less than undoing the order of the multiverse. A two-headed monster who seems as much in conflict with himself as with other beings, the Prince of Demons inspires fear and hatred among other demons and demon lords.

Demogorgon's layer is a vast wilderness of brutality and horror known as the Gaping Maw, where even powerful demons are overcome by fear. Reflecting Demogorgon's dual nature, the Gaping Maw consists of a massive primeval continent covered in dense jungle, surrounded by a seemingly endless expanse of ocean and brine flats. The Prince of Demons rules his layer from two serpentine towers, which emerge from a turbid sea. Each tower is topped with an enormous fanged skull. The spires constitute the fortress of Abysm, where echoes of Demogorgon's turbulent thoughts resound through the halls, tearing at the minds of creatures who dare to enter.

\subsubsection{Layer 113: Thanatos}
Known as the Demon Prince of Undeath and the Blood Lord, the demon lord Orcus is worshiped by Undead and by living creatures that channel the power of undeath. A brooding and nihilistic entity, Orcus yearns to make the multiverse a place of death and despair, forever unchanging except by his will, and to turn all creatures into Undead under his control.

Orcus's realm of Thanatos is a land of bleak mountains, barren moors, ruined cities, and forests of twisted black trees under a black sky. Tombs, mausoleums, gravestones, and sarcophagi litter the landscape. Undead swarm across the plane, bursting from their tombs and graves to tear apart any creatures foolish enough to journey here.

Orcus rules Thanatos from a vast palace known as Everlost, crafted of obsidian and bone. Set in a howling wasteland called Oblivion's End, the palace is surrounded by tombs and graves dug into the sheer slopes of narrow valleys, creating a tiered necropolis.

\subsubsection{Layer 222: The Slime Pits}
Also known as Shedaklah, this layer is ruled by two separate yet equally repugnant demon lords—Juiblex and Zuggtmoy—who coexist with little conflict.

Juiblex, the Demon Lord of Slimes and Oozes, is an amorphous mass that lurks in the abyssal depths. The wretched Faceless Lord cares nothing for cultists or mortal servants, and its sole desire is to turn all creatures into formless copies of its horrid self. Zuggtmoy is the Demon Queen of Fungi and the Lady of Rot and Decay. Her primary desire is to infect the living with spores, transforming them into her servants and, eventually, into decomposing hosts for the mushrooms, molds, and other fungi that she spawns.

As the name suggests, the Slime Pits is a bubbling morass of fetid sludge. The landscape is covered in vast expanses of caustic slimes, and strange organic forms rise from the oceans of ooze at Juiblex's command. Zuggtmoy's palace consists of two dozen immense mushrooms, among the largest in existence, hollowed into grand chambers and twisting corridors. The palace is surrounded by a field of acidic puffballs and poisonous vapors.

\subsubsection{Layer 422: The Death Dells}
The demon lord Yeenoghu hungers for slaughter and senseless destruction. Gnolls are his instruments on the Material Plane, and he drives them to ever-greater atrocities in his name. Delighting in sorrow and hopelessness, the Gnoll Lord yearns to turn the cosmos into a wasteland in which the last surviving gnolls tear one another apart for the right to feast upon the dead.

Yeenoghu rules a layer of ravines and badlands known as the Death Dells. Here, creatures must hunt to survive. Even plants try to snare the unwary to bathe their roots in blood. Yeenoghu's servants, helping to sate their master's hunger as he prowls his kingdom, capture creatures from the Material Plane for release in the Gnoll Lord's realm.

\subsubsection{Layer 600: The Endless Maze}
The demon lord Baphomet, also known as the Horned King and the Prince of Beasts, embodies bestial bloodlust. If he had his way, civilization would crumble and all mortals would embrace their predatory instincts.

Baphomet's layer is a never-ending dungeon with the Horned King's enormous palace at its center. A confusing jumble of crooked hallways and myriad chambers, the palace is surrounded by a mile-wide moat concealing a confounding series of submerged stairs and tunnels leading deeper into the fortress.

\subsubsection{Abyss Adventures}
The Abyss embodies the loathsome corruption of chaos and evil. A descent into the Abyss is a journey into a hostile and uncharted environment. It's also an opportunity to confront the evil in one's own heart and to resist the temptation to turn against allies in order to survive. Heroic characters might make a desperate last stand against endless hordes of demons here, or they might try to avoid detection while seeking a holy relic left behind by some lost hero who dared to venture here.

It's the nature of the Abyss to contaminate the other planes it touches. Opening a portal to the Abyss from any other plane allows the Abyss to create tiny pockets of Abyssal evil that can eventually become so corrupted that they're drawn into the Abyss. Thus, adventurers exploring a desecrated temple or fetid swamp on the Material Plane can unexpectedly find themselves in a demon-infested environment very much like the Abyss without ever leaving their home plane.

\subsection{Acheron}
Acheron is made of immense iron blocks whose metallic surfaces ring beneath the marching feet of endless armies. These blocks drift through an airy void, sometimes colliding with a fearsome clang, crushing all between them and sending shudders through the plane.

Acheron has four layers, with the largest blocks gravitating to the top layer. Some scholars have theorized that the crashing blocks of the upper layers are eventually broken down into smaller chunks of matter that sink to the lower layers. The truth is actually the opposite: the tiny shards of Ocanthus, the lowest layer, break off from an icy mire in its deepest recesses and are gradually assembled and organized, perhaps under the influence of nearby Mechanus, into the perfect cubes of Avalas.

The nature of Acheron instills \textit{bloodlust} in those who visit the plane (see "\textit{Environmental Effects}" in \textit{chapter 3}).

\begin{DndTable}[header=Layers of Acheron]{XX}

Layer & Description\\
Avalas & Spirit soldiers wage endless wars across debris-strewn battlefields.\\
Thuldanin & Pitted, hollow cubes are filled with the cast-off machinery of war.\\
Tintibulus & Jagged blocks tumble in darkness, holding frozen memories sapped by the River Styx and crystallized into fleeting images.\\
Ocanthus & Maelstroms of razor-sharp debris swirl above a mire of black ice.\\
\end{DndTable}

\subsubsection{Acheron Adventures}
Acheron is a plane of enforced order, where rigid conformity leads to crushed spirits and broken hopes. The spirits here can't conceive of anyone refusing to obey the will of their commanders. They are dedicated soldiers, forever lacking a cause.

A journey into Acheron is a confrontation with the bleak nihilism of unending conflict, the harsh reality of authoritarian rule, and the uncaring pressures of social conformity. It's also an opportunity for characters to grapple with the question of what they are fighting for, among armies that have forgotten how to even ask the question.

On a more literal level, an adventure in Acheron can involve preventing a villain from scavenging Thuldanin for some new horror of warfare to be unleashed on the battlefield. Or it might require retrieving a secret from the imprisoned thoughts or memories found in the blocks of Tintibulus.

\subsection{Arborea}
Arborea is a plane of extremes: stupendously craggy mountains; unbelievably deep gorges; forests of monstrously huge trees; and vast stretches of wheat fields, orchards, and arbors. Wild-hearted nature spirits dwell in every glade and stream, brooking no infringement. Travelers must tread lightly.

The air of Arborea seems charged with excitement. Sudden squalls brew up out of nowhere, beating the tree-lined paths with heavy winds. The storms pass within minutes and leave behind warm arcs of sunlight filtering through the forest canopy. Music always seems to be playing in the distance; sometimes it originates from groups of elf musicians, but just as often the faint tune is merely the wind curling through the boles of the great trees.

\begin{DndTable}[header=Layers of Arborea]{XX}

Layer & Description\\
Arvandor & Towering trees, colorful wildflowers, abundant grain, and delicious fruit create a lush environment.\\
Aquallor & An eternal ocean fed by the River Oceanus is home to teeming sea life and mighty storms.\\
Mithardir & A desert of white sand is abandoned by whatever powers once lived there.\\
\end{DndTable}

\subsubsection{Arborea Adventures}
Arborea is a larger-than-life place of violent moods and deep affections, of whim backed by steel, and of passions that blaze brightly until they burn out. Its good-natured inhabitants are dedicated to fighting evil, but their reckless emotions sometimes break free with devastating consequences. Rage is as common and as honored as joy in Arborea.

An adventure in Arborea can be an opportunity for characters to discover who they are when masks fall away and the honesty of unfettered emotion is revealed. The inhabitants of the plane are accustomed to this emotional honesty. Lifelong friends might share a laugh over food and wine, cross blades over a mutual lover, and write songs celebrating each other's courage and integrity, all in a single evening. For those who aren't accustomed to this candor, though, it can lead to hurt feelings and lingering resentment.

Creatures that visit Arborea and then leave sometimes experience a desperate desire to return—a yearning so intense that it can interfere with day-to-day life.

What secrets lie buried in the sands of Mithardir? An expedition might involve investigating whatever gods or Celestials once inhabited the silver desert or find some knowledge they possessed.

\subsection{Arcadia}
Arcadia thrives with orchards of perfectly lined trees, ruler-straight streams, orderly fields, immaculate roads, and cities laid out in geometrically pleasing shapes. The mountains bear no trace of erosion.

Night and day are determined by an orb that floats in the sky above both of Arcadia's layers. Half of the orb radiates sunlight and brings about the day; the other half sheds moonlight and brings on the starry night. The orb rotates evenly without fail, spreading day and night across the entire plane.

The weather in Arcadia is governed by four allied demigods called the Storm Kings: the Cloud King, the Wind Queen, the Lightning King, and the Rain Queen. Each one lives in a castle surrounded by the type of weather that ruler controls.

Arcadia is suffused with a vigorous life energy that bestows \textit{great vitality} on visitors (see "\textit{Environmental Effects}" in \textit{chapter 3}).

\begin{DndTable}[header=Layers of Arcadia]{XX}

Layer & Description\\
Abellio & Everything in these fields of plenty is dedicated to the good of all.\\
Buxenus & Military forces gather their strength, amid pleasant valleys and orchards, to reclaim the lost layer of Menausus—now part of Mechanus.\\
\end{DndTable}

\subsubsection{Arcadia Adventures}
Everything on Arcadia works toward the common good and a flawless existence. Here, purity is eternal, and nothing intrudes on harmony. Individuality is subsumed to peaceful life in community, and military might can be brought to bear on those who disrupt that peace.

An adventure in Arcadia can be an opportunity to explore the tension between individual freedom and societal responsibility for the common good. Even lawfully inclined adventurers rarely conform neatly to social expectations, and a visit to this plane can highlight that conflict.

\subsection{Astral Plane}
The Astral Plane is a realm of thought and dream, where visitors travel as disembodied souls to reach the Outer Planes. It is a great silvery sea, the same above and below, with swirling wisps of white and gray streaking among motes of light—the distant stars of far-flung Wildspace systems. Most of the Astral Sea is a vast, empty expanse. Visitors occasionally stumble upon the petrified corpse of a dead god or other chunks of rock drifting forever in the silvery void. Much more commonplace are color pools—magical pools of colored light that flicker like radiant, spinning coins.

Creatures on the Astral Plane don't age or suffer from hunger or thirst. For this reason, creatures that live on the Astral Plane (such as githyanki) establish outposts on other planes, often the Material Plane, so their children can grow to maturity.

\subsubsection{Navigating the Astral Plane}
A traveler in the Astral Plane can move by simply thinking about moving, but distance has little meaning. In combat, though, a creature has a Fly Speed (in feet) equal to 5 times its Intelligence score and can hover.

Just as movement is accomplished by the power of thought, all that is required to find one's destination is to think about it. As long as the destination is somewhere in the Astral Plane (or in Wildspace, as described below)—such as "the nearest githyanki outpost," "the nearest color pool leading to the Abyss," or "the Wildspace system of Realmspace"—thinking about a place makes the creature aware of the most direct route to that location. The creature doesn't know how long the journey will take or how perilous it will be, just which direction to go in.

The DM decides how long it takes to get to a desired destination. A trek to a specific location—a particular Wildspace system or Astral outpost, for example—might take 4d6 days. For a more general location, such as a color pool leading to a specified plane, the journey might take 1d4 × 10 hours.

\subsubsection{Dead Gods}
The Astral Plane is where the petrified remains of dead gods end up—gods who were slain by more powerful entities or who lost all their mortal worshipers and perished as a result. A dead god looks like a gigantic, nondescript stone statue that bears little resemblance to the divine entity it once was. Githyanki, mind flayers, and other residents of the Astral Plane sometimes turn these drifting hulks into outposts and cities, many of which are hollowed out beneath the surface. The githyanki city of Tu'narath is perhaps the most infamous such place.

\subsubsection{Color Pools}
Gateways leading from the Astral Plane to other planes appear as two-dimensional pools of rippling colors, 1d6 × 10 feet in diameter. ("Color," as with everything in the Astral Plane, is a matter of metaphor; since these portals are perceived by the Astral self and not by physical eyes, their colors are understood rather than seen.) Traveling to another plane requires locating a color pool that leads to the desired plane. These gateways can be identified by color, as shown on the Astral Color Pools table.

\begin{DndTable}[header=Astral Color Pools]{cXX}

1d20 & Plane & Pool Color\\
1 & Abyss & Amethyst\\
2 & Acheron & Flame red\\
3 & Arborea & Sapphire blue\\
4 & Arcadia & Saffron\\
5 & Beastlands & Emerald green\\
6 & Bytopia & Amber\\
7 & Carceri & Olive\\
8 & Elysium & Orange\\
9 & Ethereal Plane & Spiraling white\\
10 & Gehenna & Russet\\
11 & Hades & Rust\\
12 & Limbo & Jet black\\
13–14 & Material Plane & Silver\\
15 & Mechanus & Diamond blue\\
16 & Mount Celestia & Gold\\
17 & Nine Hells & Ruby\\
18 & Outlands & Leather brown\\
19 & Pandemonium & Magenta\\
20 & Ysgard & Indigo\\
\end{DndTable}

\subsubsection{Wildspace}
Bobbing in the Astral Plane like corks in an ocean are vast, airless expanses called Wildspace systems. In these systems, the Astral Plane overlaps with the Material Plane, and the stars and planets of the Material Plane are accessible from the Astral Plane. Every world of the Material Plane is situated in a Wildspace system.

As an Astral traveler approaches a Wildspace system, the silver fog of the Astral Plane gradually thins until it falls away in Wildspace. Then the sun of the Wildspace system comes into view—often millions of miles away—along with colorful gas clouds, planets, moons, and other cosmic bodies.

A Wildspace system teems with space-dwelling life-forms, including spores, plankton, and larger creatures that resemble fish and aquatic mammals. Creatures and objects in Wildspace age normally and exist on both the Astral Plane and Material Plane simultaneously. This overlap enables creatures to use spells such as \textit{Teleport} to travel from Wildspace to a nearby world, or vice versa. A creature or ship traveling from one Wildspace system to another must cross the Astral Plane unless it has some other magical means of traveling from one world to another. (See "\textit{Material Plane}" in this chapter.) \textit{Spelljammer: Adventures in Space} contains extensive information about Wildspace and Astral travel.

\subsubsection{Psychic Wind}
A psychic wind is a storm of thought that batters travelers' minds rather than their bodies. The storm is made of lost memories, forgotten ideas, and subconscious fears that went astray in the Astral Plane and conglomerated into this powerful force.

A psychic wind is first sensed as a rapid darkening of the silver-gray sky. After 1d4 minutes, the area becomes as dark as a moonless night. As the sky darkens, the traveler feels buffeting and shaking, as if the plane were rebelling against the storm. As quickly as it comes, the psychic wind passes, and the sky returns to normal in 1 minute.

A group of travelers journeying together is subjected to one location effect, determined by consulting the Psychic Wind Locations Effects table.

\begin{DndTable}[header=Psychic Wind Location Effects]{cX}

1d20 & Location Effect\\
1–8 & Diverted; add 1d6 days to travel time\\
9–12 & Blown astray; add 3d10 days to travel time\\
13–16 & Lost; at the end of the travel time, the characters arrive at a location other than the intended destination\\
17–20 & Sent through a random color pool (roll on the \textit{Astral Color Pools} table)\\
\end{DndTable}

Each traveler caught in a psychic wind makes a DC 15 Intelligence saving throw. On a failed save, the traveler suffers a random effect from the Psychic Wind Psychic Effects table as well.

\begin{DndTable}[header=Psychic Wind Psychic Effects]{cX}

1d20 & Psychic Effect\\
1–8 & You have the Stunned condition for 1 minute; you repeat the saving throw at the end of each of your turns, ending the effect on yourself on a success.\\
9–12 & You take 11 (2d10) Psychic damage.\\
13–17 & You take 22 (4d10) Psychic damage.\\
18–20 & You have the Unconscious condition for 5 (1d10) hours; the effect on you ends if you take damage or if another creature takes an action to shake you awake.\\
\end{DndTable}

\subsubsection{Astral Plane Adventures}
Characters most often visit the Astral Plane as a way of getting somewhere else—either to one of the Outer Planes or to different worlds of the Material Plane via Wildspace. En route, they might encounter fellow travelers, such as Celestials, Fiends, slaadi, modrons, or githyanki.

As a realm of thought, memory, and dream, the Astral Plane can also be an adventure destination. Characters might try to plumb the crystallized thoughts of dead gods or sift information from the torrent of a psychic wind. Or they could face Astral manifestations of their own memories, fears, and dreams.

\subsection{Beastlands}
The Beastlands is a plane of nature unbound, of forests ranging from moss-hung mangroves to snow-laden pines, of thick jungles where the branches are woven so tight that no light penetrates, of vast plains where grains and wildflowers wave in the wind with vibrant life. The plane embodies nature's wildness and beauty, but it also speaks to the animal within all living creatures—not necessarily in a fierce, predatory way but with respect to their physical substance and fundamental needs. The spirits of the dead on the Beastlands typically take animal forms or part-animal forms (such as centaurs). Among the greatest inhabitants of this plane are the primal spirits called animal lords.

Whenever a visitor slays a Beast native to the plane, the slayer must succeed on a DC 10 Charisma saving throw or shape-shift into the type of Beast that was slain. The creature's game statistics are replaced by the Beast's stat block, but the creature retains its alignment, personality, creature type, Intelligence score, and ability to speak. At the end of each Long Rest, the shape-shifted creature repeats the save. On a successful save, the creature returns to its true form. After three failed saves, the transformation can be undone only by a \textit{Remove Curse} spell or similar magic.

\begin{DndTable}[header=Layers of the Beastlands]{XX}

Layer & Description\\
Krigala & The River Oceanus is a strong torrent flowing through this wilderness of eternal summer noon.\\
Brux & A red sun hovers forever on the horizon as mists and streams roil through the trees in eternal twilight.\\
Karasuthra & A pale moon provides the only light in this wilderness of eternal night.\\
\end{DndTable}

\subsubsection{Beastlands Adventures}
The Beastlands embodies wild, beautiful nature and the vibrant power of life thriving in the natural world. Visitors to the plane feel invigorated and more vital—their minds more alert, their reflexes sharpened, and their strides quickened. Hunger pangs are acute, but food and drink taste better than ever before. Sleep is always deep and restful, and sleepers always awaken alert.

Adventures in the Beastlands might explore the ways that good philosophies, while purporting to value life, actually devalue the physical nature of life in favor of abstract concepts of law and ethics. Characters might end up in conflict with those who dismiss animals as inferior and irrelevant forms of life.

\subsection{Bytopia}
The surfaces of Bytopia's two layers face each other like the covers of a closed book. Looking up from Dothion, the "top" layer of the plane, a traveler can see Shurrock, its other layer, about a mile overhead. Both layers are idealized worlds that reflect the plane's philosophy of personal achievement alongside social interdependence.

\begin{DndTable}[header=Layers of Bytopia]{XX}

Layer & Description\\
Dothion & Farms nestled among well-tamed woods are hubs of pastoral activity and individual industry.\\
Shurrock & Small communities thrive around quarries and mills amid rough country and harsh weather.\\
\end{DndTable}

\subsubsection{Bytopia Adventures}
Bytopia is the heaven of productive work, the satisfaction of a job well done. The goodness flowing through the plane creates feelings of goodwill and happiness in creatures dwelling there. While Dothion rewards those who seek a quiet life, Shurrock is the paradise of those who continually challenge and better themselves.

The two layers of Bytopia are often referred to as the "Twin Paradises," and it's said that every action carried out on one layer has repercussions on the other—an equal and opposite reaction, though a more metaphorical than physical one. An adventure in Bytopia might challenge characters to recognize the impact their actions have on the world by witnessing the mirrored reactions to their deeds on the opposite layer.

\subsection{Carceri}
The grim inspiration for all other prisons in existence, Carceri is a plane of desolation and despair. Its six layers hold vast bogs, fetid jungles, windswept deserts, jagged mountains, frigid oceans, and black ice. All form a miserable home for the traitors and backstabbers trapped on this prison plane.

\begin{DndTable}[header=Layers of Carceri]{XX}

Layer & Description\\
Orthrys & The River Styx meanders through a layer of vast bogs and quicksand.\\
Cathrys & The stench of decay hangs over fetid jungles and scarlet plains.\\
Minethys & Stinging sand blows in unending storms, hiding the ruins of the ancient city Payratheon.\\
Colothys & Deep chasms cut between cruel mountains make foot travel nearly impossible.\\
Porphatys & Cold, acidic oceans are fed by constant black snow.\\
Agathys & Black ice streaked with red covers this frigid realm.\\
\end{DndTable}

\subsubsection{Prison Plane}
No one can leave Carceri easily. Magical efforts to leave the plane by any spell other than \textit{Wish} simply fail. Portals and gates that open onto the plane become one-way only. Secret passages out of the plane exist, but they are hidden and well guarded by traps and deadly monsters. And though the River Styx flows between Carceri and its neighbors, the passage is extremely dangerous, and ferries leading out of Carceri are both rare and expensive.

\subsubsection{Carceri Adventures}
Carceri is a sunless plane of despair, of passions and poisons, and of realm-shattering betrayals. Hatreds run like a deep, slow-moving river, and there's no telling what the flood of treachery will consume next. It is said that prisoners can escape Carceri only by becoming stronger than whatever force imprisoned them there, but that's a difficult task on a plane whose very nature breeds despair and betrayal. The same tendency toward betrayal prevents those who are imprisoned here from working together for long toward the common goal of escape.

Adventures on Carceri might explore the forces—spiritual and psychological as well as physical and outright demonic—that keep characters trapped or imprisoned there. Characters might help a prisoner escape, from an unjustly held spirit to some primeval god banished to oblivion here.

\subsection{Demiplanes}
Demiplanes are extradimensional spaces that come into being by a variety of means and boast their own physical laws. Some are created by spells. Others are natural folds of reality pinched off from the rest of the multiverse. Theoretically, a \textit{Plane Shift} spell can carry travelers to a demiplane, but the proper frequency required for the tuning fork can be extremely hard to acquire. The \textit{Gate} spell is more reliable, assuming the caster knows of the demiplane.

A demiplane can be as small as a single chamber or large enough to contain an entire realm. For example, a \textit{Mordenkainen's Magnificent Mansion} spell creates a demiplane consisting of a foyer with multiple adjoining rooms, while the land of \textit{Barovia} exists entirely within a demiplane under the sway of its vampire lord, \textit{Strahd von Zarovich}. When a demiplane is connected to the Material Plane or some other plane, entering the demiplane can be as simple as moving through a portal or passing through a wall of mist.

\subsubsection{Demiplane Adventures}
Demiplanes are limited realities shaped according to the will of whoever created them. Adventures in demiplanes might let characters explore how they would shape reality to suit their desires and ideals, or confront distortions of reality crafted by villains.

\subsection{Elemental Plane of Air}
The Plane of Air is home to constant winds of varying strength. Here and there, chunks of earth drift in the openness, many covered with lush vegetation. These earth motes serve as homes for djinn and other natives of the plane. Other creatures live on cloud banks infused with magic to become solid surfaces, supporting towns and castles.

Drifting clouds can obscure visibility in any direction in the plane. Storms are frequent, ranging from strong thunderstorms to fierce tornadoes or mighty hurricanes. The air is mild, except near the Para-elemental Planes at either end of the plane, where the temperature is more extreme. Rain and snow fall only in the part of the plane nearest to the Para-elemental Plane of Ice.

Although few landmarks distinguish one area of the Plane of Air from any other, the following are notable features of the plane:

\subsubsection{Aaqa}
Here and there are hidden realms reachable only by following a particular sequence of flowing winds. Fabled Aaqa is one such realm, a shining domain of silver spires and verdant gardens atop a fertile earth mote. The Wind Dukes of Aaqa are dedicated to law and good, and they maintain a vigilant watch against the depredations of Elemental Evil. They are served by aarakocra.

\subsubsection{Labyrinth Winds}
Much of the Plane of Air is a complex web of air streams, currents, and winds called the Labyrinth Winds. These range from stiff breezes to howling gales that can rip a creature apart. Even the most skilled flying creatures must navigate these currents carefully, flying with the winds, not against them.

\subsubsection{Mistral Reach}
Located near the Para-elemental Plane of Ice, the Mistral Reach is a region of frigid winds and driving snowstorms. Earth motes in the reach are covered with snow and ice.

\subsubsection{Sirocco Straits}
The Sirocco Straits is the region of the plane nearest to the Para-elemental Plane of Ash, where hot, dry winds scour earth motes into barren chunks of rock.

\subsubsection{Plane of Air Adventures}
The essential nature of air is movement, animation, and inspiration. Air is the breath of life, the winds of change, the fresh breeze that clears away the fog of ignorance and the stuffiness of old ideas.

When turned toward wicked ideals by cultists of Elemental Evil, elemental air represents destructive power turned to vengeful ends. Cultists of Evil Air wield howling storms to forcefully express their personal freedom or lay claim to things they feel they have been wrongfully denied.

\subsection{Elemental Plane of Earth}
The Plane of Earth is a chain of mountains rising higher than any mountain range on the Material Plane. It has no sun of its own, and no air surrounds the peaks of its highest mountains. Most visitors to the plane arrive by way of vast caverns that honeycomb the mountains.

Important features of the Plane of Earth include the following:

\subsubsection{City of Jewels}
The plane's largest cavern, called the Great Dismal Delve or the Sevenfold Mazework, is home to the City of Jewels—the capital city of the dao. The dao take great pride in their wealth and send teams across the plane in search of new veins of ore and gemstones. Thanks to these expeditions, every building and significant object in the city is made from precious stones and metals, including the slender gemstone-inlaid spires that top most buildings. The city is protected by a powerful spell that alerts the entire population if a visitor steals even a single stone.

\subsubsection{Furnaces}
The Furnaces are the mountains nearest the Para-elemental Plane of Magma. Lava seeps through their caverns, and the air reeks of sulfur. The dao have great forges and smelting furnaces here to process ores and shape precious metals.

\subsubsection{Mud Hills}
The Mud Hills abut the swampy Para-elemental Plane of Ooze. Landslides wear away the slopes of the hills, sending cascades of earth and stone into the bottomless swamp. The Plane of Earth constantly regenerates the land, pushing new hills up as the old ones erode to nothing.

\subsubsection{Plane of Earth Adventures}
Earth symbolizes stability, rigidity, stern resolve, and tradition. The plane's position opposite the Plane of Air in the ring of the Elemental Planes reflects its opposition to almost everything air represents.

Elemental Evil views earth instead as an implacable force of destruction, perfectly willing to crush venerable institutions and respected traditions in its advance. Cultists of Evil Earth crave the power to destroy the works of civilization with landslides, sinkholes, or mighty earthquakes, and they believe the earth thirsts for the blood of those who don't venerate it properly.

\subsection{Elemental Plane of Fire}
A blazing sun hangs at the zenith of a golden sky above the Plane of Fire, waxing and waning on a 24-hour cycle. It ranges from white hot at noon to deep red at midnight, so the darkest hours of the plane display a deep-red twilight. At noon, the light is intense. Most business in the City of Brass (see below) takes place during the darker hours.

The weather on the plane is marked by fierce winds and thick ash. Although the air is breathable, creatures not native to the plane must cover their mouths and eyes to avoid stinging cinders. The efreet use magic to keep the cinder storms away from the City of Brass, but elsewhere in the plane, the wind always blows, sometimes rising to hurricane force during the worst storms.

\textit{The heat} on the Plane of Fire is comparable to a hot desert on the Material Plane and poses a similar threat to travelers (see "\textit{Environmental Effects}" in \textit{chapter 3}). Sources of water are rare, so travelers must carry their own supplies or produce water by magic.

Important features of the Plane of Fire include the following:

\subsubsection{Cinder Wastes}
The Plane of Fire is dominated by a great expanse of black cinders and embers crossed by rivers of lava. Roving bands of salamanders battle each other, raid azer outposts, and avoid patrols from the City of Brass. Obsidian ruins dot the desert—remnants of forgotten civilizations.

\subsubsection{City of Brass}
Perhaps the best-known location on the Inner Planes is the City of Brass, which stands on the shores of the Sea of Fire. This is the fabled city of the efreet, and its ornate spires and spiked walls reflect the efreet's grandiose and cruel nature. True to the nature of the Plane of Fire, everything in the city seems alive with dancing flames, reflecting the vibrant energy of the place. The heart of the city is the formidable Charcoal Palace, where the tyrannical emperor of the efreet reigns supreme, surrounded by nobles and a host of servants, guardians, and sycophants.

\subsubsection{Sea of Fire}
Lava flows through the Fountains of Creation toward the Para-elemental Plane of Ash and pools into a great expanse of lava called the Sea of Fire, traversed by efreeti and azer sailors in great brass ships. Islands of obsidian and basalt jut up from the sea, dotted with ancient ruins and the lairs of powerful red dragons.

\subsubsection{Torchy's}
Built atop a tall basalt crag in the middle of a lava river is an iron-walled tavern that is most easily reached by hot-air balloon. The proprietor is a sentient Flame Tongue Mace named Torchy, who sells a fine ale and seems to have a new wielder every few months. Torchy's is a popular hangout for ballooning enthusiasts.

\subsubsection{Elemental Fire Adventures}
Fire represents vibrancy, passion, and change. At its best, fire reflects the light of inspiration, the warmth of compassion, and the flame of desire.

The cults of Elemental Evil represent fire at its worst: cruel and wantonly destructive. Cultists of Evil Fire seek the power to burn away the impurities of the world with volcanic eruptions, uncontrolled wildfires, heat waves, and droughts, transforming the Material Plane into a mirror of the Cinder Wastes.

Adventurers frequently come to the City of Brass on quests for legendary magic. If it's possible to buy magic items in your campaign, the City of Brass is the most likely place to find any item for sale. The efreet are fond of trading in favors, especially when they have the upper hand in negotiations. Perhaps a magical contagion or poison can be cured only with something that must be purchased in the markets of the city.

\subsection{Elemental Plane of Water}
The Plane of Water is an endless sea, called the Sea of Worlds, dotted with atolls and islands that rise up from enormous coral reefs that seem to stretch forever into the depths. The storms that move across the sea sometimes create temporary portals to the Material Plane and draw ships into the Plane of Water. Surviving vessels from countless worlds and navies ply these waters with little hope of ever returning home.

A warm sun arcs across the sky of the Plane of Water, seeming to rise and set from within the water at the horizon. Several times a day, the sky clouds over and releases a deluge of rain, often accompanied by spectacular shows of lightning, before clearing up again. At night, a glittering array of stars and auroras bedecks the sky.

The weather on the plane is a lesson in extremes. If the sea isn't calm, it is battered by storms. On rare occasions, a tremor in the planar firmament sends a rogue wave sweeping across the plane, swamping entire islands and driving ships down to the reefs.

Any land that rises above the surface of the sea is hotly contested by the few air breathers that live on the plane. Fleets of rafts and ships lashed together serve as solid ground where nothing else is available, although most natives of the plane never break the surface of the sea and thus ignore these habitations.

The following are among the most important features of the Plane of Water:

\subsubsection{Citadel of Ten Thousand Pearls}
The nominal emperor of the marids dwells in the Citadel of Ten Thousand Pearls, an opulent palace made of coral and studded with pearls. The palace is the glittering centerpiece of the Sea of Light (see below). Visitors are welcome to ask favors of the emperor, whose mood is as changeable as the sea.

\subsubsection{Darkened Depths}
The deeper extents of the plane, where no sunlight reaches, are called the Darkened Depths. Horrid creatures dwell here, and the absolute cold and crushing pressure mean a swift end to creatures accustomed to the surface or the Sea of Light. Krakens and other mighty leviathans claim this realm.

\subsubsection{Isle of Dread}
One of the few islands on the plane is the Isle of Dread, which is connected to the Material Plane by means of a storm that regularly sweeps over the island. Ships from different worlds of the Material Plane end up wrecked on the rocks and reefs around the island, and settlements across the island are populated by the descendants of sailors who never found a way home. Theoretically, travelers who understand the workings of the storm could use it to travel to a desired Material Plane world.

\subsubsection{Sea of Ice}
Bordering the Para-elemental Plane of Ice is the Sea of Ice. The frigid water is choked with icebergs and sheet ice, which are inhabited by cold-loving creatures from the Plane of Ice. Drifting icebergs can carry these creatures farther into the Plane of Water to threaten ships and islands in warmer seas.

\subsubsection{Sea of Light}
Life flourishes in the sunlit waters of the Sea of Light, located in the upper reaches of the Sea of Worlds. Aquatic peoples craft castles and fortresses in the coral reefs here.

\subsubsection{Silt Flats}
The region of the Plane of Water nearest to the Para-elemental Plane of Ooze is called the Silt Flats. The water is thick with soil and sludge, turning into muddy ground before giving way to the great swamp that is the Para-elemental Plane.

\subsubsection{Elemental Water Adventures}
The nature of water is to flow, not like the gusting wind or the leaping flame, but smoothly and steadily. It is the rhythm of the tide, the nectar of life, the bitter tears of mourning, and the balm of sympathy and healing. Given time, it can erode all in its path.

Elemental Evil emphasizes the erosive power of water, as well as the destruction wrought by surging tides, deadly maelstroms, and raging torrents. Cultists of Evil Water believe the seas and deep waters are eager to reclaim the water trapped in the bodies of living creatures, and think it's their duty to return others to the primal waters by drowning them or shedding their blood.

\subsection{Elysium}
Elysium is home to creatures of unfettered kindness and a refuge for planar travelers seeking a safe haven. The plane's bucolic landscapes glimmer with life and beauty.

The River Oceanus originates in the lowest layer of Elysium, Thalasia, and flows through the plane's layers before cascading onward to the Beastlands. Though illustrations of the plane's layers seem to show the river flowing "up" from each layer to the one "above" it, the experience of passing from one layer to another on the river is no more dramatic than weathering rapids on any ordinary river. Along its course, the great river splits into myriad smaller flows, recombines, and splits again.

\begin{DndTable}[header=Layers of Elysium]{XX}

Layer & Description\\
Amoria & Small towns on the banks of the River Oceanus, surrounded by lush meadows, are among the most hospitable refuges on the Outer Planes.\\
Eronia & Steep hills, craggy mountains, and white granite valleys offer a rugged home for hardy souls.\\
Belierin & Lighthouses pierce the fog and form hubs for small communities amid sprawling marshlands.\\
Thalasia & The Heroic Isles, rising from the headwaters of the River Oceanus, hold the best departed souls.\\
\end{DndTable}

\subsubsection{Elysium Adventures}
Tranquility and contentment seep into the bones and souls of those who enter Elysium. The longer a visitor remains on the plane, the less reason they find to ever leave. An adventure in Elysium can challenge characters' devotion to doing good by offering them the opportunity (or the temptation) to rest from their labors and enjoy a well-earned reward.

Belierin is said to be the prison of some deadly creature. Some tales say it's a powerful titan, perhaps the \textbf{tarrasque}, while others claim it's a deposed duke of the Nine Hells, a banished elemental prince, or even a near-dead deity. Evil creatures sometimes lurk in the marshes, seeking to free the prisoner or claim some power from it.

Characters might also venture to Elysium to seek out some ancient spirit on the Heroic Isles. When faced with the once-a-millennium task of forestalling a prophesied disaster, characters might consult with the valorous knight who accomplished the deed a thousand years ago.

\subsection{Ethereal Plane}
The Ethereal Plane is a misty, fogbound dimension. Its "shores," called the Border Ethereal, overlap the Material Plane, the Feywild, the Shadowfell, and the Inner Planes, and every location on those planes has a corresponding location on the Ethereal Plane. Visibility in the Border Ethereal is usually limited to 60 feet. The plane's depths comprise a region of swirling mist and fog called the Deep Ethereal, where visibility is usually limited to 30 feet.

Characters can use the \textit{Etherealness} spell to enter the Border Ethereal. The \textit{Plane Shift} spell allows transport to the Border Ethereal or the Deep Ethereal, but unless the intended destination is a specific location or a teleportation circle, the point of arrival could be anywhere on the plane.

\subsubsection{Border Ethereal}
From the Border Ethereal, a traveler can see into whatever plane it overlaps, but that plane appears grayish and indistinct, its colors blurring into each other and its edges turning fuzzy, limiting visibility to 30 feet into the other plane. Conversely, the Ethereal Plane is usually imperceptible to those on the overlapped planes, except with the aid of magic.

Normally, creatures in the Border Ethereal can't attack creatures on the overlapped plane, and vice versa. A traveler on the Ethereal Plane is imperceptible to someone on the overlapped plane, and solid objects on the overlapped plane don't hamper the movement of a creature in the Border Ethereal. The exceptions are certain magical effects (including anything made of magical force) and living beings. This makes the Ethereal Plane ideal for scouting, spying on opponents, and moving around without being detected. The Ethereal Plane also disobeys the laws of gravity; a creature there can freely move in any direction.

\subsubsection{Deep Ethereal}
To reach the Deep Ethereal, one typically needs a \textit{Plane Shift} spell, a \textit{Gate} spell, or a magical portal. Visitors to the Deep Ethereal are engulfed by roiling mist. Scattered throughout the plane are curtains of vaporous color, and passing through a curtain leads a traveler to a region of the Border Ethereal connected to a specific Inner Plane, the Material Plane, the Feywild, or the Shadowfell. The color of the curtain indicates the plane whose Border Ethereal the curtain conceals; see the Ethereal Curtains table. The curtains are also distinguishable by texture and temperature, each one reflecting something of the nature of the plane beyond.

\begin{DndTable}[header=Ethereal Curtains]{cXX}

1d12 & Plane & Curtain Color\\
1–2 & Material Plane & Turquoise\\
3 & Shadowfell & Dusky gray\\
4 & Feywild & Opalescent\\
5 & Elemental Plane of Air & Pale blue\\
6 & Elemental Plane of Earth & Chestnut\\
7 & Elemental Plane of Fire & Orange\\
8 & Elemental Plane of Water & Green\\
9 & Para-elemental Plane of Ash & Dark gray\\
10 & Para-elemental Plane of Ice & Aquamarine\\
11 & Para-elemental Plane of Magma & Maroon\\
12 & Para-elemental Plane of Ooze & Chocolate\\
\end{DndTable}

Traveling through the Deep Ethereal is unlike physical travel. Distance is meaningless, so although travelers feel as if they can move by a simple act of will, it's impossible to measure speed and hard to track the passage of time. A trip through the Deep Ethereal takes 1d10 × 10 hours from one curtain to another, regardless of the origin and destination. In combat, creatures move at their normal speeds.

\subsubsection{Ether Cyclones}
An ether cyclone is a serpentine column that spins through the plane. The cyclone appears abruptly, distorting and uprooting everything in its path and carrying the debris for miles. Travelers with a Passive Perception score of 15 or higher receive 1 minute of warning: a deep thrum in the ethereal matter. Travelers who can't reach a curtain or portal leading elsewhere suffer the cyclone's effect. Roll 1d20 and consult the Ether Cyclone table to determine the effect on all creatures in the vicinity.

\begin{DndTable}[header=Ether Cyclone]{cX}

1d20 & Effect\\
1–12 & Extended journey. Each character in a group traveling together makes a DC 15 Charisma saving throw. If at least half the group succeeds, travel is delayed by 1d10 hours. Otherwise, the journey's travel time is doubled.\\
13–19 & Blown to a location in the Border Ethereal overlapping a random plane (roll on the \textit{Ethereal Curtains} table)\\
20 & Hurled to a random destination on the Astral Plane\\
\end{DndTable}

\subsubsection{Radiant Citadel}
Against the unending mist and unseen terrors of the Ethereal Plane, the Radiant Citadel stands bright as a bastion of hope. It's a living relic of the ingenuity and collaboration of twenty-seven great civilizations on the Material Plane. Abandoned and lost for ages, the Radiant Citadel was resurrected from its slumber and reclaimed by descendants of those societies.

The Radiant Citadel is a nexus of diplomacy and trade, a repository of histories and secrets, and a thriving sanctuary for those seeking safety or a better life. The floating city is a miracle of architecture carved out of a single, massive fossil that snakes around a colossal gemstone shard known as the Auroral Diamond. The luminescence of the Auroral Diamond is mirrored in the constellation of fifteen structure-sized gemstones, the Concord Jewels, that orbit the city and provide transportation to the far-flung homes of the city's founding civilizations. In the haze of the Ethereal Plane, the Auroral Diamond is a scintillating beacon visible from miles away. The diamond seems to have moods, changing colors unpredictably, but it is always visible for wanderers lost and in need.

Just beyond the city whirls a massive ether cyclone known as the Keening Gloom—a looming threat that's a grim reminder of the Radiant Citadel's precarious position.

Heroes and paupers meet on equal footing in the Radiant Citadel. By common agreement, power and resources are equitably shared. Dignity is afforded to all, and great need is met with great aid.

\subsubsection{Ethereal Plane Adventures}
Adventurers typically use the Ethereal Plane to travel from one place to another, either skirting around Material Plane obstacles on the Border Ethereal or venturing into the Deep Ethereal to travel to the Inner Planes.

The Radiant Citadel can serve as a home base for any campaign built around the idea of exploring new worlds. Several such worlds are introduced in Journeys through the Radiant Citadel, an anthology of short adventures.

\subsection{Far Realm}
The Far Realm is outside the known multiverse. In fact, it might be an entirely separate universe with its own physical and magical laws. Where stray energies from the Far Realm leak onto another plane, matter is warped into alien shapes that defy understandable geometry and biology. Aberrations such as mind flayers and beholders are either from this plane or shaped by its strange influence.

The entities that abide in the Far Realm are too alien for mortal minds to accept without strain. Titanic creatures swim through nothingness there, and unspeakable beings whisper awful truths to those who dare listen. For mortals, knowledge of the Far Realm is a struggle of the mind to overcome the boundaries of matter, space, and rational thought. Some Warlocks embrace this struggle by forming pacts with entities there. Anyone who has seen the Far Realm mutters about eyes, tentacles, and horror.

The Far Realm has no well-known portals, or at least none that are still viable. Ancient elves once opened a vast portal to the Far Realm within a mountain called Firestorm Peak, but their civilization imploded in bloody terror, and the portal's location—even its home world—is long forgotten. Lost portals might still exist, marked by an alien magic that transforms the surrounding area.

\subsubsection{Far Realm Adventures}
The Far Realm is the home of entities so far beyond comprehension that mortals can't fathom their motivations. To see these beings is to become lost in their magnitude and the evidence that mortals have never, will never, and could never matter to the cosmos at large.

Adventures involving travel to the Far Realm or its influence seeping into the Material Plane might touch on fundamental questions of what it means to be a person, what mental and bodily autonomy mean and their value, and whether mortals have any control over their fate or any importance in the grand scheme of things. (It's an especially good idea to review your players' limits that might pertain to such issues before planning an adventure exploring these themes, as discussed in the "\textit{Ensuring Fun for All}" section in \textit{chapter 1}.)

\subsection{Feywild}
The Feywild, also called the Plane of Faerie, is a land of soft lights and wonder, a place of music and magic. The plane responds to unfettered emotion: flowers turn and tremble in the presence of a heated argument, grass withers under the feet of one who seethes with malice, and birds chip merrily in the presence of those who are joyous and squawk angrily at those who are dour.

Time and distance in the Feywild are mutable, as is the plane's geography. Roads are uncommon, and the ones that exist change as frequently as the land around themselves. Feywild natives are accustomed to the plane's mutability, but it can be terribly disorienting to visitors.

The Feywild exists in parallel to the Material Plane as an alternate dimension that occupies the same cosmological space. When moving from the Material Plane to the Feywild, travelers usually find themselves in a location similar to the one they left, but more marvelous and magical—and often more vibrant and colorful, too. Adventurers climbing a volcano on the Material Plane might suddenly find themselves scaling a Feywild mountain topped with skyscraper-sized crystals that glow with internal fire. Leaving behind a wide and muddy river on the Material Plane, characters might appear beside a clear and winding brook whose waters glitter like diamonds in the Feywild. In the heart of a dismal marsh might lie a portal leading to a vast bog filled with eerie lights and sinister shapes twisting in the water. And moving to the Feywild from old ruins on the Material Plane might put a traveler at the door of an archfey's castle.

\subsubsection{Domains of Delight}
Much of the Feywild is governed by powerful Fey called archfey. The area under a particular archfey's command—called a Domain of Delight—reflects the character and desires of its ruler. Some domains are bright and cheery, bathed in perpetual sunlight and awash in colorful wildflowers, while others are gloomy and drab, cast in unending twilight. Most of them change with the emotional state of their rulers.

The following sections describe a handful of the best-known Domains of Delight.

\subsubsection{Fablerise}
The domain of a story-spinning spider archfey named Yarnspinner, Fablerise is a rambling thicket of twisted roots, thorny vines, and sinuous creepers. This vegetation weaves together to form long tunnels, grand hallways, and enormous domes. Yarnspinner loves reading stories to the animals that occupy his domain.

\subsubsection{Gloaming Court}
The Queen of Air and Darkness rules the Gloaming Court, a realm of twilight, fireflies, cobwebs, and autumn leaves accompanied by the music of hooting owls and croaking frogs. The Fey of the Gloaming Court shun the formalized etiquette and rituals of the Summer Court (see below), instead prizing behavior that is intuitive and instinctual.

\subsubsection{Prismeer}
Prismeer is a large domain belonging to the archfey Zybilna. It encompasses a vast swamp called Hither; an ancient forest named Thither; and a stormy, mountainous landscape called Yon. Zybilna resides in the Palace of Heart's Desire, situated where the three portions of her realm meet. As its name suggests, the palace is fabled as a destination for anyone seeking their heart's desire. On some worlds, Zybilna is regarded as a fairy godmother of sorts, granting wishes for the lost, the forsaken, or the betrayed. Sometimes her wishes bring happiness, other times despair. (Prismeer is detailed in The Wild Beyond the Witchlight.)

\subsubsection{Summer Court}
Ruled by the archfey Queen Titania, the Summer Court is the most settled and pastoral domain in the Feywild. Wrapped in the warmth of a perpetual summer day, with fluttering butterflies and a riot of colorful flowers, the lands of the Summer Court mimic the trappings of courtly life in some realms of the Material Plane. The residents of this court wear elegant clothing and value elaborate ceremony and ritualized etiquette, and the Fey are quick to shun those who flout the Summer Court's baroque rules.

\subsubsection{Fey Crossings}
Fey crossings are places of mystery and beauty on the Material Plane that have a near-perfect mirror in the Feywild, creating a portal where the two planes touch. A traveler passes through a fey crossing by entering a clearing, wading into a pool, passing into a circle of mushrooms, or crawling under the trunk of a tree. To the traveler, it seems like simply moving into the Feywild. To an observer, the traveler is there one moment and gone the next.

Like other portals between planes, most fey crossings open infrequently. A crossing might open only during a full moon, on the dawn of a particular day, or for someone carrying a certain type of item. A fey crossing can be closed permanently if the land on either side is dramatically altered—for example, if a castle is built over the clearing on the Material Plane.

\subsubsection{Feywild Magic}
Tales speak of children kidnapped by Fey creatures and spirited away to the Feywild, only to return to their parents years later without having aged a day and with no memories of their captors or the realm they came from. Likewise, adventurers who return from an excursion to the Feywild are often alarmed to discover upon their return that time flows differently on the Plane of Faerie and that the memories of their visit are hazy. You can use these optional rules to reflect the strange magic that suffuses the plane.

\subparagraph{Memory Loss}
A creature that leaves the Feywild makes a DC 10 Wisdom saving throw. Fey creatures automatically succeed on the saving throw, as do creatures that have the Fey Ancestry trait, such as elves. On a failed save, the creature remembers nothing from its time spent in the Feywild. On a successful save, the creature's memories remain intact but are a little hazy. Any spell that can end a curse can restore the creature's lost memories.

\subparagraph{Time Warp}
While time seems to pass normally in the Feywild, characters might spend a day there and realize, upon leaving the plane, that less or more time has elapsed everywhere else in the multiverse.

Whenever a creature or group of creatures leaves the Feywild after spending at least 1 day on that plane, you can choose a time change that works best for your campaign, if any, or roll on the Feywild Time Warp table. A \textit{Wish} spell can be used to remove the effect on up to ten creatures. Some powerful Fey have the ability to grant such wishes and might do so if the beneficiaries agree to subject themselves to a \textit{Geas} spell and complete a quest after the \textit{Wish} spell is cast.

\begin{DndTable}[header=Feywild Time Warp]{cX}

1d20 & Result\\
1–2 & Days become minutes\\
3–6 & Days become hours\\
7–13 & No change\\
14–17 & Days become weeks\\
18–19 & Days become months\\
20 & Days become years\\
\end{DndTable}

\subsubsection{Feywild Adventures}
The Feywild gives physical expression to powerful emotion and excels at metaphor. When characters venture into the Feywild, they might find themselves robbed of a cherished memory or deep regret, then later find the stolen memories embodied in little figurines or lockets. A mischievous sprite might sneak up behind a character who is laughing loudly and steal their laughter, robbing the character of the ability to laugh until the sprite is found and the laughter—perhaps taking physical form as a bouquet of lovely flowers—reclaimed.

\subsection{Gehenna}
A volcanic mountain dominates each of the four layers of Gehenna, and lesser volcanic earthbergs drift in the air and smash into the greater mountains. The rocky slopes of the plane make movement difficult and dangerous. The ground inclines at least 45 degrees almost everywhere. In places, steep cliffs and deep canyons present more challenging obstacles. Hazards include volcanic fissures that vent noxious fumes or searing flames.

Gehenna is the birthplace of yugoloths, greedy and selfish Fiends that dwell here in great numbers.

\begin{DndTable}[header=Layers of Gehenna]{XX}

Layer & Description\\
Khalas & Lava illuminates clouds of volcanic ash and steam from the River Styx.\\
Chamada & Constant lava flows and eruptions make overland travel difficult. Iron zeppelins piloted by yugoloths drift through the constant gray ashfall.\\
Mungoth & Acidic ash mingles with falling snow on this freezing layer.\\
Krangath & The Dead Furnace is a great mountain suspended in silence and darkness, home to a coterie of liches.\\
\end{DndTable}

\subsubsection{Gehenna Adventures}
Gehenna is the plane of suspicion and greed, with no space for mercy or compassion. Adventures on this plane might be an opportunity to explore themes of betrayal, examining how characters behave when tensions run high and they can trust no one—perhaps not even each other. (See "\textit{Environmental Effects}" in \textit{chapter 3} for one way this atmosphere can manifest.) Characters might encounter people in need who turn out to be yugoloths in disguise, pitting the characters' growing suspicion against their empathy and compassion.

Characters might make their way to the Teardrop Palace on Khalas to purchase something they can't find elsewhere, probably at a terrifying cost. This bustling market, crowded with Fiends and occasional mortal visitors, offers all manner of forbidden and sinister goods for sale. Its name comes from its shape: the point is on the palace's uphill side so it diverts the ever-present lava flow to either side of the structure.

Or characters could try to infiltrate the Tower Arcane in search of some great secret from yugoloths' ancient history. The tower, a sinister structure adorned with blades and spikes and guarded by arcanaloths, stands somewhere on Chamada. It is rumored to hold yugoloths' history and the records of their contracts.

\subsection{Hades}
The layers of Hades are called the Three Glooms—places without joy, hope, or passion. A gray land with an ashen sky, Hades is the destination of many souls that are unclaimed by gods or Fiends. These souls become larvae and spend eternity in this place, which lacks a sun, a moon, stars, or seasons. Leaching away color and emotion, the gloom on this plane is more than most visitors can stand.

\begin{DndTable}[header=Layers of Hades]{XX}

Layer & Description\\
Oinos & A land of dead-gray ash, stunted trees, and virulent disease is stalked by roving bands of Fiends looking for a fight or recruits for the Blood War.\\
Niflheim & Gray pine trees blanket rolling hills and rocky bluffs, and thick mists coil around their trunks.\\
Pluton & Shriveled willows, olive trees, and poplars contribute to the gloom of this concentration of the deepest despair in the multiverse.\\
\end{DndTable}

\subsubsection{Plane of Gloom}
At the end of each Long Rest taken on the plane, a visitor makes a DC 10 Wisdom saving throw. On a failed save, the creature gains 1 Exhaustion level that can't be removed while the creature is in Hades. If the creature reaches 6 Exhaustion levels, it doesn't die. Instead, it permanently transforms into a \textbf{Larva}, whereupon all Exhaustion levels afflicting the creature are removed.

\subsubsection{Hades Adventures}
Hades embodies despair manifested as apathy. Pure, undiluted evil is like an inescapable force of gravity, dragging all creatures down—not in body, but in spirit. Even the consuming rage of the Abyss and the devious plotting of the Nine Hells are subjugated to hopelessness in the Gray Wastes of Hades. The plane slowly kills dreams and desires, draining hope and optimism from formerly fiery spirits.

An adventure in Hades can challenge characters to find an answer to the ever-present question that hangs over this plane: why bother? As apathetic despair saturates their hearts and spirits, they must find a way to rekindle the passion of life and the sense of purpose that drives them or else succumb to the hopelessness of the plane.

Adventurers might pursue a hag, a lich, or another evil spellcaster who comes to Hades to collect larvae for vile purposes. Once they are in the Three Glooms, the adventurers risk becoming trapped by the overwhelming despair of the place.

\subsection{Limbo}
Limbo is a plane of pure chaos, a roiling soup of impermanent matter and energy. Stone melts into water that freezes into metal, then turns into diamond that burns up into smoke that becomes snow, and on and on in an endless, unpredictable process of change. Fragments of more ordinary landscapes—bits of forest, meadow, ruined castles, and even burbling streams—drift through the disorder. The whole plane is a nightmarish riot.

Limbo has no gravity, so creatures visiting the plane float in place. A creature can move up to its Speed in any direction by merely thinking of the desired direction of travel.

Limbo has no layers—or if it does, the layers continually merge and part, each is as chaotic as the next, and distinguishing one from another is impossible.

\subsubsection{Power of the Mind}
Limbo conforms to the will of the creatures inhabiting it. Creative imaginations can create whole islands of their own invention within the plane, sometimes maintaining those places for years. A nonsapient creature such as a fish, though, might have less than a minute before the pocket of water surrounding it freezes, vanishes, or turns to glass. Slaadi live here and swim amid this chaos, creating nothing, whereas githzerai build entire monasteries with their minds.

As a Magic action, a creature in Limbo can make an Intelligence check to mentally move an object within 30 feet of itself that is on the plane and isn't being worn or carried. The DC depends on the object's size: DC 5 for Tiny, DC 10 for Small, DC 15 for Medium, DC 20 for Large, and DC 25 for Huge or larger. On a successful check, the creature moves the object 5 feet plus a number of a feet equal to how much the total exceeded the DC.

A creature can also take a Magic action to make an Intelligence check to alter a nonmagical object within 30 feet of itself that isn't being worn or carried. The DC is based on the object's size: DC 10 for Tiny, DC 15 for Small, DC 20 for Medium, and DC 25 for Large or larger. On a successful check, the creature changes the object into another nonliving form of the same size, such as turning a boulder into a ball of fire.

Finally, a creature in Limbo can take a Magic action to make a DC 20 Intelligence check to stabilize an area within a 30-foot-radius Sphere [Area of Effect] centered on a point it can see on the plane. On a successful check, the creature prevents the area from being altered by the plane for 24 hours or until the creature takes this Magic action again.

\subsubsection{Limbo Adventures}
Limbo is change. That constant churn is most easily discerned in the ever-shifting physical form of elements altering and reconfiguring in the vast expanse of the plane, but it applies just as much on a mental and emotional level. Visitors to the plane find themselves caught up in a storm of intrusive thoughts and unruly emotions, forcing them to confront the transient nature of so much of what they think of as their identity. The key to success on this plane—both in shaping the physical environment and in mastering the internal landscape of chaos—is asserting one's sense of self, identifying what is unchanging amid the storm of constant change.

The sanctuaries of the githzerai are among the few havens that adventurers can hope to find on this tumultuous plane. Although githzerai aren't generally hostile to visitors who come in peace, they don't welcome those who bring the chaos of Limbo with them: a tumultuous heart brought into a refuge can unravel the entire sanctuary.

Adventurers might also come to Limbo to explore the secrets of the Spawning Stone. Said to have been created by Primus, the overlord of the modrons, the Spawning Stone absorbs chaotic energy and makes it possible to shape enclaves of order in Limbo, but the chaotic energy it absorbs is responsible for the creation of slaadi.

\subsection{Material Plane}
Worlds of the Material Plane are infinitely diverse, but it was not always so. Some legends speak of a primordial state, a single reality called the \textit{First World}, where many of the peoples and monsters that inhabit the worlds on the Material Plane originated. After the First World was shattered by a great cataclysm, the many worlds were formed like reflections or (in some cases) distortions of that original reality.

Some myths describe a great tree that grew on the First World at the dawn of time. Planted and tended by the god Corellon, this tree was a seedling of Yggdrasil, the World Tree that connects all the Outer Planes (see "\textit{Traveling the Outer Planes}" earlier in this chapter). When the First World was destroyed, seeds from this great tree scattered into the void of the Material Plane. Legends say that these seeds sprouted and formed worlds of their own—all the myriad worlds that now constitute the Material Plane.

The most widely known worlds are the ones that have been published as official campaign settings for the D\&D game over the years, many of which are shown on the \textit{D\&D Settings table} in \textit{chapter 5}. If your campaign takes place in one of these settings, your version of it can diverge wildly from what's in print.

\subsubsection{Traveling between Worlds}
Transit between the worlds of the Material Plane is rare but not impossible and can be accomplished in a variety of ways.

\subsubsection{The Dream of Other Worlds}
Aided by magic, travelers can fall into a deep slumber and dream themselves into a new realm.

\subsubsection{The Great Journey}
Characters can undertake an epic voyage fraught with peril and obstacles to be overcome. One route leads through Wildspace and across the Astral Plane aboard a vessel powered by magic. (The "\textit{Astral Plane}" section in this chapter describes Wildspace.) It is also possible to travel through the Shadowfell or the Feywild, though such routes are less charted and no less perilous.

\subsubsection{The Leap to Another Realm}
The most direct method involves the use of spells such as \textit{Teleportation Circle} or \textit{Teleport}, or magical portals like those described in this chapter. This magic causes the user to appear in a known teleportation circle or some other location in another world.

\subsubsection{The Roots of the Worlds}
Similar to magical portals, nexus points are locations that exist in multiple worlds at the same time. These points might be located at or near the roots of the worlds—the places where the seedlings of the First World's great tree took root and grew into a new world.

A nexus point can be a geographical feature, such as an enormous tree, a mountain or mesa, a yawning cavern deep under the mountains, or a meteorite in an enormous crater. It might also be a constructed feature: a lonely tower or castle, a bustling tavern, or even a city. Normally, visitors to these places return to the same world they came from when they depart, but it's also possible to use a nexus point to travel from one world to another. Depending on the place, shifting worlds might require the use of magic, an object from the desired destination world as a sort of key, or nothing more than an act of will.

Some nexus points exist in multiple worlds—but not at the same time. They flit from world to world, disappearing from one and appearing in another according to a regular schedule. Such a place might linger on one world for anywhere from a year to an hour before moving on to another, carrying everyone inside with it to a new world.

\subsection{Mechanus}
Mechanus is where perfectly regimented order reigns supreme. It consists of equal measures of light and dark, and equal proportions of heat and cold. On Mechanus, law is reflected in a realm of gigantic clockwork gears, interlocked and turning according to their measure. The cogs seem to be engaged in a calculation so vast that no deity can fathom its purpose. Some theories hold that they are the clockwork of time throughout the cosmos—that time itself would stop if the gears ceased their turning. Other theories propose that the cogs uphold the basic rules and order of the cosmos.

Modrons are the primary inhabitants of Mechanus and maintain its intricate clockworks. The plane is also home to the creator of the modrons: a godlike being called Primus, whose realm is called Regulus.

Mechanus has no distinct layers. Each turning cog has its own force of gravity pulling toward its center, with structures built on the faces of the cogs. Some of the cogs are like small islands, while others are hundreds of miles across.

\subsubsection{Mechanus Adventures}
Mechanus embodies absolute order, and it influences those who spend time here. Individual consciousness is subordinated to the search for perfect order, and "I" is ultimately subsumed into "we."

An adventure on Mechanus might lead characters to examine their individual egos in the light of the adventuring party. It might challenge a character to set aside personal goals for the benefit of the group (or the greater cause of cosmic law), or alternatively it might encourage characters to assert their own individual identities, distinct from the party and possessing their own goals and needs.

\subsection{Mount Celestia}
The Seven Heavens of Mount Celestia rise like a mountain from a shining Silver Sea to utterly incomprehensible heights, with seven plateaus marking its seven heavenly layers. The plane is the model of justice and order, of celestial grace and endless mercy, where angels and champions of good guard against incursions of evil. It is one of the few places on the planes where travelers can let down their guard. Its inhabitants strive constantly to be as righteous as possible. Countless creatures aim to reach the highest and most sublime peak of the mountain, but only the purest souls can. That peak fills even the most jaded of travelers with awe.

The pervasive goodness of Mount Celestia \textit{bestows blessings} on creatures on the plane (see "\textit{Environmental Effects}" in \textit{chapter 3}).

\begin{DndTable}[header=Layers of Mount Celestia]{XX}

Layer & Description\\
Lunia & In the Silver Heaven, the holy water of the Silver Sea laps at the base of the celestial mountain under a starry sky.\\
Mercuria & The Golden Heaven's tame slopes and lush valleys are bathed in golden light that evokes the hope of a new dawn.\\
Venya & In the Pearly Heaven, terraced fields and tended woodlands dot the snowy slopes.\\
Solania & In the Crystal Heaven, holy shrines glitter under a silvery sky amid luminescent fog.\\
Mertion & On the sweeping plains of the Platinum Heaven, holy soldiers muster in grand citadels for battles across the planes.\\
Jovar & The Glittering Heaven, strewn with beautiful rubies and garnets, is home to the seven-tiered Heavenly City.\\
Chronias & The Illuminated Heaven is an ineffable mystery.\\
\end{DndTable}

\subsubsection{Mount Celestia Adventures}
The plane of ultimate law and good is sometimes imagined to be the most boring place in the multiverse, but in truth Mount Celestia's nature makes it the target of unrelenting attacks by evil forces. The devils of the Nine Hells, in particular, long to corrupt the goodness of the Seven Heavens. The yugoloths of the Lower Planes covet the wealth of the plane, particularly the mines of Solania and the scattered gems of Jovar. And the demons of the Abyss would like nothing better to smear their filth on the gleaming purity of Mount Celestia.

But while Fiends of all sorts launch doomed assaults on the shores of Lunia, evil's true foothold on the plane is in the hearts of those well-meaning visitors who bring their secret shame and hidden sins to the holy mountain. An adventure in Mount Celestia is an opportunity for characters to prove themselves worthy of the many blessings it offers—or to become worthy by forswearing the selfishness, greed, and hatred that lurk in their hearts.

On the edge of a clear lake in Mertion stands the city of Empyrea, renowned for the healing power of its fountains and springs. Pilgrims from across the planes seek out the healers, hospitals, and restorative magic found here.

\subsection{Negative Plane}
Cupped like a bowl beneath the other planes, the Negative Plane is the source of necrotic energy that destroys the living and animates the Undead. A lightless void without end, it is a needy, greedy plane, sucking the life out of anything that is vulnerable. Heat, fire, and life itself are all drawn into the maw of this plane, which always hungers for more.

To an observer, there's little to see on the Negative Plane. It is a dark, empty place, an eternal pit where a traveler can fall until the plane steals away all light and life. Merely entering the plane is comparable to the life-draining touch of a wraith, so only creatures that have Immunity to Necrotic damage can survive there for long.

In some locations on the Negative Plane, the intensity of the plane is so great that the negative energy folds in on itself, stabilizing into solid chunks of matter that devour light. These chunks, called voidstones, are thought to be the source of Sphere of Annihilation and similar magical effects. Anything that comes into contact with a voidstone is destroyed in seconds.

\subsubsection{Negative Plane Adventures}
An adventure on the Negative Plane is a face-to-face confrontation with annihilation, which is unlikely to end well for the adventurers. Even if their magic enables them to survive the environment of the plane, the experience tends to drain all the vitality, energy, and happiness from body and soul. The Negative Plane has all the apathy and despair of Hades and the Shadowfell, combined and concentrated in an infinite expanse of nonbeing and uncreation.

\subsection{Nine Hells}
The Nine Hells inflames the imaginations of travelers, the greed of treasure seekers, and the battle fury of all moral creatures. It is the ultimate plane of law and evil, and the epitome of premeditated cruelty. The devils of the Nine Hells are bound to obey the laws of their superiors, but they squabble within their individual castes. Most undertake any plot, no matter how foul, to advance themselves. At the very top of the hierarchy is Asmodeus, who has yet to be bested. If he were vanquished, the victor would rule the plane in turn. Such is the law of the Nine Hells.

\subsubsection{The Nine Layers}
The Nine Hells has nine layers. The first eight are ruled by archdevils who answer to Asmodeus, the archduke of Nessus, the ninth layer. Collectively, the rulers of the Hells are called the Lords of the Nine. To reach Nessus, one must descend through all eight layers above it in order. The most expeditious means of doing so is the River Styx, which plunges ever deeper as it flows from one layer to the next. Only the most courageous adventurers can withstand the torment and horror of that journey.

The Layers of the Nine Hells table summarizes each layer; detailed descriptions of these layers follow the table.

\begin{DndTable}[header=Layers of the Nine Hells]{XX}

Layer & Description\\
\textit{Avernus} & The Blood War rages across battlefields littered with corpses and the wreckage of hellish war machines.\\
\textit{Dis} & Iron roads in deep canyons lead to the dreaded Iron City of Dis.\\
\textit{Minauros} & Acid falls like rain on putrid bogs and decaying cities.\\
\textit{Phlegethos} & Obsidian fortresses bask in the heat of raging volcanoes and magma rivers.\\
\textit{Stygia} & Levistus's prison is a frigid hellscape of jagged ice and cold fire.\\
\textit{Malbolge} & An ever-crumbling mountain threatens to bury visitors.\\
\textit{Maladomini} & Swarms of hungry flies plague dead cities surrounded by utter desolation.\\
\textit{Cania} & Ice-trapped cities provide shelter in a realm cold enough to freeze the soul.\\
\textit{Nessus} & Mighty fortresses stand watch over the deepest pits of the Nine Hells.\\
\end{DndTable}

\subsubsection{Avernus}
By Asmodeus's orders, no planar portals connect directly to the lower layers of the Nine Hells. The first layer, Avernus, is the arrival point for visitors, a rocky wasteland with rivers of blood and clouds of biting flies. Fiery comets occasionally fall from the darkened sky and carve out fuming impact craters. Empty battlefields are littered with weapons and bones, showing where the legions of the Nine Hells prevailed against invading enemies.

The archdevil Zariel rules Avernus, having supplanted her rival, Bel, who fell out of Asmodeus's favor and was forced to serve as Zariel's adviser. Tiamat, the Queen of Evil Dragons, is a prisoner on this layer, ruling her own domain but confined to the Nine Hells by Asmodeus in accordance with some ancient contract (the terms of which are known only to Tiamat and the Lords of the Nine).

Zariel appears as an angel whose skin and wings are scorched. Her eyes burn with a furious white light that can cause creatures looking upon her to burst into flame. Her seat of power is a flying basalt citadel that rakes the battlefields of Avernus.

\subsubsection{Dis}
Dis, the second layer of the Nine Hells, is a labyrinth of canyons wedged between sheer mountains rich with iron ore. Iron roads span and wend through the canyons, watched over by the garrisons of iron fortresses perched atop jagged pinnacles.

The second layer takes its name from its current lord, Dispater. A manipulator and deceiver, the archduke is devilishly handsome, bearing only small horns, a tail, and a cloven left hoof to distinguish him from a human. His crimson throne stands in the heart of the Iron City of Dis, a hideous metropolis. Planar travelers come here to conspire with devils and to close deals with night hags, rakshasas, incubi, succubi, and other Fiends. Contracts signed on his layer contain special provisions that allow Dispater to collect a cut of the deal.

Dispater is one of Asmodeus's most loyal and resourceful vassals, and few beings in the multiverse can outwit him. He is more obsessed than most devils with striking deals with mortals in exchange for their souls, and his emissaries work tirelessly to foster evil schemes on the Material Plane.

\subsubsection{Minauros}
The third layer of the Nine Hells is a stench-ridden bog. Acidic rain spills from the layer's brown skies, thick layers of scum cover its putrid surface, and yawning pits lie in wait beneath the murk to engulf careless wanderers.

Cyclopean cities of ornately carved stone rise up from the bog, including the great city of Minauros, for which the layer is named. The slimy walls of the city rise hundreds of feet, protecting the flooded halls that are the lair of Mammon, the archduke of Minauros. Mammon resembles a massive serpent with the upper torso and head of a hairless, horned humanoid. Mammon's greed is legendary, and he is one of the few archdevils who will trade favors for gold instead of souls. His lair is piled high with treasures left behind by those who tried—and failed—to best him in a deal.

\subsubsection{Phlegethos}
Phlegethos, the fourth layer, is a fiery landscape whose seas of molten magma brew hurricanes of hot wind, choking smoke, and volcanic ash. Within the fire-filled caldera of Phlegethos's largest volcano rises Abriymoch, a fortress city made of obsidian and dark glass. With rivers of molten lava pouring down its outer walls, the city resembles the sculpted centerpiece of a gigantic, hellish fountain.

Abriymoch is the seat of power for the two archdevils who rule Phlegethos in tandem: Belial and Fierna, Belial's daughter. Both are handsome devils who resemble tieflings, with red skin and small horns. Belial exudes civility, even as his words carry an undercurrent of threat. His daughter is said to have the wickedest heart in the Nine Hells. The alliance of Belial and Fierna is unbreakable, for both are aware that their mutual survival hinges on it.

\subsubsection{Stygia}
The fifth layer of the Nine Hells is a freezing realm of ice within which cold flames burn. A frozen sea surrounds the layer, and its gloomy sky crackles with lightning.

Archduke Levistus once betrayed Asmodeus and is now encased deep in the ice of Stygia as punishment. He rules this layer all the same, communicating telepathically with his followers and servants, both in the Nine Hells and on the Material Plane.

Stygia is also home to its previous ruler, the serpentine archdevil Geryon, who was dismissed by Asmodeus to allow the imprisoned Levistus to regain his rule. Geryon's fall from grace has spurred much debate within the infernal courts. No one is certain whether Asmodeus had some secret cause to dismiss the archdevil or whether he is testing Geryon's allegiance for some greater purpose.

\subsubsection{Malbolge}
Malbolge, the sixth layer, has outlasted many rulers, among them Malagard the Hag Countess and the archdevil Moloch. Malagard fell out of favor and was struck down by Asmodeus in a fit of pique, while her predecessor, Moloch, still lingers somewhere on the sixth layer as an imp, plotting to regain Asmodeus's favor. Malbolge is a seemingly endless slope, like the sides of an impossibly huge mountain. Parts of the layer break off from time to time, creating deadly, booming avalanches. The inhabitants of Malbolge live in crumbling fortresses and great caves carved into the mountainside.

Malbolge's current ruler is Asmodeus's daughter, Glasya. Her cruelty and love of wicked schemes rival those of her father. The citadel that serves as her domicile on the slopes of Malbolge, called Osseia, is supported by cracked pillars and buttresses that are sturdy yet seem on the verge of collapse. Beneath the palace is a labyrinth lined with cells and torture chambers, where Glasya confines and torments those who displease her.

\subsubsection{Maladomini}
The seventh layer, Maladomini, is ruin-covered wasteland. Dead cities form a desolate urban landscape, and between them are empty quarries, crumbling roads, slag heaps, the hollow shells of empty fortresses, and swarms of hungry flies.

The archduke of Maladomini is Baalzebul, the Lord of Flies. He is a tall, powerful devil with the compound eyes of a fly. The archduke has long conspired to usurp Asmodeus, yet has failed at every turn. Asmodeus laid a curse on him that causes any deal made with him to lead to calamity. Asmodeus occasionally shows Baalzebul favor for reasons no other archduke can fathom, though some suspect that Asmodeus still respects the worthiness of this adversary.

\subsubsection{Cania}
Cania, the eighth layer of the Nine Hells, is a frozen hellscape whose ice storms can tear flesh from bone. Cities embedded in the ice provide shelter for guests and prisoners of Cania's ruler, the brilliant and conniving archdevil Mephistopheles.

Mephistopheles dwells in the ice citadel of Mephistar, where he plots to seize the throne of the Nine Hells and conquer all the planes. He is Asmodeus's greatest enemy and ally, and the archduke of Nessus appears to trust Mephistopheles's counsel. Mephistopheles knows he can't depose Asmodeus until his adversary makes a fatal miscalculation, and so both wait to discover what circumstances might turn them against each other. Mephistopheles is also a godfather of sorts to Glasya, further complicating the relationship between Mephistopheles and Asmodeus.

Mephistopheles is a tall, striking devil with impressive horns and a cool demeanor. He trades in souls, as do other archdevils, but he rarely gives his time to any creatures not worthy of his personal attention. It is said that only Asmodeus has ever deceived or thwarted him.

\subsubsection{Nessus}
The lowest layer of the Nine Hells, Nessus is a realm of dark pits whose walls are set with bleak fortresses. There, pit fiend generals loyal to Asmodeus garrison their diabolical legions and plot the conquest of the multiverse. At the center of the layer stands a vast rift of unknown depth, out of which rises the great citadel-spire of Malsheem, home to Asmodeus and his infernal court.

Malsheem resembles a gigantic hollowed-out stalagmite. The citadel is also a prison for souls that Asmodeus has locked away for safekeeping. Convincing him to release even one of those souls comes at a steep price, and it is rumored that Asmodeus has claimed whole kingdoms in the past in exchange for such favors.

Asmodeus most often appears as a handsome, bearded man with four large horns, piercing red eyes, and flowing robes. He can also assume other forms and is seldom seen without his ruby-tipped scepter in hand. Asmodeus is the most cunning and well-mannered of archdevils. On the surface, he seems warm, pleasant, and lighthearted, doling out wisdom and small acts of kindness like a caring father. The ultimate evil he represents can be seen only when he wills it so, or if he forgets himself and flies into a rage.

\subsubsection{Infernal Hierarchy}
The Nine Hells has a rigid hierarchy that defines every aspect of its society. Asmodeus is the supreme ruler of all devils and wields the power of a lesser god. Worshiped as such on the Material Plane, Asmodeus inspires evil cults and sinister Warlocks. In the Nine Hells, he commands scores of pit fiend generals, which in turn command legions of subordinates.

A supreme tyrant, a brilliant deceiver, and a master of subtlety, Asmodeus protects his throne by keeping his friends close and his enemies closer. He delegates most matters of rulership to the pit fiends and lesser archdevils that make up the infernal bureaucracy of the Nine Hells, even as he knows that those powerful devils conspire to usurp his throne. Asmodeus appoints archdevils, and he can strip any member of the infernal hierarchy of rank and status as he likes.

\subparagraph{Archdevils}
The archdevils include all the current and deposed rulers of the Nine Hells, as well as the fiendish aristocrats that make up their courts, attend them as advisers, and hope to supplant them.

\subparagraph{Promotion and Demotion}
When the soul of an evil mortal sinks into the Nine Hells, it takes on the physical form of a wretched lemure. Archdevils and greater devils can promote lemures to lesser devils. Archdevils can promote lesser devils to greater devils, and Asmodeus alone can promote a greater devil to archdevil status. All diabolic promotions involve a brief, painful transformation, with the devil's memories passing intact from one form to the next.

Low-level promotions are typically based on need, such as when a pit fiend transforms lemures into imps to gain stealthy spies under its command. High-level promotions are almost always based on merit, such as when a bone devil that distinguishes itself in battle is transformed into a horned devil by the archdevil it serves. A devil is seldom promoted more than one step at a time.

Demotion is the customary punishment for failure or disobedience among the devils. Archdevils or greater devils can demote a lesser devil to a lemure, which loses all memory of its prior existence. An archdevil can demote a greater devil to lesser devil status, but the demoted devil retains its memories—and might seek vengeance if the severity of the demotion is excessive.

No devil can promote or demote another devil that hasn't sworn fealty to it, preventing rival archdevils from demoting each other's most powerful servants. Since all devils swear fealty to Asmodeus, he can freely demote any other devil, transforming it into whatever infernal form he desires.

\subsubsection{Nine Hells Adventures}
The Nine Hells embodies the cruelty and corruption of law turned to evil ends. The devils of the Nine Hells are more cunning, more insidious, and far more dangerous than other Fiends. Their intelligence, their delight in deceit and manipulation, and their unhindered pursuit of their own agendas make them truly terrifying foes.

A descent into the Nine Hells is a journey into the heart of evil. Every shred of evil is used in the Nine Hells, and each layer specializes in some way to accommodate and exploit the vices and weaknesses of mortals. Far too many people who make such a journey discover their own hearts aren't immune to temptation and corruption, and they end up making the Nine Hells their eternal home. To avoid such a fate, good-hearted adventurers must resist the insidious manipulation, deceit, and treachery of devils, even when the devils promise to fulfill their deepest longings.

\subsection{Outlands}
The Outlands lies between the Outer Planes. It is the plane of neutrality, keeping all aspects of the planes in a paradoxical balance—simultaneously concordant and in opposition. The plane has varied terrain, with prairies, mountains, and shallow rivers.

The Outlands is a great disk. In fact, those who envision the Outer Planes as a wheel point to the Outlands as proof, calling it a microcosm of the planes. That argument might be circular, since the arrangement of the Outlands inspired the idea of the Great Wheel in the first place.

Evenly spaced around the outside edge of the circle are the gate-towns: sixteen settlements, each built around a portal leading to one of the Outer Planes. The Gate-Towns of the Outlands table lists all sixteen gate-towns and the Outer Planes they connect to. Each gate-town shares many of the characteristics of the plane where its gate leads. Planar emissaries often meet in these gate-towns, so it isn't unusual to see strange interactions, such as a Celestial and a Fiend arguing in a tavern while sharing a fine bottle of wine.

\begin{DndTable}[header=Gate-Towns of the Outlands]{XX}

Town & Gate Destination\\
Automata & Mechanus\\
Bedlam & Pandemonium\\
Curst & Carceri\\
Ecstasy & Elysium\\
Excelsior & Mount Celestia\\
Faunel & Beastlands\\
Fortitude & Arcadia\\
Glorium & Ysgard\\
Hopeless & Hades\\
Plague-Mort & Abyss\\
Ribcage & Nine Hells\\
Rigus & Acheron\\
Sylvania & Arborea\\
Torch & Gehenna\\
Tradegate & Bytopia\\
Xaos & Limbo\\
\end{DndTable}

\subsubsection{Outlands Adventures}
The Outlands is the closest the Outer Planes come to being like a world on the Material Plane. Adventurers can travel easily from one gate-town to the next, making a tremendous variety of planar-themed adventures possible within the boundaries of the Outlands.

Adventures in the Outlands often involve the conflicts between opposing planar influences. It's much easier for a slaad to wreak havoc in the gate-town of Automata than it is for even a horde of slaadi to accomplish anything in Mechanus itself. Celestial spies and Fiend assassins carry out subtle plots and deadly sabotage across the Outlands.

Despite these conflicts, the Outlands remains a plane of balance. Toward the center of the plane, away from the gate-towns, lie vast stretches of land similar to the different environments found on worlds of the Material Plane. Preserving nature's balance from the pull of powerful extremes in any direction can also be a theme of adventures here.

\textit{Planescape: Adventures in the Multiverse} includes extensive information on \textit{the Outlands}.

\subsection{Pandemonium}
Pandemonium is a plane of overwhelming chaos, a great mass of rock riddled with tunnels carved by howling winds. It is cold, noisy, and dark, with no natural light. Wind quickly extinguishes nonmagical open flames such as torches and campfires. It also makes conversation possible only by yelling, and even then only to a maximum distance of 10 feet. See "\textit{Environmental Effects}" in \textit{chapter 3} for more information about the \textit{winds of Pandemonium}.

Most of the plane's inhabitants are creatures that were banished to the plane with no hope of escape. The incessant winds force them to take shelter in places where the howls of the winds sound like distant cries of torment.

\begin{DndTable}[header=Layers of Pandemonium]{XX}

Layer & Description\\
Pandesmos & Howling winds, dark streams bound for the River Styx, and blowing snow pour through vast, desolate caverns.\\
Cocytus & Winds blowing through narrower tunnels create a stronger force and louder wails, making this the so-called "Layer of Lamentation."\\
Phlegethon & Tunnel walls absorb light while water creates intricate rock formations.\\
Agathion & Sealed-off tunnels are largely inaccessible from elsewhere, making them ideal as vaults for ancient secrets.\\
\end{DndTable}

\subsubsection{Pandemonium Adventures}
Pandemonium is the plane of last straws. The incessant howling of its winds brings everyone on the plane, sooner or later, to the edge of lashing out in frustration, breaking down in despair, or dissolving into incoherence—and then some event, force, or creature on the plane pushes them over that edge. Simply existing on the plane is exhausting; trying to accomplish even a basic conversation is aggravating.

An adventure in Pandemonium can be a way to explore what happens to characters on their worst day, when everything goes wrong and the howling wind won't let up. The trick is to convey the frustration that characters are bound to experience there without transferring that frustration to the players.

A jagged spike somewhere in Cocytus, called Howler's Crag, is rumored to have a unique magical property: anything yelled from the top of the crag is said to find the ears of its intended recipient—carried on a shrieking, frigid wind—no matter where in the multiverse that person might be.

\subsection{Para-elemental Planes}
The regions where the Elemental Planes collide and their elemental substances overlap are called Paraelemental Planes.

\subsubsection{Plane of Ash}
On the Plane of Ash, also called the Great Conflagration, howling winds from the Plane of Air mix with the cinder storms and lava of the Sea of Fire. This plane is an endless storm of flames, smoke, and ash. The thick ash obscures sight beyond a few dozen feet, and the battering winds make travel difficult. Here and there, ash clusters into floating realms where outlaws and fugitives take shelter.

\subsubsection{Plane of Ice}
The Plane of Ice, also called the Frostfell, forms the border between the Plane of Air and the Plane of Water. This plane is a seemingly endless glacier swept by constant, raging blizzards. Frozen caverns twist through the Plane of Ice, home to yetis, remorhazes, white dragons, and other creatures of cold. The inhabitants of the plane engage in a never-ending battle to prove their strength and ensure their survival.

The Frostfell's monsters and bitter cold make it a dangerous place to travel. Most planar voyagers keep to the air, braving the powerful winds and driving snow to avoid setting foot on the great glacier.

\subsubsection{Plane of Magma}
The boundary between the Plane of Earth and the Plane of Fire is a great range of volcanic mountains. The Plane of Magma, also called the Fountains of Creation, is home to azers, fire giants, and red dragons, as well as creatures from the neighboring planes. Lava flows down the slopes of these mountains toward the Plane of Fire.

\subsubsection{Plane of Ooze}
The border region between the Plane of Water and the Plane of Earth is a horrid swamp where gnarled trees and thick, stinging vines grow from the dense muck and slime. Here and there on the Plane of Ooze (also called the Swamp of Oblivion), stagnant lakes and pools play host to thickets of weeds and monstrous swarms of mosquitoes. The few settlements here consist of wooden structures suspended above the muck on platforms between trees. Visitors to the plane have sometimes tried elevating houses on poles stuck in the mud, but since no solid earth underlies the muck, even such structures eventually sink.

It is said that any object cast into the Swamp of Oblivion can't be found again for at least a century. Now and then, a desperate soul casts an Artifact of power into this place, keeping it away from the rest of the multiverse for a time. The promise of powerful magic lures adventurers to brave the monstrous insects and hags of the swamp.

\subsubsection{Para-elemental Plane Adventures}
The Para-elemental Planes are extreme environments but fundamentally similar to places found on the Material Plane—the place where all four elements mingle freely.

At a symbolic level, the Para-elemental Planes represent the interaction and sometimes the contrast between the forces and ideals embodied by their constituent elements. The Plane of Ash, for example, highlights the commonality between air and fire—the tendency to movement and change, given a destructive tone by the raging conflagration of the plane. The Plane of Ooze heightens the contrast between stable, rigid earth and steadily flowing water.

\subsection{Positive Plane}
Like a dome above the other planes, the Positive Plane is the source of radiant energy and the raw life force that suffuses all living beings. Like the heart of a star, it is a continual furnace of creation, a domain of brilliance beyond the ability of mortal eyes to comprehend. It is a vibrant plane, so alive that travelers are empowered by visiting it.

The Positive Plane has no surface and is akin to the Elemental Plane of Air with its wide-open nature. However, every bit of this plane glows brightly with innate power. This power is dangerous to mortal forms, which can't handle it for long. Only creatures that have Immunity to Radiant damage can survive there.

\subsubsection{Positive Plane Adventures}
Vibrant life, creative energy, and radiant health are the essential characteristics of the Positive Plane, though touching this fundamental energy can be just as dangerous as entering the soul-siphoning annihilation of the Negative Plane. Characters who survive an excursion into the Positive Plane often find that it leaves them with a lasting charge, making it hard to calm down, to stem the flow of ideas, and even to sleep. On the other hand, they also find themselves with a persistent resistance to disease and despair.

\subsection{Shadowfell}
The Shadowfell, also called the Plane of Shadow, is a gloomy dimension whose sky is a black vault with neither sun nor stars.

The Shadowfell overlaps the Material Plane in much the same way as the Feywild. Aside from the bleak landscape, it appears similar to the Material Plane. Travelers from the Material Plane who enter the Shadowfell often observe landmarks similar to the world they left, but distorted and often sinister. A mountain on the Material Plane might be replaced in the Shadowfell by a skull-shaped rock outcropping, a heap of rubble, or the crumbling ruin of a once-great castle. The forests of the Shadowfell hold sinister-looking trees, their branches reaching out to snare travelers' cloaks, and their roots coiling to trip those who pass by.

Shadow dragons and Undead haunt this bleak plane, as do other creatures that thrive in the gloom, including cloakers and darkmantles.

\subsubsection{Shadow Crossings}
Shadow crossings are locations where the veil between the Material Plane and the Shadowfell is so thin that creatures can pass from one plane to the other. A blot of shadow in the corner of a dusty crypt might be a shadow crossing, as might an open grave. Shadow crossings form in gloomy places where spirits or the stench of death lingers, such as battlefields, graveyards, and tombs. They manifest only in darkness, closing as soon as they feel light's kiss.

\subsubsection{Domains of Dread}
In a far-flung corner of the Shadowfell drifts a hidden expanse of roiling mist and vague semireality. At this eerie edge of the multiverse, mysterious entities known as the Dark Powers collect the most wicked beings from across ages and worlds within inescapable, mist-shrouded demiplanes. In these shadowy prisons, the villainous beings become Darklords, able to exercise great power but confined to realms that twist their desires, capturing them in cycles of dread and despair.

Mists surround each of the Domains of Dread, making it difficult to leave one domain and even harder to find a path to another. The Mists rise and fall at the whim of the Dark Powers, and they can even slip across the planes to drag people unwittingly into the dread domains. Those who live in these domains ascribe all sorts of sinister stories to the Mists—any supernatural happening, inexplicable disappearance, or malicious force can be blamed on the Mists.

The following Domains of Dread are among the most infamous. They are described in more detail in \textit{Van Richten's Guide to Ravenloft}.

\subsubsection{Barovia}
The towering spires of \textit{Castle Ravenloft} loom above the valley of \textit{Barovia}, which is ruled by \textit{Strahd von Zarovich}, the first vampire.

\subsubsection{Borca}
Amid opulent estates and impoverished villages, two Darklords—the vicious poisoner Ivana Boritsi and the childishly cruel stalker Ivan Dilisnya—pursue their obsessive schemes.

\subsubsection{Falkovnia}
Empty countryside surrounds ruined or crumbling cities, with only a few pockets of civilization fighting a losing battle against an endless plague of zombies. General Vladeska Drakov commands a fierce military force that desperately clings to power.

\subsubsection{Kalakeri}
A beautiful land of rainforests, rivers, and lakes is a quagmire of intrigue and despair as three royal heirs—transformed into monsters by their depravity and hatred—battle endlessly to claim the throne of their ancient dynasty.

\subsubsection{Lamordia}
Inventors and scholars violate both natural and moral laws amid the frozen bogs and glacial expanses of Lamordia. The worst of them is the domain's Darklord, Doctor Viktra Mordenheim, whose efforts to create life and abolish death have led to the creation of many monsters.

\subsubsection{Mordent}
Death in Mordent heralds the beginning of a haunted afterlife as a restless spirit, for this domain is the realm of ghost stories and hauntings. The dead here earn no rest, no finality, no peace—just a passage into a shadow world of wispy phantoms, mournful groaning, and clanking chains.

\subsubsection{Valachan}
The devious hunter Chajuna roams the jungles of her domain, hunting the most dangerous beasts she can find. When she grows dissatisfied with simpler prey, she draws people into a fatal contest, ensuring that the land remains steeped in blood.

\subsubsection{Shadowfell Despair}
A melancholic atmosphere pervades the Shadowfell, and extended forays to this plane can afflict characters with despair.

When you deem it appropriate, though usually not more than once per day, you can require a character not from the Shadowfell to make a DC 10 Wisdom saving throw. On a failed save, the character is affected by despair. Roll on the Shadowfell Despair table to determine the effects. You can substitute different despair effects of your own creation.

\begin{DndTable}[header=Shadowfell Despair]{cX}

1d6 & Effect\\
1–3 & Apathy. The character has Disadvantage on Death Saving Throw and Initiative rolls.\\
4–5 & Dread. The character has Disadvantage on all saving throws.\\
6 & Delusion. The character has Disadvantage on ability checks and saving throws that use Intelligence, Wisdom, or Charisma.\\
\end{DndTable}

If a character is already suffering a despair effect and fails the saving throw, the new despair effect replaces the old one. After finishing a Long Rest, a character can attempt to overcome the despair with a DC 15 Wisdom saving throw. (The DC is higher because it's harder to shake off despair once it has taken hold.) On a successful save, the despair effect ends for that character. A \textit{Calm Emotions} spell or magic that removes curses cures the despair.

\subsection{Sigil, City of Doors}
At the center of the Outlands, like the axle of a great wheel, is the Spire, a needle-shaped mountain that rises high into the sky. Above this mountain's narrow peak, not part of the Outlands but a plane in its own right, floats the ring-shaped city of Sigil, its myriad structures built on the ring's inner surface. Creatures standing on one of Sigil's streets can see the city curve up over their heads and—most disconcerting of all—the far side of the city directly overhead. Called the City of Doors, this bustling planar metropolis holds countless portals to other planes and worlds.

Sigil is a trader's paradise. Goods and information come here from across the planes. The city sustains a brisk trade in information about the planes, particularly the commands or items required for the operation of particular portals.

The city is the domain of the inscrutable Lady of Pain, a being whose purpose and goals are unknown to even the sages of her city. She appears almost human, although she most definitely isn't. She wears ornate robes that shroud her body, and a mantle of blades coated in blue-green verdigris surrounds her masklike face. No one is certain who or what exactly the Lady of Pain is, but it's widely accepted she's a being on par with deities. Is Sigil her prison? Is she the fallen creator of the multiverse? No one knows—or if they do, they aren't telling.

\textit{Planescape: Adventures in the Multiverse} includes extensive information on \textit{Sigil}.

\subsection{Ysgard}
Ysgard is a rugged realm of soaring mountains, deep fjords, and windswept battlefields, with summers that are long and hot and winters that are cold and unforgiving. Its continents float above oceans of volcanic rock, below which are enormous icy caverns that hold entire kingdoms of giants, humans, dwarves, gnomes, and other beings. Heroes come to Ysgard to test their mettle not only against the plane itself, but also against giants, dragons, and other mighty creatures across Ysgard's vast terrain.

Ysgard is the home of slain heroes who wage eternal battle on fields of glory. Any creature, other than a Construct or Undead, that is killed in combat while in Ysgard is restored to life at dawn the next day. The creature has all its Hit Points restored, and all conditions that affected it before its death are removed.

\begin{DndTable}[header=Layers of Ysgard]{XX}

Layer & Description\\
Ysgard & Immense rivers of floating earth grind together in eternal rumbling.\\
Muspelheim & The ground smokes and burns beneath the earthbergs of the top layer.\\
Nidavellir & Floating chunks of earth are closer together, giving the appearance of endless tunnels with rich mineral deposits.\\
\end{DndTable}

\subsubsection{Ysgard Adventures}
The nature of Ysgard is glory earned through heroic deeds in battle. It's the euphoria of an athlete, the exhilaration of a summer storm, and the triumphant celebration of victory. Since those who die on the plane return to life to fight again the next day, Ysgard can overlook the horrors of war and focus entirely on the glory.

An adventure in Ysgard can be an opportunity for lighthearted combat without consequences, for characters to prove their mettle against truly epic foes and perhaps even against each other. Adventurers might find themselves on the Plain of Ida on the topmost layer of Ysgard, where daily festivals let warriors and athletes show off their bravery and skill. Or they might venture into the lower layers to face greater challenges—or secrets buried in the deep caverns of the plane.

\chapter{Treasure}
Adventurers strive for many things, including glory, knowledge, and justice. Many adventurers also seek something more tangible: treasure. This chapter presents treasure in all its forms, from coins to magic items.

\section{Treasure Themes}
Monsters have treasure preferences, as explained in the \textit{Monster Manual}. These preferences are expressed as themes, which helps you determine what treasures are found in monsters' hoards, as summarized in the Treasure Themes table.

For advice on how to \textit{include treasure in an adventure}, see \textit{chapter 4}.

To randomly determine a magic item found as treasure, use the \textit{tables} at the end of this chapter.

\begin{DndTable}[header=Treasure Themes]{XX}

Theme & Appropriate Treasure\\
\textit{Arcana} & Gemstones plus magic items of an eldritch or esoteric nature\\
\textit{Armaments} & Coins or trade bars plus magic items that are useful in battle\\
\textit{Implements} & Coins, trade bars, or trade goods plus magic items that focus on utility\\
\textit{Relics} & Art objects plus magic items that have religious origins or purposes\\
\end{DndTable}

\section{Coins}
The most basic type of treasure is money, including Copper Pieces (CP), Silver Pieces (SP), Electrum Pieces (EP), Gold Pieces (GP), and Platinum Pieces (PP). See the \textit{Player's Handbook} for their Coins; Coin Values. Fifty coins of any type weigh 1 pound.

\section{Trade Bars}
Because large numbers of coins can be difficult to transport and account for, many merchants prefer to use trade bars—ingots of precious metals and alloys (usually silver). These bars are valued by weight, as shown in the Trade Bars table.

\begin{DndTable}[header=Trade Bars]{XcX}

Bar & Value & Dimensions\\
silver bar (2-pound) & 10 GP & 5 in. long × 2 in. wide × 1/2 in. thick\\
silver bar (5-pound) & 25 GP & 6 in. long × 2 in. wide × 1 in. thick\\
gold bar (5-pound) & 250 GP & 5 in. long × 2 in. wide × 3/4 in. thick\\
\end{DndTable}

\section{Trade Goods}
Merchants commonly exchange trade goods without using currency. The Trade Goods table shows the value of commonly exchanged goods.

\begin{DndTable}[header=Trade Goods]{cX}

Cost & Goods\\
1 CP & 1 lb. of \textit{wheat}\\
2 CP & 2 lb. of \textit{flour} or one \textit{chicken}\\
5 CP & 1 lb. of \textit{salt}\\
1 SP & 1 lb. of \textit{iron} or 1 sq. yd. of Canvas (1 sq. yd.)\\
5 SP & 1 lb. of \textit{copper} or 1 sq. yd. of Cotton Cloth (1 sq. yd.)\\
1 GP & 1 lb. of \textit{ginger} or one \textit{goat}\\
2 GP & 1 lb. of \textit{cinnamon} or \textit{pepper}, or one \textit{sheep}\\
3 GP & 1 lb. of \textit{cloves} or one \textit{pig}\\
5 GP & 1 lb. of \textit{silver} or 1 sq. yd. of Linen (1 sq. yd.)\\
10 GP & 1 lb. of \textit{silk} or one \textit{cow}\\
15 GP & 1 lb. of \textit{saffron} or one \textit{ox}\\
50 GP & 1 lb. of \textit{gold}\\
500 GP & 1 lb. of \textit{platinum}\\
\end{DndTable}

\section{Gemstones}
Gemstones are small, lightweight, and easily secured compared to their same value in coins.

If a treasure hoard includes gemstones, you can use the following tables to randomly determine the kind of gemstones found, based on their value. You can roll once and assume all the gems are the same or roll multiple times to create mixed collections.

\begin{DndTable}[header=10 GP Gemstones]{cX}

1d12 & Stone\\
1 & \textit{Azurite} (mottled deep blue)\\
2 & \textit{Banded agate} (striped brown, blue, white, or red)\\
3 & \textit{Blue quartz} (pale blue)\\
4 & \textit{Eye agate} (circles of gray, white, brown, blue, or green)\\
5 & \textit{Hematite} (gray black)\\
6 & \textit{Lapis lazuli} (light and dark blue with yellow flecks)\\
7 & \textit{Malachite} (striated light and dark green)\\
8 & \textit{Moss agate} (pink or yellow white with mossy gray or green markings)\\
9 & \textit{Obsidian} (black)\\
10 & \textit{Rhodochrosite} (light pink)\\
11 & \textit{Tiger eye} (brown with golden center)\\
12 & \textit{Turquoise} (light blue green)\\
\end{DndTable}

\begin{DndTable}[header=50 GP Gemstones]{cX}

1d12 & Stone\\
1 & \textit{Bloodstone} (dark gray with red flecks)\\
2 & \textit{Carnelian} (orange to red brown)\\
3 & \textit{Chalcedony} (white)\\
4 & \textit{Chrysoprase} (green)\\
5 & \textit{Citrine} (pale yellow brown)\\
6 & \textit{Jasper} (blue, black, or brown)\\
7 & \textit{Moonstone} (white with pale-blue glow)\\
8 & \textit{Onyx} (bands of black and white, or pure black or white)\\
9 & \textit{Quartz} (white, smoky gray, or yellow)\\
10 & \textit{Sardonyx} (bands of red and white)\\
11 & \textit{Star rose quartz} (rosy stone with white star-shaped center)\\
12 & \textit{Zircon} (pale blue green)\\
\end{DndTable}

\begin{DndTable}[header=100 GP Gemstones]{cX}

1d10 & Stone\\
1 & \textit{Amber} (watery gold to rich gold)\\
2 & \textit{Amethyst} (deep purple)\\
3 & \textit{Chrysoberyl} (yellow green to pale green)\\
4 & \textit{Coral} (crimson)\\
5 & \textit{Garnet} (red, brown green, or violet)\\
6 & \textit{Jade} (light green, deep green, or white)\\
7 & \textit{Jet} (deep black)\\
8 & \textit{Pearl} (lustrous white, yellow, or pink)\\
9 & \textit{Spinel} (red, red brown, or deep green)\\
10 & \textit{Tourmaline} (pale green, blue, brown, or red)\\
\end{DndTable}

\begin{DndTable}[header=500 GP Gemstones]{cX}

1d6 & Stone\\
1 & \textit{Alexandrite} (dark green)\\
2 & \textit{Aquamarine} (pale blue green)\\
3 & \textit{Black pearl} (pure black)\\
4 & \textit{Blue spinel} (deep blue)\\
5 & \textit{Peridot} (rich olive green)\\
6 & \textit{Topaz} (golden yellow)\\
\end{DndTable}

\begin{DndTable}[header=1,000 GP Gemstones]{cX}

1d8 & Stone\\
1 & \textit{Black opal} (dark green with black mottling and golden flecks)\\
2 & \textit{Blue sapphire} (medium blue)\\
3 & \textit{Emerald} (deep bright green)\\
4 & \textit{Fire opal} (fiery red)\\
5 & \textit{Opal} (pale blue with green and golden mottling)\\
6 & \textit{Star ruby} (ruby with white star-shaped center)\\
7 & \textit{Star sapphire} (blue sapphire with white star-shaped center)\\
8 & \textit{Yellow sapphire} (fiery yellow or yellow green)\\
\end{DndTable}

\begin{DndTable}[header=5,000 GP Gemstones]{cX}

1d4 & Stone\\
1 & \textit{Black sapphire} (lustrous black with glowing highlights)\\
2 & \textit{Diamond} (blue white, canary, pink, brown, or blue)\\
3 & \textit{Jacinth} (fiery orange)\\
4 & \textit{Ruby} (clear red to deep crimson)\\
\end{DndTable}

\section{Art Objects}
Idols cast of solid gold, necklaces studded with precious stones, paintings of ancient kings, bejeweled dishes—art objects include all these and more.

If a treasure hoard includes art objects, you can use the following tables to randomly determine what art objects are found, based on their value. Roll on a table as many times as there are art objects in the treasure hoard. There can be more than one of a given art object.

\begin{DndTable}[header=25 GP Art Objects]{cX}

1d10 & Object\\
1 & \textit{Silver ewer}\\
2 & \textit{Carved bone statuette}\\
3 & \textit{Gold bracelet}\\
4 & \textit{Cloth-of-gold vestments}\\
5 & \textit{Black velvet mask stitched with silver thread}\\
6 & \textit{Copper chalice with silver filigree}\\
7 & \textit{Pair of engraved bone dice}\\
8 & \textit{Handheld mirror set in a painted wooden frame}\\
9 & \textit{Embroidered silk handkerchief}\\
10 & \textit{Gold locket with a painted portrait inside}\\
\end{DndTable}

\begin{DndTable}[header=250 GP Art Objects]{cX}

1d10 & Object\\
1 & \textit{Gold ring set with bloodstones}\\
2 & \textit{Carved ivory statuette}\\
3 & \textit{Bejeweled gold bracelet}\\
4 & \textit{Silver necklace with a gemstone pendant}\\
5 & \textit{Bronze crown}\\
6 & \textit{Silk vestments with gold embroidery}\\
7 & \textit{Well-made tapestry that is 10 feet by 10 feet}\\
8 & \textit{Brass mug with jade inlay}\\
9 & \textit{Box of turquoise animal figurines}\\
10 & \textit{Gold birdcage with electrum filigree}\\
\end{DndTable}

\begin{DndTable}[header=750 GP Art Objects]{cX}

1d10 & Object\\
1 & \textit{Silver chalice set with moonstones}\\
2 & \textit{Bundle of sheet music representing the lost dirges of a famous composer}\\
3 & \textit{Carved wooden harp with ivory inlay and zircon gems}\\
4 & \textit{Gold idol}\\
5 & \textit{Gold comb shaped like a dragon with red garnets as eyes}\\
6 & \textit{Bottle stopper cork embossed with gold leaf and set with amethysts}\\
7 & \textit{Detailed, life-sized dragonborn skull cast in electrum}\\
8 & \textit{Silver and gold brooch}\\
9 & \textit{Obsidian statuette with gold fittings and inlay}\\
10 & \textit{Painted gold war mask}\\
\end{DndTable}

\begin{DndTable}[header=2,500 GP Art Objects]{cX}

1d10 & Object\\
1 & \textit{Fine gold chain set with a fire opal}\\
2 & \textit{Old masterpiece painting}\\
3 & \textit{Embroidered silk and velvet mantle set with numerous moonstones}\\
4 & \textit{Platinum bracelet set with an emerald}\\
5 & \textit{Embroidered glove set with jewel chips}\\
6 & \textit{Jeweled anklet}\\
7 & \textit{Gold music box}\\
8 & \textit{Gold circlet set with four aquamarines}\\
9 & \textit{Eye patch decorated with tiny blue sapphires and moonstones}\\
10 & A \textit{necklace string of small pink pearls}\\
\end{DndTable}

\begin{DndTable}[header=7,500 GP Art Objects]{cX}

1d10 & Object\\
1 & \textit{Jeweled gold crown}\\
2 & \textit{Jeweled platinum ring}\\
3 & \textit{Gold statuette set with rubies}\\
4 & \textit{Gold cup set with emeralds}\\
5 & \textit{Gold jewelry box with platinum filigree}\\
6 & \textit{Set of gold nesting dolls}\\
7 & \textit{Jade game board with gold playing pieces}\\
8 & \textit{Bejeweled ivory drinking horn with gold filigree}\\
9 & \textit{Gilded royal coach or funeral barge}\\
10 & \textit{Ceremonial gold armor with black pearls}\\
\end{DndTable}

\section{Magic Items}
Magic items are gleaned from the hoards of felled monsters or discovered in long-lost vaults. Such items grant capabilities a character could rarely have otherwise, or they complement their owner's capabilities in wondrous ways.

\begin{DndSidebar}[float=!b]{Magic Item Rules}
Rules for \textit{identifying}, \textit{attuning to}, and \textit{using} magic items appear in the \textit{Player's Handbook}. Additional rules are presented below.



\subsubsection{Attunement Prerequisites}
If a magic item has a class prerequisite, a creature must be a member of that class to attune to the item. If a creature must be a spellcaster to attune to an item, the creature qualifies if it can cast at least one spell using its traits or features, not by using a magic item or the like.



\subsubsection{Items Made for Specific Creatures}
Magic items that are meant to be worn tend to magically adjust themselves to the wearer. However, you can decide that a magic item doesn't adjust its size to fit any wearer. For example, a particular armorer might make items usable only by folk who are sized and shaped like dwarves.



\subsubsection{Unusual Anatomy}
Use your discretion to decide whether a creature can wear an item not made for its anatomy. A ring placed on a tentacle might work, but a yuan-ti with a snakelike tail instead of legs can't wear magic boots.



\subsubsection{Paired Items}
You can allow exceptions to the rule that paired items must both be worn. For example, a character with only one arm might be able to use a single Gloves of Missile Snaring so long as the matching glove is on their person.



\end{DndSidebar}

\subsection{Magic Item Categories}
Every magic item belongs to a category. The Magic Item Categories table lists the nine categories and provides examples. Rules for the categories appear after the table.

\begin{DndTable}[header=Magic Item Categories]{XX}

Category & Examples\\
\textit{Armor} & \textit{+1 Leather Armor}, \textit{+1 Shield}\\
\textit{Potions} & \textit{Potion of Healing}\\
\textit{Rings} & \textit{Ring of Invisibility}\\
\textit{Rods} & \textit{Immovable Rod}\\
\textit{Scrolls} & \textit{Spell Scroll}\\
\textit{Staffs} & \textit{Staff of Striking}\\
\textit{Wands} & \textit{Wand of Fireballs}\\
\textit{Weapons} & \textit{+1 Ammunition}, \textit{+1 Longsword}\\
\textit{Wondrous Items} & \textit{Bag of Holding}, \textit{Boots of Elvenkind}\\
\end{DndTable}

\subsubsection{Armor}
An item in the Armor category is typically a magical version of armor from the \textit{Player's Handbook}. Unless an armor's description notes otherwise, the armor must be worn for its magic to function.

Some suits of magic armor specify the type of armor they are, such as \textit{Chain Mail} or \textit{Plate Armor}. If no type is specified, choose the type or determine it randomly.

\subsubsection{Potions}
An item in the Potion category might be a magical brew that must be imbibed or an oil that must be applied to a creature or an object. A typical potion consists of 1 ounce of liquid in a vial.

\subparagraph{Using a Potion}
Potions are consumable items. Drinking a potion or administering it to another creature requires a Bonus Action. Applying an oil might take longer as specified in its description. Once used, a potion takes effect immediately, and it is used up.

\subparagraph{Mixing Potions}
A character might drink one potion while still under the effects of another or pour several potions into a single container. The strange ingredients used in creating potions can result in unpredictable interactions.

When a character mixes two potions together, roll on the Potion Miscibility table. If more than two are combined, roll again for each subsequent potion, combining the results. Unless the effects are immediately obvious, reveal them only when they become evident.

\begin{DndTable}[header=Potion Miscibility]{cX}

1d100 & Result\\
01 & Both potions lose their effects, and the mixture creates a magical explosion in a 5-foot-radius Sphere [Area of Effect] centered on itself. Each creature in that area takes 4d10 Force damage.\\
02–08 & Both potions lose their effects, and the mixture becomes an ingested poison of your choice (see "\textit{Poison}" in \textit{chapter 3}).\\
09–15 & Both potions lose their effects.\\
16–25 & One potion loses its effect.\\
26–35 & Both potions work, but with their numerical effects and durations halved. If a potion has no numerical effect and no duration, it instead loses its effect.\\
36–90 & Both potions work normally.\\
91–99 & Both potions work, but the numerical effects and duration of one potion are doubled. If neither potion has anything to double in this way, they work normally.\\
00 & Only one potion works, but its effects are permanent. Choose the simplest effect to make permanent or the one that seems the most fun. For example, a \textit{Potion of Healing} might increase the drinker's Hit Point maximum by 2d4 + 2, or a \textit{Potion of Invisibility} might give the drinker the Invisible condition indefinitely. At your discretion, a \textit{Dispel Magic} spell or similar magic might end this lasting effect.\\
\end{DndTable}

\subsubsection{Rings}
For its magic to function, an item in the Ring category must be worn on a finger or a similar digit unless its description notes otherwise.

\subsubsection{Rods}
An item in the Rod category is a scepter usually made of metal, wood, or bone. A typical rod weighs 2 to 5 pounds.

Unless its description notes otherwise, a rod can be used as an \textit{Arcane Focus}.

\subsubsection{Scrolls}
An item in the Scroll category is a roll of paper or parchment, sometimes attached to wooden rods and typically kept safe in a tube of ivory, jade, leather, metal, or wood. The most prevalent scroll is the \textit{Spell Scroll}, a spell stored in written form. However, some scrolls, like the \textit{Scroll of Protection}, bear an incantation that isn't a spell.

\subparagraph{Using a Scroll}
Scrolls are consumable items. Unleashing the magic in a scroll requires the user to read the scroll. When its magic has been invoked, the scroll can't be used again. Its words fade, or it crumbles into dust.

Any creature that can understand a written language can read a scroll and attempt to activate it unless its description notes otherwise.

\subsubsection{Staffs}
Items in the Staff category vary widely in appearance: some are of nearly equal diameter throughout and smooth, others are gnarled and twisted, some are made of wood, and others are composed of polished metal or crystal. A staff weighs between 2 and 7 pounds and serves well as a walking stick or cane.

Unless its description notes otherwise, a staff can be used as a nonmagical \textit{Quarterstaff} and an \textit{Arcane Focus}.

\subsubsection{Wands}
An item in the Wand category is typically 12 to 15 inches long and crafted of metal, bone, or wood. It is tipped with metal, crystal, stone, or some other material.

Unless its description notes otherwise, a wand can be used as an \textit{Arcane Focus}.

\subsubsection{Weapons}
A magic weapon is typically a magical version of a weapon from the \textit{Player's Handbook}. Some magic weapons specify the type of weapon they are in their descriptions, such as a \textit{Longsword} or \textit{Longbow}. If no weapon type is specified, you may choose the type or determine it randomly.

\subparagraph{Ammunition}
If a magic weapon has the AF property, ammunition fired from it is considered magical for the purpose of any rule that cares whether a weapon is magical or not.

\subsubsection{Wondrous Items}
Wondrous Items include wearable items such as boots, belts, capes, amulets, brooches, and circlets. Bags, carpets, figurines, horns, musical instruments, and more also fall into this category.

\subsection{Magic Item Rarity}
Every magic item has a rarity, which provides a rough measure of an item's power relative to other magic items. The rarities are shown in the Magic Item Rarities and Values table.

Common magic items, such as a \textit{Potion of Healing}, are the most plentiful. Artifacts, such as the \textit{Wand of Orcus}, are priceless, unique, and difficult to acquire.

\subsubsection{Magic Item Values by Rarity}
Common magic items can often be bought in a town or city. Uncommon and Rare magic items are usually found only in cities, and rarer magic items might be sold only in wondrous locations, such as the \textit{City of Brass} or \textit{Sigil}. If you allow characters to buy and sell magic items in your campaign, rarity can help you set prices for those items. Gold Piece values are provided in the Magic Item Rarities and Values table, though a seller might ask for a service rather than coin as payment.

If a magic item incorporates an item that has a purchase cost in the \textit{Player's Handbook} (such as a weapon or a suit of armor), add that item's cost to the magic item's value. For example, +1 Plate Armor has a value of 5,500 GP, which is the sum of a Rare magic item's value (4,000 GP) and the cost of \textit{Plate Armor} (1,500 GP).

\begin{DndTable}[header=Magic Item Rarities and Values]{cX}

Rarity & Value*\\
Common & 100 GP\\
Uncommon & 400 GP\\
Rare & 4,000 GP\\
Very Rare & 40,000 GP\\
Legendary & 200,000 GP\\
Artifact & Priceless\\
\end{DndTable}

\subsection{Awarding Magic Items}
Awarding magic items is the purview of the DM. You can award a magic item because the story calls for it or the players would be especially pleased to have it. This section helps you to determine which magic items end up in the characters' possession.

\begin{DndSidebar}[float=!b]{Are Magic Items Necessary?}
The D\&D game assumes that magic items appear sporadically and that they are a boon unless an item bears a curse. Characters and monsters are built to face each other without the help of magic items, which means that having a magic item makes a character more powerful or versatile than a generic character of the same level. As DM, you never have to worry about awarding magic items just so the characters can keep up with the campaign's threats. Magic items are truly prizes—desirable but not necessary.



\end{DndSidebar}

\subsubsection{Magic Items Awarded by Level}
The Magic Items Awarded by Level table shows the number of magic items a D\&D party typically gains during a campaign, totaling one hundred magic items by level 20. The table shows how many items of each rarity are meant to be handed out during each of the four tiers of play.

Artifacts are omitted from the table because they are most often used as plot devices in high-level adventures, and characters rarely have them for long (either because the Artifacts are meant to be destroyed or because the campaign is nearing its end).

\subparagraph{Player Wish List}
Encourage your players to keep a wish list of magic items they hope their characters will find in the course of the campaign. If you want to award a magic item but don't have a specific magic item in mind, you can pick an item of the appropriate rarity from your players' wish list.

\subparagraph{Overstocking an Adventure}
When creating or modifying an adventure, assume that the characters won't find all the magic items you place in it. An adventure usually can include a number of items that's 25 percent higher than the number in the Magic Items Awarded by Level table (round up). For example, an adventure designed to take characters from level 1 to 4 might include fourteen items rather than eleven, in the expectation that three items won't be found.

\begin{DndTable}[header=Magic Items Awarded by Level]{ccccccc}

Character Level (Tier of Play) & Common Items & Uncommon Items & Rare Items & Very Rare Items & Legendary Items & All Items\\
1–4 (tier 1) & 6 & 4 & 1 & 0 & 0 & 11\\
5–10 (tier 2) & 10 & 17 & 6 & 1 & 0 & 34\\
11–16 (tier 3) & 3 & 7 & 11 & 7 & 2 & 30\\
17–20 (tier 4) & 0 & 0 & 5 & 11 & 9 & 25\\
\textbf{Total} & \textbf{19} & \textbf{28} & \textbf{23} & \textbf{19} & \textbf{11} & \textbf{100}\\
\end{DndTable}

\subsubsection{Magic Item Tracker}
You can use the Magic Item Tracker sheet to track how many magic items the characters have acquired. Each time the characters get a magic item, put a check mark in one of the empty circles corresponding to the item's rarity and the current level range of the characters. If the characters gain a magic item of a rarity that has no unchecked circles at the current level range, check off an empty circle from a lower tier. If all lower level ranges also have no circles left, check off an empty circle from a higher level range.

\subsubsection{Random Magic Item Rarity}
When you decide that a treasure contains magic items, there are two ways to determine the rarity of those items. You can choose an appropriate rarity based on the items you've given out already (using the Magic Item Tracker sheet to keep track), or you can roll on the Magic Item Rarities table.

To use the table, find the level of the characters in the top row. Roll 1d100, and read down that column to find your roll. Then read across to the right column to find the rarity of the item.

\begin{DndTable}[header=Magic Item Rarities]{ccccc}

01–54 & 01–30 & 01–11 & — & Common\\
55–91 & 31–81 & 12–34 & — & Uncommon\\
92–00 & 82–98 & 35–70 & 01–20 & Rare\\
— & 99–00 & 71–93 & 21–64 & Very Rare\\
— & — & 94–00 & 65–00 & Legendary\\
\end{DndTable}

\begin{DndSidebar}[float=!b]{Magic Items for Starting Characters}
If you're starting a campaign for characters above level 1, the \textit{Player's Handbook} offers Starting Equipment at Higher Levels for how many magic items such characters should start with and the rarity of those items. Consider these approaches to determining the items each character receives:




\begin{description}
\item[DM Choice:]
Choose items for each character using your own judgment.
\item[Player Choice:]
Let the players choose whatever items they want, within the specified rarity.
\item[Random Determination:]
Use the tables at the end of this chapter to randomly determine starting items. Use the \textit{Arcana tables} for Sorcerers, Warlocks, and Wizards. Use the \textit{Armaments tables} for Barbarians, Fighters, Paladins, and Rangers. Use the \textit{Implements tables} for Bards, Monks, and Rogues. Use the \textit{Relics tables} for Clerics and Druids. Feel free to vary the tables you use to give each character a mix of useful items and to match character themes. For example, for a Fighter with the Eldritch Knight subclass, you might choose one item from the \textit{Arcana tables} and the rest of the character's items from the \textit{Armaments tables}.
\end{description}



\end{DndSidebar}

\subsection{Activating a Magic Item}
It usually takes a Magic action to activate a magic item. The item's user might also need to do something special. The description of each item category or individual item details how an item is activated. Certain items use the following rules for their activation.

\subsubsection{Command Word}
A command word is a word or short phrase that must be spoken or signed for an item to work. Spoken command words must be audible and fail to work in areas where all sound is suppressed, as in the area of the \textit{Silence} spell.

\subsubsection{Consumable Items}
Some items are consumed—used up, in other words—when they are activated. A \textit{Potion of Healing} must be swallowed, for example, while the writing vanishes from a scroll when it is read. Once used, a consumable item loses its magic.

\subsubsection{Spells Cast from Items}
Some magic items allow the user to cast a spell from the item. The spell is cast at the lowest possible spell and caster level, doesn't expend any of the user's spell slots, and requires no components unless the item's description notes otherwise. The spell uses its normal casting time, range, and duration, and the user of the item must concentrate if the spell requires Concentration. Many items, such as Potions, bypass the casting of a spell and confer the spell's effects with its usual duration. Certain items make exceptions to these rules, changing the casting time, duration, or other parts of a spell.

A magic item may require the user to use their own spellcasting ability when casting a spell from the item. If the user has more than one spellcasting ability, the user chooses which one to use with the item. If the user doesn't have a spellcasting ability, their spellcasting ability modifier is +0 for the item, and the user's Proficiency Bonus applies.

\subsubsection{Charges}
Some magic items have charges that must be expended to activate their properties. The number of charges an item has remaining is revealed when the \textit{Identify} spell is cast on it. A creature attuned to an item knows how many charges the item has and how many it regains.

\subsection{"The Next Dawn"}
Magic items often have charges or properties that recharge at the next dawn or some other specified time. If such an item is on a world or plane of existence where the specified event doesn't occur, the DM determines when the item recharges.

\subsection{Cursed Items}
A magic item's description specifies whether it bears a curse. Most methods of identifying items, including the \textit{Identify} spell, fail to reveal such a curse.

Attunement to a cursed item can't be ended voluntarily unless the curse is broken first, such as with a \textit{Remove Curse} spell.

\subsection{Magic Item Resilience}
A magic item is at least as durable as a nonmagical item of its kind. Most magic items, other than Potions and Scrolls, have Resistance to all damage.

An Artifact can be destroyed only in some special way. Otherwise, it is impervious to damage. Learning how to destroy an Artifact usually requires research or the completion of a quest.

\subsection{Crafting Magic Items}
The \textit{Player's Handbook} contains rules on \textit{brewing Potions of Healing} and \textit{scribing Spell Spell Scrolls}. To create other magic items, follow the rules below. In these rules, "you" refers to the character crafting the magic item.

\subsubsection{Arcana Proficiency}
To craft a magic item, you and any assistants must have proficiency in the Arcana skill.

\subsubsection{Tools}
The Magic Item Tools table lists which tool is required to make a magic item of each category. You must use the required tool to make an item and have proficiency with that tool. Any assistants must also have proficiency with it. For more \textit{information on the tools}, see the \textit{Player's Handbook}.

\begin{DndTable}[header=Magic Item Tools]{XX}

Item Category & Required Tool\\
Armor & \textit{Leatherworker's Tools}, \textit{Smith's Tools}, or \textit{Weaver's Tools} depending on the kind of armor as noted in the tools' descriptions\\
Potion & \textit{Alchemist's Supplies} or \textit{Herbalism Kit}\\
Ring & \textit{Jeweler's Tools}\\
Rod & \textit{Woodcarver's Tools}\\
Scroll & \textit{Calligrapher's Supplies}\\
Staff & \textit{Woodcarver's Tools}\\
Wand & \textit{Woodcarver's Tools}\\
Weapon & \textit{Leatherworker's Tools}, \textit{Smith's Tools}, or \textit{Woodcarver's Tools} depending on the kind of weapon as noted in the tools' descriptions\\
Wondrous Item & \textit{Tinker's Tools} or the tool required to make the nonmagical item on which the magic item is based\\
\end{DndTable}

\subsubsection{Spells}
If a magic item allows its user to cast any spells from it, you must have all those spells prepared every day you spend crafting the item.

\subsubsection{Time and Cost}
Crafting a magic item takes an amount of time and money based on the item's rarity as shown in the Magic Item Crafting Time and Cost table.

\subparagraph{Work per Day}
For each day of crafting, you must work for 8 hours. If an item requires multiple days, those days needn't be consecutive.

\subparagraph{Assistants}
Characters can combine their efforts to shorten the crafting time. Divide the time needed to create an item by the number of characters working on it. Normally, only one other character can assist you, but the DM might allow more assistants.

\subparagraph{Raw Materials}
The cost in the table represents the raw materials needed to make a magic item. The DM determines whether appropriate raw materials are available. In a city, there is a 75 percent chance that the materials are available, and in any other settlement, that chance is 25 percent. If materials aren't available, you must wait at least 7 days before checking on the availability again.

If a magic item incorporates an item that has a purchase cost (such as a weapon or a suit of armor), you must also pay that entire cost or craft that item \textit{using the rules} in the \textit{Player's Handbook}. For example, to make +1 Plate Armor, you must pay 3,500 GP or pay 2,000 GP and craft the armor.

\begin{DndTable}[header=Magic Item Crafting Time and Cost]{XXX}

Item Rarity & Time* & Cost*\\
Common & 5 days & 50 GP\\
Uncommon & 10 days & 200 GP\\
Rare & 50 days & 2,000 GP\\
Very Rare & 125 days & 20,000 GP\\
Legendary & 250 days & 100,000 GP\\
\end{DndTable}

\subsection{Magic Item Special Features}
You can add distinctiveness to a magic item by thinking about its backstory. Who made the item? Is anything unusual about its construction? Why was it made, and how was it used originally? What minor magical quirks set it apart from other items of its kind? Answering these questions can help turn a generic magic item, such as a \textit{+1 Longsword}, into a more flavorful discovery.

Use the following tables to fill in details about a magic item's history. Some table entries make more sense for certain items than for others. If you roll something that doesn't make sense, roll again, choose a more appropriate entry, or use the rolled detail as inspiration to make up your own special feature.

On the Magic Item's Minor Property table and the Magic Item's Quirk table, "you" refers to the item's bearer.

\begin{DndTable}[header=Magic Item's Creator or Intended User]{cX}

1d20 & Creator or Intended User\\
1 & Aberration. The item is ancient. At a glance, it seems to be covered with mucus.\\
2 & Celestial. The item is half the normal weight and inscribed with feathered wings, suns, and stars. Fiends find it repulsive.\\
3 & Devotees of Lolth. The item is half the normal weight. It is inscribed with spiders and webs in honor of Lolth.\\
4 & Dragon. This item incorporates precious metals and gems from a dragon's hoard. It grows warm when within 120 feet of a Dragon.\\
5–6 & Dwarf. The item is durable and has Dwarvish runes worked into its design. It might be associated with a clan that would like to see it returned to their ancestral halls.\\
7 & Elemental Air. The item is half the normal weight and feels hollow. If it's made of fabric, it is diaphanous.\\
8 & Elemental Earth. This item might be crafted from stone. Any cloth or leather elements are studded with finely polished rock.\\
9 & Elemental Fire. This item is warm to the touch, and any metal parts are crafted from black iron. Flame imagery covers its surface.\\
10 & Elemental Water. Lustrous fish scales replace leather or cloth on this item, and seashells and worked coral (as hard as any metal) replace metal portions.\\
11–12 & Elf. The item is half the normal weight. It is adorned with symbols of nature: leaves, vines, stars, and the like.\\
13 & Fey. The item is exquisitely crafted from the finest materials and glows with a pale radiance in moonlight, shedding Dim Light in a 5-foot radius. Any metal in the item is silver or mithral rather than iron or steel.\\
14 & Fiend. The item is made of iron or horn, and any cloth or leather components are crafted from the hide of Fiends. Leering faces or vile runes are engraved on its surface. Celestials find it repulsive.\\
15 & Giant. The item is larger than normal and was crafted by Giants for use by their smaller allies.\\
16 & Gnome. The item is crafted to appear ordinary and well used. It could also incorporate gears and mechanical components, even if these aren't essential to its function.\\
17–19 & Human. The item was created during the heyday of a fallen human kingdom, or it is tied to a human of legend. It might hold writing in a forgotten language or symbols whose significance is lost to the ages.\\
20 & Undead. The item incorporates symbols of death, such as bones and skulls, and it might be crafted from parts of corpses. It feels cold to the touch.\\
\end{DndTable}

\begin{DndTable}[header=Magic Item's History]{cX}

1d8 & History\\
1 & Arcane. This item was created for an ancient order of spellcasters and bears the order's symbol.\\
2 & Bane. This item was created to oppose creatures of a particular type, such as Aberrations or Dragons.\\
3 & Heroic. A great hero once wielded this item. Anyone who knows the item's history expects great deeds from the new owner.\\
4 & Ornament. The item honors a special event. Inset gemstones, gold or platinum inlays, and gold or silver filigree adorn its surface.\\
5 & Prophecy. The item features in a prophecy: its bearer is destined to play a key role in future events.\\
6 & Religious. This item was used in religious ceremonies dedicated to a particular deity. It has religious symbols worked into it.\\
7 & Sinister. This item is linked to a deed of great evil, such as a massacre or an assassination. It might have a name or be closely associated with a villain who used it.\\
8 & Symbol of Power. This item was once used as part of royal regalia or as a badge of high office.\\
\end{DndTable}

\begin{DndTable}[header=Magic Item's Minor Property]{cX}

1d20 & Minor Property\\
1–2 & Beacon. You can take a Bonus Action to cause the item to shed Bright Light in a 10-foot radius and Dim Light for an additional 10 feet, or to extinguish the light.\\
3 & Compass. You can take a Magic action to learn which way is magnetic north. Nothing happens if this property is used in a location that has no magnetic north.\\
4 & Delver. While underground, you always know the item's depth below the surface and the direction to the nearest staircase, ramp, or other path leading upward.\\
5–6 & Guardian. The item warns you, granting a +2 bonus to your Initiative rolls if you don't have the Incapacitated condition.\\
7–8 & Harmonious. Attuning to this item takes only 1 minute.\\
9 & Key. The item is used to unlock a container, chamber, vault, or door.\\
10 & Secret Message. A message is hidden somewhere on the item. It might be visible only at a certain time, under the light of one phase of the moon, or in a specific location.\\
11–12 & Sentinel. The DM chooses a kind of creature, such as mind flayers or trolls. This item glows faintly when such creatures are within 120 feet of it.\\
13 & Songcraft. Whenever this item is struck or is used to strike a foe, you hear a fragment of an ancient song.\\
14–15 & Strange Material. The item was created from a material that is bizarre given its purpose. Its durability is unaffected.\\
16 & Temperate. You are unharmed by temperatures of 0 degrees Fahrenheit or lower, and 100 degrees Fahrenheit or higher.\\
17 & Unbreakable. The item can't be broken. Special means must be used to destroy it.\\
18 & War Leader. You can take a Magic action to cause your voice or signal to carry clearly for up to 600 feet until the end of your next turn.\\
19 & Waterborne. This item floats on water and other liquids. You have Advantage on Strength (Athletics) checks to swim.\\
20 & Roll twice, rerolling any additional 20s.\\
\end{DndTable}

\begin{DndTable}[header=Magic Item's Quirk]{cX}

1d8 & Quirk\\
1 & Blissful. You feel fortunate and optimistic about what the future holds. Butterflies and other harmless creatures might frolic in the item's presence.\\
2 & Confident. The item helps you feel self-assured.\\
3 & Covetous. You become obsessed with material wealth.\\
4 & Fragile. The item crumbles, frays, chips, or cracks slightly when wielded, worn, or activated. This quirk has no effect on its properties.\\
5 & Loud. The item makes a loud noise—such as a clang, a shout, or a resonating gong—when used.\\
6 & Metamorphic. The item periodically alters its appearance in slight ways. You have no control over these minor alterations, which have no effect on the item's use.\\
7 & Painful. You experience a harmless flash of pain when using the item.\\
8 & Repulsive. You feel a sense of distaste when in contact with the item and continue to experience discomfort while bearing it.\\
\end{DndTable}

\subsection{Artifacts}
An Artifact is a unique, singularly powerful magic item with its own origin and history. It could have been created in the midst of a crisis that threatened a kingdom, a world, or the entire multiverse, and carry the weight of that pivotal moment in history.

Some Artifacts appear when they are needed most. For others, the reverse is true; when these Artifacts are discovered, the world trembles at the ramifications of the find. In either case, introducing an Artifact into a campaign requires forethought. It could be an item that opposing sides are hoping to claim, or it might be something the adventurers need to overcome their greatest challenge.

Characters don't typically find Artifacts in the normal course of adventuring. In fact, Artifacts appear only when you want them to, for they are as much plot devices as magic items. Tracking down and recovering an Artifact is often the main goal of an adventure. Characters must chase down rumors, undergo significant trials, and venture into dangerous, half-forgotten places to find the Artifact they seek. Alternatively, a major villain might already have the Artifact. Obtaining and destroying the Artifact could be the only way to ensure that its power can't be used for evil.

\subsubsection{Artifact Properties}
In addition to its defined properties, an Artifact might have other properties that are either beneficial or detrimental. You can choose such properties from the tables in this section or determine them randomly. You can also invent new beneficial and detrimental properties. These properties typically change each time an Artifact appears in the world.

An Artifact can have as many as four minor beneficial properties and two major beneficial properties. It can have as many as four minor detrimental properties and two major detrimental properties.

\begin{DndTable}[header=Minor Beneficial Properties]{cX}

1d100 & Property\\
01–20 & While attuned to the Artifact, you gain proficiency in one skill of the DM's choice.\\
21–30 & While attuned to the Artifact, you have Immunity to the Poisoned condition.\\
31–40 & While attuned to the Artifact, you have Immunity to the Charmed and Frightened conditions.\\
41–50 & While attuned to the Artifact, you have Resistance to one damage type of the DM's choice.\\
51–60 & While attuned to the Artifact, you can cast one cantrip (chosen by the DM) from it.\\
61–70 & While attuned to the Artifact, you can cast one level 1 spell (chosen by the DM) from it. After you cast the spell, roll 1d6. On a roll of 1–5, you can't cast it again in this way until the next dawn.\\
71–80 & As 61–70 above, except the spell is level 2.\\
81–90 & As 61–70 above, except the spell is level 3.\\
91–00 & While attuned to the Artifact, you gain a +1 bonus to Armor Class.\\
\end{DndTable}

\begin{DndTable}[header=Major Beneficial Properties]{cX}

1d100 & Property\\
01–20 & While attuned to the Artifact, one of your ability scores (DM's choice) increases by 2, to a maximum of 24.\\
21–30 & While attuned to the Artifact, you regain 1d6 Hit Points at the start of each of your turns if you have at least 1 Hit Point.\\
31–40 & When you hit with an attack roll while attuned to the Artifact, the target takes an extra 1d6 Force damage.\\
41–50 & While you're attuned to the Artifact, your Speed increases by 10 feet.\\
51–60 & While attuned to the Artifact, you can cast one level 4 spell (chosen by the DM) from it. After you cast the spell, roll 1d6. On a roll of 1–5, you can't cast it again in this way until the next dawn.\\
61–70 & As 51–60 above, except the spell is level 5.\\
71–80 & As 51–60 above, except the spell is level 6.\\
81–90 & As 51–60 above, except the spell is level 7.\\
91–00 & While attuned to the Artifact, you have Immunity to the Blinded, Deafened, Petrified, and Stunned conditions.\\
\end{DndTable}

\begin{DndTable}[header=Minor Detrimental Properties]{cX}

1d100 & Property\\
01–08 & While attuned to the Artifact, you have Disadvantage on any ability check or saving throw that uses Strength or Constitution.\\
09–16 & While attuned to the Artifact, you have Disadvantage on Intelligence, Wisdom, and Charisma saving throws.\\
17–24 & While attuned to the Artifact, you have Vulnerability to Poison damage.\\
25–32 & While attuned to the Artifact, you have the Blinded condition when you're more than 10 feet away from it.\\
33–40 & While attuned to the Artifact, you have the Deafened condition when you're more than 10 feet away from it.\\
41–48 & While attuned to the Artifact, you lose all sense of smell.\\
49–66 & While you're attuned to the Artifact, your appearance changes as the DM decides.\\
67–72 & While attuned to the Artifact, you emit a sour stench noticeable from up to 10 feet away.\\
73–76 & Whenever you touch a nonmagical gem or an art object while attuned to this Artifact, the value of the gem or art object is reduced by half. This affects a particular object only once.\\
77–80 & While you're attuned to the Artifact, all \textit{Holy Water} within 10 feet of you is destroyed.\\
81–84 & While you're attuned to the Artifact, nonmagical flames are extinguished within 30 feet of you.\\
85–88 & While you're attuned to the Artifact, other creatures can't take Short or Long Rests while within 300 feet of you.\\
89–92 & While attuned to the Artifact, you kill any nonmagical vegetation you touch that isn't a creature.\\
93–00 & While you're attuned to the Artifact, Beasts within 30 feet of you that have a Challenge Rating of 6 or lower are Hostile [Attitude] toward you.\\
\end{DndTable}

\begin{DndTable}[header=Major Detrimental Properties]{cX}

1d100 & Property\\
01–09 & You can't attune to other magic items while you're attuned to the Artifact. When you become attuned to the Artifact, your Attunement to other magic items ends immediately.\\
10–18 & When you become attuned to the Artifact, a random one of your ability scores is reduced by 2, to a minimum of 3. A \textit{Greater Restoration} spell restores the ability.\\
19–27 & When you become attuned to the Artifact, you take 8d10 Psychic damage.\\
28–36 & The first time you become attuned to the Artifact, it gives you a quest determined by the DM. You can't use any of the Artifact's properties until you complete the quest.\\
37–45 & Each time you become attuned to the Artifact, there is a 10 percent chance that you attract the attention of a god who sends an avatar to wrest the Artifact from you. The avatar has the same alignment as its creator and uses the \textbf{Empyrean} stat block. Once it obtains the Artifact, the avatar vanishes.\\
46–54 & Each time you become attuned to the Artifact, you must succeed on a DC 10 Constitution saving throw or die from the shock. If you die, you're instantly transformed into a \textbf{Wight} under the DM's control that must protect the Artifact.\\
55–63 & The Artifact dilutes potions within 10 feet of itself, rendering them nonmagical.\\
64–72 & The Artifact erases scrolls within 10 feet of itself, rendering them nonmagical.\\
73–81 & While you're attuned to the Artifact, creatures of a particular type other than Humanoid (chosen by the DM) are always Hostile [Attitude] toward you.\\
82–90 & While attuned to the Artifact, you have Vulnerability to all damage.\\
91–96 & The Artifact imprisons a \textbf{Death Slaad}. Each time you become attuned to the Artifact, the slaad has a 10 percent chance of escaping, whereupon it appears in an unoccupied space as close to you as possible and attacks you.\\
97–00 & While attuned to the Artifact, you can't spend Hit Point Dice or regain Hit Points.\\
\end{DndTable}

\subsection{Sentient Magic Items}
Some magic items have sentience and personality. Such an item might be possessed, haunted by the spirit of a previous owner, or self-aware thanks to the magic used to create it. A sentient item might be a cherished ally to its wielder or a continual thorn in the side.

Most sentient items are weapons, but other kinds of items can manifest sentience. Single-use items such as potions and scrolls are never sentient.

The DM controls sentient magic items and their activated properties. A bearer who maintains a good relationship with the item can access those properties. If the relationship is strained, a conflict can ensue (see "\textit{Conflict}" below).

\subsubsection{Sentient Magic Item Traits}
When you make a sentient magic item, you create the item's persona much as you would create an NPC (as described in the "\textit{Nonplayer Characters}" section of \textit{chapter 3}), with these exceptions.

\subparagraph{Abilities}
A sentient magic item has Intelligence, Wisdom, and Charisma scores. Choose the item's abilities, or determine them randomly as follows: roll 4d6 for each one, dropping the lowest roll and totaling the rest.

\subparagraph{Alignment}
A sentient magic item has an alignment. Its creator or nature might suggest an alignment. Otherwise, pick an alignment or roll on the Sentient Item's Alignment table.

\subparagraph{Communication}
A sentient item communicates by sharing its emotions, broadcasting its thoughts telepathically, or speaking aloud. You can choose how it communicates or roll on the Sentient Item's Communication table.

\subparagraph{Senses}
A sentient item can perceive its surroundings out to a limited range. You can choose its senses or roll on the Sentient Item's Senses table.

\subparagraph{Special Purpose}
You can give a sentient item an objective it pursues, perhaps to the exclusion of all else. As long as the wielder's use of the item aligns with that special purpose, the item remains cooperative. Deviating from this course might cause conflict between the wielder and the item (see "\textit{Conflict}" below). You can pick a special purpose or roll on the Sentient Item's Special Purpose table.

\begin{DndTable}[header=Sentient Item's Alignment]{cX}

1d100 & Alignment\\
01–15 & Lawful Good\\
16–35 & Neutral Good\\
36–50 & Chaotic Good\\
51–63 & Lawful Neutral\\
64–73 & Neutral\\
74–85 & Chaotic Neutral\\
86–89 & Lawful Evil\\
90–96 & Neutral Evil\\
97–00 & Chaotic Evil\\
\end{DndTable}

\begin{DndTable}[header=Sentient Item's Communication]{cX}

1d10 & Communication\\
1–6 & The item communicates by transmitting emotion to the creature carrying or wielding it.\\
7–9 & The item speaks one or more languages.\\
10 & The item speaks one or more languages. In addition, the item can communicate telepathically with any creature that carries or wields it.\\
\end{DndTable}

\begin{DndTable}[header=Sentient Item's Senses]{cX}

1d4 & Senses\\
1 & Hearing and standard vision out to 30 feet\\
2 & Hearing and standard vision out to 60 feet\\
3 & Hearing and standard vision out to 120 feet\\
4 & Hearing and Darkvision out to 120 feet\\
\end{DndTable}

\begin{DndTable}[header=Sentient Item's Special Purpose]{cX}

1d10 & Special Purpose\\
1 & Aligned. The item seeks to defeat or destroy those of a diametrically opposed alignment. Such an item is never Neutral.\\
2 & Bane. The item seeks to thwart or destroy creatures of a particular type, such as Constructs, Fiends, or Undead.\\
3 & Creator Seeker. The item seeks its creator and wants to understand why it was created.\\
4 & Destiny Seeker. The item believes it and its bearer have key roles to play in future events.\\
5 & Destroyer. The item craves destruction and goads its user to fight arbitrarily.\\
6 & Glory Seeker. The item seeks renown as the greatest magic item in the world by winning fame or notoriety for its user.\\
7 & Lore Seeker. The item craves knowledge or is determined to solve a mystery, learn a secret, or unravel a cryptic prophecy.\\
8 & Protector. The item seeks to defend a particular kind of creature, such as elves or werewolves.\\
9 & Soulmate Seeker. The item seeks another sentient magic item, perhaps one that is similar to itself.\\
10 & Templar. The item seeks to defend the servants and interests of a particular deity.\\
\end{DndTable}

\subsubsection{Conflict}
When the bearer of a sentient item acts in a manner opposed to the item's alignment or purpose, conflict can arise. When such a conflict occurs, the item's bearer makes a Charisma saving throw (DC 12 plus the item's Charisma modifier). On a failed save, the item makes one or more of the following demands:


\begin{description}
\item[Chase My Dreams.]
The item demands that its bearer pursue the item's goals to the exclusion of all other goals.
\item[Get Rid of It.]
The item demands that its bearer dispose of anything the item finds repugnant.
\item[It's Time for a Change.]
The item demands to be given to someone else.
\item[Keep Me Close.]
The item insists on being carried or worn at all times.
\end{description}

If its bearer refuses to comply with the item's demands, the item can do any of the following:


\begin{itemize}
\item Make it impossible for its bearer to attune to it.
\item Suppress one or more of its activated properties.
\item Attempt to take control of its bearer, whereupon the bearer makes a Charisma saving throw (DC 12 plus the item's Charisma modifier). On a failed save, the bearer has the Charmed condition for 1d12 hours. While Charmed in this way, the bearer must try to follow the item's commands. If the bearer takes damage, it repeats the save, ending the effect on a success. Whether or not the attempt to control its bearer succeeds, the item can't use this power again until the next dawn.
\end{itemize}

\section{Magic Items A-Z}
Magic items are presented in alphabetical order.

If a magic item description capitalizes a creature's name and presents it in \textbf{bold} type, that's a visual cue pointing you to the creature's stat block. Unless the text states otherwise, the stat block is in the \textit{Monster Manual}. How to read and use a stat block is explained in the \textit{Monster Manual} and to a lesser degree in the \textit{Player's Handbook}.


\begin{itemize}
\begin{multicols}{3}
\item \textit{Adamantine Armor}
\item \textit{Adamantine Weapon}
\item \textit{Alchemy Jug}
\item \textit{+1 Ammunition}
\item \textit{+2 Ammunition}
\item \textit{+3 Ammunition}
\item \textit{Ammunition of Slaying}
\item \textit{Amulet of Health}
\item \textit{Amulet of Proof Against Detection and Location}
\item \textit{Amulet of the Planes}
\item \textit{Animated Shield}
\item \textit{Apparatus of Kwalish}
\item \textit{+1 Armor}
\item \textit{+2 Armor}
\item \textit{+3 Armor}
\item \textit{Armor of Gleaming}
\item \textit{Armor of Invulnerability}
\item \textit{Armor of Resistance}
\item \textit{Armor of Vulnerability}
\item \textit{Arrow-Catching Shield}
\item \textit{Axe of the Dwarvish Lords}
\item \textit{Baba Yaga's Dancing Broom}
\item \textit{Bag of Beans}
\item \textit{Bag of Devouring}
\item \textit{Bag of Holding}
\item \textit{Bag of Tricks}
\item \textit{Bead of Force}
\item \textit{Bead of Nourishment}
\item \textit{Bead of Refreshment}
\item \textit{Belt of Dwarvenkind}
\item \textit{Belt of Giant Strength}
\item \textit{Berserker Axe}
\item \textit{Blackrazor}
\item \textit{Book of Exalted Deeds}
\item \textit{Book of Vile Darkness}
\item \textit{Boots of Elvenkind}
\item \textit{Boots of False Tracks}
\item \textit{Boots of Levitation}
\item \textit{Boots of Speed}
\item \textit{Boots of Striding and Springing}
\item \textit{Boots of the Winterlands}
\item \textit{Bowl of Commanding Water Elementals}
\item \textit{Bracers of Archery}
\item \textit{Bracers of Defense}
\item \textit{Brazier of Commanding Fire Elementals}
\item \textit{Brooch of Shielding}
\item \textit{Broom of Flying}
\item \textit{Candle of Invocation}
\item \textit{Candle of the Deep}
\item \textit{Cape of the Mountebank}
\item \textit{Cap of Water Breathing}
\item \textit{Carpet of Flying}
\item \textit{Cast-Off Armor}
\item \textit{Cauldron of Rebirth}
\item \textit{Censer of Controlling Air Elementals}
\item \textit{Charlatan's Die}
\item \textit{Chime of Opening}
\item \textit{Circlet of Blasting}
\item \textit{Cloak of Arachnida}
\item \textit{Cloak of Billowing}
\item \textit{Cloak of Displacement}
\item \textit{Cloak of Elvenkind}
\item \textit{Cloak of Invisibility}
\item \textit{Cloak of Many Fashions}
\item \textit{Cloak of Protection}
\item \textit{Cloak of the Bat}
\item \textit{Cloak of the Manta Ray}
\item \textit{Clockwork Amulet}
\item \textit{Clothes of Mending}
\item \textit{Crystal Ball}
\item \textit{Cube of Force}
\item \textit{Cube of Summoning}
\item \textit{Cubic Gate}
\item \textit{Daern's Instant Fortress}
\item \textit{Dagger of Venom}
\item \textit{Dancing Sword}
\item \textit{Dark Shard Amulet}
\item \textit{Decanter of Endless Water}
\item \textit{Deck of Illusions}
\item \textit{Deck of Many Things}
\item \textit{Defender}
\item \textit{Demon Armor}
\item \textit{Demonomicon of Iggwilv}
\item \textit{Dimensional Shackles}
\item \textit{Dragon Scale Mail}
\item \textit{Dragon Slayer}
\item \textit{Dread Helm}
\item \textit{Driftglobe}
\item \textit{Dust of Disappearance}
\item \textit{Dust of Dryness}
\item \textit{Dust of Sneezing and Choking}
\item \textit{Dwarven Plate}
\item \textit{Dwarven Thrower}
\item \textit{Ear Horn of Hearing}
\item \textit{Efreeti Bottle}
\item \textit{Efreeti Chain}
\item \textit{Elemental Gem}
\item \textit{Elixir of Health}
\item \textit{Elven Chain}
\item \textit{Enduring Spellbook}
\item \textit{Energy Bow}
\item \textit{Enspelled Armor}
\item \textit{Enspelled Staff}
\item \textit{Enspelled Weapon}
\item \textit{Ersatz Eye}
\item \textit{Eversmoking Bottle}
\item \textit{Executioner's Axe}
\item \textit{Eye and Hand of Vecna}
\item \textit{Eyes of Charming}
\item \textit{Eyes of Minute Seeing}
\item \textit{Eyes of the Eagle}
\item \textit{Figurine of Wondrous Power}
\item \textit{Flame Tongue}
\item \textit{Folding Boat}
\item \textit{Frost Brand}
\item \textit{Gauntlets of Ogre Power}
\item \textit{Gem of Brightness}
\item \textit{Gem of Seeing}
\item \textit{Giant Slayer}
\item \textit{Glamoured Studded Leather}
\item \textit{Gloves of Missile Snaring}
\item \textit{Gloves of Swimming and Climbing}
\item \textit{Gloves of Thievery}
\item \textit{Goggles of Night}
\item \textit{Hag Eye}
\item \textit{Hammer of Thunderbolts}
\item \textit{Hat of Disguise}
\item \textit{Hat of Many Spells}
\item \textit{Hat of Vermin}
\item \textit{Hat of Wizardry}
\item \textit{Headband of Intellect}
\item \textit{Helm of Brilliance}
\item \textit{Helm of Comprehending Languages}
\item \textit{Helm of Telepathy}
\item \textit{Helm of Teleportation}
\item \textit{Heward's Handy Haversack}
\item \textit{Heward's Handy Spice Pouch}
\item \textit{Holy Avenger}
\item \textit{Horn of Blasting}
\item \textit{Horn of Silent Alarm}
\item \textit{Horn of Valhalla}
\item \textit{Horseshoes of a Zephyr}
\item \textit{Horseshoes of Speed}
\item \textit{Immovable Rod}
\item \textit{Instrument of Illusions}
\item \textit{Instrument of Scribing}
\item \textit{Instrument of the Bards}
\item \textit{Ioun Stone}
\item \textit{Iron Bands of Bilarro}
\item \textit{Iron Flask}
\item \textit{Javelin of Lightning}
\item \textit{Keoghtom's Ointment}
\item \textit{Lantern of Revealing}
\item \textit{Lock of Trickery}
\item \textit{Luck Blade}
\item \textit{Lute of Thunderous Thumping}
\item \textit{Mace of Disruption}
\item \textit{Mace of Smiting}
\item \textit{Mace of Terror}
\item \textit{Mantle of Spell Resistance}
\item \textit{Manual of Bodily Health}
\item \textit{Manual of Gainful Exercise}
\item \textit{Manual of Golems}
\item \textit{Manual of Quickness of Action}
\item \textit{Mariner's Armor}
\item \textit{Medallion of Thoughts}
\item \textit{Mirror of Life Trapping}
\item \textit{Mithral Armor}
\item \textit{Moonblade}
\item \textit{Moon-Touched Sword}
\item \textit{Mystery Key}
\item \textit{Nature's Mantle}
\item \textit{Necklace of Adaptation}
\item \textit{Necklace of Fireballs}
\item \textit{Necklace of Prayer Beads}
\item \textit{Nine Lives Stealer}
\item \textit{Nolzur's Marvelous Pigments}
\item \textit{Oathbow}
\item \textit{Oil of Etherealness}
\item \textit{Oil of Sharpness}
\item \textit{Oil of Slipperiness}
\item \textit{Orb of Direction}
\item \textit{Orb of Dragonkind}
\item \textit{Orb of Time}
\item \textit{Pearl of Power}
\item \textit{Perfume of Bewitching}
\item \textit{Periapt of Health}
\item \textit{Periapt of Proof Against Poison}
\item \textit{Periapt of Wound Closure}
\item \textit{Philter of Love}
\item \textit{Pipe of Smoke Monsters}
\item \textit{Pipes of Haunting}
\item \textit{Pipes of the Sewers}
\item \textit{Plate Armor of Etherealness}
\item \textit{Pole of Angling}
\item \textit{Pole of Collapsing}
\item \textit{Portable Hole}
\item \textit{Potion of Animal Friendship}
\item \textit{Potion of Clairvoyance}
\item \textit{Potion of Climbing}
\item \textit{Potion of Comprehension}
\item \textit{Potion of Diminution}
\item \textit{Potion of Fire Breath}
\item \textit{Potion of Flying}
\item \textit{Potion of Gaseous Form}
\item \textit{Potion of Giant Strength}
\item \textit{Potion of Greater Invisibility}
\item \textit{Potion of Growth}
\item \textit{Potion of Healing}
\item \textit{Potion of Heroism}
\item \textit{Potion of Invisibility}
\item \textit{Potion of Invulnerability}
\item \textit{Potion of Longevity}
\item \textit{Potion of Mind Reading}
\item \textit{Potion of Poison}
\item \textit{Potion of Pugilism}
\item \textit{Potion of Resistance}
\item \textit{Potion of Speed}
\item \textit{Potion of Vitality}
\item \textit{Potion of Water Breathing}
\item \textit{Pot of Awakening}
\item \textit{Prosthetic Limb}
\item \textit{Quaal's Feather Token}
\item \textit{Quarterstaff of the Acrobat}
\item \textit{Quiver of Ehlonna}
\item \textit{Ring of Animal Influence}
\item \textit{Ring of Djinni Summoning}
\item \textit{Ring of Elemental Command}
\item \textit{Ring of Evasion}
\item \textit{Ring of Feather Falling}
\item \textit{Ring of Free Action}
\item \textit{Ring of Invisibility}
\item \textit{Ring of Jumping}
\item \textit{Ring of Mind Shielding}
\item \textit{Ring of Protection}
\item \textit{Ring of Regeneration}
\item \textit{Ring of Resistance}
\item \textit{Ring of Shooting Stars}
\item \textit{Ring of Spell Storing}
\item \textit{Ring of Spell Turning}
\item \textit{Ring of Swimming}
\item \textit{Ring of Telekinesis}
\item \textit{Ring of the Ram}
\item \textit{Ring of Three Wishes}
\item \textit{Ring of Warmth}
\item \textit{Ring of Water Walking}
\item \textit{Ring of X-Ray Vision}
\item \textit{Rival Coin}
\item \textit{Robe of Eyes}
\item \textit{Robe of Scintillating Colors}
\item \textit{Robe of Stars}
\item \textit{Robe of the Archmagi}
\item \textit{Robe of Useful Items}
\item \textit{Rod of Absorption}
\item \textit{Rod of Alertness}
\item \textit{Rod of Lordly Might}
\item \textit{Rod of Resurrection}
\item \textit{Rod of Rulership}
\item \textit{Rod of Security}
\item \textit{Rod of the Pact Keeper}
\item \textit{Rope of Climbing}
\item \textit{Rope of Entanglement}
\item \textit{Rope of Mending}
\item \textit{Ruby of the War Mage}
\item \textit{Saddle of the Cavalier}
\item \textit{Scarab of Protection}
\item \textit{Scimitar of Speed}
\item \textit{Scroll of Protection}
\item \textit{Scroll of Titan Summoning}
\item \textit{Sending Stones}
\item \textit{Sentinel Shield}
\item \textit{+1 Shield}
\item \textit{+2 Shield}
\item \textit{+3 Shield}
\item \textit{Shield of Expression}
\item \textit{Shield of Missile Attraction}
\item \textit{Shield of the Cavalier}
\item \textit{Silvered Weapon}
\item \textit{Slippers of Spider Climbing}
\item \textit{Smoldering Armor}
\item \textit{Sovereign Glue}
\item \textit{Spellguard Shield}
\item \textit{Spell Scroll}
\item \textit{Sphere of Annihilation}
\item \textit{Spirit Board}
\item \textit{Staff of Adornment}
\item \textit{Staff of Birdcalls}
\item \textit{Staff of Charming}
\item \textit{Staff of Fire}
\item \textit{Staff of Flowers}
\item \textit{Staff of Frost}
\item \textit{Staff of Healing}
\item \textit{Staff of Power}
\item \textit{Staff of Striking}
\item \textit{Staff of Swarming Insects}
\item \textit{Staff of the Adder}
\item \textit{Staff of the Magi}
\item \textit{Staff of the Python}
\item \textit{Staff of the Woodlands}
\item \textit{Staff of Thunder and Lightning}
\item \textit{Staff of Withering}
\item \textit{Stone of Controlling Earth Elementals}
\item \textit{Stone of Good Luck}
\item \textit{Sun Blade}
\item \textit{Sword of Answering}
\item \textit{Sword of Kas}
\item \textit{Sword of Life Stealing}
\item \textit{Sword of Sharpness}
\item \textit{Sword of Vengeance}
\item \textit{Sword of Wounding}
\item \textit{Sylvan Talon}
\item \textit{Talisman of Pure Good}
\item \textit{Talisman of the Sphere}
\item \textit{Talisman of Ultimate Evil}
\item \textit{Talking Doll}
\item \textit{Tankard of Sobriety}
\item \textit{Tentacle Rod}
\item \textit{Thunderous Greatclub}
\item \textit{Tome of Clear Thought}
\item \textit{Tome of Leadership and Influence}
\item \textit{Tome of the Stilled Tongue}
\item \textit{Tome of Understanding}
\item \textit{Trident of Fish Command}
\item \textit{Universal Solvent}
\item \textit{Veteran's Cane}
\item \textit{Vicious Weapon}
\item \textit{Vorpal Sword}
\item \textit{Walloping Ammunition}
\item \textit{Wand of Binding}
\item \textit{Wand of Conducting}
\item \textit{Wand of Enemy Detection}
\item \textit{Wand of Fear}
\item \textit{Wand of Fireballs}
\item \textit{Wand of Lightning Bolts}
\item \textit{Wand of Magic Detection}
\item \textit{Wand of Magic Missiles}
\item \textit{Wand of Orcus}
\item \textit{Wand of Paralysis}
\item \textit{Wand of Polymorph}
\item \textit{Wand of Pyrotechnics}
\item \textit{Wand of Secrets}
\item \textit{Wand of the War Mage}
\item \textit{Wand of Web}
\item \textit{Wand of Wonder}
\item \textit{Wave}
\item \textit{+1 Weapon}
\item \textit{+2 Weapon}
\item \textit{+3 Weapon}
\item \textit{Weapon of Warning}
\item \textit{Well of Many Worlds}
\item \textit{Whelm}
\item \textit{Wind Fan}
\item \textit{Winged Boots}
\item \textit{Wings of Flying}
\item \textit{Wraps of Unarmed Power}
\end{multicols}

\end{itemize}

\section{Random Magic Items}
Use the tables in this section to randomly determine magic items the characters find in your adventures. The tables are sorted first by treasure theme (as described in the \textit{Monster Manual}), and then by item rarity. If an item isn't associated with a theme, roll 1d4 to decide which table to roll on next: on a \textbf{1}, roll on the \textit{Arcana tables}; on a \textbf{2}, \textit{Armaments}; on a \textbf{3}, \textit{Implements}; and on a \textbf{4}, \textit{Relics}.

\subsection{Arcana Tables}
\begin{DndTable}[header=Arcana - Common]{cX}

1d100 & Item\\
01–02 & \textit{Bead of Nourishment}\\
03–04 & \textit{Bead of Refreshment}\\
05–07 & \textit{Candle of the Deep}\\
08–10 & \textit{Cloak of Billowing}\\
11–13 & \textit{Cloak of Many Fashions}\\
14–15 & \textit{Clothes of Mending}\\
16–17 & \textit{Dark Shard Amulet}\\
18–19 & \textit{Enduring Spellbook}\\
20–21 & \textit{Ersatz Eye}\\
22–24 & \textit{Hat of Vermin}\\
25–27 & \textit{Hat of Wizardry}\\
28–29 & \textit{Heward's Handy Spice Pouch}\\
30–31 & \textit{Horn of Silent Alarm}\\
32–33 & \textit{Instrument of Illusions}\\
34–35 & \textit{Instrument of Scribing}\\
36–37 & \textit{Lock of Trickery}\\
38–40 & \textit{Mystery Key}\\
41–42 & \textit{Orb of Direction}\\
43–44 & \textit{Orb of Time}\\
45–46 & \textit{Perfume of Bewitching}\\
47–49 & \textit{Pipe of Smoke Monsters}\\
50–52 & \textit{Potion of Climbing}\\
53–55 & \textit{Potion of Comprehension}\\
56–58 & \textit{Pot of Awakening}\\
59–60 & \textit{Prosthetic Limb}\\
61–64 & \textit{Rival Coin}\\
65–66 & \textit{Rope of Mending}\\
67–68 & \textit{Ruby of the War Mage}\\
69–82 & \textit{Spell Scroll} (Spell Scroll (Cantrip) or Spell Scroll (Level 1) spell)\\
83–84 & \textit{Staff of Adornment}\\
85–86 & \textit{Staff of Birdcalls}\\
87–89 & \textit{Staff of Flowers}\\
90–92 & \textit{Talking Doll}\\
93–94 & \textit{Tankard of Sobriety}\\
95–97 & \textit{Wand of Conducting}\\
98–00 & \textit{Wand of Pyrotechnics}\\
\end{DndTable}

\begin{DndTable}[header=Arcana - Uncommon]{cX}

1d100 & Item\\
01 & \textit{Amulet of Proof against Detection and Location}\\
02 & \textit{Baba Yaga's Dancing Broom}\\
03–05 & \textit{Bag of Holding}\\
06–07 & \textit{Bag of Tricks}\\
08 & \textit{Brooch of Shielding}\\
09 & \textit{Broom of Flying}\\
10 & \textit{Cap of Water Breathing}\\
11 & \textit{Circlet of Blasting}\\
12–13 & \textit{Cloak of Protection}\\
14 & \textit{Cloak of the Manta Ray}\\
15 & \textit{Decanter of Endless Water}\\
16 & \textit{Deck of Illusions}\\
17–18 & \textit{Driftglobe}\\
19–20 & \textit{Dust of Disappearance}\\
21 & \textit{Dust of Dryness}\\
22 & \textit{Dust of Sneezing and Choking}\\
23–24 & \textit{Elemental Gem}\\
25 & \textit{Enspelled Staff} (cantrip or level 1 spell)\\
26 & \textit{Eversmoking Bottle}\\
27 & \textit{Eyes of Charming}\\
28 & \textit{Eyes of Minute Seeing}\\
29–30 & Figurine of Wondrous Power, Silver Raven\\
31 & \textit{Gem of Brightness}\\
32 & \textit{Hag Eye}\\
33 & \textit{Hat of Disguise}\\
34 & \textit{Headband of Intellect}\\
35 & \textit{Helm of Comprehending Languages}\\
36 & \textit{Helm of Telepathy}\\
37 & \textit{Immovable Rod}\\
38 & \textit{Lantern of Revealing}\\
39 & \textit{Medallion of Thoughts}\\
40 & \textit{Mithral Armor}\\
41–42 & \textit{Necklace of Adaptation}\\
43 & \textit{Oil of Slipperiness}\\
44 & \textit{Pearl of Power}\\
45 & \textit{Periapt of Health}\\
46–47 & \textit{Philter of Love}\\
48–49 & \textit{Potion of Animal Friendship}\\
50–51 & \textit{Potion of Fire Breath}\\
52–53 & Potion of Hill Giant Strength\\
54–55 & \textit{Potion of Growth}\\
56–57 & \textit{Potion of Poison}\\
58–59 & \textit{Potion of Resistance}\\
60–61 & \textit{Potion of Water Breathing}\\
62 & \textit{Quaal's Feather Token} (Quaal's Feather Token, Anchor, Quaal's Feather Token, Fan, or Quaal's Feather Token, Tree)\\
63 & \textit{Ring of Mind Shielding}\\
64–65 & \textit{Robe of Useful Items}\\
66–67 & \textit{Rod of the Pact Keeper}\\
68–69 & \textit{Rope of Climbing}\\
70 & \textit{Saddle of the Cavalier}\\
71–72 & \textit{Sending Stones}\\
73–74 & \textit{Slippers of Spider Climbing}\\
75–82 & \textit{Spell Scroll} (Spell Scroll (Level 2) or Spell Scroll (Level 3) spell)\\
83 & \textit{Staff of the Adder}\\
84 & \textit{Staff of the Python}\\
85–88 & \textit{Wand of Magic Detection}\\
89–91 & \textit{Wand of Magic Missiles}\\
92–93 & \textit{Wand of Secrets}\\
94–95 & +1 Wand of the War Mage\\
96–97 & \textit{Wand of Web}\\
98–99 & \textit{Wind Fan}\\
00 & \textit{Winged Boots}\\
\end{DndTable}

\begin{DndTable}[header=Arcana - Rare]{cX}

1d100 & Item\\
01 & \textit{Bag of Beans}\\
02–03 & \textit{Bead of Force}\\
04 & \textit{Bowl of Commanding Water Elementals}\\
05–06 & \textit{Bracers of Defense}\\
07 & \textit{Brazier of Commanding Fire Elementals}\\
08–09 & \textit{Cape of the Mountebank}\\
10 & \textit{Censer of Controlling Air Elementals}\\
11–12 & \textit{Chime of Opening}\\
13–14 & \textit{Cloak of Displacement}\\
15–16 & \textit{Cloak of the Bat}\\
17 & \textit{Cube of Force}\\
18 & \textit{Cube of Summoning}\\
19 & \textit{Daern's Instant Fortress}\\
20–21 & \textit{Enspelled Staff} (level 2 or 3 spell)\\
22–23 & \textit{Figurine of Wondrous Power} (Figurine of Wondrous Power, Bronze Griffon, Figurine of Wondrous Power, Ebony Fly, Figurine of Wondrous Power, Golden Lions, Figurine of Wondrous Power, Ivory Goats, Figurine of Wondrous Power, Marble Elephant, Figurine of Wondrous Power, Onyx Dog, or Figurine of Wondrous Power, Serpentine Owl)\\
24–25 & \textit{Folding Boat}\\
26–27 & \textit{Gem of Seeing}\\
28 & \textit{Helm of Teleportation}\\
29–30 & \textit{Heward's Handy Haversack}\\
31–32 & \textit{Horseshoes of Speed}\\
33–34 & Ioun Stone, Reserve\\
35 & \textit{Iron Bands of Bilarro}\\
36 & \textit{Mantle of Spell Resistance}\\
37–38 & \textit{Necklace of Fireballs}\\
39 & \textit{Oil of Etherealness}\\
40 & \textit{Portable Hole}\\
41–42 & \textit{Potion of Clairvoyance}\\
43–44 & \textit{Potion of Diminution}\\
45–46 & \textit{Potion of Gaseous Form}\\
47 & Potion of Fire Giant Strength\\
48–49 & \textit{Potion of Giant Strength} (Potion of Frost Giant Strength or Potion of Stone Giant Strength)\\
50–51 & \textit{Potion of Heroism}\\
52–53 & \textit{Potion of Invisibility}\\
54–55 & \textit{Potion of Invulnerability}\\
56–57 & \textit{Potion of Mind Reading}\\
58–59 & \textit{Quaal's Feather Token} (Quaal's Feather Token, Bird, Quaal's Feather Token, Swan Boat, or Quaal's Feather Token, Whip)\\
60–61 & \textit{Ring of Feather Falling}\\
62 & \textit{Ring of Spell Storing}\\
63 & \textit{Ring of X-ray Vision}\\
64–65 & \textit{Robe of Eyes}\\
66 & \textit{Rod of Rulership}\\
67–68 & +2 Rod of the Pact Keeper\\
69–70 & \textit{Scroll of Protection}\\
71–75 & \textit{Spell Scroll} (Spell Scroll (Level 4) or Spell Scroll (Level 5) spell)\\
76–77 & \textit{Staff of Charming}\\
78–79 & \textit{Staff of Swarming Insects}\\
80–81 & \textit{Staff of Withering}\\
82 & \textit{Stone of Controlling Earth Elementals}\\
83–84 & \textit{Wand of Binding}\\
85–86 & \textit{Wand of Fear}\\
87–90 & \textit{Wand of Fireballs}\\
91–94 & \textit{Wand of Lightning Bolts}\\
95–96 & +2 Wand of the War Mage\\
97–98 & \textit{Wand of Wonder}\\
99–00 & \textit{Wings of Flying}\\
\end{DndTable}

\begin{DndTable}[header=Arcana - Very Rare]{cX}

1d100 & Item\\
01–02 & \textit{Amulet of the Planes}\\
03–04 & \textit{Bag of Devouring}\\
05–06 & \textit{Carpet of Flying}\\
07–08 & \textit{Cauldron of Rebirth}\\
09–10 & \textit{Cloak of Arachnida}\\
11–12 & \textit{Crystal Ball}\\
13 & \textit{Dancing Sword}\\
14 & \textit{Efreeti Bottle}\\
15–16 & \textit{Enspelled Staff} (level 4 or 5 spell)\\
17–18 & Figurine of Wondrous Power, Obsidian Steed\\
19–20 & \textit{Hat of Many Spells}\\
21–22 & \textit{Helm of Brilliance}\\
23–24 & \textit{Horseshoes of a Zephyr}\\
25–26 & \textit{Ioun Stone} (Ioun Stone, Absorption, Ioun Stone, Fortitude, Ioun Stone, Intellect, or Ioun Stone, Leadership)\\
27–28 & \textit{Manual of Golems} (Manual of Clay Golems, Manual of Flesh Golems, Manual of Iron Golems, or Manual of Stone Golems)\\
29 & \textit{Mirror of Life Trapping}\\
30–31 & \textit{Nolzur's Marvelous Pigments}\\
32–34 & \textit{Oil of Sharpness}\\
35–38 & \textit{Potion of Flying}\\
39–42 & Potion of Cloud Giant Strength\\
43–46 & \textit{Potion of Greater Invisibility}\\
47–49 & \textit{Potion of Longevity}\\
50–53 & \textit{Potion of Speed}\\
54–57 & \textit{Potion of Vitality}\\
58–59 & \textit{Ring of Regeneration}\\
60–61 & \textit{Ring of Shooting Stars}\\
62–63 & \textit{Ring of Telekinesis}\\
64–65 & \textit{Robe of Scintillating Colors}\\
66–67 & \textit{Robe of Stars}\\
68–69 & \textit{Rod of Absorption}\\
70–71 & \textit{Rod of Security}\\
72–73 & +3 Rod of the Pact Keeper\\
74–85 & \textit{Spell Scroll} (Spell Scroll (Level 6), Spell Scroll (Level 7), or Spell Scroll (Level 8) spell)\\
86–87 & \textit{Staff of Fire}\\
88–89 & \textit{Staff of Frost}\\
90 & \textit{Staff of Power}\\
91–92 & \textit{Staff of Thunder and Lightning}\\
93–94 & \textit{Tome of Clear Thought}\\
95–97 & \textit{Wand of Polymorph}\\
98–00 & +3 Wand of the War Mage\\
\end{DndTable}

\begin{DndTable}[header=Arcana - Legendary]{cX}

1d100 & Item\\
01–04 & \textit{Apparatus of Kwalish}\\
05–08 & \textit{Cloak of Invisibility}\\
09–12 & \textit{Crystal Ball of Mind Reading}\\
13–16 & \textit{Crystal Ball of Telepathy}\\
17–20 & \textit{Crystal Ball of True Seeing}\\
21–22 & \textit{Cubic Gate}\\
23 & \textit{Deck of Many Things}\\
24–27 & \textit{Enspelled Staff} (level 6, 7, or 8 spell)\\
28–31 & \textit{Ioun Stone} (Ioun Stone, Greater Absorption, Ioun Stone, Mastery, or Ioun Stone, Regeneration)\\
32–33 & \textit{Iron Flask}\\
34–41 & Potion of Storm Giant Strength\\
42–45 & \textit{Ring of Djinni Summoning}\\
46–49 & \textit{Ring of Elemental Command} (Ring of Elemental Command (Air), Ring of Elemental Command (Earth), Ring of Elemental Command (Fire), or Ring of Elemental Command (Water))\\
50–53 & \textit{Ring of Invisibility}\\
54–57 & \textit{Ring of Spell Turning}\\
58 & \textit{Ring of Three Wishes}\\
59 & \textit{Robe of the Archmagi}\\
60–61 & \textit{Scroll of Titan Summoning}\\
62–65 & \textit{Sovereign Glue}\\
66–83 & Spell Scroll (Level 9) spell)\\
84 & \textit{Sphere of Annihilation}\\
85 & \textit{Staff of the Magi}\\
86–88 & \textit{Talisman of the Sphere}\\
89–92 & \textit{Tome of the Stilled Tongue}\\
93–96 & \textit{Universal Solvent}\\
97–00 & \textit{Well of Many Worlds}\\
\end{DndTable}

\subsection{Armaments Tables}
\begin{DndTable}[header=Armaments - Common]{cX}

1d100 & Item\\
01–10 & \textit{Armor of Gleaming}\\
11–20 & \textit{Cast-Off Armor}\\
21–30 & \textit{Dread Helm}\\
31–40 & \textit{Moon-Touched Sword}\\
41–50 & \textit{Shield of Expression}\\
51–60 & \textit{Silvered Weapon}\\
61–70 & \textit{Smoldering Armor}\\
71–80 & \textit{Sylvan Talon}\\
81–90 & \textit{Veteran's Cane}\\
91–00 & \textit{Walloping Ammunition}\\
\end{DndTable}

\begin{DndTable}[header=Armaments - Uncommon]{cX}

1d100 & Item\\
01–04 & \textit{Adamantine Armor}\\
05–08 & \textit{Adamantine Weapon}\\
09–13 & +1 Ammunition\\
14–18 & \textit{Bracers of Archery}\\
19–23 & \textit{Enspelled Armor} (cantrip or level 1 spell)\\
24–28 & \textit{Enspelled Weapon} (cantrip or level 1 spell)\\
29–33 & \textit{Gauntlets of Ogre Power}\\
34–38 & \textit{Javelin of Lightning}\\
39–43 & \textit{Mariner's Armor}\\
44–48 & \textit{Mithral Armor}\\
49–53 & Potion of Hill Giant Strength\\
54–58 & \textit{Potion of Pugilism}\\
59–62 & \textit{Quiver of Ehlonna}\\
63–66 & \textit{Saddle of the Cavalier}\\
67–71 & \textit{Sentinel Shield}\\
72–76 & +1 Shield\\
77–81 & \textit{Sword of Vengeance}\\
82–85 & \textit{Trident of Fish Command}\\
86–90 & +1 Weapon\\
91–95 & \textit{Weapon of Warning}\\
96–00 & +1 Wraps of Unarmed Power\\
\end{DndTable}

\begin{DndTable}[header=Armaments - Rare]{cX}

1d100 & Item\\
01–03 & +2 Ammunition\\
04–07 & +1 Armor\\
08–10 & \textit{Armor of Resistance}\\
11–13 & \textit{Armor of Vulnerability}\\
14–15 & \textit{Arrow-Catching Shield}\\
16–18 & Belt of Hill Giant Strength\\
19–20 & \textit{Berserker Axe}\\
21–22 & \textit{Daern's Instant Fortress}\\
23–25 & \textit{Dagger of Venom}\\
26–28 & \textit{Dragon Slayer}\\
29–31 & \textit{Elven Chain}\\
32–34 & \textit{Enspelled Armor} (level 2 or 3 spell)\\
35–37 & \textit{Enspelled Weapon} (level 2 or 3 spell)\\
38–40 & \textit{Flame Tongue}\\
41–43 & \textit{Giant Slayer}\\
44–46 & \textit{Horn of Blasting}\\
47–48 & \textit{Horn of Valhalla} (Horn of Valhalla, Brass or Horn of Valhalla, Silver)\\
49–51 & \textit{Ioun Stone} (Ioun Stone, Protection)\\
52–54 & \textit{Mace of Disruption}\\
55–57 & \textit{Mace of Smiting}\\
58–60 & \textit{Mace of Terror}\\
61–63 & Potion of Fire Giant Strength\\
64–66 & \textit{Potion of Giant Strength} (Potion of Frost Giant Strength or Potion of Stone Giant Strength)\\
67–69 & \textit{Potion of Heroism}\\
70–72 & \textit{Potion of Invulnerability}\\
73–75 & \textit{Ring of Protection}\\
76–78 & \textit{Ring of the Ram}\\
79–81 & +2 Shield\\
82–84 & \textit{Shield of Missile Attraction}\\
85–86 & \textit{Sun Blade}\\
87–88 & \textit{Sword of Life Stealing}\\
89–90 & \textit{Sword of Wounding}\\
91–92 & \textit{Tentacle Rod}\\
93–94 & \textit{Vicious Weapon}\\
95–97 & +2 Weapon\\
98–00 & +2 Wraps of Unarmed Power\\
\end{DndTable}

\begin{DndTable}[header=Armaments - Very Rare]{cX}

1d100 & Item\\
01–03 & +3 Ammunition\\
04–06 & \textit{Ammunition of Slaying}\\
07–09 & \textit{Animated Shield}\\
10–12 & +2 Armor\\
13–14 & Belt of Fire Giant Strength\\
15–17 & \textit{Belt of Giant Strength} (Belt of Frost Giant Strength or Belt of Stone Giant Strength)\\
18–19 & \textit{Dancing Sword}\\
20–22 & \textit{Demon Armor}\\
23–25 & \textit{Dragon Scale Mail}\\
26–28 & \textit{Dwarven Plate}\\
29–31 & \textit{Dwarven Thrower}\\
32–34 & \textit{Energy Bow}\\
35–37 & \textit{Enspelled Armor} (level 4 or 5 spell)\\
38–40 & \textit{Enspelled Weapon} (level 4 or 5 spell)\\
41–43 & \textit{Executioner's Axe}\\
44–46 & \textit{Frost Brand}\\
47–49 & Horn of Valhalla, Bronze\\
50–52 & \textit{Ioun Stone} (Ioun Stone, Strength)\\
53–55 & \textit{Lute of Thunderous Thumping}\\
56–58 & \textit{Manual of Gainful Exercise}\\
59–61 & \textit{Nine Lives Stealer}\\
62–64 & \textit{Oathbow}\\
65–68 & \textit{Oil of Sharpness}\\
69–72 & Potion of Cloud Giant Strength\\
73–75 & \textit{Quarterstaff of the Acrobat}\\
76–78 & \textit{Scimitar of Speed}\\
79–82 & +3 Shield\\
83–85 & \textit{Shield of the Cavalier}\\
86–88 & \textit{Spellguard Shield}\\
89–91 & \textit{Sword of Sharpness}\\
92–94 & \textit{Thunderous Greatclub}\\
95–97 & +3 Weapon\\
98–00 & +3 Wraps of Unarmed Power\\
\end{DndTable}

\begin{DndTable}[header=Armaments - Legendary]{cX}

1d100 & Item\\
01–06 & +3 Armor\\
07–12 & \textit{Armor of Invulnerability}\\
13–18 & Belt of Cloud Giant Strength\\
19–21 & Belt of Storm Giant Strength\\
22–27 & \textit{Defender}\\
28–33 & \textit{Efreeti Chain}\\
34–39 & \textit{Enspelled Armor} (level 6, 7, or 8 spell)\\
40–45 & \textit{Enspelled Weapon} (level 6, 7, or 8 spell)\\
46–51 & \textit{Hammer of Thunderbolts}\\
52–56 & \textit{Holy Avenger}\\
57–62 & Horn of Valhalla, Iron\\
63–68 & \textit{Luck Blade}\\
69–72 & \textit{Moonblade}\\
73–78 & \textit{Plate Armor of Etherealness}\\
79–87 & Potion of Storm Giant Strength\\
88–90 & \textit{Rod of Lordly Might}\\
91–95 & \textit{Sword of Answering}\\
96–00 & \textit{Vorpal Sword}\\
\end{DndTable}

\subsection{Implements Tables}
\begin{DndTable}[header=Implements - Common]{cX}

1d100 & Item\\
01–02 & \textit{Bead of Nourishment}\\
03–04 & \textit{Bead of Refreshment}\\
05–06 & \textit{Boots of False Tracks}\\
07–08 & \textit{Candle of the Deep}\\
09–10 & \textit{Charlatan's Die}\\
11–13 & \textit{Cloak of Many Fashions}\\
14–15 & \textit{Clockwork Amulet}\\
16–17 & \textit{Ear Horn of Hearing}\\
18–19 & \textit{Ersatz Eye}\\
20–21 & \textit{Heward's Handy Spice Pouch}\\
22–23 & \textit{Horn of Silent Alarm}\\
24–25 & \textit{Instrument of Illusions}\\
26–27 & \textit{Instrument of Scribing}\\
28–29 & \textit{Lock of Trickery}\\
30–32 & \textit{Moon-Touched Sword}\\
33–34 & \textit{Mystery Key}\\
35–36 & \textit{Orb of Direction}\\
37–38 & \textit{Orb of Time}\\
39–40 & \textit{Perfume of Bewitching}\\
41–42 & \textit{Pipe of Smoke Monsters}\\
43–44 & \textit{Pole of Angling}\\
45–46 & \textit{Pole of Collapsing}\\
47–52 & \textit{Potion of Climbing}\\
53–58 & \textit{Potion of Comprehension}\\
59–74 & \textit{Potion of Healing}\\
75–76 & \textit{Prosthetic Limb}\\
77–78 & \textit{Rope of Mending}\\
79–80 & \textit{Staff of Birdcalls}\\
81–82 & \textit{Sylvan Talon}\\
83–84 & \textit{Talking Doll}\\
85–86 & \textit{Tankard of Sobriety}\\
87–90 & \textit{Veteran's Cane}\\
91–92 & \textit{Walloping Ammunition}\\
93–94 & \textit{Wand of Conducting}\\
95–97 & \textit{Wand of Enemy Detection}\\
98–00 & \textit{Wand of Pyrotechnics}\\
\end{DndTable}

\begin{DndTable}[header=Implements - Uncommon]{cX}

1d100 & Item\\
01–02 & \textit{Alchemy Jug}\\
03–06 & +1 Ammunition\\
07–10 & \textit{Bag of Holding}\\
11–12 & \textit{Boots of Elvenkind}\\
13–14 & \textit{Boots of Striding and Springing}\\
15–16 & \textit{Boots of the Winterlands}\\
17–18 & \textit{Broom of Flying}\\
19–20 & \textit{Cap of Water Breathing}\\
21–22 & \textit{Cloak of Elvenkind}\\
23–24 & \textit{Cloak of Protection}\\
25–26 & \textit{Cloak of the Manta Ray}\\
27 & \textit{Decanter of Endless Water}\\
28–30 & \textit{Driftglobe}\\
31–32 & \textit{Dust of Disappearance}\\
33–34 & \textit{Dust of Dryness}\\
35–36 & \textit{Dust of Sneezing and Choking}\\
37–38 & \textit{Enspelled Weapon} (cantrip or level 1 spell)\\
39–40 & \textit{Eyes of Minute Seeing}\\
41–42 & \textit{Eyes of the Eagle}\\
43–44 & \textit{Gloves of Missile Snaring}\\
45–46 & \textit{Gloves of Swimming and Climbing}\\
47–48 & \textit{Gloves of Thievery}\\
49–50 & \textit{Goggles of Night}\\
51 & \textit{Hag Eye}\\
52–54 & \textit{Helm of Comprehending Languages}\\
55 & \textit{Immovable Rod}\\
56–57 & \textit{Instrument of the Bards} (Instrument of the Bards, Doss Lute, Instrument of the Bards, Fochlucan Bandore, or Instrument of the Bards, Mac-Fuirmidh Cittern)\\
58–59 & \textit{Lantern of Revealing}\\
60–61 & \textit{Nature's Mantle}\\
62–63 & \textit{Oil of Slipperiness}\\
64–65 & \textit{Pipes of Haunting}\\
66–67 & \textit{Pipes of the Sewers}\\
68–71 & \textit{Potion of Growth}\\
72–80 & Potion of Greater Healing\\
81–84 & \textit{Potion of Water Breathing}\\
85–86 & \textit{Quaal's Feather Token} (Quaal's Feather Token, Anchor, Quaal's Feather Token, Fan, or Quaal's Feather Token, Tree)\\
87–88 & \textit{Ring of Jumping}\\
89–90 & \textit{Ring of Swimming}\\
91–92 & \textit{Ring of Warmth}\\
93–94 & \textit{Robe of Useful Items}\\
95–96 & \textit{Rope of Climbing}\\
97–98 & \textit{Stone of Good Luck}\\
99–00 & \textit{Wand of Secrets}\\
\end{DndTable}

\begin{DndTable}[header=Implements - Rare]{cX}

1d100 & Item\\
01–04 & +2 Ammunition\\
05–08 & \textit{Bag of Beans}\\
09–12 & \textit{Belt of Dwarvenkind}\\
13–16 & \textit{Boots of Levitation}\\
17–20 & \textit{Boots of Speed}\\
21–24 & \textit{Chime of Opening}\\
25–28 & \textit{Dimensional Shackles}\\
29–32 & \textit{Enspelled Weapon} (level 2 or 3 spell)\\
33–36 & \textit{Folding Boat}\\
37–40 & \textit{Glamoured Studded Leather}\\
41–44 & \textit{Heward's Handy Haversack}\\
45–48 & \textit{Horseshoes of Speed}\\
49–52 & \textit{Instrument of the Bards} (Instrument of the Bards, Canaith Mandolin or Instrument of the Bards, Cli Lyre)\\
53–56 & \textit{Ioun Stone} (Ioun Stone, Awareness)\\
57–60 & \textit{Portable Hole}\\
61–64 & \textit{Potion of Diminution}\\
65–68 & \textit{Potion of Gaseous Form}\\
69–76 & Potion of Superior Healing\\
77–80 & \textit{Quaal's Feather Token} (Quaal's Feather Token, Bird, Quaal's Feather Token, Swan Boat, or Quaal's Feather Token, Whip)\\
81–84 & \textit{Ring of Evasion}\\
85–88 & \textit{Ring of Free Action}\\
89–92 & \textit{Rope of Entanglement}\\
93–96 & \textit{Staff of Healing}\\
97–00 & \textit{Wand of Enemy Detection}\\
\end{DndTable}

\begin{DndTable}[header=Implements - Very Rare]{cX}

1d100 & Item\\
01–07 & +3 Ammunition\\
08–14 & \textit{Bag of Devouring}\\
15–21 & \textit{Carpet of Flying}\\
22–28 & \textit{Enspelled Weapon} (level 4 or 5 spell)\\
29–35 & \textit{Horseshoes of a Zephyr}\\
36–42 & Instrument of the Bards, Anstruth Harp\\
43–49 & \textit{Ioun Stone} (Ioun Stone, Agility)\\
50–56 & \textit{Lute of Thunderous Thumping}\\
57–63 & \textit{Manual of Quickness of Action}\\
64–70 & \textit{Nolzur's Marvelous Pigments}\\
71–77 & \textit{Potion of Flying}\\
78–86 & Potion of Supreme Healing\\
87–93 & \textit{Potion of Speed}\\
94–00 & \textit{Tome of Leadership and Influence}\\
\end{DndTable}

\begin{DndTable}[header=Implements - Legendary]{cX}

1d100 & Item\\
01–17 & \textit{Enspelled Weapon} (level 6, 7, or 8 spell)\\
18–34 & \textit{Instrument of the Bards} (Instrument of the Bards, Ollamh Harp)\\
35–54 & \textit{Sovereign Glue}\\
55–70 & \textit{Sphere of Annihilation}\\
71–83 & \textit{Talisman of the Sphere}\\
84–00 & \textit{Universal Solvent}\\
\end{DndTable}

\subsection{Relics Tables}
\begin{DndTable}[header=Relics - Common]{cX}

1d100 & Item\\
01–08 & \textit{Ear Horn of Hearing}\\
09–28 & \textit{Potion of Healing}\\
29–36 & \textit{Pot of Awakening}\\
37–44 & \textit{Ruby of the War Mage}\\
45–52 & \textit{Shield of Expression}\\
53–60 & \textit{Smoldering Armor}\\
61–80 & \textit{Spell Scroll} (Spell Scroll (Cantrip) or Spell Scroll (Level 1) spell)\\
81–90 & \textit{Staff of Adornment}\\
91–00 & \textit{Staff of Flowers}\\
\end{DndTable}

\begin{DndTable}[header=Relics - Uncommon]{cX}

1d100 & Item\\
01–05 & \textit{Enspelled Staff} (cantrip or level 1 spell)\\
06–10 & \textit{Keoghtom's Ointment}\\
11–15 & \textit{Mariner's Armor}\\
16–20 & \textit{Nature's Mantle}\\
21–25 & \textit{Pearl of Power}\\
26–30 & \textit{Periapt of Health}\\
31–35 & \textit{Periapt of Wound Closure}\\
36–40 & \textit{Potion of Animal Friendship}\\
41–55 & Potion of Greater Healing\\
56–60 & \textit{Potion of Resistance}\\
61–65 & \textit{Ring of Water Walking}\\
66–70 & \textit{Sending Stones}\\
71–80 & \textit{Spell Scroll} (Spell Scroll (Level 2) or Spell Scroll (Level 3) spell)\\
81–85 & \textit{Staff of the Adder}\\
86–90 & \textit{Staff of the Python}\\
91–95 & \textit{Wand of Magic Detection}\\
96–00 & +1 Wand of the War Mage\\
\end{DndTable}

\begin{DndTable}[header=Relics - Rare]{cX}

1d100 & Item\\
01–03 & \textit{Amulet of Health}\\
04–07 & +1 Armor\\
08–09 & \textit{Bowl of Commanding Water Elementals}\\
10–11 & \textit{Brazier of Commanding Fire Elementals}\\
12–13 & \textit{Censer of Controlling Air Elementals}\\
14–16 & \textit{Elixir of Health}\\
17–19 & \textit{Enspelled Staff} (level 2 or 3 spell)\\
20–22 & \textit{Horn of Blasting}\\
23–25 & \textit{Horn of Valhalla} (Horn of Valhalla, Brass or Horn of Valhalla, Silver)\\
26–28 & \textit{Ioun Stone} (Ioun Stone, Reserve or Ioun Stone, Sustenance)\\
29–31 & \textit{Mace of Disruption}\\
32–34 & \textit{Mace of Smiting}\\
35–37 & \textit{Mace of Terror}\\
38–40 & \textit{Necklace of Prayer Beads}\\
41–43 & \textit{Periapt of Proof against Poison}\\
44–51 & Potion of Superior Healing\\
52–54 & \textit{Ring of Animal Influence}\\
55–58 & \textit{Ring of Resistance}\\
59–61 & \textit{Ring of Spell Storing}\\
62–65 & \textit{Scroll of Protection}\\
66–73 & \textit{Spell Scroll} (Spell Scroll (Level 4) or Spell Scroll (Level 5) spell)\\
74–76 & \textit{Staff of Charming}\\
77–79 & \textit{Staff of Healing}\\
80–82 & \textit{Staff of Swarming Insects}\\
83–85 & \textit{Staff of the Woodlands}\\
86–88 & \textit{Staff of Withering}\\
89–90 & \textit{Stone of Controlling Earth Elementals}\\
91–93 & \textit{Tentacle Rod}\\
94–96 & \textit{Wand of Paralysis}\\
97–00 & +2 Wand of the War Mage\\
\end{DndTable}

\begin{DndTable}[header=Relics - Very Rare]{cX}

1d100 & Item\\
01–05 & +2 Armor\\
06–10 & \textit{Candle of Invocation}\\
11–15 & \textit{Cauldron of Rebirth}\\
16–20 & \textit{Enspelled Staff} (level 4 or 5 spell)\\
21–25 & Horn of Valhalla, Bronze\\
26–30 & \textit{Ioun Stone} (Ioun Stone, Insight)\\
31–35 & \textit{Manual of Bodily Health}\\
36–43 & Potion of Supreme Healing\\
44–50 & \textit{Potion of Vitality}\\
51–55 & \textit{Rod of Alertness}\\
56–65 & \textit{Spell Scroll} (Spell Scroll (Level 6), Spell Scroll (Level 7), or Spell Scroll (Level 8) spell)\\
66–70 & \textit{Spirit Board}\\
71–75 & \textit{Staff of Fire}\\
76–80 & \textit{Staff of Frost}\\
81–85 & \textit{Staff of Striking}\\
86–90 & \textit{Staff of Thunder and Lightning}\\
91–95 & \textit{Tome of Understanding}\\
96–00 & +3 Wand of the War Mage\\
\end{DndTable}

\begin{DndTable}[header=Relics - Legendary]{cX}

1d100 & Item\\
01–08 & \textit{Armor of Invulnerability}\\
09–18 & +3 Armor\\
19–27 & \textit{Enspelled Staff} (level 6, 7, or 8 spell)\\
28–36 & \textit{Holy Avenger}\\
37–45 & Horn of Valhalla, Iron\\
46–54 & \textit{Rod of Resurrection}\\
55–63 & \textit{Scarab of Protection}\\
64–72 & \textit{Scroll of Titan Summoning}\\
73–80 & \textit{Spell Scroll} (Spell Scroll (Level 9) spell)\\
81–90 & \textit{Talisman of Pure Good}\\
91–00 & \textit{Talisman of Ultimate Evil}\\
\end{DndTable}

\chapter{Bastions}
A Bastion is a location that belongs to a player character: a home, a stronghold, and a place of power that the character develops over the course of a campaign. A Bastion offers a character temporary refuge from the dangerous world of adventuring, and it provides opportunities for a character to craft magic items, conduct research, harvest poisons, build ships, and carry out a range of other activities.

As DM, you decide whether Bastions are available in a campaign. Bastions are best suited to campaigns that allow characters to return to their Bastions during intervals when they're not actively adventuring. Not every character needs to have a Bastion. It's fine for some players in your campaign to opt in to Bastion ownership and others to opt out.

There's no need to choose between going on adventures and commanding a Bastion; a character can do both at once. A Bastion has special facilities that generate benefits, and these facilities can also undertake projects while the character is otherwise occupied.

Most importantly, a Bastion is a creative playground for a player and a shared storytelling space in the campaign. Be as permissive as you can with the stories players tell in their Bastions, but players should know their control might be limited by the campaign's larger story, which you strive to make fun for everyone.

\section{Gaining a Bastion}
If you allow Bastions in your campaign, characters acquire their Bastions when they reach level 5. You and the players can decide together how these Bastions come into being. A character might inherit or receive a parcel of land on which to build their Bastion (see "\textit{Marks of Prestige}" in \textit{chapter 3}), or they might take a preexisting structure and refurbish it. It's fair to assume that work has been going on behind the scenes of the campaign during a character's early adventuring career, so the Bastion is ready when the character reaches level 5.

The shape, style, and function of a character's Bastion are up to the player to determine. For example, a Wizard might build a tower, a Cleric might establish a shrine, a Fighter might build a fortified keep or similar stronghold, and a Rogue might establish a guildhall or lodge. Characters of other classes might choose one of these forms or combine them—a Paladin's Bastion might be similar to a Cleric's shrine but as fortified as a Fighter's stronghold. And multiple characters can combine their Bastions to form a single large structure (see "\textit{Bastion Map}").

Every Bastion has facilities that serve basic needs as well as special ones, such as libraries, menageries, and workshops (see "\textit{Basic Facilities}" and "\textit{Special Facilities}").

\section{Bastion Turns}
As time passes in the campaign, players take Bastion turns to reflect the activity occurring in their Bastions, whether or not the characters are present. On a Bastion turn, a character can issue orders to the special facilities in their Bastion or issue the \textit{Maintain} order to the entire Bastion (see "\textit{Orders}").

\subsection{Frequency of Bastion Turns}
By default, a Bastion turn occurs every 7 days of in-game time. Here are common examples of players taking one or more Bastion turns:


\begin{itemize}
\item The characters are on a long journey away from their Bastions. After the characters have been away for 7 days, you might say, "Time for a Bastion turn. Since you're not there, we'll assume you issue the \textit{Maintain} order for your Bastions." Then roll for events as described in the "\textit{Bastion Events}" section in this chapter.
\item The characters spend 7 days or more in their Bastions between adventures. You might say, "You have six weeks to spend in your Bastions. Go ahead and resolve six Bastion turns."
\item The characters return to their Bastions in the midst of an adventure. You might say, "You have just enough time to take a Bastion turn before you leave again in the morning."
\item The characters are adventuring near their Bastions and resting in their Bastions at night. You might say, "It's been a week since your last Bastion turn, so you can take one now."
\end{itemize}

You can slow the frequency of Bastion turns to better serve the needs of your players and your campaign. For example, if the characters have months between adventures, you can call for a Bastion turn every month instead of every 7 days, so the characters aren't issuing so many orders or reaping too many benefits at once.

\section{Bastion Map}
Encourage players to create floor plans of their characters' Bastions, configuring its facilities as they see fit and using the same techniques you use to create dungeon maps (see "\textit{Dungeons}" in \textit{chapter 3} and "\textit{Adventure Maps}" in \textit{chapter 4}).

In addition to \textit{basic} and \textit{special facilities} (described later in this chapter), a Bastion can have the following features:


\begin{description}
\item[Closets.]
A basic or special facility can have one or more closets, washrooms, or similar enclosures. The closets must be inside the facility and can't increase the facility's area in squares. These features are free.
\item[Corridors, Ramps, and Staircases.]
A facility can contain one or more corridors, ramps, or staircases leading to other facilities in the Bastion. These features are free.
\item[Defensive Walls.]
A character can add defensive walls around their Bastion. A defensive wall is 20 feet high and may include a walkway along its top, with a means to access it (such as a ladder or lift). Each 5-foot square of defensive wall takes 10 days to build and costs 250 GP. If a character's Bastion is completely enclosed by defensive walls and it comes under attack (see "\textit{Bastion Events}" at the end of this chapter), reduce by 2 the number of dice rolled to determine how many Bastion Defenders are lost in the attack.
\item[Doors and Windows.]
Each facility comes with one or more doors and shuttered windows, placed wherever the player sees fit. See "\textit{Doors}" in \textit{chapter 3} for kinds of doors to choose from, including locked doors, secret doors, and portcullises. These features are free.
\end{description}

\subsection{Combining Bastions}
Two or more players can combine their characters' Bastions into a single structure. Doing so doesn't change the number of special facilities each Bastion can have, how those special facilities work, or who issues orders to each Bastion. Each Bastion retains its own hirelings, which can't be sent to or shared with another Bastion. Bastion Defenders are handled differently: if some event deprives one character's Bastion of defenders, another character can apply all or some of those losses to their Bastion instead, provided the two Bastions are combined.

\subsection{Facility Space}
The amount of space in a basic or special facility determines its maximum area in 5-foot squares, as shown in the Facility Space table. A player can configure the squares of a facility as they please on their Bastion's map. The squares can be stacked so that a facility's area is distributed over multiple levels or stories.

\begin{DndTable}[header=Facility Space]{XX}

Space & Maximum Area\\
Cramped & 4 squares\\
Roomy & 16 squares\\
Vast & 36 squares\\
\end{DndTable}

\section{Basic Facilities}
A character's Bastion starts with two free basic facilities, which the character's player chooses from the Basic Facilities list below. One of the chosen facilities is Cramped, and the other is Roomy (see the \textit{Facility Space} table). A Bastion can have more than one of each basic facility.

\subsubsection{Basic Facilities}

\begin{description}
\begin{multicols}{2}
\item Bedroom
\item Dining Room
\item Parlor
\item Courtyard
\item Kitchen
\item Storage
\end{multicols}

\end{description}

A basic facility comes with nonmagical furnishings and decor appropriate for that facility.

Basic facilities don't have any game effects, but they can inspire roleplaying opportunities and enhance a Bastion's verisimilitude. A Bastion with a kitchen is functionally the same as one without, but the former gives you and your players a fun setting to start game sessions, have in-character discussions, or introduce new NPCs.

A character can add new basic facilities or enlarge existing ones by spending money and time, as discussed in the sections that follow. Any number of basic facilities can be added or enlarged at the same time. A character doesn't need to be in their Bastion while basic facilities are being added or enlarged.

\subsection{Adding Basic Facilities}
A character can add a basic facility to their Bastion by spending money and time. The cost of adding a basic facility and the time required depend on the facility's space, as shown on the table below.


\begin{DndTable}{XXX}

Facility Space & Cost & Time Required\\
Cramped & 500 GP & 20 days\\
Roomy & 1,000 GP & 45 days\\
Vast & 3,000 GP & 125 days\\
\end{DndTable}

\subsection{Enlarging Basic Facilities}
There is no in-game benefit to enlarging a basic facility, but a character might enlarge a facility for cosmetic reasons or to increase the Bastion's size.

A character can spend money and time to increase the space of a basic facility in their Bastion by one category, as shown on the table below.


\begin{DndTable}{XXX}

Space Increase & Cost & Time Required\\
Cramped to Roomy & 500 GP & 25 days\\
Roomy to Vast & 2,000 GP & 80 days\\
\end{DndTable}

\section{Special Facilities}
Special facilities are Bastion locations where certain activities yield game benefits. A character's Bastion initially has two special facilities of the character's choice for which they qualify. Each special facility can be chosen only once unless its description says otherwise.

Unlike basic facilities, special facilities can't be bought; a character gains them through level advancement. At level 9, a character gains two additional special facilities of their choice for which they qualify; they gain one additional facility at level 13 and another at level 17. The Special Facility Acquisition table shows the total number of special facilities in a character's Bastion. Each new special facility immediately becomes part of the character's Bastion when the character reaches the level.

Each time a character gains a level, that character can replace one of their Bastion's special facilities with another for which the character qualifies.

\begin{DndTable}[header=Special Facility Acquisition]{cc}

Level & Special Facilities\\
5 & 2\\
9 & 4\\
13 & 5\\
17 & 6\\
\end{DndTable}

\subsection{Requirements}
Each special facility has a level. A character must be that level or higher to gain that facility. A special facility might also have a prerequisite the character must meet to gain that facility. For example, only a character who can use an Arcane Focus or a tool as a Spellcasting Focus can have an \textit{Arcane Study}.

\subsection{Space}
A special facility occupies a certain amount of space (see "\textit{Facility Space}"). A player can configure the squares of a special facility as desired on the Bastion's map. A special facility can be enlarged to grant additional benefits if its description says so.

\subsection{Hirelings}
A special facility comes with one or more hirelings who work in the facility, maintain it, and execute Bastion orders there, as described in the next section. A player can assign names and personalities to hirelings in their character's Bastion using the same tools DMs use to \textit{create NPCs} (see \textit{chapter 3}).

Each special facility in a Bastion generates enough income to pay the salary of its hirelings. Hirelings follow the orders they're given and are loyal to the Bastion's owner.

\subsection{Orders}
On a Bastion turn, a character in their Bastion can issue special orders—called Bastion orders—to one or more of their Bastion's special facilities. A character needn't issue orders to all the special facilities in their Bastion on a given Bastion turn.

The Maintain order is unusual; it is issued to the whole Bastion rather than to one or more special facilities. If a character isn't in their Bastion on a given Bastion turn, the Bastion acts as though it were issued the Maintain order on that turn unless the owner can communicate with the Bastion hirelings using the \textit{Sending} spell or similar magic.

The orders are as follows:

\subsubsection{Craft}
Hirelings in the special facility craft an item that can be made in that facility. During the time required to craft an item, the facility can't be used to craft anything else, even if a special ability allows the facility to carry out two orders at once. The hirelings have proficiencies with Artisan's Tools as specified in the facility's description.

\subsubsection{Empower}
The special facility confers a temporary empowerment to you or someone else.

\subsubsection{Harvest}
Hirelings gather a resource produced in the special facility. During the time required to harvest a resource, the facility can't be used to harvest anything else, even if a special ability allows the facility to carry out two orders at once.

\subsubsection{Maintain}
All the Bastion's hirelings focus on maintaining the Bastion rather than executing orders in special facilities. Issuing this order prohibits other orders from being issued to the Bastion on the current Bastion turn. Each time the Maintain order is issued, the DM rolls once on the \textit{Bastion Events} table at the end of this chapter. Bastion events are resolved before the next Bastion turn.

\subsubsection{Recruit}
Hirelings recruit creatures to the Bastion. These creatures might include one or more Bastion Defenders, whose main purpose is to defend the Bastion if it is attacked (see "\textit{Bastion Events}" at the end of this chapter). The Bastion generates enough income to provide for the needs of its Bastion Defenders.

\subsubsection{Research}
Hirelings in the special facility gather information.

\subsubsection{Trade}
Hirelings buy and sell goods or services stored or produced in this special facility.

\section{Special Facility Descriptions}
Special facilities are presented in alphabetical order. The Special Facilities table lists all the special facilities presented in this section, along with their prerequisites and orders. Some facilities confer additional benefits, specified in their descriptions.

\begin{DndTable}[header=Special Facilities]{cXXX}

Level & Special Facility & Prerequisite & Order\\
5 & Arcane Study & Ability to use an \textit{Arcane Focus} or a tool as a Spellcasting Focus & Craft\\
5 & Armory & None & Trade\\
5 & Barrack & None & Recruit\\
5 & Garden & None & Harvest\\
5 & Library & None & Research\\
5 & Sanctuary & Ability to use a \textit{Holy Symbol} or \textit{Druidic Focus} as a Spellcasting Focus & Craft\\
5 & Smithy & None & Craft\\
5 & Storehouse & None & Trade\\
5 & Workshop & None & Craft\\
9 & Gaming Hall & None & Trade\\
9 & Greenhouse & None & Harvest\\
9 & Laboratory & None* & Craft\\
9 & Sacristy & Ability to use a \textit{Holy Symbol} or \textit{Druidic Focus} as a Spellcasting Focus & Craft\\
9 & Scriptorium & None* & Craft\\
9 & Stable & None & Trade\\
9 & Teleportation Circle & None & Recruit\\
9 & Theater & None & Empower\\
9 & Training Area & None & Empower\\
9 & Trophy Room & None & Research\\
13 & Archive & None & Research\\
13 & Meditation Chamber & None & Empower\\
13 & Menagerie & None & Recruit\\
13 & Observatory & Ability to use a Spellcasting Focus & Empower\\
13 & Pub & None & Research\\
13 & Reliquary & Ability to use a \textit{Holy Symbol} or \textit{Druidic Focus} as a Spellcasting Focus & Harvest\\
17 & Demiplane & Ability to use an \textit{Arcane Focus} or a tool as a Spellcasting Focus & Empower\\
17 & Guildhall & Expertise in a skill & Recruit\\
17 & Sanctum & Ability to use a \textit{Holy Symbol} or \textit{Druidic Focus} as a Spellcasting Focus & Empower\\
17 & War Room & Fighting Style feature or Unarmored Defense feature & Recruit\\
\end{DndTable}

Arcane Study
Archive
Armory
Barrack
Demiplane
Gaming Hall
Garden
Greenhouse
Guildhall
Laboratory
Library
Meditation Chamber
Menagerie
Observatory
Pub
Reliquary
Sacristy
Sanctuary
Sanctum
Scriptorium
Smithy
Stable
Storehouse
Teleportation Circle
Theater
Training Area
Trophy Room
War Room
Workshop
\section{Bastion Events}
Immediately after a character issues the \textit{Maintain} order to their Bastion, the DM rolls once on the Bastion Events table to determine what event, if any, befalls the Bastion before the next Bastion Turn. If an event occurs, the DM reads the event aloud to the player whose character controls that Bastion. (All the events are described in the sections following the table.) The event is resolved immediately, with the player and DM working together to expand story details as needed. If multiple characters issue the Maintain order on the same Bastion turn, roll once on the table for each of them, resolving each event separately even if the Bastions are combined.

Bastion events occur only when a Bastion is operating under the Maintain order, which often means that the Bastion's owner isn't present in the Bastion at the time. That means these events can be opportunities for the player to take on the role of the Bastion's hirelings and roleplay their reactions to these events. The DM can even turn a Bastion event into a cutscene where each player takes on the role of one of the Bastion's hirelings (under the guidance of the player whose character owns the Bastion).

\begin{DndTable}[header=Bastion Events]{cX}

1d100 & Event\\
01–50 & \textit{All Is Well}\\
51–55 & \textit{Attack}\\
56–58 & \textit{Criminal Hireling}\\
59–63 & \textit{Extraordinary Opportunity}\\
64–72 & \textit{Friendly Visitors}\\
73–76 & \textit{Guest}\\
77–79 & \textit{Lost Hirelings}\\
80–83 & \textit{Magical Discovery}\\
84–91 & \textit{Refugees}\\
92–98 & \textit{Request for Aid}\\
99–00 & \textit{Treasure}\\
\end{DndTable}

\subsection{Event Descriptions}
The events from the Bastion Events table are detailed here in alphabetical order.

\subsubsection{All Is Well}
Nothing significant happens. Roll on the following table, fleshing out the details as you see fit.


\begin{DndTable}{cX}

1d8 & Details\\
1 & Accident reports are way down.\\
2 & The leak in the roof has been fixed.\\
3 & No vermin infestations to report.\\
4 & You-Know-Who lost their spectacles again.\\
5 & One of your hirelings adopted a stray dog.\\
6 & You received a lovely letter from a friend.\\
7 & Some practical joker has been putting rotten eggs in people's boots.\\
8 & Someone thought they saw a ghost.\\
\end{DndTable}

\begin{DndSidebar}[float=!b]{Bastion Tracker}
Players can use the Bastion Tracker sheet as a record of the facilities and other characteristics of their Bastions. Encourage your players to list their hirelings on this sheet and develop them more fully as NPCs on separate sheets (perhaps using the \textit{NPC Tracker} in \textit{chapter 3}). Similarly, the Bastion Tracker provides space to indicate the space of each special facility, but encourage players to draw maps of their Bastions (see "\textit{Bastion Map}" earlier in this chapter).



\end{DndSidebar}

\subsubsection{Attack}
A hostile force attacks your Bastion but is defeated.

Roll 6d6; for each die that rolls a 1, one Bastion Defender dies. Remove these Bastion Defenders from your Bastion's roster. If the Bastion has zero Bastion Defenders, one of the Bastion's special facilities (determined randomly) is damaged and forced to shut down.

A special facility that shuts down can't be used on your next Bastion turn, after which it is repaired and made operational again at no cost to you.

\subsubsection{Criminal Hireling}
One of your Bastion's hirelings has a criminal past that comes to light when officials or bounty hunters visit your Bastion with a warrant for the hireling's arrest. You can retain the hireling by paying a bribe of 1d6 × 100 GP. Otherwise, the hireling is arrested and taken away. If this loss leaves one of your facilities without any hirelings, that facility can't be used on your next Bastion turn. The hireling is then replaced at no cost to you.

\subsubsection{Extraordinary Opportunity}
Your Bastion is given the opportunity to host an important festival or celebration, fund the research of a powerful spellcaster, or appease a domineering noble. Work with the DM to determine the details.

If you seize the opportunity, you must pay 500 GP to cover costs. In return, your Bastion gains a sudden influx of recognition or attention, prompting the DM to roll again on the \textit{Bastion Events} table (rerolling this result if it comes up again).

If you decline the opportunity, you don't pay the money and nothing else happens.

\subsubsection{Friendly Visitors}
Friendly visitors come to your Bastion, seeking to use one of your special facilities. They offer 1d6 × 100 GP for the brief use of that facility. For example, a knight might want your \textit{Smithy} to replace a horseshoe or repair a damaged weapon or suit of armor, or sages might need your Arcane Study to help them settle a dispute. Their use of the facility doesn't interrupt any orders you've issued to it.

\subsubsection{Guest}
A Friendly guest comes to stay at your Bastion. Determine the guest by rolling on the following table, and work with your DM to flesh out the details.


\begin{DndTable}{cX}

1d4 & Guest\\
1 & The guest is an individual of great renown who stays for 7 days. At the end of their stay, the guest gives you a letter of recommendation (see "\textit{Marks of Prestige}" in \textit{chapter 3}).\\
2 & The guest requests sanctuary while avoiding persecution for their beliefs or crimes. They depart 7 days later, but not before offering you a gift of 1d6 × 100 GP.\\
3 & The guest is a mercenary, giving you one additional Bastion Defender. The guest doesn't require a facility to house them, and they stay until you send them away or they're killed.\\
4 & The guest is a Friendly [Attitude] monster, such as a brass dragon or a treant. If your Bastion is attacked while this monster is your guest, it defends your Bastion, and you lose no Bastion Defenders. The monster leaves after it defends your Bastion once or when you send it away.\\
\end{DndTable}

\subsubsection{Lost Hirelings}
One of your Bastion's special facilities (determined randomly) loses its hirelings. The cause of their departure is up to you. The facility can't be used on your next Bastion turn, but the hirelings are replaced at no cost to you at that point.

\subsubsection{Magical Discovery}
Your hirelings discover or accidentally create an Uncommon magic item of your choice at no cost to you. The magic item must be a Potion or Scroll.

\subsubsection{Refugees}
A group of 2d4 refugees fleeing from a monster attack, a natural disaster, or some other calamity seeks refuge in your Bastion. If your Bastion lacks a basic facility large enough to house them, the refugees camp right outside the Bastion. The refugees offer you 1d6 × 100 GP as payment for your hospitality and protection. They stay until you find them a new home or a hostile force attacks your Bastion.

\subsubsection{Request for Aid}
Your Bastion is called on to help a local leader. Perhaps there's a search on for a missing person, or brigands are plaguing the area. If you help, you must dispatch one or more Bastion Defenders. Roll 1d6 for each Bastion Defender you send. If the total is 10 or higher, the problem is solved and you earn a reward of 1d6 × 100 GP. If the total is less than 10, the problem is still solved, but the reward is halved and one of your Bastion Defenders is killed. Remove that Bastion Defender from your Bastion's roster.

\subsubsection{Treasure}
Your Bastion acquires an art object or a magic item determined by rolling on the table below and then rolling on the specified table in \textit{chapter 7}. How the Bastion acquires this treasure is up to you. It might represent an inheritance, a gift from a guest or an admirer, a theft, or a fortunate discovery. If you're in the Bastion, you can claim the treasure immediately; otherwise, it is placed in storage until you can claim it.


\begin{DndTable}{cX}

1d100 & Treasure\\
01–40 & Roll on the 25 GP \textit{Art Objects} table.\\
41–63 & Roll on the 250 GP \textit{Art Objects} table.\\
64–73 & Roll on the 750 GP \textit{Art Objects} table.\\
74–75 & Roll on the 2,500 GP \textit{Art Objects} table.\\
76–90 & Roll on a Common Magic Items table of your choice (\textit{Arcana}, \textit{Armaments}, \textit{Implements}, or \textit{Relics}).\\
91–98 & Roll on an Uncommon Magic Items table of your choice (\textit{Arcana}, \textit{Armaments}, \textit{Implements}, or \textit{Relics}).\\
99–00 & Roll on a Rare Magic Items table of your choice (\textit{Arcana}, \textit{Armaments}, \textit{Implements}, or \textit{Relics}).\\
\end{DndTable}

\section{Fall of a Bastion}
A player character can lose their Bastion in the following ways:


\begin{description}
\item[Divestiture.]
A character can give up their Bastion anytime, releasing the Bastion's hirelings and abandoning the location. The divested Bastion is quickly vacated, is eventually looted, and might even be burned to the ground.
\item[Neglect.]
If a character issues no orders to their Bastion for a number of consecutive Bastion turns equal to the character's level (typically because the character is dead or otherwise out of commission), the hirelings abandon the Bastion and the site is eventually looted. If the character returns later, they can start a new Bastion, perhaps building it amid the ruins of the old one.
\item[Ruination.]
Drawing the Ruin card from the \textit{Deck of Many Things} (as described in \textit{chapter 7}) instantly deprives a character of their Bastion. When such an event occurs, the player can decide what terrible fate befalls the Bastion. The Bastion might be sacked by enemies or destroyed by an earthquake, for example.
\end{description}

Regardless of how the Bastion falls, the player can work with the DM to establish a new Bastion and determine how it comes into being. Use the \textit{Special Facility Acquisition} table to determine how many special facilities come with it. The new Bastion also starts with two basic facilities (one Cramped and one Roomy) of the player's choice.

\appendix
\chapter{Lore Glossary}
This appendix provides brief descriptions for many of the D\&D game's most famous heroes, villains, creatures, locations, and materials. The entries are presented in alphabetical order.

\section{Acererak}
Acererak \textit{(ah-SAIR-er-rack)} is a powerful lich who travels between worlds and takes pleasure in devouring the souls of adventurers, whom he lures into trap-ridden dungeons to suffer horrible deaths. The most famous of such dungeons is the Tomb of Horrors, hidden in the Vast Swamp in the Greyhawk setting (see "\textit{Greyhawk Gazetteer}" in \textit{chapter 5}); another lies under the lost city of Omu in the jungles of Chult in the Forgotten Realms setting (described in the adventure Tomb of Annihilation).

\section{Adamantine}
Adamantine is one of the hardest substances in existence, a dark metal found in meteorites and extraordinary mineral veins. (See the \textit{Adamantine Armor} and \textit{Adamantine Weapon} magic items in \textit{chapter 7}.)

\section{Alustriel Silverhand}
Alustriel \textit{(ah-LOOSE-tree-ell)} Silverhand is the second of seven daughters of Mystra, a deity of magic in the Forgotten Realms. (\textit{Laeral Silverhand} is Alustriel's younger sister.) She ruled the city of Silverymoon for centuries but stepped down from that position a little over a century ago to promote goodness and compassion in the multiverse. She has befriended adventurers such as \textit{Drizzt Do'Urden} and worked alongside the \textit{Harpers} in pursuit of these aims.

\section{Ashardalon}
The terrible red dragon Ashardalon \textit{(ah-SHAR-duh-lawn)} is legendary across the multiverse for his rampages, which turned lush grasslands into ashen plains and caused mighty citadels to be swallowed into the earth. As he grew aware of worlds beyond his home, he recruited a balor demon named Ammet, the Eater of Souls, to help him extend his reach and increase his power. When an adventurer dealt a mortal blow to Ashardalon's heart, the dragon replaced the injured organ with the balor. Ashardalon's features then took on a fiendish cast, and his eyes now glow with demonic flame. A secretive cult of Ashardalon's followers and spawn serves him on the Material Plane, furthering the dragon's schemes to attain godhood.

\section{Baba Yaga}
Baba Yaga \textit{(BAH-bah YAH-gah)} is an arch-hag known as the Mother of All Witches. She is also the adoptive mother of \textit{Iggwilv}. Baba Yaga is famous for her chicken-legged hut, in which she travels across the planes of the multiverse. This impatient, foul-tempered hag is a font of knowledge about all things magical.

\section{Bahamut}
Bahamut \textit{(ba-HA-mutt)} is one of the primordial dragons who (along with \textit{Tiamat}) is said to have created the \textit{First World}. For practical purposes, he is a god—ageless and immortal—who has dwelled in \textit{Mount Celestia} (see \textit{chapter 6}) since the destruction of the First World. In the Dragonlance setting, where he is called Paladine \textit{(PAL-a-deen)}, he is the greatest of the gods of good. On other worlds, he is revered as a god of justice and nobility and is favored by Paladins.

To metallic dragons, Bahamut is more like a king than a god. Individual dragons might owe Bahamut allegiance, respect him, pay tribute to him, and strive to emulate him, but they don't worship him.

\section{Baldur's Gate}
The city of Baldur's Gate \textit{(BAWL-durz GATE)}, in the Forgotten Realms setting, is a teeming metropolis haunted by the lingering influence of three evil gods (Bane, Bhaal, and Myrkul) who refuse to stay dead. It's a place where a sword for hire can find a rich patron, join a secret guild, pursue killers for a bounty, or aid desperate citizens. The city offers opportunities for good-hearted champions to fight against corruption and bring murderers to justice, while less moral mercenaries find a good price for their services. Baldur's Gate hosts the most reliable and ruthless market on the \textit{Sword Coast}. Information, treasures, secrets, and souls can be bought or sold for the right price.

\section{Barovia}
Barovia \textit{(buh-ROVE-ee-ah)} is a \textit{Domain of Dread} sequestered in the \textit{Shadowfell} (see \textit{chapter 6})—a sort of spiritual prison for \textit{Strahd von Zarovich} in his \textit{Castle Ravenloft}. It's also the name of a village in that domain.

\section{Bigby}
Bigby \textit{(BIG-bee)} is a former apprentice of \textit{Mordenkainen}. Though he began his career determined to use magic to dominate and control others, Bigby eventually changed his ways and has worked hard to make amends for his past villainy. Mordenkainen welcomed Bigby into the adventuring company known as the Citadel of Eight (which later became the \textit{Circle of Eight}, described in \textit{chapter 5}). Shy and soft-spoken, Bigby was often eclipsed by his mentor, who taught Bigby how to control his ambitions. After years of adventuring, Bigby crafted a handful of spells that gained him widespread renown. Of these spells, \textit{Bigby's Hand} is his undisputed magnum opus.

A recent misadventure led to Bigby's untimely demise. After he was crushed to death by a frost giant's boulder, Bigby was the target of a \textit{Reincarnate} spell. The spell transformed Bigby from a human into a gnome. Before setting off on his next adventure, Bigby was overheard saying how curious he was to experience the multiverse from a gnome's perspective.

\section{Boo}
Boo is a hamster. More precisely, he is a miniature giant space hamster, a species native to \textit{Wildspace} (see \textit{chapter 6}) that is both sapient and telepathic. Boo's adventures with \textit{Minsc}, as well as the hamster's ferocity, have given Boo legendary status, particularly in the city of \textit{Baldur's Gate}.

\section{Castle Ravenloft}
Castle Ravenloft \textit{(RAY-ven-loft)} is the heart of the domain of \textit{Barovia} and the home of the vampire \textit{Strahd von Zarovich}. \textit{Ravenloft} is also one of the most famous dungeon adventures in the history of the D\&D game, providing a horror-themed experience that inspired the creation of the Ravenloft setting and the adventure Curse of Strahd.

\section{Companions of the Hall, The}
Named for the dwarf stronghold of Mithral Hall in the Forgotten Realms setting, the Companions of the Hall is a group of heroic adventurers whose exploits have spanned centuries and even crossed the bounds of death. \textit{Drizzt Do'Urden} is the central figure in this party; all his companions have died and been reincarnated in new forms in recent years. These other companions are Drizzt's wife, Cattie-Brie; his adoptive father, Bruenor Battlehammer; and his friends Wulfgar and Regis.

\section{Company of Seven, The}
The Company of Seven was a group of adventurers active hundreds of years ago in the Greyhawk setting. Its members included Heward, Keoghtom, Murlynd, Nolzur, Quaal, Tasha (see \textit{Iggwilv}), and Zagig (see \textit{Zagyg}). Some of these adventurers achieved near-divine status, and most of them are remembered for magic items that carry their names. The group inspired the formation of the \textit{Circle of Eight} (described in \textit{chapter 5}).

\section{Corellon}
At the dawn of the multiverse, Corellon \textit{(core-ELL-on or CORE-eh-lawn)} danced from world to world and plane to plane. A being of consummate mutability and infinite grace, Corellon is a whimsical shape-shifter, able to take the form of a chuckling stream, a teasing breeze, an incandescent beam, a school of fish, or a flock of birds. Corellon's flamboyant, mercurial personality infuses every form the god adopts. Corellon loves wholeheartedly and takes pleasure from every encounter with other divine beings of the multiverse.

According to legend, the first elves emerged from Corellon's shed blood, and they shared the god's changeable and audacious nature. Many elves, along with members of other species, worship Corellon.

\section{Delina}
Delina \textit{(dell-EEN-ah)} is a young elf Sorcerer who wields the unpredictable power of wild magic. Finding herself in trouble in the city of \textit{Baldur's Gate}, she accidentally reawakened the ancient hero of the city, \textit{Minsc}, and got thrown into further adventures pursuing her lost twin brother.

\section{Diancastra}
Diancastra \textit{(DIE-ann-CAST-rah)} is a demigod and a daughter of the divine ancestor of giants, Annam. She is a trickster, an adventurer, and a scholar of magic who enjoys wandering the Material Plane in search of new curiosities and spells to learn. She longs to see the descendants of Annam—storm, cloud, fire, frost, stone, and hill giants—restored to the position of honor and respect they held in ancient times.

\section{Drizzt Do'Urden}
Drizzt Do'Urden \textit{(DRIZT doh-UR-den)} is a drow exile from the city of \textit{Menzoberranzan} and a fugitive from the wrath of \textit{Lolth} (see \textit{chapter 6}) and her priests. He wandered the surface world and gathered a circle of friends who became known as the \textit{Companions of the Hall}.

\section{Elder Evils}
The Elder Evils are a variety of entities whose existence dates to the beginnings of the multiverse—or possibly predates it. Some Elder Evils are creatures of the \textit{Far Realm} (see \textit{chapter 6}), while others are akin to gods or primordial beings of the Elemental Planes. Some are thought to be imprisoned, while others are said to be slumbering until they awaken in some apocalyptic cataclysm.

The names given to these terrible entities include such strange descriptions as Atropus, the World Born Dead; Dendar, the Night Serpent; \textit{Hadar}, the Dark Hunger; Haemnathuun, the Blood Lord; Ityak-Ortheel, the Elf-Eater; Kezef, the Chaos Hound; \textit{Kyuss}, the Worm That Walks; the Queen of Chaos; \textit{Tharizdun}, the Chained God; Tyranthraxus, the Flamed One; and \textit{Zargon}, the Returner. They are all forces of corruption and evil. Nothing good can come from their influence. Bargains made with them end in catastrophe or death.

\section{Elminster}
Elminster \textit{(el-MIN-ster)} is a powerful and ancient archmage in the Forgotten Realms setting. As one of Mystra's Chosen, divinely called and empowered by a deity of magic, Elminster fosters magic and protects the fabric of magic in the world. Though this responsibility demands a certain amount of neutrality and dispassionate judgment, Elminster has a fundamentally kind and compassionate heart.

\section{Euryale}
One card in the \textit{Deck of Many Things} (see \textit{chapter 7}) bears a person's proper name, and the card's namesake, Euryale \textit{(YUR-ee-ale or yur-EYE-a-lee)}, is the subject of much speculation. Often assumed to be a fearsome demigod (perhaps the first medusa) or the wielder of a destructive curse, Euryale is actually a key part of the story of the magical deck's creation. After befriending a princess named Asteria and spending many years in her dear friend's company, Euryale was captured and sentenced to death by Asteria's father. Asteria pleaded with the gods to save her friend, and Istus (a god from the Greyhawk setting; see \textit{chapter 5}) intervened to help the pair rewrite their story, creating the \textit{Deck of Many Things}. The magic of the deck helped the two escape, and they adventured together across the multiverse. Eventually, Euryale—having become an ancient, wise, and powerful Druid—settled in the \textit{Outlands} (see \textit{chapter 6}), where she still sometimes acts as a patron, mentor, or ally for adventurers.

\section{Fallbacks, The}
Tessalynde is a young elf Rogue who dreams of leading the foremost adventuring party of the Forgotten Realms setting. While the crew she's gathered isn't the stuff of legend yet, she's confident her guidance can get them there. Called the Fallbacks, the team includes Anson, a human Fighter too stubborn to stay down; Cazrin, a self-taught, human Wizard determined to test her theoretical mettle against the real world; Baldric, a dwarf Cleric who refuses to tie himself to a single deity when he can trade favors with them all; Lark, a tiefling Bard with as many secrets as songs; and Uggie, a pet otyugh.

\section{First World, The}
Scholars speak of a primordial state, a single reality they call the First World, which preceded the Material Plane. Many of the peoples and monsters that inhabit the worlds on the Material Plane originated there. After the First World was shattered by a great cataclysm—giving birth to the worlds that came in its wake—the progeny of the first elves, dwarves, beholders, and other iconic creatures took root on world after world, like seeds scattered by a cosmic wind. If the musings of these great sages are true, every world of the Material Plane is a reflection—or, in some cases, a distortion—of the First World.

\section{Fizban}
In the Dragonlance setting, \textit{Bahamut}—who is known there as Paladine—dwelled among mortals in human guise for a time, aiding the forces of good against Takhisis (see \textit{Tiamat}). He appeared as a bumbling old mage named Fizban \textit{(FIZZ-ban)} the Fabulous.

\section{Great Modron March, The}
When the gears of the plane of \textit{Mechanus} (see \textit{chapter 6}) complete seventeen cycles—once every 289 years—the modron \textit{(MOE-dron)} leader, Primus, sends a vast army of modrons across the Outer Planes. The purpose of this march is unclear. Most believe it to be a data-gathering mission meant to ascertain the current state of the cosmos, but some see it as reconnaissance aimed at some future act of conquest. The march is long and dangerous, and only a small number of modrons returns to Mechanus.

\section{Gruumsh}
Gruumsh \textit{(GROOMSH)} is a warring god who is often described as the creator or patron of the orc people. Some orcs attribute their tenacity and toughness to Gruumsh's lingering influence. Some myths describe a primeval conflict between Gruumsh and \textit{Corellon}, which resulted in Gruumsh losing one eye and Corellon's spilled blood becoming the first elves.

\section{Hadar}
Hadar \textit{(HAY-dar or ha-DARR)}, the Dark Hunger, is an ancient stellar entity originating from the \textit{Far Realm} (see \textit{chapter 6}). It appears as a cinder-red dying star, barely visible in the night sky, and it siphons life from its minions to avert its own demise. Two widely used Warlock spells invoke Hadar's power (see the \textit{Arms of Hadar} and \textit{Hunger of Hadar} spells in the \textit{Player's Handbook}), and a few Warlocks claim this \textit{Elder Evil} as their Great Old One patron.

\section{Harpers, The}
The Harpers is a scattered network of spellcasters and spies in the Forgotten Realms setting. Its members advocate equality and covertly oppose the abuse of power, magical or otherwise.

The Harpers has risen, been shattered, and risen again several times. The faction's longevity and resilience are largely due to its decentralized, grassroots, secretive nature and the autonomy of its members. It has small cells and lone operatives throughout the Forgotten Realms, although they interact and share information with one another as the need arises. The Harpers' ideology is noble, and members pride themselves on their ingenuity and incorruptibility. Harpers don't seek power or glory, only fair treatment for all.

\section{Heroes of the Lance, The}
The adventurers known as the Heroes of the Lance, whose deeds helped prevent Takhisis (see \textit{Tiamat}) and her evil dragons from conquering the world of the Dragonlance setting, began as a small group of young adventurers and their aging dwarf mentor. This original group (known as the Innfellows) consisted of Tanis Half-Elven, the brothers Raistlin and Caramon Majere, Sturm Brightblade, a kender (similar to a halfling) named Tasslehoff Burrfoot, the dwarf Flint Fireforge, and Kitiara Uth Matar, Raistlin and Caramon's older half-sister. While Kitiara eventually joined Takhisis's service, the heroes were joined by others: Riverwind and Goldmoon, from the nomadic peoples of the plains; Gilthanas and Laurana, elf friends from Tanis's childhood; and Tika Waylan, a young barmaid whom the original Innfellows remembered as just a kid.

\section{Heroes of the Realm, The}
The so-called "heroes of the Realm" are a group of young adventurers—Bobby, Diana, Eric, Hank, Presto, and Sheila—who traveled from Earth in the 1980s into a world in the D\&D multiverse. Equipped with powerful magic items, they foiled the schemes of foes such as \textit{Venger} and \textit{Tiamat} while seeking some means to return home. Eventually, as the heroes' mastery of adventuring skills increased, they discovered the secrets of traveling between worlds, though they still haven't found a way home.

The heroes of the Realm weren't the only kids transported from Earth to the worlds of D\&D. Other young adventurers, including Niko the Cleric, are still exploring the vast D\&D multiverse.

\section{Icewind Dale}
Icewind Dale is the northernmost settled region of the Forgotten Realms setting. Freezing wind sweeps across the tundra, finding its way through every crack and draining any hint of warmth. Between the ice cliffs of the eastern glacier, the snowcapped peaks of the mountains to the south, and the Sea of Moving Ice to the north and west, ten small towns cluster around three icy lakes.

\section{Iggwilv}
Before she changed her name and conquered enough of Eastern Oerik (the Greyhawk setting) to rightfully call herself the Witch Queen of Perrenland, Iggwilv \textit{(IGG-wilv)} was known as Tasha, a human mage who began her career as the apprentice of Zagig Yragerne (see \textit{Zagyg}). Later, as an adventurer, she created several new spells, including \textit{Tasha's Hideous Laughter} and \textit{Tasha's Bubbling Cauldron}, leaning on the teachings of her adoptive mother, the arch-hag \textit{Baba Yaga}. As Tasha grew in power and made powerful enemies, she changed her name to Iggwilv. In this guise, she became enchanted with the power of the Abyss and wrote the definitive treatise on demonkind: the \textit{Demonomicon of Iggwilv} (see \textit{chapter 7}). She also bound and trapped the demon lord \textit{Graz'zt} (see \textit{chapter 6}).

Iggwilv ruled Perrenland as a tyrant. When Graz'zt escaped his magical prison, Iggwilv went into hiding. Iggwilv's current location is unknown, but she left behind a cambion son (\textit{Iuz}), who has his mother's tyrannical bent, and a daughter (Drelnza), who is now a vampire and lairs in the Lost Caverns of Tsojcanth (see "\textit{Greyhawk Gazetteer}" in \textit{chapter 5}), not far from her mother's old haunts.

\section{Iuz}
Iuz \textit{(EYE-ooze or eye-OOZE)} is a cambion and the son of \textit{Iggwilv} and \textit{Graz'zt} (see \textit{chapter 6}). He is every bit as evil as his father and as bent on conquest as his mother at her very worst. He rules a significant portion of Eastern Oerik (in the Greyhawk setting), and some fear that he aspires to conquer even more territory. See "\textit{Greyhawk's Premise}" in \textit{chapter 5}.

\section{Jallarzi Sallavarian}
The youngest and most recently appointed member of the \textit{Circle of Eight} (see \textit{chapter 5}), Jallarzi Sallavarian \textit{(juh-LAR-zee sal-ah-VAIR-ee-en)} is a Warlock with a celestial patron. She is more inclined toward good than the more neutral members of the Circle of Eight and enjoys wandering the Free City of Greyhawk (see \textit{chapter 5}) in disguise, monitoring sinister elements of society there.

\section{Kas the Betrayer}
Kas \textit{(KOSS)} is a vampire, legendary swordfighter, and ruthless warlord. He once served as the leader of \textit{Vecna}'s armies and the lich's most trusted lieutenant, and he wielded a sword made for him (the \textit{Sword of Kas}, described in \textit{chapter 7}) by his liege. But the evil sword convinced Kas to betray Vecna, and now Kas is driven primarily by his hatred for his former lord.

\section{Keraptis}
According to legend, Keraptis \textit{(kuh-RAP-tiss)} was an evil archmage who long ago seized power as a petty tyrant somewhere in the northern mountains of the Greyhawk setting. When his demands on the local populace became too extreme, a revolt drove him into hiding. He disappeared into the tunnels beneath White Plume Mountain (see "\textit{Greyhawk's Premise}" in \textit{chapter 5}). It's conceivable that he became a lich as the grim climax of his magical research.

\section{Kyuss}
Variously identified as an \textit{Elder Evil}, as a demigod, or as a mortal priest of the demon lord \textit{Orcus} (see \textit{chapter 6}), Kyuss \textit{(KYE-uss)} is a mysterious figure best known as the Worm That Walks. Kyuss manifests on the Material Plane as a colossal mass of maggots and worms animated by a single evil will.

\section{Laeral Silverhand}
Laeral \textit{(LAIR-ull)} Silverhand is the fifth of seven daughters of Mystra, a deity of magic in the Forgotten Realms setting. (\textit{Alustriel Silverhand} is Laeral's older sister.) She is over 700 years old, and after a period of withdrawal from public life, she now serves as the Open Lord of \textit{Waterdeep}—the only one of the twenty or so Lords of Waterdeep whose identity is publicly known. Though her power has waxed and waned depending on Mystra's involvement with the world and has declined significantly in recent years, Laeral remains one of the most formidable spellcasters in a world known for powerful spellcasters.

\section{League of Malevolence, The}
The League of Malevolence is an odious assemblage of villains united in one purpose: the accumulation of power. Its founding member, Kelek the Sorcerer, expects his confederates to work together for their mutual benefit, but he also encourages them to pursue their own evil schemes. Other members include the remorseless killer for hire Warduke, the assassin Zarak, and the evil Zargash, a servant of Orcus. Skylla, a Warlock whose power derives from a pact with \textit{Baba Yaga}, maintains a tenuous alliance with the league but despises Kelek.

\section{Lord Soth}
Lord Soth is a death knight in the Dragonlance setting—a former Paladin who allowed his pride to lead him into evil. Offered an opportunity to redeem himself by averting a cataclysmic event, his pride again led him to abandon his quest and allow the devastation of his world. His mortal life ended in that cataclysm, but he rose from the ashes into undeath and retreated to his cursed castle, Dargaard Keep. He ignored the outside world for many centuries until the Dragon Queen Takhisis (see \textit{Tiamat}) called him forth to serve her in the War of the Lance.

\section{Menzoberranzan}
The \textit{Underdark} city of Menzoberranzan \textit{(men-zoh-buh-RAN-zan)}, lying far beneath the \textit{Sword Coast} region of the Forgotten Realms, is ruled by tyrannical priests of the Demon Queen of Spiders, \textit{Lolth} (see \textit{chapter 6}). Most of its people are drow, whose noble houses are locked in a constant struggle for Lolth's favor and the power that comes with it.

\section{Minsc}
Minsc \textit{(MINSK)} is a heroic Ranger from the land of Rashemen in the Forgotten Realms setting. Known for his good cheer, his hearty optimism, and his pet miniature giant space hamster, \textit{Boo}, Minsc is particularly adored in the city of \textit{Baldur's Gate}.

\section{Mithral}
Mithral \textit{(MITH-ral)} is a light, flexible metal that resembles silver but is much more durable. Its most common use is for crafting lightweight armor that provides excellent protection without bulk or burden (see the \textit{Mithral Armor} magic item in \textit{chapter 7}).

\section{Moradin}
Myths told on many worlds describe Moradin \textit{(MORE-ah-din)} crafting all the peoples of the Material Plane in his great workshop in \textit{Mount Celestia} (see \textit{chapter 6}), causing them to spring to life from inert metal when he cooled the heated castings with his breath. These myths are the source of Moradin's appellation, "All-Father." Many peoples across the multiverse, including many dwarves, worship Moradin.

\section{Mordenkainen}
Mordenkainen \textit{(mor-den-KIGH-nen or mor-den-KAY-nen)} is a human archmage who maintains a residence in the Free City of Greyhawk and another hidden somewhere in the Yatil Mountains between Perrenland and Ket (see \textit{chapter 5}). He is the founder and leader of the \textit{Circle of Eight}, which he created to implement his philosophy of balance in the world. Among Mordenkainen's celebrated accomplishments are several name-branded spells, including ones that create extradimensional mansions, swords of spectral force, and phantom watchdogs (detailed in the \textit{Player's Handbook}).

Mordenkainen has visited numerous worlds and planes of existence, making friends and fomenting rivalries on many of them. He maintains a close friendship with his former apprentice \textit{Bigby} and the famous archmage \textit{Elminster} of Shadowdale. Mordenkainen's greatest rival is \textit{Iggwilv}, a former adventuring companion.

\section{Neverwinter}
Once known in the Forgotten Realms setting as the Jewel of the North, the city of Neverwinter was badly damaged when a nearby volcano erupted about fifty years ago. Now, its people work furiously to rebuild. Some of its outer walls still lie in ruins, and several of its neighborhoods remain abandoned.

Dagult Neverember is the Lord Protector of Neverwinter, ruling in the absence of a true heir to the crown. His dream is to see Neverwinter eclipse nearby \textit{Waterdeep} in wealth and prosperity.

\section{Otiluke}
Otiluke \textit{(AW-teh-luke)} is an impulsive, aggressive Wizard who is the main agent of the \textit{Circle of Eight} in the Free City of Greyhawk (see \textit{chapter 5}). He holds a position of political power and keeps his membership in the Circle of Eight secret. He is best known for his spells that create magical spheres (detailed in the \textit{Player's Handbook}).

\section{Otto}
Otto \textit{(AW-toe)} is an affable dwarf Bard with a taste for fine food, good music, and expensively tailored clothes. His sociable and outgoing personality masks the fact that he's also a member of the \textit{Circle of Eight} (see \textit{chapter 5}) and committed to the goals of that organization. He is well-known across the multiverse for his creation of the \textit{Otto's Irresistible Dance} spell.

\section{Phandalin}
The frontier town of Phandalin \textit{(FAN-duh-lin or fan-DAY-lin)} in the \textit{Sword Coast} region of the Forgotten Realms setting is built on the ruins of a much older settlement. Hundreds of years ago, old Phandalin was a thriving town until it was sacked by bandits and lay abandoned for centuries.

In the past few years, hardy folk from the cities of \textit{Neverwinter} and \textit{Waterdeep} have begun settling atop the ruins of old Phandalin. It's now home to farmers, woodcutters, fur traders, and prospectors drawn by stories of gold and platinum in the foothills of the Sword Mountains—as well as more than a few ruffians and bandits.

\section{Prince of Frost, The}
The Prince of Frost is a son of Titania, the \textit{Summer Queen}. He was once known as the Sun Prince, but his heart grew cold when he was spurned by the fey noble he loved. He is now a creature of wrath and winter, ruling from his Fortress of Frozen Tears in the Vale of Long Night in the \textit{Feywild} (see \textit{chapter 6}). He detests mortals of the Material Plane and dreams of covering their many worlds with perpetual winter.

\section{Queen of Air and Darkness, The}
The mysterious Queen of Air and Darkness is the archfey ruler of the \textit{Gloaming Court} in the \textit{Feywild} (see \textit{chapter 6}). Though she is often described as malicious and evil, her motivations are as obscure and whimsical as any other archfey's. She gives supernatural boons and terrible curses with equal whimsy, with little regard for the elaborate customs surrounding gifts and favors in other courts of the Feywild. Though she is said to have a spectral form resembling a terrifying elf, to the unaided eye she appears as a strange, black diamond, hovering in the air like the slitted pupil of an unseen cat.

\section{Raven Queen, The}
The Raven Queen is a being of mystery. Those who claim to have encountered her have described an array of disturbing images: a terrible shadow that clawed at their innermost thoughts, a pale and regal elf who exploded into an untold number of ravens, a shambling tangle of slick roots and sticks that overwhelmed them with dread, or an unknown presence that pulled them screaming into the gloom.

Despite all attempts to demystify her, the Raven Queen has remained enigmatic and aloof. She rules from her Raven Throne within the Fortress of Memories, a mazelike castle deep within the bleakness of the \textit{Shadowfell} (see \textit{chapter 6}). From there she sends out her raven servants to find interesting souls she can pluck from the planes of existence. Once these souls are in the Shadowfell, she watches as they attempt to unravel the mystery of their being.

\section{Rock of Bral, The}
Bral is a city built on an asteroid that drifts through \textit{Wildspace} (see \textit{chapter 6}). Its inhabitants, who hail from many worlds across the Material Plane, typically refer to Bral as the Rock. There is no other place quite like it in Wildspace—a teeming hive of business that spans the breadth of the Astral Sea.

Bral is populated by an outlandish collection of traders, scoundrels, mercenaries, pirates, nobles, and entrepreneurs. Generally, law enforcement is sporadic, which means that order is elusive. Most folks who call the Rock home adhere to two principles: mind your own business whenever possible, and enough gold can fix anything.

\section{Rudolph van Richten}
A scholar and monster hunter, Rudolph van Richten has traveled to dozens of \textit{Domains of Dread} in the \textit{Shadowfell} (see \textit{chapter 6}), investigating reports of monstrous beings and documenting them in a series of published guides, the best known of which is \textit{Van Richten's Guide to Vampires}.

In kinder days, Rudolph lived with his wife, Ingrid, and son, Erasmus, but he lost them both to a wicked vampire. In the process of hunting the vampire, Rudolph came under a terrible curse: "Live you always among monsters, and see everyone you love die beneath their claws."

In the decades since, van Richten has hunted monsters and armed others with the knowledge they need to confront evil. Though he's made many devoted allies, he keeps them at arm's length, fearing the threat of his curse.

Van Richten is closely associated with other heroes who prowl the Domains of Dread, including the detective Alanik Ray and his partner Arthur Sedgwick, the explorer Ez d'Avenir, and the twins Gennifer and Laurie Weathermay-Foxgrove.

\section{Strahd von Zarovich}
A brilliant thinker and capable warrior in life, Strahd von Zarovich \textit{(STRAWD von zuh-ROH-vitch or ZAR-oh-vitch)} fought in countless battles for his people. After a long life of warfare and slaughter, he settled in the remote valley of \textit{Barovia} and built a castle on a towering pinnacle. His brother Sergei came to live with him in \textit{Castle Ravenloft}, becoming Strahd's adviser and constant companion.

In his brother, Strahd saw everything he was not. Resentment colored their relationship and eventually turned into hatred. Strahd's beloved, Tatyana, spurned Strahd for Sergei, whom she pledged to marry.

In a desperate attempt to regain Tatyana's love, Strahd forged a pact with dark powers that made him immortal. Strahd confronted his brother at Sergei and Tatyana's wedding and killed him. Tatyana fled and flung herself from Ravenloft's walls. Guards shot Strahd with arrows, but he didn't die. He became a vampire—the first vampire, according to many sages.

In the centuries since his transformation, Strahd's lust for life has grown. He broods in his castle, cursing the living for all he has lost and never admitting to his role in the tragedy.

\section{Summer Queen, The}
Titania, the Summer Queen, is the regal and charismatic ruler of the \textit{Summer Court} in the \textit{Feywild} (see \textit{chapter 6}). Perhaps the mightiest of the archfey, she can ripen a crop with a smile and summon wildfires with the merest crinkling of her brow.

\section{Sword Coast, The}
The Sword Coast is the western edge of the continent of the Forgotten Realms setting, running along the Sea of Swords. It's a narrow band of territory dominated by several major city-states, from \textit{Neverwinter} in the north to \textit{Baldur's Gate} in the south, with \textit{Waterdeep} in between.

\section{Szass Tam}
Szass Tam \textit{(ZASS TAM)} is a powerful lich and necromancer in the Forgotten Realms setting. As leader of the Red Wizards of Thay, he has countless spies, assassins, and other agents at his service, and his schemes span centuries and reach to the ends of the world.

\section{Tasha}
See \textit{Iggwilv}.

\section{Tharizdun}
For a being known as the Chained God, the \textit{Elder Evil} Tharizdun \textit{(thuh-RIZZ-dun or thair-izz-DOON)} has managed to extend his baleful influence from the Greyhawk setting through many worlds of the Material Plane. He is an ancient force of entropy, the end of all things and the extinction of life. His worshipers are nihilists who seek to end all worlds by liberating their god. Tharizdun is often linked to cults of \textit{Elemental Evil} (see \textit{chapter 5}).

\section{Tiamat}
Tiamat \textit{(TEE-a-mat)} is one of the primordial dragons who (along with \textit{Bahamut}) is said to have created the \textit{First World}. Since the destruction of the First World, she has made her home in the \textit{Nine Hells} (see \textit{chapter 6}), where she is served by devils and enjoys the worship of mortals across the multiverse. In the Dragonlance setting, where she is known as Takhisis \textit{(ta-KEE-sis)}, she is the greatest of the gods of evil. On many worlds, she is known as a god of greed, wealth, and vengeance.

Chromatic dragons might fear, respect, envy, and appease Tiamat as a sovereign, but they don't worship her. Their devotion to her rarely supersedes their devotion to their own goals.

\section{Underdark, The}
The Underdark is a vast subterranean realm of natural caverns and passages extending far beneath the surface of most D\&D worlds. Many creatures inhabit the Underdark, including drow, duergar, svirfneblin, and other species that have adapted to the world below the surface.

\section{Undermountain}
Undermountain is the largest, deepest dungeon in the Forgotten Realms setting. It's a series of interconnected dungeon levels sprawling far beneath the city of \textit{Waterdeep}. Its tunnels connect to preexisting \textit{mithral} mines and natural caverns of the \textit{Underdark}. The best-known entrance to Undermountain is through the \textit{Yawning Portal} tavern in Waterdeep, but other routes also exist. A reclusive and cantankerous archmage named Halaster Blackcloak claims Undermountain as his domain, although his overlordship isn't widely acknowledged by the dungeon's other denizens.

\section{Vajra Safahr}
Vajra Safahr \textit{(VAWJ-rah sah-FAR)} is a human Wizard in her midthirties, making her the youngest person ever to hold the position of Blackstaff—the chief mage of \textit{Waterdeep} and the head of Blackstaff Academy, a school for mages. Several of the older and more seasoned mages of Waterdeep consider Vajra an upstart, but they are smart enough not to challenge her. As the badge of her office, she carries the \textit{Blackstaff}, an Artifact that holds the spirits of all her predecessors.

\section{Valor's Call}
The noble adventuring party known as Valor's Call was founded by Strongheart, a resolute human Paladin committed to destroying evil wherever it rears its head. Strongheart alone determines who can become a member of this prestigious group, and he is always looking for courageous heroes willing to devote themselves to a good cause.

Only good-aligned characters can join Valor's Call. The group enjoys the patronage of \textit{Yolande}, the queen of Celene, and carries out missions on her behalf both in the Greyhawk setting (see \textit{chapter 5}) and in the \textit{Feywild} (see \textit{chapter 6}).

Prominent members of Valor's Call include the dwarf Fighter Elkhorn, the human Cleric Mercion, the human Rogue Molliver, and the human Wizard Ringlerun.

\section{Vecna}
Vecna \textit{(VECK-nah)} had humble beginnings in the Greyhawk setting, where an order of Wizards used him as a bootblack and scribe. He studied magic in secret until he amassed enough power to slaughter the order, and then he turned his efforts toward scribing the \textit{Book of Vile Darkness} (described in \textit{chapter 7}). Armed with that dread tome, he forged a kingdom to rule, with the vampire \textit{Kas} as his lieutenant. But Kas betrayed and killed him, leaving only one hand and one eye intact (the \textit{Eye and Hand of Vecna} are described in \textit{chapter 7}).

Vecna's evil will was so great that he persisted beyond death and eventually became a demigod of secrets and evil magic. His ambition drives him to pursue greater divine power across the multiverse.

\section{Venger}
Venger \textit{(VEN-jur)} was once both human and good-hearted, but some fiendish force transformed him into a winged, one-horned villain. His schemes have spanned centuries and worlds, as he seeks to seize ever-increasing magical power. He has a particular loathing for \textit{Tiamat} and longs to overthrow her. He is often seen astride a flying black steed, his shadowy adviser in tow.

\section{Vi}
Vi \textit{(VYE)} is a gnome artificer, entrepreneur, and traveler of many worlds who originally hails from Eberron. After helping to create terribly destructive weapons during the Last War, Vi sought to atone for this guilt by establishing the Fixers, an organization dedicated to solving apparently hopeless problems. From the Fixers' headquarters in Sigil, Vi has brought her genius, warm heart, and love for a stiff drink to many worlds across the multiverse.

\section{Vlaakith}
Vlaakith \textit{(VLAH-kith)} is the lich-queen of the githyanki city of Tu'narath, adrift in the \textit{Astral Plane} (see \textit{chapter 6}). She is ancient, having fought alongside a leader named Gith to win her people's freedom from the mind flayers that enslaved them. Githyanki of Tu'narath swear unquestioning obedience to the lich-queen, and she promises access to paradise to those who serve her loyally. In truth, she feasts on the souls of any subjects who grow powerful enough to dream of challenging her.

\section{Waterdeep}
Waterdeep is the most famous and cosmopolitan city in the Forgotten Realms setting. It's a center of wealth and influence where people come to make their dreams come true. The city is home to the \textit{Yawning Portal}, the best-known point of access to the sprawling dungeon of \textit{Undermountain}.

\section{Xanathar}
Xanathar \textit{(ZAN-a-thar)} is an eccentric beholder crime lord dwelling beneath \textit{Waterdeep}. Desiring to know all there is to know, Xanathar collects knowledge from across the multiverse, but its most prized possession is its goldfish, Sylgar.

\section{Yawning Portal, The}
The Yawning Portal is a tavern in \textit{Waterdeep}, built around an entrance to the infamous dungeon of \textit{Undermountain}. Adventurers throughout the Forgotten Realms setting and elsewhere in the multiverse visit the Yawning Portal to exchange knowledge about Undermountain and other dungeons. Most visitors are content to swap stories by the hearth, but some adventurers pay the toll for entry into Undermountain (collected by the mysterious owner and bartender, Durnan).

\section{Yolande}
Yolande \textit{(yoh-LAWND}), known as the Faerie Queen, is the benevolent and beloved elf monarch of Celene (see \textit{chapter 5}). Raised in the court of the \textit{Summer Queen}, Yolande had no wish to rule. She preferred the life of an adventuring magic-user. She built her reputation on triumphs, such as her capture of the fomorian brigand Solgna and the theft of the \textit{Prince of Frost}'s sentient sword, \textit{Winterflash}. Yolande was among the first elves to migrate from the \textit{Feywild} to the Greyhawk setting.

\section{Zagyg}
The archmage Zagig Yragerne \textit{(ZAG-igg EE-rag-airn)} was an adventurer in the Greyhawk setting and a member of the \textit{Company of Seven}. At the climax of a long and prosperous adventuring career, he built a fortress known as Castle Greyhawk outside the Free City of Greyhawk (see \textit{chapter 5}). From this stronghold, he took an increasingly powerful role in the politics of the city, contributing to its reputation as "the Gem of the Flanaess." He obsessively delved deeper under his castle and withdrew from the outside world until he managed to entrap nine demigods in magical prisons and claim a fragment of their divine power. He ascended to a minor form of godhood, took the name Zagyg, and took a place in the court of Boccob, a god of magic.

\section{Zargon}
Zargon \textit{(ZAR-gawn)} is an \textit{Elder Evil}—an undying abomination from eons past with an insatiable appetite. A tentacled, slime-covered horror with a cyclopic red eye and an indestructible horn, Zargon corrupts creatures it doesn't devour, transforming its victims into amorphous servants. Zargon is imprisoned on the Material Plane in a prison deep in the earth. This prison is described in the Quests from the Infinite Staircase adventure anthology.

\chapter{Maps}
\section{Barrow Crypt}
\section{Caravan Encampment}
\section{Crossroads Village}
\section{Dragon's Lair}
\section{Dungeon Hideout}
\section{Farmstead}
\section{Keep}
\section{Manor}
\section{Mine}
\section{Roadside Inn}
\section{Ship}
\section{Spooky House}
\section{Underdark Warren}
\section{Volcanic Caves}
\section{Wizard's Tower}
\chapter{Tracking Sheets}
You can find all the tracking sheets found throughout the \textit{Dungeon Master's Guide} below. A combined sheet is also available.

\section{Game Expectations}
\section{Travel Planner}
\section{NPC Tracker}
\section{Settlement Tracker}
\section{Campaign Journal}
\section{DM's Character Tracker}
\section{Campaign Conflicts}
\section{Magic Item Tracker}
\section{Bastion Tracker}
\chapter{Credits}

\begin{description}
\begin{multicols}{2}

\begin{description}
\item[Lead Designers:]
Christopher Perkins, James Wyatt
\item[Designers:]
Jeremy Crawford, F. Wesley Schneider, Ray Winninger
\item[Rules Developers:]
Jeremy Crawford (lead), Makenzie De Armas, Ben Petrisor
\item[Editors:]
Adrian Ng (lead), Judy Bauer, Janica Carter
\item[Playtest Analysts:]
Ron Lundeen (lead), Dan Dillon, Patrick Renie
\item[Art Directors:]
Kate Irwin (lead), Josh Herman
\item[Graphic Designers:]
Trish Yochum (lead), Matt Cole
\item[Cover Illustrators:]
Tyler Jacobson, Simen Meyer, Olena Richards
\item[Interior Illustrators:]
Helder Almeida, Joy Ang, David Astruga, Alfven Ato, Tom Babbey, Helge C. Balzer, Luca Bancone, Mark Behm, Eric Belisle, Olivier Bernard, Zoltan Boros, Bruce Brenneise, Aleksi Briclot, Ekaterina Burmak, Filip Burburan, Paul Scott Canavan, Dawn Carlos, Clint Cearley, Diana Cearley, Sidharth Chaturvedi, David René Christensen, Conceptopolis, Harry Conway, CoupleOfKooks, Daarken, Kent Davis, Nikki Dawes, Axel Defois, Alex Diaz, Simon Dominic, Olga Drebas, Wayne England, Aurore Folny, Jessica Fong, Vallez Gax, Justyna Gil, Ilse Gort, Alexandre Honoré, Ralph Horsley, Jason Juta, Sam Keiser, Julian Kok, Katerina Ladon, Abigail Larson, Olly Lawson, Linda Lithen, Titus Lunter, Andrew Mar, Raluca Marinescu, Viko Menezes, Brynn Metheney, Robson Michel, Calder Moore, Martin Mottet, Jodie Muir, Scott Murphy, David Auden Nash, Irina Nordsol, William O'Connor, Robin Olausson, Claudio Pozas, Livia Prima, April Prime, Noor Rahman, Chris Rallis, Chris Seaman, Andrea Sipl, Craig J Spearing, Chase Stone, Joel Thomas, Beth Trott, Brian Valeza, Randy Vargas, Richard Whitters, Kieran Yanner, Zuzanna Wuzyk
\item[Cartographers:]
Francesca Baerald, Dyson Logos, Mike Schley
\item[Concept Art Director:]
Josh Herman
\item[Concept Artists:]
Even Amundsen, Carlo Arellano, Michael Broussard, John Grello, Tyler Jacobson, 9B Collective, Noor Rahman
\end{description}


\begin{description}
\item[Consultants:]
Zac Clay, Jennifer Kretchmer, James Mendez, Matthew Mercer, Jonathan Tomhave, Alyssa Visscher, Deborah Ann Woll
\item[Additional Consultation:]
Justice Ramin Arman, Kyle Brink, Amanda Hamon, Jay Jani, Carl Sibley, Emi Tanji, Jason Tondro
\end{description}


\begin{description}
\item[Producers:]
Dan Tovar (lead), Bill Benham, Siera Bruggeman, Robert Hawkey
\item[Print Operations Engineers:]
Basil Hale, Scott West
\item[Imaging Technicians:]
Daniel Corona, Meagan Kenreck, Kevin Yee
\item[Prepress Specialist:]
Jefferson Dunlap
\end{description}


\begin{description}
\item \textbf{Based on the \textit{Dungeon Master's Guide (2014)} by} Jeremy Crawford (co-lead), Christopher Perkins (co-lead), James Wyatt (co-lead), Peter Lee, Mike Mearls, Robert J. Schwalb, Rodney Thompson
\item \textbf{Building on the original game created by} Gary Gygax and Dave Arneson and then developed by many others over the past 50 years
\end{description}

\end{multicols}

\end{description}

\chapter{Magic Items}
\DndItemHeader{+1 Wand of the War Mage}{Uncommon (requires attunement by a spellcaster)}
While holding this wand, you gain a +1 bonus to spell attack rolls. In addition, you ignore Cover when making a spell attack roll.

\DndItemHeader{+1 Wraps of Unarmed Power}{Wondrous item, Uncommon}
While wearing these wraps, you have a +1 bonus to attack rolls and damage rolls made with your Unarmed Strikes. Those strikes deal your choice of Force damage or their normal damage type.

\DndItemHeader{+2 Rod of the Pact Keeper}{Rare (requires attunement by a warlock)}
While holding this rod, you gain a +2 bonus to spell attack rolls and to the saving throw DCs of your Warlock spells.

In addition, you can regain one spell slot as a Magic action while holding the rod. You can't use this property again until you finish a Long Rest.

\DndItemHeader{+2 Wand of the War Mage}{Rare (requires attunement by a spellcaster)}
While holding this wand, you gain a +2 bonus to spell attack rolls. In addition, you ignore Cover when making a spell attack roll.

\DndItemHeader{+2 Wraps of Unarmed Power}{Wondrous item, Rare}
While wearing these wraps, you have a +2 bonus to attack rolls and damage rolls made with your Unarmed Strikes. Those strikes deal your choice of Force damage or their normal damage type.

\DndItemHeader{+3 Rod of the Pact Keeper}{Very Rare (requires attunement by a warlock)}
While holding this rod, you gain a +3 bonus to spell attack rolls and to the saving throw DCs of your Warlock spells.

In addition, you can regain one spell slot as a Magic action while holding the rod. You can't use this property again until you finish a Long Rest.

\DndItemHeader{+3 Wand of the War Mage}{Very Rare (requires attunement by a spellcaster)}
While holding this wand, you gain a +3 bonus to spell attack rolls. In addition, you ignore Cover when making a spell attack roll.

\DndItemHeader{+3 Wraps of Unarmed Power}{Wondrous item, Very Rare}
While wearing these wraps, you have a +3 bonus to attack rolls and damage rolls made with your Unarmed Strikes. Those strikes deal your choice of Force damage or their normal damage type.

\DndItemHeader{Alchemist's Supplies}{}

\begin{description}
\item[Ability:]
Intelligence

\item[Utilize:]
Identify a substance (DC 15), or start a fire (DC 15)

\item[Craft:]
\textit{Acid}, \textit{Alchemist's Fire}, \textit{Component Pouch}, \textit{Oil}, \textit{Paper}, \textit{Perfume}

\end{description}

\DndItemHeader{Alchemy Jug}{Wondrous item, Uncommon}
This ceramic jug appears to be able to hold a gallon of liquid and weighs 12 pounds whether full or empty. The jug sloshes when it is shaken, even if the jug is empty.

You can take a Magic action and name one liquid from the Alchemy Jug Liquids table to cause the jug to produce the chosen liquid. Afterward, you can uncork the jug as a Utilize action and pour that liquid out, up to 2 gallons per minute. The maximum amount of liquid the jug can produce depends on the liquid you named.

Once the jug starts producing a liquid, it can't produce a different one, or more of one that has reached its maximum, until the next dawn.

\begin{DndTable}[header=Alchemy Jug Liquids]{XX}

Liquid & Max. Amount\\
Acid & 8 ounces\\
Basic Poison & 4 ounces\\
Beer & 4 gallons\\
Honey & 1 gallon\\
Mayonnaise & 2 gallons\\
Oil & 1 quart\\
Vinegar & 2 gallons\\
Water, fresh & 8 gallons\\
Water, salt & 12 gallons\\
Wine & 1 gallon\\
\end{DndTable}

\DndItemHeader{Alexandrite}{}
A dark green gemstone.

\DndItemHeader{Amber}{}
A watery gold to rich gold gemstone.

\DndItemHeader{Amethyst}{}
A deep purple gemstone.

\DndItemHeader{Amulet of Health}{Wondrous item, Rare (requires attunement)}
Your Constitution score is 19 while you wear this amulet. It has no effect on you if your Constitution is already 19 or higher without it.

\DndItemHeader{Amulet of Proof against Detection and Location}{Wondrous item, Uncommon (requires attunement)}
While wearing this amulet, you can't be targeted by Divination spells or perceived through magical scrying sensors unless you allow it.

\DndItemHeader{Amulet of Proof against Detection and Location}{Wondrous item, Uncommon (requires attunement)}
While wearing this amulet, you can't be targeted by Divination spells or perceived through magical scrying sensors unless you allow it.

\DndItemHeader{Amulet of the Planes}{Wondrous item, Very Rare (requires attunement)}
While wearing this amulet, you can take a Magic action to name a location that you are familiar with on another plane of existence. Then make a DC 15 Intelligence (Arcana) check. On a successful check, you cast \textit{Plane Shift}. On a failed check, you and each creature and object within 15 feet of you travel to a random destination determined by rolling 1d100 and consulting the following table.


\begin{DndTable}{cX}

1d100 & Destination\\
01-60 & Random location on the plane you named\\
61-70 & Random location on an Inner Plane determined by rolling 1d6: on a \textbf{1}, the Plane of Air; on a \textbf{2}, the Plane of Earth; on a \textbf{3}, the Plane of Fire; on a \textbf{4}, the Plane of Water; on a \textbf{5}, the Feywild; on a \textbf{6}, the Shadowfell\\
71-80 & Random location on an Outer Plane determined by rolling 1d8: on a \textbf{1}, Arborea; on a \textbf{2}, Arcadia; on a \textbf{3}, [Tooltip Not Found]; on a \textbf{4}, Bytopia; on a \textbf{5}, Elysium; on a \textbf{6}, Mechanus; on a \textbf{7}, Mount Celestia; on an \textbf{8}, Ysgard\\
81-90 & Random location on an Outer Plane determined by rolling 1d8: on a \textbf{1}, [Tooltip Not Found]; on a \textbf{2}, Acheron; on a \textbf{3}, Carceri; on a \textbf{4}, Gehenna; on a \textbf{5}, Hades; on a \textbf{6}, Limbo; on a \textbf{7}, [Tooltip Not Found]; on an \textbf{8}, Pandemonium\\
91-00 & Random location on the Astral Plane\\
\end{DndTable}

\DndItemHeader{Animated Shield}{Very Rare (requires attunement)}
While holding this Shield, you can take a Bonus Action to cause it to animate. The Shield leaps into the air and hovers in your space to protect you as if you were wielding it, leaving your hands free. The Shield remains animate for 1 minute, until you take a Bonus Action to end this effect, or until you die or have the Incapacitated condition, at which point the Shield falls to the ground or into your hand if you have one free.

\DndItemHeader{Apparatus of Kwalish}{Wondrous item, Legendary}
This item first appears to be a sealed iron barrel weighing 500 pounds. The barrel has a hidden catch, which can be found with a successful DC 20 Intelligence (Investigation) check. Releasing the catch unlocks a hatch at one end of the barrel, allowing two Medium or smaller creatures to crawl inside. Ten levers are set in a row at the far end, each in a neutral position, able to move up or down. When certain levers are used, the apparatus transforms to resemble a giant lobster.

The Apparatus of Kwalish is a Large object with the following statistics: \textbf{AC} 20; \textbf{HP} 200; \textbf{Speed} 30 ft., Swim 30 ft. (or 0 ft. for both if the legs aren't extended); \textbf{Immunity} to Poison and Psychic damage.

To be used as a vehicle, the apparatus requires one pilot. While the apparatus's hatch is closed, the compartment is airtight and watertight. The compartment holds enough air for 10 hours of breathing, divided by the number of breathing creatures inside.

The apparatus floats on water. It can also go underwater to a depth of 900 feet. Below that, the vehicle takes 2d6 Bludgeoning damage each minute from pressure.

A creature in the compartment can take a Utilize action to move as many as two of the apparatus's levers up or down. After each use, a lever goes back to its neutral position. Each lever, from left to right, functions as shown in the Apparatus of Kwalish Levers table.

\begin{DndTable}[header=Apparatus of Kwalish Levers]{cXX}

Lever & Up & Down\\
1 & Legs extend, allowing the apparatus to walk and swim. & Legs retract, reducing the apparatus's Speed and Swim Speed to 0 and making it unable to benefit from bonuses to speed.\\
2 & Forward window shutter opens. & Forward window shutter closes.\\
3 & Side window shutters open (two per side). & Side window shutters close (two per side).\\
4 & Two claws extend from the front side of the apparatus. & The claws retract.\\
5 & Each extended claw makes the following melee attack: +8 to hit, reach 5 ft. \textit{Hit:} 7 (2d6) Bludgeoning damage. & Each extended claw makes the following melee attack: +8 to hit, reach 5 ft. \textit{Hit:} The target has the Grappled condition (escape DC 15).\\
6 & The apparatus walks or swims forward provided its legs are extended. & The apparatus walks or swims backward provided its legs are extended.\\
7 & The apparatus turns 90 degrees counterclockwise provided its legs are extended. & The apparatus turns 90 degrees clockwise provided its legs are extended.\\
8 & Eyelike fixtures emit Bright Light in a 30-foot radius and Dim Light for an additional 30 feet. & The light turns off.\\
9 & The apparatus sinks up to 20 feet if it's in liquid. & The apparatus rises up to 20 feet if it's in liquid.\\
10 & The rear hatch unseals and opens. & The rear hatch closes and seals.\\
\end{DndTable}

\DndItemHeader{Aquamarine}{}
A pale blue green gemstone.

\DndItemHeader{Armor of Invulnerability}{Legendary (requires attunement)}
You have Resistance to Bludgeoning, Piercing, and Slashing damage while you wear this armor.

\subparagraph{Metal Shell}
You can take a Magic action to give yourself Immunity to Bludgeoning, Piercing, and Slashing damage for 10 minutes or until you are no longer wearing the armor. Once this property is used, it can't be used again until the next dawn.

\DndItemHeader{Armor of Resistance}{Rare (requires attunement)}
You have Resistance to {{getFullImmRes item.resist}} damage while you wear this armor.

\DndItemHeader{Arrow-Catching Shield}{Rare (requires attunement)}
You gain a +2 bonus to Armor Class against ranged attack rolls while you wield this Shield. This bonus is in addition to the Shield's normal bonus to AC.

Whenever an attacker makes a ranged attack roll against a target within 5 feet of you, you can take a Reaction to become the target of the attack instead.

\DndItemHeader{Assassin's Blood}{}
A creature subjected to Assassin's Blood makes a DC 10 Constitution saving throw. On a failed save, the creature takes 6 (1d12) Poison damage and has the Poisoned condition for 24 hours. On a successful save, the creature takes half as much damage only.

\DndItemHeader{Axe of the Dwarvish Lords}{Artifact (requires attunement)}
A young dwarf prince set out to forge a weapon that would be regarded as a symbol of unity among his people. Venturing deep under the mountains, deeper than any dwarf had ever delved, the prince came to the blazing heart of a great volcano. With the aid of Moradin, a god of creation, he first crafted four mighty tools: the Starmetal Pick, the Earthheart Forge, the Anvil of Songs, and the Shaping Hammer. With these tools, he forged the Axe of the Dwarvish Lords.

Armed with the Artifact, the prince brought peace to the dwarf clans, ending grudges and answering slights. The clans became allies, and they threw back their enemies and enjoyed an era of prosperity. This young dwarf is remembered as the First King. When he became old, he passed the weapon, which had become his badge of office, to his heir. The rightful inheritors passed the axe on for many generations.

Later, in an era marked by treachery and wickedness, the axe was lost in a bloody civil war fomented by greed for its power and the status it bestowed. Centuries later, the dwarves still search for the axe, and many adventurers have made careers of chasing after rumors and plundering old vaults to find it.

\subparagraph{Magic Weapon}
The Axe of the Dwarvish Lords is a magic weapon that grants a +3 bonus to attack rolls and damage rolls made with it.

When you attack a creature with the axe and roll a 20 on the d20 for the attack roll, the axe deals an extra 20 Slashing damage.

The axe has T with a normal range of 20 feet and a long range of 60 feet. When you hit with a ranged attack using this weapon, it deals an extra 1d8 Force damage, or an extra 2d8 Force damage if the target is a creature of the Giant type. Immediately after hitting or missing, the weapon flies back to your hand.

\subparagraph{Blessings of Moradin}
While attuned to the axe, you gain the following benefits:


\begin{description}
\item[Darkvision.]
You gain Darkvision with a range of 60 feet. If you already have Darkvision, its range increases by 60 feet.

\item[Fortitude of Stone.]
Your Constitution increases by 2, to a maximum of 20.

\item[Gifts of the Creator.]
You have proficiency with \textit{Brewer's Supplies}, \textit{Mason's Tools}, and \textit{Smith's Tools}.

\item[One with the Forge.]
You have Immunity to Poison damage and Resistance to Fire damage.

\item[Sunder.]
When you hit an object with the axe, the object takes the maximum amount of damage possible.

\end{description}

\subparagraph{Conjure Earth Elemental}
While holding the axe, you can take a Magic action to summon an \textbf{Earth Elemental}. It appears in an unoccupied space you choose within 30 feet of yourself, understands your languages, obeys your commands, and takes its turn immediately after you on your Initiative count. The elemental disappears after 24 hours, when it dies, or when you dismiss it as a Bonus Action. You can't use this property again until the next dawn.

\subparagraph{Random Properties}
The axe has the following random properties:


\begin{itemize}
\item 2 Artifact Properties; Minor Beneficial Properties properties
\item 1 Artifact Properties; Major Beneficial Properties property
\item 2 Artifact Properties; Minor Detrimental Properties properties
\end{itemize}

\subparagraph{Travel the Depths}
You can take a Magic action to touch the axe to a fixed piece of dwarven stonework and cast \textit{Teleport} from the axe. If your intended destination is underground, there is no chance of a mishap or arriving somewhere unexpected. You can't use this property again until 3 days have passed.

\subparagraph{Destroying the Axe}
The only way to destroy the axe is to melt it down in the Earthheart Forge, where it was created. It must remain in the burning forge for 50 years before it finally succumbs to the fire and is consumed.

\DndItemHeader{Azurite}{}
A mottled deep blue gemstone.

\DndItemHeader{Baba Yaga's Dancing Broom}{Wondrous item, Uncommon (requires attunement)}
The archfey \textit{Baba Yaga} crafted many of these magic brooms. No two appear exactly alike. While holding the broom, you can take a Magic action to transform it into an \textbf{Animated Broom} under your control. The broom then moves into an unoccupied space as close to you as possible. The broom acts immediately after you on your Initiative count and remains animate until you take a Bonus Action and use a command word to render it inanimate.

On your turn, you can mentally command the animated broom if it is within 30 feet of you and you don't have the Incapacitated condition (no action required). You decide what action the broom takes and where it moves during its next turn, or you can issue it a general command, such as to attack your enemies or guard a location.

If the broom is reduced to 0 Hit Points, it shatters and is destroyed. If the broom reverts to its inanimate form before losing all its Hit Points it regains all of them.

\DndItemHeader{Bag of Beans}{Wondrous item, Rare}
This heavy cloth bag contains 3d4 dry beans when found. The bag weighs half a pound regardless of how many beans it contains and becomes a nonmagical item when it no longer contains any beans.

If you dump one or more beans out of the bag, they explode in a 10-foot-radius Sphere [Area of Effect] centered on them. All the dumped beans are destroyed in the explosion, and each creature in the Sphere [Area of Effect], including you, makes a DC 15 Dexterity saving throw, taking 5d4 Force damage on a failed save or half as much damage on a successful one.

If you remove a bean from the bag, plant it in dirt or sand, and then water it, the bean disappears as it produces an effect 1 minute later from the ground where it was planted. The DM can choose an effect from the following table or determine it randomly.


\begin{DndTable}{cX}

1d100 & Effect\\
01 & 5d4 toadstools sprout. If a creature eats a toadstool, roll any die. On an odd roll, the eater must succeed on a DC 15 Constitution saving throw or take 5d6 Poison damage and have the Poisoned condition for 1 hour. On an even roll, the eater gains 5d6 Temporary Hit Points for 1 hour.\\
02-10 & A geyser erupts and spouts water, beer, mayonnaise, tea, vinegar, wine, or oil (DM's choice) 30 feet into the air for 1d4 minutes.\\
11-20 & A \textbf{Treant} sprouts. Roll any die. On an odd roll, the treant is Chaotic Evil. On an even roll, the treant is Chaotic Good.\\
21-30 & An animate but immobile stone statue in your likeness rises and makes verbal threats against you. If you leave it and others come near, it describes you as the most heinous of villains and directs the newcomers to find and attack you. If you are on the same plane of existence as the statue, it knows where you are. The statue becomes inanimate after 24 hours.\\
31-40 & A campfire with green flames springs forth and burns for 24 hours or until it is extinguished.\\
41-50 & Three \textbf{Shrieker Fungi} sprout.\\
51-60 & 1d4 + 4 bright-pink toads crawl forth. Whenever a toad is touched, it transforms into a Large or smaller monster of the DM's choice that acts in accordance with its alignment and nature. The monster remains for 1 minute, then disappears in a puff of bright-pink smoke.\\
61-70 & A hungry \textbf{Bulette} burrows up and attacks.\\
71-80 & A fruit tree grows. It has 1d10 + 20 fruit, 1d8 of which act as randomly determined potions. The tree vanishes after 1 hour. Picked fruit remains, retaining any magic for 30 days.\\
81-90 & A nest of 1d4 + 3 rainbow-colored eggs springs up. Any creature that eats an egg makes a DC 20 Constitution saving throw. On a successful save, a creature permanently increases its lowest ability score by 1, randomly choosing among equally low scores. On a failed save, the creature takes 10d6 Force damage from an internal explosion.\\
91-95 & A pyramid with a 60-foot-square base bursts upward. Inside is a burial chamber containing a \textbf{Mummy}, a \textbf{Mummy Lord}, or some other Undead of the DM's choice. Its sarcophagus contains treasure of the DM's choice.\\
96-00 & A giant beanstalk sprouts, growing to a height of the DM's choice. The top leads where the DM chooses, such as to a great view, a cloud giant's castle, or another plane of existence.\\
\end{DndTable}

\DndItemHeader{Bag of Devouring}{Wondrous item, Very Rare}
This bag resembles a \textit{Bag of Holding} but is a feeding orifice for a gigantic extradimensional creature. Turning the bag inside out closes the orifice.

The extradimensional creature attached to the bag can sense whatever is placed inside the bag. Animal or vegetable matter placed wholly in the bag is devoured and lost forever. When part of a living creature is placed in the bag, as happens when someone reaches inside it, there is a 50 chance that the creature is pulled inside the bag. A creature inside the bag can take an action to try to escape, doing so with a successful DC 15 Strength (Athletics) check. Another creature can take an action to reach into the bag to pull a creature out, doing so with a successful DC 20 Strength (Athletics) check, provided the puller isn't pulled inside the bag first. Any creature that starts its turn inside the bag is devoured, its body destroyed.

Inanimate objects can be stored in the bag, which can hold a cubic foot of such material. However, once each day, the bag swallows any objects inside it and spits them out into another plane of existence. The DM determines the time and plane.

If the bag is pierced or torn, it is destroyed, and anything contained within it is transported to a random location on the Astral Plane.

\DndItemHeader{Bag of Holding}{Wondrous item, Uncommon}
This bag has an interior space considerably larger than its outside dimensions—roughly 2 feet square and 4 feet deep on the inside. The bag can hold up to 500 pounds, not exceeding a volume of 64 cubic feet. The bag weighs 5 pounds, regardless of its contents. Retrieving an item from the bag requires a Utilize action.

If the bag is overloaded, pierced, or torn, it is destroyed, and its contents are scattered in the Astral Plane. If the bag is turned inside out, its contents spill forth unharmed, but the bag must be put right before it can be used again. The bag holds enough air for 10 minutes of breathing, divided by the number of breathing creatures inside.

Placing a Bag of Holding inside an extradimensional space created by a \textit{Heward's Handy Haversack}, Portable Hole, or similar item instantly destroys both items and opens a gate to the Astral Plane. The gate originates where the one item was placed inside the other. Any creature within a 10-foot-radius Sphere [Area of Effect] centered on the gate is sucked through it to a random location on the Astral Plane. The gate then closes. The gate is one-way and can't be reopened.

\DndItemHeader{Banded Agate}{}
A striped brown, blue, white, or red gemstone.

\DndItemHeader{Bead of Force}{Wondrous item, Rare}
This small black sphere measures 3/4 of an inch in diameter and weighs an ounce. Typically, 1d4 + 4 Beads of Force are found together.

You can take a Magic action to throw the bead up to 60 feet. The bead explodes in a 10-foot-radius Sphere [Area of Effect] on impact and is destroyed. Each creature in the Sphere [Area of Effect] must succeed on a DC 15 Dexterity saving throw or take 5d4 Force damage. A sphere of transparent force then encloses the area for 1 minute. Any creature that failed the save and is completely within the area is trapped inside this sphere. Creatures that succeeded on the save or are partially within the area are pushed away from the center of the sphere until they are no longer inside it. Only breathable air can pass through the sphere's wall. No attack or other effect can pass through.

An enclosed creature can take a Utilize action to push against the sphere's wall, moving the sphere up to half the creature's Speed. The sphere can be picked up, and its magic causes it to weigh only 1 pound, regardless of the weight of creatures inside.

\DndItemHeader{Bead of Nourishment}{Wondrous item, Common}
This flavorless, gelatinous bead dissolves on your tongue and provides as much nourishment as 1 day of Rations.

\DndItemHeader{Bead of Refreshment}{Wondrous item, Common}
This flavorless, gelatinous bead dissolves in liquid, transforming up to a pint of the liquid into fresh, cold drinking water. The bead has no effect on magical liquids or harmful substances such as poison.

\DndItemHeader{Belt of Cloud Giant Strength}{Wondrous item, Rare (requires attunement)}
While wearing this belt, your Strength score changes to 27. The item has no effect on you if your Strength without the belt is equal to or greater than the belt's score.

\DndItemHeader{Belt of Dwarvenkind}{Wondrous item, Rare (requires attunement)}
While wearing this belt, you gain the following benefits:


\begin{description}
\item[Dwarvish.]
You know Dwarvish.

\item[Friend of Dwarvenkind.]
You have Advantage on Charisma (Persuasion) checks made to interact with dwarves and duergar.

\item[Toughness.]
Your Constitution increases by 2, to a maximum of 20.

In addition, while attuned to the belt, you have a 50 chance each day at dawn of growing a full beard if you can grow one, or a thicker beard if you already have one.

If you aren't a dwarf or duergar, you gain the following additional benefits while wearing the belt:

\item[Darkvision.]
You have Darkvision with a range of 60 feet.

\item[Resilience.]
You have Resistance to Poison damage. You also have Advantage on saving throws you make to avoid or end the Poisoned condition.

\end{description}

\DndItemHeader{Belt of Fire Giant Strength}{Wondrous item, Very Rare (requires attunement)}
While wearing this belt, your Strength score changes to 25. The item has no effect on you if your Strength without the belt is equal to or greater than the belt's score.

\DndItemHeader{Belt of Frost Giant Strength}{Wondrous item, Very Rare (requires attunement)}
While wearing this belt, your Strength score changes to 23. The item has no effect on you if your Strength without the belt is equal to or greater than the belt's score.

\DndItemHeader{Belt of Hill Giant Strength}{Wondrous item, Very Rare (requires attunement)}
While wearing this belt, your Strength score changes to 21. The item has no effect on you if your Strength without the belt is equal to or greater than the belt's score.

\DndItemHeader{Belt of Stone Giant Strength}{Wondrous item, Legendary (requires attunement)}
While wearing this belt, your Strength score changes to 23. The item has no effect on you if your Strength without the belt is equal to or greater than the belt's score.

\DndItemHeader{Belt of Storm Giant Strength}{Wondrous item, Legendary (requires attunement)}
While wearing this belt, your Strength score changes to 29. The item has no effect on you if your Strength without the belt is equal to or greater than the belt's score.

\DndItemHeader{Black Opal}{}
A dark green with black mottling and golden flecks gemstone.

\DndItemHeader{Black Pearl}{}
A pure black gemstone.

\DndItemHeader{Black Sapphire}{}
A lustrous black with glowing highlights gemstone.

\DndItemHeader{Blackrazor}{Artifact (requires attunement)}
Hidden in the dungeon of White Plume Mountain, Blackrazor shines like a piece of night sky filled with stars. Its black scabbard is decorated with pieces of cut obsidian.

You gain a +3 bonus to attack rolls and damage rolls made with this magic weapon. If you hit an Undead with this weapon, you take 1d10 Necrotic damage, and the target regains 1d10 Hit Points. If this Necrotic damage reduces you to 0 Hit Points, Blackrazor devours your soul (see "Devour Soul" below).

While you hold this weapon, you have Immunity to the Charmed and Frightened conditions, and you have Blindsight with a range of 30 feet.

\subparagraph{Devour Soul}
Whenever you use Blackrazor to reduce a creature to 0 Hit Points, the sword slays the creature and devours its soul unless it is a Construct or an Undead. A creature whose soul has been devoured by Blackrazor can be restored to life only by a \textit{Wish} spell.

When Blackrazor devours a soul that isn't yours, you gain Temporary Hit Points equal to the slain creature's Hit Points maximum.

\subparagraph{Haste}
Blackrazor can cast \textit{Haste} on you, after which it can't cast this spell again until the next dawn. Blackrazor decides when to cast the spell, which takes effect at the start of your turn. The spell lasts for 1 minute (no Concentration required) or until Blackrazor decides to end it, which it can do at the end of any of your turns.

\subparagraph{Sentience}
Blackrazor is a sentient Chaotic Neutral weapon with an Intelligence of 17, a Wisdom of 10, and a Charisma of 19. It has hearing and Darkvision out to 120 feet.

The weapon speaks Common and can communicate with its wielder telepathically. Its voice is deep and echoing. While you are attuned to it, Blackrazor also understands every language you know.

\subparagraph{Personality}
Blackrazor speaks with an imperious tone, as though accustomed to being obeyed.The sword's purpose is to consume souls. It doesn't care whose souls it eats, including the wielder's. The sword believes that all matter and energy sprang from a void of negative energy and will one day return to it. Blackrazor is meant to hurry that process along.

Despite its nihilism, Blackrazor feels a strange kinship to Wave and Whelm, two other weapons locked away under White Plume Mountain. It wants the three weapons to be reunited and wielded together in combat, even though it violently disagrees with Whelm and finds Wave tedious.

Blackrazor's hunger for souls must be regularly fed. If the sword goes 3 days or more without consuming a soul, a conflict between it and its wielder occurs at the next sunset.

\subparagraph{Destroying Blackrazor}
Blackrazor can be destroyed by crushing it in the great gears of Mechanus. Primus, the creator of the modrons, also knows a series of musical tones that Blackrazor can't stand to hear, causing the sword to shatter.

\DndItemHeader{Bloodstone}{}
A dark gray with red flecks gemstone.

\DndItemHeader{Blue Quartz}{}
A pale blue gemstone.

\DndItemHeader{Blue Sapphire}{}
A medium blue gemstone.

\DndItemHeader{Blue Spinel}{}
A deep blue gemstone.

\DndItemHeader{Bomb}{}
As an action, you can light a Bomb and throw it at a point up to 60 feet away, where it explodes. Each creature in a 5-foot-radius Sphere [Area of Effect] centered on that point makes a DC 12 Dexterity saving throw, taking 3d6 Fire damage on a failed save or half as much damage on a successful one.

\DndItemHeader{Book of Exalted Deeds}{Wondrous item, Artifact (requires attunement)}
The definitive treatise on all that is good in the multiverse, the Book of Exalted Deeds figures prominently in many religions. Rather than being a scripture devoted to a particular faith, the book's authors filled the pages with their own visions of true virtue, providing guidance for defeating evil.

The Book of Exalted Deeds rarely lingers in one place. As soon as the book is read, it vanishes to some other corner of the multiverse where its moral guidance can bring hope to an endangered world. Although attempts have been made to copy the work, efforts to do so fail to capture its magical nature or translate the benefits it offers to those pure of heart and firm of purpose.

A heavy clasp, wrought to look like angel wings, keeps the book's contents secure. Only a creature that is attuned to the book can release the clasp that holds it shut. Once the book is opened, the attuned creature must spend 80 hours reading and studying the book to digest its contents and gain its benefits.

Other creatures that peruse the book's open pages can read the text but glean no deeper meaning and reap no benefits. A Fiend, an Undead, or a servant of a god from the Lower Planes that tries to read from the book takes 24d6 Radiant damage. This damage ignores Resistance and Immunity, and it can't be reduced or avoided by any means. A creature reduced to 0 Hit Points by this damage disappears in a flash and is destroyed, leaving its possessions behind. The book then vanishes, and the creature's Attunement to it ends.

Benefits granted by the Book of Exalted Deeds last only as long as you strive to do good. If you fail to perform at least one act of kindness or generosity within the span of 10 days, or if you willingly perform an evil act, you lose all the benefits granted by the book.

\subparagraph{Celestial Calm}
While attuned to the book, you have Immunity to the Charmed and Frightened conditions and Resistance to Psychic damage. These benefits become permanent after you spend the requisite amount of time reading and studying the book.

\subparagraph{Divine Wisdom}
After you spend the requisite amount of time reading and studying the book, your Wisdom increases by 2, to a maximum of 24. You can't gain this benefit from the book more than once.

\subparagraph{Enlightened Magic}
After you spend the requisite amount of time reading and studying the book, any spell slot you expend to cast a spell counts as a spell slot of one level higher.

\subparagraph{Halo}
After you spend the requisite amount of time reading and studying the book, you gain a protective halo. This halo sheds Bright Light in a 10-foot radius and Dim Light for an additional 10 feet. You can dismiss or manifest the halo as a Bonus Action. While present, the halo gives you Advantage on Charisma (Persuasion) checks. In addition, Fiends and Undead within the halo's Bright Light make attack rolls against you with Disadvantage.

\subparagraph{Random Properties}
The Book of Exalted Deeds has the following random properties:


\begin{itemize}
\item 2 Artifact Properties; Minor Beneficial Properties properties
\item 2 Artifact Properties; Major Beneficial Properties properties
\end{itemize}

\subparagraph{Destroying the Book}
The Book of Exalted Deeds can't be destroyed. However, drowning the book in the River Styx removes all writing and imagery from its pages and renders the book powerless for 1d100 years.

\DndItemHeader{Book of Vile Darkness}{Wondrous item, Artifact (requires attunement)}
The contents of this foul manuscript are the meat and drink of the wicked. It contains knowledge so horrid that to even glimpse the scrawled pages invites doom.

Most believe the lich-god \textit{Vecna} authored the Book of Vile Darkness. He recorded in its pages every horrid idea, every corrupt thought, and every example of foul magic he came across or devised.

Other practitioners of evil have added their own input to the book's catalog of vile knowledge. Their additions are clear, for the writers of later works stitched whatever they were writing into the tome or, in some cases, made notations and additions to existing text. There are places where pages are missing, torn, or covered so completely with ink, blood, and scratches that the original text can't be divined.

Nature can't abide the book's presence. Ordinary plants wither in its presence, common animals are unwilling to approach it, and the book gradually destroys whatever it touches. Even stone cracks and turns to powder if the book rests on it long enough.

Whenever a creature that isn't a Fiend or an Undead attunes to the Book of Vile Darkness, that creature makes a DC 17 Charisma saving throw. On a failed save, the creature is magically transformed into a \textbf{Larva} under the DM's control. Only a \textit{Wish} spell can reverse this vile transformation.

A creature attuned to the book must spend 80 hours reading and studying it to digest its contents and use its Adjusted Ability Scores, Tireless Form, Spells, Vile Lore, and Vile Speech properties.

The Book of Vile Darkness remains with you only as long as you strive to work evil in the world. If you fail to perform at least one evil act within the span of 10 days, or if you willingly perform a good act, the book disappears, your Attunement to it ends immediately, and you lose all benefits granted by it. If you die while attuned to the book, an entity of great evil claims your soul. You can't be restored to life by any means while your soul remains imprisoned.

\subparagraph{Adjusted Ability Scores}
One ability score of your choice increases by 2, to a maximum of 24. Another ability score of your choice decreases by 2, to a minimum of 3. The book can't adjust your ability scores again.

\subparagraph{Tireless Form}
While the book is on your person, you have Immunity to the Exhaustion condition.

\subparagraph{Random Properties}
The Book of Vile Darkness has the following random properties:


\begin{itemize}
\item 3 Artifact Properties; Minor Beneficial Properties properties
\item 1 Artifact Properties; Major Beneficial Properties property
\item 3 Artifact Properties; Minor Detrimental Properties properties
\item 2 Artifact Properties; Major Detrimental Properties properties
\end{itemize}

\subparagraph{Spells}
While holding the book and holding it, you can cast the following spells (save DC 18) from it:


\begin{itemize}
\item \textit{Animate Dead}
\item \textit{Circle of Death}
\item \textit{Dominate Monster}
\item \textit{Finger of Death}
\end{itemize}

Once you use the book to cast a spell, you can't cast that spell again from it until the next dawn.

\subparagraph{Vile Lore}
You can reference the Book of Vile Darkness whenever you make an Intelligence check to recall information about some aspect of evil, such as lore about demons. When you do so, you have Advantage on that check.

At the DM's discretion, the book might reveal secrets no mortal should know, such as the true names of powerful Fiends, foul rites that allow one to transform into a death knight or lich, or long-lost spells crafted by beings so evil their names ought never to be spoken aloud.

\subparagraph{Vile Speech}
While the book is on your person, you can take a Magic action to recite words from its pages in a foul, dead language. Each time you do so, you take 1d12 Psychic damage, and each creature within 15 feet of you takes 3d6 Psychic damage unless the creature is a Fiend or an Undead.

\subparagraph{Destroying the Book}
The Book of Vile Darkness allows pages to be torn from it, but any evil lore contained on those pages finds its way back into the book eventually, usually when a new author adds pages to the tome.

If a \textbf{solar} tears the book in two, the book is destroyed for 1d100 years, after which it re-forms in some far corner of the multiverse.

A creature attuned to the book for 100 years can unearth a phrase hidden in the original text that, when translated to Celestial and spoken aloud, destroys both the speaker and the book in a flash of radiance. However, as long as evil exists in the multiverse, the book re-forms 1d10 × 100 years later.

\DndItemHeader{Boots of Elvenkind}{Wondrous item, Uncommon}
While you wear these boots, your steps make no sound, regardless of the surface you are moving across. You also have Advantage on Dexterity (Stealth) checks.

\DndItemHeader{Boots of False Tracks}{Wondrous item, Common}
While wearing these boots, you can have them leave tracks like those of any kind of Humanoid of your size.

\DndItemHeader{Boots of Levitation}{Wondrous item, Rare (requires attunement)}
While you wear these boots, you can cast \textit{Levitate} on yourself.

\DndItemHeader{Boots of Speed}{Wondrous item, Rare (requires attunement)}
While you wear these boots, you can take a Bonus Action to click the boots' heels together. If you do, the boots double your Speed, and any creature that makes an Opportunity Attack against you has Disadvantage on the attack roll. If you click your heels together again, you end the effect.

When you've used the boots' property for a total of 10 minutes, the magic ceases to function for you until you finish a Long Rest.

\DndItemHeader{Boots of Striding and Springing}{Wondrous item, Uncommon (requires attunement)}
While you wear these boots, your Speed becomes 30 feet unless your Speed is higher, and your Speed isn't reduced by you carrying weight in excess of your carrying capacity or wearing Heavy Armor.

Once on each of your turns, you can jump up to 30 feet by spending only 10 feet of movement.

\DndItemHeader{Boots of the Winterlands}{Wondrous item, Uncommon (requires attunement)}
These furred boots are snug and feel warm. While wearing them, you gain the following benefits.

\subparagraph{Cold Resistance}
You have Resistance to Cold damage and can tolerate temperatures of 0 degrees Fahrenheit or lower without any additional protection.

\subparagraph{Winter Strider}
You ignore Difficult Terrain created by ice or snow.

\DndItemHeader{Bowl of Commanding Water Elementals}{Wondrous item, Rare}
While this bowl is filled with water and you are within 5 feet of it, you can take a Magic action to summon a \textbf{Water Elemental}. The elemental appears in an unoccupied space as close to the bowl as possible, understands your languages, obeys your commands, and takes its turn immediately after you on your Initiative count. The elemental disappears after 1 hour, when it dies, or when you dismiss it as a Bonus Action. The bowl can't be used this way again until the next dawn.

The bowl is about 1 foot in diameter and half as deep. It holds about 3 gallons.

\DndItemHeader{Bracers of Archery}{Wondrous item, Uncommon (requires attunement)}
While wearing these bracers, you have proficiency with the Longbow and Shortbow, and you gain a +2 bonus to damage rolls made with such weapons.

\DndItemHeader{Bracers of Defense}{Wondrous item, Rare (requires attunement)}
While wearing these bracers, you gain a +2 bonus to Armor Class if you are wearing no armor and using no Shield.

\DndItemHeader{Brazier of Commanding Fire Elementals}{Wondrous item, Rare}
While you are within 5 feet of this brazier, you can take a Magic action to summon a \textbf{Fire Elemental}. The elemental appears in an unoccupied space as close to the brazier as possible, understands your languages, obeys your commands, and takes its turn immediately after you on your Initiative count. The elemental disappears after 1 hour, when it dies, or when you dismiss it as a Bonus Action. The brazier can't be used this way again until the next dawn.

\DndItemHeader{Brooch of Shielding}{Wondrous item, Uncommon (requires attunement)}
While wearing this brooch, you have Resistance to Force damage, and you have Immunity to damage from the \textit{Magic Missile} spell.

\DndItemHeader{Broom of Flying}{Wondrous item, Uncommon (requires attunement)}
This wooden broom functions like a mundane broom until you stand astride it and take a Magic action to make it hover beneath you, at which time it can be ridden in the air. It has a Fly Speed of 50 feet. It can carry up to 400 pounds, but its Fly Speed becomes 30 feet while carrying over 200 pounds. The broom stops hovering when you land or when you're no longer riding it.

As a Magic action, you can send the broom to travel alone to a destination within 1 mile of you if you name the location and are familiar with it. The broom comes back to you when you take a Magic action and use a command word if the broom is still within 1 mile of you.

\DndItemHeader{Burnt Othur Fumes}{}
A creature subjected to Burnt Othur Fumes must succeed on a DC 13 Constitution saving throw or take 10 (3d6) Poison damage, and it must repeat the save at the start of each of its turns. On each successive failed save, the creature takes 3 (1d6) Poison damage. After three successful saves, the poison ends.

\DndItemHeader{Calligrapher's Supplies}{}

\begin{description}
\item[Ability:]
Dexterity

\item[Utilize:]
Write text with impressive flourishes that guard against forgery (DC 15)

\item[Craft:]
\textit{Ink}, \textit{Spell Scroll}

\end{description}

\DndItemHeader{Candle of Invocation}{Wondrous item, Very Rare (requires attunement)}
This candle's magic is activated when the candle is lit, which requires a Magic action. After burning for 4 hours, the candle is destroyed. You can snuff it out early for use at a later time. Deduct the time it burned in increments of 1 minute from its total burn time.

While lit, the candle sheds Dim Light in a 30-foot radius. While you are within that light, you have Advantage on D20 Tests. In addition, a Cleric or Druid in the light can cast level 1 spells they have prepared without expending spell slots.

Alternatively, when you light the candle for the first time, you can cast \textit{Gate} with it. Doing so destroys the candle. The portal created by the spell links to a particular Outer Plane chosen by the DM or determined by rolling on the following table.


\begin{DndTable}{cX}

1d100 & Outer Plane\\
01-05 & \textit{Abyss}\\
06-10 & \textit{Acheron}\\
11-17 & \textit{Arborea}\\
18-25 & \textit{Arcadia}\\
26-33 & \textit{Beastlands}\\
34-41 & \textit{Bytopia}\\
42-46 & \textit{Carceri}\\
47-54 & \textit{Elysium}\\
55-59 & \textit{Gehenna}\\
60-64 & \textit{Hades}\\
65-69 & \textit{Limbo}\\
70-77 & \textit{Mechanus}\\
78-85 & \textit{Mount Celestia}\\
86-90 & \textit{Nine Hells}\\
91-95 & \textit{Pandemonium}\\
96-00 & \textit{Ysgard}\\
\end{DndTable}

\DndItemHeader{Candle of the Deep}{Wondrous item, Common}
The flame of this candle isn't extinguished when immersed in water. It gives off light and heat like a normal candle.

\DndItemHeader{Cap of Water Breathing}{Wondrous item, Uncommon}
While wearing this cap underwater, you can take a Magic action to create a bubble of air around your head. This bubble allows you to breathe normally underwater. This bubble stays with you until the cap is removed or you are no longer underwater.

\DndItemHeader{Cape of the Mountebank}{Wondrous item, Rare}
This cape smells faintly of brimstone. While wearing it, you can use it to cast \textit{Dimension Door} as a Magic action. This property can't be used again until the next dawn.

When you teleport with that spell, you leave behind a cloud of smoke. The space you left is Lightly Obscured by that smoke until the end of your next turn.

\DndItemHeader{Carnelian}{}
A orange to red brown gemstone.

\DndItemHeader{Carpet of Flying}{Wondrous item, Very Rare}
You can make this carpet hover and fly by taking a Magic action and using the carpet's command word. It moves according to your directions if you are within 30 feet of it.

Four sizes of Carpet of Flying exist. The DM chooses the size of a given carpet or determines it randomly by rolling on the following table. A carpet can carry up to twice the weight shown on the table, but its Fly Speed is halved if it carries more than its normal capacity.

\begin{DndTable}[header=Applicable Carpet Sizes]{cXcc}

1d100 & Size & Capacity & Fly Speed\\
01-20 & \textit{Carpet of Flying, 3 ft. × 5 ft.} & 200 lb. & 80 feet\\
21-55 & \textit{Carpet of Flying, 4 ft. × 6 ft.} & 400 lb. & 60 feet\\
56-80 & \textit{Carpet of Flying, 5 ft. × 7 ft.} & 600 lb. & 40 feet\\
81-00 & \textit{Carpet of Flying, 6 ft. × 9 ft.} & 800 lb. & 30 feet\\
\end{DndTable}

\DndItemHeader{Carrion Crawler Mucus}{}
A creature subjected to Carrion Crawler Mucus must succeed on a DC 13 Constitution saving throw or have the Poisoned condition for 1 minute. The creature also has the Paralyzed condition while Poisoned in this way. The creature repeats the save at the end of each of its turns, ending the effect on itself on a success.

\DndItemHeader{Cauldron of Rebirth}{Wondrous item, Very Rare (requires attunement by a druid or warlock)}
This Tiny pot bears relief scenes of heroes on its cast-iron sides.

You can use the cauldron as a Spellcasting Focus for your spells, and it functions as a suitable component for the \textit{Scrying} spell.

\subparagraph{Brew Potion}
When you finish a Long Rest, you can use the cauldron to create a Potion of Greater Healing, which takes 1 minute. The potion lasts for 24 hours, then loses its magic if not consumed.

\subparagraph{Raise Dead}
As a Magic action, you can cause the cauldron to grow large enough for a Medium creature to crouch within. You can revert the cauldron to its normal size as a Magic action, harmlessly shunting anything that can't fit inside to the nearest unoccupied space.

If you place the corpse of a Humanoid into the cauldron and cover the corpse with 200 pounds of salt (which costs 10 GP) for at least 8 hours, the salt is consumed and the creature returns to life as if by \textit{Raise Dead} at the next dawn. Once used, this property can't be used again for 7 days.

\DndItemHeader{Censer of Controlling Air Elementals}{Wondrous item, Rare}
While gently swinging this censer, you can take a Magic action to summon an \textbf{Air Elemental}. The elemental appears in an unoccupied space as close to the censer as possible, understands your languages, obeys your commands, and takes its turn immediately after you on your Initiative count. The elemental disappears after 1 hour, when it dies, or when you dismiss it as a Bonus Action. The censer can't be used this way again until the next dawn.

\DndItemHeader{Chalcedony}{}
A white gemstone.

\DndItemHeader{Charlatan's Die}{Wondrous item, Common (requires attunement)}
Whenever you roll this six-sided die, you can control which number it rolls.

\DndItemHeader{Chime of Opening}{Wondrous item, Rare}
This hollow metal tube measures about 1 foot long and weighs 1 pound. As a Magic action, you can strike the chime to cast \textit{Knock}. The spell's customary knocking sound is replaced by the clear, ringing tone of the chime, which is audible out to 300 feet. The chime can be used 10 times. After the tenth time, it cracks and becomes useless.

\DndItemHeader{Chrysoberyl}{}
A yellow green to pale green gemstone.

\DndItemHeader{Chrysoprase}{}
A green gemstone.

\DndItemHeader{Circlet of Blasting}{Wondrous item, Uncommon}
While wearing this circlet, you can cast \textit{Scorching Ray} with it (5 to hit). The circlet can't cast this spell again until the next dawn.

\DndItemHeader{Citrine}{}
A pale yellow brown gemstone.

\DndItemHeader{Cloak of Arachnida}{Wondrous item, Very Rare (requires attunement)}
This fine garment is made of black silk interwoven with faint, silvery threads. While wearing it, you gain the following benefits.

\subparagraph{Poison Resistance}
You have Resistance to Poison damage.

\subparagraph{Spider Climb}
You have a Climb Speed equal to your Speed and can move up, down, and across vertical surfaces and along ceilings, while leaving your hands free.

\subparagraph{Spider Walk}
You can't be caught in webs of any sort and can move through webs as if they were Difficult Terrain.

\subparagraph{Web}
You can cast \textit{Web} (save DC 13). The web created by the spell fills twice its normal area. Once used, this property can't be used again until the next dawn.

\DndItemHeader{Cloak of Billowing}{Wondrous item, Common}
While wearing this cloak, you can take a Bonus Action to make it billow dramatically for 1 minute.

\DndItemHeader{Cloak of Displacement}{Wondrous item, Rare (requires attunement)}
While you wear this cloak, it magically projects an illusion that makes you appear to be standing in a place near your actual location, causing any creature to have Disadvantage on attack rolls against you. If you take damage, the property ceases to function until the start of your next turn. This property is suppressed while your Speed is 0.

\DndItemHeader{Cloak of Elvenkind}{Wondrous item, Uncommon (requires attunement)}
While you wear this cloak, Wisdom (Perception) checks made to perceive you have Disadvantage, and you have Advantage on Dexterity (Stealth) checks.

\DndItemHeader{Cloak of Invisibility}{Wondrous item, Legendary (requires attunement)}
This cloak has 3 charges and regains 1d3 expended charges daily at dawn. While wearing the cloak, you can take a Magic action to pull its hood over your head and expend 1 charge to give yourself the Invisible condition for 1 hour. The effect ends early if you pull the hood down (no action required) or cease wearing the hood.

\DndItemHeader{Cloak of Many Fashions}{Wondrous item, Common}
While wearing this cloak, you can take a Bonus Action to change the style, color, and apparent quality of the garment. The cloak's weight doesn't change. Regardless of its appearance, the cloak can't be anything but a cloak. Although it can duplicate the appearance of other magic cloaks, it doesn't gain their magical properties.

\DndItemHeader{Cloak of Protection}{Wondrous item, Uncommon (requires attunement)}
You gain a +1 bonus to Armor Class and saving throws while you wear this cloak.

\DndItemHeader{Cloak of the Bat}{Wondrous item, Rare (requires attunement)}
While wearing this cloak, you have Advantage on Dexterity (Stealth) checks. In an area of Dim Light or Darkness, you can grip the edges of the cloak and use it to gain a Fly Speed of 40 feet. If you ever fail to grip the cloak's edges while flying in this way, or if you are no longer in Dim Light or Darkness, you lose this Fly Speed.

While wearing the cloak in an area of Dim Light or Darkness, you can cast \textit{Polymorph} on yourself, shape-shifting into a \textbf{Bat}. While in that form, you retain your Intelligence, Wisdom, and Charisma scores. The cloak can't be used this way again until the next dawn.

\DndItemHeader{Cloak of the Manta Ray}{Wondrous item, Uncommon (requires attunement)}
While wearing this cloak, you can breathe underwater, and you have a Swim Speed of 60 feet.

\DndItemHeader{Clockwork Amulet}{Wondrous item, Common}
This copper amulet contains tiny interlocking gears and is powered by magic from Mechanus, a plane of clockwork predictability. Faint ticking and whirring noises emanate from within.

When you make an attack roll while wearing the amulet, you can forgo rolling the d20 to get a 10 on the die. Once used, this property can't be used again until the next dawn.

\DndItemHeader{Clothes of Mending}{Wondrous item, Common}
This elegant outfit magically mends itself to counteract daily wear and tear. Pieces of the outfit that are destroyed can't be repaired in this way.

\DndItemHeader{Coral}{}
A crimson gemstone.

\DndItemHeader{Crystal Ball}{Wondrous item, Very Rare (requires attunement)}
The typical crystal ball, a very rare item, is about 6 inches in diameter. While touching it, you can cast the \textit{scrying} spell (save DC 17) with it.

\DndItemHeader{Crystal Ball of Mind Reading}{Wondrous item, Legendary (requires attunement)}
While touching this crystal orb, you can cast \textit{Scrying} (save DC 17) with it. In addition, you can cast \textit{Detect Thoughts} (save DC 17) targeting creatures you can see within 30 feet of the spell's sensor. You don't need to concentrate on this \textit{Detect Thoughts} spell to maintain it during its duration, but it ends if the \textit{Scrying} spell ends.

\DndItemHeader{Crystal Ball of Telepathy}{Wondrous item, Legendary (requires attunement)}
While touching this crystal orb, you can cast \textit{Scrying} (save DC 17) with it. In addition, you can communicate telepathically with creatures you can see within 30 feet of the spell's sensor. You can also cast \textit{Suggestion} (save DC 17) through the sensor on one of those creatures. You don't need to concentrate on this \textit{Suggestion} to maintain it during its duration, but it ends if \textit{Scrying} ends. You can't cast \textit{Suggestion} in this way again until the next dawn.

\DndItemHeader{Crystal Ball of True Seeing}{Wondrous item, Legendary (requires attunement)}
While touching this crystal orb, you can cast \textit{Scrying} (save DC 17) with it. In addition, you have Truesight with a range of 120 feet centered on the spell's sensor.

\DndItemHeader{Cube of Force}{Wondrous item, Rare (requires attunement)}
This cube is about an inch across. Each face has a distinct marking on it. You can press one of those faces, expend the number of charges required for it, and thereby cast the spell associated with it, as shown in the Cube of Force Faces table.

The cube starts with 10 charges, and it regains 1d6 expended charges daily at dawn.

\begin{DndTable}[header=Cube of Force Faces]{Xc}

Spell & Charge Cost\\
\textit{Mage Armor} & 1\\
\textit{Shield} & 1\\
\textit{Leomund's Tiny Hut} & 3\\
\textit{Mordenkainen's Private Sanctum} & 4\\
\textit{Otiluke's Resilient Sphere} & 4\\
\textit{Wall of Force} & 5\\
\end{DndTable}

\DndItemHeader{Cube of Summoning}{Wondrous item, Rare}
This Tiny cube looks like a jack-in-the-box. When you wind its crank as a Magic action, a merry tune emits from the box, the lid pops open, a creature appears in the nearest unoccupied space, and the lid closes. The lid can't otherwise be opened.

Roll on the Cube [Area of Effect] of Summoning table to determine which spell the cube casts to summon the creature. The spell is cast at level 5 and doesn't require Concentration, and the summoned creature views you as its summoner.

Once the cube summons a creature, it can't do so again until the next dawn.

\begin{DndTable}[header=Cube of Summoning]{cX}

1d6 & Spell\\
1 & \textit{Summon Aberration}\\
2 & \textit{Summon Beast}\\
3 & \textit{Summon Construct}\\
4 & \textit{Summon Dragon}\\
5 & \textit{Summon Elemental}\\
6 & \textit{Summon Fey}\\
\end{DndTable}

\DndItemHeader{Cubic Gate}{Wondrous item, Legendary}
This cube is 3 inches across and radiates palpable magical energy. The six sides of the cube are each keyed to a different plane of existence, one of which is the Material Plane. The other sides are linked to planes determined by the DM.

The cube has 3 charges and regains 1d3 expended charges daily at dawn. As a Magic action, you can expend 1 of the cube's charges to cast one of the following spells using the cube.

\subparagraph{Gate}
Pressing one side of the cube, you cast \textit{Gate}, opening a portal to the plane of existence keyed to that side.

\subparagraph{Plane Shift}
Pressing one side of the cube twice, you cast \textit{Plane Shift}, transporting the targets to the plane of existence keyed to that side.

\DndItemHeader{Daern's Instant Fortress}{Wondrous item, Rare (requires attunement)}
As a Magic action, you can place this 1-inch adamantine statuette on the ground and, using a command word, cause it to grow rapidly into a square adamantine tower. Repeating the command word causes the tower to revert to statuette form, which works only if the tower is empty. Each creature in the area where the tower appears is pushed to an unoccupied space outside but next to the tower. Objects in the area that aren't being worn or carried are also pushed clear of the tower.

The tower is 20 feet on a side and 30 feet high, with arrow slits on all sides and a battlement atop it. Its interior is divided into two floors, with a ladder, staircase, or ramp (your choice) connecting them. This ladder, staircase, or ramp ends at a trapdoor leading to the roof. When created, the tower has a single door at ground level on the side facing you. The door opens only at your command, which you can issue as a Bonus Action. It is immune to the \textit{Knock} spell and similar magic.

Magic prevents the tower from being tipped over. The roof, the door, and the walls each have AC 20; HP 100; Immunity to Bludgeoning, Piercing, and Slashing damage except that which is dealt by siege equipment; and Resistance to all other damage. Shrinking the tower back down to statuette form doesn't repair damage to the tower. Only a \textit{Wish} spell can repair the tower (this use of the spell counts as replicating a spell of level 8 or lower). Each casting of Wish causes the tower to regain all its Hit Points.

\DndItemHeader{Dagger of Venom}{Rare}
You gain a +1 bonus to attack rolls and damage rolls made with this magic weapon.

You can take a Bonus Action to magically coat the blade with poison. The poison remains for 1 minute or until an attack using this weapon hits a creature. That creature must succeed on a DC 15 Constitution saving throw or take 2d10 Poison damage and have the Poisoned condition for 1 minute. The weapon can't be used this way again until the next dawn.

\DndItemHeader{Dark Shard Amulet}{Wondrous item, Common (requires attunement by a warlock)}
This amulet is fashioned from a shard of resilient material originating from an otherworldly realm. While you are wearing it, you gain the following benefits.

\subparagraph{Spellcasting Focus}
You can use the amulet as a Spellcasting Focus for your Warlock spells.

\subparagraph{Unknown Spell}
As a Magic action, you can try to cast a cantrip that you don't know. The cantrip must be on the Warlock spell list and have a casting time of an action, and you make a DC 10 Intelligence (Arcana) check. On a successful check, you cast the spell. On a failed check, the spell fails, and the action used to cast it is wasted. In either case, you can't use this property again until you finish a Long Rest.

\DndItemHeader{Decanter of Endless Water}{Wondrous item, Uncommon}
This stoppered flask sloshes when shaken, as if it contains water. The decanter weighs 2 pounds.

You can take a Magic action to remove the stopper and issue one of three command words, whereupon an amount of fresh water or salt water (your choice) pours out of the flask. The water stops pouring out at the start of your next turn. Choose from the following command words:


\begin{description}
\item[Splash.]
The decanter produces 1 gallon of water.

\item[Fountain.]
The decanter produces 5 gallons of water.

\item[Geyser.]
The decanter produces 30 gallons of water that gushes forth in a Line [Area of Effect] 30 feet long and 1 foot wide. If you're holding the decanter, you can aim the geyser in one direction (no action required). One creature of your choice in the Line [Area of Effect] must succeed on a DC 13 Strength saving throw or take 1d4 Bludgeoning damage and have the Prone condition. Instead of a creature, you can target one object in the Line [Area of Effect] that isn't being worn or carried and that weighs no more than 200 pounds. The object is knocked over by the geyser.

\end{description}

\DndItemHeader{Deck of Illusions}{Wondrous item, Uncommon}
This box contains a set of cards. A full deck has 34 cards: 32 depicting specific creatures and two with a mirrored surface. A deck found as treasure is usually missing 1d20 - 1 cards.

The magic of the deck functions only if its cards are drawn at random. You can take a Magic action to draw a card at random from the deck and throw it to the ground at a point within 30 feet of yourself. An illusion of a creature, determined by rolling on the Deck of Illusions table, forms over the thrown card and remains until dispelled. The illusory creature created by the card looks and behaves like a real creature of its kind, except that it can do no harm. While you are within 120 feet of the illusory creature and can see it, you can take a Magic action to move it anywhere within 30 feet of its card.

Any physical interaction with the illusory creature reveals it to be false, because objects pass through it. A creature that takes a Study action to visually inspect the illusory creature identifies it as an illusion with a successful DC 15 Intelligence (Investigation) check. The illusion lasts until its card is moved or the illusion is dispelled (using a \textit{Dispel Magic} spell or a similar effect). When the illusion ends, the image on its card disappears, and that card can't be used again.

\begin{DndTable}[header=Deck of Illusions]{cX}

1d100 & Illusion\\
01-03 & \textbf{Adult Red Dragon}\\
04-06 & \textbf{Archmage}\\
07-09 & \textbf{Assassin}\\
10-12 & \textbf{Bandit Captain}\\
13-15 & \textbf{Beholder}\\
16-18 & \textbf{Berserker}\\
19-21 & \textbf{Bugbear Warrior}\\
22-24 & \textbf{Cloud Giant}\\
25-27 & \textbf{Druid}\\
28-30 & \textbf{Erinyes}\\
31-33 & \textbf{Ettin}\\
34-36 & \textbf{Fire Giant}\\
37-39 & \textbf{Frost Giant}\\
40-42 & \textbf{Gnoll Warrior}\\
43-45 & \textbf{Goblin Warrior}\\
46-48 & \textbf{Guardian Naga}\\
49-51 & \textbf{Hill Giant}\\
52-54 & \textbf{Hobgoblin Warrior}\\
55-57 & \textbf{Incubus}\\
58-60 & \textbf{Iron Golem}\\
61-63 & \textbf{Knight}\\
64-66 & \textbf{Kobold Warrior}\\
67-69 & \textbf{Lich}\\
70-72 & \textbf{Medusa}\\
73-75 & \textbf{Night Hag}\\
76-78 & \textbf{Ogre}\\
79-81 & \textbf{Oni}\\
82-84 & \textbf{Priest}\\
85-87 & \textbf{Succubus}\\
88-90 & \textbf{Troll}\\
91-93 & \textbf{Warrior Veteran}\\
94-96 & \textbf{Wyvern}\\
97-00 & The card drawer\\
\end{DndTable}

\DndItemHeader{Deck of Many Things}{Wondrous item, Legendary}
Usually found in a box or pouch, this deck contains a number of cards made of ivory or vellum. Most (75 percent) of these decks have thirteen cards, but some have twenty-two. Use the appropriate column of the Deck of Many Things table when randomly determining cards drawn from the deck.

Before you draw a card, you must declare how many cards you intend to draw and then draw them randomly. Any cards drawn in excess of this number have no effect. Otherwise, as soon as you draw a card from the deck, its magic takes effect. You must draw each card no more than 1 hour after the previous draw. If you fail to draw the chosen number, the remaining number of cards fly from the deck on their own and take effect all at once.

Once a card is drawn, it disappears. Unless the card is the Fool or Jester, the card reappears in the deck, making it possible to draw the same card twice. (Once the Fool or Jester has left the deck, reroll on the table if that card comes up again.)

\begin{DndTable}[header=Deck of Many Things]{ccX}

1d100 (13-Card Deck) & 1d100 (22-Card Deck) & Card\\
— & 01-05 & Balance\\
— & 06-10 & Comet\\
— & 11-14 & Donjon\\
01-08 & 15-18 & Euryale\\
— & 19-23 & Fates\\
09-16 & 24-27 & Flames\\
— & 28-31 & Fool\\
— & 32-36 & Gem\\
17-24 & 37-41 & Jester\\
25-32 & 42-46 & Key\\
33-40 & 47-51 & Knight\\
41-48 & 52-56 & Moon\\
— & 57-60 & Puzzle\\
49-56 & 61-64 & Rogue\\
57-64 & 65-68 & Ruin\\
— & 69-73 & Sage\\
65-72 & 74-77 & Skull\\
73-80 & 78-82 & Star\\
81-88 & 83-87 & Sun\\
— & 88-91 & Talons\\
89-96 & 92-96 & Throne\\
97-00 & 97-00 & Void\\
\end{DndTable}

Each card's effect is described below.

\subparagraph{Balance}
You can increase one of your ability scores by 2, to a maximum of 22, provided you also decrease another one of your ability scores by 2. You can't decrease an ability that has a score of 5 or lower. Alternatively, you can choose not to adjust your ability scores, in which case this card has no effect.

\subparagraph{Comet}
The next time you enter combat against one or more Hostile [Attitude] creatures, you can select one of them as your foe when you roll Initiative. If you reduce your foe to 0 Hit Points during that combat, you have Advantage on Death Saving Throws for 1 year. If someone else reduces your chosen foe to 0 Hit Points or you don't choose a foe, this card has no effect.

\subparagraph{Donjon}
You disappear and become entombed in a state of suspended animation in an extradimensional sphere. Everything you're wearing and carrying disappears with you except for Artifacts, which stay behind in the space you occupied when you disappeared. You remain imprisoned until you are found and removed from the sphere. You can't be located by any Divination magic, but a \textit{Wish} spell can reveal the location of your prison. You draw no more cards.

\subparagraph{Euryale}
The card's medusa-like visage curses you. You take a -2 penalty to saving throws while cursed in this way. Only a god or the magic of the Fates card can end this curse.

\subparagraph{Fates}
Reality's fabric unravels and spins anew, allowing you to avoid or erase one event as if it never happened. You can use the card's magic as soon as you draw the card or at any other time before you die.

\subparagraph{Flames}
A powerful devil becomes your enemy. The devil seeks your ruin and torments you, savoring your suffering before attempting to slay you. This enmity lasts until either you or the devil dies.

\begin{DndSidebar}[float=!b]{A Question of Enmity}
Two of the cards in the Deck of Many Things can earn a character the enmity of another being. With the Flames card, the enmity is overt. The character should experience the devil's malevolent efforts on multiple occasions. Seeking out the fiend shouldn't be a simple task, and the adventurer should clash with the devil's allies and followers a few times before being able to confront the devil.



In the case of the Rogue card, the enmity is secret and should come from someone thought to be a friend or an ally. As Dungeon Master, you should wait for a dramatically appropriate moment to reveal this enmity, leaving the adventurer guessing who is likely to become a betrayer.



\end{DndSidebar}

\subparagraph{Fool}
You have Disadvantage on D20 Tests for the next 72 hours. Draw another card; this draw doesn't count as one of your declared draws.

\subparagraph{Gem}
Twenty-five pieces of jewelry worth 2,000 GP each or fifty gems worth 1,000 GP each appear at your feet.

\subparagraph{Jester}
You have Advantage on D20 Tests for the next 72 hours, or you can draw two additional cards beyond your declared draws.

\subparagraph{Key}
A Rare or rarer magic weapon with which you are proficient appears on your person. The DM chooses the weapon.

\subparagraph{Knight}
You gain the service of a \textbf{Knight}, who magically appears in an unoccupied space you choose within 30 feet of yourself. The knight has the same alignment as you and serves you loyally until death, believing the two of you have been drawn together by fate. Work with your DM to create a name and backstory for this NPC. The DM can use a different stat block to represent the knight, as desired.

\subparagraph{Moon}
You gain the ability to cast \textit{Wish} 1d3 times.

\subparagraph{Puzzle}
Permanently reduce your Intelligence or Wisdom by 1d4 + 1 (to a minimum score of 1). You can draw one additional card beyond your declared draws.

\subparagraph{Rogue}
An NPC of the DM's choice becomes Hostile [Attitude] toward you. You don't know the identity of this NPC until they or someone else reveals it. Nothing less than a \textit{Wish} spell or divine intervention can end the NPC's hostility toward you.

\subparagraph{Ruin}
All forms of wealth that you carry or own, other than magic items, are lost to you. Portable property vanishes. Businesses, buildings, and land you own are lost in a way that alters reality the least. If you have a Bastion (see the \textit{Dungeon Master's Guide}), it is destroyed by some calamity beyond your control. Any documentation that proves you should own something lost to this card also disappears.

\subparagraph{Sage}
At any time you choose within one year of drawing this card, you can ask a question in meditation and mentally receive a truthful answer to that question.

\subparagraph{Skull}
An \textbf{Avatar of Death} appears in an unoccupied space as close to you as possible. The avatar targets only you with its attacks, appearing as a ghostly skeleton clad in a tattered black robe and carrying a spectral scythe. The avatar disappears when it drops to 0 Hit Points or you die. If an ally of yours deals damage to the avatar, that ally summons another Avatar of Death. The new avatar appears in an unoccupied space as close to that ally as possible and targets only that ally with its attacks. You and your allies can each summon only one avatar as a consequence of this draw. A creature slain by an avatar can't be restored to life.

\subparagraph{Star}
Increase one of your ability scores by 2, to a maximum of 24.

\subparagraph{Sun}
A magic item (chosen by the DM) appears on your person. In addition, you gain 10 Temporary Hit Points daily at dawn until you die.

\subparagraph{Talons}
Every magic item you wear or carry disintegrates. Artifacts in your possession vanish instead.

\subparagraph{Throne}
You gain proficiency and Expertise in your choice of History, Insight, Intimidation, or Persuasion. In addition, you gain rightful ownership of a small keep somewhere in the world. However, the keep is currently home to one or more monsters, which must be cleared out before you can claim the keep as yours.

\subparagraph{Void}
Your soul is drawn from your body and contained in an object in a place of the DM's choice. One or more powerful beings guard the place. While your soul is trapped in this way, your body is inert, ceases aging, and requires no food, air, or water. A \textit{Wish} spell can't return your soul to your body, but the spell reveals the location of the object that holds your soul. You draw no more cards.

\DndItemHeader{Demonomicon of Iggwilv}{Wondrous item, Artifact (requires attunement)}
This treatise, composed by \textit{Iggwilv} the archmage, documents [Tooltip Not Found] layers and inhabitants and is widely regarded as the most thorough and blasphemous tome of demonology in the multiverse. The tome recounts both the oldest and most current profanities of the Abyss and demons. Demons have attempted to censor the text, and while sections have been ripped from the book's spine, the general chapters remain, ever revealing demonic secrets. Caged behind lines of script roils a secret piece of the Abyss itself, which keeps the book up-to-date, no matter how many pages are removed, and it longs to be more than mere reference material.

\subparagraph{Abyssal Lore}
You can reference the Demonomicon whenever you make an Intelligence check to discern information about demons or a Wisdom (Survival) check related to the Abyss. When you do so, you gain Advantage on the check.

\subparagraph{Containment}
The first ten pages of the Demonomicon are blank. As a Magic action while holding the book, you can target a Fiend that you can see that is trapped within the area of a \textit{Magic Circle} spell. The Fiend must succeed on a DC 20 Charisma saving throw with Disadvantage or become trapped within one of the Demonomicon's blank pages, which fills with writing detailing the trapped creature's widely known name and depravities. Once used, this action can't be used again until the next dawn.

When you finish a Long Rest, if you and the Demonomicon are on the same plane of existence, one trapped creature within the book can attempt to possess you. You make a DC 20 Charisma saving throw. On a failed save, you are possessed by the creature, which controls you like a puppet. As a Magic action, the possessing creature can release you and appear in the closest unoccupied space to you. On a successful save, the Fiend can't try to possess you again for 7 days (but another Fiend trapped in the book can certainly try).

When the tome is discovered, it has 1d4 Fiends occupying its pages—typically an assortment of demons.

\subparagraph{Ensnarement}
While carrying the book, whenever you cast \textit{Magic Circle} naming only Fiends or cast \textit{Planar Binding} targeting a Fiend, the spell is cast at level 9, regardless of what level spell slot you used, if any. Additionally, the Fiend has Disadvantage on its saving throw against the spell.

\subparagraph{Fiendish Scourging}
While carrying the book, when you make a damage roll for a spell you cast against a Fiend, you use the maximum possible result instead of rolling.

\subparagraph{Random Properties}
The Artifact has the following random properties:


\begin{itemize}
\item 2 Artifact Properties; Minor Beneficial Properties properties
\item 1 Artifact Properties; Minor Detrimental Properties property
\item 1 Artifact Properties; Major Detrimental Properties property
\end{itemize}

\subparagraph{Spells}
The book has 8 charges and regains 1d8 expended charges daily at dawn. While holding the book, you can take a Magic action to cast one of the spells (save DC 20) on the following table. The table indicates how many charges you must expend to cast the spell.


\begin{DndTable}{Xc}

Spell & Charge Cost\\
\textit{Magic Circle} & 1\\
\textit{Magic Jar} & 3\\
\textit{Planar Ally} & 3\\
\textit{Planar Binding} & 2\\
\textit{Plane Shift} (to [Tooltip Not Found] only) & 3\\
\textit{Summon Fiend} & 3\\
\textit{Tasha's Hideous Laughter} & 0\\
\end{DndTable}

\subparagraph{Destroying the Demonomicon}
To destroy the book, six different demon lords must each tear out a sixth of the book's pages. If this occurs, the pages reappear after 24 hours. Before all those hours pass, anyone who opens the book's remaining binding is transported to a nascent layer of the Abyss that lies hidden within the book. At the heart of this deadly, semisentient domain lies a long-lost Artifact, Fraz-Urb'luu's Staff. If the staff is dragged from the pocket plane, the tome is reduced to a mundane and out-of-date copy of the Tome of Zyx, the work that served as the foundation of the Demonomicon of Iggwilv. The Tome of Zyx can be destroyed like any ordinary book. Once the staff emerges, the demon lord Fraz-Urb'luu knows instantly.

\DndItemHeader{Diamond}{}
A blue white, canary, pink, brown, or blue gemstone.

\DndItemHeader{Dimensional Shackles}{Wondrous item, Rare}
You can take a Utilize action to place these shackles on a creature that has the Incapacitated condition. The shackles adjust to fit a creature of Small to Large size. The shackles prevent a creature bound by them from using any method of extradimensional movement, including teleportation or travel to a different plane of existence. They don't prevent the creature from passing through an interdimensional portal.

You and any creature you designate when you use the shackles can take a Utilize action to remove them. Once every 30 days, the bound creature can make a DC 30 Strength (Athletics) check. On a successful check, the creature breaks free and destroys the shackles.

\DndItemHeader{Dragon Scale Mail}{Very Rare (requires attunement)}
Dragon Scale Mail is made of the scales of one kind of dragon. Sometimes dragons collect their cast-off scales and gift them. Other times, hunters carefully preserve the hide of a dead dragon. In either case, Dragon Scale Mail is highly valued.

While wearing this armor, you gain a +1 bonus to Armor Class, you have Advantage on saving throws against the breath weapons of Dragons, and you have Resistance to {{getFullImmRes item.resist}} damage.

Additionally, you can focus your senses as a Magic action to discern the distance and direction to the closest {{item.detail1}} dragon within 30 miles of yourself. This action can't be used again until the next dawn.

\DndItemHeader{Dread Helm}{Wondrous item, Common}
While you're wearing this fearsome steel helm, your eyes glow red and the rest of your face is hidden in shadow.

\DndItemHeader{Driftglobe}{Wondrous item, Uncommon}
This small sphere of thick glass weighs 1 pound. If you are within 60 feet of it, you can command it to emanate light equivalent to that of the \textit{Light} or \textit{Daylight} spell (your choice). Once used, the Daylight effect can't be used again until the next dawn.

You can issue another command as a Magic action to make the illuminated globe rise into the air and float no more than 5 feet off the ground. The globe hovers in this way until you or another creature grasps it. If you move more than 60 feet from the hovering globe, it follows you until it is within 60 feet of you. It takes the shortest route to do so. If prevented from moving, the globe sinks gently to the ground and becomes inactive, and its light winks out.

\DndItemHeader{Dust of Disappearance}{Wondrous item, Uncommon}
This powder resembles fine sand. There is enough of it for one use. When you take a Utilize action to throw the dust into the air, you and each creature and object within a 10-foot Emanation [Area of Effect] originating from you have the Invisible condition for 2d4 minutes. The duration is the same for all subjects, and the dust is consumed when its magic takes effect. Immediately after an affected creature makes an attack roll, deals damage, or casts a spell, the Invisible condition ends for that creature.

\DndItemHeader{Dust of Dryness}{Wondrous item, Uncommon}
This small packet contains 1d6 + 4 pinches of dust. As a Utilize action, you can sprinkle a pinch of the dust over water, turning up to a 15-foot Cube [Area of Effect] of water into one marble-sized pellet, which floats or rests near where the dust was sprinkled. The pellet's weight is negligible. A creature can take a Utilize action to smash the pellet against a hard surface, causing the pellet to shatter and release the water the dust absorbed. Doing so destroys the pellet and ends its magic.

As a Utilize action, you can sprinkle a pinch of the dust on an Elemental within 5 feet of yourself that is composed mostly of water (such as a \textbf{Water Elemental} or a \textbf{Water Weird}). Such a creature exposed to a pinch of the dust makes a DC 13 Constitution saving throw, taking 10d6 Necrotic damage on a failed save or half as much damage on a successful one.

\DndItemHeader{Dust of Sneezing and Choking}{Wondrous item, Uncommon}
Found in a small container, this powder resembles \textit{Dust of Disappearance}, and \textit{Identify} reveals it to be such. There is enough of it for one use.

As a Utilize action, you can throw the dust into the air, forcing yourself and every creature in a 30- foot Emanation [Area of Effect] originating from you to make a DC 15 Constitution saving throw. Constructs, Elementals, Oozes, Plants, and Undead succeed on the save automatically.

On a failed save, a creature begins sneezing uncontrollably; it has the Incapacitated condition and is suffocating. The creature repeats the save at the end of each of its turns, ending the effect on itself on a success. The effect also ends on any creature targeted by a \textit{Lesser Restoration} spell.

\DndItemHeader{Dwarven Thrower}{Very Rare (requires attunement by a dwarf)}
You gain a +3 bonus to attack rolls and damage rolls made with this magic weapon. It has T with a normal range of 20 feet and a long range of 60 feet. When you hit with a ranged attack using this weapon, it deals an extra 1d8 Force damage, or an extra 2d8 Force damage if the target is a Giant. Immediately after hitting or missing, the weapon flies back to your hand.

\DndItemHeader{Dynamite Stick}{}
An an action, you can light a Dynamite Stick and throw it at a point up to 60 feet away, where it explodes. Each creature in a 5-foot-radius Sphere [Area of Effect] centered on that point makes a DC 12 Dexterity saving throw, taking 3d6 Force damage on a failed save or half as much damage on a successful one.

It takes 1 minute to bind two or more Dynamite Sticks together so they explode at the same time. Each stick after the first increases the damage by 1d6 (to a maximum of 10d6) and the effect's radius by 5 feet (to a maximum of 20 feet).

It takes 1 minute to rig dynamite with a longer fuse so it explodes after a longer period of time, such as 1 minute or 10 minutes.

\DndItemHeader{Ear Horn of Hearing}{Wondrous item, Common}
While held up to your ear, this horn suppresses the effects of the Deafened condition on you.

\DndItemHeader{Efreeti Bottle}{Wondrous item, Very Rare}
When you take a Magic action to remove the stopper of this painted brass bottle, a cloud of thick smoke flows out of it. At the end of your turn, the smoke disappears with a flash of harmless fire, and an \textbf{Efreeti} appears in an unoccupied space within 30 feet of you.

The first time the bottle is opened, the DM rolls on the following table to determine what happens.


\begin{DndTable}{cX}

1d10 & Effect\\
1 & The \textbf{efreeti} attacks you. After fighting for 5 rounds, the efreeti disappears, and the bottle loses its magic.\\
2-9 & The \textbf{efreeti} understands your languages and obeys your commands for 1 hour, after which it returns to the bottle, and a new stopper contains it. The stopper can't be removed for 24 hours. The next two times the bottle is opened, the same effect occurs. If the bottle is opened a fourth time, the efreeti escapes and disappears, and the bottle loses its magic.\\
10 & The \textbf{efreeti} understands your languages and can cast \textit{Wish} once for you. It disappears when it grants the wish or after 1 hour, and the bottle loses its magic.\\
\end{DndTable}

\DndItemHeader{Elixir of Health}{Rare}
When you drink this potion, it cures any disease afflicting you, and it removes the blinded, deafened, paralyzed, and poisoned conditions. The clear red liquid has tiny bubbles of light in it.

\DndItemHeader{Emerald}{}
A deep bright green gemstone.

\DndItemHeader{Enduring Spellbook}{Wondrous item, Common}
This spellbook, along with anything written on its pages, can't be damaged by fire or water. In addition, the spellbook doesn't deteriorate with age.

\DndItemHeader{Ersatz Eye}{Wondrous item, Common}
This magical eye replaces a real one that was lost or removed. While the Ersatz Eye is embedded in your eye socket, you can see through the tiny orb as though it were your natural eye. You can insert or remove the Ersatz Eye as a Magic action, and it can't be removed against your will while you are alive.

\DndItemHeader{Essence of Ether}{}
A creature subjected to Essence of Ether must succeed on a DC 15 Constitution saving throw or have the Poisoned condition for 8 hours. The creature also has the Unconscious condition while Poisoned in this way. The creature wakes up if it takes damage or if another creature takes an action to shake it awake.

\DndItemHeader{Eversmoking Bottle}{Wondrous item, Uncommon}
As a Magic action, you can open or close this bottle.

Opening the bottle causes thick smoke to billow out, forming a cloud that fills a 60-foot Emanation [Area of Effect] originating from the bottle. The area within the smoke is Heavily Obscured.

Each minute the bottle remains open, the size of the Emanation [Area of Effect] increases by 10 feet until it reaches its maximum size of 120 feet.

Closing the bottle causes the cloud to become fixed in place until it disperses after 10 minutes. A strong wind (such as that created by the \textit{Gust of Wind} spell) disperses the cloud after 1 minute.

\DndItemHeader{Eye Agate}{}
A circles of gray, white, brown, blue, or green gemstone.

\DndItemHeader{Eye and Hand of Vecna}{Wondrous item, Artifact (requires attunement)}
\textit{Vecna} was a mighty wizard who, through magic and conquest, forged a terrible empire. For all his power, however, Vecna feared death and took steps to prevent his demise by becoming a lich.

A treacherous lieutenant named Kas brought Vecna's rule to an end in a terrible battle. Of Vecna, all that remained were one hand and one eye, grisly Artifacts that still seek to work Vecna's will in the world.

The \textit{Eye of Vecna} and the \textit{Hand of Vecna} are separate Artifacts that might be found together or separately. The eye looks like a bloodshot organ torn free from the socket. The hand is a shriveled left extremity.

\subparagraph{Random Properties of the Eye and Hand}
The \textit{Eye of Vecna} and the \textit{Hand of Vecna} each have the following random properties


\begin{itemize}
\item 1 Artifact Properties; Minor Beneficial Properties property
\item 1 Artifact Properties; Major Beneficial Properties property
\item 1 Artifact Properties; Minor Detrimental Properties property
\end{itemize}

\subparagraph{Attuning to the Eye}
To attune to the eye, you must press it into your empty socket. The eye grafts itself to your head and remains there until you die. If the eye is ever removed, you die.

\subparagraph{Properties of the Eye}
While you are attuned to the eye, your alignment is Neutral Evil, and you gain the following benefits:


\begin{description}
\item[Truesight.]
You have Truesight out to 240 feet.

\item[Spellcasting.]
The eye has 8 charges and regains 1d4 + 4 expended charges daily at dawn. You can cast a spell on the \textit{Eye of Vecna} Spells table from the eye (save DC 18). The table indicates how many charges you must expend to cast the spell. Each time you cast a spell from the eye, there is a 5 chance that Vecna tears your soul from your body, devours it, and then takes control of the body like a puppet. If that happens, you become an NPC under the DM's control.

\begin{DndTable}[header=Eye of Vecna Spells]{Xc}

Spell & Charge Cost\\
\textit{Clairvoyance} & 2\\
\textit{Crown of Madness} & 1\\
\textit{Disintegrate} & 4\\
\textit{Dominate Monster} & 5\\
\textit{Eyebite} & 4\\
\end{DndTable}

\item[X-ray Vision.]
You can take a Magic action to gain X-ray vision with a range of 30 feet for 1 minute. To you, solid objects within that radius appear transparent and don't prevent light from passing through themselves. The vision can penetrate 1 foot of stone, 1 inch of common metal, or up to 3 feet of wood or dirt. Thicker substances block the vision, as does a thin sheet of lead.

\end{description}

\subparagraph{Attuning to the Hand}
To attune to the hand, you must press it against the stump where your left hand was. The hand grafts itself to your arm and becomes a functioning appendage. If the hand is ever removed, you die.

\subparagraph{Properties of the Hand}
When you are attuned to the hand, your alignment is Neutral Evil, and you gain the following benefits:


\begin{description}
\item[Great Strength.]
Your Strength becomes 20 unless it is already 20 or higher.

\item[Icy Touch.]
Any melee spell attack you make with the hand and any melee attack made with a weapon held by it deals an extra 2d8 Cold damage on a hit.

\item[Spellcasting.]
The hand has 8 charges and regains 1d4 + 4 expended charges daily at dawn. You can cast a spell on the \textit{Hand of Vecna} Spells table from the hand (save DC 18). The table indicates how many charges you must expend to cast the spell. Each time you cast a spell from it, the hand casts \textit{Suggestion} on you (save DC 18; no Concentration required), demanding that you commit an evil act. The hand might have a specific act in mind or leave it up to you.

\begin{DndTable}[header=Hand of Vecna Spells]{Xc}

Spell & Charge Cost\\
\textit{Finger of Death} & 5\\
\textit{Sleep} & 1\\
\textit{Slow} & 2\\
\textit{Teleport} & 3\\
\end{DndTable}

\end{description}

\subparagraph{Properties of the Eye and Hand}
While attuned to both the hand and eye, you gain the following additional benefits:


\begin{description}
\item[Danger Sense.]
You have Advantage on Initiative rolls.

\item[Necrotic Reduction.]
As a Magic action, you can target one creature you can see within 5 feet of yourself. The target makes a DC 18 Constitution saving throw, taking 7d6 Necrotic damage on a failed save or half as much damage on a successful one. A creature reduced to 0 Hit Points by this damage is transformed into green slime (see chapter 3) that covers the ground in its space, each 5-foot square of slime representing a separate patch. Nonmagical objects worn or carried by the target that are made of metal or organic material are destroyed by the slime.

\item[Poison Immunity.]
You have Immunity to Poison damage and the Poisoned condition.

\item[Regeneration.]
If you start your turn with at least 1 Hit Points, you regain 1d10 Hit Points.

\item[Wish.]
You can cast \textit{Wish}. Once used, this property can't be used again until 30 days have passed.

\end{description}

\subparagraph{Destroying the Eye and Hand}
If the \textit{Eye of Vecna} and the \textit{Hand of Vecna} are both attached to the same creature and that creature is slain by the \textit{Sword of Kas}, both the eye and the hand burst into flame, turn to ash, and are destroyed. Any other attempt to destroy the eye or hand seems to work, but the Artifact reappears in one of Vecna's many hidden vaults, where it waits to be rediscovered.

\DndItemHeader{Eyes of Charming}{Wondrous item, Uncommon (requires attunement)}
These crystal lenses fit over the eyes. They have 3 charges. While wearing them, you can expend 1 or more charges to cast \textit{Charm Person} (save DC 13). For 1 charge, you cast the level 1 version of the spell. You increase the spell's level by one for each additional charge you expend. The lenses regain all expended charges daily at dawn.

\DndItemHeader{Eyes of Minute Seeing}{Wondrous item, Uncommon}
These crystal lenses fit over the eyes. While wearing them, your vision improves significantly out to a range of 1 foot, granting you Darkvision within that range and Advantage on Intelligence (Investigation) checks made to examine something within that range.

\DndItemHeader{Eyes of the Eagle}{Wondrous item, Uncommon}
These crystal lenses fit over the eyes. While wearing them, you have Advantage on Wisdom (Perception) checks that rely on sight. In conditions of clear visibility, you can make out details of even extremely distant creatures and objects as small as 2 feet across.

\DndItemHeader{Figurine of Wondrous Power, Bronze Griffon}{Wondrous item, Rare}
A \textit{Figurine of Wondrous Power} is a statuette small enough to fit in a pocket. If you take a Magic action to throw the figurine to a point on the ground within 60 feet of yourself, the figurine becomes a living creature specified in the figurine's description below. If the space where the creature would appear is occupied by other creatures or objects, or if there isn't enough space for the creature, the figurine doesn't become a creature.

The creature is Friendly [Attitude] to you and your allies. It understands your languages, obeys your commands, and takes its turn immediately after you on your Initiative count. If you issue no commands, the creature defends itself but takes no other actions.

The creature exists for a duration specific to each figurine. At the end of the duration, the creature reverts to its figurine form. It reverts to a figurine early if its creature form drops to 0 Hit Points or if you take a Magic action while touching the creature to make it revert to figurine form. When the creature becomes a figurine again, its property can't be used again until a certain amount of time has passed, as specified below.

\subparagraph{Bronze Griffon (Rare)}
This bronze statuette is of a griffon rampant. It can become a \textbf{Griffon} for up to 6 hours. Once it has been used, it can't be used again until 5 days have passed.

\DndItemHeader{Figurine of Wondrous Power, Ebony Fly}{Wondrous item, Rare}
A \textit{Figurine of Wondrous Power} is a statuette small enough to fit in a pocket. If you take a Magic action to throw the figurine to a point on the ground within 60 feet of yourself, the figurine becomes a living creature specified in the figurine's description below. If the space where the creature would appear is occupied by other creatures or objects, or if there isn't enough space for the creature, the figurine doesn't become a creature.

The creature is Friendly [Attitude] to you and your allies. It understands your languages, obeys your commands, and takes its turn immediately after you on your Initiative count. If you issue no commands, the creature defends itself but takes no other actions.

The creature exists for a duration specific to each figurine. At the end of the duration, the creature reverts to its figurine form. It reverts to a figurine early if its creature form drops to 0 Hit Points or if you take a Magic action while touching the creature to make it revert to figurine form. When the creature becomes a figurine again, its property can't be used again until a certain amount of time has passed, as specified below.

\subparagraph{Ebony Fly (Rare)}
This ebony statuette, carved in the likeness of a horsefly, can become a \textbf{Giant Fly} for up to 12 hours and can be ridden as a mount. Once it has been used, it can't be used again until 2 days have passed.

\DndItemHeader{Figurine of Wondrous Power, Golden Lions}{Wondrous item, Rare}
A \textit{Figurine of Wondrous Power} is a statuette small enough to fit in a pocket. If you take a Magic action to throw the figurine to a point on the ground within 60 feet of yourself, the figurine becomes a living creature specified in the figurine's description below. If the space where the creature would appear is occupied by other creatures or objects, or if there isn't enough space for the creature, the figurine doesn't become a creature.

The creature is Friendly [Attitude] to you and your allies. It understands your languages, obeys your commands, and takes its turn immediately after you on your Initiative count. If you issue no commands, the creature defends itself but takes no other actions.

The creature exists for a duration specific to each figurine. At the end of the duration, the creature reverts to its figurine form. It reverts to a figurine early if its creature form drops to 0 Hit Points or if you take a Magic action while touching the creature to make it revert to figurine form. When the creature becomes a figurine again, its property can't be used again until a certain amount of time has passed, as specified below.

\subparagraph{Golden Lions (Rare)}
These gold statuettes of lions are always created in pairs. You can use one figurine or both simultaneously. Each can become a \textbf{Lion} for up to 1 hour. Once a lion has been used, it can't be used again until 7 days have passed.

\DndItemHeader{Figurine of Wondrous Power, Ivory Goats}{Wondrous item, Rare}
A \textit{Figurine of Wondrous Power} is a statuette small enough to fit in a pocket. If you take a Magic action to throw the figurine to a point on the ground within 60 feet of yourself, the figurine becomes a living creature specified in the figurine's description below. If the space where the creature would appear is occupied by other creatures or objects, or if there isn't enough space for the creature, the figurine doesn't become a creature.

The creature is Friendly [Attitude] to you and your allies. It understands your languages, obeys your commands, and takes its turn immediately after you on your Initiative count. If you issue no commands, the creature defends itself but takes no other actions.

The creature exists for a duration specific to each figurine. At the end of the duration, the creature reverts to its figurine form. It reverts to a figurine early if its creature form drops to 0 Hit Points or if you take a Magic action while touching the creature to make it revert to figurine form. When the creature becomes a figurine again, its property can't be used again until a certain amount of time has passed, as specified below.

\subparagraph{Ivory Goats (Rare)}
These ivory statuettes of goats are always created in sets of three. Each goat looks unique and functions differently from the others. Their properties are as follows:


\begin{description}
\item[Goat of Terror.]
This figurine can become a \textbf{Giant Goat} for up to 3 hours. The goat can't attack, but you can (harmlessly) remove its horns and use them as weapons. One horn becomes a \textit{+1 Lance}, and the other becomes a \textit{+2 Longsword}. Removing a horn requires a Magic action, and the weapons disappear and the horns return when the goat reverts to figurine form. While you ride the goat, any Hostile [Attitude] creature that starts its turn within a 30-foot Emanation [Area of Effect] originating from the goat must succeed on a DC 15 Wisdom saving throw or have the Frightened condition for 1 minute, until you are no longer riding the goat, or until the goat reverts to figurine form. The Frightened creature repeats the save at the end of each of its turns, ending the effect on itself on a success. Once it succeeds on the save, a creature is immune to this effect for the next 24 hours. Once the figurine has been used, it can't be used again until 15 days have passed.

\item[Goat of Traveling.]
This figurine can become a Large goat with the same statistics as a \textbf{Riding Horse}. It has 24 charges, and each hour or portion thereof it spends in goat form costs 1 charge. While it has charges, you can use it as often as you wish. When it runs out of charges, it reverts to a figurine and can't be used again until 7 days have passed, when it regains all expended charges.

\item[Goat of Travail.]
This figurine can become a \textbf{Giant Goat} for up to 3 hours. Once it has been used, it can't be used again until 30 days have passed.

\end{description}

\DndItemHeader{Figurine of Wondrous Power, Marble Elephant}{Wondrous item, Rare}
A \textit{Figurine of Wondrous Power} is a statuette small enough to fit in a pocket. If you take a Magic action to throw the figurine to a point on the ground within 60 feet of yourself, the figurine becomes a living creature specified in the figurine's description below. If the space where the creature would appear is occupied by other creatures or objects, or if there isn't enough space for the creature, the figurine doesn't become a creature.

The creature is Friendly [Attitude] to you and your allies. It understands your languages, obeys your commands, and takes its turn immediately after you on your Initiative count. If you issue no commands, the creature defends itself but takes no other actions.

The creature exists for a duration specific to each figurine. At the end of the duration, the creature reverts to its figurine form. It reverts to a figurine early if its creature form drops to 0 Hit Points or if you take a Magic action while touching the creature to make it revert to figurine form. When the creature becomes a figurine again, its property can't be used again until a certain amount of time has passed, as specified below.

\subparagraph{Marble Elephant (Rare)}
This marble statuette resembles a trumpeting elephant. It can become an \textbf{Elephant} for up to 24 hours. Once it has been used, it can't be used again until 7 days have passed.

\DndItemHeader{Figurine of Wondrous Power, Obsidian Steed}{Wondrous item, Very Rare}
A \textit{Figurine of Wondrous Power} is a statuette small enough to fit in a pocket. If you take a Magic action to throw the figurine to a point on the ground within 60 feet of yourself, the figurine becomes a living creature specified in the figurine's description below. If the space where the creature would appear is occupied by other creatures or objects, or if there isn't enough space for the creature, the figurine doesn't become a creature.

The creature is Friendly [Attitude] to you and your allies. It understands your languages, obeys your commands, and takes its turn immediately after you on your Initiative count. If you issue no commands, the creature defends itself but takes no other actions.

The creature exists for a duration specific to each figurine. At the end of the duration, the creature reverts to its figurine form. It reverts to a figurine early if its creature form drops to 0 Hit Points or if you take a Magic action while touching the creature to make it revert to figurine form. When the creature becomes a figurine again, its property can't be used again until a certain amount of time has passed, as specified below.

\subparagraph{Obsidian Steed (Very Rare)}
This polished obsidian horse can become a \textbf{Nightmare} for up to 24 hours. The nightmare fights only to defend itself. Once it has been used, it can't be used again until 5 days have passed.

The figurine has a 10 chance each time you use it to ignore your orders, including a command to revert to figurine form. If you mount the nightmare while it is ignoring your orders, you and the nightmare are instantly transported to a random location on the plane of Hades, where the nightmare reverts to figurine form.

\DndItemHeader{Figurine of Wondrous Power, Onyx Dog}{Wondrous item, Rare}
A \textit{Figurine of Wondrous Power} is a statuette small enough to fit in a pocket. If you take a Magic action to throw the figurine to a point on the ground within 60 feet of yourself, the figurine becomes a living creature specified in the figurine's description below. If the space where the creature would appear is occupied by other creatures or objects, or if there isn't enough space for the creature, the figurine doesn't become a creature.

The creature is Friendly [Attitude] to you and your allies. It understands your languages, obeys your commands, and takes its turn immediately after you on your Initiative count. If you issue no commands, the creature defends itself but takes no other actions.

The creature exists for a duration specific to each figurine. At the end of the duration, the creature reverts to its figurine form. It reverts to a figurine early if its creature form drops to 0 Hit Points or if you take a Magic action while touching the creature to make it revert to figurine form. When the creature becomes a figurine again, its property can't be used again until a certain amount of time has passed, as specified below.

\subparagraph{Onyx Dog (Rare)}
This onyx statuette of a dog can become a \textbf{Mastiff} for up to 6 hours. The mastiff has an Intelligence of 8 and can speak Common. It also has Blindsight with a range of 60 feet. Once it has been used, it can't be used again until 7 days have passed.

\DndItemHeader{Figurine of Wondrous Power, Serpentine Owl}{Wondrous item, Rare}
A \textit{Figurine of Wondrous Power} is a statuette small enough to fit in a pocket. If you take a Magic action to throw the figurine to a point on the ground within 60 feet of yourself, the figurine becomes a living creature specified in the figurine's description below. If the space where the creature would appear is occupied by other creatures or objects, or if there isn't enough space for the creature, the figurine doesn't become a creature.

The creature is Friendly [Attitude] to you and your allies. It understands your languages, obeys your commands, and takes its turn immediately after you on your Initiative count. If you issue no commands, the creature defends itself but takes no other actions.

The creature exists for a duration specific to each figurine. At the end of the duration, the creature reverts to its figurine form. It reverts to a figurine early if its creature form drops to 0 Hit Points or if you take a Magic action while touching the creature to make it revert to figurine form. When the creature becomes a figurine again, its property can't be used again until a certain amount of time has passed, as specified below.

\subparagraph{Serpentine Owl (Rare)}
This serpentine statuette of an owl can become a \textbf{Giant Owl} for up to 8 hours. The owl can communicate telepathically with you at any range if you and it are on the same plane of existence. Once it has been used, it can't be used again until 2 days have passed.

\DndItemHeader{Figurine of Wondrous Power, Silver Raven}{Wondrous item, Uncommon}
A \textit{Figurine of Wondrous Power} is a statuette small enough to fit in a pocket. If you take a Magic action to throw the figurine to a point on the ground within 60 feet of yourself, the figurine becomes a living creature specified in the figurine's description below. If the space where the creature would appear is occupied by other creatures or objects, or if there isn't enough space for the creature, the figurine doesn't become a creature.

The creature is Friendly [Attitude] to you and your allies. It understands your languages, obeys your commands, and takes its turn immediately after you on your Initiative count. If you issue no commands, the creature defends itself but takes no other actions.

The creature exists for a duration specific to each figurine. At the end of the duration, the creature reverts to its figurine form. It reverts to a figurine early if its creature form drops to 0 Hit Points or if you take a Magic action while touching the creature to make it revert to figurine form. When the creature becomes a figurine again, its property can't be used again until a certain amount of time has passed, as specified below.

\subparagraph{Silver Raven (Uncommon)}
This silver statuette of a raven can become a \textbf{Raven} for up to 12 hours. Once it has been used, it can't be used again until 2 days have passed. While in raven form, the figurine grants you the ability to cast \textit{Animal Messenger} on it.

\DndItemHeader{Fire Opal}{}
A fiery red gemstone.

\DndItemHeader{Folding Boat}{Wondrous item, Rare}
This object appears as a wooden box that measures 12 inches long, 6 inches wide, and 6 inches deep. It weighs 4 pounds and floats. It can be opened to store items inside. This item also has three command words, each requiring a Magic action to use:


\begin{description}
\item[First Command Word.]
The box unfolds into a Rowboat.

\item[Second Command Word.]
The box unfolds into a Keelboat.

\item[Third Command Word.]
The Folding Boat folds back into a box if no creatures are aboard. Any objects in the vessel that can't fit inside the box remain outside the box as it folds. Any objects in the vessel that can fit inside the box do so.

\end{description}

When the box becomes a vessel, its weight becomes that of a normal vessel its size, and anything that was stored in the box remains in the boat.

Statistics for the \textit{Rowboat} and \textit{Keelboat} appear in the \textit{Player's Handbook}. If either vessel is reduced to 0 Hit Points, the Folding Boat is destroyed.

\DndItemHeader{Fragmentation Grenade}{}
As an action, you can either throw a grenade at a point up to 60 feet away or use a Grenade Launcher to propel the grenade to a point up to 1,000 feet away. The grenade explodes at that point, creating a particular effect in a 20-foot-radius Sphere [Area of Effect].

Each creature in the Sphere [Area of Effect] makes a DC 15 Dexterity saving throw, taking 17 (5d6) Piercing damage on a failed save or half as much damage on a successful one.

\DndItemHeader{Garnet}{}
A red, brown green, or violet gemstone.

\DndItemHeader{Gauntlets of Ogre Power}{Wondrous item, Uncommon (requires attunement)}
Your Strength score is 19 while you wear these gauntlets. They have no effect on you if your Strength is 19 or higher without them.

\DndItemHeader{Gem of Brightness}{Wondrous item, Uncommon}
This prism has 50 charges. While you are holding it, you can take a Magic action and use one of three command words to cause one of the following effects:


\begin{description}
\item[First Command Word.]
The gem sheds Bright Light in a 30-foot radius and Dim Light for an additional 30 feet. This effect doesn't expend a charge. It lasts until you take a Bonus Action to repeat the command word or until you use another function of the gem.

\item[Second Command Word.]
You expend 1 charge and cause the gem to fire a brilliant beam of light at one creature you can see within 60 feet of yourself. The creature must succeed on a DC 15 Constitution saving throw or have the Blinded condition for 1 minute. The creature repeats the save at the end of each of its turns, ending the effect on itself on a success.

\item[Third Command Word.]
You expend 5 charges and cause the gem to flare with intense light in a 30-foot Cone [Area of Effect]. Each creature in the Cone [Area of Effect] makes a saving throw as if struck by the beam created with the second command word.

\end{description}

When all of the gem's charges are expended, the gem becomes a nonmagical jewel worth 50 GP.

\DndItemHeader{Gem of Seeing}{Wondrous item, Rare (requires attunement)}
This gem has 3 charges. As a Magic action, you can expend 1 charge. For the next 10 minutes, you have Truesight out to 120 feet when you peer through the gem.

The gem regains 1d3 expended charges daily at dawn.

\DndItemHeader{Glamoured Studded Leather}{Rare}
While wearing this armor, you gain a +1 bonus to Armor Class. You can also take a Bonus Action to cause the armor to assume the appearance of a normal set of clothing or some other kind of armor. You decide what it looks like—including color, style, and accessories—but the armor retains its normal bulk and weight. The illusory appearance lasts until you use this property again or doff the armor.

\DndItemHeader{Gloves of Missile Snaring}{Wondrous item, Uncommon (requires attunement)}
If you're hit by an attack roll made with a Ranged or Thrown weapon while wearing these gloves, you can take a Reaction to reduce the damage by 1d10 plus your Dexterity modifier if you have a free hand. If you reduce the damage to 0, you can catch the ammunition or weapon if it is small enough for you to hold in that hand.

\DndItemHeader{Gloves of Swimming and Climbing}{Wondrous item, Uncommon (requires attunement)}
While wearing these gloves, climbing and swimming don't cost you extra movement, and you gain a +5 bonus to Strength (Athletics) checks made to climb or swim.

\DndItemHeader{Gloves of Thievery}{Wondrous item, Uncommon}
These gloves are invisible while worn. While wearing them, you gain a +5 bonus to Dexterity (Sleight of Hand) checks and Dexterity checks made to pick locks.

\DndItemHeader{Goggles of Night}{Wondrous item, Uncommon}
While wearing these dark lenses, you have Darkvision out to 60 feet. If you already have Darkvision, wearing the goggles increases its range by 60 feet.

\DndItemHeader{Grenade Launcher}{}
As an action, you can either throw a grenade at a point up to 60 feet away or use a Grenade Launcher to propel the grenade to a point up to 1,000 feet away. The grenade explodes at that point, creating a particular effect in a 20-foot-radius Sphere [Area of Effect].

\DndItemHeader{Gunpowder (keg)}{}
Setting fire to a keg full of Gunpowder causes it to explode. When a keg explodes, each creature in a 10-foot-radius Sphere [Area of Effect] centered on the keg makes a DC 12 Dexterity saving throw, taking 24 (7d6) Fire damage on a failed save or half as much damage on a successful one.

\DndItemHeader{Gunpowder (powder horn)}{}
Setting fire to a powder horn full of Gunpowder causes it to explode. When a powder horn explodes, each creature in a 10-foot-radius Sphere [Area of Effect] centered on the powder horn makes a DC 12 Dexterity saving throw, taking 10 (3d6) Fire damage on a failed save or half as much damage on a successful one.

\DndItemHeader{Hag Eye}{Wondrous item, Uncommon}
A Hag Eye has 3 charges. While wearing or holding this item, you can expend 1 charge to cast \textit{Darkvision} (targeting yourself only) or \textit{See Invisibility}. The Hag Eye regains all expended charges daily at dawn.

\subparagraph{Coven Sensor}
The Hag Eye is usually entrusted to a hag's minion for safekeeping and transport. As a Magic action, a hag who belongs to the coven that created the Hag Eye can see what the Hag Eye sees if the hag and the Hag Eye are on the same plane of existence. This effect lasts as long as the hag maintains Concentration. Multiple hags in the coven can see through the Hag Eye simultaneously.

\subparagraph{Creating a Hag Eye}
Only a hag coven can craft this item, which is made from a real eye coated in varnish and often fitted to a pendant or another wearable item. A hag coven can have only one Hag Eye at a time, and creating a new one requires all three members of the coven to perform a special rite. This rite takes 1 hour, and the hags can't perform it if one or more of them has the Incapacitated condition. If the hags take any other actions during this rite, the rite fails and ends.

\DndItemHeader{Hat of Disguise}{Wondrous item, Uncommon (requires attunement)}
While wearing this hat, you can cast the \textit{Disguise Self} spell. The spell ends if the hat is removed.

\DndItemHeader{Hat of Many Spells}{Wondrous item, Very Rare (requires attunement by a wizard)}
This pointed hat has the following properties.

\subparagraph{Spellcasting Focus}
While holding the hat, you can use it as a Spellcasting Focus for your Wizard spells. Any spell you cast using the hat gains a special Somatic component: you must reach into the hat and "pull" the spell out of it.

\subparagraph{Unknown Spell}
While holding the hat, you can try to cast a level 1+ spell you don't know. The spell must be on the Wizard spell list, it must be of a level you can cast, and it can't have Material components costing more than 1,000 GP. Once you decide on the spell, you must expend a spell slot of the spell's level. Then, to determine whether you cast the spell, make an Intelligence (Arcana) check (DC 10 plus the spell's level). On a successful check, you cast the spell using its normal casting time, and you can't use this property again until you finish a Short Rest or Long Rest. On a failed check, you fail to cast the spell and a random effect occurs instead, determined by rolling on the following table.

Any spell you cast from the hat uses your spell save DC and spell attack bonus.


\begin{DndTable}{cX}

1d100 & Effect\\
01-50 & You cast a random spell determined by rolling 1d10: on a \textbf{1}, \textit{Enlarge/Reduce} (enlarge effect); on a \textbf{2}, \textit{Enlarge/Reduce} (reduce effect); on a \textbf{3}, \textit{Faerie Fire}; on a \textbf{4}, \textit{Fireball}; on a \textbf{5}, \textit{Gust of Wind}; on a \textbf{6}, \textit{Invisibility} (cast on yourself); on a \textbf{7}, \textit{Lightning Bolt}; on an \textbf{8}, \textit{Phantasmal Force}; on a \textbf{9}, \textit{Polymorph}; on a \textbf{10}, \textit{Stinking Cloud}.\\
51-55 & You have the Stunned condition until the end of your next turn, believing something awesome just happened.\\
56-60 & A harmless swarm of butterflies fills a 10-foot Cube [Area of Effect] within 30 feet of yourself. The swarm disperses after 1 minute.\\
61-65 & You pull a nonmagical object out of the hat. Roll 1d4 to determine the object: on a \textbf{1}, a vial of Acid; on a \textbf{2}, a flask of Alchemist's Fire; on a \textbf{3}, a Crowbar; on a \textbf{4}, a lit Torch.\\
66-70 & You suffer a bout of "magic sickness" and have the Poisoned condition for 1 hour.\\
71-75 & You have the Petrified condition until the end of your next turn.\\
76-80 & You pull a nonmagical object out of the hat. Roll 1d4 to determine the object: on a \textbf{1}, a \textit{Dagger}; on a \textbf{2}, a \textit{Rope} with a \textit{Grappling Hook} tied to one end; on a \textbf{3}, a bag of \textit{Caltrops}; on a \textbf{4}, a gem worth 50 GP.\\
81-85 & A creature appears in an unoccupied space as close to you as possible. The creature isn't under your control and acts as it normally would, and it disappears after 1 hour or when it drops to 0 Hit Points. Roll 1d4 to determine the creature: on a \textbf{1}, a \textbf{Camel}; on a \textbf{2}, a \textbf{Constrictor Snake}; on a \textbf{3}, an \textbf{Elephant}; on a \textbf{4}, a \textbf{Mule}.\\
86-90 & A Hostile [Attitude] \textbf{Swarm of Bats} flies out of the hat, occupies your space, and attacks you.\\
91-95 & A vertical, 10-foot-diameter, two-way portal to another plane of existence opens in an unoccupied space within 30 feet of you and remains open until the end of your next turn. The DM determines where it leads.\\
96-00 & You pull a magic item out of the hat. Roll 1d6 to determine the item's rarity: on a \textbf{1-3}, Common; on a \textbf{4-5}, Uncommon; on a \textbf{6}, Rare. The DM chooses the item, which disappears after 1 hour if it's not consumed or destroyed before then.\\
\end{DndTable}

\DndItemHeader{Hat of Vermin}{Wondrous item, Common}
This hat has 3 charges. While holding the hat, you can take a Magic action to expend 1 charge and summon your choice of a \textbf{Bat}, a \textbf{Frog}, or a \textbf{Rat}. The summoned creature magically appears in the hat and tries to get away from you as quickly as possible. The creature is Indifferent [Attitude] toward you and other creatures, and it isn't under your control. It behaves as an ordinary creature of its kind and disappears after 1 hour or when it drops to 0 Hit Points. The hat regains all expended charges daily at dawn.

\DndItemHeader{Hat of Wizardry}{Wondrous item, Common (requires attunement by a wizard)}
This cone-shaped hat is adorned with moons and stars. While you are wearing it, you gain the following benefits.

\subparagraph{Spellcasting Focus}
You can use the hat as a Spellcasting Focus for your Wizard spells.

\subparagraph{Unknown Spell}
As a Magic action, you can try to cast a cantrip that you don't know. The cantrip must be on the Wizard spell list and have a casting time of an action, and you make a DC 10 Intelligence (Arcana) check. On a successful check, you cast the spell. On a failed check, the spell fails, and the action used to cast the spell is wasted. In either case, you can't use this property again until you finish a Long Rest.

\DndItemHeader{Headband of Intellect}{Wondrous item, Uncommon (requires attunement)}
Your Intelligence score is 19 while you wear this headband. It has no effect on you if your Intelligence is 19 or higher without it.

\DndItemHeader{Helm of Brilliance}{Wondrous item, Very Rare (requires attunement)}
This helm is set with 1d10 diamonds, 2d10 rubies, 3d10 fire opals, and 4d10 opals. Any gem pried from the helm crumbles to dust. When all the gems are removed or destroyed, the helm loses its magic.

You gain the following benefits while wearing the helm.

\subparagraph{Diamond Light}
As long as it has at least one diamond, the helm emits a 30-foot Emanation [Area of Effect]. When at least one Undead is within that area, the Emanation [Area of Effect] is filled with Dim Light. Any Undead that starts its turn in that area takes 1d6 Radiant damage.

\subparagraph{Fire Opal Flames}
As long as the helm has at least one fire opal, you can take a Magic action to cause one weapon you are holding to burst into flames. The flames emit Bright Light in a 10-foot radius and Dim Light for an additional 10 feet. The flames are harmless to you and the weapon. When you hit with an attack using the blazing weapon, the target takes an extra 1d6 Fire damage. The fl ames last until you take a Bonus Action to extinguish them or until you drop or stow the weapon.

\subparagraph{Ruby Resistance}
As long as the helm has at least one ruby, you have Resistance to Fire damage.

\subparagraph{Spells}
You can cast one of the following spells (save DC 18), using one of the helm's gems of the specified type as a component: \textit{Daylight} (opal), \textit{Fireball} (fire opal), \textit{Prismatic Spray} (diamond), or \textit{Wall of Fire} (ruby). The gem is destroyed when the spell is cast and disappears from the helm.

\subparagraph{Taking Fire Damage}
Roll 1d20 if you are wearing the helm and take Fire damage as a result of failing a saving throw against a spell. On a roll of 1, the helm emits beams of light from its remaining gems and is then destroyed. Each creature within a 60-foot Emanation [Area of Effect] originating from you must succeed on a DC 17 Dexterity saving throw or be struck by a beam, taking Radiant damage equal to the number of gems in the helm.

\DndItemHeader{Helm of Comprehending Languages}{Wondrous item, Uncommon}
While wearing this helm, you can cast \textit{Comprehend Languages} from it.

\DndItemHeader{Helm of Telepathy}{Wondrous item, Uncommon (requires attunement)}
While wearing this helm, you have telepathy with a range of 30 feet, and you can cast \textit{Detect Thoughts} or \textit{Suggestion} (save DC 13) from the helm. Once either spell is cast from the helm, that spell can't be cast from it again until the next dawn.

\DndItemHeader{Helm of Teleportation}{Wondrous item, Rare (requires attunement)}
This helm has 3 charges. While wearing it, you can expend 1 charge to cast \textit{Teleport} from it. The helm regains 1d3 expended charges daily at dawn.

\DndItemHeader{Hematite}{}
A gray black gemstone.

\DndItemHeader{Herbalism Kit}{}

\begin{description}
\item[Ability:]
Intelligence

\item[Utilize:]
Identify a plant (DC 10)

\item[Craft:]
\textit{Antitoxin}, \textit{Candle}, \textit{Healer's Kit}, \textit{Potion of Healing}

\end{description}

\DndItemHeader{Heward's Handy Haversack}{Wondrous item, Rare}
This backpack has a central pouch and two side pouches, each of which is an extradimensional space. Each side pouch can hold up to 200 pounds of material, not exceeding a volume of 25 cubic feet. The central pouch can hold up to 500 pounds of material, not exceeding a volume of 64 cubic feet. The haversack always weighs 5 pounds, regardless of its contents.

Retrieving an item from the haversack requires a Utilize action or a Bonus Action (your choice). When you reach into the haversack for a specific item, the item is always magically on top.

If any of its pouches is overloaded, pierced, or torn, the haversack ruptures and is destroyed. If the haversack is destroyed, its contents are lost forever, although an Artifact always turns up again somewhere. If the haversack is turned inside out, its contents spill forth unharmed, and the haversack must be put right before it can be used again.

Each pouch of the haversack holds enough air for 10 minutes of breathing, divided by the number of breathing creatures inside.

Placing the haversack inside an extradimensional space created by a \textit{Bag of Holding}, \textit{Portable Hole}, or similar item instantly destroys both items and opens a gate to the Astral Plane. The gate originates where the one item was placed inside the other. Any creature within 10 feet of the gate and not behind Cover is sucked through it and deposited in a random location on the Astral Plane. The gate then closes. The gate is one-way only and can't be reopened.

\DndItemHeader{Heward's Handy Spice Pouch}{Wondrous item, Common}
This belt pouch appears empty and has 10 charges. While holding the pouch, you can take a Magic action to expend 1 charge, name any nonmagical food seasoning (such as salt, pepper, saffron, or cilantro), and remove a pinch of the desired seasoning from the pouch. A pinch is enough to season a single meal. The pouch regains 1d6 + 4 expended charges daily at dawn.

\DndItemHeader{Holy Water}{}
When you take the Attack action, you can replace one of your attacks with throwing a flask of Holy Water. Target one creature you can see within 20 feet of yourself. The target must succeed on a Dexterity saving throw (DC 8 plus your Dexterity modifier and Proficiency) or take 2d8 Radiant damage if it is a Fiend or an Undead.

\DndItemHeader{Horn of Blasting}{Wondrous item, Rare}
You can take a Magic action to blow the horn, which emits a thunderous blast in a 30-foot Cone [Area of Effect] that is audible out to 600 feet. Each creature in the Cone [Area of Effect] makes a DC 15 Constitution saving throw. On a failed save, a creature takes 5d8 Thunder damage and has the Deafened condition for 1 minute. On a successful save, a creature takes half as much damage only. Glass or crystal objects in the Cone [Area of Effect] that aren't being worn or carried take 10d8 Thunder damage.

Each use of the horn's magic has a 20 chance of causing the horn to explode. The explosion deals 10d6 Force damage to the user and destroys the horn.

\DndItemHeader{Horn of Silent Alarm}{Wondrous item, Common}
This horn has 4 charges and regains 1d4 expended charges daily at dawn. As a Magic action, you can blow the horn while expending 1 charge. One creature of your choice hears the horn's blare, provided that creature is within 600 feet of the horn. No other creature hears the horn.

\DndItemHeader{Horn of Valhalla}{Wondrous item, Varies}
You can take a Magic action to blow this horn. In response, warrior spirits from the plane of Ysgard appear in unoccupied spaces within 60 feet of you. Each spirit uses the \textbf{Berserker} stat block and returns to Ysgard after 1 hour or when it drops to 0 Hit Points. The spirits look like living, breathing warriors, and they have Immunity to the Charmed and Frightened conditions. Once you use the horn, it can't be used again until 7 days have passed.

Four types of Horn of Valhalla are known to exist, each made of a different metal. The horn's type determines how many spirits it summons, as well as the requirement for its use. The DM chooses the horn's type or determines it randomly by rolling on the following table. If you blow the horn without meeting its requirement, the summoned spirits attack you.

If you meet the requirement, they are Friendly [Attitude] to you and your allies and follow your commands.


\begin{DndTable}{cccX}

1d100 & Horn Type & Spirits & Requirement\\
01-40 & Horn of Valhalla, Silver & 2 & None\\
41-75 & Horn of Valhalla, Brass & 3 & Proficiency with all Simple weapons\\
76-90 & Horn of Valhalla, Bronze & 4 & Training with all Medium armor\\
91-00 & Horn of Valhalla, Iron & 5 & Proficiency with all Martial weapons\\
\end{DndTable}

\DndItemHeader{Horn of Valhalla, Brass}{Wondrous item, Rare}
You can take a Magic action to blow this horn. In response, warrior spirits from the plane of Ysgard appear in unoccupied spaces within 60 feet of you. Each spirit uses the \textbf{Berserker} stat block and returns to Ysgard after 1 hour or when it drops to 0 Hit Points. The spirits look like living, breathing warriors, and they have Immunity to the Charmed and Frightened conditions. Once you use the horn, it can't be used again until 7 days have passed.

A brass horn summons 3 Berserkers. To use the brass horn, you must have Proficiency with all Simple weapons.

If you blow the horn without meeting its requirement, the summoned Berserkers attack you. If you meet the requirement, they are Friendly [Attitude] to you and your allies and follow your commands.

\DndItemHeader{Horn of Valhalla, Bronze}{Wondrous item, Very Rare}
You can take a Magic action to blow this horn. In response, warrior spirits from the plane of Ysgard appear in unoccupied spaces within 60 feet of you. Each spirit uses the \textbf{Berserker} stat block and returns to Ysgard after 1 hour or when it drops to 0 Hit Points. The spirits look like living, breathing warriors, and they have Immunity to the Charmed and Frightened conditions. Once you use the horn, it can't be used again until 7 days have passed.

A bronze horn summons 4 Berserkers. To use the bronze horn, you must have training with all Medium armor.

If you blow the horn without meeting its requirement, the summoned Berserkers attack you. If you meet the requirement, they are Friendly [Attitude] to you and your allies and follow your commands.

\DndItemHeader{Horn of Valhalla, Iron}{Wondrous item, Legendary}
You can take a Magic action to blow this horn. In response, warrior spirits from the plane of Ysgard appear in unoccupied spaces within 60 feet of you. Each spirit uses the \textbf{Berserker} stat block and returns to Ysgard after 1 hour or when it drops to 0 Hit Points. The spirits look like living, breathing warriors, and they have Immunity to the Charmed and Frightened conditions. Once you use the horn, it can't be used again until 7 days have passed.

A iron horn summons 5 Berserkers. To use the iron horn, you must have Proficiency with all Martial weapons.

If you blow the horn without meeting its requirement, the summoned Berserkers attack you. If you meet the requirement, they are Friendly [Attitude] to you and your allies and follow your commands.

\DndItemHeader{Horn of Valhalla, Silver}{Wondrous item, Rare}
You can take a Magic action to blow this horn. In response, warrior spirits from the plane of Ysgard appear in unoccupied spaces within 60 feet of you. Each spirit uses the \textbf{Berserker} stat block and returns to Ysgard after 1 hour or when it drops to 0 Hit Points. The spirits look like living, breathing warriors, and they have Immunity to the Charmed and Frightened conditions. Once you use the horn, it can't be used again until 7 days have passed.

A silver horn summons 2 Berserkers. They are Friendly [Attitude] to you and your allies and follow your commands.

\DndItemHeader{Horseshoes of Speed}{Wondrous item, Rare}
These horseshoes come in a set of four. As a Magic action, you can touch one of the horseshoes to the hoof of a horse or similar creature, whereupon the horseshoe affixes itself to the hoof. Removing a horseshoe also takes a Magic action.

While all four horseshoes are attached to the same creature, its Speed is increased by 30 feet.

\DndItemHeader{Horseshoes of a Zephyr}{Wondrous item, Very Rare}
These horseshoes come in a set of four. As a Magic action, you can touch one of the horseshoes to the hoof of a horse or similar creature, whereupon the horseshoe affixes itself to the hoof. Removing a horseshoe also takes a Magic action.

While all four shoes are affixed to the hooves of a horse or similar creature, they allow the creature to move normally while floating 4 inches above a surface. This effect means the creature can cross or stand above nonsolid or unstable surfaces, such as water or lava. The creature leaves no tracks and ignores Difficult Terrain. In addition, the creature can travel for up to 12 hours a day without gaining Exhaustion levels from extended travel.

\DndItemHeader{Immovable Rod}{Uncommon}
This iron rod has a button on one end. You can take a Utilize action to press the button, which causes the rod to become magically fixed in place. Until you or another creature takes a Utilize action to push the button again, the rod doesn't move, even if it defies gravity. The rod can hold up to 8,000 pounds of weight. More weight causes the rod to deactivate and fall. A creature can take a Utilize action to make a DC 30 Strength (Athletics) check, moving the fixed rod up to 10 feet on a successful check.

\DndItemHeader{Instrument of Illusions}{Wondrous item, Common}
While you are playing this \textit{musical instrument}, you can take a Magic action to create harmless, illusory visual effects within a 5-foot Emanation [Area of Effect] originating from the instrument. If you are a Bard, the size of the Emanation [Area of Effect] increases to 15 feet. Sample visual effects include luminous musical notes, a spectral dancer, butterflies, and gently falling snow. The magical effects have neither substance nor sound, and they are obviously illusory. The effects end when you stop playing.

\DndItemHeader{Instrument of Scribing}{Wondrous item, Common}
This \textit{musical instrument} has 3 charges and regains all expended charges daily at dawn. While you are playing it, you can take a Magic action to expend 1 charge and write a magical message on a nonmagical object or surface that you can see within 30 feet of yourself. The message can be up to six words long and is written in a language you know. If you are a Bard, you can scribe an additional seven words and make the message glow faintly, allowing it to be seen in nonmagical Darkness. Casting the \textit{Dispel Magic} spell on the message erases it. Otherwise, the message fades away after 24 hours.

\DndItemHeader{Instrument of the Bards}{Wondrous item, Varies (requires attunement by a bard)}
An Instrument of the Bards is superior to an ordinary instrument in every way. Seven types of these instruments exist, each named after a bard college. A creature that attempts to play the instrument without being attuned to it must succeed on a DC 15 Wisdom saving throw or take 2d4 Psychic damage.

You can play the instrument to cast one of its spells. Once the instrument has been used to cast a spell, it can't be used to cast that spell again until the next dawn. The spells use your spellcasting ability and spell save DC.

All Instrument of the Bards can be used to cast the following spells: \textit{Fly}, \textit{Invisibility}, \textit{Levitate}, \textit{Protection from Evil and Good}

\DndItemHeader{Instrument of the Bards, Anstruth Harp}{Wondrous item, Very Rare (requires attunement by a bard)}
An Instrument of the Bards is superior to an ordinary instrument in every way. Seven types of these instruments exist, each named after a bard college. A creature that attempts to play the instrument without being attuned to it must succeed on a DC 15 Wisdom saving throw or take 2d4 Psychic damage.

You can play Anstruth Harp to cast one of the following spells: \textit{Fly}, \textit{Invisibility}, \textit{Levitate}, \textit{Protection from Evil and Good}, \textit{Cure Wounds} (level 5), \textit{Ice Storm}, and \textit{Wall of Thorns}. Once the Anstruth Harp has been used to cast a spell, it can't be used to cast that spell again until the next dawn. The spells use your spellcasting ability and spell save DC.

\DndItemHeader{Instrument of the Bards, Canaith Mandolin}{Wondrous item, Rare (requires attunement by a bard)}
An Instrument of the Bards is superior to an ordinary instrument in every way. Seven types of these instruments exist, each named after a bard college. A creature that attempts to play the instrument without being attuned to it must succeed on a DC 15 Wisdom saving throw or take 2d4 Psychic damage.

You can play the Canaith Mandolin to cast one of the following spells: \textit{Fly}, \textit{Invisibility}, \textit{Levitate}, \textit{Protection from Evil and Good}, \textit{Cure Wounds} (level 3), \textit{Dispel Magic}, and \textit{Protection from Energy} (Lightning damage only). Once the Canaith Mandolin has been used to cast a spell, it can't be used to cast that spell again until the next dawn. The spells use your spellcasting ability and spell save DC.

\DndItemHeader{Instrument of the Bards, Cli Lyre}{Wondrous item, Rare (requires attunement by a bard)}
An Instrument of the Bards is superior to an ordinary instrument in every way. Seven types of these instruments exist, each named after a bard college. A creature that attempts to play the instrument without being attuned to it must succeed on a DC 15 Wisdom saving throw or take 2d4 Psychic damage.

You can play the Cli \textit{Lyre} to cast one of the following spells: \textit{Fly}, \textit{Invisibility}, \textit{Levitate}, \textit{Protection from Evil and Good}, \textit{Stone Shape}, \textit{Wall of Fire}, and \textit{Wind Wall}. Once the Cli \textit{Lyre} has been used to cast a spell, it can't be used to cast that spell again until the next dawn. The spells use your spellcasting ability and spell save DC.

\DndItemHeader{Instrument of the Bards, Doss Lute}{Wondrous item, Uncommon (requires attunement by a bard)}
An Instrument of the Bards is superior to an ordinary instrument in every way. Seven types of these instruments exist, each named after a bard college. A creature that attempts to play the instrument without being attuned to it must succeed on a DC 15 Wisdom saving throw or take 2d4 Psychic damage.

You can play the Doss \textit{Lute} to cast one of the following spells: \textit{Fly}, \textit{Invisibility}, \textit{Levitate}, \textit{Protection from Evil and Good}, \textit{Animal Friendship}, \textit{Protection from Energy} (Fire damage only), and \textit{Protection from Poison}. Once the Doss \textit{Lute} has been used to cast a spell, it can't be used to cast that spell again until the next dawn. The spells use your spellcasting ability and spell save DC.

\DndItemHeader{Instrument of the Bards, Fochlucan Bandore}{Wondrous item, Uncommon (requires attunement by a bard)}
An Instrument of the Bards is superior to an ordinary instrument in every way. Seven types of these instruments exist, each named after a bard college. A creature that attempts to play the instrument without being attuned to it must succeed on a DC 15 Wisdom saving throw or take 2d4 Psychic damage.

You can play the Fochlucan Bandore to cast one of the following spells: \textit{Fly}, \textit{Invisibility}, \textit{Levitate}, \textit{Protection from Evil and Good}, \textit{Entangle}, \textit{Faerie Fire}, \textit{Shillelagh}, and \textit{Speak with Animals}. Once the Fochlucan Bandore has been used to cast a spell, it can't be used to cast that spell again until the next dawn. The spells use your spellcasting ability and spell save DC.

\DndItemHeader{Instrument of the Bards, Mac-Fuirmidh Cittern}{Wondrous item, Uncommon (requires attunement by a bard)}
An Instrument of the Bards is superior to an ordinary instrument in every way. Seven types of these instruments exist, each named after a bard college. A creature that attempts to play the instrument without being attuned to it must succeed on a DC 15 Wisdom saving throw or take 2d4 Psychic damage.

You can play the Mac-Fuirmidh Cittern to cast one of the following spells: \textit{Fly}, \textit{Invisibility}, \textit{Levitate}, \textit{Protection from Evil and Good}, \textit{Barkskin}, \textit{Cure Wounds}, and \textit{Fog Cloud}. Once the Mac-Fuirmidh Cittern has been used to cast a spell, it can't be used to cast that spell again until the next dawn. The spells use your spellcasting ability and spell save DC.

\DndItemHeader{Instrument of the Bards, Ollamh Harp}{Wondrous item, Legendary (requires attunement by a bard)}
An Instrument of the Bards is superior to an ordinary instrument in every way. Seven types of these instruments exist, each named after a bard college. A creature that attempts to play the instrument without being attuned to it must succeed on a DC 15 Wisdom saving throw or take 2d4 Psychic damage.

You can play the Ollamh Harp to cast one of the following spells: \textit{Fly}, \textit{Invisibility}, \textit{Levitate}, \textit{Protection from Evil and Good}, \textit{Confusion}, \textit{Control Weather}, and \textit{Fire Storm}. Once the Ollamh Harp has been used to cast a spell, it can't be used to cast that spell again until the next dawn. The spells use your spellcasting ability and spell save DC.

\DndItemHeader{Ioun Stone}{Wondrous item, Varies (requires attunement)}
Roughly marble sized, \textit{Ioun Stones} are named after Ioun, a god of knowledge and prophecy revered on some worlds. Many types of \textit{Ioun Stones} exist, each type a distinct combination of shape and color.

When you take a Magic action to toss an \textit{Ioun Stone} into the air, the stone orbits your head at a distance of 1d3 feet, conferring its benefit to you while doing so. You can have up to three \textit{Ioun Stones} orbiting your head at the same time.

Each \textit{Ioun Stone} orbiting your head is considered to be an object you are wearing. The orbiting stone avoids contact with other creatures and objects, adjusting its orbit to avoid collisions and thwarting all attempts by other creatures to attack or snatch it.

As a Utilize action, you can seize and stow any number of \textit{Ioun Stones} orbiting your head. If your Attunement to an Ioun Stone ends while it's orbiting your head, the stone falls as though you had dropped it.

\DndItemHeader{Ioun Stone, Absorption}{Wondrous item, Very Rare (requires attunement)}
Roughly marble sized, \textit{Ioun Stones} are named after Ioun, a god of knowledge and prophecy revered on some worlds. Many types of \textit{Ioun Stones} exist, each type a distinct combination of shape and color.

When you take a Magic action to toss an \textit{Ioun Stone} into the air, the stone orbits your head at a distance of 1d3 feet, conferring its benefit to you while doing so. You can have up to three \textit{Ioun Stones} orbiting your head at the same time.

Each \textit{Ioun Stone} orbiting your head is considered to be an object you are wearing. The orbiting stone avoids contact with other creatures and objects, adjusting its orbit to avoid collisions and thwarting all attempts by other creatures to attack or snatch it.

As a Utilize action, you can seize and stow any number of \textit{Ioun Stones} orbiting your head. If your Attunement to an Ioun Stone ends while it's orbiting your head, the stone falls as though you had dropped it.

While this pale lavender ellipsoid orbits your head, you can take a Reaction to cancel a spell of level 4 or lower cast by a creature you can see. A canceled spell has no effect, and any resources used to cast it are wasted. Once the stone has canceled 20 levels of spells, it burns out, turns dull gray, and loses its magic.

\DndItemHeader{Ioun Stone, Agility}{Wondrous item, Very Rare (requires attunement)}
Roughly marble sized, \textit{Ioun Stones} are named after Ioun, a god of knowledge and prophecy revered on some worlds. Many types of \textit{Ioun Stones} exist, each type a distinct combination of shape and color.

When you take a Magic action to toss an \textit{Ioun Stone} into the air, the stone orbits your head at a distance of 1d3 feet, conferring its benefit to you while doing so. You can have up to three \textit{Ioun Stones} orbiting your head at the same time.

Each \textit{Ioun Stone} orbiting your head is considered to be an object you are wearing. The orbiting stone avoids contact with other creatures and objects, adjusting its orbit to avoid collisions and thwarting all attempts by other creatures to attack or snatch it.

As a Utilize action, you can seize and stow any number of \textit{Ioun Stones} orbiting your head. If your Attunement to an Ioun Stone ends while it's orbiting your head, the stone falls as though you had dropped it.

Your Dexterity increases by 2, to a maximum of 20, while this deep-red sphere orbits your head.

\DndItemHeader{Ioun Stone, Awareness}{Wondrous item, Rare (requires attunement)}
Roughly marble sized, \textit{Ioun Stones} are named after Ioun, a god of knowledge and prophecy revered on some worlds. Many types of \textit{Ioun Stones} exist, each type a distinct combination of shape and color.

When you take a Magic action to toss an \textit{Ioun Stone} into the air, the stone orbits your head at a distance of 1d3 feet, conferring its benefit to you while doing so. You can have up to three \textit{Ioun Stones} orbiting your head at the same time.

Each \textit{Ioun Stone} orbiting your head is considered to be an object you are wearing. The orbiting stone avoids contact with other creatures and objects, adjusting its orbit to avoid collisions and thwarting all attempts by other creatures to attack or snatch it.

As a Utilize action, you can seize and stow any number of \textit{Ioun Stones} orbiting your head. If your Attunement to an Ioun Stone ends while it's orbiting your head, the stone falls as though you had dropped it.

While this dark-blue rhomboid orbits your head, you have Advantage on Initiative rolls and Wisdom (Perception) checks.

\DndItemHeader{Ioun Stone, Fortitude}{Wondrous item, Very Rare (requires attunement)}
Roughly marble sized, \textit{Ioun Stones} are named after Ioun, a god of knowledge and prophecy revered on some worlds. Many types of \textit{Ioun Stones} exist, each type a distinct combination of shape and color.

When you take a Magic action to toss an \textit{Ioun Stone} into the air, the stone orbits your head at a distance of 1d3 feet, conferring its benefit to you while doing so. You can have up to three \textit{Ioun Stones} orbiting your head at the same time.

Each \textit{Ioun Stone} orbiting your head is considered to be an object you are wearing. The orbiting stone avoids contact with other creatures and objects, adjusting its orbit to avoid collisions and thwarting all attempts by other creatures to attack or snatch it.

As a Utilize action, you can seize and stow any number of \textit{Ioun Stones} orbiting your head. If your Attunement to an Ioun Stone ends while it's orbiting your head, the stone falls as though you had dropped it.

Your Constitution increases by 2, to a maximum of 20, while this pink rhomboid orbits your head.

\DndItemHeader{Ioun Stone, Greater Absorption}{Wondrous item, Legendary (requires attunement)}
Roughly marble sized, \textit{Ioun Stones} are named after Ioun, a god of knowledge and prophecy revered on some worlds. Many types of \textit{Ioun Stones} exist, each type a distinct combination of shape and color.

When you take a Magic action to toss an \textit{Ioun Stone} into the air, the stone orbits your head at a distance of 1d3 feet, conferring its benefit to you while doing so. You can have up to three \textit{Ioun Stones} orbiting your head at the same time.

Each \textit{Ioun Stone} orbiting your head is considered to be an object you are wearing. The orbiting stone avoids contact with other creatures and objects, adjusting its orbit to avoid collisions and thwarting all attempts by other creatures to attack or snatch it.

As a Utilize action, you can seize and stow any number of \textit{Ioun Stones} orbiting your head. If your Attunement to an Ioun Stone ends while it's orbiting your head, the stone falls as though you had dropped it.

While this marbled lavender and green ellipsoid orbits your head, you can take a Reaction to cancel a spell of level 8 or lower cast by a creature you can see. A canceled spell has no effect, and any resources used to cast it are wasted. Once the stone has canceled 20 levels of spells, it burns out, turns dull gray, and loses its magic.

\DndItemHeader{Ioun Stone, Insight}{Wondrous item, Very Rare (requires attunement)}
Roughly marble sized, \textit{Ioun Stones} are named after Ioun, a god of knowledge and prophecy revered on some worlds. Many types of \textit{Ioun Stones} exist, each type a distinct combination of shape and color.

When you take a Magic action to toss an \textit{Ioun Stone} into the air, the stone orbits your head at a distance of 1d3 feet, conferring its benefit to you while doing so. You can have up to three \textit{Ioun Stones} orbiting your head at the same time.

Each \textit{Ioun Stone} orbiting your head is considered to be an object you are wearing. The orbiting stone avoids contact with other creatures and objects, adjusting its orbit to avoid collisions and thwarting all attempts by other creatures to attack or snatch it.

As a Utilize action, you can seize and stow any number of \textit{Ioun Stones} orbiting your head. If your Attunement to an Ioun Stone ends while it's orbiting your head, the stone falls as though you had dropped it.

Your Wisdom increases by 2, to a maximum of 20, while this incandescent blue sphere orbits your head.

\DndItemHeader{Ioun Stone, Intellect}{Wondrous item, Very Rare (requires attunement)}
Roughly marble sized, \textit{Ioun Stones} are named after Ioun, a god of knowledge and prophecy revered on some worlds. Many types of \textit{Ioun Stones} exist, each type a distinct combination of shape and color.

When you take a Magic action to toss an \textit{Ioun Stone} into the air, the stone orbits your head at a distance of 1d3 feet, conferring its benefit to you while doing so. You can have up to three \textit{Ioun Stones} orbiting your head at the same time.

Each \textit{Ioun Stone} orbiting your head is considered to be an object you are wearing. The orbiting stone avoids contact with other creatures and objects, adjusting its orbit to avoid collisions and thwarting all attempts by other creatures to attack or snatch it.

As a Utilize action, you can seize and stow any number of \textit{Ioun Stones} orbiting your head. If your Attunement to an Ioun Stone ends while it's orbiting your head, the stone falls as though you had dropped it.

Your Intelligence increases by 2, to a maximum of 20, while this marbled scarlet and blue sphere orbits your head.

\DndItemHeader{Ioun Stone, Leadership}{Wondrous item, Very Rare (requires attunement)}
Roughly marble sized, \textit{Ioun Stones} are named after Ioun, a god of knowledge and prophecy revered on some worlds. Many types of \textit{Ioun Stones} exist, each type a distinct combination of shape and color.

When you take a Magic action to toss an \textit{Ioun Stone} into the air, the stone orbits your head at a distance of 1d3 feet, conferring its benefit to you while doing so. You can have up to three \textit{Ioun Stones} orbiting your head at the same time.

Each \textit{Ioun Stone} orbiting your head is considered to be an object you are wearing. The orbiting stone avoids contact with other creatures and objects, adjusting its orbit to avoid collisions and thwarting all attempts by other creatures to attack or snatch it.

As a Utilize action, you can seize and stow any number of \textit{Ioun Stones} orbiting your head. If your Attunement to an Ioun Stone ends while it's orbiting your head, the stone falls as though you had dropped it.

Your Charisma increases by 2, to a maximum of 20, while this marbled pink and green sphere orbits your head.

\DndItemHeader{Ioun Stone, Mastery}{Wondrous item, Legendary (requires attunement)}
Roughly marble sized, \textit{Ioun Stones} are named after Ioun, a god of knowledge and prophecy revered on some worlds. Many types of \textit{Ioun Stones} exist, each type a distinct combination of shape and color.

When you take a Magic action to toss an \textit{Ioun Stone} into the air, the stone orbits your head at a distance of 1d3 feet, conferring its benefit to you while doing so. You can have up to three \textit{Ioun Stones} orbiting your head at the same time.

Each \textit{Ioun Stone} orbiting your head is considered to be an object you are wearing. The orbiting stone avoids contact with other creatures and objects, adjusting its orbit to avoid collisions and thwarting all attempts by other creatures to attack or snatch it.

As a Utilize action, you can seize and stow any number of \textit{Ioun Stones} orbiting your head. If your Attunement to an Ioun Stone ends while it's orbiting your head, the stone falls as though you had dropped it.

Your Proficiency increases by 1 while this pale green prism orbits your head.

\DndItemHeader{Ioun Stone, Protection}{Wondrous item, Rare (requires attunement)}
Roughly marble sized, \textit{Ioun Stones} are named after Ioun, a god of knowledge and prophecy revered on some worlds. Many types of \textit{Ioun Stones} exist, each type a distinct combination of shape and color.

When you take a Magic action to toss an \textit{Ioun Stone} into the air, the stone orbits your head at a distance of 1d3 feet, conferring its benefit to you while doing so. You can have up to three \textit{Ioun Stones} orbiting your head at the same time.

Each \textit{Ioun Stone} orbiting your head is considered to be an object you are wearing. The orbiting stone avoids contact with other creatures and objects, adjusting its orbit to avoid collisions and thwarting all attempts by other creatures to attack or snatch it.

As a Utilize action, you can seize and stow any number of \textit{Ioun Stones} orbiting your head. If your Attunement to an Ioun Stone ends while it's orbiting your head, the stone falls as though you had dropped it.

You gain a +1 bonus to Armor Class while this dusty-rose prism orbits your head.

\DndItemHeader{Ioun Stone, Regeneration}{Wondrous item, Legendary (requires attunement)}
Roughly marble sized, \textit{Ioun Stones} are named after Ioun, a god of knowledge and prophecy revered on some worlds. Many types of \textit{Ioun Stones} exist, each type a distinct combination of shape and color.

When you take a Magic action to toss an \textit{Ioun Stone} into the air, the stone orbits your head at a distance of 1d3 feet, conferring its benefit to you while doing so. You can have up to three \textit{Ioun Stones} orbiting your head at the same time.

Each \textit{Ioun Stone} orbiting your head is considered to be an object you are wearing. The orbiting stone avoids contact with other creatures and objects, adjusting its orbit to avoid collisions and thwarting all attempts by other creatures to attack or snatch it.

As a Utilize action, you can seize and stow any number of \textit{Ioun Stones} orbiting your head. If your Attunement to an Ioun Stone ends while it's orbiting your head, the stone falls as though you had dropped it.

You regain 15 Hit Points at the end of each hour this pearly white spindle orbits your head if you have at least 1 Hit Points.

\DndItemHeader{Ioun Stone, Reserve}{Wondrous item, Rare (requires attunement)}
Roughly marble sized, \textit{Ioun Stones} are named after Ioun, a god of knowledge and prophecy revered on some worlds. Many types of \textit{Ioun Stones} exist, each type a distinct combination of shape and color.

When you take a Magic action to toss an \textit{Ioun Stone} into the air, the stone orbits your head at a distance of 1d3 feet, conferring its benefit to you while doing so. You can have up to three \textit{Ioun Stones} orbiting your head at the same time.

Each \textit{Ioun Stone} orbiting your head is considered to be an object you are wearing. The orbiting stone avoids contact with other creatures and objects, adjusting its orbit to avoid collisions and thwarting all attempts by other creatures to attack or snatch it.

As a Utilize action, you can seize and stow any number of \textit{Ioun Stones} orbiting your head. If your Attunement to an Ioun Stone ends while it's orbiting your head, the stone falls as though you had dropped it.

This vibrant purple prism stores spells cast into it, holding them until you use them. The stone can store up to 4 levels of spells at a time. When found, it contains 1d4 levels of stored spells chosen by the DM.

Any creature can cast a spell of level 1 through 4 into the stone by touching it as the spell is cast. The spell has no effect, other than to be stored in the stone. If the stone can't hold the spell, the spell is expended without effect. The level of the slot used to cast the spell determines how much space it uses.

While this stone orbits your head, you can cast any spell stored in it. The spell uses the slot level, spell save DC, spell attack bonus, and spellcasting ability of the original caster but is otherwise treated as if you cast the spell. The spell cast from the stone is no longer stored in it, freeing up space. Strength (Very Rare). Your Strength increases by 2, to a maximum of 20, while this pale blue rhomboid orbits your head.

\DndItemHeader{Ioun Stone, Strength}{Wondrous item, Very Rare (requires attunement)}
Roughly marble sized, \textit{Ioun Stones} are named after Ioun, a god of knowledge and prophecy revered on some worlds. Many types of \textit{Ioun Stones} exist, each type a distinct combination of shape and color.

When you take a Magic action to toss an \textit{Ioun Stone} into the air, the stone orbits your head at a distance of 1d3 feet, conferring its benefit to you while doing so. You can have up to three \textit{Ioun Stones} orbiting your head at the same time.

Each \textit{Ioun Stone} orbiting your head is considered to be an object you are wearing. The orbiting stone avoids contact with other creatures and objects, adjusting its orbit to avoid collisions and thwarting all attempts by other creatures to attack or snatch it.

As a Utilize action, you can seize and stow any number of \textit{Ioun Stones} orbiting your head. If your Attunement to an Ioun Stone ends while it's orbiting your head, the stone falls as though you had dropped it.

Your Strength increases by 2, to a maximum of 20, while this pale blue rhomboid orbits your head.

\DndItemHeader{Ioun Stone, Sustenance}{Wondrous item, Rare (requires attunement)}
Roughly marble sized, \textit{Ioun Stones} are named after Ioun, a god of knowledge and prophecy revered on some worlds. Many types of \textit{Ioun Stones} exist, each type a distinct combination of shape and color.

When you take a Magic action to toss an \textit{Ioun Stone} into the air, the stone orbits your head at a distance of 1d3 feet, conferring its benefit to you while doing so. You can have up to three \textit{Ioun Stones} orbiting your head at the same time.

Each \textit{Ioun Stone} orbiting your head is considered to be an object you are wearing. The orbiting stone avoids contact with other creatures and objects, adjusting its orbit to avoid collisions and thwarting all attempts by other creatures to attack or snatch it.

As a Utilize action, you can seize and stow any number of \textit{Ioun Stones} orbiting your head. If your Attunement to an Ioun Stone ends while it's orbiting your head, the stone falls as though you had dropped it.

You don't need to eat or drink while this clear spindle orbits your head.

\DndItemHeader{Iron Bands of Bilarro}{Wondrous item, Rare}
This rusty iron sphere measures 3 inches in diameter and weighs 1 pound. You can take a Magic action to throw the sphere at a Huge or smaller creature you can see within 60 feet of yourself. As the sphere moves through the air, it opens into a tangle of metal bands.

Make a ranged attack roll with an attack bonus equal to your Dexterity modifier plus your Proficiency. On a hit, the target has the Restrained condition until you take a Bonus Action to issue a command that releases it. Doing so or missing with the attack causes the bands to contract and become a sphere once more.

A creature that can touch the bands, including the one Restrained, can take an action to make a DC 20 Strength (Athletics) check to break the iron bands. On a successful check, the item is destroyed, and the Restrained creature is freed. On a failed check, any further attempts made by that creature automatically fail until 24 hours have elapsed.

Once the bands are used, they can't be used again until the next dawn.

\DndItemHeader{Iron Flask}{Wondrous item, Legendary}
While holding this brass-stoppered iron flask, you can take a Magic action to target a creature that you can see within 60 feet of yourself. If the flask is empty and the target is native to a plane of existence other than the one you're on, the target must succeed on a DC 17 Wisdom saving throw or be trapped in the flask. If the target has been trapped by the flask before, it has Advantage on the save. Once trapped, a creature remains in the flask until released. The flask can hold only one creature at a time. A creature trapped in the flask doesn't age and doesn't need to breathe, eat, or drink.

You can take a Magic action to remove the flask's stopper and release the creature in the flask. The creature then obeys your commands for 1 hour, understanding those commands even if it doesn't know the language in which the commands are given. If you issue no commands or give the creature a command that is likely to result in its death or imprisonment, it defends itself but otherwise takes no actions. At the end of the duration, the creature acts in accordance with its normal disposition and alignment.

An \textit{Identify} spell reveals if the flask contains a creature, but the only way to determine the type of creature is to open the flask. A newly discovered Iron Flask might already contain a creature chosen by the DM or determined randomly by rolling on the following table (see the \textit{Monster Manual} for the creature's stat block).


\begin{DndTable}{cX}

1d100 & Contents\\
01-50 & No creature\\
51 & \textbf{Arcanaloth}\\
52-54 & \textbf{Bone Devil}\\
55-56 & \textbf{Cambion}\\
57-58 & \textbf{Dao}\\
59 & \textbf{Deva}\\
60-61 & \textbf{Djinni}\\
62-63 & \textbf{Efreeti}\\
64-65 & \textbf{Erinyes}\\
66-67 & \textbf{Fomorian}\\
68 & \textbf{Githyanki Knight}\\
69 & \textbf{Githzerai Zerth}\\
70-71 & \textbf{Glabrezu}\\
72-74 & \textbf{Hezrou}\\
75 & \textbf{Incubus}\\
76-77 & \textbf{Invisible Stalker}\\
78-79 & \textbf{Marid}\\
80 & \textbf{Marilith}\\
81-82 & \textbf{Mezzoloth}\\
83-84 & \textbf{Nalfeshnee}\\
85-86 & \textbf{Night Hag}\\
87-88 & \textbf{Nycaloth}\\
89 & \textbf{Planetar}\\
90-91 & \textbf{Red Slaad}\\
92-93 & \textbf{Salamander}\\
94 & \textbf{Solar}\\
95 & \textbf{Succubus}\\
96 & \textbf{Ultroloth}\\
97-99 & \textbf{Vrock}\\
00 & \textbf{Xorn}\\
\end{DndTable}

\DndItemHeader{Jacinth}{}
A fiery orange gemstone.

\DndItemHeader{Jade}{}
A light green, deep green, or white gemstone.

\DndItemHeader{Jasper}{}
A blue, black, or brown gemstone.

\DndItemHeader{Javelin of Lightning}{Uncommon}
Each time you make an attack roll with this magic weapon and hit, you can have it deal Lightning damage instead of Piercing damage.

\subparagraph{Lightning Bolt}
When you throw this weapon at a target no farther than 120 feet from you, you can forgo making a ranged attack roll and instead turn the weapon into a bolt of lightning. This bolt forms a 5-foot-wide Line [Area of Effect] between you and the target. The target and each other creature in the Line [Area of Effect] (excluding you) makes a DC 13 Dexterity saving throw, taking 4d6 Lightning damage on a failed save or half as much damage on a successful one. Immediately after dealing this damage, the weapon reappears in your hand. This property can't be used again until the next dawn.

\DndItemHeader{Jet}{}
A deep black gemstone.

\DndItemHeader{Jeweler's Tools}{}

\begin{description}
\item[Ability:]
Intelligence

\item[Utilize:]
Discern a gem's value (DC 15)

\item[Craft:]
\textit{Arcane Focus}, \textit{Holy Symbol}

\end{description}

\DndItemHeader{Keoghtom's Ointment}{Wondrous item, Uncommon}
This glass jar, 3 inches in diameter, contains 1d4 + 1 doses of a thick mixture that smells faintly of aloe. The jar and its contents weigh 1/2 pound.

As a Utilize action, you can swallow one dose of the ointment or apply it to a creature within 5 feet of yourself. The creature that receives it regains 2d8 + 2 Hit Points and ceases to have the Poisoned condition.

\DndItemHeader{Lantern of Revealing}{Wondrous item, Uncommon}
While lit, this hooded lantern burns for 6 hours on 1 pint of oil, shedding Bright Light in a 30-foot radius and Dim Light for an additional 30 feet. Invisible creatures and objects are visible as long as they are in the lantern's Bright Light. You can take a Utilize action to lower the hood, reducing the lantern's light to Dim Light in a 5-foot radius.

\DndItemHeader{Lapis Lazuli}{}
A light and dark blue with yellow flecks gemstone.

\DndItemHeader{Leatherworker's Tools}{}

\begin{description}
\item[Ability:]
Dexterity

\item[Utilize:]
Add a design to a leather item (DC 10)

\item[Craft:]
\textit{Sling}, \textit{Whip}, \textit{Hide Armor}, \textit{Leather Armor}, \textit{Studded Leather Armor}, \textit{Backpack}, \textit{Crossbow Bolt Case}, \textit{Map or Scroll Case}, \textit{Parchment}, \textit{Pouch}, \textit{Quiver}, \textit{Waterskin}

\end{description}

\DndItemHeader{Lock of Trickery}{Wondrous item, Common}
This lock appears to be an ordinary \textit{Lock} (of the type described in chapter 6 of the \textit{Player's Handbook}) and comes with a single key. The tumblers in this lock magically adjust to thwart burglars. Dexterity checks made to pick the lock have Disadvantage.

\DndItemHeader{Lolth's Sting}{}
A creature subjected to Lolth's Sting must succeed on a DC 13 Constitution saving throw or have the Poisoned condition for 1 hour. If the creature fails the save by 5 or more, the creature also has the Unconscious condition while Poisoned in this way. The creature wakes up if it takes damage or if another creature takes an action to shake it awake.

\DndItemHeader{Lute of Thunderous Thumping}{Very Rare}
This reinforced \textit{lute} can be wielded as a magic Club that deals an extra 2d8 Thunder damage on a hit.

\subparagraph{Sing and Swing}
If you're a Bard, you can use your Charisma modifier instead of your Strength modifier when making a melee attack roll with the \textit{lute}, provided you sing or hum while making the attack.

\DndItemHeader{Mace of Disruption}{Rare (requires attunement)}
When you hit a Fiend or an Undead with this magic weapon, that creature takes an extra 2d6 Radiant damage. If the target has 25 Hit Points or fewer after taking this damage, it must succeed on a DC 15 Wisdom saving throw or be destroyed. On a successful save, the creature has the Frightened condition until the end of your next turn.

\subparagraph{Light}
While you hold this weapon, it sheds Bright Light in a 20-foot radius and Dim Light for an additional 20 feet.

\DndItemHeader{Mace of Smiting}{Rare}
You gain a +1 bonus to attack rolls and damage rolls made with this magic weapon. The bonus increases to +3 when you use the weapon to attack a Construct.

When you roll a 20 on an attack roll made with this weapon, the target takes an extra 7 Bludgeoning damage, or 14 Bludgeoning damage if it's a Construct. If a Construct has 25 Hit Points or fewer after taking this damage, it is destroyed.

\DndItemHeader{Mace of Terror}{Rare (requires attunement)}
This magic weapon has 3 charges and regains 1d3 expended charges daily at dawn. While holding the weapon, you can take a Magic action and expend 1 charge to release a wave of terror from it. Each creature of your choice within 30 feet of you must succeed on a DC 15 Wisdom saving throw or have the Frightened condition for 1 minute. While Frightened in this way, a creature must spend its turns trying to move as far away from you as it can, and it can't make Opportunity Attacks. For its action, it can use only the Dash action or try to escape from an effect that prevents it from moving. If it has nowhere it can move, the creature can take the Dodge action. At the end of each of its turns, a creature repeats the save, ending the effect on itself on a success.

\DndItemHeader{Malachite}{}
A striated light and dark green gemstone.

\DndItemHeader{Malice}{}
A creature subjected to Malice must succeed on a DC 15 Constitution saving throw or have the Poisoned condition for 1 hour. The creature also has the Blinded condition while Poisoned in this way.

\DndItemHeader{Mantle of Spell Resistance}{Wondrous item, Rare (requires attunement)}
You have Advantage on saving throws against spells while you wear this cloak.

\DndItemHeader{Manual of Bodily Health}{Wondrous item, Very Rare}
This book contains health and diet tips, and its words are charged with magic. If you spend 48 hours over a period of 6 days or fewer studying the book's contents and practicing its guidelines, your Constitution score increases by 2, as does your maximum for that score. The manual then loses its magic, but regains it in a century.

\DndItemHeader{Manual of Clay Golems}{Wondrous item, Very Rare}
This tome contains information and incantations necessary to make a \textbf{clay golem}. To decipher and use the manual, you must be a spellcaster with at least two 5th-level spell slots. A creature that can't use a \textit{manual of golems} and attempts to read it takes 6d6 psychic damage.

To create a clay golem, you must spend 30 days, working without interruption with the manual at hand and resting no more than 8 hours per day. You must also pay 65,000 gp to purchase supplies.

Once you finish creating the golem, the book is consumed in eldritch flames. The golem becomes animate when the ashes of the manual are sprinkled on it. It is under your control, and it understands and obeys your spoken commands.

\DndItemHeader{Manual of Flesh Golems}{Wondrous item, Very Rare}
This tome contains information and incantations necessary to make a \textbf{flesh golem}. To decipher and use the manual, you must be a spellcaster with at least two 5th-level spell slots. A creature that can't use a \textit{manual of golems} and attempts to read it takes 6d6 psychic damage.

To create a flesh golem, you must spend 60 days, working without interruption with the manual at hand and resting no more than 8 hours per day. You must also pay 50,000 gp to purchase supplies.

Once you finish creating the golem, the book is consumed in eldritch flames. The golem becomes animate when the ashes of the manual are sprinkled on it. It is under your control, and it understands and obeys your spoken commands.

\DndItemHeader{Manual of Gainful Exercise}{Wondrous item, Very Rare}
This book describes fitness exercises, and its words are charged with magic. If you spend 48 hours over a period of 6 days or fewer studying the book's contents and practicing its guidelines, your Strength score increases by 2, as does your maximum for that score. The manual then loses its magic, but regains it in a century.

\DndItemHeader{Manual of Golems}{Wondrous item, Very Rare}
This tome contains information and incantations necessary to make a particular type of golem. To decipher and use the manual, you must be a spellcaster with at least two 5th-level spell slots. A creature that can't use a manual of golems and attempts to read it takes 6d6 psychic damage.


\begin{DndTable}{cXcX}

d20 & Golem & Time & Cost\\
1-5 & \textit{Manual of Clay Golems} & 30 days & 65,000 GP\\
6-17 & \textit{Manual of Flesh Golems} & 60 days & 50,000 GP\\
18 & \textit{Manual of Iron Golems} & 120 days & 100,000 GP\\
19-20 & \textit{Manual of Stone Golems} & 90 days & 80,000 GP\\
\end{DndTable}

To create a golem, you must spend the time shown on the table, working without interruption with the manual at hand and resting no more than 8 hours per day. You must also pay the specified cost to purchase supplies.

Once you finish creating the golem, the book is consumed in eldritch flames. The golem becomes animate when the ashes of the manual are sprinkled on it. It is under your control, and it understands and obeys your spoken commands.

\DndItemHeader{Manual of Iron Golems}{Wondrous item, Very Rare}
This tome contains information and incantations necessary to make a \textbf{iron golem}. To decipher and use the manual, you must be a spellcaster with at least two 5th-level spell slots. A creature that can't use a \textit{manual of golems} and attempts to read it takes 6d6 psychic damage.

To create an iron golem, you must spend 120 days, working without interruption with the manual at hand and resting no more than 8 hours per day. You must also pay 100,000 gp to purchase supplies.

Once you finish creating the golem, the book is consumed in eldritch flames. The golem becomes animate when the ashes of the manual are sprinkled on it. It is under your control, and it understands and obeys your spoken commands.

\DndItemHeader{Manual of Quickness of Action}{Wondrous item, Very Rare}
This book contains coordination and balance exercises, and its words are charged with magic. If you spend 48 hours over a period of 6 days or fewer studying the book's contents and practicing its guidelines, your Dexterity score increases by 2, as does your maximum for that score. The manual then loses its magic, but regains it in a century.

\DndItemHeader{Manual of Stone Golems}{Wondrous item, Very Rare}
This tome contains information and incantations necessary to make a \textbf{stone golem}. To decipher and use the manual, you must be a spellcaster with at least two 5th-level spell slots. A creature that can't use a \textit{manual of golems} and attempts to read it takes 6d6 psychic damage.

To create a stone golem, you must spend 90 days, working without interruption with the manual at hand and resting no more than 8 hours per day. You must also pay 80,000 gp to purchase supplies.

Once you finish creating the golem, the book is consumed in eldritch flames. The golem becomes animate when the ashes of the manual are sprinkled on it. It is under your control, and it understands and obeys your spoken commands.

\DndItemHeader{Medallion of Thoughts}{Wondrous item, Uncommon (requires attunement)}
The medallion has 5 charges. While wearing it, you can expend 1 charge to cast \textit{Detect Thoughts} (save DC 13) from it. The medallion regains 1d4 expended charges daily at dawn.

\DndItemHeader{Midnight Tears}{}
A creature that ingests Midnight Tears suffers no effect until the stroke of midnight. Any effect that ends the Poisoned condition neutralizes this poison. If the poison hasn't been neutralized before midnight, the creature makes a DC 17 Constitution saving throw, taking 31 (9d6) Poison damage on a failed save or half as much damage on a successful one.

\DndItemHeader{Mirror of Life Trapping}{Wondrous item, Very Rare}
When this 4-foot-tall, 2-foot-wide mirror is viewed indirectly, its surface shows faint images of creatures. The mirror weighs 50 pounds, and it has AC 11, HP 10, Immunity to Poison and Psychic damage, and Vulnerability to Bludgeoning damage. It shatters and is destroyed when reduced to 0 Hit Points.

If the mirror is hanging on a vertical surface and you are within 5 feet of it, you can take a Magic action and use a command word to activate it. It remains activated until you take a Magic action and repeat the command word to deactivate it.

Any creature other than you that sees its reflection in the activated mirror while within 30 feet of the mirror must succeed on a DC 15 Charisma saving throw or be trapped, along with anything it is wearing or carrying, in one of the mirror's twelve extradimensional cells. A creature that knows the mirror's nature makes the save with Advantage, and Constructs succeed on the save automatically.

An extradimensional cell is an infinite expanse filled with thick fog that reduces visibility to 10 feet. Creatures trapped in the mirror's cells don't age, and they don't need to eat, drink, or sleep. A creature trapped within a cell can escape using magic that permits planar travel. Otherwise, the creature is confined to the cell until freed.

If the mirror traps a creature but its twelve extradimensional cells are already occupied, the mirror frees one trapped creature at random to accommodate the new prisoner. A freed creature appears in an unoccupied space within sight of the mirror but facing away from it. If the mirror is shattered, all creatures it contains are freed and appear in unoccupied spaces near it.

While within 5 feet of the mirror, you can take a Magic action to name one creature trapped in it or call out a particular cell by number. The creature named or contained in the named cell appears as an image on the mirror's surface. You and the creature can then communicate.

In a similar way, you can take a Magic action and use a second command word to free one creature trapped in the mirror. The freed creature appears, along with its possessions, in the unoccupied space nearest to the mirror and facing away from it.

Placing the mirror inside an extradimensional space created by a \textit{Bag of Holding}, Portable Hole, or similar item instantly destroys both items and opens a gate to the Astral Plane. The gate originates where the one item was placed inside the other. Any creature within 10 feet of the gate and not behind Cover is sucked through it to a random location on the Astral Plane. The gate then closes. The gate is one-way only and can't be reopened.

\DndItemHeader{Moonstone}{}
A white with pale-blue glow gemstone.

\DndItemHeader{Moss Agate}{}
A pink or yellow white with mossy gray or green markings gemstone.

\DndItemHeader{Mystery Key}{Wondrous item, Common}
A question mark is worked into the head of this key. The key has a 5 chance of unlocking any lock into which it's inserted. Once it unlocks something, the key disappears.

\DndItemHeader{Nature's Mantle}{Wondrous item, Uncommon (requires attunement by a druid or ranger)}
This cloak shifts color and texture to blend with the terrain surrounding you. While wearing the cloak, you can use it as a Spellcasting Focus for your Druid and Ranger spells.

While you are in an area that is Lightly Obscured, you can Hide as a Bonus Action even if you are being directly observed.

\DndItemHeader{Necklace of Adaptation}{Wondrous item, Uncommon (requires attunement)}
While wearing this necklace, you can breathe normally in any environment, and you have Advantage on saving throws made to avoid or end the Poisoned condition.

\DndItemHeader{Necklace of Fireballs}{Wondrous item, Rare}
This necklace has 1d6 + 3 beads hanging from it. You can take a Magic action to detach a bead and throw it up to 60 feet away. When it reaches the end of its trajectory, the bead detonates as a level 3 \textit{Fireball} (save DC 15).

You can hurl multiple beads, or even the whole necklace, at one time. When you do so, increase the damage of the Fireball by 1d6 for each bead after the first (maximum 12d6).

\DndItemHeader{Necklace of Prayer Beads}{Wondrous item, Rare (requires attunement by a cleric, druid, or paladin)}
This necklace has 1d4 + 2 magic beads made from aquamarine, black pearl, or topaz. It also has many nonmagical beads made from stones such as amber, bloodstone, citrine, coral, jade, pearl, or quartz. If a magic bead is removed from the necklace, that bead loses its magic.

Six types of magic beads exist. The DM decides the type of each bead on the necklace or determines it randomly by rolling on the table below. A necklace can have more than one bead of the same type. To use one, you must be wearing the necklace. Each bead contains a spell that you can cast from it as a Bonus Action (using your spell save DC if a save is necessary). Once a magic bead's spell is cast, that bead can't be used again until the next dawn.


\begin{DndTable}{cXX}

1d20 & Bead & Spell\\
1-6 & Bead of Blessing & \textit{Bless}\\
7-12 & Bead of Curing & \textit{Cure Wounds} (level 2 version)\\
13-16 & Bead of Favor & \textit{Greater Restoration}\\
17-18 & Bead of Smiting & \textit{Shining Smite}\\
19 & Bead of Summons & \textit{Guardian of Faith}\\
20 & Bead of Wind Walking & \textit{Wind Walk}\\
\end{DndTable}

\DndItemHeader{Nolzur's Marvelous Pigments}{Wondrous item, Very Rare}
This fine wooden box contains 1d4 pots of pigment and a brush (weighing 1 pound in total). Using the brush and expending 1 pot of pigment, you can paint any number of three-dimensional objects and terrain features (such as walls, doors, trees, flowers, weapons, webs, and pits), provided these elements are all confined to a 20-foot Cube [Area of Effect]. The effort takes 10 minutes (regardless of the number of elements you create), during which time you must remain in the Cube [Area of Effect], and requires Concentration. If your Concentration is broken or you leave the Cube [Area of Effect] before the work is done, all the painted elements vanish, and the pot of pigment is wasted.

When the work is done, all the painted objects and terrain features become real. Thus, painting a door on a wall creates an actual door, which can be opened to whatever is beyond. Painting a pit creates a real pit, the entire depth of which must lie within the 20-foot Cube [Area of Effect].

No object created by a pot of pigment can have a value greater than 25 GP, and the total value of all objects created by a pot of pigment can't exceed 500 GP. If you paint objects of greater value (such as a large pile of gold), they look authentic, but close inspection reveals they're made from paste, cookies, or some other worthless material.

If you paint a form of energy such as fire or lightning, the energy dissipates as soon as you complete the painting, doing no harm.

\DndItemHeader{Obsidian}{}
A black gemstone.

\DndItemHeader{Oil of Etherealness}{Rare}
One vial of this oil can cover one Medium or smaller creature, along with the equipment it's wearing and carrying (one additional vial is required for each size category above Medium). Applying the oil takes 10 minutes. The affected creature then gains the effect of the \textit{Etherealness} spell for 1 hour.

Beads of this cloudy, gray oil form on the outside of its container and quickly evaporate.

\DndItemHeader{Oil of Sharpness}{Very Rare}
One vial of this oil can coat one Melee weapon or twenty pieces of ammunition, but only ammunition and Melee weapons that are nonmagical and deal Slashing or Piercing damage are affected. Applying the oil takes 1 minute, after which the oil magically seeps into whatever it coats, turning the coated weapon into a +3 Weapon or the coated ammunition into +3 Ammunition.

This clear, gelatinous oil sparkles with tiny, ultrathin silver shards.

\DndItemHeader{Oil of Slipperiness}{Uncommon}
One vial of this oil can cover one Medium or smaller creature, along with the equipment it's wearing and carrying (one additional vial is required for each size category above Medium). Applying the oil takes 10 minutes. The affected creature then gains the effect of the \textit{Freedom of Movement} spell for 8 hours.

Alternatively, the oil can be poured on the ground as a Magic action, where it covers a 10-foot square, duplicating the effect of the \textit{Grease} spell in that area for 8 hours.

This sticky, black unguent is thick and heavy, but it flows quickly when poured.

\DndItemHeader{Oil of Taggit}{}
A creature subjected to Oil of Taggit must succeed on a DC 13 Constitution saving throw or have the Poisoned condition for 24 hours. The creature also has the Unconscious condition while Poisoned in this way. It wakes up if it takes damage.

\DndItemHeader{Onyx}{}
A bands of black and white, or pure black or white gemstone.

\DndItemHeader{Opal}{}
A pale blue with green and golden mottling gemstone.

\DndItemHeader{Orb of Direction}{Wondrous item, Common}
This orb can be used as an Arcane Focus.

While holding this orb, you can take a Magic action to determine which way is magnetic north. Nothing happens if the orb is used in a location that has no magnetic north.

\DndItemHeader{Orb of Dragonkind}{Wondrous item, Artifact (requires attunement)}
Long ago, in the Dragonlance setting, elves and humans waged a terrible war against chromatic dragons. When the world seemed doomed, the wizards of the Towers of High Sorcery came together and forged five Orbs of Dragonkind to help defeat the dragons. One orb was taken to each of the five towers, and there they were used to speed the war toward a victorious end. The wizards used the orbs to lure dragons to them, then destroyed the dragons with powerful magic.

As the Towers of High Sorcery fell in later ages, the orbs were destroyed or faded into legend, and only three are thought to survive. Their magic has been warped over the centuries. Their primary purpose of calling dragons still functions, but they also allow some measure of control over dragons.

Each orb contains the essence of an evil dragon, a presence that resents any attempt to coax magic from it. Those who try to wield an orb's magic but lack sufficient force of personality might find themselves under the orb's control.

An orb is an etched crystal globe about 10 inches in diameter. When used, it grows to about 20 inches in diameter, and mist swirls inside it.

While attuned to an orb, you can take a Magic action to peer into the orb's depths. You must then make a DC 15 Charisma saving throw. On a successful save, you control the orb for as long as you remain attuned to it. On a failed save, the orb imposes the Charmed condition on you for as long as you remain attuned to it.

While you are Charmed by the orb, you can't voluntarily end your Attunement to it, and the orb casts \textit{Suggestion} on you at will (save DC 18), urging you to work toward the evil ends it desires. The dragon essence within the orb might want many things: the annihilation of a particular society or organization, freedom from the orb, to spread suffering in the world, to advance the worship of \textit{Tiamat}, or something else the DM decides.

\subparagraph{Random Properties}
An Orb of Dragonkind has the following random properties:


\begin{itemize}
\item 2 Artifact Properties; Minor Beneficial Properties properties
\item 1 Artifact Properties; Minor Detrimental Properties property
\item 1 Artifact Properties; Major Detrimental Properties property
\end{itemize}

\subparagraph{Spells}
The orb has 7 charges and regains 1d4 + 3 expended charges daily at dawn. If you control the orb, you can cast one of the spells on the following table from it. The table indicates how many charges you must expend to cast the spell.


\begin{DndTable}{Xc}

Spell & Charge Cost\\
\textit{Cure Wounds} (level 9 version) & 4\\
\textit{Daylight} & 1\\
\textit{Death Ward} & 2\\
\textit{Detect Magic} & 0\\
\textit{Scrying} (save DC 18) & 3\\
\end{DndTable}

\subparagraph{Call Dragons}
While you control the orb, you can take a Magic action to cause the orb to issue a telepathic call that extends in all directions for 40 miles. Chromatic dragons in range feel compelled to come to the orb as soon as possible by the most direct route. Dragon deities such as \textit{Tiamat} are unaffected by this call. Chromatic dragons drawn to the orb might be Hostile [Attitude] toward you for compelling them against their will. Once you have used this property, it can't be used again for 1 hour.

\subparagraph{Destroying an Orb}
An Orb of Dragonkind has AC 20 and is destroyed if it takes damage from a +3 Weapon or a \textit{Disintegrate} spell. Nothing else can harm it.

\DndItemHeader{Orb of Time}{Wondrous item, Common}
This orb can be used as an Arcane Focus.

While holding the orb, you can take a Magic action to determine whether it is morning, afternoon, evening, or nighttime. This property functions only on the Material Plane.

\DndItemHeader{Pale Tincture}{}
A creature subjected to Pale Tincture must succeed on a DC 16 Constitution saving throw or take 3 (1d6) Poison damage and have the Poisoned condition. The Poisoned creature repeats the save every 24 hours, taking 3 (1d6) Poison damage on a failed save. The damage the poison deals can't be healed by any means while the creature remains Poisoned. After seven successful saves against the poison, the creature is no longer Poisoned.

\DndItemHeader{Pearl}{}
A lustrous white, yellow, or pink gemstone.

\DndItemHeader{Pearl of Power}{Wondrous item, Uncommon (requires attunement by a spellcaster)}
While this pearl is on your person, you can take a Magic action to regain one expended spell slot of level 3 or lower. Once you use the pearl, it can't be used again until the next dawn.

\DndItemHeader{Perfume of Bewitching}{Wondrous item, Common}
This tiny vial contains magic perfume, enough for one use. You can take a Magic action to apply the perfume to yourself, and its effect lasts 1 hour. For the duration, you have Advantage on all Charisma (Deception and Persuasion) checks made to influence a creature within 5 feet of yourself.

\DndItemHeader{Periapt of Health}{Wondrous item, Uncommon (requires attunement)}
While wearing this pendant, you can take a Magic action to regain 2d4 + 2 Hit Points. Once used, this property can't be used again until the next dawn.

In addition, you have Advantage on saving throws to avoid or end the Poisoned condition while you wear this pendant.

\DndItemHeader{Periapt of Proof against Poison}{Wondrous item, Rare (requires attunement)}
This delicate silver chain has a brilliant-cut black gem pendant. While you wear it, you have Immunity to the Poisoned condition and Poison damage

\DndItemHeader{Periapt of Proof against Poison}{Wondrous item, Rare (requires attunement)}
This delicate silver chain has a brilliant-cut black gem pendant. While you wear it, you have Immunity to the Poisoned condition and Poison damage

\DndItemHeader{Periapt of Wound Closure}{Wondrous item, Uncommon (requires attunement)}
While wearing this pendant, you gain the following benefits.

\subparagraph{Life Preservation}
Whenever you make a Death Saving Throw, you can change a roll of 9 or lower to a 10, turning a failed save into a successful one.

\subparagraph{Natural Healing Boost}
Whenever you roll a Hit Points Die to regain Hit Points, double the number of Hit Points it restores.

\DndItemHeader{Peridot}{}
A rich olive green gemstone.

\DndItemHeader{Philter of Love}{Uncommon}
The next time you see a creature within 10 minutes after drinking this philter, you are charmed by that creature and have the Charmed condition for 1 hour.

This rose-hued, effervescent liquid contains one easy-to-miss bubble shaped like a heart.

\DndItemHeader{Pipe of Smoke Monsters}{Wondrous item, Common}
While smoking this pipe, you can take a Magic action to exhale a puff of smoke that takes the form of a creature, such as a dragon, a flumph, or a slaad. The form must be small enough to fit in a 1-foot cube and loses its shape after a few seconds, becoming an ordinary puff of smoke.

\DndItemHeader{Pipes of Haunting}{Wondrous item, Uncommon}
These pipes have 3 charges and regain 1d3 expended charges daily at dawn. You can take a Magic action to play them and expend 1 charge to create an eerie, spellbinding tune. Each creature of your choice within 30 feet of you must succeed on a DC 15 Wisdom saving throw or have the Frightened condition for 1 minute. A creature that fails the save repeats it at the end of each of its turns, ending the effect on itself on a success. A creature that succeeds on its save is immune to the effect of these pipes for 24 hours.

\DndItemHeader{Pipes of the Sewers}{Wondrous item, Uncommon (requires attunement)}
While these pipes are on your person, ordinary rats and giant rats are Indifferent [Attitude] toward you and won't attack you unless you threaten or harm them.

The pipes have 3 charges and regain 1d3 expended charges daily at dawn. If you play the pipes as a Magic action, you can take a Bonus Action to expend 1 to 3 charges, calling forth one \textbf{Swarm of Rats} with each expended charge if enough rats are within half a mile of you to be called in this fashion (as determined by the DM). If there aren't enough rats to form a swarm, the charge is wasted. Called swarms move toward the music by the shortest available route but aren't under your control otherwise.

Whenever a Swarm of Rats that isn't under another creature's control comes within 30 feet of you while you are playing the pipes, the swarm makes a DC 15 Wisdom saving throw. On a successful save, the swarm behaves as it normally would and can't be swayed by the pipes' music for the next 24 hours. On a failed save, the swarm is swayed by the pipes' music and becomes Friendly [Attitude] to you and your allies for as long as you continue to play the pipes each round as a Magic action. A Friendly [Attitude] swarm obeys your commands. If you issue no commands to a Friendly [Attitude] swarm, it defends itself but otherwise takes no actions. If a Friendly [Attitude] swarm starts its turn more than 30 feet away from you, your control over that swarm ends, and the swarm behaves as it normally would and can't be swayed by the pipes' music for the next 24 hours.

\DndItemHeader{Poisoner's Kit}{}

\begin{description}
\item[Ability:]
Intelligence

\item[Utilize:]
Detect a poisoned object (DC 10)

\item[Craft:]
\textit{Basic Poison}

\end{description}

\DndItemHeader{Pole of Angling}{Wondrous item, Common}
This item functions as a Pole. While holding it, you can take a Magic action to cause it to transform into a fishing pole with a hook, a line, and a reel, or have the fishing pole revert to a Pole.

\DndItemHeader{Pole of Collapsing}{Wondrous item, Common}
This item functions as a Pole. While holding it, you can take a Magic action to collapse it into a 1-foot-long rod for ease of storage (the pole's weight doesn't change) or cause the 1-foot-long rod to revert to a Pole. The rod elongates only as far as the surrounding space allows.

\DndItemHeader{Portable Hole}{Wondrous item, Rare}
This fine black cloth, soft as silk, is folded up to the dimensions of a handkerchief. It unfolds into a circular sheet 6 feet in diameter.

You can take a Magic action to unfold a Portable Hole and place it on or against a solid surface, whereupon the Portable Hole creates an extradimensional hole 10 feet deep. The cylindrical space within the hole exists on a different plane of existence, so it can't be used to create open passages. Any creature inside an open Portable Hole can exit the hole by climbing out of it.

You can take a Magic action to close a Portable Hole by taking hold of the edges of the cloth and folding it up. Folding the cloth closes the hole, and any creatures or objects within remain in the extradimensional space. No matter what's in it, the hole weighs next to nothing.

If the hole is folded up, a creature within the hole's extradimensional space can take an action to make a DC 10 Strength (Athletics) check. On a successful check, the creature forces its way out and appears within 5 feet of the Portable Hole. A closed Portable Hole holds enough air for 1 hour of breathing, divided by the number of breathing creatures inside.

Placing a Portable Hole inside an extradimensional space created by a \textit{Bag of Holding}, \textit{Heward's Handy Haversack}, or similar item instantly destroys both items and opens a gate to the Astral Plane. The gate originates where the one item was placed inside the other. Any creature within 10 feet of the gate and not behind Cover is sucked through it and deposited in a random location on the Astral Plane. The gate then closes. The gate is one-way only and can't be reopened.

\DndItemHeader{Pot of Awakening}{Wondrous item, Common}
If you plant an ordinary shrub in this 10-pound clay pot and let it grow for 30 days, the shrub magically transforms into an \textbf{Awakened Shrub} at the end of that time. When the shrub awakens, its roots break the pot, destroying it.

The awakened shrub is Friendly [Attitude] toward you and obeys your commands. Absent commands from you, it does nothing.

\DndItemHeader{Potion of Animal Friendship}{Uncommon}
When you drink this potion, you can cast the level 3 version of the \textit{Animal Friendship} spell (save DC 13).

Agitating this potion's muddy liquid brings little bits into view: a fish scale, a hummingbird feather, a cat claw, or a squirrel hair.

\DndItemHeader{Potion of Clairvoyance}{Rare}
When you drink this potion, you gain the effect of the \textit{Clairvoyance} spell (no Concentration required).

An eyeball bobs in this potion's yellowish liquid but vanishes when the potion is opened.

\DndItemHeader{Potion of Climbing}{Common}
When you drink this potion, you gain a Climb Speed equal to your Speed for 1 hour. During this time, you have Advantage on Strength (Athletics) checks to climb.

This potion is separated into brown, silver, and gray layers resembling bands of stone. Shaking the bottle fails to mix the colors.

\DndItemHeader{Potion of Cloud Giant Strength}{Very Rare}
When you drink this potion, your Strength score changes to 27 for 1 hour. The potion has no effect on you if your Strength is equal to or greater than that score.

This potion's transparent liquid has floating in it a sliver of fingernail from a cloud giant.

\DndItemHeader{Potion of Comprehension}{Common}
When you drink this potion, you gain the effect of the \textit{Comprehend Languages} spell for 1 hour.

This potion's liquid is a clear concoction with bits of salt and soot swirling in it.

\DndItemHeader{Potion of Diminution}{Rare}
When you drink this potion, you gain the "reduce" effect of the \textit{Enlarge/Reduce} spell for 1d4 hours (no Concentration required).

The red in the potion's liquid continuously contracts to a tiny bead and then expands to color the clear liquid around it. Shaking the bottle fails to interrupt this process.

\DndItemHeader{Potion of Fire Breath}{Uncommon}
After drinking this potion, you can take a Bonus Action to exhale fire at a target within 30 feet of yourself. The target makes a DC 13 Dexterity saving throw, taking 4d6 Fire damage on a failed save or half as much damage on a successful one. The effect ends after you exhale the fire three times or when 1 hour has passed.

This potion's orange liquid flickers, and smoke fills the top of the container and wafts out whenever it is opened.

\DndItemHeader{Potion of Fire Giant Strength}{Rare}
When you drink this potion, your Strength score changes to 25 for 1 hour. The potion has no effect on you if your Strength is equal to or greater than that score.

This potion's transparent liquid has floating in it a sliver of fingernail from a fire giant.

\DndItemHeader{Potion of Flying}{Very Rare}
When you drink this potion, you gain a Fly Speed equal to your Speed for 1 hour and can hover. If you're in the air when the potion wears off, you fall unless you have some other means of staying aloft.

This potion's clear liquid floats at the top of its container and has cloudy white impurities drifting in it.

\DndItemHeader{Potion of Frost Giant Strength}{Rare}
When you drink this potion, your Strength score changes to 23 for 1 hour. The potion has no effect on you if your Strength is equal to or greater than that score.

This potion's transparent liquid has floating in it a sliver of fingernail from a frost giant.

\DndItemHeader{Potion of Gaseous Form}{Rare}
When you drink this potion, you gain the effect of the \textit{Gaseous Form} spell for 1 hour (no Concentration required) or until you end the effect as a Bonus Action.

This potion's container seems to hold fog that moves and pours like water.

\DndItemHeader{Potion of Greater Healing}{Uncommon}
You regain 4d4 + 4 Hit Points when you drink this potion. The potion's red liquid glimmers when agitated.

\DndItemHeader{Potion of Greater Invisibility}{Very Rare}
This potion's container looks empty but feels as though it holds liquid. When you drink the potion, you have the Invisible condition for 1 hour.

\DndItemHeader{Potion of Growth}{Uncommon}
When you drink this potion, you gain the "enlarge" effect of the \textit{Enlarge/Reduce} spell for 10 minutes (no Concentration required).

The red in the potion's liquid continuously expands from a tiny bead to color the clear liquid around it and then contracts. Shaking the bottle fails to interrupt this process.

\DndItemHeader{Potion of Healing}{Common}
This potion is a magic item. As a Bonus Action, you can drink it or administer it to another creature within 5 feet of yourself. The creature that drinks the magical red fluid in this vial regains 2d4 + 2 Hit Points. The potion's red liquid glimmers when agitated.

\DndItemHeader{Potion of Heroism}{Rare}
When you drink this potion, you gain 10 Temporary Hit Points that last for 1 hour. For the same duration, you are under the effect of the \textit{Bless} spell (no Concentration required).

This potion's blue liquid bubbles and steams as if boiling.

\DndItemHeader{Potion of Hill Giant Strength}{Uncommon}
When you drink this potion, your Strength score changes to 21 for 1 hour. The potion has no effect on you if your Strength is equal to or greater than that score.

This potion's transparent liquid has floating in it a sliver of fingernail from a hill giant.

\DndItemHeader{Potion of Invisibility}{Rare}
This potion's container looks empty but feels as though it holds liquid. When you drink the potion, you have the Invisible condition for 1 hour. The effect ends early if you make an attack roll, deal damage, or cast a spell.

\DndItemHeader{Potion of Invulnerability}{Rare}
For 1 minute after you drink this potion, you have Resistance to all damage.

This potion's syrupy liquid looks like liquefied iron.

\DndItemHeader{Potion of Longevity}{Very Rare}
When you drink this potion, your physical age is reduced by 1d6 + 6 years, to a minimum of 13 years. Each time you subsequently drink a Potion of Longevity, there is 10 percent cumulative chance that you instead age by 1d6 + 6 years.

Suspended in this amber liquid is a tiny heart that, against all reason, is still beating. These ingredients vanish when the potion is opened.

\DndItemHeader{Potion of Mind Reading}{Rare}
When you drink this potion, you gain the effect of the \textit{Detect Thoughts} spell (save DC 13) for 10 minutes (no Concentration required).

This potion's dense, purple liquid has an ovoid cloud of pink floating in it.

\DndItemHeader{Potion of Poison}{Uncommon}
This concoction looks, smells, and tastes like a \textit{Potion of Healing} or another beneficial potion. However, it is actually poison masked by illusion magic. \textit{Identify} reveals its true nature.

If you drink this potion, you take 4d6 Poison damage and must succeed on a DC 13 Constitution saving throw or have the Poisoned condition for 1 hour.

\DndItemHeader{Potion of Pugilism}{Uncommon}
After you drink this potion, each Unarmed Strike you make deals an extra 1d6 Force damage on a hit. This effect lasts 10 minutes.

This potion is a thick green fluid that tastes like spinach.

\DndItemHeader{Potion of Resistance}{Uncommon}
When you drink this potion, you have Resistance to {{getFullImmRes item.resist}} damage for 1 hour.

\DndItemHeader{Potion of Speed}{Very Rare}
When you drink this potion, you gain the effect of the \textit{Haste} spell for 1 minute (no Concentration required) without suffering the wave of lethargy that typically occurs when the effect ends.

This potion's yellow fluid is streaked with black and swirls on its own.

\DndItemHeader{Potion of Stone Giant Strength}{Rare}
When you drink this potion, your Strength score changes to 23 for 1 hour. The potion has no effect on you if your Strength is equal to or greater than that score.

This potion's transparent liquid has floating in it a sliver of fingernail from a stone giant.

\DndItemHeader{Potion of Storm Giant Strength}{Legendary}
When you drink this potion, your Strength score changes to 29 for 1 hour. The potion has no effect on you if your Strength is equal to or greater than that score.

This potion's transparent liquid has floating in it a sliver of fingernail from a storm giant.

\DndItemHeader{Potion of Superior Healing}{Rare}
You regain 8d4 + 8 Hit Points when you drink this potion. The potion's red liquid glimmers when agitated.

\DndItemHeader{Potion of Supreme Healing}{Very Rare}
You regain 10d4 + 20 Hit Points when you drink this potion. The potion's red liquid glimmers when agitated.

\DndItemHeader{Potion of Vitality}{Very Rare}
When you drink this potion, it removes any Exhaustion levels you have and ends the Poisoned condition on you. For the next 24 hours, you regain the maximum number of Hit Points for any Hit Points Die you spend.

This potion's crimson liquid regularly pulses with dull light, calling to mind a heartbeat.

\DndItemHeader{Potion of Water Breathing}{Uncommon}
You can breathe underwater for 24 hours after drinking this potion.

This potion's cloudy green fluid smells of the sea and has a jellyfish-like bubble floating in it.

\DndItemHeader{Prosthetic Limb}{Wondrous item, Common}
This magic item replaces a lost limb—a hand, an arm, a foot, a leg, or a similar body part. While the prosthetic is attached, it functions identically to the part it replaces. You can detach or reattach it as a Magic action, and it can't be removed against your will while you are alive.

\DndItemHeader{Purple Worm Poison}{}
A creature subjected to Purple Worm Poison makes a DC 21 Constitution saving throw, taking 35 (10d6) Poison damage on a failed save or half as much damage on a successful one.

\DndItemHeader{Quaal's Feather Token}{Wondrous item, Varies}
This object looks like a feather. Different types of feather tokens exist, each with a different single-use effect. The DM chooses the kind of token or determines it randomly by rolling on the Quaal's Feather Tokens table. The type of token determines its rarity.

\begin{DndTable}[header=Quaal's Feather Tokens]{ccc}

1d100 & Token & Rarity\\
01-20 & Quaal's Feather Token, Anchor & Uncommon\\
21-35 & Quaal's Feather Token, Bird & Rare\\
36-50 & Quaal's Feather Token, Fan & Uncommon\\
51-65 & Quaal's Feather Token, Swan boat & Rare\\
66-90 & Quaal's Feather Token, Tree & Uncommon\\
91-00 & Quaal's Feather Token, Whip & Rare\\
\end{DndTable}

\DndItemHeader{Quaal's Feather Token, Anchor}{Wondrous item, Uncommon}
This object looks like a feather.

You can take a Magic action to touch the token to a boat or ship. For the next 24 hours, the vessel can't be moved by any means. Touching the token to the vessel again ends the effect. When the effect ends, the token disappears.

\DndItemHeader{Quaal's Feather Token, Bird}{Wondrous item, Rare}
This object looks like a feather.

You can take a Magic action to toss the token 5 feet into the air. The token disappears and an enormous, multicolored bird takes its place. The bird has the statistics of a \textbf{Roc}, but it can't attack. It obeys your simple commands and can carry up to 500 pounds while flying at its maximum speed (16 miles per hour for a maximum of 144 miles per day, with a 1-hour rest for every 3 hours of flying) or 1,000 pounds at half that speed. The bird disappears after flying its maximum distance for a day or if it drops to 0 Hit Points. You can dismiss the bird as a Magic action.

\DndItemHeader{Quaal's Feather Token, Fan}{Wondrous item, Uncommon}
This object looks like a feather.

If you are on a boat or ship, you can take a Magic action to toss the token up to 10 feet in the air. The token disappears, and a giant flapping fan takes its place. The fan floats and creates a strong wind. This wind can fill the sails of one ship, increasing its speed by 5 miles per hour for 8 hours. You can dismiss the fan as a Magic action.

\DndItemHeader{Quaal's Feather Token, Swan Boat}{Wondrous item, Rare}
This object looks like a feather.

You can take a Magic action to touch the token to a body of water at least 60 feet in diameter. The token disappears, and a 50-foot-long, 20-foot-wide boat shaped like a swan takes its place. The boat is self-propelled and moves across water at a speed of 6 miles per hour. You can take a Magic action while on the boat to command it to move or to turn up to 90 degrees. The boat remains for 24 hours and then disappears. You can dismiss the boat as a Magic action.

\DndItemHeader{Quaal's Feather Token, Tree}{Wondrous item, Uncommon}
This object looks like a feather.

You must be outdoors to use this token. You can take a Magic action to touch it to an unoccupied space on the ground. The token disappears, and in its place a nonmagical oak tree springs into existence. The tree is 60 feet tall and has a 5-foot-diameter trunk, and its branches at the top spread out in a 20-foot radius.

\DndItemHeader{Quaal's Feather Token, Whip}{Wondrous item, Rare}
This object looks like a feather.

You can take a Magic action to throw the token to a point within 10 feet of yourself. The token disappears, and a floating whip takes its place. You can then take a Bonus Action to make a melee spell attack against a creature within 10 feet of the whip, with an attack bonus of +9. On a hit, the target takes 1d6 + 5 Force damage.

As a Bonus Action, you can direct the whip to fly up to 20 feet and repeat the attack against a creature within 10 feet of the whip. The whip disappears after 1 hour, when you take a Magic action to dismiss it, or when you die or have the Incapacitated condition.

\DndItemHeader{Quarterstaff of the Acrobat}{Very Rare (requires attunement)}
You have a +2 bonus to attack rolls and damage rolls made with this magic weapon.

While holding this weapon, you can cause it to emit green Dim Light out to 10 feet, either as a Bonus Action or after you roll Initiative, or you can extinguish the light as a Bonus Action.

While holding this weapon, you can take a Bonus Action to alter its form, turning it into a 6-inch rod (for ease of storage) or a 10-foot pole, or reverting it a Quarterstaff; the weapon will elongate only as far as the surrounding space allows.

In certain forms, the weapon has the following additional properties.

\subparagraph{Acrobatic Assist (Quarterstaff and 10-Foot Pole Forms Only)}
While holding this weapon, you have Advantage on Dexterity (Acrobatics) checks.

\subparagraph{Attack Deflection (Quarterstaff Form Only)}
When you are hit by an attack while holding the weapon, you can take a Reaction to twirl the weapon around you, gaining a +5 bonus to your Armor Class against the triggering attack, potentially causing the attack to miss you. You can't use this property again until you finish a Short Rest or Long Rest.

\subparagraph{Ranged Weapon (Quarterstaff Form Only)}
This weapon has T with a normal range of 30 feet and a long range of 120 feet. Immediately after you make a ranged attack with the weapon, it flies back to your hand.

\DndItemHeader{Quartz}{}
A white, smoky gray, or yellow gemstone.

\DndItemHeader{Quiver of Ehlonna}{Wondrous item, Uncommon}
Each of the quiver's three compartments connects to an extradimensional space that allows the quiver to hold numerous items while never weighing more than 2 pounds. The shortest compartment can hold up to 60 Arrow, Bolt, or similar objects. The midsize compartment holds up to 18 Javelin or similar objects. The longest compartment holds up to 6 long objects, such as bows, Quarterstaff, or Spear.

You can draw any item the quiver contains as if doing so from a regular quiver or scabbard.

\DndItemHeader{Rations}{}
Rations consist of travel-ready food, including jerky, dried fruit, hardtack, and nuts. See "Malnutrition" for the risks of not eating.

\DndItemHeader{Rhodochrosite}{}
A light pink gemstone.

\DndItemHeader{Ring of Animal Influence}{Rare}
This ring has 3 charges, and it regains 1d3 expended charges daily at dawn. While wearing the ring, you can expend 1 charge to cast one of the following spells (save DC 13) from it:


\begin{itemize}
\item \textit{Animal Friendship}
\item \textit{Fear} (affects Beasts only)
\item \textit{Speak with Animals}
\end{itemize}

\DndItemHeader{Ring of Djinni Summoning}{Legendary (requires attunement)}
While wearing this ring, you can take a Magic action to summon a particular \textbf{Djinni} from the Elemental Plane of Air. The djinni appears in an unoccupied space you choose within 120 feet of yourself. It remains as long as you maintain Concentration, to a maximum of 1 hour, or until it drops to 0 Hit Points.

While summoned, the djinni is Friendly [Attitude] to you and your allies, and it obeys your commands. If you fail to command it, the djinni defends itself against attackers but takes no other actions.

After the djinni departs, it can't be summoned again for 24 hours, and the ring becomes nonmagical if the djinni dies.

Rings of Djinni Summoning are often created by the djinn they summon and given to mortals as gifts of friendship or tokens of esteem.

\DndItemHeader{Ring of Elemental Command (Air)}{Legendary (requires attunement)}
The Ring of Elemental Command (air) is linked to the Elemental Plane of Air. Every Ring of Elemental Command has the following two properties:


\begin{itemize}
\item{Elemental Bane.}
While wearing the ring, you have Advantage on attack rolls against Elementals and they have Disadvantage on attack rolls against you.

\item{Elemental Compulsion.}
While wearing the ring, you can take a Magic action to try to compel an Elemental you see within 60 feet of yourself. The Elemental makes a DC 18 Wisdom saving throw. On a failed save, the Elemental has the Charmed condition until the start your next turn, and you determine what it does with its move and action on its next turn.

\end{itemize}

\subparagraph{Elemental Focus}
While wearing the ring, you benefit from additional properties corresponding to the ring's linked Elemental Plane:


\begin{description}
\item[Air.]
You know Auran, you have Resistance to Lightning damage, and you have a Fly Speed equal to your Speed and can hover.

\end{description}

\subparagraph{Spellcasting}
The ring has 5 charges and regains 1d4 + 1 expended charges daily at dawn. While wearing the ring, you can cast a spell from it. A spell cast from the ring has a save DC of 18. Choose the spell from the following list: \textit{Chain Lightning} (3 charges), \textit{Feather Fall} (0 charges), \textit{Gust of Wind} (2 charges), \textit{Wind Wall} (1 charge)

\DndItemHeader{Ring of Elemental Command (Earth)}{Legendary (requires attunement)}
The Ring of Elemental Command (earth) is linked to the Elemental Plane of Earth. Every Ring of Elemental Command has the following two properties:


\begin{itemize}
\item{Elemental Bane.}
While wearing the ring, you have Advantage on attack rolls against Elementals and they have Disadvantage on attack rolls against you.

\item{Elemental Compulsion.}
While wearing the ring, you can take a Magic action to try to compel an Elemental you see within 60 feet of yourself. The Elemental makes a DC 18 Wisdom saving throw. On a failed save, the Elemental has the Charmed condition until the start your next turn, and you determine what it does with its move and action on its next turn.

\end{itemize}

\subparagraph{Elemental Focus}
While wearing the ring, you benefit from additional properties corresponding to the ring's linked Elemental Plane:


\begin{description}
\item[Earth.]
You know Terran, and you have Resistance to Acid damage. Terrain composed of rubble, rocks, or dirt isn't Difficult Terrain for you. In addition, you can move through solid earth or rock as if those areas were Difficult Terrain without disturbing the matter through which you pass. If you end your turn in solid earth or rock, you are shunted out to the nearest unoccupied space you last occupied.

\end{description}

\subparagraph{Spellcasting}
The ring has 5 charges and regains 1d4 + 1 expended charges daily at dawn. While wearing the ring, you can cast a spell from it. A spell cast from the ring has a save DC of 18. Choose the spell from the following list: \textit{Earthquake} (5 charges), \textit{Stone Shape} (2 charges), \textit{Stoneskin} (3 charges), \textit{Wall of Stone} (3 charges)

\DndItemHeader{Ring of Elemental Command (Fire)}{Legendary (requires attunement)}
The Ring of Elemental Command (fire) is linked to the Elemental Plane of Fire. Every Ring of Elemental Command has the following two properties:


\begin{itemize}
\item{Elemental Bane.}
While wearing the ring, you have Advantage on attack rolls against Elementals and they have Disadvantage on attack rolls against you.

\item{Elemental Compulsion.}
While wearing the ring, you can take a Magic action to try to compel an Elemental you see within 60 feet of yourself. The Elemental makes a DC 18 Wisdom saving throw. On a failed save, the Elemental has the Charmed condition until the start your next turn, and you determine what it does with its move and action on its next turn.

\end{itemize}

\subparagraph{Elemental Focus}
While wearing the ring, you benefit from additional properties corresponding to the ring's linked Elemental Plane:


\begin{description}
\item[Fire.]
You know Ignan, and you have Immunity to Fire damage.

\end{description}

\subparagraph{Spellcasting}
The ring has 5 charges and regains 1d4 + 1 expended charges daily at dawn. While wearing the ring, you can cast a spell from it. A spell cast from the ring has a save DC of 18. Choose the spell from the following list: \textit{Burning Hands} (1 charge), \textit{Fireball} (2 charges), \textit{Fire Storm} (4 charges), \textit{Wall of Fire} (3 charges)

\DndItemHeader{Ring of Elemental Command (Water)}{Legendary (requires attunement)}
The Ring of Elemental Command (water) is linked to the Elemental Plane of Water. Every Ring of Elemental Command has the following two properties:


\begin{itemize}
\item{Elemental Bane.}
While wearing the ring, you have Advantage on attack rolls against Elementals and they have Disadvantage on attack rolls against you.

\item{Elemental Compulsion.}
While wearing the ring, you can take a Magic action to try to compel an Elemental you see within 60 feet of yourself. The Elemental makes a DC 18 Wisdom saving throw. On a failed save, the Elemental has the Charmed condition until the start your next turn, and you determine what it does with its move and action on its next turn.

\end{itemize}

\subparagraph{Elemental Focus}
While wearing the ring, you benefit from additional properties corresponding to the ring's linked Elemental Plane:


\begin{description}
\item[Water.]
You know Aquan, you gain a Swim Speed of 60 feet, and you can breathe underwater.

\end{description}

\subparagraph{Spellcasting}
The ring has 5 charges and regains 1d4 + 1 expended charges daily at dawn. While wearing the ring, you can cast a spell from it. A spell cast from the ring has a save DC of 18. Choose the spell from the following list: \textit{Create or Destroy Water} (1 charge), \textit{Ice Storm} (2 charges), \textit{Tsunami} (5 charges), \textit{Wall of Ice} (3 charges), \textit{Water Walk} (2 charges)

\DndItemHeader{Ring of Evasion}{Rare (requires attunement)}
This ring has 3 charges, and it regains 1d3 expended charges daily at dawn. When you fail a Dexterity saving throw while wearing the ring, you can take a Reaction to expend 1 charge to succeed on that save instead.

\DndItemHeader{Ring of Feather Falling}{Rare (requires attunement)}
When you fall while wearing this ring, you descend 60 feet per round and take no damage from falling.

\DndItemHeader{Ring of Free Action}{Rare (requires attunement)}
While you wear this ring, Difficult Terrain doesn't cost you extra movement. In addition, magic can neither reduce any of your Speeds nor cause you to have the Paralyzed or Restrained condition.

\DndItemHeader{Ring of Invisibility}{Legendary (requires attunement)}
While wearing this ring, you can take a Magic action to give yourself the Invisible condition. You remain Invisible until the ring is removed or until you take a Bonus Action to become visible again.

\DndItemHeader{Ring of Jumping}{Uncommon (requires attunement)}
While wearing this ring, you can cast \textit{Jump} from it, but can target only yourself when you do so.

\DndItemHeader{Ring of Mind Shielding}{Uncommon (requires attunement)}
While wearing this ring, you are immune to magic that allows other creatures to read your thoughts, determine whether you are lying, know your alignment, or know your creature type. Creatures can telepathically communicate with you only if you allow it.

You can take a Magic action to cause the ring to become imperceptible until you take another Magic action to make it perceptible, until you remove the ring, or until you die.

If you die while wearing the ring, your soul enters it, unless it already houses a soul. You can remain in the ring or depart for the afterlife. As long as your soul is in the ring, you can telepathically communicate with any creature wearing it. A wearer can't prevent this telepathic communication.

\DndItemHeader{Ring of Protection}{Rare (requires attunement)}
You gain a +1 bonus to Armor Class and saving throws while wearing this ring.

\DndItemHeader{Ring of Regeneration}{Very Rare (requires attunement)}
While wearing this ring, you regain 1d6 Hit Points every 10 minutes if you have at least 1 Hit Points. If you lose a body part, the ring causes the missing part to regrow and return to full functionality after 1d6 + 1 days if you have at least 1 Hit Points the whole time.

\DndItemHeader{Ring of Resistance}{Rare (requires attunement)}
You have Resistance to {{getFullImmRes item.resist}} damage while wearing this ring. The ring is set with {{item.detail1}}.

\DndItemHeader{Ring of Shooting Stars}{Very Rare (requires attunement)}
You can cast \textit{Dancing Lights} or \textit{Light} from the ring.

The ring has 6 charges and regains 1d6 expended charges daily at dawn. You can expend its charges to use the properties below.

\subparagraph{Faerie Fire}
You can expend 1 charge to cast \textit{Faerie Fire} from the ring.

\subparagraph{Lightning Spheres}
You can expend 2 charges as a Magic action to create up to four 3-foot-diameter spheres of lightning.

Each sphere appears in an unoccupied space you can see within 120 feet of yourself. The spheres last as long as you maintain Concentration, up to 1 minute. Each sphere sheds Dim Light in a 30-foot radius.

As a Bonus Action, you can move each sphere up to 30 feet, but no farther than 120 feet away from yourself. The first time the sphere comes within 5 feet of a creature other than you that isn't behind Cover, the sphere discharges lightning at that creature and disappears. That creature makes a DC 15 Dexterity saving throw. On a failed save, the creature takes Lightning damage based on the number of spheres you created, as shown in the following table. On a successful save, the creature takes half as much damage.


\begin{DndTable}{cX}

Number of Spheres & Lightning Damage\\
1 & 4d12\\
2 & 5d4\\
3 & 2d6\\
4 & 2d4\\
\end{DndTable}

\subparagraph{Shooting Stars}
You can expend 1 to 3 charges as a Magic action. For every charge you expend, you launch a glowing mote of light from the ring at a point you can see within 60 feet of yourself. Each creature in a 15-foot Cube [Area of Effect] originating from that point is showered in sparks and makes a DC 15 Dexterity saving throw, taking 5d4 Radiant damage on a failed save or half as much damage on a successful one.

\DndItemHeader{Ring of Spell Storing}{Rare (requires attunement)}
This ring stores spells cast into it, holding them until the attuned wearer uses them. The ring can store up to 5 levels worth of spells at a time. When found, it contains 1d6 - 1 levels of stored spells chosen by the DM.

Any creature can cast a spell of level 1 through 5 into the ring by touching the ring as the spell is cast. The spell has no effect other than to be stored in the ring. If the ring can't hold the spell, the spell is expended without effect. The level of the slot used to cast the spell determines how much space it uses.

While wearing this ring, you can cast any spell stored in it. The spell uses the slot level, spell save DC, spell attack bonus, and spellcasting ability of the original caster but is otherwise treated as if you cast the spell. The spell cast from the ring is no longer stored in it, freeing up space.

\DndItemHeader{Ring of Spell Turning}{Legendary (requires attunement)}
While wearing this ring, you have Advantage on saving throws against spells. If you succeed on the save for a spell of level 7 or lower, the spell has no effect on you. If that spell targeted only you and didn't create an area of effect, you can take a Reaction to deflect the spell back at the spell's caster; the caster must make a saving throw against the spell using their own spell save DC.

\DndItemHeader{Ring of Swimming}{Uncommon}
You have a Swim Speed of 40 feet while wearing this ring.

\DndItemHeader{Ring of Telekinesis}{Very Rare (requires attunement)}
While wearing this ring, you can cast \textit{Telekinesis} from it.

\DndItemHeader{Ring of Three Wishes}{Legendary}
While wearing this ring, you can expend 1 of its 3 charges to cast \textit{Wish} from it. The ring becomes nonmagical when you use the last charge.

\DndItemHeader{Ring of Warmth}{Uncommon (requires attunement)}
If you take Cold damage while wearing this ring, the ring reduces the damage you take by 2d8.

In addition, while wearing this ring, you and everything you wear and carry are unharmed by temperatures of 0 degrees Fahrenheit or lower.

\DndItemHeader{Ring of Water Walking}{Uncommon}
While wearing this ring, you cast \textit{Water Walk} from it, targeting only yourself.

\DndItemHeader{Ring of X-ray Vision}{Rare (requires attunement)}
While wearing this ring, you can take a Magic action to gain X-ray vision with a range of 30 feet for 1 minute. To you, solid objects within that radius appear transparent and don't prevent light from passing through them. The vision can penetrate 1 foot of stone, 1 inch of common metal, or up to 3 feet of wood or dirt. Thicker substances or a thin sheet of lead block the vision.

Whenever you use the ring again before taking a Long Rest, you must succeed on a DC 15 Constitution saving throw or gain 1 Exhaustion level.

\DndItemHeader{Ring of X-ray Vision}{Rare (requires attunement)}
While wearing this ring, you can take a Magic action to gain X-ray vision with a range of 30 feet for 1 minute. To you, solid objects within that radius appear transparent and don't prevent light from passing through them. The vision can penetrate 1 foot of stone, 1 inch of common metal, or up to 3 feet of wood or dirt. Thicker substances or a thin sheet of lead block the vision.

Whenever you use the ring again before taking a Long Rest, you must succeed on a DC 15 Constitution saving throw or gain 1 Exhaustion level.

\DndItemHeader{Ring of the Ram}{Rare (requires attunement)}
This ring has 3 charges and regains 1d3 expended charges daily at dawn. While wearing the ring, you can take a Magic action to expend 1 to 3 charges to make a ranged spell attack against one creature you can see within 60 feet of yourself. The ring produces a spectral ram's head and makes its attack roll with a +7 bonus. On a hit, for each charge you spend, the target takes 2d10 Force damage and is pushed 5 feet away from you.

Alternatively, you can expend 1 to 3 of the ring's charges as a Magic action to try to break a nonmagical object you can see within 60 feet of yourself that isn't being worn or carried. The ring makes a Strength check with a +5 bonus for each charge you spend.

\DndItemHeader{Rival Coin}{Wondrous item, Common}
This gold coin has a creature embossed on each side. The two depicted creatures must be famous rivals or enemies of each other. For example, a Rival Coin might show \textit{Iggwilv} on one side and \textit{Mordenkainen} on the other, or \textit{Venger} on one side and \textit{Tiamat} on the other. One of these figures is on the "heads" side of the coin, the other on the "tails" side.

The coin has 1 charge and regains its expended charge daily at dawn. You can take a Magic action to toss the coin, expending its charge. Roll any die to determine whether the coin comes up heads (on an even number) or tails (on an odd number). The roll also determines the effect:


\begin{description}
\item[Heads.]
Target one creature you can see within 60 feet of yourself. The target makes a DC 13 Wisdom saving throw. On a failed save, the target takes 2d4 Psychic damage and has Disadvantage on the next attack roll it makes before the end of its next turn. On a successful save, the target takes half as much damage only.

\item[Tails.]
You take 1d4 Psychic damage.

\end{description}

\DndItemHeader{Robe of Eyes}{Wondrous item, Rare (requires attunement)}
This robe is adorned with eyelike patterns. While you wear the robe, you gain the following benefits:


\begin{description}
\item[All-Around Vision.]
The robe gives you Advantage on Wisdom (Perception) checks that rely on sight.

\item[Special Senses.]
You have Darkvision and Truesight, both with a range of 120 feet.

\end{description}

\subparagraph{Drawbacks}
A \textit{Light} spell cast on the robe or a \textit{Daylight} spell cast within 5 feet of the robe gives you the Blinded condition for 1 minute. At the end of each of your turns, you make a Constitution saving throw (DC 11 for Light or DC 15 for Daylight), ending the condition on yourself on a success.

\DndItemHeader{Robe of Scintillating Colors}{Wondrous item, Very Rare (requires attunement)}
This robe has 3 charges, and it regains 1d3 expended charges daily at dawn. While you wear it, you can take a Magic action and expend 1 charge to cause the garment to display a shifting pattern of dazzling hues until the end of your next turn. During this time, the robe sheds Bright Light in a 30-foot radius and Dim Light for an additional 30 feet, and creatures that can see you have Disadvantage on attack rolls against you. Any creature in the Bright Light that can see you when the robe's power is activated must succeed on a DC 15 Wisdom saving throw or have the Stunned condition until the effect ends.

\DndItemHeader{Robe of Stars}{Wondrous item, Very Rare (requires attunement)}
This black or dark-blue robe is embroidered with small white or silver stars. You gain a +1 bonus to saving throws while you wear it.

Six stars, located on the robe's upper-front portion, are particularly large. While wearing this robe, you can take a Magic action to remove one of the stars and expend it to cast the level 5 version of \textit{Magic Missile}. Daily at dusk, 1d6 removed stars reappear on the robe.

While you wear the robe, you can take a Magic action to enter the Astral Plane along with everything you are wearing and carrying. You remain there until you take a Magic action to return to the plane you were on. You reappear in the last space you occupied or, if that space is occupied, the nearest unoccupied space.

\DndItemHeader{Robe of Useful Items}{Wondrous item, Uncommon}
This robe has cloth patches of various shapes and colors covering it. While wearing the robe, you can take a Magic action to detach one of the patches, causing it to become the object or creature it represents. Once the last patch is removed, the robe becomes an ordinary garment.

The robe has two of each of the following patches:


\begin{itemize}
\item \textit{Bullseye Lantern} (filled and lit)
\item \textit{Dagger}
\item \textit{Mirror}
\item \textit{Pole}
\item \textit{Rope} (coiled)
\item \textit{Sack}
\end{itemize}

In addition, the robe has 4d4 other patches. The DM chooses the patches or determines them randomly by rolling on the following table.


\begin{DndTable}{cX}

1d100 & Patch\\
01-08 & Bag of 100 GP\\
09-15 & Silver coffer (1 foot long, 6 inches wide and deep) worth 500 GP\\
16-22 & Iron door (up to 10 feet wide and 10 feet high, barred on one side of your choice), which you can place in an opening you can reach; it conforms to fit the opening, attaching and hinging itself\\
23-30 & 10 gems worth 100 GP each\\
31-44 & Wooden ladder (24 feet long)\\
45-51 & \textbf{Riding Horse} with a \textit{Riding Saddle}\\
52-59 & Open pit (a 10-foot Cube [Area of Effect]), which you can place on the ground within 10 feet of yourself\\
60-68 & 4 Potion of Healing\\
69-75 & \textit{Rowboat} (12 feet long)\\
76-83 & \textit{Spell Scroll} containing one spell of level 1, 2, or 3 (your choice)\\
84-90 & 2 Mastiffs\\
91-96 & Window (2 feet by 4 feet, up to 2 feet deep), which you can place on a vertical surface you can reach\\
97-00 & \textit{Portable Ram}\\
\end{DndTable}

\DndItemHeader{Robe of the Archmagi}{Wondrous item, Legendary (requires attunement by a sorcerer, warlock, or wizard)}
This elegant garment is made from exquisite cloth and adorned with runes.

You gain these benefits while wearing the robe.

\subparagraph{Armor}
If you aren't wearing armor, your base Armor Class is 15 plus your Dexterity modifier.

\subparagraph{Magic Resistance}
You have Advantage on saving throws against spells and other magical effects.

\subparagraph{War Mage}
Your spell save DC and spell attack bonus each increase by 2.

\DndItemHeader{Rod of Absorption}{Very Rare (requires attunement)}
While holding this rod, you can take a Reaction to absorb a spell that is targeting only you and doesn't create an area of effect. The absorbed spell's effect is canceled, and the spell's energy—not the spell itself—is stored in the rod. The energy has the same level as the spell when it was cast. A canceled spell dissipates with no effect, and any resources used to cast it are wasted. The rod can absorb and store up to 50 levels of energy over the course of its existence. Once the rod absorbs 50 levels of energy, it can't absorb more. If you are targeted by a spell that the rod can't store, the rod has no effect on that spell.

When you become attuned to the rod, you know how many levels of energy the rod has absorbed over the course of its existence and how many levels of spell energy it currently has stored.

If you are a spellcaster holding the rod, you can convert energy stored in it into spell slots to cast spells you have prepared or know. You can create spell slots only of a level equal to or lower than your own spell slots, up to a maximum of level 5. You use the stored levels in place of your slots but otherwise cast the spell as normal. For example, you can use 3 levels stored in the rod as a level 3 spell slot.

A newly found rod typically has 1d10 levels of spell energy stored in it. A rod that can no longer absorb spell energy and has no energy remaining becomes nonmagical.

\DndItemHeader{Rod of Alertness}{Very Rare (requires attunement)}
This rod has the following properties.

\subparagraph{Alertness}
While holding the rod, you have Advantage on Wisdom (Perception) checks and on Initiative rolls. Spells. While holding the rod, you can cast the following spells from it:


\begin{itemize}
\item \textit{Detect Evil and Good}
\item \textit{Detect Magic}
\item \textit{Detect Poison and Disease}
\item \textit{See Invisibility}
\end{itemize}

\subparagraph{Protective Aura}
As a Magic action, you can plant the haft end of the rod in the ground, whereupon the rod's head sheds Bright Light in a 60-foot radius and Dim Light for an additional 60 feet. While in that Bright Light, you and your allies gain a +1 bonus to Armor Class and saving throws and can sense the location of any Invisible creature that is also in the Bright Light.

The rod's head stops glowing and the effect ends after 10 minutes or when a creature takes a Magic action to pull the rod from the ground. Once used, this property can't be used again until the next dawn.

\DndItemHeader{Rod of Lordly Might}{Legendary (requires attunement)}
This rod has a flanged head, and it functions as a magic Mace that grants a +3 bonus to attack rolls and damage rolls made with it. The rod has properties associated with six different buttons that are set in a row along the haft. It has three other properties as well, detailed below.

\subparagraph{Buttons}
You can press one of the following buttons as a Bonus Action; a button's effect lasts until you push a different button or until you push the same button again, which causes the rod to revert to its normal form:


\begin{description}
\item[Button 1.]
A fiery blade sprouts from the end opposite the rod's flanged head. The flames shed Bright Light in a 40-foot radius and Dim Light for an additional 40 feet, and the blade functions as a magic Longsword or Shortsword (your choice) that deals an extra 2d6 Fire damage on a hit.

\item[Button 2.]
The rod's flanged head folds down and two crescent-shaped blades spring out, transforming the rod into a magic Battleaxe that grants a +3 bonus to attack rolls and damage rolls made with it.

\item[Button 3.]
The rod's flanged head folds down, a spear point springs from the rod's tip, and the rod's handle lengthens into a 6-foot haft, transforming the rod into a magic Spear that grants a +3 bonus to attack rolls and damage rolls made with it.

\item[Button 4.]
The rod transforms into a climbing pole up to 50 feet long (you specify the length), though the rod's buttons remain within your reach. In surfaces as hard as granite, a spike at the bottom and three hooks at the top anchor the pole. Horizontal bars 3 inches long fold out from the sides, 1 foot apart, forming a ladder. The pole can bear up to 4,000 pounds. More weight or lack of solid anchoring causes the rod to revert to its normal form.

\item[Button 5.]
The rod transforms into a handheld battering ram and grants its user a +10 bonus to Strength (Athletics) checks made to break through doors, barricades, and other barriers.

\item[Button 6.]
The rod assumes or remains in its normal form and indicates magnetic north. (Nothing happens if this function of the rod is used in a location that has no magnetic north.) The rod also gives you knowledge of your approximate depth beneath the ground or your height above it.

\end{description}

\subparagraph{Drain Life}
When you hit a creature with a melee attack using the rod, you can force the target to make a DC 17 Constitution saving throw. On a failed save, the target takes an extra 4d6 Necrotic damage, and you regain a number of Hit Points equal to half that Necrotic damage. Once used, this property can't be used again until the next dawn.

\subparagraph{Paralyze}
When you hit a creature with a melee attack using the rod, you can force the target to make a DC 17 Constitution saving throw. On a failed save, the target has the Paralyzed condition for 1 minute. The target repeats the save at the end of each of its turns, ending the effect on a success. Once used, this property can't be used again until the next dawn.

\subparagraph{Terrify}
While holding the rod, you can take a Magic action to force each creature you can see within 30 feet of yourself to make a DC 17 Wisdom saving throw. On a failed save, a target has the Frightened condition for 1 minute. A Frightened target repeats the save at the end of each of its turns, ending the effect on itself on a success. Once used, this property can't be used again until the next dawn.

\DndItemHeader{Rod of Resurrection}{Legendary (requires attunement by a cleric, druid, or paladin)}
The rod has 5 charges. While you hold it, you can cast one of the following spells from it: \textit{Heal} (expends 1 charge) or \textit{Resurrection} (expends 5 charges).

The rod regains 1 expended charge daily at dawn. If you expend the last charge, roll 1d20. On a 1, the rod disappears in a harmless burst of radiance.

\DndItemHeader{Rod of Rulership}{Rare (requires attunement)}
You can take a Magic action to present the rod and command obedience from each creature of your choice that you can see within 120 feet of yourself. Each target must succeed on a DC 15 Wisdom saving throw or have the Charmed condition for 8 hours. While Charmed in this way, the creature regards you as its trusted leader. If harmed by you or your allies or commanded to do something contrary to its nature, a target ceases to be Charmed in this way. Once used, this property can't be used again until the next dawn.

\DndItemHeader{Rod of Security}{Very Rare}
While holding this rod, you can take a Magic action to activate it. The rod then instantly transports you and up to 199 other willing creatures you can see to a demiplane. You choose the form the demiplane takes. It could be a tranquil garden, a cheery tavern, an immense palace, a tropical island, a fantastic carnival, or whatever else you can imagine. Regardless of its nature, the demiplane contains enough water and food to sustain its visitors, and the demiplane's environment can't harm its occupants. Everything else that can be interacted with there can exist only there. For example, a flower picked from a garden there disappears if it is taken outside the demiplane.

For each hour spent in the demiplane, a visitor regains Hit Points as if it had spent 1 Hit Points Die. Also, creatures don't age while there, although time passes normally. Visitors can remain there for up to 200 days divided by the number of creatures present (round down).

When the time runs out or you take a Magic action to end the effect, all visitors reappear in the location they occupied when you activated the rod or an unoccupied space nearest that location. Once used, this property can't be used again until 10 days have passed.

\DndItemHeader{Rope of Climbing}{Wondrous item, Uncommon}
This 60-foot length of rope can hold up to 3,000 pounds. While holding one end of the rope, you can take a Magic action to command the other end of the rope to animate and move toward a destination you choose, up to the rope's length away from you. That end moves 10 feet on your turn when you first command it and 10 feet at the start of each of your subsequent turns until reaching its destination or until you tell it to stop. You can also tell the rope to fasten itself securely to an object or to unfasten itself, to knot or unknot itself, or to coil itself for carrying.

If you tell the rope to knot, large knots appear at 1-foot intervals along the rope. While knotted, the rope shortens to a 50-foot length and grants Advantage on ability checks made to climb using the rope.

The rope has AC 20, HP 20, and Immunity to Poison and Psychic damage. It regains 1 Hit Points every 5 minutes as long as it has at least 1 Hit Points. If the rope drops to 0 Hit Points, it is destroyed.

\DndItemHeader{Rope of Entanglement}{Wondrous item, Rare}
This rope is 30 feet long. While holding one end of the rope, you can take a Magic action to command the other end to dart forward and entangle one creature you can see within 20 feet of yourself. The target must succeed on a DC 15 Dexterity saving throw or have the Restrained condition. You can release the target by letting go of your end of the rope (causing the rope to coil up in the target's space) or by using a Bonus Action to repeat the command (causing the rope to coil up in your hand).

A target Restrained by the rope can take an action to make its choice of a DC 15 Strength (Athletics) or Dexterity (Acrobatics) check. On a successful check, the target is no longer Restrained by the rope. If you're still holding onto the rope when a target escapes from it, you can take a Reaction to command the rope to coil up in your hand; otherwise, the rope coils up in the target's space.

The rope has AC 20, HP 20, and Immunity to Poison and Psychic damage. It regains 1 Hit Points every 5 minutes as long as it has at least 1 Hit Points. If the rope drops to 0 Hit Points, it is destroyed.

\DndItemHeader{Rope of Mending}{Wondrous item, Common}
This 50-foot coil of rope can repair itself when cut into any number of smaller pieces. As a Magic action, you can cause all pieces of the rope that are in contact with each other and not otherwise in use to knit back together. A Rope of Mending is forever shortened if a section of it is lost or destroyed.

\DndItemHeader{Ruby}{}
A clear red to deep crimson gemstone.

\DndItemHeader{Ruby of the War Mage}{Wondrous item, Common (requires attunement by a spellcaster)}
Etched with eldritch runes, this 1-inch-diameter ruby allows you to use a Simple or Martial weapon as a Spellcasting Focus for your spells. For this property to work, you must attach the ruby to the weapon by pressing the ruby against it for at least 10 minutes. Thereafter, the ruby can't be removed unless you detach it as a Magic action, the weapon is destroyed, or your Attunement to the ruby ends.

\DndItemHeader{Saddle of the Cavalier}{Wondrous item, Uncommon}
While in this saddle on a mount, you can't be dismounted against your will if you're conscious, and attack rolls against the mount have disadvantage.

\DndItemHeader{Sardonyx}{}
A bands of red and white gemstone.

\DndItemHeader{Scarab of Protection}{Wondrous item, Legendary (requires attunement)}
This beetle-shaped medallion provides three benefits while it is on your person.

\subparagraph{Defense}
You gain a +1 bonus to Armor Class.

\subparagraph{Preservation}
The scarab has 12 charges. If you fail a saving throw against a Necromancy spell or a harmful effect originating from an Undead, you can take a Reaction to expend 1 charge and turn the failed save into a successful one. The scarab crumbles into powder and is destroyed when its last charge is expended.

\subparagraph{Spell Resistance}
You have Advantage on saving throws against spells.

\DndItemHeader{Scimitar of Speed}{Very Rare (requires attunement)}
You gain a +2 bonus to attack rolls and damage rolls made with this magic weapon. In addition, you can make one attack with it as a Bonus Action on each of your turns.

\DndItemHeader{Scroll of Protection}{Rare}
Using a Magic action to read the scroll creates a 5-foot Emanation [Area of Effect] originating from you. For 5 minutes, {{item.detail2}} can't enter or affect anything in the area. However, if you move in such a way that an {{item.detail1}} would be inside the area, the effect ends.

As a Magic action, a creature within 5 feet of the Emanation [Area of Effect] can attempt to overcome it, which forces the creature to make a DC 15 Charisma saving throw. On a successful save, the creature ceases to be affected by the Emanation [Area of Effect].

\DndItemHeader{Scroll of Titan Summoning}{Legendary}
When you take a Magic action to read this scroll, a particular titan named in the scroll appears in an unoccupied space on the ground or in water that you can see within 1 mile of yourself. The DM picks a suitable titan or determines it randomly by rolling on the table below.

The titan is Hostile [Attitude] toward all other creatures and disappears when it drops to 0 Hit Points. If the titan is summoned into a space that isn't large enough to contain it, the summoning fails, and the scroll is wasted.


\begin{DndTable}{cX}

1d100 & Titan\\
01-15 & Scroll of Titan Summoning (Animal Lord)\\
16-30 & Scroll of Titan Summoning (Blob of Annihilation)\\
31-45 & Scroll of Titan Summoning (Colossus)\\
46-60 & Scroll of Titan Summoning (Elemental Cataclysm)\\
61-75 & Scroll of Titan Summoning (Empyrean)\\
76-90 & Scroll of Titan Summoning (Kraken) (a kraken requires a body of water large enough to contain it, or the summoning fails and the scroll is wasted)\\
91-100 & Scroll of Titan Summoning (Tarrasque)\\
\end{DndTable}

\DndItemHeader{Sending Stones}{Wondrous item, Uncommon}
Sending Stones come in pairs, with each stone carved to match the other so the pairing is easily recognized. While you touch one stone, you can cast \textit{Sending} from it. The target is the bearer of the other stone. If no creature bears the other stone, you know that fact as soon as you use the stone, and you don't cast the spell.

Once Sending is cast using either stone, the stones can't be used again until the next dawn. If one of the stones in a pair is destroyed, the other one becomes nonmagical.

\DndItemHeader{Sentinel Shield}{Uncommon}
While holding this Shield, you have Advantage on Initiative rolls and Wisdom (Perception) checks. The Shield is emblazoned with a symbol of an eye.

\DndItemHeader{Serpent Venom}{}
A creature subjected to Serpent Venom must succeed on a DC 11 Constitution saving throw, taking 10 (3d6) Poison damage on a failed save or half as much damage on a successful one.

\DndItemHeader{Shield of Expression}{Common}
The front of this Shield is shaped in the likeness of a face. While bearing the Shield, you can take a Bonus Action to alter the face's expression.

\DndItemHeader{Shield of Missile Attraction}{Rare (requires attunement)}
While holding this Shield, you have Resistance to damage from attacks made with Ranged weapons.

\subparagraph{Curse}
This Shield is cursed. Attuning to it curses you until you are targeted by a \textit{Remove Curse} spell or similar magic. Removing the Shield fails to end the curse on you. Whenever an attack with a Ranged weapon targets a creature within 10 feet of you, the curse causes you to become the target instead.

\DndItemHeader{Shield of the Cavalier}{Very Rare (requires attunement)}
While holding this Shield, you have a +2 bonus to Armor Class. This bonus is in addition to the Shield's normal bonus to AC.

The Shield has the following additional properties that you can use while holding it.

\subparagraph{Forceful Bash}
When you take the Attack, you can make one of the attack rolls using the Shield against a target within 5 feet of yourself. Apply your Proficiency and Strength modifier to the attack roll. On a hit, the Shield deals Force damage to the target equal to 2d6 + 2 plus your Strength modifier, and if the target is a creature, you can push it up to 10 feet directly away from yourself. If the creature is your size or smaller, you can also knock it down, giving it the Prone condition.

\subparagraph{Protective Field}
As a Reaction, when you or an ally you can see within 5 feet of you is targeted by an attack or makes a saving throw against an area of effect, you can use the Shield to create an immobile 5-foot Emanation [Area of Effect] originating from you. When the Emanation [Area of Effect] appears, any creatures or objects not fully contained within it are pushed into the nearest unoccupied spaces outside it. The attack or area of effect that triggered the Reaction has no effect on creatures and objects inside the Emanation [Area of Effect], which lasts as long as you maintain Concentration, up to 1 minute. Nothing can pass into or out of the Emanation [Area of Effect]. A creature or object inside the Emanation [Area of Effect] can't be damaged by attacks or effects originating from outside, nor can a creature inside the Emanation [Area of Effect] damage anything outside it. Once this property is used, it can't be used again until the next dawn.

\DndItemHeader{Slippers of Spider Climbing}{Wondrous item, Uncommon (requires attunement)}
While you wear these light shoes, you can move up, down, and across vertical surfaces and along ceilings, while leaving your hands free. You have a Climb Speed equal to your Speed. However, the slippers don't allow you to move this way on a slippery surface, such as one covered by ice or oil.

\DndItemHeader{Smith's Tools}{}

\begin{description}
\item[Ability:]
Strength

\item[Utilize:]
Pry open a door or container (DC 20)

\item[Craft:]
Any Melee weapon (except \textit{Club}, \textit{Greatclub}, \textit{Quarterstaff}, and \textit{Whip}), Medium armor (except Hide Armor), Heavy armor, \textit{Ball Bearings}, \textit{Bucket}, \textit{Caltrops}, \textit{Chain}, \textit{Crowbar}, Firearm Bullets (10), \textit{Grappling Hook}, \textit{Iron Pot}, \textit{Iron Spikes}, Sling Bullets (20)

\end{description}

\DndItemHeader{Smoke Grenade}{}
As an action, you can either throw a grenade at a point up to 60 feet away or use a Grenade Launcher to propel the grenade to a point up to 1,000 feet away. The grenade explodes at that point, creating a particular effect in a 20-foot-radius Sphere [Area of Effect].

The area of the Sphere [Area of Effect] is Heavily Obscured by smoke for 1 minute. A strong wind (such as the \textit{Gust of Wind} spell) disperses the smoke.

\DndItemHeader{Sovereign Glue}{Wondrous item, Legendary}
This viscous, milky-white substance can form a permanent adhesive bond between any two objects. It must be stored in a jar or flask that has been coated inside with \textit{Oil of Slipperiness}. When found, a container contains 1d6 + 1 ounces.

One ounce of the glue can cover a 1-foot square surface. Applying an ounce of Sovereign Glue takes a Utilize action, and the applied glue takes 1 minute to set. Once it has done so, the bond it creates can be broken only by the application of Universal Solvent or \textit{Oil of Etherealness}, or with a \textit{Wish} spell.

\DndItemHeader{Spell Scroll (Cantrip)}{Common}
A Spell Scroll bears the words of a single spell, written in a mystical cipher. If the spell is on your spell list, you can read the scroll and cast its spell without Material components. Otherwise, the scroll is unintelligible. Casting the spell by reading the scroll requires the spell's normal casting time. Once the spell is cast, the scroll crumbles to dust. If the casting is interrupted, the scroll isn't lost.

If the spell is on your spell list but of a higher level than you can normally cast, you make a DC 10 ability check using your spellcasting ability to determine whether you cast the spell. On a failed check, the spell disappears from the scroll with no other effect.

If the spell requires a saving throw or an attack roll, the spell save DC is 13, and the attack bonus is 5.

\subparagraph{Copying a Scroll into a Spellbook}
A Wizard spell on a Spell Scroll can be copied into a spellbook. When a cantrip is copied in this way, the copier must succeed on a DC 10 Intelligence (Arcana). On a successful check, the spell is copied. Whether the check succeeds or fails, the Spell Scroll is destroyed.

\DndItemHeader{Spell Scroll (Level 1)}{Common}
A Spell Scroll bears the words of a single spell, written in a mystical cipher. If the spell is on your spell list, you can read the scroll and cast its spell without Material components. Otherwise, the scroll is unintelligible. Casting the spell by reading the scroll requires the spell's normal casting time. Once the spell is cast, the scroll crumbles to dust. If the casting is interrupted, the scroll isn't lost.

If the spell is on your spell list but of a higher level than you can normally cast, you make a DC 11 ability check using your spellcasting ability to determine whether you cast the spell. On a failed check, the spell disappears from the scroll with no other effect.

If the spell requires a saving throw or an attack roll, the spell save DC is 13, and the attack bonus is 5.

\subparagraph{Copying a Scroll into a Spellbook}
A Wizard spell on a Spell Scroll can be copied into a spellbook. When a level 1 spell is copied in this way, the copier must succeed on a DC 11 Intelligence (Arcana). On a successful check, the spell is copied. Whether the check succeeds or fails, the Spell Scroll is destroyed.

\DndItemHeader{Spell Scroll (Level 2)}{Uncommon}
A Spell Scroll bears the words of a single spell, written in a mystical cipher. If the spell is on your spell list, you can read the scroll and cast its spell without Material components. Otherwise, the scroll is unintelligible. Casting the spell by reading the scroll requires the spell's normal casting time. Once the spell is cast, the scroll crumbles to dust. If the casting is interrupted, the scroll isn't lost.

If the spell is on your spell list but of a higher level than you can normally cast, you make a DC 12 ability check using your spellcasting ability to determine whether you cast the spell. On a failed check, the spell disappears from the scroll with no other effect.

If the spell requires a saving throw or an attack roll, the spell save DC is 13, and the attack bonus is 5.

\subparagraph{Copying a Scroll into a Spellbook}
A Wizard spell on a Spell Scroll can be copied into a spellbook. When a level 2 spell is copied in this way, the copier must succeed on a DC 12 Intelligence (Arcana). On a successful check, the spell is copied. Whether the check succeeds or fails, the Spell Scroll is destroyed.

\DndItemHeader{Spell Scroll (Level 3)}{Uncommon}
A Spell Scroll bears the words of a single spell, written in a mystical cipher. If the spell is on your spell list, you can read the scroll and cast its spell without Material components. Otherwise, the scroll is unintelligible. Casting the spell by reading the scroll requires the spell's normal casting time. Once the spell is cast, the scroll crumbles to dust. If the casting is interrupted, the scroll isn't lost.

If the spell is on your spell list but of a higher level than you can normally cast, you make a DC 13 ability check using your spellcasting ability to determine whether you cast the spell. On a failed check, the spell disappears from the scroll with no other effect.

If the spell requires a saving throw or an attack roll, the spell save DC is 15, and the attack bonus is 7.

\subparagraph{Copying a Scroll into a Spellbook}
A Wizard spell on a Spell Scroll can be copied into a spellbook. When a level 3 spell is copied in this way, the copier must succeed on a DC 13 Intelligence (Arcana). On a successful check, the spell is copied. Whether the check succeeds or fails, the Spell Scroll is destroyed.

\DndItemHeader{Spell Scroll (Level 4)}{Rare}
A Spell Scroll bears the words of a single spell, written in a mystical cipher. If the spell is on your spell list, you can read the scroll and cast its spell without Material components. Otherwise, the scroll is unintelligible. Casting the spell by reading the scroll requires the spell's normal casting time. Once the spell is cast, the scroll crumbles to dust. If the casting is interrupted, the scroll isn't lost.

If the spell is on your spell list but of a higher level than you can normally cast, you make a DC 14 ability check using your spellcasting ability to determine whether you cast the spell. On a failed check, the spell disappears from the scroll with no other effect.

If the spell requires a saving throw or an attack roll, the spell save DC is 15, and the attack bonus is 7.

\subparagraph{Copying a Scroll into a Spellbook}
A Wizard spell on a Spell Scroll can be copied into a spellbook. When a level 4 spell is copied in this way, the copier must succeed on a DC 14 Intelligence (Arcana). On a successful check, the spell is copied. Whether the check succeeds or fails, the Spell Scroll is destroyed.

\DndItemHeader{Spell Scroll (Level 5)}{Rare}
A Spell Scroll bears the words of a single spell, written in a mystical cipher. If the spell is on your spell list, you can read the scroll and cast its spell without Material components. Otherwise, the scroll is unintelligible. Casting the spell by reading the scroll requires the spell's normal casting time. Once the spell is cast, the scroll crumbles to dust. If the casting is interrupted, the scroll isn't lost.

If the spell is on your spell list but of a higher level than you can normally cast, you make a DC 15 ability check using your spellcasting ability to determine whether you cast the spell. On a failed check, the spell disappears from the scroll with no other effect.

If the spell requires a saving throw or an attack roll, the spell save DC is 17, and the attack bonus is 9.

\subparagraph{Copying a Scroll into a Spellbook}
A Wizard spell on a Spell Scroll can be copied into a spellbook. When a level 5 spell is copied in this way, the copier must succeed on a DC 15 Intelligence (Arcana). On a successful check, the spell is copied. Whether the check succeeds or fails, the Spell Scroll is destroyed.

\DndItemHeader{Spell Scroll (Level 6)}{Very Rare}
A Spell Scroll bears the words of a single spell, written in a mystical cipher. If the spell is on your spell list, you can read the scroll and cast its spell without Material components. Otherwise, the scroll is unintelligible. Casting the spell by reading the scroll requires the spell's normal casting time. Once the spell is cast, the scroll crumbles to dust. If the casting is interrupted, the scroll isn't lost.

If the spell is on your spell list but of a higher level than you can normally cast, you make a DC 16 ability check using your spellcasting ability to determine whether you cast the spell. On a failed check, the spell disappears from the scroll with no other effect.

If the spell requires a saving throw or an attack roll, the spell save DC is 17, and the attack bonus is 9.

\subparagraph{Copying a Scroll into a Spellbook}
A Wizard spell on a Spell Scroll can be copied into a spellbook. When a level 6 spell is copied in this way, the copier must succeed on a DC 16 Intelligence (Arcana). On a successful check, the spell is copied. Whether the check succeeds or fails, the Spell Scroll is destroyed.

\DndItemHeader{Spell Scroll (Level 7)}{Very Rare}
A Spell Scroll bears the words of a single spell, written in a mystical cipher. If the spell is on your spell list, you can read the scroll and cast its spell without Material components. Otherwise, the scroll is unintelligible. Casting the spell by reading the scroll requires the spell's normal casting time. Once the spell is cast, the scroll crumbles to dust. If the casting is interrupted, the scroll isn't lost.

If the spell is on your spell list but of a higher level than you can normally cast, you make a DC 17 ability check using your spellcasting ability to determine whether you cast the spell. On a failed check, the spell disappears from the scroll with no other effect.

If the spell requires a saving throw or an attack roll, the spell save DC is 18, and the attack bonus is 10.

\subparagraph{Copying a Scroll into a Spellbook}
A Wizard spell on a Spell Scroll can be copied into a spellbook. When a level 7 spell is copied in this way, the copier must succeed on a DC 17 Intelligence (Arcana). On a successful check, the spell is copied. Whether the check succeeds or fails, the Spell Scroll is destroyed.

\DndItemHeader{Spell Scroll (Level 8)}{Very Rare}
A Spell Scroll bears the words of a single spell, written in a mystical cipher. If the spell is on your spell list, you can read the scroll and cast its spell without Material components. Otherwise, the scroll is unintelligible. Casting the spell by reading the scroll requires the spell's normal casting time. Once the spell is cast, the scroll crumbles to dust. If the casting is interrupted, the scroll isn't lost.

If the spell is on your spell list but of a higher level than you can normally cast, you make a DC 18 ability check using your spellcasting ability to determine whether you cast the spell. On a failed check, the spell disappears from the scroll with no other effect.

If the spell requires a saving throw or an attack roll, the spell save DC is 18, and the attack bonus is 10.

\subparagraph{Copying a Scroll into a Spellbook}
A Wizard spell on a Spell Scroll can be copied into a spellbook. When a level 8 spell is copied in this way, the copier must succeed on a DC 18 Intelligence (Arcana). On a successful check, the spell is copied. Whether the check succeeds or fails, the Spell Scroll is destroyed.

\DndItemHeader{Spell Scroll (Level 9)}{Legendary}
A Spell Scroll bears the words of a single spell, written in a mystical cipher. If the spell is on your spell list, you can read the scroll and cast its spell without Material components. Otherwise, the scroll is unintelligible. Casting the spell by reading the scroll requires the spell's normal casting time. Once the spell is cast, the scroll crumbles to dust. If the casting is interrupted, the scroll isn't lost.

If the spell is on your spell list but of a higher level than you can normally cast, you make a DC 19 ability check using your spellcasting ability to determine whether you cast the spell. On a failed check, the spell disappears from the scroll with no other effect.

If the spell requires a saving throw or an attack roll, the spell save DC is 19, and the attack bonus is 11.

\subparagraph{Copying a Scroll into a Spellbook}
A Wizard spell on a Spell Scroll can be copied into a spellbook. When a level 9 spell is copied in this way, the copier must succeed on a DC 19 Intelligence (Arcana). On a successful check, the spell is copied. Whether the check succeeds or fails, the Spell Scroll is destroyed.

\DndItemHeader{Spellguard Shield}{Very Rare (requires attunement)}
While holding this Shield, you have Advantage on saving throws against spells and other magical effects, and spell attack rolls have Disadvantage against you.

\DndItemHeader{Sphere of Annihilation}{Wondrous item, Legendary}
This 2-foot-diameter black sphere is a hole in the multiverse, hovering in space and stabilized by a magical field surrounding it.

The sphere obliterates all matter it passes through and all matter that passes through it. Artifacts are the exception. Unless an Artifact is susceptible to damage from a Sphere of Annihilation, it passes through the sphere unscathed. Anything else that touches the sphere but isn't wholly engulfed and obliterated by it takes 8d10 Force damage.

\subparagraph{Controlling the Sphere}
A Sphere of Annihilation is stationary until someone takes control of it. If you are within 60 feet of a sphere, you can take a Magic action to make a DC 25 Intelligence (Arcana) check. On a successful check, you control the sphere until the start of your next turn, and if it was under another creature's control, that creature loses control of the sphere. On a failed check, the sphere moves 10 feet toward you in a straight line.

While in control of the sphere, you can take a Bonus Action to cause it to move in one direction of your choice, up to a number of feet equal to 5 times your Intelligence modifier (minimum 5 feet). Any creature whose space the sphere enters must succeed on a DC 19 Dexterity saving throw or be touched by it, taking 8d10 Force damage. A creature reduced to 0 Hit Points by this damage is obliterated, leaving its possessions behind but no other physical remains.

\subparagraph{Sphere Interactions}
If the sphere comes into contact with a planar portal (such as that created by the \textit{Gate} spell) or an extradimensional space (such as that within a Portable Hole), the DM determines randomly what happens using the following table.


\begin{DndTable}{cX}

1d100 & Result\\
01-50 & The sphere is destroyed.\\
51-85 & The sphere moves through the portal or into the extradimensional space.\\
86-00 & A spatial rift sends the sphere and each creature and object within 180 feet of the sphere to a random plane of existence.\\
\end{DndTable}

\DndItemHeader{Spinel}{}
A red, red brown, or deep green gemstone.

\DndItemHeader{Spirit Board}{Wondrous item, Very Rare}
This ornate wooden board has the letters of the Common alphabet printed on one side, alongside the words "Yes" and "No" and symbols representing "Weal" and "Woe." The board comes with a heartshaped, wooden planchette. This planchette must be resting on the lettered side of the board for the board's magic to function.

This board has 3 charges and regains 1 expended charge daily at dawn. While touching the planchette, you can take 1 minute to cast one of the spells on the table below. The table indicates how many charges you must expend to cast the spell. As you cast the spell, you call on the spirits of the dead to help guide the planchette across the board's surface, answering your questions by pointing to the letters or words on the board.


\begin{DndTable}{Xc}

Spell & Charge Cost\\
\textit{Augury} & 1\\
\textit{Commune} & 3\\
\end{DndTable}

\DndItemHeader{Staff of Adornment}{Common}
If you place a Tiny object weighing no more than 1 pound (such as a shard of crystal, an egg, or a stone) above the tip of this staff while holding it, the object floats an inch from the staff's tip and remains there until it is removed or until the staff is no longer in your possession. The staff can have up to three such objects floating over its tip at any given time. While holding the staff, you can make one or more of the objects slowly spin or turn in place.

\DndItemHeader{Staff of Birdcalls}{Common}
This wooden staff is decorated with bird carvings. It has 10 charges. While holding it, you can take a Magic action to expend 1 charge from the staff and cause it to create one of the following sounds, which can be heard out to 120 feet: a finch's chirp, a raven's caw, a duck's quack, a chicken's cluck, a goose's honk, a loon's call, a turkey's gobble, a seagull's cry, an owl's hoot, or an eagle's shriek.

\subparagraph{Regaining Charges}
The staff regains 1d6 + 4 expended charges daily at dawn. If you expend the last charge, roll 1d20. On a 1, the staff explodes in a harmless cloud of bird feathers and is lost forever.

\DndItemHeader{Staff of Charming}{Rare (requires attunement by a bard, cleric, druid, sorcerer, warlock, or wizard)}
This staff has 10 charges. While holding the staff, you can use any of its properties:


\begin{description}
\item[Cast Spell.]
You can expend 1 of the staff's charges to cast \textit{Charm Person}, \textit{Command}, or \textit{Comprehend Languages} from it using your spell save DC.

\item[Reflect Enchantment.]
If you succeed on a saving throw against an Enchantment spell that targets only you, you can take a Reaction to expend 1 charge from the staff and turn the spell back on its caster as if you had cast the spell.

\item[Resist Enchantment.]
If you fail a saving throw against an Enchantment spell that targets only you, you can turn your failed save into a successful one. You can't use this property of the staff again until the next dawn.

\end{description}

\subparagraph{Regaining Charges}
The staff regains 1d8 + 2 expended charges daily at dawn. If you expend the last charge, roll 1d20. On a 1, the staff crumbles to dust and is destroyed.

\DndItemHeader{Staff of Fire}{Very Rare (requires attunement by a druid, sorcerer, warlock, or wizard)}
You have Resistance to Fire damage while you hold this staff.

\subparagraph{Spells}
The staff has 10 charges. While holding the staff, you can cast one of the spells on the following table from it, using your spell save DC. The table indicates how many charges you must expend to cast the spell.


\begin{DndTable}{Xc}

Spell & Charge Cost\\
\textit{Burning Hands} & 1\\
\textit{Fireball} & 3\\
\textit{Wall of Fire} & 4\\
\end{DndTable}

\subparagraph{Regaining Charges}
The staff regains 1d6 + 4 expended charges daily at dawn. If you expend the last charge, roll 1d20. On a 1, the staff crumbles into cinders and is destroyed.

\DndItemHeader{Staff of Flowers}{Common}
This wooden staff has 10 charges. While holding it, you can take a Magic action to expend 1 charge from the staff and cause a flower to sprout from a patch of earth or soil within 5 feet of yourself, or from the staff itself. Unless you choose a specific kind of flower, the staff creates a mild-scented daisy. The flower is harmless and nonmagical, and it grows or withers as a normal flower would.

\subparagraph{Regaining Charges}
The staff regains 1d6 + 4 expended charges daily at dawn. If you expend the last charge, roll 1d20. On a 1, the staff turns into flower petals and is lost forever.

\DndItemHeader{Staff of Frost}{Very Rare (requires attunement by a druid, sorcerer, warlock, or wizard)}
You have Resistance to Cold damage while you hold this staff.

\subparagraph{Spells}
The staff has 10 charges. While holding the staff, you can cast one of the spells on the following table from it, using your spell save DC. The table indicates how many charges you must expend to cast the spell.


\begin{DndTable}{Xc}

Spell & Charge Cost\\
\textit{Cone of Cold} & 5\\
\textit{Fog Cloud} & 1\\
\textit{Ice Storm} & 4\\
\textit{Wall of Ice} & 4\\
\end{DndTable}

\subparagraph{Regaining Charges}
The staff regains 1d6 + 4 expended charges daily at dawn. If you expend the last charge, roll 1d20. On a 1, the staff turns to water and is destroyed.

\DndItemHeader{Staff of Healing}{Rare (requires attunement by a bard, cleric, or druid)}
This staff has 10 charges. While holding the staff, you can cast one of the spells on the following table from it, using your spellcasting ability modifier. The table indicates how many charges you must expend to cast the spell.


\begin{DndTable}{Xc}

Spell & Charge Cost\\
\textit{Cure Wounds} & 1 charge per spell level (maximum 4 for a level 4 spell)\\
\textit{Lesser Restoration} & 2\\
\textit{Mass Cure Wounds} & 5\\
\end{DndTable}

\subparagraph{Regaining Charges}
The staff regains 1d6 + 4 expended charges daily at dawn. If you expend the last charge, roll 1d20. On a 1, the staff vanishes in a flash of light, lost forever.

\DndItemHeader{Staff of Power}{Very Rare (requires attunement by a sorcerer, warlock, or wizard)}
This staff has 20 charges and can be wielded as a magic Quarterstaff that grants a +2 bonus to attack rolls and damage rolls made with it. While holding it, you gain a +2 bonus to Armor Class, saving throws, and spell attack rolls.

\subparagraph{Spells}
While holding the staff , you can cast one of the spells on the following table from it, using your spell save DC. The table indicates how many charges you must expend to cast the spell.


\begin{DndTable}{Xc}

Spell & Charge Cost\\
\textit{Cone of Cold} & 5\\
\textit{Fireball} (level 5 version) & 5\\
\textit{Globe of Invulnerability} & 6\\
\textit{Hold Monster} & 5\\
\textit{Levitate} & 2\\
\textit{Lightning Bolt} (level 5 version) & 5\\
\textit{Magic Missile} & 1\\
\textit{Ray of Enfeeblement} & 1\\
\textit{Wall of Force} & 5\\
\end{DndTable}

\subparagraph{Regaining Charges}
The staff regains 2d8 + 4 expended charges daily at dawn. If you expend the last charge, roll 1d20. On a 1, the staff retains its +2 bonus to attack rolls and damage rolls but loses all other properties. On a 20, the staff regains 1d8 + 2 charges.

\subparagraph{Retributive Strike}
You can take a Magic action to break the staff over your knee or against a solid surface. The staff is destroyed and releases its magic in an explosion that fills a 30-foot Emanation [Area of Effect] originating from itself. You have a 50 chance to instantly travel to a random plane of existence, avoiding the explosion. If you fail to avoid the effect, you take Force damage equal to 16 times the number of charges in the staff . Each other creature in the area makes a DC 17 Dexterity saving throw. On a failed save, a creature takes Force damage equal to 4 times the number of charges in the staff . On a successful save, a creature takes half as much damage.

\DndItemHeader{Staff of Striking}{Very Rare (requires attunement)}
This staff can be wielded as a magic Quarterstaff that grants a +3 bonus to attack rolls and damage rolls made with it.

The staff has 10 charges. When you hit with a melee attack using it, you can expend up to 3 charges. For each charge you expend, the target takes an extra 1d6 Force damage.

\subparagraph{Regaining Charges}
The staff regains 1d6 + 4 expended charges daily at dawn. If you expend the last charge, roll 1d20. On a 1, the staff becomes a nonmagical Quarterstaff.

\DndItemHeader{Staff of Swarming Insects}{Rare (requires attunement by a bard, cleric, druid, sorcerer, warlock, or wizard)}
This staff has 10 charges.

\subparagraph{Insect Cloud}
While holding the staff , you can take a Magic action and expend 1 charge to cause a swarm of harmless flying insects to fill a 30-foot Emanation [Area of Effect] originating from you. The insects remain for 10 minutes, making the area Heavily Obscured for creatures other than you. A strong wind (like that created by \textit{Gust of Wind}) disperses the swarm and ends the effect.

\subparagraph{Spells}
While holding the staff, you can cast one of the spells on the following table from it, using your spell save DC and spell attack modifier. The table indicates how many charges you must expend to cast the spell.


\begin{DndTable}{Xc}

Spell & Charge Cost\\
\textit{Giant Insect} & 4\\
\textit{Insect Plague} & 5\\
\end{DndTable}

\subparagraph{Regaining Charges}
The staff regains 1d6 + 4 expended charges daily at dawn. If you expend the last charge, roll 1d20. On a 1, a swarm of insects consumes and destroys the staff, then disperses.

\DndItemHeader{Staff of Thunder and Lightning}{Very Rare (requires attunement)}
This staff can be wielded as a magic Quarterstaff that grants a +2 bonus to attack rolls and damage rolls made with it. It also has the following additional properties. Once one of these properties is used, it can't be used again until the next dawn.

\subparagraph{Lightning}
When you hit with a melee attack using the staff, you can cause the target to take an extra 2d6 Lightning damage (no action required).

\subparagraph{Thunder}
When you hit with a melee attack using the staff, you can cause the staff to emit a crack of thunder audible out to 300 feet (no action required). The target you hit must succeed on a DC 17 Constitution saving throw or have the Stunned condition until the end of your next turn.

\subparagraph{Thunder and Lightning}
Immediately after you hit with a melee attack using the staff, you can take a Bonus Action to use the Lightning and Thunder properties (see above) at the same time. Doing so doesn't expend the daily use of those properties, only the use of this one.

\subparagraph{Lightning Strike}
You can take a Magic action to cause a bolt of lightning to leap from the staff's tip in a Line [Area of Effect] that is 5 feet wide and 120 feet long. Each creature in that Line [Area of Effect] makes a DC 17 Dexterity saving throw, taking 9d6 Lightning damage on a failed save or half as much damage on a successful one.

\subparagraph{Thunderclap}
You can take a Magic action to cause the staff to produce a thunderclap audible out to 600 feet. Every creature within a 60-foot Emanation [Area of Effect] originating from you makes a DC 17 Constitution saving throw. On a failed save, a creature takes 2d6 Thunder damage and has the Deafened condition for 1 minute. On a successful save, a creature takes half as much damage only.

\DndItemHeader{Staff of Withering}{Rare (requires attunement by a cleric, druid, or warlock)}
This staff has 3 charges and regains 1d3 expended charges daily at dawn.

The staff can be wielded as a magic Quarterstaff. On a hit, it deals damage as a normal Quarterstaff , and you can expend 1 charge to deal an extra 2d10 Necrotic damage to the target and force it to make a DC 15 Constitution saving throw. On a failed save, the target has Disadvantage for 1 hour on any ability check or saving throw that uses Strength or Constitution.

\DndItemHeader{Staff of the Adder}{Uncommon (requires attunement)}
As a Bonus Action, you can turn the head of this staff into that of an animate, venomous snake for 1 minute or revert the staff to its inanimate form.

When you take the Attack action, you can make one of the attack rolls using the animated snake head, which has a reach of 5 feet. Apply your Proficiency and Wisdom modifier to the attack roll. On a hit, the target takes 1d6 Piercing damage and 3d6 Poison damage.

The snake head can be attacked while it is animate. It has AC 15, HP 20, and Immunity to Poison and Psychic damage. If the head drops to 0 Hit Points, the staff is destroyed. As long as it's not destroyed, the staff regains all lost Hit Points when it reverts to its inanimate form.

\DndItemHeader{Staff of the Magi}{Legendary (requires attunement by a sorcerer, warlock, or wizard)}
This staff has 50 charges and can be wielded as a magic Quarterstaff that grants a +2 bonus to attack rolls and damage rolls made with it. While you hold it, you gain a +2 bonus to spell attack rolls.

\subparagraph{Spell Absorption}
While holding the staff , you have Advantage on saving throws against spells. In addition, you can take a Reaction when another creature casts a spell that targets only you. If you do, the staff absorbs the magic of the spell, canceling its effect and gaining a number of charges equal to the absorbed spell's level. However, if doing so brings the staff's total number of charges above 50, the staff explodes as if you activated its Retributive Strike (see below).

\subparagraph{Spells}
While holding the staff, you can cast one of the spells on the following table from it, using your spell save DC. The table indicates how many charges you must expend to cast the spell.


\begin{DndTable}{Xc}

Spell & Charge Cost\\
\textit{Arcane Lock} & 0\\
\textit{Conjure Elemental} & 7\\
\textit{Detect Magic} & 0\\
\textit{Dispel Magic} & 3\\
\textit{Enlarge/Reduce} & 0\\
\textit{Fireball} (level 7 version) & 7\\
\textit{Flaming Sphere} & 2\\
\textit{Ice Storm} & 4\\
\textit{Invisibility} & 2\\
\textit{Knock} & 2\\
\textit{Light} & 0\\
\textit{Lightning Bolt} (level 7 version) & 7\\
\textit{Mage Hand} & 0\\
\textit{Passwall} & 5\\
\textit{Plane Shift} & 7\\
\textit{Protection from Evil and Good} & 0\\
\textit{Telekinesis} & 5\\
\textit{Wall of Fire} & 4\\
\textit{Web} & 2\\
\end{DndTable}

\subparagraph{Regaining Charges}
The staff regains 4d6 + 2 expended charges daily at dawn. If you expend the last charge, roll 1d20. On a 20, the staff regains 1d12 + 1 charges.

\subparagraph{Retributive Strike}
You can take a Magic action to break the staff over your knee or against a solid surface. The staff is destroyed and releases its magic in an explosion that fills a 30-foot Emanation [Area of Effect] originating from itself. You have a 50 chance to instantly travel to a random plane of existence, avoiding the explosion. If you fail to avoid the effect, you take Force damage equal to 16 times the number of charges in the staff. Each other creature in the area makes a DC 17 Dexterity saving throw. On a failed save, a creature takes Force damage equal to 6 times the number of charges in the staff. On a successful save, a creature takes half as much damage.

\DndItemHeader{Staff of the Python}{Uncommon (requires attunement)}
As a Magic action, you can throw this staff so that it lands in an unoccupied space within 10 feet of you, causing the staff to become a \textbf{Giant Constrictor Snake} in that space. The snake is under your control and shares your Initiative count, taking its turn immediately after yours.

On your turn, you can mentally command the snake (no action required) if it is within 60 feet of you and you don't have the Incapacitated condition. You decide what action the snake takes and where it moves during its turn, or you can issue it a general command, such as to attack your enemies or guard a location. Absent commands from you, the snake defends itself.

As a Bonus Action, you can command the snake to revert to staff form in its current space, and you can't use the staff's property again for 1 hour. If the snake is reduced to 0 Hit Points, it dies and reverts to its staff form; the staff then shatters and is destroyed. If the snake reverts to staff form before losing all its Hit Points, it regains all of them.

\DndItemHeader{Staff of the Woodlands}{Rare (requires attunement by a druid)}
This staff has 6 charges and can be wielded as a magic Quarterstaff that grants a +2 bonus to attack rolls and damage rolls made with it. While holding it, you have a +2 bonus to spell attack rolls.

\subparagraph{Spells}
While holding the staff, you can cast one of the spells on the following table from it, using your spell save DC. The table indicates how many charges you must expend to cast the spell.


\begin{DndTable}{Xc}

Spell & Charge Cost\\
\textit{Animal Friendship} & 1\\
\textit{Awaken} & 5\\
\textit{Barkskin} & 2\\
\textit{Locate Animals or Plants} & 2\\
\textit{Pass without Trace} & 2\\
\textit{Speak with Animals} & 1\\
\textit{Speak with Plants} & 3\\
\textit{Wall of Thorns} & 6\\
\end{DndTable}

\subparagraph{Tree Form}
You can take a Magic action to plant one end of the staff in earth in an unoccupied space and expend 1 charge to transform the staff into a healthy tree. The tree is 60 feet tall and has a 5-foot-diameter trunk, and its branches at the top spread out in a 20-foot radius. The tree appears ordinary but radiates a faint aura of Transmutation magic that can be discerned with the \textit{Detect Magic} spell. While touching the tree and using a Magic action, you return the staff to its normal form. Any creature in the tree falls when the tree reverts to a staff.

\subparagraph{Regaining Charges}
The staff regains 1d6 expended charges daily at dawn. If you expend the last charge, roll 1d20. On a 1, the staff loses its properties and becomes a nonmagical Quarterstaff.

\DndItemHeader{Star rose quartz}{}
A rosy stone with white star-shaped center gemstone.

\DndItemHeader{Star Ruby}{}
A ruby with white star-shaped center gemstone.

\DndItemHeader{Star Sapphire}{}
A blue sapphire with white star-shaped center gemstone.

\DndItemHeader{Stone of Controlling Earth Elementals}{Wondrous item, Rare}
While touching this 5-pound stone to the ground, you can take a Magic action to summon an \textbf{Earth Elemental}. The elemental appears in an unoccupied space you choose within 30 feet of yourself, obeys your commands, and takes its turn immediately after you on your Initiative count. The elemental disappears after 1 hour, when it dies, or when you dismiss it as a Bonus Action. The stone can't be used this way again until the next dawn.

\DndItemHeader{Stone of Good Luck}{Wondrous item, Uncommon (requires attunement)}
While this polished agate is on your person, you gain a +1 bonus to ability checks and saving throws.

\DndItemHeader{Sun Blade}{Rare (requires attunement)}
This item appears to be a sword hilt.

\subparagraph{Blade of Radiance}
While grasping the hilt, you can take a Bonus Action to cause a blade of pure radiance to spring into existence or make the blade disappear. While the blade exists, this magic weapon functions as a Longsword with F. If you are proficient with Longswords or Shortswords, you are proficient with the Sun Blade.

You gain a +2 bonus to attack rolls and damage rolls made with this weapon, which deals Radiant damage instead of Slashing damage. When you hit an Undead with it, that target takes an extra 1d8 Radiant damage.

\subparagraph{Sunlight}
The sword's luminous blade emits Bright Light in a 15-foot radius and Dim Light for an additional 15 feet. The light is sunlight. While the blade persists, you can take a Magic action to expand or reduce its radius of Bright Light and Dim Light by 5 feet each, to a maximum of 30 feet each or a minimum of 10 feet each.

\DndItemHeader{Sword of Answering}{Legendary (requires attunement)}
You gain a +3 bonus to attack rolls and damage rolls made with this sword. In addition, while you hold the sword, you can take a Reaction to make one melee attack with it against any creature in your reach that deals damage to you. You have Advantage on the attack roll, and any damage dealt with this special attack ignores any Immunity or Resistance the target has to that damage.

\DndItemHeader{Sword of Kas}{Artifact (requires attunement)}
Kas was a powerful warrior who served \textit{Vecna} and whose loyalty was rewarded with this sword. As Kas's power grew, so did his hubris. The sword urged Kas to destroy Vecna and usurp his throne. Legend says Vecna's destruction came at Kas's hand, but Vecna also wrought his rebellious lieutenant's doom, leaving only Kas's sword behind.

\subparagraph{Bloodthirst}
The sword thirsts for blood. If the sword doesn't taste blood on its blade within 1 minute of being drawn from its scabbard, its wielder makes a DC 15 Charisma saving throw. On a successful save, the wielder takes 3d6 Psychic damage. On a failed save, the wielder is dominated by the sword, as if by the \textit{Dominate Monster} spell, and the sword demands blood. The spell effect ends when the sword's demand is met.

\subparagraph{Magic Weapon}
You gain a +3 bonus to attack rolls and damage rolls made with the sword, which scores a Critical Hit on a roll of 19 or 20 on the d20 and deals an extra 2d10 Slashing damage to Undead.

\subparagraph{Random Properties}
The sword has the following random properties:


\begin{itemize}
\item 1 Artifact Properties; Minor Beneficial Properties property
\item 1 Artifact Properties; Major Beneficial Properties property
\item 1 Artifact Properties; Minor Detrimental Properties property
\item 1 Artifact Properties; Major Detrimental Properties property
\end{itemize}

\subparagraph{Spells}
While the sword is on your person, you can cast the following spells (save DC 18) from it:


\begin{itemize}
\item \textit{Call Lightning}
\item \textit{Divine Word}
\item \textit{Finger of Death}
\end{itemize}

Once you use the sword to cast a spell, you can't cast that spell again from it until the next dawn.

\subparagraph{Spirit of Kas}
While the sword is on your person, you gain the following benefits:


\begin{description}
\item[Battle Hunger.]
You add 1d10 to your Initiative rolls.

\item[Blade of Defense.]
When you take an action to attack with the sword, you can transfer some or all of its attack bonus to your Armor Class instead. The adjusted bonuses remain in effect until the start of your next turn.

\item[Necrotic Resistance.]
You have Resistance to Necrotic damage.

\end{description}

\subparagraph{Sentience}
The Sword of Kas is a sentient Chaotic Evil weapon with an Intelligence of 15, a Wisdom of 13, and a Charisma of 16. It has hearing and Darkvision out to 120 feet.

The weapon communicates telepathically with its wielder and speaks Common.

\subparagraph{Personality}
The sword's purpose is to bring ruin to Vecna. Killing Vecna's worshipers, destroying the lich's works, and foiling his machinations all help to fulfill this goal.

The Sword of Kas also seeks to destroy anyone corrupted by the \textit{Eye and Hand of Vecna}.

\subparagraph{Destroying the Sword}
A creature attuned to both the \textit{Eye of Vecna} and the \textit{Hand of Vecna} can use the \textit{Wish} property of those combined Artifacts to unmake the Sword of Kas, provided the sword is within 30 feet of the spell's caster. Upon casting Wish, the creature makes a DC 18 Charisma saving throw. On a failed save, nothing happens, and the \textit{Wish} spell is wasted. On a successful save, the Sword of Kas is destroyed.

\DndItemHeader{Talisman of Pure Good}{Wondrous item, Legendary (requires attunement by a cleric or paladin)}
This talisman is a mighty symbol of goodness. A Fiend or an Undead that touches the talisman takes 8d6 Radiant damage and takes the damage again each time it ends its turn holding or carrying the talisman.

\subparagraph{Holy Symbol}
You can use the talisman as a Holy Symbol. You gain a +2 bonus to spell attack rolls while you wear or hold it.

\subparagraph{Pure Rebuke}
The talisman has 7 charges. While wearing or holding the talisman, you can take a Magic action to expend 1 charge and target one creature you can see on the ground within 120 feet of yourself. A flaming fissure opens under the target, and the target makes a DC 20 Dexterity saving throw. If the target is a Fiend or an Undead, it has Disadvantage on the save. On a failed save, the target falls into the fissure and is destroyed, leaving no remains. On a successful save, the target isn't cast into the fissure but takes 4d6 Psychic damage from the ordeal. In either case, the fissure then closes, leaving no trace of its existence. When you expend the last charge, the talisman disperses into motes of golden light and is destroyed.

\DndItemHeader{Talisman of Ultimate Evil}{Wondrous item, Legendary (requires attunement)}
This item symbolizes unrepentant evil. A creature that isn't a Fiend or an Undead that touches the talisman takes 8d6 Necrotic damage and takes the damage again each time it ends its turn holding or carrying the talisman.

\subparagraph{Holy Symbol}
You can use the talisman as a Holy Symbol. You gain a +2 bonus to spell attack rolls while you wear or hold it.

\subparagraph{Ultimate End}
The talisman has 6 charges. While wearing or holding the talisman, you can take a Magic action to expend 1 charge and target one creature you can see on the ground within 120 feet of yourself. A flaming fissure opens under the target, and the target makes a DC 20 Dexterity saving throw. If the target is a Celestial, it has Disadvantage on the save. On a failed save, the target falls into the fissure and is destroyed, leaving no remains. On a successful save, the target isn't cast into the fissure but takes 4d6 Psychic damage from the ordeal. In either case, the fissure then closes, leaving no trace of its existence. When you expend the last charge, the talisman dissolves into foul-smelling slime and is destroyed.

\DndItemHeader{Talisman of the Sphere}{Wondrous item, Legendary (requires attunement)}
While holding or wearing this talisman, you have Advantage on any Intelligence (Arcana) check you make to control a \textit{Sphere of Annihilation}. In addition, when you start your turn in control of a \textit{Sphere of Annihilation}, you can take a Magic action to move it 10 feet plus a number of additional feet equal to 10 times your Intelligence modifier. This movement doesn't have to be in a straight line.

\DndItemHeader{Talking Doll}{Wondrous item, Common (requires attunement)}
While this doll is within 5 feet of you, you can spend a Short Rest telling it to say up to six phrases, none of which can be more than six words long, and set a condition under which the doll speaks each phrase. You can also replace old phrases with new ones. Whatever the condition, it must occur within 5 feet of the doll to make it speak. For example, whenever someone picks up the doll, it might say, "I want a piece of candy." The doll's phrases are lost when your Attunement to the doll ends.

\DndItemHeader{Tankard of Sobriety}{Wondrous item, Common}
This tankard has a stern face sculpted into one side. You can drink ale, wine, or any other nonmagical alcoholic beverage poured into it without becoming inebriated. The tankard has no effect on magical liquids or harmful substances such as poison.

\DndItemHeader{Tentacle Rod}{Rare (requires attunement)}
This rod ends in three rubbery tentacles. While holding the rod, you can take a Magic action to direct the tentacles to stretch outward, each one attacking a creature you can see within 15 feet of yourself. For each tentacle, make a melee attack roll with a +9 bonus. A tentacle deals 1d6 Psychic damage on a hit. If you hit the same target with all three tentacles, the target must succeed on a DC 15 Dexterity saving throw or have the Restrained condition until you have the Incapacitated condition, until you take a Bonus Action to release the target, or until the target is no longer within 15 feet of you. While Restrained in this way, the target takes 3d6 Psychic damage at the start of each of its turns. At the end of each of its turns, the target repeats the save, ending the effect on itself on a success.

\DndItemHeader{Thieves' Tools}{}

\begin{description}
\item[Ability:]
Dexterity

\item[Utilize:]
Pick a lock (DC 15), or disarm a trap (DC 15)

\end{description}

\DndItemHeader{Thunderous Greatclub}{Very Rare (requires attunement)}
While you are attuned to this magic weapon, your Strength is 20 unless your Strength is already equal to or greater than that score. The weapon deals an extra 1d8 Thunder damage to any creature it hits and an extra 3d8 Thunder damage to objects it hits that aren't being worn or carried.

The weapon has the following additional properties.

\subparagraph{Clap of Thunder}
As a Magic action, you can strike the weapon against a hard surface to create a loud clap of thunder audible out to 300 feet. You also create a 30-foot Cone [Area of Effect] of thunderous energy. Each creature in the Cone [Area of Effect] must succeed on a DC 15 Strength saving throw or have the Prone condition. Nonmagical objects in the Cone [Area of Effect] that aren't being worn or carried take 3d8 Thunder damage.

\subparagraph{Earthquake}
As a Magic action, you can strike the weapon against the ground to create an intense seismic disturbance in a 50-foot-radius circle centered on the point of impact. Structures in contact with the ground in that area take 50 Bludgeoning damage, and each creature on the ground in that area must succeed on a DC 20 Dexterity saving throw or have the Prone condition. If that creature is also Concentration, it must succeed on a DC 20 Constitution saving throw or its Concentration is broken. In addition, you can cause a 30-foot-deep, 10-foot-wide fissure to open up on the ground anywhere in the area. Any creature on a spot where the fissure opens must make a DC 20 Dexterity saving throw, falling into the fissure on a failed save or moving with the fissure's edge on a successful one. Any structure on a spot where the fissure opens collapses into the fissure. Once you use this property, it can't be used again until the next dawn.

\DndItemHeader{Tiger Eye}{}
A brown with golden center gemstone.

\DndItemHeader{Tinker's Tools}{}

\begin{description}
\item[Ability:]
Dexterity

\item[Utilize:]
Assemble a Tiny item composed of scrap, which falls apart in 1 minute (DC 20)

\item[Craft:]
\textit{Musket}, \textit{Pistol}, \textit{Bell}, \textit{Bullseye Lantern}, \textit{Flask}, \textit{Hooded Lantern}, \textit{Hunting Trap}, \textit{Lock}, \textit{Manacles}, \textit{Mirror}, \textit{Shovel}, \textit{Signal Whistle}, \textit{Tinderbox}

\end{description}

\DndItemHeader{Tome of Clear Thought}{Wondrous item, Very Rare}
This book contains memory and logic exercises, and its words are charged with magic. If you spend 48 hours over a period of 6 days or fewer studying the book's contents and practicing its guidelines, your Intelligence score increases by 2, as does your maximum for that score. The manual then loses its magic, but regains it in a century.

\DndItemHeader{Tome of Leadership and Influence}{Wondrous item, Very Rare}
This book contains guidelines for influencing and charming others, and its words are charged with magic. If you spend 48 hours over a period of 6 days or fewer studying the book's contents and practicing its guidelines, your Charisma score increases by 2, as does your maximum for that score. The manual then loses its magic, but regains it in a century.

\DndItemHeader{Tome of Understanding}{Wondrous item, Very Rare}
This book contains intuition and insight exercises, and its words are charged with magic. If you spend 48 hours over a period of 6 days or fewer studying the book's contents and practicing its guidelines, your Wisdom score increases by 2, as does your maximum for that score. The manual then loses its magic, but regains it in a century.

\DndItemHeader{Tome of the Stilled Tongue}{Wondrous item, Legendary (requires attunement by a wizard)}
This book has a desiccated tongue pinned to its front cover. Five of these tomes exist, and it's unknown which one is the original. The tongue on the first Tome of the Stilled Tongue belonged to a treacherous former servant of the lich \textit{Vecna}. The tongues pinned to the covers of the four copies came from other spellcasters who crossed Vecna. The first few pages of each tome are filled with indecipherable scrawls. The remaining pages are blank.

While attuned to this item, you can use it as a Spellbook and an Arcane Focus. In addition, while holding the tome, you can take a Bonus Action to cast a spell you have written in this tome, without expending a spell slot or using any Verbal or Somatic components. Once used, this property of the tome can't be used again until the next dawn.

Only you can remove the tongue from the book's cover. If you do so, all spells written in the book are permanently erased.

Vecna watches anyone using this tome and can write cryptic messages in it. These messages typically fade away after they are read.

\DndItemHeader{Topaz}{}
A golden yellow gemstone.

\DndItemHeader{Torpor}{}
A creature subjected to Torpor poison must succeed on a DC 15 Constitution saving throw or have the Poisoned condition for 4d6 hours. The creature's Speed is halved while the creature is Poisoned in this way.

\DndItemHeader{Tourmaline}{}
A pale green, blue, brown, or red gemstone.

\DndItemHeader{Trident of Fish Command}{Uncommon (requires attunement)}
This magic weapon has 3 charges, and it regains 1d3 expended charges daily at dawn. While you carry it, you can expend 1 charge to cast \textit{Dominate Beast} (save DC 15) from it on a Beast that has a Swim Speed.

\DndItemHeader{Truth Serum}{}
A creature subjected to Truth Serum must succeed on a DC 11 Constitution saving throw or have the Poisoned condition for 1 hour. The Poisoned creature can't knowingly communicate a lie.

\DndItemHeader{Turquoise}{}
A light blue green gemstone.

\DndItemHeader{Universal Solvent}{Wondrous item, Legendary}
This tube holds milky liquid with a strong alcohol smell. When found, a tube contains 1d6 + 1 ounces.

You can take a Utilize action to pour 1 or more ounces of solvent from the tube onto a surface within reach. Each ounce instantly dissolves up to 1 square foot of adhesive it touches, including Sovereign Glue.

\DndItemHeader{Veteran's Cane}{Wondrous item, Common}
As a Bonus Action, you can transform this walking cane into an ordinary Longsword or change the Longsword back into a walking cane. In either case, you must be holding the item.

\DndItemHeader{Wand of Binding}{Rare}
This wand has 7 charges.

\subparagraph{Spells}
While holding the wand, you can cast one of the spells (save DC 17) on the following table from it. The table indicates how many charges you must expend to cast the spell.


\begin{DndTable}{Xc}

Spell & Charge Cost\\
\textit{Hold Monster} & 5\\
\textit{Hold Person} & 2\\
\end{DndTable}

\subparagraph{Regaining Charges}
The wand regains 1d6 + 1 expended charges daily at dawn. If you expend the wand's last charge, roll 1d20. On a 1, the wand crumbles into ashes and is destroyed.

\DndItemHeader{Wand of Conducting}{Common}
This wand has 3 charges. While holding it, you can take a Magic action to expend 1 charge and create orchestral music by waving it around. The music can be heard out to 120 feet and ends when you stop waving the wand.

\subparagraph{Regaining Charges}
The wand regains all expended charges daily at dawn. If you expend the wand's last charge, roll 1d20. On a 1, a sad tuba sound plays as the wand crumbles into dust and is destroyed.

\DndItemHeader{Wand of Enemy Detection}{Rare}
This wand has 7 charges. While holding it, you can take a Magic action to expend 1 charge. For 1 minute, you know the direction of the nearest creature Hostile [Attitude] to you within 60 feet, but not its distance from you. The wand can sense the presence of Hostile [Attitude] creatures that are Invisible, ethereal, disguised, or hidden, as well as those in plain sight. The effect ends if you stop holding the wand.

\subparagraph{Regaining Charges}
The wand regains 1d6 + 1 expended charges daily at dawn. If you expend the wand's last charge, roll 1d20. On a 1, the wand crumbles into ashes and is destroyed.

\DndItemHeader{Wand of Fear}{Rare}
This wand has 7 charges.

\subparagraph{Spells}
While holding the wand, you can cast one of the spells (save DC 15) on the following table from it. The table indicates how many charges you must expend to cast the spell.


\begin{DndTable}{Xc}

Spell & Charge Cost\\
\textit{Command} ("flee" or "grovel" only) & 1\\
\textit{Fear} (60-foot Cone [Area of Effect]) & 3\\
\end{DndTable}

\subparagraph{Regaining Charges}
The wand regains 1d6 + 1 expended charges daily at dawn. If you expend the wand's last charge, roll 1d20. On a 1, the wand crumbles into ashes and is destroyed.

\DndItemHeader{Wand of Fireballs}{Rare (requires attunement by a spellcaster)}
This wand has 7 charges. While holding it, you can expend no more than 3 charges to cast \textit{Fireball} (save DC 15) from it. For 1 charge, you cast the level 3 version of the spell. You can increase the spell's level by 1 for each additional charge you expend.

\subparagraph{Regaining Charges}
The wand regains 1d6 + 1 expended charges daily at dawn. If you expend the wand's last charge, roll 1d20. On a 1, the wand crumbles into ashes and is destroyed.

\DndItemHeader{Wand of Lightning Bolts}{Rare (requires attunement by a spellcaster)}
This wand has 7 charges. While holding it, you can expend no more than 3 charges to cast \textit{Lightning Bolt} (save DC 15) from it. For 1 charge, you cast the level 3 version of the spell. You can increase the spell's level by 1 for each additional charge you expend.

\subparagraph{Regaining Charges}
The wand regains 1d6 + 1 expended charges daily at dawn. If you expend the wand's last charge, roll 1d20. On a 1, the wand crumbles into ashes and is destroyed.

\DndItemHeader{Wand of Magic Detection}{Uncommon}
This wand has 3 charges. While holding it, you can expend 1 charge to cast \textit{Detect Magic} from it. The wand regains 1d3 expended charges daily at dawn.

\DndItemHeader{Wand of Magic Missiles}{Uncommon}
This wand has 7 charges. While holding it, you can expend no more than 3 charges to cast \textit{Magic Missile} from it. For 1 charge, you cast the level 1 version of the spell. You can increase the spell's level by 1 for each additional charge you expend.

\subparagraph{Regaining Charges}
The wand regains 1d6 + 1 expended charges daily at dawn. If you expend the wand's last charge, roll 1d20. On a 1, the wand crumbles into ashes and is destroyed.

\DndItemHeader{Wand of Orcus}{Artifact (requires attunement)}
Crafted and wielded by \textbf{Orcus}, this ghastly wand slips from the demon lord's grasp from time to time. When it does, it magically appears wherever the demon lord senses an opportunity to achieve some fell goal.

The wand is topped with a skull that once belonged to a human hero slain by Orcus. The wand can magically change in size to better conform to the grip of its user. All Holy Water within 10 feet of the wand is destroyed.

Any creature besides Orcus that tries to attune to the wand makes a DC 17 Constitution saving throw. On a successful save, the creature takes 10d6 Necrotic damage. On a failed save, the creature dies and, if it is a Humanoid, turns into a \textbf{Zombie}.

\subparagraph{Magic Weapon}
You can wield the wand as a magic Mace that grants a +3 bonus to attack rolls and damage rolls made with it. The wand deals an extra 2d12 Necrotic damage on a hit.

\subparagraph{Random Properties}
The Wand of Orcus has the following random properties:


\begin{itemize}
\item 2 Artifact Properties; Minor Beneficial Properties properties
\item 1 Artifact Properties; Major Beneficial Properties property
\item 2 Artifact Properties; Minor Detrimental Properties properties
\item 1 Artifact Properties; Major Detrimental Properties property
\end{itemize}

The detrimental properties of the Wand of Orcus are suppressed while the wand is attuned to Orcus.

\subparagraph{Protection}
You gain a +3 bonus to Armor Class while holding the wand.

\subparagraph{Spells}
The wand has 7 charges and regains 1d4 + 3 expended charges daily at dawn. While holding the wand, you can cast one of the spells on the following table from it (save DC 18). The table indicates how many charges you must expend to cast the spell.


\begin{DndTable}{Xc}

Spell & Charge Cost\\
\textit{Animate Dead} & 1\\
\textit{Blight} & 2\\
\textit{Circle of Death} & 3\\
\textit{Finger of Death} & 3\\
\textit{Power Word Kill} & 4\\
\textit{Speak with Dead} & 1\\
\end{DndTable}

While attuned to the wand, Orcus or a follower blessed by him can cast each of the wand's spells using 2 fewer charges (minimum of 0).

\subparagraph{Call Undead}
While holding the wand, you can take a Magic action to conjure 15 Skeletons and 15 Zombies. These Undead magically rise up from the ground or otherwise form in unoccupied spaces within 300 feet of you and obey your commands until they are destroyed or until the next dawn, when they collapse into inanimate piles of bones and rotting corpses. Once you use this property, you can't use it again until the next dawn.

While holding the wand, Orcus can summon any kind of Undead, not just skeletons and zombies. These Undead don't perish at dawn the following day, remaining until Orcus dismisses them.

\subparagraph{Sentience}
The Wand of Orcus is a sentient Chaotic Evil item with an Intelligence of 16, a Wisdom of 12, and a Charisma of 16. It has hearing and Darkvision out to 120 feet.

The wand communicates telepathically with its wielder and speaks Abyssal and Common.

\subparagraph{Personality}
The wand's purpose is to help satisfy Orcus's desire to slay everything in the multiverse. The wand is cruel, nihilistic, and bereft of humor.

To further Orcus's goals, the wand feigns devotion to its current user and makes grandiose promises that it has no intention of fulfilling, such as vowing to help its user overthrow Orcus.

\subparagraph{Destroying the Wand}
Destroying the Wand of Orcus requires that it be taken to the Positive Energy Plane by the ancient hero whose skull surmounts it. For this to happen, the long-lost hero must first be restored to life—no easy task, given the fact that Orcus has imprisoned the hero's soul and keeps it hidden and well guarded.

Bathing the wand in positive energy (such as that which permeates the Positive Plane) causes it to crack and explode, but unless the above conditions are met, the wand instantly re-forms on Orcus's layer of the Abyss.

\DndItemHeader{Wand of Paralysis}{Rare (requires attunement by a spellcaster)}
This wand has 7 charges. While holding it, you can take a Magic action to expend 1 charge to cause a thin blue ray to streak from the tip toward a creature you can see within 60 feet of yourself. The target must succeed on a DC 15 Constitution saving throw or have the Paralyzed condition for 1 minute. At the end of each of the target's turns, it repeats the save, ending the effect on itself on a success.

\subparagraph{Regaining Charges}
The wand regains 1d6 + 1 expended charges daily at dawn. If you expend the wand's last charge, roll 1d20. On a 1, the wand crumbles into ashes and is destroyed.

\DndItemHeader{Wand of Polymorph}{Very Rare (requires attunement by a spellcaster)}
This wand has 7 charges. While holding it, you can expend 1 charge to cast \textit{Polymorph} (save DC 15) from it.

\subparagraph{Regaining Charges}
The wand regains 1d6 + 1 expended charges daily at dawn. If you expend the wand's last charge, roll 1d20. On a 1, the wand crumbles into ashes and is destroyed.

\DndItemHeader{Wand of Pyrotechnics}{Common}
This wand has 7 charges. While holding it, you can take a Magic action to expend 1 charge and create a harmless burst of multicolored light at a point you can see up to 120 feet away. The burst of light is accompanied by a crackling noise that can be heard up to 300 feet away. The light is as bright as a torch flame but lasts only a second.

\subparagraph{Regaining Charges}
The wand regains 1d6 + 1 expended charges daily at dawn. If you expend the wand's last charge, roll 1d20. On a 1, the wand erupts in a harmless pyrotechnic display and is destroyed.

\DndItemHeader{Wand of Secrets}{Uncommon}
This wand has 3 charges and regains 1d3 expended charges daily at dawn. While holding it, you can take a Magic action to expend 1 charge, and if a secret door or trap is within 60 feet of you, the wand pulses and points at the one nearest to you.

\DndItemHeader{Wand of Web}{Uncommon (requires attunement by a spellcaster)}
This wand has 7 charges. While holding it, you can expend 1 charge to cast \textit{Web} (save DC 13) from it.

\subparagraph{Regaining Charges}
The wand regains 1d6 + 1 expended charges daily at dawn. If you expend the wand's last charge, roll 1d20. On a 1, the wand crumbles into ashes and is destroyed.

\DndItemHeader{Wand of Wonder}{Rare (requires attunement)}
This wand has 7 charges. While holding it, you can take a Magic action to expend 1 charge while choosing a point within 120 feet of yourself. That location becomes the point of origin of a spell or other magical effect determined by rolling on the Wand of Wonder Effects table. Spells cast from the wand have a save DC of 15. If a spell's maximum range is normally less than 120 feet, it becomes 120 feet when cast from the wand. If an effect has multiple possible subjects, the DM determines randomly which among them are affected.

\subparagraph{Regaining Charges}
The wand regains 1d6 + 1 expended charges daily at dawn. If you expend the wand's last charge, roll 1d20. On a 1, the wand crumbles into dust and is destroyed.

\begin{DndTable}[header=Wand of Wonder Effects]{cX}

1d100 & Effect\\
01-20 & You cast a spell originating from the chosen point. Roll 1d10 to determine the spell: on a \textbf{1-2}, \textit{Darkness}; on a \textbf{3-4}, \textit{Faerie Fire}; on a \textbf{5-6}, \textit{Fireball}; on a \textbf{7-8}, \textit{Slow}; on a \textbf{9-10}, \textit{Stinking Cloud}.\\
21-25 & Nothing happens at the chosen point of origin. Instead, you have the Stunned condition until the start of your next turn, believing something awesome just happened.\\
26-30 & You cast \textit{Gust of Wind}. The Line [Area of Effect] created by the spell extends from you to the chosen point of origin.\\
31-35 & Nothing happens at the chosen point of origin. Instead, you take 1d6 Psychic damage.\\
36-40 & Heavy rain falls for 1 minute in a 120-foot-high, 60-foot-radius Cylinder [Area of Effect] centered on the chosen point of origin. During that time, the area of effect is Lightly Obscured.\\
41-45 & A cloud of 600 oversized butterflies fills a 60-foot-high, 30-foot-radius Cylinder [Area of Effect] centered on the chosen point of origin. The butterflies remain for 10 minutes, during which time the area of effect is Heavily Obscured.\\
46-50 & You cast \textit{Lightning Bolt}. The Line [Area of Effect] created by the spell extends from you to the chosen point of origin.\\
51-55 & The creature closest to the chosen point of origin is enlarged as if you had cast \textit{Enlarge/Reduce} on it. If the target isn't you and can't be affected by that spell, you become the target instead.\\
56-60 & A magically formed creature appears in an unoccupied space as close to the chosen point of origin as possible. The creature isn't under your control, acts as it normally would, and disappears after 1 hour or when it drops to 0 Hit Points. Roll 1d4 to determine which creature appears. On a \textbf{1}, a \textbf{Rhinoceros} appears; on a \textbf{2}, an \textbf{Elephant} appears; and on a \textbf{3-4}, a \textbf{Rat} appears.\\
61-64 & Grass covers a 60-foot-radius circle of ground, with the center of that circle as close to the chosen point of origin as possible. Grass that's already there grows to ten times its normal size and remains overgrown for 1 minute.\\
65-68 & An object of the DM's choice disappears into the Ethereal Plane. The object must be neither worn nor carried, within 120 feet of the chosen point of origin, and no larger than 10 feet in any dimension. If there are no such objects in range, nothing happens.\\
69-72 & Nothing happens at the chosen point of origin. Instead, you shrink as if you had cast \textit{Enlarge/Reduce} on yourself and remain in that state for 1 minute.\\
73-77 & Leaves grow from the creature nearest to the chosen point of origin. Unless they are picked off, the leaves turn brown and fall off after 24 hours.\\
78-82 & Nothing happens at the chosen point of origin. Instead, a burst of colorful, shimmering light extends from you in a 30-foot Emanation [Area of Effect]. Each creature in the area must succeed on a DC 15 Constitution saving throw or have the Blinded condition for 1 minute. A creature repeats the save at the end of each of its turns, ending the effect on itself on a success.\\
83-87 & Nothing happens at the chosen point of origin. Instead, you cast \textit{Invisibility} on yourself.\\
88-92 & Nothing happens at the chosen point of origin. Instead, a stream of 1d4 × 10 gems, each worth 1 GP, shoots from the wand's tip in a Line [Area of Effect] 30 feet long and 5 feet wide toward the chosen point of origin. Each gem deals 1 Bludgeoning damage, and the total damage of the gems is divided equally among all creatures in the Line [Area of Effect].\\
93-97 & You cast \textit{Polymorph}, targeting the creature closest to the chosen point of origin. Roll 1d4 to determine the target's new form. On a \textbf{1}, the new form is a \textbf{Black Bear}; on a \textbf{2}, the new form is a \textbf{Giant Wasp}; on a \textbf{3-4}, the new form is a \textbf{Frog}.\\
98-00 & The creature closest to the chosen point of origin makes a DC 15 Constitution saving throw. On a failed save, the creature has the Restrained condition and begins to turn to stone. While Restrained in this way, the creature repeats the save at the end of its next turn. On a successful save, the effect ends. On a failed save, the creature has the Petrified condition instead of the Restrained condition. The petrification lasts until the creature is freed by the \textit{Greater Restoration} spell or similar magic.\\
\end{DndTable}

\DndItemHeader{Wave}{Artifact (requires attunement)}
Held in the dungeon of White Plume Mountain, Wave is engraved with images of waves, shells, and sea creatures.

You gain a +3 bonus to attack rolls and damage rolls made with this magic weapon. When you roll a 20 on the d20 for an attack roll with this weapon, the target takes an extra 21 Necrotic damage.

While holding Wave, you gain the following benefits:


\begin{description}
\item[Combat Ready.]
You have Advantage on Initiative rolls.

\item[Underwater Adaptation.]
A bubble of air forms around your head while you are underwater, allowing you to breathe normally in that environment.

\end{description}

\subparagraph{Aquatic Command}
Wave has 3 charges and regains 1d3 expended charges daily at dawn. While you carry it, you can expend 1 charge to cast \textit{Dominate Beast} (save DC 20) from it on a Beast that has a Swim Speed.

\subparagraph{Globe of Invulnerability}
While holding Wave, you can cast the level 9 version of \textit{Globe of Invulnerability} from it. Once used, this property can't be used again until the next dawn.

\subparagraph{Sentience}
Wave is a sentient weapon of Neutral alignment, with an Intelligence of 14, a Wisdom of 10, and a Charisma of 18. It has hearing and Darkvision out to 120 feet.

The weapon communicates telepathically with its wielder and speaks Aquan.

\subparagraph{Personality}
Wave zealously encourages mortals to worship sea gods and has a habit of humming sea chanteys. Conflict arises if the wielder fails to further the weapon's objectives in the world.

\subparagraph{Destroying Wave}
Wave can be destroyed only on the island of Thunderforge, where it was forged. The weapon must be melted down by a \textbf{storm giant} or someone imbued with a storm giant's strength. Destroying Wave angers a god of the sea, who sends powerful agents to attack the island and punish the destroyers.

\DndItemHeader{Weaver's Tools}{}

\begin{description}
\item[Ability:]
Dexterity

\item[Utilize:]
Mend a tear in clothing (DC 10), or sew a Tiny design (DC 10)

\item[Craft:]
\textit{Padded Armor}, \textit{Basket}, \textit{Bedroll}, \textit{Blanket}, \textit{Fine Clothes}, \textit{Net}, \textit{Robe}, \textit{Rope}, \textit{Sack}, \textit{String}, \textit{Tent}, \textit{Traveler's Clothes}

\end{description}

\DndItemHeader{Well of Many Worlds}{Wondrous item, Legendary}
This fine black cloth, soft as silk, is folded up to the dimensions of a handkerchief. It unfolds into a circular sheet 6 feet in diameter.

You can take a Magic action to unfold the Well of Many Worlds and place it on a solid surface, whereupon it forms a two-way, 6-foot-diameter, circular portal to another world or plane of existence. Each time the item opens a portal, the DM decides where it leads. The portal remains open until a creature within 5 feet of it takes a Magic action to close it by taking hold of the edges of the cloth and folding it up.

Once the Well of Many Worlds has opened a portal, it can't do so again for 1d8 hours.

\DndItemHeader{Whelm}{Artifact (requires attunement by a dwarf or a creature attuned to a belt of dwarvenkind)}
Whelm is a powerful weapon forged by dwarves and lost in the dungeon of White Plume Mountain.

You gain a +3 bonus to attack rolls and damage rolls made with this magic weapon.

\subparagraph{Hurl}
Whelm has T with a normal range of 60 feet and a long range of 180 feet. When you hit with a ranged attack roll using Whelm, the target takes an extra 1d8 Force damage, or an extra 4d8 Force damage if the target is a Construct, an Elemental, or a Giant. Immediately after hitting or missing, the weapon flies back to your hand.

\subparagraph{Shock Wave}
You can take a Magic action to strike the ground with Whelm and send a shock wave out from the point of impact. Each creature of your choice on the ground within 60 feet of that point must succeed on a DC 20 Constitution saving throw or have the Stunned condition for 1 minute. A creature repeats the save at the end of each of its turns, ending the effect on itself on a success. Once used, this property can't be used again until the next dawn.

\subparagraph{Supernatural Awareness}
While you are holding the weapon, it alerts you to the location of any secret or concealed doors within 30 feet of you. In addition, you can cast \textit{Detect Evil and Good} or \textit{Locate Object} from the weapon. Once you cast either spell, you can't cast it from the weapon again until the next dawn.

\subparagraph{Sentience}
Whelm is a sentient, Lawful Neutral weapon with an Intelligence of 15, a Wisdom of 12, and a Charisma of 15. It has hearing and Darkvision out to 120 feet.

The weapon communicates telepathically with its wielder and speaks Dwarvish, Giant, and Goblin.

\subparagraph{Personality}
Whelm has ties to the dwarf clan that created it, called the Dankil or the Mightyhammer clan. It longs to be returned to that clan. Whelm's purpose is to protect dwarves. Conflict arises if the wielder doesn't share this goal.

\subparagraph{Destroying Whelm}
Whelm can be dissolved in the acidic bile of a recently slain \textbf{ancient black dragon}. It can also be melted down in the forges of the Mightyhammer dwarf clan, but only by the rightful leader of that clan.

\DndItemHeader{Wind Fan}{Wondrous item, Uncommon}
While holding this fan, you can cast \textit{Gust of Wind} (save DC 13) from it. Each subsequent time the fan is used before the next dawn, it has a cumulative 20 chance of not working; if the fan fails to work, it tears into useless, nonmagical tatters.

\DndItemHeader{Winged Boots}{Wondrous item, Uncommon (requires attunement)}
These boots have 4 charges and regain 1d4 expended charges daily at dawn. While wearing the boots, you can take a Magic action to expend 1 charge, gaining a Fly Speed of 30 feet for 1 hour. If you are flying when the duration expires, you descend at a rate of 30 feet per round until you land.

\DndItemHeader{Wings of Flying}{Wondrous item, Rare (requires attunement)}
While wearing this cloak, you can take a Magic action to turn the cloak into a pair of wings on your back. The wings lasts for 1 hour or until you end the effect early as a Magic action. The wings give you a Fly Speed of 60 feet. If you are aloft when the wings disappear, you fall. When the wings disappear, you can't use them again for 1d12 hours.

\DndItemHeader{Woodcarver's Tools}{}

\begin{description}
\item[Ability:]
Dexterity

\item[Utilize:]
Carve a pattern in wood (DC 10)

\item[Craft:]
\textit{Club}, \textit{Greatclub}, \textit{Quarterstaff}, Ranged weapons (except \textit{Pistol}, \textit{Musket}, and  \textit{Sling}), \textit{Arcane Focus}, Arrows (20), Bolts (20), \textit{Druidic Focus}, \textit{Ink Pen}, Needles (50)

\end{description}

\DndItemHeader{Wyvern Poison}{}
A creature subjected to Wyvern Poison makes a DC 14 Constitution saving throw, taking 24 (7d6) Poison damage on a failed save or half as much damage on a successful one.

\DndItemHeader{Yellow Sapphire}{}
A fiery yellow or yellow green gemstone.

\DndItemHeader{Zircon}{}
A pale blue green gemstone.

\end{document}
